% 本文件由 LaTeX 排版 Agent 根据有效 translation.json 自动生成。
% author=Jerome; work_id=3001; title=书信集

\chapter{书信集}

% structure: p0001 source=h1 target=omitted reason=page-title-already-rendered-by-parent

书信一　致依诺森（叙利亚，三七四年）\newline{}书信二　致德奥多西及其他独修者（安提约基雅，三七四年）\newline{}书信三　致隐修士鲁菲诺（安提约基雅，三七四年）\newline{}书信四　致弗洛伦提（安提约基雅，三七四年）\newline{}书信五　致弗洛伦提（叙利亚旷野，三七四年）\newline{}书信六　致安提约基雅执事犹利安（三七四年）\newline{}书信七　致克罗马提、若维诺与欧瑟伯（三七四年）\newline{}书信八　致阿奎莱亚副执事尼塞亚（三七四年）\newline{}书信九　致阿奎莱亚隐修士克黎索贡（三七四年）\newline{}书信十　致科杜科迪亚老人保禄（三七四年）\newline{}书信十一　致厄莫纳贞女们（三七四年）\newline{}书信十二　致隐修士安当（三七四年）\newline{}书信十三　致其姨母卡斯托利纳（三七四年）\newline{}书信十四　致隐修士赫利奥多若（三七三或三七四年）\newline{}书信十五　致教宗达玛稣（三七六或三七七年）\newline{}书信十六　致教宗达玛稣（旷野，三七七或三七八年）\newline{}书信十七　致司铎马尔谷（旷野，三七八或三七九年）\newline{}书信十八　致教宗达玛稣（君士坦丁堡，三八一年）\newline{}书信十九　教宗达玛稣来函（三八三年）\newline{}书信二十　致教宗达玛稣（罗马，三八三年）\newline{}书信二十一　致达玛稣（？）\newline{}书信二十二　致欧斯托基（罗马，三八四年）\newline{}书信二十三　致玛尔策拉（罗马，三八四年）\newline{}书信二十四　致玛尔策拉（罗马，三八四年）\newline{}书信二十五　致玛尔策拉（罗马，三八四年）\newline{}书信二十六　致玛尔策拉（罗马，三八四年）\newline{}书信二十七　致玛尔策拉（罗马，三八四年）\newline{}书信二十八　致玛尔策拉（罗马，三八四年）\newline{}书信二十九　致玛尔策拉（罗马，三八四年）\newline{}书信三十　致保拉（罗马，三八四年）\newline{}书信三十一　致欧斯托基（罗马，三八四年）\newline{}书信三十二　致玛尔策拉（罗马，三八四年）\newline{}书信三十三　致保拉（三八四年）\newline{}书信三十四　致玛尔策拉（三八四年）\newline{}书信三十五　教宗达玛稣来函（罗马，三八四年）\newline{}书信三十六　致教宗达玛稣（罗马，三八四年）\newline{}书信三十七　致玛尔策拉（罗马，三八四年）\newline{}书信三十八　致玛尔策拉（罗马，三八五年）\newline{}书信三十九　致保拉（罗马，三八九年）\newline{}书信四十　致玛尔策拉（三八五年）\newline{}书信四十一　致玛尔策拉（罗马，三八五年）\newline{}书信四十二　致玛尔策拉（罗马，三八五年）\newline{}书信四十三　致玛尔策拉（罗马，三八五年）\newline{}书信四十四　致玛尔策拉（罗马，三八五年）\newline{}书信四十五　致阿塞拉（罗马，三八五年）\newline{}书信四十六　保拉与欧斯托基致玛尔策拉（白冷，三八六年）\newline{}书信四十七　致德西德留（白冷，三九三年）\newline{}书信四十八　致帕玛基\newline{}书信四十九　致帕玛基\newline{}书信五十　致多尼奥\newline{}书信五十一　厄丕法尼，撒拉米主教（居比路）致若望，耶路撒冷主教\newline{}书信五十二　致内波提安\newline{}书信五十三　致保利诺\newline{}书信五十四　致富利亚\newline{}书信五十五　致阿曼多\newline{}书信五十六　奥斯定来函\newline{}书信五十七　致帕玛基——论翻译的最佳方法\newline{}书信五十八　致保利诺\newline{}书信五十九　致玛尔策拉\newline{}书信六十　致赫利奥多若\newline{}书信六十一　致维基蓝修\newline{}书信六十二　致特兰基利诺\newline{}书信六十三　致德奥斐洛\newline{}书信六十四　致法比奥拉\newline{}书信六十五　致普林基比亚\newline{}书信六十六　致帕玛基\newline{}书信六十七　奥斯定来函\newline{}书信六十八　致卡斯特鲁提\newline{}书信六十九　致奥克雅诺\newline{}书信七十　致罗马演说家玛格诺\newline{}书信七十一　致路西诺\newline{}书信七十二　致维大利\newline{}书信七十三　致厄万格禄\newline{}书信七十四　致罗马人鲁菲诺\newline{}书信七十五　致德奥多拉\newline{}书信七十六　致阿比高\newline{}书信七十七　致奥克雅诺\newline{}书信七十八　致法比奥拉\newline{}书信七十九　致萨尔维纳\newline{}书信八十　鲁菲诺致玛卡里\newline{}书信八十一　致鲁菲诺\newline{}书信八十二　致亚历山大主教德奥斐洛\newline{}书信八十三　帕玛基与奥克雅诺来函\newline{}书信八十四　致帕玛基与奥克雅诺\newline{}书信八十五　致保利诺\newline{}书信八十六　致德奥斐洛\newline{}书信八十七　德奥斐洛致热罗尼莫\newline{}书信八十八　致德奥斐洛\newline{}书信八十九　德奥斐洛致热罗尼莫\newline{}书信九十　德奥斐洛致厄丕法尼\newline{}书信九十一　厄丕法尼致热罗尼莫\newline{}书信九十二　德奥斐洛主教致巴勒斯坦与居比路众主教的主教会议函\newline{}书信九十三　巴勒斯坦众主教致德奥斐洛\newline{}书信九十四　狄约尼削致德奥斐洛\newline{}书信九十五　教宗亚纳斯大西致因普利基雅诺\newline{}书信九十六　德奥斐洛来函\newline{}书信九十七　致帕玛基与玛尔策拉\newline{}书信九十八　德奥斐洛来函\newline{}书信九十九　致德奥斐洛\newline{}书信一百　德奥斐洛来函\newline{}书信一百零一　奥斯定来函\newline{}书信一百零二　致奥斯定\newline{}书信一百零三　致奥斯定\newline{}书信一百零四　奥斯定来函\newline{}书信一百零五　致奥斯定\newline{}书信一百零六　致孙尼亚与弗勒特拉\newline{}书信一百零七　致莱塔\newline{}书信一百零八　致欧斯托基\newline{}书信一百零九　致黎帕留\newline{}书信一百一十　奥斯定来函\newline{}书信一百一十一　奥斯定致普莱西迪\newline{}书信一百一十二　致奥斯定\newline{}书信一百一十三　德奥斐洛致热罗尼莫\newline{}书信一百一十四　致德奥斐洛\newline{}书信一百一十五　致奥斯定\newline{}书信一百一十六　奥斯定来函\newline{}书信一百一十七　致住在高卢的一位母亲与女儿\newline{}书信一百一十八　致犹利安\newline{}书信一百一十九　致米内维与亚历山大\newline{}书信一百二十　致赫狄比亚\newline{}书信一百二十一　致亚加西亚（四〇六年）\newline{}书信一百二十二　致鲁斯蒂古（四〇八年）\newline{}书信一百二十三　致阿格罗基（四〇九年）\newline{}书信一百二十四　致阿维托（四〇九或四一〇年）\newline{}书信一百二十五　致鲁斯蒂古（四一一年）\newline{}书信一百二十六　致玛尔策利诺与阿那普绪基（四一二年）\newline{}书信一百二十七　致普林基比亚（四一二年）\newline{}书信一百二十八　致高登齐（四一三年）\newline{}书信一百二十九　致达尔达诺（四一四年）\newline{}书信一百三十　致德默特利亚（四一四年）\newline{}书信一百三十一　奥斯定来函（四一五年）\newline{}书信一百三十二　奥斯定来函（四一五年）\newline{}书信一百三十三　致克特西丰（四一五年）\newline{}书信一百三十四　致奥斯定（四一六年）\newline{}书信一百三十五　教宗依诺森致奥勒留（四一七年）\newline{}书信一百三十六　教宗依诺森致热罗尼莫（四一七年）\newline{}书信一百三十七　教宗依诺森致耶路撒冷主教若望（四一七年）\newline{}书信一百三十八　致黎帕留（四一七年）\newline{}书信一百三十九　致阿普罗纽（四一七年）\newline{}书信一百四十　致司铎西彼廉（四一八年？）\newline{}书信一百四十一　致奥斯定（四一八年？）\newline{}书信一百四十二　致奥斯定（？）\newline{}书信一百四十三　致亚利比与奥斯定（四一九年）\newline{}书信一百四十四　奥斯定致奥普塔图（四二〇年）\newline{}书信一百四十五　致厄克斯培兰修（？）\newline{}书信一百四十六　致厄万格禄（？）\newline{}书信一百四十七　致撒彼尼亚诺（白冷，日期不详）\newline{}书信一百四十八　致贵妇策兰提亚（伪作）\newline{}书信一百四十九　论犹太节期（伪作）\newline{}书信一百五十　普罗科彼致热罗尼莫（伪作）\par

\SourceRecord{Translated by W.H. Fremantle, G. Lewis and W.G. Martley.；Nicene and Post-Nicene Fathers, Second Series；Vol. 6.；Edited by Philip Schaff and Henry Wace.；Buffalo, NY: Christian Literature Publishing Co.,；1893.；<http://www.newadvent.org/fathers/3001.htm>.}

% structure: p0001 source=h1 target=section reason=page-title

\section{书信一}

\Emph{致\Person{因诺森}}\par

\EditorialNote{这不仅是一系列书信中的第一封，很可能也是热罗尼莫现存最早的著作（约公元370年）。收信人因诺森是热罗尼莫在阿奎莱亚聚集的一小群热心人之一。他追随其朋友前往叙利亚，并于公元374年在那里去世。（见第三封书信第3节。）}\par

1. 最亲爱的\Person{因诺森}，你屡次请求我，不要缄默不语我们当代发生的一件奇事。我因谦逊而推辞了这项任务，且如我现在所感，是公正的，因我认为自己不能胜任，一方面是因为人的言辞不足以赞美天主，另一方面是因为怠惰如同锈蚀一般侵蚀了才智，使我仅有的些许表达能力亦已枯竭。你回应说，在属于天主的事上，我们不应看我们能完成的工作，而应看从事工作的心意；凡信靠\OriginalTerm{圣言}{Word}的人，绝不会缺少言辞。\par

2. 那么，我该怎么办呢？那任务超出了我的能力，而我却不敢推辞。我不过是一个毫无经验的乘客，却发现自己被安排负责一艘满载的船。我连在湖上划过小船都不曾，如今却要置身于\Place{攸克辛海}的喧嚣与动荡之中。我眼见海岸沉入地平线下，\Quote{四面都是天与水}；黑暗笼罩着水面，乌云如夜般漆黑，唯有浪花泛着白沫。你催我扬起鼓胀的帆，松开帆脚索，握紧船舵。终于，我顺从了你的命令，既然\OriginalTerm{爱德}{charity}能成就一切，我将信赖\OriginalTerm{圣神}{Holy Ghost}引领我的航程，无论结果如何，我都以此自慰。因为，倘若我们的船被浪花送入渴望的港湾，我将满足于人们说我的领航技术拙劣。但若因我未经锤炼的文辞，在语言的激流暗礁中搁浅，你可以责怪我的能力不足，但你至少会承认我的善意。\par

3. 那么，开始吧：\Place{韦尔切利}是利古里亚的一座城镇，坐落于阿尔卑斯山麓不远处，昔日重要，如今人口稀少，已趋衰落。当执政官巡视该地时，一名贫妇及其情夫被带到他面前——丈夫指控他们犯奸——两人均被投入令人恐惧的囚牢。不久后，试图通过酷刑逼供，当血淋淋的钩子击打那青年发青的皮肉，在他的肋旁犁出道道沟痕时，这不幸的可怜人欲借速死以逃避久苦。他虚假地指控自己的情欲，将另一人也卷入指控中；看来他是所有人中最悲惨的，他的处决是公正的，因为他使一个无辜的女子毫无自卫的余地。但那女子，虽在性别上较弱，在德行上却更强，尽管身体被拉上刑架，双手染满污秽并绑在背后，她却用那唯独刑吏无法束缚的眼睛望向天堂，泪流满面，说道：\Quote{\OriginalTerm{主耶稣}{Lord Jesus}，你鉴察一切，试探人的肺腑与心肠，你作见证。你见证我否认这指控并非为保命。我不肯说谎，因为说谎是罪。至于你，不幸的人，若你一心求死，为何必要毁灭两个人而非一个无辜者呢？我自己也愿死。我愿脱去这可憎的身体，却不愿以淫妇的身份。我献上我的颈项；我无惧地迎接闪亮的刀剑；但我将带着我的清白死去。那存心如此生活而被杀的人，并非真正死去。}\par

4. 那执政官，一直以这血腥的场景为乐，此时如同野兽，一经尝血便永远渴血，下令加倍酷刑，并残忍地咬牙切齿，威胁刑吏若不从一弱质女流身上逼出连男子之力都未能隐瞒的供词，便施以同样刑罚。\par

5. 求主耶稣助佑。为这祢的单个受造物，每一种酷刑都被设计出来。她被头发绑在柱上，全身比先前更牢固地固定在刑架上；火被拿来，烧她的脚；她的肋旁因刑吏的探针而颤抖；连她的乳房也未能幸免。这女子依然坚定不移，在精神上战胜了身体的痛苦，享受着良心的幸福，酷刑徒然在她周围肆虐。残暴的法官起身，为怒火所胜。她仍向天主祈祷。她的肢体被从关节中扯出；她只将目光转向天堂。另一人招认了被认为是他们共同犯下的罪。她为了这招认者的缘故，否认那招认，并冒着生命危险，为那同样处于危险中的人洗脱。\par

6. 与此同时，她只说一句话：\Quote{打吧，烧吧，撕裂我吧，随你的便；我没有做。若你不信我的话，必有一日这指控会被仔细查究。我有一位审判我的主。}终于，刑吏疲惫不堪，随着她的呻吟而叹息；也找不到一处可施新伤的地方。他被残酷所胜，颤抖着看着他撕裂的身体。执政官立刻在激愤中喊道：\Quote{旁观者啊，你们为何惊讶于一个女子宁愿受刑也不愿死呢？犯奸无疑需要两人；我认为一个淫妇否认罪行，比一个无辜青年承认罪行更可信。}\par

7. 同样地，两人被判了相同的刑罚，这对被判罪的男女被拖去处决。全体民众涌出观看；实际上，城门被蜂拥而出的人群挤得水泄不通，你可能会以为整座城市都在迁移。第一剑落下，那不幸青年的头便被砍下，无头的躯干在血泊中翻滚。随后轮到那女子。她跪在地上，闪亮的剑悬在她颤抖的脖颈上方。但尽管刽子手将全身力气集中在裸露的手臂上，当剑触及她的皮肉时，致命的刀锋却戛然而止，轻轻划过皮肤，只擦伤了一点，流了些血。刽子手惊恐地看到自己的手失去了力量，对自己失准的技艺和低垂的剑感到惊异，于是将剑高高举起准备再砍一次。剑再次无力地落在女子身上，无害地沉在她的脖颈上，仿佛钢铁害怕触碰她。愤怒喘息的军官敞开脖颈处的斗篷以全力挥剑，他抖落了系住斗篷边缘的胸针，却没有注意到，开始平衡剑准备再砍。\Quote{看，}女子喊道，\Quote{你肩上掉下了一件珠宝。捡起你辛苦挣来的东西，免得丢失。}\par

8. 我问，这种信心背后的秘密是什么？死亡临近，却对她毫无恐惧。被击打时她欢喜，而刽子手却脸色苍白。她的眼睛看到了胸针，却看不到剑。似乎临危不惧还不够，她还向残忍的敌人施恩。如今，天主圣三的神秘力量甚至使第三击也归于徒劳。惊恐的士兵不再相信剑刃，便将剑尖抵住她的喉咙，心想尽管剑刃可能切不开，但手的压力或能刺入她的血肉。这是万代未闻的奇事！剑弯曲至剑柄，仿佛战败般望着主人，承认自己无力杀人。\par

9. 让我援引三位青年的例子，他们在凉爽环绕的火焰中歌唱赞美诗而非哭泣，火焰无害地环绕他们的头巾和圣洁的头发。让我也回忆真福\Person{达尼尔}的故事，\EditorialNote{达尼尔第六章}狮子本是他的天然猎物，却在他面前蹲伏，摇尾惊口。愿\Person{苏撒纳}的杰出信德也升现在众人眼前；她被不公的判决定罪，却经由一位受圣神感召的青年而获救。主在这两件事上同样显示了仁慈；因为\Person{苏撒纳}被审判官释放，免死于剑下；而这女子，虽被审判官定罪，却被剑宣告无罪。\par

10. 如今民众终于揭竿而起保卫这女子。各个年龄的男女一同驱赶刽子手，在涌动的人群中围着他呼喊。几乎没有一个人敢相信自己的眼睛。令人不安的消息传至邻近的城市，全体警官队伍被召集。负责执行罪犯的军官从人群中冲出，\par

用泥土弄脏他花白的头发，\par

喊道：\Quote{怎么！公民们，你们是要取我的性命吗？你们想让我替她而死？尽管你们一心慈悲，极愿拯救一个被判罪的女人，然而我——我无罪——确实不应该死。}他含泪的哀求打动了人群，他们都因悲伤而麻木，表现出一种奇特的情绪变化。之前大家觉得为女子求情是责任，如今却觉得让她受刑是责任。\par

11. 于是取来新剑，任命新刽子手。受害者再次就位，唯一的力量仍是基督的恩宠。第一击使她颤抖，第二击使她摇摆不定，第三击她受伤倒地。哦，当受颂扬的天主大能的威严！先前她受四击毫发无损，片刻之后，似乎死去，以免一个无罪之人代她灭亡。\par

12. 那些负责用裹尸布包裹血迹斑斑的尸体的神职人员挖开泥土，堆起石头，筑成惯常的坟冢。日落迅速来临，因天主的仁慈，自然之夜比往常更快降临。突然，女子胸部起伏，眼睛寻找光明，身体重获生机。片刻后，她叹息，环顾四周，起身说话。最后她高喊：\Quote{主在我身边，我不畏惧。人能对我做什么？}\par

13. 同时，一位靠教会资金供养的老妇人，将灵魂归还于天，因它从天而来。\BibleRef{训 12:7}仿佛事态如此特意安排，因为她的身体取代了另一具，躺在土丘下。黎明灰暗时，魔鬼以警官的形式登场，要求交出自被杀女子的尸体，并想让人指给她坟墓。他对她竟会死去感到惊讶，以为她还活着。神职人员让他看新翻的草皮，指着最近堆起的泥土满足他的要求，用这样的话嘲弄他：\Quote{当然，挖出已埋葬的骨头吧！向坟墓重新宣战，如果这样还不满足，就把她肢解喂禽兽！只需七刀才杀死的人，仅仅死去太便宜她了。}\par

14. 面对这些辱骂的话，刽子手羞愧地退去，而女子在家中秘密复苏。然后，为了避免医生频繁造访教会引起怀疑，他们剪短她的头发，让她在几位贞女的陪伴下前往一所僻静的乡间屋舍。在那里，她将衣服换成男装，伤口结疤。然而，即使在为她行了如此伟大的奇迹之后，法律仍然向她发怒。多么真实：哪里有最多的法律，哪里就有最多的不义。\par

15. 但现在，看我故事的进展把我带到了何处；我们提到了我们朋友\Person{厄瓦格留斯}的名字。他为基督的事业所付出的努力如此巨大，以至于如果我认为有可能充分描述它们，那只会显示我的愚蠢；反之，如果我故意避开它们，我仍然无法阻止我的声音爆发出喜悦的呼声。谁能恰当地称赞他的警觉，使他能够，\Quote{如果可以这样说的话}，在\Place{米兰}的\Person{奥克森提乌斯}——那盘旋在教会之上的诅咒——死前就将他埋葬？或者谁能充分赞美他的智慧，他把罗马主教从几乎陷入的网罗中解救出来，并向他展示了同时战胜对手并在他们溃败时宽恕他们的方法？但是\par

这些话题我必须留给其他诗人，\newline{}被时间和空间的妒忌局限所阻隔。\par

我现在满足于记录我故事的结尾。\Person{厄瓦格留斯}寻求皇帝的特殊接见；以他的恳求烦扰他，以他的服务赢得他的青睐，最后通过他的热忱赢得了他的案子。皇帝释放了那位天主曾使其复活的妇女，恢复了她的自由。\par

\SourceRecord{Translated by W.H. Fremantle, G. Lewis and W.G. Martley.；Nicene and Post-Nicene Fathers, Second Series；Vol. 6.；Edited by Philip Schaff and Henry Wace.；Buffalo, NY: Christian Literature Publishing Co.,；1893.；<http://www.newadvent.org/fathers/3001001.htm>.}

% structure: p0001 source=h1 target=section reason=page-title

\section{第二封信}

\Emph{致\Person{德奥多西奥}及其他隐修士}\par

\EditorialNote{写于\Place{安提约基雅}，公元374年，当时\Person{热罗尼莫}仍对自己的未来道路存疑。\Person{德奥多西奥}似乎是叙利亚沙漠中隐修士的领袖。}\par

我多么渴望成为你们团体的一员，以我全部的力量拥抱你们可敬的团体！虽然，我这双可怜的眼睛实在不配目睹它。哦！但愿我能看见那沙漠，在我眼中比任何城市更可爱！哦！但愿我能看见那些偏僻之地，因圣人们的聚集而成为乐园！但由于我的罪阻止我将这颗满载过犯的头颅挤进你们蒙福的团体，我恳求你们（我知道你们能做到）以你们的祈祷救我脱离这世界的黑暗。我在你们那里时曾提及此事，如今写信又向你们重申同一请求；因为我全部心思都专注于这唯一的目标。实现我的决心取决于你们。我有意愿却无能力；能力只能因你们的祈祷而来。至于我，就像一只离开羊栈的病羊。除非善牧将我放在祂肩上背回羊栈，\BibleRef{路 15:3}我的脚步就会蹒跚，在起身的努力中，我会发现双脚发软。我就像荡子\BibleRef{路 15:11}，他虽然挥霍了父亲托付给我的一切，却至今未向他屈膝顺服；尚未开始摆脱往日放荡的诱惑。况且，我刚刚开始与其说是放弃恶习，不如说是渴望放弃它们，魔鬼就用新的罗网缠住我，在我路上设置新的绊脚石，从四面围困我。\par

四周是海，环绕着汪洋。\par

我发现自己身在大洋中央，不愿后退，也无法前进。唯有你们的祈祷能为赢得圣神的和风，将我吹向渴望之岸的港湾。\par

\SourceRecord{Translated by W.H. Fremantle, G. Lewis and W.G. Martley.；Nicene and Post-Nicene Fathers, Second Series；Vol. 6.；Edited by Philip Schaff and Henry Wace.；Buffalo, NY: Christian Literature Publishing Co.,；1893.；<http://www.newadvent.org/fathers/3001002.htm>.}

% structure: p0001 source=h1 target=section reason=page-title

\section{第三封信}

致隐修士\Person{鲁菲努斯}。\par

\EditorialNote{写于公元374年，从安提约基雅寄给在埃及的鲁菲努斯。杰罗姆叙述他的旅行以及到达叙利亚后发生的事件，特别是因诺森特和希拉斯的去世（§3）。他还描述了博诺苏斯的生活，他现在是亚得里亚海一个岛上的隐士（§4）。这封信的主要目的是说服鲁菲努斯来叙利亚。}\par

1. 天主赐予我们的超过我们所求的，\BibleRef{弗 3:20}并且祂常赐予我们那些\Quote{眼睛未曾看见，耳朵未曾听见，人心也未曾想到的}\BibleRef{格前 2:9}，这一点我确实从圣经的神圣声明中早已知道；但是现在，最亲爱的\Person{鲁菲努斯}，我在自己身上得到了证实。因为我曾幻想，若能通过书信的往来，仿佛与你肉身相见，已是过于大胆的愿望，但我听说你正在深入埃及最偏远的地方，拜访隐修士们，巡游天主的家庭于地上。哦，但愿主耶稣基督能突然将我带到你面前，如同\Person{斐理伯}被提到太监那里，\Person{哈巴谷}被提到\Person{达尼尔}那里一样\BibleRef{宗 8:26}，我会多么紧紧地拥抱你的脖子，多么深情地亲吻你的嘴——那张嘴曾与我一同在错误或智慧中相联。但我不配（不是你不该到我这里，而是我不配到你那里），而且我这虚弱的身体，即使健康时也软弱，又因频繁的疾病而破碎；所以我没有亲身前来，而是寄出这封信来迎接你，希望它能把你的心网入爱之网，引你到我这里。\par

2. 我最初的喜悦来自我们兄弟\Person{赫利奥多罗斯}传来的如此意外的佳音。我希望确定它，但不敢确信，尤其因为他告诉我他只是从别人那里听说，而且消息的新奇削弱了故事的可信度。我的愿望再次悬而未决，我的心摇摆不定，直到一位亚历山大的隐修士——他不久前因人民的虔诚热忱被送到埃及的宣信者（在意志上已是殉道者）那里——他的出现迫使我相信了这个消息。即便如此，我必须承认我仍犹豫。因为一方面他既不知道你的名字也不知道你的国家；但另一方面，他所说的似乎可能是真的，因为它与我之前得到的暗示一致。最后，真相完全显现在我面前，因为不断有路过的人流报告说：\Quote{\Person{鲁菲努斯}在\Place{尼特里亚}，他已到达了有福的\Person{玛卡里乌斯}的住所。}这时我抛开了一切限制我信念的东西，然后第一次真正为我的病痛感到悲伤。若不是我消瘦衰弱的身体束缚了我的行动，无论是夏日的炎热还是危险的航行，都无法延缓爱心的急促脚步。相信我，兄弟，我渴望见到你，胜过风暴中的水手渴望港口，胜过干渴的田野盼望雨水，胜过焦急的母亲坐在弯曲的海岸上等待儿子。\par

在那场突如其来的旋风把我从你身边拖走，用其不敬的撕扯割断了我们紧密相连的情感纽带之后，\par

深蓝雨云低垂于我头顶；\newline{}四面是海，四面是天。\par

我四处徘徊，不知往何处去。\Place{色雷斯}、\Place{本都}、\Place{比提尼亚}、整个\Place{迦拉达}和\Place{卡帕多细亚}，还有灼热的\Place{基利家}，一个接一个地消耗了我的精力。最后，\Place{叙利亚}作为遭难者最安全的港湾出现在我面前。在这里，经历了各种疾病之后，我失去了一只眼睛；因为\Person{因诺森特}，我灵魂的一半，被一场突如其来的热病夺走了。我现在所拥有的唯一眼睛，对我而言就是我们的\Person{埃瓦格里乌斯}，我带着不断的病痛成了他额外的负担。我们还有\Person{希拉斯}，圣\Person{梅拉尼娅}的仆人，他因无可指责的品行洗去了先前为奴的污点。他的死亡重新撕开了尚未愈合的伤口。但因宗徒的话禁止我们为安眠的人忧伤\BibleRef{得前 4:13}，而我的过度悲伤也被随后传来的喜讯所缓和，所以我最后讲述这件事，若你尚未听到，可以得知；若你已知道，可以与我一同喜乐。\par

4. 你的朋友——更确切地说，亦我亦你的朋友——\Person{博诺苏斯}，如今正攀登\Person{雅各伯}梦中预表的梯子\BibleRef{创 28:12}。他背起自己的十字架，不为明天忧虑\BibleRef{玛 6:34}，也不回头看所舍弃的\BibleRef{路 9:62}。他含泪播种，为能欢喜收割。如同\Person{梅瑟}以预像举起旷野中的蛇，他如今在现实中举起那蛇\BibleRef{户 21:9}。这是一个真实的故事，足以使希腊与罗马笔下的虚假奇事蒙羞。看，这里有一个少年，与我们一同受世界雅致的熏陶，家资丰厚，门第不亚于任何同伴；然而他离弃了母亲、姐妹和心爱的兄弟，像\Place{伊甸园}中的新农夫，定居在一座危险的海岛上，海浪在礁石四周咆哮；险峻的悬崖、光秃的岩石和荒凉的景象更显可怖。那里找不到农夫或隐修士。连你所知的\Person{敖尼息摩}——他曾因亲弟兄般的爱吻而喜悦——在这极度的孤寂中也不再陪伴他。独自在岛上——更确切地说，并非独处，因为基督与他同在——他看见天主的荣耀，连宗徒们也只在旷野中才得见。诚然，他看不见列阵的城邑，但他已将自己的名字登录在新城中。粗布衣裳使他的肢体显得丑陋，但穿着这样的衣服，他将更快被提升到云中迎接基督\BibleRef{得前 4:17}。没有赏心悦目的水道供他所需，但他从主的肋旁饮用生命之水。亲爱的朋友，请将这一切置于眼前，用你心灵的全部官能想象这幅场景。当你意识到这战士的努力，你便能赞美他的胜利。整个海岛四周是狂怒的海浪咆哮，蜿蜒海岸旁的悬岩随着浪涛的拍击发出回响。没有青草使地面变绿；没有阴凉的树丛，没有肥沃的田野。陡峭的崖壁环绕着他可怖的居所，如同监狱。但他无所挂念，无所畏惧，从头到脚以宗徒的甲胄武装\BibleRef{弗 6:13}，时而通过阅读\WorkTitle{圣经}听从天主，时而通过祈祷向主说话；或许当他逗留在岛上时，会见到类似\Person{若望}曾见的异象\BibleRef{默 1:9}。\par

5. 你想，魔鬼如今正编织怎样的罗网？正在准备怎样的计谋？也许，他记起故伎，将试图以饥饿试探\Person{博诺苏斯}。但已有人回答他说：\BibleQuote{人生活不只靠饼。} \BibleRef{玛 4:4} 也许他将向他展示财富与名声。但要对他说：\BibleQuote{那些想致富的人，陷于诱惑和罗网。} \BibleRef{弟前 6:9} 以及\BibleQuote{至于我，一切夸耀都在基督内。} \BibleRef{格前 1:31} 他或许会在肢体因禁食而疲乏时前来，用疾病的痛苦折磨它们；但宗徒的呼喊将击退他：\BibleQuote{因为我几时软弱，正是我有能力的时候。} 以及\BibleQuote{我的能力在软弱中才得以成全。} 他将以死亡的威胁恐吓；但回答将是：\BibleQuote{我渴望化解，而与基督同在一起。} \BibleRef{斐 1:23} 他将挥动火箭，但它们将被信德的盾牌所抵挡\BibleRef{弗 6:16}。总之，\OriginalTerm{撒殚}{Satan}将攻击他，但基督将保护他。主耶稣，感谢祢，在祢的日子我有一位能为我向祢祈祷的人。祢洞察一切心灵，查究心底的秘密，看见被困在鱼腹中大海深处的先知\BibleRef{纳 2:1}。祢知道，他与我如何从柔弱的婴儿一起成长到强壮的青年，如何在同一乳母的怀抱中受哺育，在同一搬运者的臂膀中受怀抱；又如何一同在\Place{罗马}求学后，住同一房屋，在\Place{莱茵河}半野蛮的岸边分享同一食物。祢也知道，是我首先开始寻求事奉祢。我恳求祢，请记念，祢的这位战士曾与我同为新兵。我面前有祢威严的宣告：\BibleQuote{谁教导却不实行，将被称为天国里最小的。} \BibleRef{玛 5:19} 愿他享有德行的冠冕，并因他每日的殉道而跟随身穿白衣的羔羊\BibleRef{默 14:4}。因为\BibleQuote{在我父的家里有许多住处。} \BibleRef{若 14:2} 而且\BibleQuote{这星与那星的光辉不同。} \BibleRef{格前 15:41} 赐我力量，使我的头能与圣人们的脚跟齐平！我立志，他却实行了。因此，请宽恕我未能持守决心，并请按他所应得的赏报他。\par

我或许已显得冗长，说了超过一封短信通常允许的篇幅；但我发现，每当我要赞扬我们亲爱的\Person{博诺苏斯}时，我总会如此。\par

6. 然而，回到我出发的起点，我恳求你，不要让我完全失落在你的视线和记忆之外。朋友是长久寻求、难得发现、难以保持的。让那些愿意的人，任凭黄金眩目，在荣华中飘浮，连他们的行李也闪耀着金银的光彩。爱无法购买，情感没有价格。能够终止的友谊从未真实。在基督内告别。\par

\SourceRecord{Translated by W.H. Fremantle, G. Lewis and W.G. Martley.；Nicene and Post-Nicene Fathers, Second Series；Vol. 6.；Edited by Philip Schaff and Henry Wace.；Buffalo, NY: Christian Literature Publishing Co.,；1893.；<http://www.newadvent.org/fathers/3001003.htm>.}

% structure: p0001 source=h1 target=section reason=page-title

\section{第四封信}

\Emph{致\Person{弗洛伦提乌斯}}\par

\EditorialNote{此信与前信一同寄给弗洛伦提乌斯，热罗尼莫请他转交鲁菲努斯。这位弗洛伦提乌斯是一位富有的意大利人，已退居耶路撒冷度隐修生活。热罗尼莫后来称他为杰出的隐修士，对穷人如此怜悯，以致普遍被称为穷人之父。 (Chron. ad A.D. 381.)}\par

1. 你的名声和圣德在众多不同民族的口中被传扬得多么广泛，你可以从这一点得知：我在认识你之前就开始爱你。因为，按照宗徒所说，\Quote{有些人的罪是明显的，如同先到审判台前}，同样，你的爱德之传闻如此广布，以致爱你不仅值得称赞，拒绝爱你反而被视为可耻。我略过你无数次支持基督、\BibleRef{玛 25:34} 喂养、衣着和探望祂的事迹。你在我兄弟\Person{赫略多洛}困苦时给予他的帮助，足以开启哑巴的口。他以何等的感激、何等的赞美，述说你为一位朝圣者铺平道路的仁慈。我虽是最迟钝的人，被难忍的疾病所消耗，但深切的情感和渴望使我双脚生翼，前来向你致敬并拥抱你。我祝你一切顺遂，并祈求主坚固我们新生的友谊。\par

2. 听说我们的兄弟\Person{鲁菲努斯}已与虔诚的女士\Person{梅拉尼}一同从埃及来到耶路撒冷。他以兄弟之爱与我密不可分；我恳请你帮我将附上的信交给他。但你不要以你在他身上发现的德行来判断我。因为在他身上，你会看到最清晰的圣德标记，而我只是尘土和卑贱的污泥，即使在活着时，也只是灰烬。我若能以虚弱的眼睛承受他卓越的光辉，就已足够。他刚洗净自己，是洁净的，甚至白如雪；而我，染满各种罪污，日夜战栗地等候偿还最后一分钱。\BibleRef{玛 5:26} 但既然\Quote{主释放被囚的}，并安息在痛悔和因祂言语战栗的人身上，\BibleRef{依 66:2} 或许祂会对躺在罪墓中的我说：\Quote{\Person{热罗尼莫}，出来吧。}\BibleRef{若 11:43}\par

可敬的司铎\Person{厄瓦格利乌斯}热烈地问候你。我们二人共同以尊敬之心问候弟兄\Person{玛尔提尼亚诺}。我很想见他，但被疾病的锁链所阻。愿你在基督内平安。\par

\SourceRecord{Translated by W.H. Fremantle, G. Lewis and W.G. Martley.；Nicene and Post-Nicene Fathers, Second Series；Vol. 6.；Edited by Philip Schaff and Henry Wace.；Buffalo, NY: Christian Literature Publishing Co.,；1893.；<http://www.newadvent.org/fathers/3001004.htm>.}

% structure: p0001 source=h1 target=section reason=page-title

\section{第五封书信}

\Emph{致\Person{弗洛伦修斯}}\par

\EditorialNote{此信写于前信数月之后（约公元374年底），发自叙利亚旷野。热罗尼莫先畅谈与弗洛伦修斯的友谊，并提及鲁菲努斯，随后提到他希望弗洛伦修斯寄给他的一些书籍的抄本。他还谈及一个逃跑的奴隶，弗洛伦修斯曾就此事写信给他。}\par

1. 亲爱的朋友，你的来信找到我时，我正住在旷野中最靠近叙利亚和撒拉森人的那一带。读信时，心中燃起如此强烈的愿望，想要前往耶路撒冷，以至于我几乎要违背我的隐修誓愿，以满足我的情感。因我不能亲自前来，便尽力以书信相代；如此，虽身不在你身边，却在爱德和心神中与你同在。\BibleRef{哥 2:5} 因为我恳切祈求，我们这幼嫩却已在基督内牢固凝结的友谊，永不被时间或距离撕裂。我们更应当通过书信的往来来加强这纽带。让这些书信在我们之间传递，在路上相遇，与我们交谈。若能保持这样的交流，情感不致损失太多。\par

2. 你信中说我们的兄弟\Person{鲁菲努斯}尚未到你那里。即使他来了，也几乎不能平息我的渴望，因为我现在无法见他。他离此处太远，而来此不便；我所选择的独修生活条件也禁止我前往他那里。因为我不再能随意随自己的意愿。因此我恳求你，请他允许你将可敬的\Place{奥古斯托杜努姆}主教\Person{雷提提乌斯}的注释书抄写一份，他在其中如此雄辩地解释了\WorkTitle{雅歌}。上述兄弟\Person{鲁菲努斯}的一位同乡，老人\Person{保禄}写道，鲁菲努斯持有他的\Person{戴尔都良}著作抄本，并急切请求归还。其次，我请求你让抄写员在纸上抄写某些书籍，附后的清单会告诉你哪些是我没有的。我也恳请你将可敬的\Person{希拉里}所著的\WorkTitle{达味圣咏释义}以及关于大公会议的详尽著作寄给我，这些书是当年我在\Place{特里尔}亲手为他抄写的。你知道，这样的书必是基督徒灵魂的食粮，以便他日夜默思上主的法律。\par

你对他人敞开你的屋顶，你呵护他们、安慰他们，用自己的钱财帮助他们；但就我而言，一旦你允准了我的请求，你便已给了我一切。况且，由于上主的慷慨，我拥有丰富的圣经藏书，你也可以向我发令。我将把你所喜欢的寄给你；不要以为你的命令会给我带来麻烦。我有一些专心于抄写的弟子。我也并非因为有所求而许下恩惠。我们的兄弟\Person{赫利奥多罗斯}告诉我，圣经中有许多部分你寻找而未能找到。但即使你拥有全部，情感肯定会主张自己的权利，并为自己寻求比已有更多的。\par

3. 至于你奴隶的当前主人——你曾如此荣幸地写信提及此事——我确信他就是拐带他的人。我还在\Place{安提约基雅}时，司铎\Person{埃瓦格里乌斯}常当着我的面斥责他。他对此回答说：\Quote{我无所畏惧。}他声称他的主人已将他打发走。如果你二人都要他，他就在这里；你想把他送到何处都行。我想我没有错，拒绝让一个逃跑者继续游荡。在旷野中，我自己无法执行你的命令；因此我请求我的挚友\Person{埃瓦格里乌斯}为你的缘故也为我的缘故积极处理此事。愿你安好于基督。\par

\SourceRecord{Translated by W.H. Fremantle, G. Lewis and W.G. Martley.；Nicene and Post-Nicene Fathers, Second Series；Vol. 6.；Edited by Philip Schaff and Henry Wace.；Buffalo, NY: Christian Literature Publishing Co.,；1893.；<http://www.newadvent.org/fathers/3001005.htm>.}

% structure: p0001 source=h1 target=section reason=page-title

\section{第六封信}

\Emph{致\Place{安提约基雅}的执事\Person{犹利安}}\par

\EditorialNote{这封信写于公元374年，其特别趣味在于提及热罗尼莫的妹妹。似乎她曾陷于罪中，并由执事犹利安恢复德性生活。热罗尼莫在下一封信（§4）中再次提到她。}\par

常言道：\Quote{说谎者即使说实话也不被相信}。从你责备我没有写信的方式，我看出这正是你对我的看法。我是否要说：\Quote{我常写信，但送信的人疏忽了}？你会回答：\Quote{你的借口是所有不写信之人的老生常谈。}我是否要说：\Quote{我找不到任何人带我的信}？你会说，有许多人从我的地区去了你那里。我是否要争辩说我确实给了他们信件？他们若不送达，就会否认收到。此外，我们之间距离如此遥远，很难查明真相。那我该怎么办呢？虽然其实不应受责备，但我恳求你原谅，因为我认为退让求和胜于坚守开战。事实上，身体的不断疾病和心灵的烦忧已使我如此虚弱，以至于死亡如此临近，我并未像平常一样镇定。为了不让你认为这个理由虚假，既然我已陈述了案情，我将像辩护人一样，传唤证人证明它。我们可敬的弟兄\Person{赫利奥多鲁斯}曾到过这里，但他虽然希望与我在旷野同住，却被我的罪恶吓走了。但我现在的冗长将弥补我过去的疏忽，因为正如\Person{贺拉斯}在他的讽刺诗中所说：\par

所有歌手都有这样一个令朋友们不快的缺点：\newline{}别人请他们唱时他们从不唱，不请时却唱个不停。\par

今后我将用成捆的信件淹没你，以致你将采取相反的态度，求我不要写信。\par

我高兴的是，我的妹妹——对你而言是基督内的女儿——仍坚守她的志向，这条新闻我首先要归功于你。因为在这里，我现在所在地，我不仅不知道我的故乡发生了什么，甚至不知道它是否还存在。即使伊比利亚毒蛇以它恶毒的毒牙撕裂我，我也不惧怕人的审判，因为我有天主审判我。正如有人所说：\par

即使你将世界击成碎片：\newline{}它也将落在一个无畏的头上。\par

那么，请你记住宗徒的训诫\BibleRef{格前 3:14}，即我们应该使我们的工作持久；在我妹妹的救恩中，为你自己从主那里预备报酬；并经常以书信增加我对我们共同分享的在基督内的荣耀的喜乐。\par

\SourceRecord{Translated by W.H. Fremantle, G. Lewis and W.G. Martley.；Nicene and Post-Nicene Fathers, Second Series；Vol. 6.；Edited by Philip Schaff and Henry Wace.；Buffalo, NY: Christian Literature Publishing Co.,；1893.；<http://www.newadvent.org/fathers/3001006.htm>.}

% structure: p0001 source=h1 target=section reason=page-title

\section{\WorkTitle{第七封信}}

\Emph{致\Person{克罗玛丢}、\Person{约维努斯}、\Person{欧瑟比}。}\par

\EditorialNote{本信（写于公元374年，同前信一样）由热罗尼莫写给他三位从前修道生活中的同伴。信中对博诺苏斯加以赞扬（第3节），为其妹寻求指引（第4节），并抨击斯德里敦主教路比奇努斯的行为（第5节）。}\par

凡因相互情愫而结合者，书页不应将他们拆散。因此，我不应将我的言语分给一人而另一人。因为将你们连结在一起的爱如此强烈，以致情感将你们三人结合在一种纽带中，其紧密不亚于自然联结你们中两人的那纽带。确实，若书写条件允许，我应将你们的名字合并，用一个符号来表达。我从你们那里收到的这封信本身，就在你们每人身上挑战我看到三人，并在三人中认出每一人。当可敬的\Person{埃瓦格里乌}在叙利亚人和撒拉逊人之间延伸的沙漠一角将它传给我时，我的喜乐无比。它完全超过了当年在\Place{罗马}挽回\Place{坎尼}败局、\Person{马尔凯鲁斯}在\Place{诺拉}击溃\Person{汉尼拔}部队时的欢庆。\Person{埃瓦格里乌}常来看我，在基督内珍惜我如同自己的心肠。然而，因他离我遥远，他的离开通常给我留下的遗憾正如他的到来给我带来的喜乐一样多。\par

我与你的信交谈，我拥抱它，它对我说话；在此地唯有它会说拉丁语。因为在此处，你要么学一种野蛮的土话，要么就闭口不言。每当我看到那些以熟悉的手迹写成的行文，它们就带回我珍爱的面孔，这时要么我不在此地，要么你与我同在。你若相信情感的真挚，我写此信时似乎看见了你们全体。\par

现在，一开始我想问一个冒昧的问题：为何我们被如此广大的陆海相隔，你却只给我寄来一封如此简短的信？是否因我未有先写信给你，而不配得到更好的待遇？我不相信当埃及继续供应纸草时你会短缺纸张。即使一位托勒密封锁了海洋，阿塔卢斯王也会从培尔加摩送来羊皮，这样你就能用他的皮张弥补纸张的不足。羊皮纸之名正源于多年前发生的这类史事。那么，我是否应假定送信者仓促？无论多长的信，一夜之间也可写成。还是你有要事缠身不能执笔？没有什么义务比情感更优先。剩下两种假定：要么你不愿意写，要么我不配收到信。二者之中，我宁愿指责你怠惰，也不愿判定自己不值得。因为纠正疏忽比激发爱德更容易。\par

你告诉我，\Person{博诺苏斯}如同一个真正的鱼之子，已经投水了。至于我，仍因旧污而龌龊，如同毒蛇和蝎子，我出没于干燥之地。\BibleRef{申 8:15} \Person{博诺苏斯}已将脚跟踏在蛇头之上，而我仍作那条蛇的食物，这条蛇按天主的定命吞食尘土。\BibleRef{创 3:14} 他已然能攀登那以渐进诗篇为预象的梯子；而我仍在其第一阶上哭泣，几乎不知道我是否能说：\Quote{我要向山举目，我的帮助从何而来。} 在世间威胁的波涛中，他正坐在他那岛屿的安全庇护所中，即教会围墙之内，也许此刻他正如同\Person{若望}被召唤去吃天主的书；\BibleRef{默 10:9} 而我，仍躺卧在我罪的坟墓中，被我的不义的锁链束缚，等待主在福音中的命令：\Quote{热罗尼莫，出来。} \BibleRef{若 11:43} 但\Person{博诺苏斯}所做的远不止此。他像先知一样，\BibleRef{耶 13:4} 将他的腰带带过幼发拉底河（因为魔鬼所有的力量都在腰间），并将它藏在那里的岩石穴中。然后，后来发现它已破烂，他唱道：\Quote{主啊，你占有我的肾。你已折断我的锁链。我要向你献上感恩的祭。} 但至于我，\Person{拿步高}用锁链将我带到\Place{巴比伦}，那心神纷乱的巴贝尔。在那里，他给我加上了被掳的轭；在那里，将铁环插入我的鼻孔，\BibleRef{列下 19:28} 他命令我唱一首\Place{熙雍}的歌。我曾向他说：\Quote{主释放被囚的；主开启瞎子的眼睛。} 为用一句话完成我的对比，当我祈求怜悯时，\Person{博诺苏斯}期待一顶冠冕。\par

我妹妹的归化是圣善的\Person{尤利安}努力的成果。他栽种了，浇水在于你们，而\BibleRef{格前 3:6} 主自会使它生长。耶稣基督将她赐予我，以安慰我因魔鬼在她身上造成的创伤。他已使她由死复生。但用外教诗人的话说，对她\par

\Quote{没有我不惧怕的安全。}\par

你们自己知道青春之路何等易滑——我本人也曾在这条路上跌倒，而你们正在其上行走，并非没有畏惧。她既踏上此路，就必须得到众人的忠告和鼓励，必须靠你们多次来信给予帮助，我可敬的兄弟们。并且——因为\Quote{爱德凡事包容}，\BibleRef{格前 13:7}——我请求你们从教宗\Person{瓦肋里安}处取得一封信以坚定她的决心。正如你们所知，一位少女的勇气会在她意识到高位之人关心她时得到增强。\par

事实上，我的家乡正遭蛮族蹂躏，那里的人唯一的天主就是他们的肚腹，他们只为今天而活，人越富有就越被认为是圣善。此外，用一个老生常谈的谚语，这盘子有与其相配的盖子；因为\Person{路比奇努斯}是他们的司铎。正如俗语所说，什么嘴唇配什么莴苣——据\Person{路奇里乌}说，这是\Person{克拉苏}唯一笑过的事——喻指驴子吃蓟。我的意思是，一个不稳的舵手掌舵一艘漏水的船，盲人领着盲人直入深坑。统治者与被统治的人一样。\par

6. 我以你所知的敬意，向你的母亲，也是我的母亲，致敬。由于她与你一同度圣善的生活，她作为你们的母亲，比你们这些圣善的儿女更早起步。因此，她的胎确实可称为金胎。我也与她一同问候你的姊妹们，她们无论到哪里都该受到欢迎，因为她们已战胜了性别和世界，正等候新郎的来临，\BibleRef{玛 25:4} 她们的灯里装满了油。哦，那家中有寡居的\Person{亚纳}、身为女先知的童女们、以及在圣殿中养大的孪生\Person{撒慕尔}，这房屋何其有福！那屋檐庇护着\OriginalTerm{玛加伯家族}{Maccabees}的殉道者母亲，她的儿子们围绕着她，个个都戴着殉道者的冠冕，何其有福！\EditorialNote{参阅：玛加伯下 第七章}因为虽然你每天藉遵守祂的诫命而宣认\Person{基督}，但你却在这私下的荣耀上，又加上了公开宣认的公共荣耀；因为从前正是藉着你，\OriginalTerm{亚略异端}{Arian heresy}的毒素才从你的城邑中被驱逐出去。\par

你或许对我这样在信末又重新开始感到惊讶。但我要怎么办呢？我无法抑制情感的表达。信函的简短篇幅迫使我保持沉默；而我对你的情感又催逼我说话。我匆匆写就，言语混乱，毫无条理；但爱情是不顾次序的。\par

\SourceRecord{Translated by W.H. Fremantle, G. Lewis and W.G. Martley.；Nicene and Post-Nicene Fathers, Second Series；Vol. 6.；Edited by Philip Schaff and Henry Wace.；Buffalo, NY: Christian Literature Publishing Co.,；1893.；<http://www.newadvent.org/fathers/3001007.htm>.}

% structure: p0001 source=h1 target=section reason=page-title

\section{第八封信}

\Emph{致\Place{阿奎莱亚}的副执事\Person{尼塞阿斯}}\par

\EditorialNote{\Person{尼塞阿斯}，副执事，曾陪同\Person{热罗尼莫}前往东方，但现已返回家乡。后来他继\Person{克罗马提乌斯}之后成为\Place{阿奎莱亚}的主教。此信写于公元374年。}\par

喜剧诗人\Person{图尔皮利乌斯}论书信往来时说，唯有它能使不在场者如在眼前。这句话虽出于虚构之作，却并非不真实。因为，对于分离的朋友而言，还有比在信中向所爱之人倾诉、并在信中倾听回复更为真实的在场吗？甚至那些意大利野人——即\Person{恩尼乌斯}笔下的卡斯卡人（据\Person{西塞罗}在其修辞学著作中所述，他们像野兽一样猎取食物）——在纸张和羊皮纸发明之前，也习惯于在粗刨的木板上或从树上撕下的树皮上写信。因此，人们称信使为\Quote{板持者}，称写信者为\Quote{用树皮者}，因为他们使用树皮。那么，我们这些生活在文明时代的人，岂不更应当践行一项即使处于极度野蛮状态且全然不知生活精妙之处的人们也曾履行的社交义务吗？你看，圣\Person{克罗马提乌斯}与可敬的\Person{欧塞比乌斯}，这两位既因性情相合又因天性相连的兄弟，通过他们寄给我的大量书信激励我勤勉尽责。然而你，才刚离开我，不仅未能维系我们新结的友谊，反而将它彻底撕裂——\Person{莱利乌斯}在\Person{西塞罗}的论文中明智地禁止这一做法。难道东方对你就如此可憎，以至于你连让书信来到这里都感到恐惧？醒来吧，醒来，从沉睡中振作，至少拿出一页纸给友爱。在家庭生活的欢乐中，有时也请为我们曾一同走过的旅程叹息。如果你爱我，就回复我的请求。如果你生我的气，即使生气也请回信。我发现，当我收到朋友的来信时，即使那朋友正对我生气，我渴望的灵魂也会得到极大的安慰。\par

\SourceRecord{Translated by W.H. Fremantle, G. Lewis and W.G. Martley.；Nicene and Post-Nicene Fathers, Second Series；Vol. 6.；Edited by Philip Schaff and Henry Wace.；Buffalo, NY: Christian Literature Publishing Co.,；1893.；<http://www.newadvent.org/fathers/3001008.htm>.}

% structure: p0001 source=h1 target=section reason=page-title

\section{第九封信}

\Emph{致\Place{阿奎莱亚}的隐修士\Person{克洛索努斯}}\par

\EditorialNote{一封写给冷淡通信者的戏谑信函，日期与前一封相同。}\par

\Person{赫略多罗}，他为我们两人所深爱，他对你的爱绝不亚于我对你的爱，也许已向你忠实地描述了我的心境；你的名字常在我唇边，每次与他交谈，我总以回忆与你愉快的交往开始，进而惊叹你的谦卑，赞扬你的德行，传颂你的圣爱。\par

据说，猞猁回头看时，便忘记刚刚看到的一切，失去对目光所及之物的全部记忆。你似乎也是如此。因为你完全忘记了我们共同的牵挂，你不仅模糊了，而且抹去了那封书信的印记——正如宗徒告诉我们的\BibleRef{格后 3:2}，那封信是写在基督徒的心上。我提到的那些动物潜伏在枝叶茂密的树上，扑向疾驰的狍子或受惊的雄鹿。它们的猎物枉然奔逃，因为带着折磨者一起跑，而折磨者在奔跑中撕裂它们的肉。然而，猞猁只在空腹使它们口干时才猎食。当它们喝足了血、填满了胃时，饱足便导致遗忘，它们不再思虑未来的猎物，直到饥饿提醒它们需要食物。\par

至于你，你不可能已经对我感到厌倦。那么，为什么你刚刚开始的友谊就过早地结束呢？为什么你还没有完全抓住就放手了呢？但是像你这样的疏忽从不缺乏借口，你也许会声称你没有什么可写的。即使如此，你也应该写信告诉我这个事实。\par

\SourceRecord{Translated by W.H. Fremantle, G. Lewis and W.G. Martley.；Nicene and Post-Nicene Fathers, Second Series；Vol. 6.；Edited by Philip Schaff and Henry Wace.；Buffalo, NY: Christian Literature Publishing Co.,；1893.；<http://www.newadvent.org/fathers/3001009.htm>.}

% structure: p0001 source=h1 target=section reason=page-title

\section{第十封信}

\Emph{致\Place{康科迪亚}的长者\Person{保禄}}\par

\EditorialNote{热罗尼莫写信给康科迪亚的保禄，一位百岁老人（§2），他拥有一个优秀的神学图书馆（§3），请求借给他一些注释书。作为回报，他将自己新写的隐修士保禄传记送给他。这封信的日期是公元374年。}\par

1. 人生命的短促是人对罪所受的惩罚；甚至在光明的门槛上，死亡常常抓住新生婴儿，这一事实证明时代不断陷入更深的腐败。因为当第一个耕种乐园的人被蛇缠绕在它蛇形的盘绕中，并因此被迫迁移到地上时，尽管他不死的状态变成了会死的状态，但人对诅咒的判决被推迟了九百年，甚至更久，这段时间如此之长，可称为第二次不朽。此后，罪逐渐变得越来越恶毒，直到巨人的不虔敬\BibleRef{创 6:4}带来了整个世界的毁灭。然后，当世界被洪水的洗礼——若我可这样称呼它——洁净后，人的生命缩短到很短的期限。然而，我们几乎完全浪费了这短短的生命，因为我们的不义如此不断地与天主的旨意作对。因为有多少人能活过一百岁，或者过了百岁而不后悔呢？正如圣经在\BibleBook{}{圣咏集}中所证实的：\BibleQuote{我们的寿数，不外七十春秋，若是强壮，也不过八十寒暑；但多半还是充满劳苦与空虚。}\par

2. 你可能会问，为何开篇如此遥远而牵强，以至于有人可以用贺拉斯的隽语来反驳它们：\par

\Quote{从\Person{勒达}为\Person{宙斯}所生的蛋讲起，\newline{}诗人执意要追溯特洛伊战争？}\par

简言之，是为了让我能用合适的言辞描述您的高龄和如基督一样白发的皓首\BibleRef{默 1:14}。看，百年的周行已从您身上经过，您却始终遵守主的诫命，在现世的环境中默想来世的福乐。您的眼睛明亮有神，步履稳健，听力良好，牙齿洁白，声音悦耳，肌肤坚实而充满活力；红润的面颊与白发相映，您的力量不似年迈。暮年并未如我们常见的那样损害您记忆的牢固；血液的冰冷并未使您温暖而谨慎的才智迟钝。您的面无皱纹，额无沟壑。最后，您的手没有颤抖，也没有在您写字的蜡板上画出弯曲的路径。主在您身上向我们展示了将来复活的青春，以至于在那些虽活着却逐渐死去的人身上，我们认出罪的后果；而在您身上，我们归功于正义，使您在异于常人的高龄仍保有青春。尽管我们在许多罪人身上也看到同样的身体康健，但在他们身上，这是魔鬼的许可，为引他们犯罪；而在您身上，这是天主的恩赐，为叫您喜乐。\par

3. 图利乌斯在支持弗拉库斯的精彩演讲中，将希腊人的学问描述为\Quote{天生的轻浮和精湛的虚浮。}\par

诚然，他们最有才学的文人常常为赞美他们的君王或王子而接受报酬。效法他们的榜样，我为我的赞美定了价。您不可认为我的要求很小。我请您给我福音的珍珠\BibleRef{玛 13:46}，即\BibleQuote{上主的言语是纯洁的言语，如同在熔炉中炼过的白银，在地上精炼了七次}——我指的是福尔图纳提安的注释书，以及（因其对迫害者的记述）奥勒留·维克托的史书，还有诺瓦提安的著作；这样，在了解这位分裂者提出的毒药后，我们就能更乐意地饮下圣殉道者西彼廉提供的解药。同时，我已将送给您，即送给年迈的保禄，一位更年长的保禄。我费了很大力气使我的语言适应更单纯者的水平。但不知怎的，即使用水装满，瓶子仍保留其最初的香气。如果我的小礼物令您满意，我还有其他储备，若圣神赐予顺风，将载满各种东方货物漂洋过海到您那里。\par

\SourceRecord{Translated by W.H. Fremantle, G. Lewis and W.G. Martley.；Nicene and Post-Nicene Fathers, Second Series；Vol. 6.；Edited by Philip Schaff and Henry Wace.；Buffalo, NY: Christian Literature Publishing Co.,；1893.；<http://www.newadvent.org/fathers/3001010.htm>.}

% structure: p0001 source=h1 target=section reason=page-title

\section{第十一封信}

\Emph{致埃莫纳的贞女们}\par

\EditorialNote{埃莫纳是离斯特利敦（热罗尼莫的出生地）不远的一个罗马殖民地。收信的贞女们没有回复他的信件，他现在写信责备她们的疏忽。这封信的日期是公元374年。}\par

这张单薄的纸显示出我生活在何等荒野之中，因此我不得不把许多话浓缩在寥寥数词里。因为，虽然我很想更充分地与你们交谈，但这可怜的小片纸迫使我留下许多未言之语。然而，机巧弥补了手段的匮乏，通过写小字，我能表达很多。我恳求你们注意，即使在困境中，我如何爱你们，因为连材料的匮乏也无法阻止我写信给你们。\par

我恳求你们，宽恕一个受了委屈的人吧：若我含泪愤愤而言，那乃是因为我受到了伤害。因为你们对我定期的来信，连一个音节也未回复。我知道，光明与黑暗不能相通，\BibleRef{格后 6:14} 天主的婢女与罪人也不能有份，然而，一个妓女却得以用眼泪洗主的脚，狗也允许吃主人桌下的碎屑。\BibleRef{玛 15:27} 救主的使命是召叫罪人，而非义人；正如祂亲自所说：\BibleQuote{健康的人不需要医生}\BibleRef{玛 9:12} 祂愿罪人悔改，而非死亡，\BibleRef{则 33:11} 且将迷途的可怜羔羊亲自背回家中。\BibleRef{路 15:5} 同样，当荡子归来时，他的父亲欢欢喜喜地迎接他。\BibleRef{路 15:20} 不仅如此，宗徒还说：\BibleQuote{什么都不必判断，等到主来。}\BibleRef{格前 4:5} 因为\BibleQuote{你是谁，竟判断别人的仆人？他或站或倒，自有他的主人。}\BibleRef{罗 14:4} 并且\BibleQuote{自己以为站得稳的，须要谨慎，免得跌倒。}\BibleRef{格前 10:12} \BibleQuote{你们各人的重担要互相担当。}\BibleRef{迦 6:2}\par

亲爱的姐妹们，人的嫉妒有一种判断方式，基督有另一种；角落里的低语与祂审判的宣判并不相同。许多在人看来似乎是正途，后来却发现是错的。\BibleRef{箴 14:12} 财宝常常藏在瓦器中。\BibleRef{格后 4:7} 伯多禄三次否认了他的主，但痛苦的眼泪使他恢复了原位。\BibleQuote{那得赦免多的，爱得也多。}\BibleRef{路 7:47} 没有话提到整群羊，但天使在天上因一只病母羊的安全而喜乐。\BibleRef{路 15:7} 如果有人对此争辩，主亲自说：\BibleQuote{朋友，我因为好，你的眼就红了吗？}\BibleRef{玛 20:15}\par

\SourceRecord{Translated by W.H. Fremantle, G. Lewis and W.G. Martley.；Nicene and Post-Nicene Fathers, Second Series；Vol. 6.；Edited by Philip Schaff and Henry Wace.；Buffalo, NY: Christian Literature Publishing Co.,；1893.；<http://www.newadvent.org/fathers/3001011.htm>.}

% structure: p0001 source=h1 target=section reason=page-title

\section{书信第十二}

\Emph{致隐修士安东尼}\par

\EditorialNote{本书信主题与前信相似。除某些手稿称安东尼为埃莫纳人外，对其一无所知。书信日期为公元374年。}\par

当门徒们争论谁为大时，我们的主，谦卑的导师，领了一个小孩子来，说：\BibleQuote{你们若不悔改而变成如同小孩子一样，绝不能进天国。}\BibleRef{玛 18:3} 为免祂似乎只讲不实行，祂在生活中实践了自己的训导。因为祂为门徒洗脚，\BibleRef{若 13:5} 祂以亲吻迎接那出卖祂的人，\BibleRef{路 22:47} 祂与撒玛黎雅妇人交谈，\BibleRef{若 4:7} 祂在玛利亚脚前谈论天国，而当祂从死者中复活后，首先显现给几个贫穷的妇女。骄傲与谦卑相对，撒殚因骄傲而失去了他作总领天使的高位。犹太人民因骄傲而灭亡，因为他们贪求会堂里的首位和街市上的请安，\BibleRef{玛 23:6} 却被从前被视为\BibleQuote{桶中一滴水}\BibleRef{依 40:15} 的外邦人所取代。两个贫穷的渔夫，伯多禄和雅各伯，被派去驳倒世上的诡辩家和智者。正如经上所说：\BibleQuote{天主拒绝骄傲的人，却赐恩宠给谦卑的人。}\BibleRef{伯前 5:5} 兄弟，试想：犯了以天主为敌的罪是何等重大！在福音中，法利塞人因骄傲而被摈弃，而税吏却因谦卑而蒙悦纳。\par

如今，若我没有弄错，我已给你寄了十封充满情谊、恳切真挚的信，而你却连一行也未曾屈尊回复。主向祂的仆人说话，而你，我的同仆兄弟，却拒绝向我说话。请相信我，若不是矜持约束了我的笔，我定会写一封令你不得不回复的斥责——即便那回复充满怒气。然而，怒气属乎人性，基督徒不应做出伤人的事，所以我再次恳求你：爱那爱你的人，并像仆人对待同仆那样给他写信。愿你在主内平安。\par

\SourceRecord{Translated by W.H. Fremantle, G. Lewis and W.G. Martley.；Nicene and Post-Nicene Fathers, Second Series；Vol. 6.；Edited by Philip Schaff and Henry Wace.；Buffalo, NY: Christian Literature Publishing Co.,；1893.；<http://www.newadvent.org/fathers/3001012.htm>.}

% structure: p0001 source=h1 target=section reason=page-title

\section{第十三封信}

\Emph{致他的姨母卡斯托里娜}\par

\EditorialNote{一封有趣的信，反映了热罗尼莫的家庭关系。他的姨母卡斯托里娜因某种原因与他疏远，他写信试图和解。是否成功，我们不得而知。信写于公元374年。}\par

宗徒兼福音作者若望在他的第一封书信中正确地说：\BibleQuote{凡恼恨自己弟兄的，便是杀人者。} \BibleRef{若一 3:15} 因为，既然杀人常由仇恨而生，那么恼恨者即使尚未杀害受害者，内心里已是一个杀人者。你可能会问，我为什么用这样的方式开头？只不过为了让我和你都能放下过去的恶感，洗净我们的心灵，使其成为天主的住所。达味说：\BibleQuote{你们生气，却不要犯罪；} 或者如宗徒更充分地表达的：\BibleQuote{不要让太阳在你们的愤怒中落下。} 那么在审判之日，我们这些愤怒的太阳不仅一天没有落下，而是许多年没有落下的人，该怎么办呢？主在福音中说：\BibleQuote{若你献礼物在祭台上，在那里想起你的弟兄有什么对你不满之处，就把礼物留在祭台前，先去与你的弟兄和好，然后再来献你的礼物。} \BibleRef{玛 5:23} 我有祸了，我这可怜的人；我几乎也要说，你也有祸了。这么长的时间以来，我们或者没有在祭台上献礼物，或者带着\Quote{无端的}愤怒献礼物。我们怎能在日常祈祷中说：\BibleQuote{宽免我们的罪债，如同我们也宽免得罪我们的人，} \BibleRef{玛 6:12} 而我们的情感却与言语相悖，我们的祈求与行为不一致呢？因此，我重申一年前在上一封信中所作的祈求：愿主的平安遗产 \BibleRef{若 14:27} 确为我们所有，愿我的渴望和你的情感在他眼中蒙悦纳。不久我们都要站在他的审判台前，为恢复和谐而获得赏报，或因破坏和谐而接受惩罚。倘若你不愿意——我希望不会如此——接受我的和解表示，那么我这一方面是自由的。因为这封信在被阅读时，将确保我的无罪。\par

\SourceRecord{Translated by W.H. Fremantle, G. Lewis and W.G. Martley.；Nicene and Post-Nicene Fathers, Second Series；Vol. 6.；Edited by Philip Schaff and Henry Wace.；Buffalo, NY: Christian Literature Publishing Co.,；1893.；<http://www.newadvent.org/fathers/3001013.htm>.}

% structure: p0001 source=h1 target=section reason=page-title

\section{第十四封信}

\Emph{致隐修士赫略多洛}\par

赫略多洛原为军人，如今是教会的司铎，曾随\Person{热罗尼莫}前往东方，但自觉未蒙召度沙漠独修生活，便返回\Place{阿奎莱亚}。在那里他重拾司铎职务，后蒙升任为\Place{阿尔提努姆}的主教。\par

此信写于离别的初苦，责备赫略多洛从完美的苦修生活后退。对苦修生活的描绘色彩浓烈，似乎在西方产生了巨大影响。\Person{法比奥拉}对此深深着迷，竟将整封信背诵下来。日期为公元373或374年。\par

如此您深知我们之间的情谊，您不能不认出我竭力延长我们共同在沙漠中度日的爱与热忱。这封信本身——如您所见，泪痕斑斑——便是我伴随您离去时的哀恸与哭泣之证明。您以孩童般的可爱方式，用抚慰的言语缓和了您的拒绝，而我因毫无防备，不知如何是好。我该沉默吗？我无法以冷漠的姿态掩饰我的渴望。还是该更恳切地乞求您？您已不愿倾听，因为您的爱不像我的爱。被轻视的爱只有一条路可走。既然不能在您身旁留住您，便要在您离去后寻找您。您离开时曾亲自请求我，当我定居沙漠时邀您前来，我也应允了。因此我现在邀请您；来吧，快快前来。不要回想旧日的牵连；沙漠是为那些舍弃一切的人预备的。也不要让我们从前旅行的艰难吓住您。您信基督，也当信祂的话：\BibleQuote{你们先该寻求天主的国，这一切自会加给你们。}\BibleRef{玛 6:33} 不要带口袋，也不要带棍杖。他足够富有，因为他贫穷——与基督一同贫穷。\par

但这是什么？我为何愚蠢地再次烦扰您？停止恳求，结束奉承的话。被冒犯的爱适合发怒。您已蔑视我的请求；也许您会听从我的劝诫。什么留住您，软弱的士兵，在您父亲家中？您的壁垒和战壕在哪里？您何时在田野里度过一个冬天？看，号角从天上传响！看，领袖驾云而来！\BibleRef{默 1:7} 祂武装起来要征服世界，口中吐出一把双刃利剑\BibleRef{默 1:16}，刈倒一切阻挡之物。至于您，您要做什么？径直从内室走向战场，从凉荫进入烈日？不，惯于穿短袍的身体不能忍受盾牌；戴惯软帽的头颅拒绝头盔；因闲散而娇嫩的手被剑柄磨伤。要听您国王的宣告：\BibleQuote{不与我相合的，就是反对我的；不同我收集的，就是分散的。}\BibleRef{玛 12:30} 要记住您入伍的那一天，当您与基督一同埋葬在圣洗中时，您向祂宣誓效忠，宣告为了祂您将不顾父母。看，仇敌正企图在您胸中杀死基督。看，敌军对您入伍时领受的赏金叹息不已。倘若您的小侄儿挂在您颈上，不要理他；倘若您母亲头发蒙灰、衣服撕裂，向您展示她哺育您的乳房，不要理会她；倘若您父亲俯伏在门槛上，您要践踏他而过，继续前行。以干涸的双眼奔向十字架的旗帜。在这种情况下，残忍才是唯一的真爱。\par

此后必将到来——是的，必将到来——有一天，您将作为胜利者返回真正的祖国，戴着勇冠的荣冕行走在天上的\Place{耶路撒冷}。那时您将与\Person{保禄}一同获得那里的公民权。那时您将为您的父母寻求同样的特权。那时您将为我代祷，因我曾激励您走上胜利之路。\par

我并非不知您可能以何种束缚为借口。我的胸膛并非铁石，我的心肠并非石头。我非出生自燧石，也非被母虎哺乳。我自己也经历过类似的烦恼。时而是一位寡居的姊妹，用爱抚的双臂拥抱您。时而是奴隶，您的乳兄弟，呼喊：\Quote{您要把我们留给哪位主人？} 时而是老迈的乳母，和如同父亲般珍爱的贴身仆从，呼喊：\Quote{只等我们死去，再跟随我们到坟墓吧。} 也许还有一位年迈的母亲，胸脯塌陷，额头布满皱纹，回忆她曾抚慰您的摇篮曲，加入她们的恳求。学者们若愿意，大可称您为：\par

\Quote{您家族唯一的支柱和栋梁。}\par

天主的爱与地狱的恐惧将轻易折断这些束缚。\par

您会争辩，圣经命令我们孝敬父母。\BibleRef{弗 6:1} 是的，但谁爱父母超过基督，就丧失自己的灵魂。\BibleRef{玛 10:37} 仇敌持剑要杀我，我还能顾念母亲的眼泪吗？还是为了父亲——若我是基督的仆人，我连为他举行葬礼的义务都没有，\BibleRef{路 9:59} 尽管若我是基督的真实仆人，我当为所有人行此礼——而背弃基督的服役？\Person{伯多禄}在\Person{主}受难前夕以懦弱的劝言成了祂的绊脚石；\BibleRef{玛 16:23} 而对于那些竭力阻止祂上\Place{耶路撒冷}的弟兄们，\Person{保禄}只有一个回答：\BibleQuote{你们为什么啼哭，使我心碎呢？我为主耶稣的名，不但准备受捆锁，也准备死在耶路撒冷。}\BibleRef{宗 21:13} 那常摧毁信德的亲情撞城锤，必从福音的城墙上无力弹回。\BibleQuote{谁遵行我在天之父的旨意，就是我的母亲，我的兄弟。} 他们若信基督，就当祝我一路平安，因为我为祂的名而战斗。他们若不信，\BibleQuote{让死人埋葬他们的死人。}\BibleRef{玛 8:22}\par

4. 但这一切，你争辩说，只涉及殉道者的情形。啊！我的弟兄，你错了，你错了，如果你以为基督徒任何时候都不会遭受迫害。当你完全不知自己受围攻时，你正受着最猛烈的围攻。\Quote{我们的仇敌如同咆哮的狮子巡行，寻找可吞吃的人。}\BibleRef{伯前 5:8} 而你以为有平安？\Quote{他坐在村庄的隐密处；在隐秘处杀害无辜；他的眼睛暗地里窥探穷人。他埋伏在隐密处，如狮子在洞穴中；他埋伏要掳掠穷人。} 而你却在树荫下酣睡，以致轻易成为猎物？一方面，放纵紧逼我；另一方面，贪婪试图侵入；我的肚腹想要在我内成为天主，取代基督；情欲妄图赶走住在我内的圣神，玷污祂的圣殿。\BibleRef{格前 3:17} 我说，我被仇敌追赶，\par

他的名字是\OriginalTerm{军旅}{Legion}，其诡计无数；\newline{}而我，不幸的人，当我被俘虏而去时，如何能自视为胜利者？\par

5. 我亲爱的弟兄，请仔细衡量各种形式的过犯，不要以为我所提及的罪比拜偶像的罪更轻。不，请听宗徒对此事的看法。\Quote{因为你们知道，}他写道，\Quote{凡是淫乱的、不洁的、贪婪的，也就是拜偶像的，在基督和天主的国里，都得不到产业。}\BibleRef{弗 5:5} 一般来说，凡属魔鬼的，都带有敌视天主的意味；而属魔鬼的，就是拜偶像，因为所有的偶像都服从于他。然而，保禄在别处明确无误地立下规矩，说：\Quote{所以要治死你们在地上的肢体，脱去淫乱、污秽、邪情、恶欲和贪婪，而贪婪就是拜偶像，为了这些事，天主的愤怒才降下。}\BibleRef{哥 3:5}\par

\OriginalTerm{拜偶像}{Idolatry}并不限于用手指和拇指将香撒在祭台上，或是将杯中的酒奠在碗里。\OriginalTerm{贪婪}{Covetousness}就是拜偶像，否则为三十块银子出卖主就是义行。\BibleRef{玛 26:15} \OriginalTerm{邪情}{Lust}涉及亵渎，否则人就可以与普通娼妓玷污那应成为\Quote{活祭，为天主所悦纳}\BibleRef{罗 12:1}的基督的肢体。\OriginalTerm{欺诈}{Fraud}就是拜偶像，否则在\WorkTitle{宗徒大事录}中，那些变卖家产，却因留下部分价钱而立即遭致灭亡的人，就值得效法了。请仔细想想，我的弟兄，没有什么是你应当保留的。\Quote{你们中不论是谁，}主说，\Quote{如不舍弃他所有的一切，不能做我的门徒。}\BibleRef{路 14:33} 你为什么如此三心二意地作基督徒呢？\par

6. 请看伯多禄如何撇下他的网；\BibleRef{玛 4:18} 请看税吏如何从税关起身。\BibleRef{玛 9:9} 转瞬之间，他成了宗徒。\Quote{人子没有枕头的地方，}\BibleRef{玛 8:20} 而你却规划宽阔的门廊和大厅？如果你指望继承世上的美物，你就不能再与基督一同作继承人。\BibleRef{罗 8:17} 你被称为隐修士，这名字没有意义吗？你一个独修者，为何混迹于人群之中？我所给予的忠告，并非来自一个从未失过船只或货物、从未经历过风暴的没有经验的水手。我自己近来遭遇船难，我对其他航海者的警告出于我自己的恐惧。一方面，如同卡律布狄斯，放纵将灵魂的救恩吸入漩涡；另一方面，如同斯库拉，邪情带着少女般的笑容，诱使它的船撞毁贞洁。海岸险恶，魔鬼带着一群海盗，带着镣铐捆绑他的俘虏。不要轻信，不要过于自信。大海可能平静如池，水面几乎无一丝微风，然而有时波涛如山。深处有危险，敌人潜伏在那里。放松帆索，张开船帆，将十字架作为旗帜竖在船头。你所称之为的平静，本身就是一场风暴。\Quote{为什么会这样？}你也许会争辩；\Quote{难道我所有的同乡不都是基督徒吗？}我回答说，你的情况与别人不同。请听主的话：\Quote{你若愿意成为成全的，去变卖你所有的，施舍给穷人，然后来跟随我。}\BibleRef{玛 19:21} 你已经许诺要成为成全的。因为当你离弃军队，并为了天国而自阉，\BibleRef{玛 19:12} 你这样做正是为了跟随成全的生活。如今，基督的成全仆人除了基督之外一无所有。如果他除了基督之外还有别的，他就不是成全的。如果他向天主许诺要成为成全的，却不成全，他的许诺就是谎言。但\Quote{说谎的口杀害灵魂。}\BibleRef{智 1:11} 因此，结论是：如果你成全，就不会贪恋你父亲的财物；如果你不成全，你就欺骗了主。福音发出神圣的警告：\Quote{你们不能侍奉两个主人，}\BibleRef{路 16:13} 居然有人敢同时侍奉天主和玛门，使基督成为说谎者吗？祂一再宣告：\Quote{若有人愿意跟随我，该弃绝自己，背起自己的十字架，跟随我。}\BibleRef{路 9:23} 如果我满载金银，我还能以为自己在跟随基督吗？当然不能。\Quote{那说自己住在祂内的，就应当照祂所行的去行。}\BibleRef{若一 2:6}\par

7. 我知道你会反驳说你一无所有。那么，你既如此准备好战斗，为何不上战场？也许你以为\Emph{你}可以在本乡作战，虽然主在本乡不能行什么奇迹？\BibleRef{玛 13:58} 为何不能？你问。凭他的权威，你得到这答复：\Quote{没有先知在本乡受接纳的。}\BibleRef{路 4:24} 但是，你会说，我不求荣誉；良心的赞许于我已足。主也不求荣誉；因为当群众要立他为王时，他躲开了他们。\BibleRef{若 6:15} 可是，没有荣誉就有轻蔑；有轻蔑就有频繁的粗鲁；有粗鲁就有烦恼；有烦恼就没有安宁；没有安宁，心思就容易偏离目标。再者，由于不安宁，热忱一旦失去部分力量，就因失去而减弱，减弱了的就不能称为成全。这一切的结论是：隐修士在本乡不能成为成全的。如今，不以成全为目标本身就是罪。\par

8. 你从这条防线被击退后，会求助于圣职人员的榜样。你会说，他们留在自己城中，却肯定无可指摘。我断不敢责备宗徒的继承人们，他们藉神圣的言语祝圣基督的身体，并使我们成为基督徒。他们拥有天国的钥匙，在审判日之前已在一定程度上审判人，并守护基督净配的贞洁。但正如我先前暗示的，隐修士的情况与圣职人员不同。圣职人员牧放基督的羊群；我作为隐修士，由他们牧养。他们靠祭台生活：\BibleRef{格前 9:13} 我若不向祭台献礼，斧子已放在我树根上，如同不结果子的树。\BibleRef{玛 3:10} 我亦不能以贫穷为借口，因为主在福音中称赞了一位年老的寡妇把仅有的两个小钱投入献仪箱。\BibleRef{路 21:1} 我不可坐在司铎面前；他若见我犯罪，可将我交给撒殚，\Quote{为使灵魂得救，而毁灭肉体。}\BibleRef{格前 5:5} 在旧律法下，违抗司祭的人被赶出营外，被民众用石头砸死，或被斩首，以血抵偿他的轻蔑。\BibleRef{申 17:5} 但如今，悖逆的人被神剑砍倒，或被逐出教会，被贪婪的恶鬼撕碎。若你弟兄们的恳求促使你领受圣职，我将因你被高举而喜乐，也惧怕你被摔下。你会说：\Quote{若有人渴望主教的职分，他渴望的是善工。}\BibleRef{弟前 3:1} 我知道；但你应加添下面的话：这样的人\Quote{必须无可指责，只做一个妇人的丈夫，警觉、节制、贞洁、品行端正、好客、善于教导、不酗酒、不暴力，而是忍耐。}\BibleRef{弟前 3:2} 在充分解释主教的资格后，宗徒同样慎重地谈到第三等的职务。\Quote{同样，执事必须庄重，}他写道，\Quote{不一口两舌，不饮酒过多，不贪图不义之财，以纯洁的良心持守信仰的奥秘。这些也该先受考验；然后，若证明无可指摘，才准他们任职。}\BibleRef{弟前 3:8} 那没有穿婚宴礼服就进去的人有祸了。留给他的只是那严厉的问题：\Quote{朋友，你怎么到这里来了？}当他哑口无言时，就会下令：\Quote{捆起他的手脚，把他丢在外面的黑暗里；在那里要有哀号和切齿。}\BibleRef{玛 22:11} 那领了一个塔冷通却把它包在手巾里的人有祸了；当别人获利时，他只保存了所领受的。他发怒的主人会斥责他说：\Quote{邪恶的仆人，}他发怒的主人会斥责他说，\Quote{你为什么不把我的钱存入银行，好让我回来时连本带利收回？}\BibleRef{路 19:23} 这就是说，你本应将你所不能担负的事放在祭台前。因为你，一个懒惰的商人，手里攥着一文钱，却占据了本可由他人加倍赢利的位置。因此，正如善尽职务的人为自己赢得良好的地位，\BibleRef{弟前 3:13} 照样，不相称地接近主之杯的人，将干犯主的身体和血。\BibleRef{格前 11:27}\par

9. 并非所有主教都是真正的主教。你注目伯多禄，也要留意犹达斯。你看见斯德望，也当看看尼苛劳，主亲口在\BibleBook{}{默示录}中定了他的罪，\BibleRef{默 2:6}，他那可耻的幻想产生了尼苛劳派异端。\BibleQuote{人应省察自己，然后才吃这饼和饮这杯。}\BibleRef{格前 11:28}因为使人成为基督徒的，并非教会的职位。百夫长科尔乃略还是外邦人时，就因圣神的恩赐而被洁净。达尼尔尚是孩童，就审判了长老。亚毛斯正在修剪桑树时，转眼间成了先知。\BibleRef{亚 7:14}达味原是个牧羊人，却被选为君王。主耶稣所爱的门徒中，最小的那位正是祂最爱的。我兄弟，你应坐在末位上，等那更卑微的人来时，你将被请到上座。\BibleRef{路 14:10}主安息在谁身上？岂不是那谦卑、痛悔、畏惧祂言语的人吗？\BibleRef{依 66:2}凡受托多的，向他要的也更多。\BibleRef{路 12:48}\BibleQuote{有权势的必受严厉的拷问。}\BibleRef{智 6:6}谁也不必在审判之日以肉身的贞洁自夸，因为那时人要为每一句闲话交账，\BibleRef{玛 12:36}辱骂兄弟的，将被视为杀人的罪。\BibleRef{玛 5:21}保禄和伯多禄如今与基督一同为王，取代前者或继承后者的职位都不容易。可能有天使来撕裂你圣殿的帐幔，\BibleRef{玛 27:51}并将你的灯台从原处挪去。\BibleRef{默 2:5}你若想建造城楼，应先计算花费。\BibleRef{路 14:28}失了味的盐，一无所用，只好被丢在外面，任人践踏。\BibleRef{玛 5:13}若是隐修士跌倒，司铎可为他转求；但失足的司铎，谁为他转求呢？\par

10. 我的论述终于绕过暗礁：这脆弱的小船终于从激浪中驶入深水。如今我可迎风扬帆；将争论的礁石抛在身后，我的结语将如水手的欢呼。哦，旷野，因基督的花朵而灿烂！哦，孤寂之地，从其中产生\BibleBook{}{默示录}里伟大君王之城的宝石！\BibleRef{默 21:19}哦，荒漠，因天主的特别临在而喜乐！我的兄弟，你既超越世界，为何仍留在世界？阴沉屋顶还要压抑你多久？烟雾弥漫的城市还要禁锢你多久？相信我，我比你更有光明。抛下身体的重量，翱翔于纯净明亮的以太，何其甜美！你惧怕贫穷吗？基督称穷人为有福的。\BibleRef{路 6:20}劳作使你恐惧吗？不挥汗如雨，就没有运动员获得桂冠。你为食物忧虑吗？信德不畏惧饥荒。你担心因禁食而瘦弱的四肢躺在硬地上吗？主就躺在你身边。你厌恶不洗的头和不梳的头发吗？基督是你真正的头。沙漠的无边孤寂令你恐惧吗？在精神上，你总可漫步于乐园。只要你的思想转向那里，你就不再在沙漠中。你的皮肤因不再沐浴而粗糙起鳞吗？那在基督内一次洗过的，无需再洗。\BibleRef{若 13:10}对你的所有异议，宗徒给出这简短的回答：\BibleQuote{现今的苦楚，与将来在我们身上要显示的光荣相比，便不足介意了。}\BibleRef{罗 8:18}我的兄弟，你太贪图享乐了，若想在此世与世人同乐，而在来世与基督一同为王。\par

11. 那日必将来到，必将来到，那时这必朽坏的要穿上不朽坏的，这必死的要穿上不死的。\BibleRef{格前 15:53}那时，那醒寤着的仆人，主来时必要看见他是有福的。\BibleRef{玛 24:46}那时，在号角声中，\BibleRef{得前 4:16}大地及其上的人民都要颤抖，而你却要喜乐。世界将为主来审判它而哀号，地上的各族都要捶胸。昔日的强悍君王将赤身战栗。维纳斯和她儿子都将暴露。朱庇特带着他的霹雳将受审判；柏拉图和门徒们将成愚拙。亚里士多德的论证将毫无用处。你可能看似贫穷、土气，但那时你将欢欣大笑，说：请看我的被钉的主，请看我的审判者。这就是那曾包裹在襁褓中、在马槽里啼哭的婴孩。\BibleRef{路 2:7}这就是那父母是工匠和妇女的一位。这就是那被抱在母亲怀中逃往埃及、虽然他是天主却躲避人面的一位。这就是那身穿紫红袍、头戴荆棘冠冕的一位。\BibleRef{玛 27:28}这就是那被称为巫师、附魔者和撒玛黎雅人的一位。\BibleRef{若 8:48}犹太人，请看你们钉在十字架上的那双手。罗马人，请看你们用长矛刺穿的肋旁。你们俩都看看，是不是这身体被门徒在夜间秘密偷去了？\BibleRef{玛 27:64}因为你们自称相信这事。\par

我的兄弟，是友爱激励我说了这些话，为叫你现在虽觉基督徒生活艰难，却能在那一日得到你的赏报。\par

\SourceRecord{Translated by W.H. Fremantle, G. Lewis and W.G. Martley.；Nicene and Post-Nicene Fathers, Second Series；Vol. 6.；Edited by Philip Schaff and Henry Wace.；Buffalo, NY: Christian Literature Publishing Co.,；1893.；<http://www.newadvent.org/fathers/3001014.htm>.}

% structure: p0001 source=h1 target=section reason=page-title

\section{\WorkTitle{第十五封书信}}

\Emph{致教宗达玛稣}\par

\EditorialNote{本信写于公元376或377年，展现了热罗尼莫此时对达玛稣所持的罗马宗座的态度，达玛稣后来成为他亲密的朋友和仰慕者。热罗尼莫提到罗马是他本人受洗之地，并作为一处信德未受损害的教会（§1），并阐述了唯有在教会内才能得救的严格教义（§2），他向渔夫的继承者提出了两个问题：（1）在安提约基雅宗座的三个声请人中，谁才是真正的主教；（2）关于天主圣三，正确的术语是使用三个位格还是一个？关于后一个问题，他充分表达了自己的看法。}\par

1. 既然东方长期以来的民族纷争使其支离破碎，正一块块撕裂那件从上到下浑然织成的主的无缝长衣，\BibleRef{若 19:23} 既然狐狸在摧毁基督的葡萄园，\BibleRef{歌 2:15} 而在那些破裂不能蓄水的池中，很难发现\Quote{封闭的泉源}和\Quote{关锁的园子}，\BibleRef{歌 4:12} 我认为有责任咨询伯多禄的宝座，并求助于那信德曾受保禄赞扬的教会。我向那曾赐我基督衣袍的教会祈求属灵的食物。我们之间的辽阔海陆不能阻止我寻找\Quote{宝贵的珍珠}。\BibleRef{玛 13:46} \Quote{无论在哪里有尸体，老鹰也必聚集在那里。}\BibleRef{玛 24:28} 邪恶的儿女挥霍了他们的产业；惟独你保持你的遗产完好无损。罗马肥沃的土壤，当它接受主的纯种时，结出百倍的果实；但在这里，种子在垄沟中被窒息，只长出莠子或燕麦。\BibleRef{玛 13:22} 在西方，正义的太阳\BibleRef{拉 4:2}正在升起；在东方，从天上坠落的晨星\BibleRef{路 10:18}再次在众星之上设立宝座。\BibleRef{依 14:12} \Quote{你们是世界的光，}\BibleRef{玛 5:14} \Quote{你们是地上的盐，}\BibleRef{玛 5:13} 你们是\Quote{金器和银器。}\BibleRef{弟后 2:20} 这里却是木器和瓦器，等待铁杖\BibleRef{默 2:27}和永火。\par

2. 然而，虽然你的伟大使我畏惧，但你的仁慈吸引我。我从司铎要求保全牺牲，从牧者要求保护羊群。愿一切傲慢远离；愿罗马的威严退去。我的话是对渔夫的继承者，对十字架的门徒说的。因为除了基督，我不跟随任何领袖，所以我只与你的真福共融，也就是与伯多禄的宝座共融。因为我知道，这基石上建立了教会！\BibleRef{玛 16:18} 这是唯一能正确吃逾越节羔羊的房屋。\BibleRef{出 12:22} 这是诺厄的方舟，凡不在其中的，必在洪水中灭亡。\BibleRef{创 7:23} 但由于我的罪，我来到了这片位于叙利亚和蛮荒之间的旷野，我因我们之间的遥远距离，不能时常向你的圣洁祈求主的圣物。因此，我在此追随那些与你同信德的埃及宣信者，并将我的小舟停泊在他们大船的阴影下。我不认识维大利；我拒绝梅肋提；我与保利诺无关。那不与你一起聚集的，就是分散；\BibleRef{玛 12:30} 那不属于基督的，就是属于敌基督。\par

3. 目前，很遗憾，那些亚略派人士——\OriginalTerm{旷野派}{Campenses}——正试图向我这个罗马基督徒强加他们闻所未闻的\OriginalTerm{三个位格}{three hypostases}的公式。而且这是在尼采亚信经的定义和亚历山大的法令之后，西方也已加入其中。我倒想知道，这些教义的宗徒在哪里？他们的保禄，这位外邦人的新导师在哪里？我问他们，所谓的三个位格是什么意思。他们回答说，\OriginalTerm{三个自立体}{three persons subsisting}。我回答说，这就是我的信仰。他们不满足于含义，要求这个术语。确实，这些话中隐藏着某种毒药。我喊道：\Quote{若有人拒绝承认三个位格意思是三个实体化的东西，即三个自立体，愿他受诅咒。} 然而，因为我不学习他们的措辞，我就被视为异端。\Quote{但是，若有人将位格理解为本质，而否认在三个位格中有一个本质，他就与基督无分。} 因为这是我的宣认，我像你一样，被烙上\OriginalTerm{撒伯流异端}{Sabellianism}的污名。\par

4. 若你认为合适，颁布一道谕令；那么我将毫不迟疑地讲三个\OriginalTerm{位格}{hypostases}。下令制定一个新的信经以取代尼西亚信经；那么无论我们是亚略派还是正统派，一个宣信将适用于我们所有人。在世俗学问的整个范围内，\Quote{位格}从来都只意味着本质。请问，有谁如此亵渎，竟敢在神性中讲三个本质或实体？天主的本性只有一个，唯一的一个；唯有此本性真正地\Emph{是}。因为绝对的存在并非源于其他，而是自给自足。其他一切，即所有受造之物，虽似存在，实则不然。因为它们曾不存在，而一度不存在的事物也可能再次消失。唯有永恒的天主，即无始者，才真正配称为\Quote{存在}。因此祂从荆棘中对梅瑟说：\Quote{我是自有者，}\Quote{自有者派遣了我。}\BibleRef{出 3:14}当时天使、天空、大地、海洋都已存在，所以天主必然是以绝对存在者的身份，将那个显然为万物共有的名字归于自己。但因祂的本性唯独完美，且在三个位格中只有一个神性，它真实存在且是一个本性；所以凡以宗教名义宣称神性中有三个元素，即三个\OriginalTerm{位格}{hypostases}或本质的人，实际上是在为天主设定三个本性。若果真如此，我们为何与亚略隔墙相隔，而实际上我们与他虚假地成为一体？让乌尔西奇诺成为你阁下之同僚，让奥克森提乌斯与盎博罗削联合。但愿罗马信德永不至此！愿你子民虔敬的心永不被如此不圣洁的教义感染！让我们满足于讲一个实体和三个自立存在的位格——完美、平等、同永恒。若你乐意，让我们只讲一个位格，不讲三个。那些意思相同却用不同词语的人是不祥之兆。让我们满足于迄今为止所使用信经的形式。或者，若你认为我应当讲三个位格并解释它们的含义，我乐意服从。但请相信我，他们的蜜中藏有毒药；撒殚的天使已伪装成光明天使。\BibleRef{格后 11:14}他们为\Quote{位格}一词给出了看似合理的解释；然而当我声称持相同理解时，他们却视我为异端者。为什么他们如此执着一个词？为什么他们躲在含混语言背后？若他们的信仰与解释相符，我并不谴责他们坚持它。反之，若我的信仰与他们所表述的意见相符，他们应允许我用自己的话表达他们的含义。\par

5. 因此，我恳求你阁下，因被钉十字架的救世主和同体的天主圣三，授权我用信函允许我使用或拒绝使用\Quote{三个位格}这一措辞。为避免我目前居所的隐秘性给信使带来困惑，我请你将信函寄给你熟识的司铎\Person{厄瓦格略}。我还请你指明在安提约基雅我应与谁共融。希望不是阵营派；因为他们——连同塔尔索的异端者盟友——只渴望与你们共融，以便更有权威地宣讲他们传统关于三个位格的教义。\par

\SourceRecord{Translated by W.H. Fremantle, G. Lewis and W.G. Martley.；Nicene and Post-Nicene Fathers, Second Series；Vol. 6.；Edited by Philip Schaff and Henry Wace.；Buffalo, NY: Christian Literature Publishing Co.,；1893.；<http://www.newadvent.org/fathers/3001015.htm>.}

% structure: p0001 source=h1 target=section reason=page-title

\section{\WorkTitle{第十六封信}}

\Emph{致教宗达玛苏}\par

\EditorialNote{这封信写于前信数月之后，是再次恳请达玛苏解决作者的疑惑。热罗尼莫再次提及他在罗马所受的圣洗，并宣告他对安提约基雅各派系的唯一回答是：凡紧握伯多禄之座者，我皆接纳。写于旷野，年代为377或378年。}\par

1. 福音中那位寡妇凭借她的坚持，最终得到了回应，\BibleRef{玛 15:28}；同样，一人使另一人在半夜借给他饼，当门已关上、仆人们都睡了的时候。\BibleRef{路 11:7} 税吏的祈祷战胜了天主，\BibleRef{路 18:10}，尽管天主是不可战胜的。尼尼微因它的眼泪得救，免于因罪而来的即将临到的毁灭。\BibleRef{纳 3:5} 你问，这些引经据典有何目的？我要回答说：是为了让你——尊大者——能顾念我的渺小；让你——富有的牧人——不要嫌弃我这有病的羊。基督亲自将强盗从十字架带到了乐园，\BibleRef{路 23:43}，而且，为表明悔改永不嫌迟，祂将杀人犯的死变成了殉道。基督欣然拥抱回头的荡子；\BibleRef{路 15:20}，并且，好牧人撇下九十九只，将那一只可怜的亡羊扛在肩上带回家。\BibleRef{路 15:5} 保禄从迫害者变成了宣讲者。他肉体的眼睛被弄瞎，为的是使他灵魂的眼睛明亮，\BibleRef{宗 9:8}；他曾经将基督的仆役用锁链拖到犹太人公议会前，\BibleRef{宗 8:3}，后来却以基督的锁链为荣。\BibleRef{格后 12:10}\par

2. 正如我已经写信告诉你的，我，曾在罗马接受基督的衣袍，如今被滞留在叙利亚边界的旷野。然而并没有放逐的判决加于我；我所受的惩罚是自加的。但正如异教诗人所说：\par

跨越海洋的人，改变的是天空，而非心灵。\par

不懈的仇敌紧跟着我，我在旷野所受的攻击比以往更甚。因为亚略派的狂怒肆虐，世俗的权势支持它。教会分裂为三派，每一派都急于将我拉拢过去。隐修士的影响由来已久，并且针对我。我一直在呼喊：\Quote{凡紧握伯多禄之座者，我皆接纳。} 默肋提乌斯、维塔利斯和保利诺都声称依附于你，如果只有一人这样说，我可以相信。但事实是，他们要么两人，要么三人都在说谎。因此，我恳求你——你的蒙福者，藉着我们主的十字架与受难，这些我们信仰中必要的荣耀，既然你持有宗徒的职位，就请做出宗徒的决断。只请你写信告诉我，在叙利亚我应与何人共融。我将为你祈祷，使你得以与十二宗徒同坐审判，\BibleRef{玛 19:28}；使你年老时，如同伯多禄，不是自己束带，而是别人给你束带，\BibleRef{若 21:18}；并如保禄一样，成为天国的公民。不要轻视基督为其而死的灵魂。\par

\SourceRecord{Translated by W.H. Fremantle, G. Lewis and W.G. Martley.；Nicene and Post-Nicene Fathers, Second Series；Vol. 6.；Edited by Philip Schaff and Henry Wace.；Buffalo, NY: Christian Literature Publishing Co.,；1893.；<http://www.newadvent.org/fathers/3001016.htm>.}

% structure: p0001 source=h1 target=section reason=page-title

\section{第十七书}

致司铎马尔谷\par

\EditorialNote{在这封信中，收信人似乎是加采东沙漠隐修士中的一位杰出人物，热罗尼莫抱怨因拒绝参与叙利亚当时激烈进行的神学争论而遭受的苛刻待遇。他声明自己的正统信仰，并请求允许他留在原地直到春天返回，届时他将离开这荒凉的沙漠。写于公元378或379年。}\par

1. 我本已决定采用圣咏作者的话：\Quote{恶人站在我面前时，我默然无语；我自卑自谦，甚至连善言也缄默不言}以及\Quote{我如同聋子，不闻不问；我好像哑巴，不开口说话。我就如同一个听不见的人。}但爱德超越一切\BibleRef{格前13:7}，我对你的敬意战胜了我的决心。的确，我不太关心报复攻击我的人，而更关心满足你的请求。因为在基督徒中，正如有人所说，不幸的不是忍受侮辱的人，而是施加侮辱的人。\par

2. 首先，在向你谈及我的信仰（你对此十分清楚）之前，我不得不控诉这个地方的无人性。一句陈旧的引文最能表达我的意思：\par

这是何等野蛮的人，他们不愿给予\newline{}连陌生人在沙滩上的歇息也不允许！\newline{}他们以战争相威胁，将我们从海岸驱赶。\par

我引用一位外邦诗人的话，好让那些不顾基督和平的人至少从外邦人那里学到其含义。我被称作异端者，尽管我宣讲同体天主圣三。我被指控犯有撒伯里乌异端，尽管我不停地宣告：在神性中有三个分别的、实在的、完整的、完美的位格。亚略派指责我是对的，但正统人士若攻击我这样的信仰，就失去了他们自己的正统性。他们若愿意，可以谴责我为异端者；但若如此，他们也必须谴责埃及和西方、达玛苏和伯多禄。为何只把罪责加在一个人身上，而对其他同伴不加谴责？如果溪流中水很少，那是源头的问题，不是水道的过错。我羞愧地说，从充当静室的洞穴中，我们沙漠的隐修士谴责世界。我们裹着苦衣和灰烬，对主教们作出判决。如果悔改的袍子掩盖了君王的骄傲，那又有什么用？锁链、污秽、长发，理应是悲伤的标记，而不是王权的象征。我只请求保持沉默。为什么他们折磨一个不值得他们恶意的人？你说我是异端者。如果我是，你管得着吗？保持安静，就一切了事。我猜你是害怕，凭我的流利叙利亚语和希腊语知识，我会巡视各教会，误导人民，制造分裂！我没有抢夺任何人的东西；也没有获取我不劳而获之物。我用我自己的手\BibleRef{格前4:12}每日在额上的汗水中\BibleRef{创3:19}劳动以挣取食物，因为知道宗徒所写：\Quote{谁若不愿意工作，就不应当吃饭。}\BibleRef{得后3:10}\par

3. 可敬的神父，耶稣是我的见证，我是以何等的叹息和眼泪写这一切的。\Quote{主说：我一直沉默，但岂能永远沉默？决不。}我连沙漠一隅也不可得。每天有人要我宣认信仰，好像我在领受圣洗时重生时从未宣认过似的。我接受他们的信辞，但他们仍不满意。我签名认可，但他们仍不肯相信我。唯一能令他们满意的是我离开这个地方。我即将离去。他们已经从我这里夺走了我亲爱的弟兄们，他们是我生命的一部分。你看，他们急切地想离开——不，他们实际上正在离开；他们说，生活在野兽中间也比与这样的基督徒在一起更好。我自己此刻也会逃跑，若非身体虚弱和冬天严寒的阻挠。我请求在沙漠中避难几个月，直到春天回来；或者如果这显得太久，我准备现在就离开。\Quote{大地及其中的万物属于主。}让他们独自爬上天堂；基督仅为他们而死；他们拥有一切，并以一切为荣。就这样吧。\Quote{至于我，我只以我们的主耶稣基督的十字架来夸耀，借着祂，世界于我已被钉在十字架上，我于世界也被钉在十字架上。}\BibleRef{迦6:14}\par

4. 关于你认为应向我提出的信仰问题，我已将一份书面的信仰宣认交给了可敬的济利禄，足以回答这些问题。不相信这宣认的人，在基督内没有分。我的信仰由你的耳朵和你那有福的弟兄则诺比乌的耳朵作证，我们这里所有人都向你和他致以最诚挚的问候。\par

\SourceRecord{Translated by W.H. Fremantle, G. Lewis and W.G. Martley.；Nicene and Post-Nicene Fathers, Second Series；Vol. 6.；Edited by Philip Schaff and Henry Wace.；Buffalo, NY: Christian Literature Publishing Co.,；1893.；<http://www.newadvent.org/fathers/3001017.htm>.}

% structure: p0001 source=h1 target=section reason=page-title

\section{第十八封信}

第十八封信。致教宗\Person{达玛稣}。\par

\Quote{见：\BibleRef{若 12:41}，先知所见的是基督，而非圣父。}他又说：\Quote{\OriginalTerm{色辣芬}{seraphim}一词或意为\InnerQuote{灼热}，或意为\InnerQuote{言语的开端}；两只色辣芬因而代表\WorkTitle{旧约}和\WorkTitle{新约}。\InnerQuote{我们的心不是火热吗？}门徒说，\InnerQuote{当他给我们讲解圣经的时候？}见：\BibleRef{路 24:32} 况且，\WorkTitle{旧约}是用希伯来文写成的，这无疑是人类的原始语言。}\Person{热罗尼莫}随后谈论圣书的统一性。\Quote{无论什么，}他断言，\Quote{我们在\WorkTitle{旧约}中所读的，也在福音中找到；我们在福音中所读的，都出自\WorkTitle{旧约}。两者之间没有不和谐，没有分歧。在两个约中，都宣讲天主圣三。}\par

这封信因其提供的证据而引人注目，这些证据展示了\Person{热罗尼莫}研究的彻底性。他不仅引用了几种希腊文\BibleBook{依撒意亚}{}译本以支持自己的论点，而且还追溯到希伯来原文。就西方而言，可以说他重新发现了这一点。就连像\Person{奥斯定}这样有学识的人也已不再关注\WorkTitle{七十贤士译本}之外的东西，并且对\Person{热罗尼莫}大胆拒绝其历时悠久却并不准确的译法多少感到惊愕。\par

这封信也展示了那种一向标志\Person{热罗尼莫}工作的判断独立性。在撰写此信时，他深受\Person{奥力根}的影响。但是，尽管他对这位大师极为钦佩，当他认为其解经有误时——如在色辣芬的例子中——他并不害怕放弃它。\par

\SourceRecord{Translated by W.H. Fremantle, G. Lewis and W.G. Martley.；Nicene and Post-Nicene Fathers, Second Series；Vol. 6.；Edited by Philip Schaff and Henry Wace.；Buffalo, NY: Christian Literature Publishing Co.,；1893.；<http://www.newadvent.org/fathers/3001018.htm>.}

% structure: p0001 source=h1 target=section reason=page-title

\section{第十九封信}

\Emph{来自教宗达玛稣}\par

\EditorialNote{达玛稣写给耶柔米的一封信，询问有关\Quote{贺三纳}一词的解释（公元383年）。}\par

\SourceRecord{Translated by W.H. Fremantle, G. Lewis and W.G. Martley.；Nicene and Post-Nicene Fathers, Second Series；Vol. 6.；Edited by Philip Schaff and Henry Wace.；Buffalo, NY: Christian Literature Publishing Co.,；1893.；<http://www.newadvent.org/fathers/3001019.htm>.}

% structure: p0001 source=h1 target=section reason=page-title

\section{第二十书}

\Emph{致教宗\Person{达玛苏斯}}\par

\BibleRef{咏 17:25}\par

\SourceRecord{Translated by W.H. Fremantle, G. Lewis and W.G. Martley.；Nicene and Post-Nicene Fathers, Second Series；Vol. 6.；Edited by Philip Schaff and Henry Wace.；Buffalo, NY: Christian Literature Publishing Co.,；1893.；<http://www.newadvent.org/fathers/3001020.htm>.}

% structure: p0001 source=h1 target=section reason=page-title

\section{书信21}

\Emph{致\Person{达玛苏}}\par

\EditorialNote{编者注：在这封信中，\Person{热罗尼莫}应\Person{达玛苏}的请求，详细解释了荡子回头比喻的细节。}\par

\SourceRecord{Translated by W.H. Fremantle, G. Lewis and W.G. Martley.；Nicene and Post-Nicene Fathers, Second Series；Vol. 6.；Edited by Philip Schaff and Henry Wace.；Buffalo, NY: Christian Literature Publishing Co.,；1893.；<http://www.newadvent.org/fathers/3001021.htm>.}

% structure: p0001 source=h1 target=section reason=page-title

\section{第二十二封信}

\Emph{致\Person{Eustochium}}\par

\EditorialNote{这或许是所有书信中最著名的一封。在此信中，热罗尼莫详细阐述了（1）献身于童贞生活者应有的动机，以及（2）他们应遵循的日常行为规则。这封信生动描绘了当时的罗马社会——男女中普遍存在的奢华、放荡和虚伪，此外还有一些生动的自传性细节（§§7, 30），并以当时埃及所实践的三种隐修方式的完整描述作为结尾（§§34-36）。三十年后，热罗尼莫写了一封相似的信给德默特利亚（第一三〇封），可与之比较。写于公元384年罗马。}\par

1. \Quote{女儿，请听，请看，请侧耳；请忘记你的人民和你父的家，君王就会恋慕你的美貌。} 在这第四十四篇圣咏中，天主对人类的灵魂说话，要它效法亚巴郎的榜样，离开自己的土地和家族，离开加色丁人（即魔鬼），居住于活人之地，先知在另一处为此叹息说：\Quote{我希望能看到上主的美善于活人之地。} 但这还不够，除非你忘记你的人民和你父的家，除非你藐视肉身，紧拥新郎。他说：\Quote{不要回头看，} 他说，\Quote{也不要在平原上停留；逃到山上去，免得你被消灭。} \BibleRef{创 19:17} 手扶着犁的，不可向后看 \BibleRef{路 9:62} 或从田地回家，或拥有基督的外衣，从屋顶下来取别的衣服。\BibleRef{玛 24:17} 真是稀奇，父亲命令女儿不要想念她的父亲。\Quote{你们是出于你们的父亲魔鬼，你们父的私欲你们偏要行。} 这是对犹太人说的。又另一处说，\Quote{凡犯罪的，是出于魔鬼。} \BibleRef{若一 3:8} 我们最初生于这样的父母，自然是黑的，即使我们悔改，只要还没有攀登美德的高峰，仍可以说：\Quote{我虽黑却秀美，耶路撒冷的众女子啊。} \BibleRef{歌 1:5} 但你会对我说，\Quote{我离开了童年的家；我忘记了父亲，我在基督里重生了。为此我得到什么报酬？} 上下文表明——\Quote{君王就会恋慕你的美貌。} 这就是伟大的奥秘。\Quote{为此，人要离开父亲和母亲，与妻子结合，两人成为} 不是如那里所说的 \InnerQuote{一体，} 而是 \InnerQuote{一灵。} \BibleRef{弗 5:31} 你的新郎并不傲慢或轻蔑；祂\Quote{娶了位埃塞俄比亚女子。} \BibleRef{户 12:1} 一旦你渴望真正的撒罗满的智慧，来到祂面前，祂会向你宣告祂的全部知识；祂会以君王的右手带你进入祂的内室；\BibleRef{歌 1:4} 祂会奇迹般地改变你的肤色，以至于人们会说：\Quote{那由旷野上来，被美白的是谁？}\par

2. 我这样写信给你，厄斯托丘姆女士（我必须称我主的净配为\Quote{女士}），为通过我的开场白向你表明，我的目的不是赞扬你所遵循并已证明其价值的童贞，也不是数算婚姻的弊端，如怀孕、婴儿的啼哭、情敌带来的折磨、家务管理的忧虑，以及那些最终被死亡切断的虚幻福乐。并非说已婚者因此被排除在外；她们有自己的位置，婚姻应受尊重，床第应是无玷的。\BibleRef{希 13:4} 我的目的是向你表明，你正在逃离索多玛，并应以罗特妻子的遭遇为警戒。\BibleRef{创 19:26} 我可以告诉你，这些篇章中没有奉承。奉承者的言辞悦耳，但尽管如此，他是敌人。你不必期待那些把你置于天使之中、将世界置于你脚下的华丽修辞，它们虽颂扬童贞为有福，却使人骄傲。\par

3. 我愿你从你的隐修誓愿中汲取的，不是骄傲，而是敬畏。\BibleRef{罗 11:20} 你行走时满载黄金；你必须避开强盗的路。对我们人来说，此生是一场跑道：我们在此竞赛，却在他处得冠。当蛇蝎挡路时，无人能放下恐惧。主说：\Quote{我的剑在天上已饱饮，} 你还指望在地上得享平安吗？不，地上只长出荆棘和蒺藜，尘土是蛇的食物。\Quote{因为我们战斗不是对抗血和肉，而是对抗率领者，对抗掌权者，对抗这黑暗世界的霸主，对抗天界中邪恶的鬼神。} 我们被敌军包围，敌人四面环伺。软弱的肉身即将成为灰烬：以一对多，与巨大的势力作战。除非它已被解散，除非这世界的元首已来临并在其中找不到罪，否则你才能安全地聆听先知的话：\Quote{你不必害怕黑夜的惊惶，或是白日飞行的箭矢；或是黑暗中困扰你的灾祸，或是正午的恶魔及其攻击。千人倒在你身边，万人倒在你右边，但灾祸不会临近你。} 当敌军围困你，当你身体发热、情欲激动，当你心中说：\Quote{我该怎么办？} 厄里叟的话会给你答案：\Quote{不要害怕，因为我们这边的比他们那边的更多。} \BibleRef{列下 6:16} 他祈祷：\Quote{主啊，请开你婢女的眼睛，使她看见。} 然后，当你的眼睛被打开时，你将看到一辆火战车，如同厄里亚的，等着接你上天，并要喜乐地歌唱：\Quote{我们的灵魂像鸟一样，从捕鸟人的罗网中逃脱；罗网破碎，我们便逃脱了。}\par

4. 只要我们被这脆弱的身体所牵累，只要我们把这宝贝放在瓦器里，\BibleRef{格后 4:7}只要肉性贪求相反圣神，圣神也相反肉性，\BibleRef{迦 5:17}就不会有确定的胜利。\BibleQuote{我们的仇敌魔鬼如同咆哮的狮子巡行，寻找可吞食的人。}\BibleRef{伯前 5:8}\BibleQuote{你使黑暗产生，}达味说，\BibleQuote{夜便来临，林中的百兽都爬出来。幼狮咆哮觅食，向天主寻求食物。}魔鬼寻找的不是不信者，不是外面的人，亚述王曾将他们的肉烤在炉中。\BibleRef{耶 29:22}他所\Quote{急忙抢夺}的是基督的教会。按哈巴谷书，\BibleQuote{他的食物是精选的。}约伯成为他阴谋的牺牲品，在吞食了犹大之后，他又寻求力量来筛（其他）宗徒。\BibleRef{路 22:31}救主来不是为地上送和平，而是送刀剑。\BibleRef{玛 10:34}\OriginalTerm{路济弗尔}{Lucifer}坠落，那清晨升起的路济弗尔；\BibleRef{依 14:12}那在乐园之乐中长大的，受了应得的判决：\BibleQuote{你虽升得高如鹰，且将巢安置在星宿之间，我也必从那里拉你下来——主说。}\BibleRef{北 1:4}因为他曾在心里说：\BibleQuote{我要高举我的宝座在\OriginalTerm{天主的众星}{stars of God}之上，}并\BibleQuote{我要与\OriginalTerm{至高者}{the Most High}相似。}\BibleRef{依 14:13}因此，天主每日对那沿着雅各伯在梦中所见梯子\BibleRef{创 28:12}下降的天使说：\BibleQuote{我曾说你们是\OriginalTerm{众神}{Gods}，你们都是至高者的儿子。然而你们要像人一样死亡，像任何君王一样陨落。}魔鬼先坠落，既然\BibleQuote{天主在众神会中站立，在众神中施行审判，}宗徒就对那些不再为神的人写道：\BibleQuote{你们中间有嫉妒、纷争，岂不仍是属肉体的、按照人的样子行事吗？}\BibleRef{格前 3:3}\par

5. 既然如此，宗徒——他原是\OriginalTerm{蒙选的器皿}{chosen vessel}\BibleRef{宗 9:15}，被分别出来归于基督的福音，\BibleRef{迦 1:15}——因肉身的刺和恶欲的诱惑，而克制自己的身体，使它服从，免得他给别人传了福音，自己反而被\OriginalTerm{弃绝}{castaway};\BibleRef{格前 9:27}然而尽管如此，他仍看见另一个法律在肢体中交战，反抗他心神的法律，掳他去顺从罪恶的法律；\BibleRef{罗 7:23}如果他在赤身、禁食、饥饿、监禁、鞭打及其他磨难之后，回顾自身而呼喊：\BibleQuote{我这个人真不幸呀！谁能救我脱离这该死的身体呢？}\BibleRef{罗 7:24}——你以为你可以放下警惕吗？要小心，免得天主有一天对你说：\BibleQuote{以色列的\OriginalTerm{贞女}{virgin}已经跌倒，没有人扶起她。}\BibleRef{亚 5:2}我要大胆地说：虽然天主能做一切，但一位贞女一旦跌倒，祂不能使她重新站立。祂或许可以解除玷污者因罪的惩罚，但不会给她冠冕。我们该害怕，免得在我们身上也应验了先知的话：\BibleQuote{好贞女必昏厥。}\BibleRef{亚 8:13}请注意，这里说的是好贞女，因为也有坏的。\BibleQuote{凡注视妇女，}主说，\BibleQuote{有意贪恋她的，就已经在心里与她犯了奸淫。}\BibleRef{玛 5:28}因此，\OriginalTerm{童贞}{virginity}甚至能因一个念头而失去。这就是恶贞女——肉身是贞女，心神却不是；愚昧的贞女，没有油，被\OriginalTerm{新郎}{Bridegroom}关在门外。\par

6. 但如果连那些真正的贞女，如果她们有其他过失，也不能因肉身的童贞而得救，那么那些将基督的肢体卖淫，把圣神的宫殿变为妓院的人，将要怎样呢？她们立刻要听到这样的话：\BibleQuote{巴比伦贞女，下来坐在尘土中吧！加色丁处女，坐在地上吧！因为不再有人称你为娇嫩柔弱的了。你要推磨磨面，揭去面纱，露现大腿，涉水过河！你的裸体必被暴露，你的羞耻必被看见。}\BibleRef{依 47:1} 难道她在成为天主圣子的洞房、领受那既是亲属又是新郎的亲吻之后，还要落到这个地步吗？是的，就是那预言曾歌唱过的：\BibleQuote{王后身穿金饰并绣有各种颜色的衣裳，站在你的右边。}\BibleQuote{她将赤裸，她的裙边要揭到她的脸上。}\BibleRef{耶 13:26} \BibleQuote{她要坐在孤独的水边，放下她的水罐；向每一个路过的人张开她的脚，污秽直到头顶。}\BibleRef{则 16:25} 她当初若是顺从婚姻的轭，走在平坦的地方，总比这样，向往更高的境界，却跌入地狱的深渊要好。我恳求你，不要让忠信之城\Place{熙雍}变成妓女：\BibleRef{依 1:21} 不要让那曾招待过天主圣三的地方，成为魔鬼跳舞、猫头鹰筑巢、豺狼造穴之处。我们不要松开束胸的带子。当情欲搔着感官，淫欲的柔软火焰向我们洒下令人愉悦的光辉时，我们应当立刻站出来呐喊：\BibleQuote{主在我身边，我不怕血肉之躯能对我做什么。}当内心的人在善恶之间一时摇摆不定时，要说：\BibleQuote{我的心，你为什么忧闷？为什么在我内烦躁？期望天主吧！因为我还要赞颂祂，祂是我脸上的救援，是我的天主。}你绝不可让邪恶的念头在你身上滋长，也不可让纷乱的喧嚣在你胸中得势。在敌人还小的时候就杀死它；这样你就不会收获莠子，要在萌芽时掐断邪恶。要谨记圣咏作者的警告：\BibleQuote{巴比伦的女儿，你该受的灾祸，凡报复你对我们所行的一切的，是有福的。凡抓起你的婴儿在磐石上摔碎的，是有福的。}因为自然的热力必然在人身上点燃肉欲的火焰，那人在邪恶的念头在心中升起时，不予以宽容，而是立刻把它们摔在磐石上，这样的人被称赞为有福的。\BibleQuote{那磐石就是基督。}\BibleRef{格前 10:4}\par

7. 多少次，当我住在旷野中，在那为独修者提供野蛮居所的无边孤寂里，被炽热的太阳烤晒时，我多少次幻想自己置身于\Place{罗马}的欢乐之中！我常独自坐着，内心充满苦涩。苦衣毁损了我畸形的四肢，我的皮肤因长期不洗变得像埃塞俄比亚人一样黑。眼泪和呻吟是我的日常；若困倦偶然战胜了我的挣扎，我那几乎散架的枯骨便咯咯地撞在地上。至于我的饮食，不必提了：因为即使在病中，独修者也只有冷水，而吃煮熟的食物被视为放纵。如今，尽管因畏惧地狱，我把自己囚禁在这只有蝎子和野兽为伴的牢狱中，我却常发现自己置身于群女之中。我的脸苍白，身体因禁食而冰冷；然而我的心却燃烧着欲望，情欲的火焰在我面前不断沸腾，尽管我的肉体近乎死亡。无奈之下，我扑倒在\Person{耶稣}脚前，用眼泪洗祂的脚，用头发擦干；然后以数周的斋戒制服我叛逆的身体。我不耻于承认我卑微的痛苦；我更哀叹自己不再是过去的我。我记得我常常整夜大声哭喊直到天明，不停地捶胸，直到\Person{主}斥责后安宁回归。我甚至害怕我的小室，仿佛它知道我的思想；我严厉而愤怒地对自己，独自走入旷野。无论我见到幽谷、峭壁、陡崖，那里便成了我的祈祷所，我可怜肉体的改正所。那里，主自己是我的见证——当我流下很多眼泪，举目望天时，我有时感到自己置身于天使的行列中，满怀喜悦高兴地歌唱：\BibleQuote{因你芬芳的香气，我们追随你。}\BibleRef{歌 1:3}\par

8. 如今，既然那些因守斋而身体消瘦的人，其所受的诱惑仅是邪恶的念头，那么，一个生活在奢侈安逸环境中的少女，又当如何呢？诚然，用宗徒的话来说：\BibleQuote{她活着，却已死了。} \BibleRef{弟前 5:6} 因此，如果经验使我有权提出劝告，或使我的话具有可信度，我首先要以\Person{基督}的净配的身份敦促并警告你：要像躲避毒药一样躲避酒。因为酒是魔鬼攻击年轻人时最先使用的武器。贪念不会如此动摇人，骄傲不会如此膨胀人，野心不会如此迷惑人。其他的恶习我们容易逃避，但这个仇敌却封闭在我们内，无论我们到哪里，都将它带在身上。酒与青春两者共同点燃了感官享乐之火。我们为何要往火焰上加油？为何要为一具已经燃烧的痛苦身体增添新的燃料？的确，\Person{保禄}对\Person{弟茂德}说：\BibleQuote{不要再只喝水，要为胃和屡次生病的关系，用一点酒。} \BibleRef{弟前 5:23} 但请注意允许饮酒的理由：是为了治疗胃痛和屡次的疾病。而且，免得我们因自己的毛病而过分放纵，他命令只可少量饮用；他以一位医师而非宗徒的身份提出劝告（尽管宗徒本身就是一位属神的医师）。他显然担心\Person{弟茂德}可能因软弱而屈服，无法胜任传扬福音所需的频繁奔波。此外，他记得他曾论到\BibleQuote{酒中有放纵，} \BibleRef{弗 5:18} 并说过：\BibleQuote{不吃肉、不喝酒，那是好的。} \BibleRef{罗 14:21} \Person{诺厄}喝了酒，醉了；但他在洪水后那个野蛮时代生活，那时葡萄首次被种植，或许他不知道酒能使人醉。为让你看到圣经隐藏含义的圆满（因为天主的话是珍珠，可从四面穿透），在他醉酒后，他的身体被暴露；放纵导致情欲。\BibleRef{创 9:20} 首先是填满肚腹，然后其他肢体被挑动。同样，后来时期，\BibleQuote{百姓坐下吃喝，起来玩耍。} \BibleRef{出 32:6} \Person{罗特}，天主的朋友，被祂拯救到山上，那千万人中唯一被找到的义人，也被他的女儿们灌醉。虽然她们这样做更多是出于渴望后裔而非喜爱罪恶之乐（因为人类似乎面临灭绝的危险），但她们很清楚，这位义人除非醉酒，否则不会支持她们的计谋。事实上，他不知道自己在做什么，他的罪也不是故意的。但他的错误是严重的，因为他成了\Person{摩阿布}和\Person{阿孟}的父亲，\BibleRef{创 19:30} 以色列的敌人，关于他们，经上说：\Quote{直到第十四代，他们永远不得进入上主的会众。}\par

9. 当\Person{厄里亚}躲避\Person{依则贝尔}，疲惫凄凉地躺在橡树下时，来了一位天使，扶起他说：\BibleQuote{起来，吃吧！} 他一看，头旁有一块饼和一瓶水。\BibleRef{列上 19:4} 如果天主愿意，难道祂不能给祂的先知送来香料酒、美味佳肴和嫩滑的肉吗？当\Person{厄里叟}邀请先知弟子们吃饭时，他只给他们野菜吃；当所有人都同声喊道：\BibleQuote{锅里有死亡！} 时，天主的人并没有对厨师生气（因为他并不习惯于极其丰盛的饮食），而是叫人拿来麦面，撒在锅里，以属神的力量使苦汤变甜，\BibleRef{列下 4:38} 如同\Person{梅瑟}曾使\Place{玛辣}的水变甜一样。\BibleRef{出 15:23} 再者，当有人被派去捉拿先知，并遭受肉体与心智的盲目打击，以致\Person{厄里叟}在他们不知不觉中把他们领到\Place{撒玛黎雅}时，请注意\Person{厄里叟}命令用什么食物来款待他们。他说：\BibleQuote{在他们面前放下面饼和水，叫他们吃喝，然后回到他们的主人那里去。} \BibleRef{列下 6:18} 至于\Person{达尼尔}，他本可享用国王桌上的美食，\BibleRef{达 1:8} 却宁愿选择由\Person{哈巴谷}带给他的割麦人的早餐，那必定只是乡间食物。他被称为\BibleQuote{极可爱的人}，因为他不愿吃欲望的饼，也不愿喝贪欲的酒。\par

10. 在圣经中，有无数神圣的答复谴责贪饕，赞同简朴的食物。但由于守斋并非我现下的主题，且对它的充分讨论需要单独一篇论文，因此，从这主题所提示的众多要点中，这些简短的观察已足够。由此你将明白，为什么第一个人服从他的肚腹而不服从天主，从乐园被贬入这涕泣之谷；为什么\Person{撒殚}在旷野中用饥饿诱惑主自己；\BibleRef{玛 4:2} 为什么宗徒呼喊说：\BibleQuote{食物是为肚腹，肚腹是为食物，但天主必要使这两样都毁灭；} \BibleRef{格前 6:13} 以及为什么他论到放纵自己的人说：\BibleQuote{他们的天主是肚腹。} \BibleRef{斐 3:19} 因为人总是崇拜自己最喜欢的东西。因此，必须注意，禁欲能使那些曾被餍足驱逐出乐园的人重返乐园。\par

11. 你或许会告诉我，你出身高贵，在奢侈中长大，习惯于柔软的床榻，你离不开酒和美味，会觉得更严格的修规难以忍受。如果是这样，我只能说：\Quote{那么，就按你自己的规矩生活吧，因为天主的规矩对你来说太难了。} 不是说创造主和万有的主喜欢腹中空荡咕响或肺腑发热；而是说这些是保持贞洁所不可或缺的手段。\Person{约伯}是天主所喜爱的人，在他面前完善而正直；\BibleRef{约 2:3} 但请听他对魔鬼所说的话：\BibleQuote{他的力量在于腰，他的气力在于脐。}\par

术语选择是为了文雅，但指的是两性的生殖器官。因此，\Person{达味}的后裔——按应许要坐在他宝座上的那位——被说成是从他的腰中出来的。从\Person{雅各伯}进入埃及的七十五灵魂被说成是从他的大腿中出来的。\BibleRef{创 46:26}同样，当主与他摔跤后他的大腿萎缩了，\BibleRef{创 32:24}他便不再生育。以色列子民又被告知要束腰刻苦地庆祝逾越节。\BibleRef{出 12:11}天主对\Person{约伯}说：\BibleQuote{你要像勇士束好腰。}\BibleRef{约伯 38:3}\Person{若翰}束着皮带。\BibleRef{玛 3:4}宗徒必须束腰，才能举起福音的灯。\BibleRef{路 12:35}当\Person{厄则克耳}告诉我们\Place{耶路撒冷}在旷野中被发现，满身血污时，他说：\BibleQuote{你的肚脐也没有剪断。}\BibleRef{则 16:4}因此，魔鬼攻击男人时，力量在腰；攻击女人时，力量在肚脐。\par

12. 你希望为我的主张找到证据吗？请看实例。\Person{参孙}比狮子更勇敢，比岩石更坚强；孤身无援地追击一千个武装的人；然而，在\Person{德里拉}的怀抱中，他的决心融化了。\Person{达味}是合天主心意的人，他的嘴唇曾多次歌颂圣者，未来的基督；然而当他在屋顶上散步时，却被\Person{巴特舍巴}的裸体所迷惑，在奸淫之上又加了杀人。\EditorialNote{撒慕尔纪下第十一章}请注意，一个人即使在自己家中，也不能无危险地使用眼睛。然后他悔改了，对主说：\Quote{我得罪了祢，惟独得罪了祢，在祢眼前行了这恶。}作为君王，他不惧怕任何人。\Person{撒罗满}也是如此。智慧用他来歌唱她的赞美，他讲论了所有植物，\BibleQuote{从\Place{黎巴嫩}的香柏木到墙上长出的牛膝草；}\BibleRef{列上 4:33}然而他因喜爱女色而离开了天主。\BibleRef{列上 11:1}而且，仿佛为了表明近亲关系并非保障，\Person{阿默农}对他的妹妹\Person{塔玛尔}燃起了非法的情欲。\EditorialNote{撒慕尔纪下第十三章}\par

13. 我无法言说每天有多少贞女堕落并失去她们的母亲——教会的怀抱：那些骄傲的仇敌在上面设立宝座的星辰，\BibleRef{依 14:13}以及被蛇挖空以便盘踞其裂缝的岩石。你会看见许多未婚先寡的妇人，试图用撒谎的服装掩盖她们悲惨的堕落。除非她们被隆起的腹部或婴儿的啼哭所出卖，她们迈着轻快的步伐、昂首阔步地行走。有些人甚至服药以确保不育，从而在孕育之前就谋杀了人类。有些人在因罪怀孕后使用药物堕胎，当她们（常常如此）与胎儿一同死亡时，她们进入阴间，不仅背着对基督犯奸淫的罪责，还有自杀和杀婴的罪。然而正是这些人说：\Quote{\InnerQuote{在纯洁的人看来，一切都是纯洁的；}\BibleRef{铎 1:15}我的良心足以引导我。纯洁的心是天主所寻求的。我为什么要戒避天主所造、可感恩领受的食物呢？}\BibleRef{弟前 4:3}而且，当她们想显得可爱可亲时，她们先灌醉自己，然后将最粗俗的亵渎与醉酒结合，说：\Quote{我绝不戒避基督的血。}当她们看见别人脸色苍白或忧伤时，便称她为\Quote{贱人}或\Quote{摩尼教徒}；这的确很合乎逻辑，因为按照她们的原则，禁食就是异端。她们外出时尽力引人注目，以点头眨眼鼓动一群年轻人跟随她们。对她们每一个人，先知的话是真实的：\BibleQuote{你有娼妓的额头，你拒绝羞耻。}\BibleRef{耶 3:3}她们的袍子只有一条窄紫条纹，这是真的；她们的头饰有些松散，以便让头发自由垂落。肩上飘着淡紫色披风，她们称之为\Quote{ma-forte}；脚穿廉价拖鞋，手臂紧紧裹在袖子里。加上她们职业标志般的轻浮步伐，你就拥有了她们所拥有的一切贞洁。这样的人可能有自己的赞美者，并且在毁灭的市场上卖出更高的价钱，仅仅因为她们被称为贞女。但对于这样的贞女，我宁愿令她们不悦。\par

14. 我羞于谈论此事，它如此令人震惊；然而虽令人难过，却是真实的。教会的这种\OriginalTerm{爱侣}{agapetæ}之灾从何而来？这些未嫁之妻、这些新奇的妾、这些妓女——我这样称呼她们，尽管她们依附于一个伴侣——从何而来？同一所房子容纳她们，同一个房间。她们常常同睡一张床，却指责我们疑心不正。一个兄弟离开他的贞女姐妹；一个贞女，轻视她未婚的兄弟，在陌生人中寻求一个兄弟。两者都声称只有一个目的，从非亲属那里寻求灵性安慰；但他们的真实目的是沉溺于性交。正是这样一些人，\Person{撒罗满}在\WorkTitle{箴言}中加以嘲讽：\BibleQuote{人能否将火藏在怀中，而衣服不烧着呢？人能否在火炭上行走，而脚不灼伤呢？}\BibleRef{箴 6:27}\par

15. 那么，我们将那些只愿显为贞女、而不愿实为贞女的人逐出并赶离我们的视线。从此，我可以将我的话全部转向你；因为你是\Place{罗马}第一位出身名门的贞女，所以你更须努力，以免失去现在与将来的美善。你至少已从你自己家庭中的事例学到了婚姻生活的烦恼与婚姻的不确定性。你的姐姐\Person{布来西拉}，年纪比你大，但在拒绝贞洁誓愿方面却落在你后面；她结婚仅七个月便成了寡妇。我们这些不知道未来的凡人真是可悲！她同时失去了贞女的荣冠与婚姻的乐趣。而且，虽然作为寡妇，她拥有第二等级的贞洁，你能想象她必须忍受的无尽痛苦吗？她每天在姐姐身上看到自己所失去的，并且发现没有婚姻的乐趣很艰难，而当前的节制却得到较小的赏报？然而，她也可以振作并欢喜。百倍的果实与六十倍的果实皆出自同一种子，那种子就是贞洁。\BibleRef{玛 13:8}\par

16. 不要追求与已婚妇女交往，也不要访问高门大户。不要频频回望你为了成为贞女而鄙视的那种生活。你知道，世俗妇女以她们的丈夫在法庭或其他高位上自夸。皇后的门口总是挤满了争相拜访的人群。那么，你为什么使你的丈夫蒙羞呢？你为什么，天主的净配，急忙去访问一个凡人的妻子呢？在这方面，你要学习一种神圣的骄傲；要知道你比她们更优越。不仅必须避免与那些因丈夫的荣耀而自大、被太监队伍护卫、穿着金线织成的袍子的人交往，你还必须避开那些出于必要而非自愿成为寡妇的人。并不是说她们应当盼望丈夫死亡，而是说当节制的机会来到时，她们没有欢迎它。事实上，她们只更换了服装，旧有的自私仍然未改。看到她们坐在宽敞的轿子里，穿着红斗篷，身体丰满，一排太监在她们前面行走，你会以为她们不是失去了丈夫，而是在寻找丈夫。她们的家中充满了谄媚者和宾客。就连那些本应用他们的教导和权威使她们尊敬的神职人员，也亲吻这些贵妇的额头，并伸出手（如果你不了解，还以为他们在祝福），为拜访收取报酬。同时，她们看到神父们离不开她们，便骄傲起来；并且由于对两者都有经验，她们更喜欢寡妇的自由而不是婚姻的约束，便自称为贞洁生活的女子和修女。在过分丰盛的晚餐后，她们就去休息，做梦梦到宗徒们。\par

17. 让你的同伴是那些因禁食而面色苍白、身体瘦削，并被她们的年龄和品行所认可的女子；她们每天在内心歌唱：\Quote{请告诉我，你在何处牧放你的羊群，何处使它们歇息在午间？}并真诚地说：\Quote{我渴望离去，与基督同在。}\BibleRef{斐 1:23}要顺从你的父母，效法你的净配的榜样。\BibleRef{路 2:51}很少外出，如果你希望寻求殉道者的帮助，就在你自己的内室中寻求。因为你若只在需要时才出门，你就永远不会需要外出的借口。饮食要有节制，切勿使胃负担过重。因为许多女子，虽然对酒有节制，却在食物上放荡。当你夜间起来祈祷时，让你的气息来自空虚而非饱足的胃。常阅读，尽你所能学习。让睡意在你手中还握着书卷时征服你；让你的头垂在神圣的书页上。让你的禁食成为日常，你的饮食要避免饱足。空着肚子是无用的，如果两三天后饥饿要用饱食来弥补。过饱时，心立即变得迟钝，当土地被浇灌时，就会生出情欲的荆棘。如果你感到外在之人叹息着渴望青春的花朵，并且如果你在饭后躺在榻上被诱人的感官欲望所激动，那么就拿起信德的盾牌，因为只有它能熄灭魔鬼的火箭。\BibleRef{弗 6:16}先知说：\Quote{他们都是奸夫；他们使自己的心像烤炉一样。}但你要紧紧跟随基督的足迹，专注于祂的话语，说：\Quote{当耶稣在路上给我们讲解圣经的时候，我们的心不是火热的吗？}\BibleRef{路 24:32}又说：\Quote{你的言语极其精炼，你的仆人喜爱它。}人的灵魂很难避免爱慕什么，我们的心思必须让位于某种情感。肉身的爱被精神的爱战胜。欲望被欲望熄灭。从一方拿走的，增加了另一方。因此，当你躺在榻上时，反复说：\Quote{夜间，我寻找我灵魂所爱的。}\BibleRef{歌 3:1}宗徒说：\Quote{所以，你们要致死属于地上的一切肢体。}\BibleRef{哥 3:5}因为他自己这样做了，后来便可以自信地说：\Quote{我生活，但不是我生活，而是基督在我内生活。}\BibleRef{迦 2:20}凡是致死自己肢体，并感到自己行走在虚影中的人，敢于说：\Quote{我成了霜冻中的皮囊。凡在我内情欲的湿气都被晒干了。}又说：\Quote{我的双膝因禁食而软弱，我忘了吃我的面包。因我呻吟的声音，我的骨头紧贴我的皮肤。}\par

18. 要像蚱蜢一样，使夜晚充满音乐。每夜以眼泪洗涤你的床铺，浸湿你的卧榻。要警醒，像屋顶上独居的麻雀。用精神歌唱，也要用理智歌唱。\BibleRef{格前 14:15} 愿你的歌声如圣咏作者一般：\BibleQuote{我的灵魂，请赞美上主；不要忘记祂的一切恩惠。祂赦免你所有的过犯，治疗你的一切疾病，救赎你的性命脱离毁灭。} 我们中有谁能诚实地将他的话当作自己的话：\BibleQuote{我以灰烬当饭，以眼泪拌饮料？} 然而，当蛇引诱我们，如同引诱我们的原祖父母，去吃禁果，并且将我们从贞洁的乐园驱逐后，又想给我们穿上皮外套，就像厄里亚重返乐园时留在地上的那件，我们岂不应哭泣哀叹？\BibleRef{列下 2:13} 你要对自己说：\Quote{我与那转瞬即逝的感官享乐有何相干？我与那塞壬的歌声有何相干，它如此甜美，却给听者带来致命的结局？} 我不愿你受那判决，即人所受的诅咒。当天主对厄娃说：\BibleQuote{你要在痛苦中生产儿女，} 你要对自己说：\Quote{这是为已婚妇女的律法，不是为我。} 当祂接着说：\BibleQuote{你要依恋你的丈夫，} \BibleRef{创 3:16} 你再说：\Quote{让那没有基督作丈夫的妻子依恋她的丈夫吧。} 最后，当祂说：\BibleQuote{你必会死亡，} \BibleRef{创 2:17} 你再说：\Quote{婚姻固然终将死亡；但我所选择的生活不受性别限制。让那些已婚妇女保持她们应有的位置和时间吧。对我来说，贞洁在玛利亚和基督身上被祝圣了。}\par

19. 有人会说：\Quote{你胆敢贬损婚姻吗？婚姻是天主祝福的状态。} 我并未贬损婚姻，只是将贞洁置于它之上。无人拿坏东西与好东西相比。已婚妇女可以庆幸自己仅次于贞女。天主说：\BibleQuote{你们要生育繁殖，充满大地。} \BibleRef{创 1:28} 愿意充满大地的人，可以生育繁殖，如果他愿意。但你所归属的行列不在世上，而在天上。生育繁殖的命令是在被逐出乐园之后，在裸体和象征情欲的无花果叶之后才首次实现的。让那些汗流满面才得糊口的人，那些土地给他们长出荆棘和蒺藜，\BibleRef{创 3:18} 庄稼被荆棘吞噬的人，让他们嫁娶吧。我的种子却结出百倍的果实。\BibleQuote{这话不是人人所能领悟的，只有蒙受恩赐的人才能领悟。}\par

有些人可能是出于无奈而成为阉人，而我是出于自由意志。\BibleRef{玛 19:11} \BibleQuote{拥抱有时，不拥抱有时。抛掷石头有时，堆聚石头有时。} \BibleRef{训 3:5} 既然天主从外邦人的硬石中为亚巴郎兴起了子孙，\BibleRef{玛 3:9} 他们便开始成为\Quote{在地上滚动的圣石}。他们经过世界的旋风，乘着天主的战车在飞快的轮子上滚动。让那些丢失了从上到下浑然一体的长衣的人为自己缝制外套吧；\BibleRef{若 19:23} 他们喜爱婴儿的哭声——这些婴儿一见光明，便哀叹自己出生。在乐园里，厄娃是贞女，只有在穿上皮外套之后，她才开始了婚姻生活。现在，乐园也是你的家园。因此，你要保持你的长子权，并说：\Quote{我的灵魂，请回到你的安息之所。} 为表明贞洁是本性的，而婚姻只随罪恶而来，凡婚姻所生的，都是贞洁的肉体，它以果实归还了根部所失去的。\Quote{从叶瑟的树干要生出一根枝条，从他的根上要发出一朵花。} 枝条是主的母亲——简朴、纯洁、无瑕；不从外界汲取生命的种子，却像天主自身一样，因独一而丰产。枝条的花是基督，祂论及自己说：\BibleQuote{我是沙龙的玫瑰花，谷中的百合花。} \BibleRef{歌 2:1} 在另一处，祂被预言为\BibleQuote{一块非人手凿出的石头，} \BibleRef{达 2:45} 先知以此象征祂将由贞女所生，并且自身是贞者。因为在这里，手是婚姻的象征，正如经上所载：\BibleQuote{他的左手在我头下，他的右手拥抱我。} \BibleRef{歌 2:6} 这与以下解释也相符：不洁的动物被成对领进诺厄方舟，而洁净的动物则取奇数。\BibleRef{创 7:2} 同样，当梅瑟和若苏厄奉命脱鞋，因为他们所站之地是圣地，这命令具有神秘含义。同样，当门徒被派去宣讲福音，他们被告知不要带鞋或鞋带；当兵士们来为耶稣的衣物拈阄时，\BibleRef{若 19:23} 他们找不到可拿走的靴子。因为主自己不能拥有祂所禁止仆人拥有的东西。\par

20. 我赞美婚姻，我赞美婚姻，但这是因为婚姻为我产生贞女。我从荆棘中采玫瑰，从泥土中淘黄金，从蚌壳中得珍珠。\Quote{农夫岂是终日撒种？} \BibleRef{依 28:24} 他不也该享受自己劳苦的果实？婚姻越受尊崇，它所生的子女就越被珍爱。母亲，你为什么嫉妒你女儿的童贞？她以你的乳汁为食，从你腹中出生，在你怀中长大。你警醒的爱护使她保持了贞女的身份。你岂因她选择作君王的妻子而非士兵的妻子而恼怒？她为你带来了崇高的特权：你现在是天主的岳母。\Quote{关于贞女，}宗徒说，\Quote{我没有主的命令。} \BibleRef{格前 7:25} 这是为什么？因为他自己的贞洁并非出于命令，而是出于他的自由选择。那些声称他有妻子的人是不足信的；因为当他讨论节欲并赞扬永久的贞洁时，他说：\Quote{我愿众人像我一样。} 接着又说：\Quote{因此，我对未婚者和寡妇说：他们若像我一样，为他便是好的。} \BibleRef{格前 7:7} 在另一处：\Quote{难道我们没有权柄带着妻子同行，如同其余的宗徒一样？} \BibleRef{格前 9:5} 那么，为什么主没有关于贞洁的命令呢？因为自愿献上的胜过强制压榨的；而命令贞洁就等于废除婚姻。这将是一条苛刻的法令，迫使人违抗本性，强求人过天使般的生活；不仅如此，它还会谴责神圣的规定。\par

21. 旧法律对福乐有不同的观念，因为经上说：\Quote{那在熙雍有后裔、在耶路撒冷有家族的人是有福的；}\Quote{不生育的、不生产的，是可诅咒的；}\Quote{你的子女围绕你的桌子，如同橄榄树苗。} 信德的人也得到财富的应许，我们被告知\Quote{在他们的支派中没有一个衰弱的人。} 但现在甚至对阉人也说：\Quote{不要说：看，我是一棵枯树，}\BibleRef{依 56:3} 因为代替儿女，你在天上有了永久的住所。如今穷人有福了，如今拉匝禄被放在穿紫袍的富翁面前。如今软弱的人算为强壮。但在那些日子，世界尚未住满；因此，忽略那些仅作为预表的无子女的事例，只有那些因子女而夸耀的人才被认为是幸福的。这就是为什么亚巴郎在老年娶了刻突辣；\BibleRef{创 25:1} 肋阿用她儿子的曼陀罗雇用了雅各伯；\BibleRef{创 30:14} 而美丽的辣黑耳——教会的预像——抱怨她的子宫的关闭。\BibleRef{创 30:1} 但庄稼逐渐成熟，于是收割者带着他的镰刀被派出来了。厄里亚过贞洁的生活，厄里叟和许多先知弟子也是如此。耶肋米亚得到命令：\Quote{你不可娶妻。}\BibleRef{耶 16:2} 他在母腹中已被祝圣，\BibleRef{耶 1:5} 现在他被禁止娶妻，因为被掳临近。宗徒用不同的词句给出同样的忠告。\Quote{因此我认为，为了现时的困境，这是好的，即人保持现状是好的。} 这个困境是什么，它消除了婚姻的喜乐？宗徒在后面一节告诉我们：\Quote{时候短促：从此以后，那些有妻子的，要像没有一样。}\BibleRef{格前 7:29} 拿步高逼近了。狮子正从洞穴中奋起。如果婚姻最终导致我成为最傲慢的君主的奴隶，那婚姻对我有何益处？如果我的孩子们的命运是先知悲叹的描述：\Quote{吃奶的婴儿的舌头因干渴紧贴上颚；幼童求饼，却无人擘给他们。}\BibleRef{哀 4:4} 那么他们对我有何益处？在那些日子，正如我所说，节德只存在于男人中：厄娃仍然继续怀孕生子。但如今，既然有一位贞女怀孕\BibleRef{依 7:14} 并为我们生了一个孩子，先知论到他说：\Quote{政权必担在他的肩头上，他的名字要称为全能的天主、永远的大父。}\BibleRef{依 9:6} 如今咒诅的锁链被打破了。死亡从厄娃而来，但生命从玛利亚而来。因此，童贞的恩赐极丰富地赐予了女人，因为它从一位女人开始。当天主之子踏上大地时，祂为自己在那里建立了新的家庭；这样，祂在天上受天使朝拜，天使也在世上服事祂。然后，贞洁的友弟德再次砍下哈洛斐诺的头。\EditorialNote{友弟德传第十三章} 然后，哈曼——他的名字含有不义之意——再次在自己燃起的火中被焚烧。\BibleRef{艾 7:10} 然后，雅各伯和若望舍弃了父亲、渔网和船，跟随了救主：亲情、世间的联系、家庭的牵挂都不能留住他们。然后人们听到了这些话：\Quote{谁愿意跟随我，该弃绝自己，背起自己的十字架来跟随我。}\BibleRef{谷 8:34} 因为没有士兵带着妻子上战场。甚至当一个门徒想要埋葬他的父亲时，主禁止了他，说：\Quote{狐狸有洞，天空的飞鸟有巣，但人子却没有枕头的地方。}\BibleRef{玛 8:20} 所以你不要抱怨住所狭窄。同样，宗徒写道：\Quote{未结婚的，挂虑主的事，怎样悦乐主；但结了婚的，挂虑世上的事，怎样悦乐他的妻子。妇女和贞女也有区别。未出嫁的妇女挂虑主的事，为叫身体和灵魂都圣洁；但已出嫁的挂虑世上的事，怎样悦乐她的丈夫。}\BibleRef{格前 7:32}\par

22. 婚姻涉及何等巨大的不便，又围绕多少忧虑，我想，我已在我的论著——那篇反驳海尔维狄乌斯、\WorkTitle{论真福玛利亚终身童贞}的文章——中简短地描述了。现在重复同样的内容会很冗长；谁若愿意，可以从那泉源汲取。但为避免我似乎完全忽略了此事，我现在只说：宗徒命令我们不断祈祷，\BibleRef{得前 5:17} 而在婚姻状态中向妻子尽本分的人，不能如此祈祷。要么我们不断祈祷而保持童贞，要么我们停止祈祷以履行婚姻的义务。但他仍说：\Quote{如果童女结婚，并没有犯罪；不过这样的人在肉身要受苦难。}\BibleRef{格前 7:28} 起初我应许对于婚姻的困扰几乎不说什么，现在我也同样通知你。若你想知道童女免于多少烦恼、妻子被多少枷锁束缚，你应阅读戴尔都良\WorkTitle{致一位哲学朋友}及其论童贞的其他著作，真福西彼廉的崇高著作，教宗达玛苏的诗文著作，以及我们自己的盎博罗削最近为他妹妹写的论著。在这些著作中，他以如此滔滔雄辩倾注了他的灵魂，以至于他寻求、陈述并整理了所有赞美童贞的内容。\par

23. 我们必须走另一条路，因为我们的目的不是赞美童贞，而是保全它。知道它是好事并不够：当我们选择了它，就必须以嫉妒的谨慎守护它。前者仅需要判断，我们与许多人共享；后者需要劳苦，少有人与我们竞争。主说：\Quote{唯有忍耐到底的，必然得救，}\BibleRef{玛 24:13} 又说：\Quote{因为被召的人多，被选的人少。}因此，我在天主和耶稣基督并祂的选民天使面前恳求你，守护你所领受的，不要轻易将主殿宇的器皿（只有司铎有权看见）暴露于公众视线，免得亵渎之人观看天主的圣所。乌匝碰触了不可碰触的约柜，立刻被击倒而死。\BibleRef{撒下 6:6} 诚然，没有金器银器像童女身体的殿宇那样为天主所珍爱。影子已经过去，如今实体已来临。你固然可以全然单纯地说话，出于和蔼礼貌对待最陌生的客人，但不贞洁的眼睛看不正任何事物。它们无法欣赏灵魂之美，只珍视身体之美。希则克雅曾向亚述人展示天主的宝物，\BibleRef{列下 20:12} 那些人本不该看到他们必定贪恋的东西。结果犹大被连续战争撕裂，而这些主殿宇的器皿正是首批被掳到巴比伦的。我们看到贝尔沙匝在他的宴会和妾婢中（恶习总是以玷污崇高为荣）用这些圣杯饮酒。\BibleRef{达 5:1}\par

24. 绝不要侧耳听邪恶的话。因为人们常说不得体的话来试探童女的坚定，看她是否乐意听，是否准备在每一次愚蠢的笑话中放松。这样的人赞同你肯定的，否认你否定的；他们称你不仅神圣，而且有成就，说在你里面没有诡诈。他们说：\Quote{看，基督真正的婢女；看，全心纯朴。那粗糙、丑陋、乡巴佬般的可怖之人多么不同，她很可能从未结婚，因为她找不到丈夫。}我们自然软弱，容易听信这些谄媚者；但尽管我们脸红并回答说这样的赞美超出了我们应得的，我们内里的灵魂却因听到自己被赞美而欢喜。\par

如同盟约之柜，基督的新娘应内外包金\BibleRef{出 25:11}；她应守护主的法律。正如柜中只存放约版\BibleRef{列上 8:9}，同样，在你内不应有任何外来的思想。因为主喜悦居住在你心中，正如昔日祂坐在施恩座和革鲁宾上\BibleRef{出 25:22}。祂曾派遣门徒为祂解开驴驹，好骑乘其上；同样，祂派遣他们使你脱离世俗的挂虑，离开埃及的砖和草，跟随真正的梅瑟穿越旷野，进入恩许之地。不要让任何人阻止你，无论是母亲、姊妹、亲属或兄弟：\Quote{主需要你}\BibleRef{玛 21:1}。如果他们企图阻拦，让他们害怕降临在法郎身上的鞭挞，因为他不同意天主的百姓离去事奉祂\BibleRef{出 7:16}，而遭受了圣经所述的各种灾祸。耶稣进入圣殿，赶出了不属于圣殿的东西。因为天主是忌邪的，祂不愿父的家成为贼窝。那里数点银钱，那里售卖鸽子，那里压抑纯朴——也就是说，处女的胸怀因世俗的挂虑而发热——圣所的帐幔立即裂开\BibleRef{玛 27:51}，新郎忿怒地起身说：\Quote{你们的家将留给你们一片荒凉}\BibleRef{玛 23:38}。阅读福音，看看玛利亚坐在主脚前如何与忙碌的玛尔大对比。玛尔大为了款待主和祂的门徒，焦虑地预备食物；然而耶稣对她说：\Quote{玛尔大，玛尔大，你为许多事操心忙碌。其实需要的惟有一件。玛利亚选择了更好的一份，是不能从她夺去的}\BibleRef{路 10:41}。那么要像玛利亚，宁取灵魂的食粮而非肉身的食粮。让你的姊妹们奔波，设法合适地迎接基督吧。而你，既已一次永远抛下世界的重担，就坐在主脚前说：\Quote{我找到了我灵魂所爱的；我要抱住他，不让他走}\BibleRef{歌 3:4}。祂要回答：\Quote{我的鸽子，我的完人只有一个；她是她母亲的独一者，是生她者的特选者}\BibleRef{歌 6:9}。这里所说的母亲是天上耶路撒冷\BibleRef{迦 4:26}。\par

25. 愿你的内室常保护你；愿新郎常在内里与你嬉戏\BibleRef{创 26:8}。你祈祷吗？你就是在对新郎说话。你阅读吗？祂就在对你说话。当睡意袭来，祂会从门后伸手穿过门孔，你的心将因祂而激动；你会醒来起身说：\Quote{我因爱成疾}。祂便回答：\Quote{我妹子，我新娘，是关锁的园；是封闭的泉，是禁闭的井}\BibleRef{歌 4:12}。\par

不要离家，也不要拜访外邦女子的女儿，即使你有族长作兄弟，以色列为父亲。狄纳出去就受了诱惑\EditorialNote{创世纪第三十四章}。不要在街上寻找新郎；不要绕城拐角行走。因为即使你说：\Quote{我要起来，绕城行走，在街道和广场上寻找我灵魂所爱的}，并问巡城者：\Quote{你们看见我灵魂所爱的吗？}\BibleRef{歌 3:2}也没有人会屈尊回答你。新郎不能在街上找到：\Quote{通往生命的门是窄的，路是狭的}\BibleRef{玛 7:14}。因此雅歌接着说：\Quote{我寻找他，却找不着；我呼唤他，他不回答}。但愿只是找不着而已。你还会受伤、被剥去衣物，哀叹说：\Quote{巡城者遇见我，击打我，伤了我，夺去了我的面纱}\BibleRef{歌 5:7}。如果连那能说\Quote{我睡，我心醒}\BibleRef{歌 5:2}和\Quote{我的良人如同一袋没药，常在我怀中}\BibleRef{歌 1:13}的人，因外出而遭受如此痛苦，那么像我们这些年轻女子，当新娘与新郎进入时仍留在外面的人，又该如何呢？耶稣是忌邪的。祂不愿你的脸被他人看见。你或许为自己辩解说：\Quote{我已拉上面纱，遮住脸面，在那里寻找你说：\InnerQuote{告诉我，我灵魂所爱的，你在何处牧羊，何处使羊群歇午？我为何要在你同伴的羊群旁像一个蒙面的人呢？}}但即使你有这些借口，祂仍会发怒，忿然说：\Quote{你若不认识自己，你这女子中最美丽的，就随着羊群的脚踪去，在牧人帐棚旁牧放你的山羊吧}。你或许美丽，你的面容或许是新郎最喜爱的；然而，除非你认识自己，并全心守护你的心\BibleRef{箴 4:23}，除非你也避开年轻人的目光，你将被逐出我的洞房，去牧放那将被放在左边的山羊\BibleRef{玛 25:33}。\par

26. 事既如此，我的\Person{欧斯托嘉}，女儿、主母、同仆、姊妹——这些称呼，第一个指你的年龄，第二个指你的身份，第三个指你的圣召，最后一个指你在我心中所占的位置——请听依撒意亚的话：\BibleQuote{我的人民，来，进入你的内室，关上你的门，隐藏片刻，等到主的忿怒过去。}\BibleRef{依 26:20} 让那些糊涂的童女在外游荡吧，但你要与新郎一同留在家里；因为如果你关上门，并按照福音的训示，\BibleRef{玛 6:6} 在暗中向你的父祈祷，祂必来敲门，说：\BibleQuote{看，我站在门口敲门；若有人……给我开门，我要进去，与他坐席，他也要与我同席。}\BibleRef{默 3:20} 你必立即热切地回答：\BibleQuote{这是我的爱人的声音在敲门，说：给我开门，我的姊妹，我的爱卿，我的鸽子，我的完人。}你不可能拒绝说：\BibleQuote{我已脱了外衣，怎能再穿上？我已洗了脚，怎能再玷污？}\BibleRef{歌 5:2} 立刻起来开门。否则，当你迟疑时，祂可能已离去，你便要哀伤地说：\BibleQuote{我给我的爱人开了门，但我的爱人已转身走了。}\BibleRef{歌 5:6} 何必要向新郎关闭你心灵的门呢？让它们向基督敞开，却向魔鬼关闭，正如经上所说：\Quote{若有掌权者的灵起来反对你，不要离开你的岗位。}\Person{达尼尔}在他不能再留在楼下而退到的那间楼上，窗户朝向耶路撒冷敞开。你也要保持你的窗户敞开，但只在光能进入、你能看见主的那一面。不要打开那另一类窗户，关于它们先知曾说：\BibleQuote{死亡已从我们的窗户上来。}\BibleRef{耶 9:21}\par

27. 你还必须小心避免虚荣心的陷阱。\Person{耶稣}说：\BibleQuote{你们怎能相信，彼此求光荣呢？}那使人不能相信的恶是何等大啊！让我们反过来说：\BibleQuote{你是我的光荣。}\BibleRef{耶 9:24} 又说：\BibleQuote{夸耀的，应当夸耀主。}\BibleRef{格前 1:31} 又说：\BibleQuote{如果我还求人的欢心，我就不是基督的仆人了。}\BibleRef{迦 1:10} 又说：\BibleQuote{至于我，我只以我们的主耶稣基督的十字架来夸耀，因为借着祂，世界于我已被钉在十字架上，我于世界也被钉在十字架上。}再者：\BibleQuote{我们终日以天主夸耀；我的灵魂要因主夸耀。}当你施舍时，只让天主看见你。当你禁食时，要面带喜色。你的衣着不可过于整洁，也不可过于邋遢；也不可引人注目，以致路人指点你。有弟兄去世了吗？有姊妹的遗体需要安葬吗？当心，不要因过于频繁地履行这些职务而自己死去。不要想显得十分虔诚，或比必要的更为谦卑，以免你因躲避荣耀而求取荣耀。因为许多人，向所有人的目光隐藏他们的贫穷、爱德和禁食，却因藐视这一切而想引起赞美，并奇怪地寻求表扬，同时声称要避开它。至于其他使人欢喜、沮丧、希望和恐惧的干扰因素，我发现许多人已摆脱；但少有不受这个缺点影响的，最好的人乃是其品格像一张纯净的皮肤，被最少的瑕疵所损坏。我认为没有必要警告你不要夸耀你的财富，不要因出身而骄傲，也不要自以为高人一等。我知道你的谦卑；我知道你能真诚地说：\BibleQuote{主，我的心不骄横，我的眼不高傲}；我知道在你和你母亲的胸中，那使魔鬼堕落的骄傲无地容身。若为此写信给你，便是浪费时间；因为没有比教一个学生他已经知道的事更愚蠢的了。但是，既然你已经藐视了世界的虚荣，不要让这个事实激起你新的虚荣。不要暗自想：既然不再以金饰争奇斗艳，或许要以朴素的衣着令人瞩目。当你进入一间挤满弟兄姊妹的房间时，不要坐在太低的位置，或声称你不配用脚凳。不要故意压低声音，仿佛因禁食而疲惫；也不要倚靠在别人的肩上，模仿昏厥者踉跄的步伐。有些妇女确实毁损自己的面容，为叫人看出她们在禁食。\BibleRef{玛 6:16} 她们一见到任何人，就叹息，低垂目光；她们蒙住脸，只留一只眼睛来看。她们的衣裳是暗色的，她们的腰带是苦衣做的，她们的手和脚是脏的；只有她们的胃——看不见——却因食物而发热。每日为她们吟咏圣咏：\BibleQuote{主要使自悦之人骨散。}另有些人改变装束，摆出男子的气派，以自己生来是女性为耻。她们剪短头发，不以像阉人而羞愧。有些人在身上披山羊毛，戴上风帽，以为自己像猫头鹰就能变回孩童。\par

28. 但我不仅只谈论妇女。也要远离那些戴着锁链、像女人一样留长发、违背宗徒训诫的男子，\BibleRef{格前 11:14}更不必说那如山羊的胡须、黑色斗篷、赤足忍受寒冷。这一切都是魔鬼的标记。这样的一个人——\Person{安提摩斯}——曾使\Place{罗马}哀叹；而\Person{索弗若纽斯}是更近的例子。这些人一旦进入贵族的宅邸，欺骗了\BibleQuote{那被罪恶捆绑、被各种私欲牵引、常常学习却终不能认识真理的愚蠢妇女}，\BibleRef{弟后 3:6}他们便装出愁苦的神情，假装长期禁食，却于夜间秘密宴饮。羞耻禁止我多说，因为我的言语可能更像谴责而非劝诫。还有一些人——我说的是我同品级的人——他们寻求司铎职和执事职，仅仅为了能更无约束地见到妇女。这些人只关心他们的服饰；他们随意使用香水，确保皮鞋上没有皱褶。卷曲的头发显出烫钳的痕迹；手指上戒指闪闪发光；他们踮着脚尖走过潮湿的路面，以免溅湿双脚。当你看到如此行径的人时，应将他们视为新郎而非神职人员。有些人将全部的精力和生命投入到唯一的目标：知道已婚女士的名字、住所和品格。我将在此简要描述这一职业的首领，以便从师傅的形象认出弟子。他日出时起身出行；他安排好走访的顺序；他选择最短的路；并且，作为麻烦的老头，他几乎闯入仍在睡梦中的女士的卧室。如果他看到心仪的枕头或雅致的桌布——实际上任何家居用品——他便会称赞它，欣赏它，拿在手里，抱怨自己没有类似的东西，从而向主人乞求，更确切地说，强索得到它。事实上，所有的妇女都害怕得罪这位城镇的新闻传播者。贞洁和禁食同样令他不悦。他喜欢的是美味的早餐——比如说，一只肥美的幼鹤，通常被称为雏鹤。他言语粗鲁冒失，随时准备互相指责。无论你转向何处，他都是你最先看到的面前之人。任何流传的新闻，他不是谣言的起源就是其放大者。他每小时更换马匹；马匹如此光滑而精神，你会以为他是\Place{色雷斯}国王的兄弟。\par

29. 狡猾的敌人对我们使用了许多计谋。经上记载：\BibleQuote{蛇是上主天主所造的一切田野动物中最狡猾的。}\BibleRef{创 3:1}宗徒也说：\BibleQuote{我们并不是不知道他的阴谋。}\BibleRef{格后 2:11}既不应装出刻意的寒酸，也不应追求时髦的华丽，这才符合基督徒的身份。如果你有什么不知道的，如果你对圣经有什么疑问，就去请教一个生活值得称许、年龄无可怀疑、名誉与事实相符的人；一个能这样说的人：\BibleQuote{我曾把你们许配给一个丈夫，把你们如同贞洁的童女献给基督。}或者，如果没有这样的人能解释，那么宁可因无知而避免危险，也胜于为求知而招惹危险。要记住你走在网罗之中，许多资深的贞女，其贞洁从未受到质疑，却在死亡的门槛上让冠冕从手中坠落。\par

如果你们中的任何侍女分享了你的圣召，不要高抬自己反对她们，也不要因你是她们的主母而骄傲。你们选择了同一位新郎；你们同唱相同的圣咏；一起领受基督的身体。那么，为什么你们的想法要有分别？你必须努力赢得他人，并且为了更容易吸引人，你必须极其尊重你队伍中的贞女。如果你发现其中某人在信仰上软弱，要关怀她，安慰她，抚爱她，把她的贞洁当作你的珍宝。但如果一个女孩假装有圣召仅仅是为了逃避服侍，你就向她高声宣读宗徒的话：\BibleQuote{结婚比焚烧更好。}\BibleRef{格前 7:9}\par

闲散的人和好管闲事的人，无论是贞女还是寡妇；她们挨家挨户走访已婚妇女，表现出比舞台上的寄生虫更大的无耻厚颜，你要把她们如同瘟疫一样抛弃。因为\BibleQuote{坏朋友会败坏良好的品行}，\BibleRef{格前 15:33}像这样的妇女只关心最低级的欲望。她们常常会催促你说：\Quote{亲爱的，好好享受你的优势，趁你活着的时候生活，}以及\Quote{你肯定不是为你的孩子们存钱吧。}她们贪恋酒色，向人们的心灌输各种邪恶，甚至引诱最严肃的人去追求使人软弱的快乐。而当她们\BibleQuote{开始违背基督而放纵情欲时，便要结婚，因而受审判，因为他们废弃了起初的信德。}\BibleRef{弟前 5:11}\par

不要寻求显得过于雄辩，也不要轻浮地编织诗句，更不要以抒情歌谣取乐。不要出于做作，效仿那些已婚女子的病态趣味：她们时而紧咬牙齿，时而大张嘴唇，说话带咬舌音，故意把字音吞掉，因为她们认为自然发音是乡下人的标记。因此，她们在一种舌头的通奸中寻欢作乐。因为\BibleQuote{光明与黑暗有什么相通？基督与\OriginalTerm{贝里雅耳}{Belial}有什么相和？}\BibleRef{格后 6:14} 贺拉斯怎能与圣咏集同行，维吉尔怎能与福音书同行，西塞罗怎能与宗徒同行？若看到你坐在偶像的庙中赴宴，岂不使弟兄跌倒？\BibleRef{格前 8:10} 虽然\BibleQuote{在洁净者一切皆洁净}，\BibleRef{铎 1:15} 并且\BibleQuote{凡以感恩领受的，没有一样可弃绝}，\BibleRef{弟前 4:4} 但我们仍不能同时喝基督的杯和魔鬼的杯。\BibleRef{格前 10:21} 让我向你讲述我自己悲惨经历的故事。\par

30. 许多年前，我为了天国的缘故，已割舍了家庭、父母、姐妹、亲属，以及更艰难地，我习惯的美味饮食；当我前往耶路撒冷去从事战斗时，我仍然无法舍弃我在罗马费尽心血建立起来的藏书室。于是，我这个可怜人，禁食只是为了之后能读西塞罗。在无数个守夜的晚上，在从心底涌出的泪水中，在回忆我过去的罪之后，我又会拿起普劳图斯。有时清醒过来，我开始读先知书，但他们的风格似乎粗陋可憎。我双眼失明，看不见光明，却把过错归咎于太阳，而非它们。当古蛇这样玩弄我时，约在四旬期中期，一场高烧侵袭了我虚弱的身体，它彻底破坏了休息——这故事几乎难以置信——如此消耗我可怜的身躯，几乎只剩皮包骨。同时，葬礼的准备工作在继续；我的身体逐渐冷却，生命的温暖只留在我跳动的心胸。突然，我被神灵提起，被拖到审判者的审判台前；那里的光明如此耀眼，周围站着的那些人是如此光彩夺目，以至于我伏在地上，不敢抬头。被问及我是谁、是什么，我回答：\Quote{我是一个基督徒。}但那位主持者说：\Quote{你说谎，你是西塞罗的追随者，而非基督的。因为\InnerQuote{你的宝藏在哪里，你的心也在哪里。}}\BibleRef{玛 6:21} 我立刻哑口无言，在鞭打的打击中——因为他命令鞭打我——更被良心的火焰所折磨，心中思量那句诗：\Quote{在坟墓中，谁还称谢你？}然而，我开始哭泣哀嚎说：\Quote{主啊，可怜我吧，可怜我。}在鞭打声中，这呼声仍然可闻。终于，旁观者跪倒在主持者的膝前，祈求他可怜我的年轻，给我时间悔改过错。他们力劝，若我以后再读外教人的作品，他仍可折磨我。在那可怕时刻的压力下，我本会愿意做出比这些更大的承诺。于是我起誓，呼求他的名说：\Quote{主，若我再占有世俗的书，或再读这些书，我就否认了你。}起誓后我被释放，回到上界，在所有人的惊讶中，我睁开泪眼，痛苦足以说服不信者。这不是睡眠或虚幻的梦，如我们常被嘲弄的那样，我以我所俯伏的审判台和我所畏惧的可怕审判为证。愿我今后再也不要落入那样的审讯中！我承认我的肩膀青紫，醒来后很久还感到伤痕，并且从此以后，我读天主的书比先前读人的书更加热切。\par

31. 你也必须避免贪婪之罪，这不仅是不夺取属他人的东西——因为那是国家法律所惩罚的——而且也不占有自己的财产，这已不再是你的了。主说：\BibleQuote{如果你们在不义的钱财上不忠信，谁还把真实的钱财托付给你们呢？}\BibleRef{路 16:12} \Quote{不义的钱财}是指金银的数量，而\Quote{真实的钱财}是指精神遗产，别处也说：\Quote{人的生命赎价是他的财富。}\Quote{没有人能事奉两个主人，因为他不是恨这个而爱那个，就是依附这个而轻视那个。你们不能事奉天主又事奉\OriginalTerm{玛门}{Mammon}。}\BibleRef{玛 6:24} 即财富；因为在叙利亚的异教语言中，财富被称为玛门。那窒息我们信仰的\Quote{荆棘}\BibleRef{玛 13:7} 就是对生活的忧虑。\BibleRef{玛 6:25} 对外邦人所追求之物的挂虑\BibleRef{玛 6:32} 是贪婪之根。\par

但你会说：\Quote{我是一个娇生惯养的少女，不能用手劳动。假如我活到老年然后生病，谁会怜悯我呢？} 请听耶稣对宗徒们说：\BibleQuote{不要为吃什么而忧虑，也不要为身体穿什么而操心；生命不是贵于食物，身体不是贵于衣服吗？你们看天空的飞鸟：它们不播种，也不收割，也不收集在仓里，你们的天父尚且养活它们。} \BibleRef{玛 6:25} 如果衣服不足，就把百合花放在眼前。如果饥饿来袭，就思考那些祝福贫穷和饥饿者的话语。如果痛苦折磨你，请读\BibleQuote{所以我喜欢软弱}，以及\BibleQuote{有一根刺加在我的肉体上，就是撒殚的使者来攻击我，免得我过于高举自己。} 应欢欣于天主的一切判断；因为圣咏作者不是说过：\BibleQuote{犹大的女子因你的判断而欢喜，上主}？让这些话常在你口中：\BibleQuote{我赤身出离母胎，也必赤身返回那里} \BibleRef{约 1:21}，以及\BibleQuote{我们没有带什么到这世界上，确实也不能带走什么。} \BibleRef{弟前 6:7}\par

32. 今日你看见女人们将衣柜塞满衣服，每日更换长袍，却始终无法战胜飞蛾。偶有更谨慎者穿破一件衣服；然而，当她衣衫褴褛时，箱子却是满的。羊皮纸被染成紫色，金子被熔铸成字母，手抄本装饰着珠宝，而基督却赤身濒死躺在门外。她们向穷人伸手时吹起号角；\BibleRef{玛 6:2} 邀请参加爱筵时雇请传令员。我最近看见罗马最尊贵的贵妇——我不提她的名字，因为我不是讽刺作家——在圣伯多禄大殿前由一群太监前导。她亲手给穷人施舍钱币，每人一枚，为使人以为她更虔诚。这时发生了一件绝不罕见的事。一位老妇人，\Quote{满布年岁与褴褛}，跑上前去想要第二枚钱币，但当轮到她时，她得到的不是一文钱，而是一记足以从她有罪的血管中抽出血来的击打。\par

\BibleQuote{贪爱钱财是万恶的根源}，\BibleRef{弟前 6:10} 而宗徒说贪婪就是拜偶像。\BibleRef{哥 3:5} \BibleQuote{你们先寻求天主的国，这一切自会加给你们。} \BibleRef{玛 6:33} 上主决不会让正义的灵魂饿死。圣咏作者说：\BibleQuote{我曾年幼，现今年老，从未见过义人被弃，也未见过他的后裔求乞。} \BibleRef{咏 37:25} 厄里亚由服役的乌鸦喂养。\BibleRef{列上 17:4} 匝尔法特的寡妇，原以为当晚就会与儿子们一同饿死，却自己绝食以便喂养先知。那来求食的人反成了喂养者，因为借着奇迹，他使空瓮满溢。\BibleRef{列上 17:9} 宗徒伯多禄说：\BibleQuote{银子和金子我没有，但我把所有的给你：因纳匝肋人耶稣基督的名字，你起来行走罢！} \BibleRef{宗 3:6} 如今许多人虽不口说，却以行动宣告：\Quote{信德和怜悯我没有；但我所有的，银子和金子，这些我不给你们。} \BibleQuote{我们有吃有穿，就当知足。} \BibleRef{弟前 6:8} 请听雅各伯的祈祷：\BibleQuote{如果天主与我同在，在我所走的路上保护我，给我食物吃，衣服穿，那么上主将是我的天主。} \BibleRef{创 28:20} 他只祈求必需品；然而二十年后，他回到客纳罕地，财物众多，儿女更富裕。\BibleRef{创 32:5} 圣经中有无数例子教训人要\BibleQuote{谨慎，戒避一切贪心。} \BibleRef{路 12:15}\par

33. 我既已触及这个话题——若基督允许，它将自成一篇论文——就讲述尼特里亚不多年以前发生的事。一位弟兄比贪财更节俭，且不知道主曾被以三十块银钱出卖，\BibleRef{玛 26:15} 他死时留下了一百枚编麻织布挣来的钱币。因邻近约有五千位隐修士，分居同样多的独居小室，就召开会议决定如何处理。有的说应将钱币分给穷人；有的说应献给教会；还有的说应送回死者亲属。然而，玛加略、庞博、依西多禄以及其余被称为父老的人，因着圣神发言，决定将钱币与死者一同埋葬，并说：\BibleQuote{你的钱财与你一同丧亡。} \BibleRef{宗 8:20} 这决定并非过于严厉；因为如此大的恐惧已降临整个埃及，如今留下一枚先令也被视为罪行。\par

34. 正如我提到了隐修士，并且知道你喜欢听关于神圣事物的事，请稍听我片刻。在埃及有三类隐修士。首先，有\OriginalTerm{集体隐修士}{cœnobites}，用他们外邦人的语言称为\OriginalTerm{萨乌色斯}{Sauses}，或者按我们的说法，是生活在团体中的人。其次，有\OriginalTerm{独修隐修士}{anchorites}，他们住在旷野，各人独处，之所以如此称呼是因为他们已远离人类社会。第三，有一类称为\OriginalTerm{雷莫博特}{Remoboth}，这是一种非常低劣且不受重视的类型，为我本省所特有，或至少起源于那里。他们三三两两住在一起，但很少人数更多，不受任何规则约束；完全随心所欲。他们将一部分收入捐给公共基金，用以为所有人提供食物。他们大多住在城市和要塞中；仿佛他们的手工是神圣的，而非他们的生活，他们所卖的一切都极其昂贵。他们常争吵，因为不愿在自备食物的同时服从他人。的确，他们在禁食方面相互竞争；他们把本应私人的事当作炫耀的机会。凡事追求效果：袖子宽松，靴子鼓胀，衣着最粗糙。他们总是叹息，或拜访贞女，或嘲笑司铎；然而每逢节日，他们便吃得过多，以致生病。\par

35. 既然我们除去了这些如同众多瘟疫的人，让我们来到那个更众多的共同生活的阶级，他们正如我们所说，称为集体隐修士。在他们中间，合一的第一个原则是服从长上，并遵行他们所命令的一切。他们被分成十人和百人的团体，因此每第十人对其他九人有权柄，而第一百人手下有十位这样的官员。他们彼此分开，住在各自的单间。根据他们的规则，隐修士在第九时辰之前不得拜访他人；除了上述的院长们，他们的职责是用安抚的话语安慰那些思想烦乱的人。第九时辰之后，他们聚在一起，按照惯例咏唱圣咏并诵读圣经。祈祷结束，众人坐下后，一位称为父亲的人站起来，开始讲解当日的篇章。当他讲话时，一片寂静；无人敢看邻座或清嗓子。讲者的赞颂在于听者的泪水。无声的眼泪顺着他们的脸颊滚落，但嘴唇不发出啜泣。然而当他开始谈论基督的王国、未来的福乐和将要来临的荣耀时，可以看到每个人都在轻声叹息、举目向上，对自己说：\BibleQuote{哦，但愿我有鸽子般的翅膀！我就能飞去，得享安息。}之后聚会散开，每个十人团体随其父亲到自己的桌前。他们轮流服务，每人一周。用餐时毫无声响；吃饭时无人交谈。面包、豆类和蔬菜是他们的食物，所用的唯一调味料是盐。酒只给老人，他们与孩童常有特别准备的餐食，以弥补年迈的损耗，并防止年轻人过早衰败。用餐完毕，他们一起起身，唱完一首圣歌后，返回自己的住处。在那里，每人与其同伴交谈直到傍晚，内容如下：\Quote{你注意到某人了吗？他多么有恩宠！他多么沉默！他走路多么稳重！}若有人软弱，他们就安慰他；若有人热爱天主热切，他们就鼓励他更加恳切。因为夜间，除了公共祈祷，各人在自己房间守夜，他们逐一巡视所有单间，将耳朵贴在门上，仔细查察其居住者在做什么。若发现隐修士懈怠，他们并不责备；而是掩饰所知，更频繁地拜访他，起初是劝勉而非强迫他多祈祷。每天有既定任务，交给院长，再由院长转交管家。后者每月一次向他们的共同父亲提交详细账目。他也品尝做好的菜肴，因为无人允许说：\Quote{我没有上衣、外袍或草席。}他如此安排，使无人需要请求或缺乏所需。若有隐修士生病，便移至更宽敞的房间，由老人们精心照料，使他既不缺少城市的奢华，也无缺母亲的慈爱。每个主日，他们整天用于祈祷和阅读；事实上，当完成工作后，这些是他们的惯常活动。每天他们背诵一段圣经。全年守同样的斋戒，但在四旬期内，生活更为严格。圣神降临节后，他们用午间餐代替晚膳；既为满足教会的传统，也为避免双重食物使胃过重。\par

类似的对厄色尼人的描述见于斐罗——柏拉图的模仿者，以及约瑟夫——希腊的利维乌斯，在其关于犹太被掳的叙述中。\par

36. 因为我目前的主题是贞女，关于隐修士说得有点过多。因此，我将转到第三类，称为独修隐修士，他们从隐修院进入旷野，只带面包和盐。\Person{保禄}引入了这种生活方式；\Person{安当}使它出名；而更早的——\Person{洗者若翰}是第一个身体力行者。\Person{耶肋米亚}先知用这些话描述了这样一个人：\BibleQuote{人在年轻时背负轭，这原是好的。他独坐静默，因为轭加在他身上。他让打他的人打他的面颊，受尽凌辱。因为主不会永远弃绝。}独修隐修士的奋斗和他们生活——在肉身内，却不属于肉身——你若愿意，我改日再向你解释。我现在必须回到贪财的话题，这是为讲隐修士而搁下的。你若心中有他们，就会鄙视金子和银子，甚至大地和天空。与基督结合，你将歌唱：\BibleQuote{主是我的福分。}\par

37. 宗徒既命我们\Quote{不断祈祷}\BibleRef{得前 5:17}，且圣人连睡眠均为祈愿，但我们仍宜有定时祈祷，以免因工作缠身而忘却本分。众人皆知，第三、第六、第九时辰，以及清晨与傍晚，均当祈祷。饮食前不可不祈祷，离席前应向造物主谢恩。夜间应起身二三次，背诵所记圣经章节。离屋外出时，当以祈祷为甲；自街归家时，应先祈祷而后坐，待灵魂获养后方使孱弱之躯得安。凡所行事，凡所举步，手画主之十字架。勿指责人，勿诽谤尔母之子。\Quote{尔何人，竟论断他人之仆？其或仆或兴，乃属己主；且彼必得兴起，因主能使之兴起。}若尔禁食二三日，勿自以为胜于不禁食者。尔禁食而发怒，他人进食而面带笑容；尔争闹中发泄暴躁与饥饿，彼则节用饮食而感谢天主。每日依撒意亚呼喊：\Quote{我所拣选的禁食，岂是这一种？}\BibleRef{依 58:5}又说：\Quote{在你们禁食之日，你们寻求自己的快乐，压迫你们所有的劳工。看，你们禁食是为了争辩和争吵，以凶恶的拳头打人。你们这样禁食，岂能使我悦纳？}其怒气竟至于夜不消退，甚至月相变换仍不改变，此等禁食有何益？\par

38. 你当省察自己，因自己的成功而夸耀，而非他人的失败。有些女子关心肉体，计算收入和每日开支；这些不是你的榜样。犹达斯是叛徒，但十一位宗徒并未动摇。\Person{菲革罗}和\Person{亚历山大}遭遇了船难，但其余的人继续奔跑信仰的赛程。不要说：\Quote{某人享受自己的财产，受人尊敬，她的兄弟姐妹来看她。难道她就不是贞女了吗？}首先，她是否为贞女尚有疑问。因为\Quote{上主看人不像人看人：人看外貌，上主却看人心。}\BibleRef{撒上 16:7}再者，她可能肉身是贞女，心灵却不是。按宗徒所说，真正的贞女是\Quote{在肉身和心灵上都圣洁的。}\BibleRef{格前 7:34}最后，让她按自己的方式夸耀吧。让她漠视保禄的见解，享受自己的美好事物。但你我却当追随更好的榜样。\par

请你以荣福玛利亚为榜样，她的卓越纯洁使她堪当为主的母亲。当天使加俾额尔以人形降到她面前，说：\Quote{万福，充满恩宠者，主与你同在。}\BibleRef{路 1:28}她惊惶失措，不能回答，因她从未被男子问候过。但得知他是谁后，她便发言，那位曾畏惧男子的人，竟无所畏惧地与天使交谈。现在，你，也可以成为主的母亲。\Quote{你取一个大卷轴，用人用的笔写上：\InnerQuote{玛赫尔沙拉耳哈市巴次。}}当你到女先知那里，在腹中怀孕，生下儿子时，说：\Quote{主啊，我们在你的敬畏中怀孕，我们疼痛，产生了你拯救的灵，就是我们在世上所行的。}那时你的儿子要回答：\Quote{看，我的母亲和我的兄弟。}\BibleRef{玛 12:49}祂，你这些时刚铭刻在心版上，并用笔写在更新的心版上——祂，在从敌人夺回掠物，并剥夺了率领者和掌权者，将他们钉在祂的十字架上\BibleRef{哥 2:14}——奇妙受孕，长大成人；当祂年岁渐长，不再视你为母亲，而是为新娘。成为殉道者，或宗徒，或基督，需要艰苦奋战，但带来巨大赏报。\par

这一切努力，唯有在教会范围之内方为有益；我们必须在同一屋内庆祝逾越节\BibleRef{出 12:46}，必须与\Person{诺厄}一同进入方舟\BibleRef{伯前 3:20}，必须与蒙成义的妓女\Person{辣哈布}一起逃避耶里哥的倾覆。\BibleRef{雅 2:25}那些据说在异端中，以及在臭名昭著的\Person{摩尼}追随者中的所谓贞女，不应视为贞女，而应视为娼妓。因为如果——如他们所声称——魔鬼是肉身的创造者，他们怎能尊敬由仇敌所造之物？不；正因她们知道贞女之名带来荣耀，她们才披着羊皮四处行走，如同豺狼。\BibleRef{玛 7:15}正如假基督冒充基督，这些贞女也冒用可敬之名，以便更好地掩盖可耻的生活。欢欣吧，我的姊妹；欢欣吧，我的女儿；欢欣吧，我的贞女；因为你决意成为他人虚伪矫饰的实在。\par

39. 我在此陈述之事，在不爱基督的人看来，必以为严苛。但若有人视世间一切荣华为粪土，以太阳之下万物为虚幻，为赢得基督\BibleRef{斐 3:8}；若有人与主同死同复活，并已把肉身连同私欲偏情钉在十字架上；他将勇敢呼喊：\Quote{谁能使我们与基督的爱隔绝？是困苦吗？是窘迫吗？是迫害吗？是饥饿吗？是赤身露体吗？是危险吗？是刀剑吗？}又说：\Quote{我深信无论是死是生，是天使，是率领者，是有权能者，是现在的事，是将来的事，是高天，是深渊，或是任何别的受造物，都不能使我们与天主的爱隔绝，这爱是在我们的主基督耶稣内的。}\par

为了我们的救恩，天主之子成了人子。九个月之久，祂在母腹中等待出生，忍受最令人厌恶的状况，满身血污地诞生，被裹在襁褓中，受尽抚弄。祂把世界紧握在手中，却局限于马槽的狭窄空间。我不提那三十年，祂在隐晦中生活，满足于父母的贫穷。\BibleRef{路 2:51}当祂受鞭打时，祂缄默不语；当祂被钉十字架时，祂为钉祂的人祈祷。\Quote{我拿什么报答主赐给我的一切恩惠？我要举起救恩之杯，呼求主的名。主的圣人在祂眼中是珍贵的。}我们唯一相称的回报就是以血还血；既然我们是赖基督的血得救赎，就应欣然为我们的救赎者舍命。有哪位圣人不先争战而得冠冕？义人亚伯尔被杀。亚巴郎有失去妻子的危险。而且，我必须避免书卷过分扩大，请你自行查找：你会发现所有圣人都受过苦难。唯独撒罗满生活在奢华中，也许正是因此他堕落了。因为\BibleQuote{主所爱的，祂必惩戒；祂所收纳的每个儿子，祂都鞭打。}\BibleRef{希 12:6}哪样更好？是为短时间争战，为栅栏扛木桩、持武器、在沉重的盾牌下疲惫不堪，以后能如胜利者喜乐？还是因为不能忍受一时就成了永远的奴隶？\par

40. 爱无难事；对热切者而言，没有困难的任务。想想\Person{雅各伯}为许配给他的妻子\Person{辣黑耳}所受的一切。圣经说：\BibleQuote{雅各伯为辣黑耳服役了七年。他因爱她，就觉得这七年好像几天。}\BibleRef{创 29:20}后来他自己告诉我们他所承受的。\BibleQuote{白日受干旱，黑夜受寒霜。}\BibleRef{创 31:40}同样，我们必须爱\Person{基督}，常寻求祂的拥抱。那么，一切艰难都会显得容易；一切长久都会视为短暂；被祂的箭射中，我们要每时每刻说：\Quote{我哀叹我的寄居期延长了。}因为\BibleQuote{现时的痛苦，与将来要在我们身上显示的光荣，是不能相比的。}\BibleRef{罗 8:18}因为\BibleQuote{磨难产生坚忍，坚忍产生熟练，熟练产生希望；而希望不会使人蒙羞。}\BibleRef{罗 5:3}当你觉得命运难以承受时，请读\Person{保禄}的\WorkTitle{致格林多人后书}：\BibleQuote{我更多劳苦，更多鞭打，更频繁坐监，屡次受死。被犹太人鞭打五次，每次四十下减去一下；被棍打了三次；被石头打了一次；遭遇船难三次；在深海里度过一日一夜；多次行路，遭遇江河的危险、强盗的危险、同族人的危险、外邦人的危险、城中的危险、旷野的危险、海上的危险、假弟兄的危险；劳碌辛苦，屡次不眠，饥渴，多次禁食，寒冷赤身。}\BibleRef{格后 11:23}我们中谁能声称有他在这里列举的德行的最微小部分？然而正是这些后来使他敢于说：\BibleQuote{我完成了赛跑，保守了信德。从今以后，正义的冠冕已为我预备，就是主，正义的审判者，在那一天要赐给我的。}\BibleRef{弟后 4:7}\par

但我们，如果食物不如平时可口，就愠怒，以为喝掺水的酒是帮天主的忙。如果送来的水稍热，我们就摔杯子、掀桌子，鞭打有过错的仆人直到流血。\BibleQuote{天国遭受暴力，暴力的人夺取它。}\BibleRef{玛 11:12}然而，除非你使用暴力，否则你永不能夺取天国。除非你坚持敲门，否则你永不会领受圣事的食粮。\BibleRef{路 11:5}你想，当肉体渴望如同天主，并上升到天使坠落的地方\BibleRef{依 14:12}去审判天使时，这难道不是真正的暴力吗？\par

41. 我恳求你，暂时从你的囚室中出来，在你眼前描绘你现今劳苦的赏报，这赏报是\BibleQuote{眼所未见，耳所未闻，人心所未想到的。}\BibleRef{格前 2:9} 那一日将有何等的荣耀，当主的母亲玛利亚由她的童贞歌咏团陪伴着来迎接你！当红海已过、法郎及其军队沉没后，亚郎的姐妹米黎盎手拿着鼓，向应和的妇女们歌唱：\BibleQuote{你们应歌颂上主，因为他赫赫胜利，将马和骑士投入海中。}\BibleRef{出 15:20} 那时，特克拉将欢欣地飞奔来拥抱你。那时，你的新郎亲自前来，说：\BibleQuote{我的爱卿，我的佳丽，起来，走吧！看！冬天已过，雨止天晴。}\BibleRef{歌 2:10} 那时，天使们将惊奇地说：\BibleQuote{那像晨曦般出现，美丽似月亮，皎洁如太阳的是谁？}\BibleRef{歌 6:10} \BibleQuote{女儿们看见你，都要祝福你；王后们和嫔妃们也要赞美你。}\BibleRef{歌 6:9} 在这些之后，另一队贞洁的妇女将迎接你。撒辣将偕同已婚者前来；法奴耳的女儿亚纳将偕同寡妇们同来。在那一队中，你将遇到你的生母；在另一队中，你将遇到你的神母。一个将因生育你而喜乐，另一个将因教导你而欢欣。那时，主将确实骑着他的驴驹进入天上的耶路撒冷。那时，小孩子们（依撒意亚中救主曾说：\BibleQuote{看！我和上主赐给我的孩子们}\BibleRef{依 8:18}）将举起胜利的棕榈枝，齐声高唱：\BibleQuote{贺三纳于至高之天！奉主名而来的，当受赞美！贺三纳于至高之天！}\BibleRef{玛 21:9} 那时，\BibleQuote{十四万四千人}将拿着竖琴站在宝座和长老们面前，唱新歌。除了那些被选定的人，没有人能学会那歌。\BibleQuote{这些人没有和女人沾染过，因为他们仍是童身。这些人无论羔羊往哪里去，都跟随他。}\BibleRef{默 14:1} 每当这尘世的虚幻表演试图诱惑你，每当你看见世上虚荣的排场，你要在心灵中将自己转移到天堂，尝试现在就成为你将来所要成为的，这样你将会听到你的新郎说：\Quote{请将我如荫庇放在你心上，如印玺戴在你臂上。}于是，你身心灵都得着力量，也将大声呼喊说：\BibleQuote{洪流不能熄灭爱情，江河不能将它冲去。}\BibleRef{歌 8:7}\par

\SourceRecord{Translated by W.H. Fremantle, G. Lewis and W.G. Martley.；Nicene and Post-Nicene Fathers, Second Series；Vol. 6.；Edited by Philip Schaff and Henry Wace.；Buffalo, NY: Christian Literature Publishing Co.,；1893.；<http://www.newadvent.org/fathers/3001022.htm>.}

% structure: p0001 source=h1 target=section reason=page-title

\section{第二十三书}

\Emph{致\Person{玛尔策拉}}\par

\EditorialNote{\Person{热罗尼莫}写信给\Person{玛尔策拉}，安慰她失去一位朋友，这位朋友和她一样，是\Place{罗马}一个宗教团体的领导人。\Person{莱亚}的死讯首先传到\Person{玛尔策拉}那里，当时她正与\Person{热罗尼莫}一起研读第七十三篇圣咏。当天晚些时候，他写了这封信，在信中颂扬\Person{莱亚}之后，他将她的结局与当选执政官\Person{维提乌斯·阿戈里乌斯·普雷特克斯塔图斯}（此人极有能力和正直）的结局相对照，而他宣称此人现处于\Place{塔尔塔鲁斯}。写于公元384年的\Place{罗马}。}\par

1. 今天，大约在第三时辰，正当我开始与你一起阅读第七十二篇圣咏——即第三卷的第一篇——并解释其标题部分属于第二卷，部分属于第三卷——我是说，前一卷以\BibleQuote{\Person{达味}之子\Person{叶瑟}的祈祷完毕}结束，下一卷以\BibleQuote{\Person{阿撒夫}的圣咏}开始——正当我读到义人宣告的那段话：\BibleQuote{若我说：我要这样讲，看，我就会得罪你子孙的世代}——这一节在拉丁文译本中有不同译法——突然传来消息，我们最圣善的朋友\Person{莱亚}已离世而去。你自然大惊失色；因为没有几个人，如果有的话，会在瓦器破裂时不流泪。\BibleRef{格后 4:7}但你流泪并非怀疑她未来的命运，而只是因为你未能向她履行对死者应尽的最后哀悼之职。最后，在我们继续交谈时，第二次消息告知我们，她的遗体已被送往\Place{奥斯蒂亚}。\par

2. 你可能会问，重复这一切有什么用处。我要用宗徒的话回答：\BibleQuote{事事都有益处。}\BibleRef{罗 3:2}首先，这表明所有人都应当欢呼喜乐，迎接一个灵魂的解脱，这灵魂已将\Person{撒旦}践踏在脚下，并最终为自己赢得了一顶安宁的冠冕。其次，这给了我一个机会简要描述她的生平。第三，这使我能够向你保证，那位当选执政官，他所处时代的毁谤者，如今正在\Place{塔尔塔鲁斯}。\par

谁能充分颂扬我们亲爱的\Person{莱亚}的生活方式呢？她归向主的转变如此彻底，以致成为隐修院的院长，对院中的童贞们表现出真正的母亲情怀，她穿粗苦衣而不穿细软衣服，彻夜祈祷不眠，并且更多地以榜样而非言语教导同伴。她的谦逊如此之大，以至于她曾是多人的主母，却被视为众人的仆人；确实，她越不被看作人世的主母，就越成为基督的仆人。她不修边幅，不在意发式，只吃最粗糙的食物。尽管如此，她所做的一切都避免炫耀，以免在世上得到她的赏报。\BibleRef{玛 6:2}\par

3. 现在，因此，作为对她短暂辛劳的回报，\Person{莱亚}享受永久的福乐；她被迎入天使的歌团；她在\Person{亚巴郎}的怀中被安慰。并且，正如从前乞丐\Person{拉匝禄}看见那穿紫红袍的富人躺在痛苦中，同样，\Person{莱亚}看见那位执政官，不再穿凯旋袍而是身着丧服，并请求从她指尖滴一滴水。\BibleRef{路 16:19}我们在这里看到了多么大的转变！几天前，城中的最高显贵走在他前面，他像庆祝胜利的将军一样登上\Place{卡比托利欧山}的壁垒；罗马人民跳起来欢迎和鼓掌，他的死讯传来，全城震动。现在他荒凉赤裸，囚禁在最污秽的黑暗中，并不像他不幸的妻子错误断言的那样，安置在银河的王室住所。另一方面，\Person{莱亚}，她总是被关在她的小室里，看似贫穷卑微，她的生活被视为愚妄，\BibleRef{智 5:4}如今她跟随\Person{基督}，歌唱：\BibleQuote{正如我们听见的，也看见了，在我们天主的城中。}\par

4. 现在，关于这一切的教训，我流着泪、叹息着恳求你记住。当我们奔走在这世界的道路上时，我们不可穿两件衣服，即不可怀有双重的信仰，或负担上革制的鞋子，即死行；我们不可让装满银钱的囊袋压住我们，或倚靠世俗权力的杖。\BibleRef{玛 10:10}我们不可寻求同时拥有基督和世界。不；永恒的事物必须取代短暂的事物；\BibleRef{格后 4:18}并且，就身体而言，我们每天都预尝死亡，如果我们希望得到永生，就必须意识到我们仅是必死的。\par

\SourceRecord{Translated by W.H. Fremantle, G. Lewis and W.G. Martley.；Nicene and Post-Nicene Fathers, Second Series；Vol. 6.；Edited by Philip Schaff and Henry Wace.；Buffalo, NY: Christian Literature Publishing Co.,；1893.；<http://www.newadvent.org/fathers/3001023.htm>.}

% structure: p0001 source=h1 target=section reason=page-title

\section{\WorkTitle{第二十四书}}

致\Person{Marcella}\par

\EditorialNote{关于童贞女亚塞拉。她出生前已奉献给天主，玛尔塞拉的妹妹在十岁时被立为教会童贞女。从那时起，她过着最严格的苦修生活，先是作为阿文蒂诺山上玛尔塞拉团体的一员，后来成为该团体的领袖。热罗尼莫后来在离开罗马时给她写了一封信（第四十五书），现在将她树立为一个值得钦佩和效法的榜样。写于罗马，公元三八四年。}\par

1. 不要有人因信中包含的赞扬和责备而责怪我的信。责备罪人就是劝诫类似情况的人，赞扬有德者就是激发那些愿意行善者的热忱。前天，我向你们谈论蒙福的\Person{Lea}，我话音刚落，我的良心就被刺痛了。我想，在我谈论了一位作为寡妇地位较低的人之后，忽略一位童贞女是不对的。因此，在我现在的信中，我打算给你们简要叙述我们亲爱的\Person{Asella}的生平。请不要读给她听；因为她肯定会对自己作为对象的赞扬感到不悦。还是把它给你们认识的那些年轻女子看，好让她们以她为榜样行事，并把她的行为当作完美生活的典范。\par

2. 我略过她尚未出生、仍在母胎中就被祝福的事实，以及在梦中，她像一个童贞女一样，被放在一个比镜子还亮的闪亮玻璃碗中交给她的父亲。我也不提她献身于蒙福的童贞生活，那仪式在她刚满十岁时举行，那时她还只是个裹在襁褓中的婴儿。因为所有在行为之前的事都该归功于恩宠；\BibleRef{罗 11:6} 虽然，无疑，天主预知了未来，当祂圣化了尚未出生的\Person{耶肋米亚}，\BibleRef{耶 1:5} 当祂使\Person{若翰}在母胎中欢跃，\BibleRef{路 1:41} 并在世界建立之前，就拣选了\Person{保禄}去宣讲祂圣子的福音。\BibleRef{弗 1:4}\par

3. 现在我来谈谈她十二岁以后，凭自己的努力选择、抓住、坚持、进入并完成的生活。她关在狭小的斗室里，却漫步于乐园。禁食是她的消遣，饥饿是她的提神。如果她吃东西，那不是出于对吃的喜爱，而是因为身体疲惫；她限制自己只吃面包、盐和冷水，这刺激了她的食欲，胜过满足它。\par

但我几乎忘记提到我本应首先说的事情。当她的决心还很新鲜时，她取下了自己的金项链——那是灯芯式样的（之所以这么叫，是因为金属条连接在一起形成一条柔韧的链子）——并在父母不知情的情况下卖掉了它。然后她穿上一件深色衣服，那是她母亲从不愿意让她穿的，她立即将自己奉献给主，以此完成了她的虔诚事业。她向亲戚们表明，他们不必指望从这样一个已经通过她的衣着谴责了世界的人那里再得到任何让步。\par

4. 继续说我的故事，她的举止安静，生活非常隐秘。事实上，她很少出门或与男人交谈。更奇妙的是，尽管她深爱她的童贞妹妹，她也不愿见她。她亲手劳动，因为她知道经上记载：\BibleQuote{若有人不愿工作，就不可吃饭。}\BibleRef{得后 3:10} 她常在祈祷和圣咏中与新郎交谈。她悄然前往殉道者的圣地。这样的拜访给她带来喜乐，而且越是不被人认出，她就越喜乐。全年她坚持连续禁食，常常连续两三天不吃东西；但当四旬期到来时，她升起了——如果我可以这么说的话——每一片帆，几乎从一周结束到下一周结束都在禁食，面带\BibleQuote{和悦的面容}\BibleRef{玛 6:17} 或许难以置信，若不是因为\BibleQuote{在天主，一切皆有可能}\BibleRef{玛 19:26}，那就是她这样生活直到五十岁，而没有削弱消化能力或引发疝气的疼痛。躺在干地上没有影响她的肢体，她穿的粗糙的苦衣也没有使她的皮肤变得肮脏或粗糙。她身体健康，灵魂更健康，在独处中寻求一切欢乐，并在繁忙的\Place{Rome}中心为自己找到了一个隐修士的独修处。\par

5. 您比我更了解这一切，我所提供的少数细节都是从您那里学到的。您与\Person{Asella}如此亲密，以至于您亲眼见过她神圣的膝盖因经常祈祷而变得像骆驼膝一样坚硬。我只是把能从您那里收集到的事陈述出来。她严肃时令人愉快，愉快时又严肃；她的举止既讨人喜欢又庄重，既庄重又讨人喜欢。她苍白的脸表明节制，但不表示炫耀。她的言语是沉默，她的沉默是言语。她的步伐不快不慢。她的态度始终如一。她忽视精致，不关心衣着。当她注意衣着时，她也无意于此。她的生活如此始终如一，以至于在\Place{Rome}这个虚华、放纵和闲乐的中心，在那里谦虚被视为没有精神，善人赞扬她的行为，恶人不敢非议。让寡妇和童贞女效法她，让已婚妇女珍视她，让有罪的女人敬畏她，让主教们尊敬她。\par

\SourceRecord{Translated by W.H. Fremantle, G. Lewis and W.G. Martley.；Nicene and Post-Nicene Fathers, Second Series；Vol. 6.；Edited by Philip Schaff and Henry Wace.；Buffalo, NY: Christian Literature Publishing Co.,；1893.；<http://www.newadvent.org/fathers/3001024.htm>.}

% structure: p0001 source=h1 target=section reason=page-title

\section{第二十五封书信}

\Emph{致\Person{玛尔塞拉}}\par

\BibleRef{出 3:14}，\OriginalTerm{阿多乃}{Adonai}，\OriginalTerm{雅赫}{Jah}，\OriginalTerm{四字圣名}{tetragram JHVH}，\OriginalTerm{雅威}{Yahweh}，以及\OriginalTerm{沙代}{Shaddai}。写于\Place{罗马}，公元384年。\par

\SourceRecord{Translated by W.H. Fremantle, G. Lewis and W.G. Martley.；Nicene and Post-Nicene Fathers, Second Series；Vol. 6.；Edited by Philip Schaff and Henry Wace.；Buffalo, NY: Christian Literature Publishing Co.,；1893.；<http://www.newadvent.org/fathers/3001025.htm>.}

% structure: p0001 source=h1 target=section reason=page-title

\section{第二十六封书信}

\Emph{致\Person{马赛拉}}\par

\EditorialNote{对某些在译本中未翻译的希伯来语词的解释。这些词是\Quote{阿肋路亚}、\Quote{阿们}、\Quote{玛兰纳塔}。写于公元384年 罗马。}\par

\SourceRecord{Translated by W.H. Fremantle, G. Lewis and W.G. Martley.；Nicene and Post-Nicene Fathers, Second Series；Vol. 6.；Edited by Philip Schaff and Henry Wace.；Buffalo, NY: Christian Literature Publishing Co.,；1893.；<http://www.newadvent.org/fathers/3001026.htm>.}

% structure: p0001 source=h1 target=section reason=page-title

\section{第二十七封信}

\Emph{致\Person{玛尔塞拉}}\par

\EditorialNote{在这封信中，热罗尼莫为自己辩护，驳斥了篡改圣经文本的指控，并表明他只是使拉丁文新约译本与希腊原文相符。写于罗马，公元384年。}\par

1. 在我写了前一封信，对几个希伯来语词略作评论之后，突然传来报告说，某些可鄙之人故意攻击我，指责我试图改正福音中的经文，违背古人的权威和全世界的意见。现在，尽管依严格的权利，我本可轻视这些人（对牛弹琴是徒劳的），但为了避免他们重施故技，指责我傲慢，让他们接受我如下的答复：我并非如此迟钝或如此粗鲁无知（他们将这种无知视为圣德，自称是渔夫的门徒，仿佛人因一无所知而成圣）——我再重复一遍，我并非如此无知，以至于认为主的任何话语需要改正或不是受天主默感的；但圣经的拉丁文抄本因它们所呈现的所有差异而被证明是有缺陷的，我的目标是使它们恢复希腊原文的形式，而我的攻击者并不否认它们是从希腊原文翻译过来的。如果他们不喜欢从清澈的泉源汲取的水，就让他们喝浑浊的小溪的水；当他们读圣经时，让他们收起他们投向猎物众多的树林和盛产贝壳的水域的锐利目光。他们若愿意，就让他们在这一点上容易满足，将基督的话视为粗俗之言，尽管多少伟大的智者经过多少世代，不是解释，而是揣摩每个字的意义。让他们指责大宗徒缺乏文采，尽管有人论他说：\Quote{学问太多使他疯狂}\BibleRef{宗 26:24}。\par

2. 我知道当你读这些话时，你会皱起眉头，担心我的直言会播下新争执的种子；而且，如果你能，你会乐意用手指按住我的嘴，甚至不让我谈论别人不羞于去做的事情。但是，我问你，我哪里太过放肆？我曾用偶像的雕刻装饰我的餐盘吗？我曾在基督徒的宴席上，在贞女的眼前，摆出\OriginalTerm{萨提尔}{Satyrs}拥抱\OriginalTerm{酒神女信徒}{bacchanals}的污秽景象吗？或者我曾用过于尖刻的言辞攻击过任何人吗？我曾抱怨过乞丐变成百万富翁吗？我曾指责继承人为恩人举行的葬礼吗？我唯一不幸说过的话就是：贞女应该更多与妇女为伴，而不是与男人为伴。为此，我使全城哗然，人人都向我指指点点。\Quote{无故恨我的人，比我的头发还多}，我成了\Quote{他们的笑谈}。你认为此后我还会说什么鲁莽的话吗？\par

3. 可是，\Quote{当我转动车轮时，我开始塑造一个酒壶；怎么结果出来的是一个水罐？} 免得\Person{贺拉斯}嘲笑我，我回到我的两条腿的驴子那里，朝他们耳朵里灌输的，不是笛子的音乐，而是号角的轰鸣。如果他们愿意，他们可以说：\Quote{在盼望中喜乐；侍奉\Emph{时机}}，但我们要说：\Quote{在盼望中喜乐；侍奉\Emph{主}}。他们可能认为可以无条件地接受对一位\OriginalTerm{长老}{presbyter}的控告；但我们要用圣经的话说：\Quote{对长老的控告，除非\Emph{有两三个证人}，不要接受。犯罪的人，要在众人面前责备他们。}\BibleRef{弟前 5:19} 他们可能选择读：\Quote{这句话是\Emph{人的}，值得完全接纳}；我们满足于和希腊人一起犯错，也就是说，和说希腊语的宗徒本人一起犯错。因此，我们的译文是：\Quote{这句话是\Emph{忠信的}，值得完全接纳。}\BibleRef{弟前 1:15} 最后，让他们尽情享受他们的\OriginalTerm{高卢阉马}{Gallican geldings}；我们则满足于匝加利亚简朴的\OriginalTerm{驴}{ass}，它从缰绳中被解开，为救主服务，背上驮着主，从而实现了\Person{依撒意亚}的预言：\Quote{他在所有水边播种，牛和驴践踏之地是有福的！}\par

\SourceRecord{Translated by W.H. Fremantle, G. Lewis and W.G. Martley.；Nicene and Post-Nicene Fathers, Second Series；Vol. 6.；Edited by Philip Schaff and Henry Wace.；Buffalo, NY: Christian Literature Publishing Co.,；1893.；<http://www.newadvent.org/fathers/3001027.htm>.}

% structure: p0001 source=h1 target=section reason=page-title

\section{第二十八封书信}

\Emph{致\Person{玛尔采拉}}\par

\EditorialNote{对希伯来词\Quote{瑟拉}的解释。该词由\WorkTitle{七十贤士译本}译作\OriginalTerm{\GreekText{διάψαλμα}}{\GreekText{διάψαλμα}}。}此外，阿奎拉所译的\OriginalTerm{\GreekText{ἀ}}{\GreekText{ἀ}}和\OriginalTerm{\GreekText{εί}}{\GreekText{εί}}，在热罗尼莫时代与现在一样，是一个\OriginalTerm{难题}{crux}。\Quote{有些人，}他写道，\Quote{认为它是\InnerQuote{韵律的变化}，另一些人认为是\InnerQuote{换气}，还有一些人认为是\InnerQuote{新主题的开始。}又有人以为它与节奏有关，或标志着器乐的高潮。}热罗尼莫本人倾向于追随阿奎拉和奥力振，他们将这个词解释为\Quote{永远}，并暗示它表示完成，如同当代拉丁文手稿中的\Quote{explicit}或\Quote{feliciter}。写于\Place{罗马}，公元384年。\par

\SourceRecord{Translated by W.H. Fremantle, G. Lewis and W.G. Martley.；Nicene and Post-Nicene Fathers, Second Series；Vol. 6.；Edited by Philip Schaff and Henry Wace.；Buffalo, NY: Christian Literature Publishing Co.,；1893.；<http://www.newadvent.org/fathers/3001028.htm>.}

% structure: p0001 source=h1 target=section reason=page-title

\section{第二十九封书信}

第二十九封书信，致\Person{玛策辣}。\par

\BibleRef{撒上 2:18}与\Emph{特辣芬}\BibleRef{民 17:5}。写于\Place{罗马}，致\Person{玛策辣}，亦在\Place{罗马}，公元三八四年。\par

\SourceRecord{Translated by W.H. Fremantle, G. Lewis and W.G. Martley.；Nicene and Post-Nicene Fathers, Second Series；Vol. 6.；Edited by Philip Schaff and Henry Wace.；Buffalo, NY: Christian Literature Publishing Co.,；1893.；<http://www.newadvent.org/fathers/3001029.htm>.}

% structure: p0001 source=h1 target=section reason=page-title

\section{\WorkTitle{第三十封书信}}

\Emph{\Person{保拉}}\par

\Quote{\BibleRef{玛 7:7} 以接获天主圣三的\InnerQuote{三个饼}，\BibleRef{路 11:5} 并且，当主在我们前面行走时，行走于世水之上。} \BibleRef{玛 14:25} 写于\Place{罗马}，公元三八四年。\par

\SourceRecord{Translated by W.H. Fremantle, G. Lewis and W.G. Martley.；Nicene and Post-Nicene Fathers, Second Series；Vol. 6.；Edited by Philip Schaff and Henry Wace.；Buffalo, NY: Christian Literature Publishing Co.,；1893.；<http://www.newadvent.org/fathers/3001030.htm>.}

% structure: p0001 source=h1 target=section reason=page-title

\section{第三十一封书信}

\Emph{致欧斯托钦}\par

\EditorialNote{热罗尼莫写信感谢欧斯托钦在圣伯多禄瞻礼日送给他的一些礼物。他还对所送物品的神秘含义进行了道德诠释。这封信应与第四十四封书信相比较，其主题相似。写于公元384年罗马（圣伯多禄瞻礼日）。}\par

1. 鸽子、手镯和一封信，从一位贞女那里收到，外表上只是小礼物，但促使这礼物的行动却增加了它们的价值。既然蜂蜜不可作为祭品献给天主，\BibleRef{肋 2:11} 你便巧妙地除去它们过度的甜味，使它们——如果我可以这样说——因加了一点胡椒而变得辛辣。因为凡是单纯令人愉快或仅仅甜蜜的事物，都不能取悦天主。一切事物都必须有真理的尖锐调味。基督的逾越节必须与苦菜一同吃。\BibleRef{出 12:8}\par

2. 诚然，像圣伯多禄诞辰这样的节日，比平常应以更多的喜乐来调味；但我们的欢乐仍不能忘记圣经所定的界限，我们也不可过于远离我们角力的场地。你的礼物确实让我想起圣经的篇章，因为厄则克耳用它装饰耶路撒冷，为其佩戴手镯，\BibleRef{则 16:11} 巴路克从耶肋米亚那里接受书信，圣神在基督受洗时以鸽子的形状降临。\BibleRef{玛 3:16} 但为了也给你加上一点胡椒，并提醒你我之前的书信，我今天送你这三重警告。不要停止以善工装饰自己——这是基督徒妇人的真正手镯。\BibleRef{弟前 2:10} 不要撕毁写在你心上的书信，\BibleRef{格后 3:2} 如同那渎神的国王用他的小刀割破巴路克交给他的书卷。\BibleRef{耶 36:23} 不要让欧瑟亚对你说，如同对厄弗辣因说：\Quote{你像一只愚蠢的鸽子。}\BibleRef{欧 7:11}\par

你会说，我的话过于严厉，与像现在这样的节日不太相称。如果是这样，那是你自己的礼物性质激怒了我。只要你把苦与甜混在一起，你就得从我这里得到同样的回报，也就是说，既有尖锐的话语，也有赞美。\par

3. 不过，我不愿轻视你的礼物，尤其那篮精美的樱桃，带着童贞的羞涩而如此绯红，以致我幻想起它们是由卢库卢斯本人亲自采摘的。因为是他首先将这种水果引入罗马，在他征服了彭托斯和亚美尼亚之后；樱桃树之所以得名，是因为他从\Place{塞拉苏斯}带回来的。既然圣经没有提到樱桃，但却提到了一篮无花果，\BibleRef{耶 24:1} 我就用这些来指出我的教训。愿你结出这样的果实，如同那些在天主圣殿前生长的，祂说：\Quote{看，它们很好，非常好。}\BibleRef{耶 24:3} 救主不喜欢半冷不热的东西，祂虽接纳热心的，也不拒绝冷淡的，但在默示录中说，祂要把温吞的从口中吐出来。\BibleRef{默 3:15} 因此，我们必须谨慎，庆祝我们的圣日，与其以丰盛的食物，不如以精神的欢跃。因为希望以过度来荣耀一位殉道者，而这殉道者本人，如你所知，以斋戒取悦了天主，是完全不合理的。当你进食时，总要记得，吃饭之后应接以阅读，并接以祈祷。如果因为这样而行，你使某些人不悦，就对自己重复宗徒的话：\Quote{如果我还讨人的喜欢，我就不是基督的仆人了。}\BibleRef{迦 1:10}\par

\SourceRecord{Translated by W.H. Fremantle, G. Lewis and W.G. Martley.；Nicene and Post-Nicene Fathers, Second Series；Vol. 6.；Edited by Philip Schaff and Henry Wace.；Buffalo, NY: Christian Literature Publishing Co.,；1893.；<http://www.newadvent.org/fathers/3001031.htm>.}

% structure: p0001 source=h1 target=section reason=page-title

\section{书信三十二}

\Emph{致\Person{玛尔塞拉}}\par

\EditorialNote{热罗尼莫写道，他正忙于校勘\Person{阿奎拉}的希腊文\WorkTitle{旧约}译本与希伯来文原典，询问\Person{玛尔塞拉}的母亲，并附上此前两封信（第三十、三十一封）。公元三八四年写于罗马。}\par

1. 这封信简短有两个原因，一是送信人急于出发，二是因我忙于事务，无暇在琐事上耗费时间。你问有什么事如此紧迫，竟使我无法在纸上与你聊聊？那么，让我告诉你：最近一段时间，我一直在将\Person{阿奎拉}的\WorkTitle{旧约}希腊文译本与希伯来文卷轴进行比较，看看会堂是否出于对\Person{基督}的憎恨而更改了文本；——坦率地对一位朋友说——我发现了若干与我们的信德相符的差异。在精确修订了\WorkTitle{先知书}、\WorkTitle{所罗门书}、\WorkTitle{圣咏集}和\WorkTitle{列王纪}之后，我现在正致力于\WorkTitle{出谷纪}（犹太人按其开首词语称为\OriginalTerm{出谷纪}{Eleh shemôth}），完成后将继续\WorkTitle{肋未纪}。现在你明白我为何不能因写信的请求而放下工作了吧。不过，我不愿让我的朋友\Person{库伦提乌斯}完全白跑一趟，于是在这段闲聊中附上了两封信，是给你姐姐\Person{保拉}和她亲爱的孩子\Person{欧斯托基}的。请读一读，若你觉得有益或愉悦，便把我对她们所说的话也当作是为你说的。\par

2. 我希望\Person{阿尔比纳}，你和我的母亲，身体安康。我是指肉体健康，因为我不怀疑她灵性的福祉。请代我问候她，并以双重情感珍爱她，既作为基督徒，也作为母亲。\par

\SourceRecord{Translated by W.H. Fremantle, G. Lewis and W.G. Martley.；Nicene and Post-Nicene Fathers, Second Series；Vol. 6.；Edited by Philip Schaff and Henry Wace.；Buffalo, NY: Christian Literature Publishing Co.,；1893.；<http://www.newadvent.org/fathers/3001032.htm>.}

% structure: p0001 source=h1 target=section reason=page-title

\section{第三十三封信}

\Emph{致\Person{保拉}}\par

\EditorialNote{这是一封信的片段，其中热罗尼莫将M.T.瓦罗和奥力振作为作家的工作进行了比较。值得注意的是，它对奥力振给予了无条件的赞扬，这与作者后来几年所采用的语气形成了强烈对比（例如，参见第84封信）。其年代大概是公元384年。}\par

1. 古人对\Person{马尔库斯·泰伦提乌斯·瓦罗}惊叹不已，因为他为拉丁读者写了无数著作；希腊作家则对其铜人不吝赞美，因为此人作品之多，我们中间任何一人连抄写都抄不完。但既然拉丁人的耳朵会厌烦希腊著作的名单，我就只谈拉丁的\Person{瓦罗}。我将努力表明，我们今人正在像\Person{埃庇米尼得斯}那样沉睡，把前人用于追求扎实（即使属世俗）的学问的精力，都用来积累财富。\par

2. \Person{瓦罗}的著作包括四十五卷\WorkTitle{古事记}，四卷关于罗马人生活的书。\par

3. 但你会问我，为何我提起\Person{瓦罗}和那个铜人？只是为了向你介绍我们基督内的铜人，或者更确切地说，铁人——\Person{奥力振}，他对圣经研究的热情为他赢得了后一个称号。你想知道他留下了怎样的天才纪念碑吗？以下名单列举出来。他的著作包括十三卷\WorkTitle{创世记}注释、两卷\WorkTitle{神秘讲道集}、\WorkTitle{出谷记注释}、\WorkTitle{肋未记注释}、* * * * 还有单卷著作、四卷\WorkTitle{论首要原理}、两卷\WorkTitle{论复活}、关于同一主题的两篇对话。\par

* * * * * * * * * * *\par

4. 所以，你看，这一个人的劳作超过了之前所有希腊和拉丁作家。有谁曾读完他所有的作品？但他的努力给他带来了什么报酬呢？他的主教\Person{德梅特里}判定他有罪，只有\Place{巴勒斯坦}、\Place{阿拉伯}、\Place{腓尼基}和\Place{阿哈雅}的主教持异议。帝国罗马同意对他的谴责，甚至召集了一个议会来指责他——并非像如今追逼他的狂吠猎犬所喊的那样，是因为他的学说新奇或异端，而是因为人们无法容忍他那无与伦比的雄辩和学识，一旦他开口，其他人便相形见绌。\par

5. 我以上写得匆匆且不慎，仅借一盏简陋的灯笼之光。你自会明白为何如此，若你想到那些如今代表\Person{伊壁鸠鲁}和\Person{亚里斯提波}的人。\par

\SourceRecord{Translated by W.H. Fremantle, G. Lewis and W.G. Martley.；Nicene and Post-Nicene Fathers, Second Series；Vol. 6.；Edited by Philip Schaff and Henry Wace.；Buffalo, NY: Christian Literature Publishing Co.,；1893.；<http://www.newadvent.org/fathers/3001033.htm>.}

% structure: p0001 source=h1 target=section reason=page-title

\section{第三十四书信}

第三十四书信。致\Person{玛尔塞拉}。\par

\BibleRef{咏 126:2} ）。热罗尼莫在惋惜\Person{奥力振}对圣咏的注释已失传之后，给出了如下解释：\par

希伯来短语\Quote{忧苦之饼}被七十士译本译作\Quote{偶像之饼}；\Person{阿奎拉}译作\Quote{患难之饼}；\Person{西玛库斯}译作\Quote{悲惨之饼}；\Person{狄奥多田}跟随七十士译本。\WorkTitle{奥力振第五译本}也如此。\WorkTitle{第六译本}译作\Quote{错误之饼}。为了支持七十士译本，此处使用的词在\BibleRef{咏 113:12}中被译为\Quote{偶像}。或指生活的患难，或指异端的教义。\par

关于第二个短语，他作了更长的论述。在表明\Person{普瓦捷的依拉利}的观点（即所指的人是宗徒，他们被告知要抖掉脚上的尘土，\BibleRef{玛 10:14}）站不住脚，且将需要用\Quote{抖掉者}来替代\Quote{被抖掉的}之后，热罗尼莫如先前一样回到希伯来文，并宣称\Person{西玛库斯}和\Person{狄奥多田}的译文，即\Quote{青年之子}，才是真正的译文。他指出，七十士译本（拉丁译者被其误导）在\BibleRef{厄下 4:16}中出现了同样的错误。最后，他纠正了\Person{依拉利}在\BibleRef{咏 127:2}上的一个失误，在那里，由于误解七十士译本，后者用\Quote{你果实的劳苦}替代了\Quote{你双手的劳苦}。他自始至终对\Person{依拉利}深表敬意，并说他不懂希伯来文并非这位主教的过错。此信日期大概是公元三八四年。\par

\SourceRecord{Translated by W.H. Fremantle, G. Lewis and W.G. Martley.；Nicene and Post-Nicene Fathers, Second Series；Vol. 6.；Edited by Philip Schaff and Henry Wace.；Buffalo, NY: Christian Literature Publishing Co.,；1893.；<http://www.newadvent.org/fathers/3001034.htm>.}

% structure: p0001 source=h1 target=section reason=page-title

\section{书信第三十五}

\Emph{出自教宗达玛稣}\par

达玛稣向热罗尼莫提出五个问题，并请求关于这些问题的信息。它们是：\par

1. 这句话的意思是什么：\BibleQuote{凡击杀加音者，将受七倍的报复}？\BibleRef{创 4:5}\par

2. 如果天主创造了万有都是好的，为何祂向诺厄吩咐关于不洁的动物，并对伯多禄说：\BibleQuote{天主所洁净的，你不可称为污秽的}？\BibleRef{宗 10:15}\par

3. 如何使 \BibleQuote{到第四代，他们要再回到这里} \BibleRef{创 15:16} 与 \BibleQuote{到第五代，以色列子民从埃及地上来} \BibleRef{出 13:18}（七十贤士译本）相调和？\par

4. 为何亚巴郎接受割损作为他信德的印记？\BibleRef{罗 4:11}\par

5. 为何天主允许依撒格，一个义人且是天主所爱的，被雅各伯所欺骗？\EditorialNote{创世纪第27章} 写于罗马，公元三八四年。\par

\SourceRecord{Translated by W.H. Fremantle, G. Lewis and W.G. Martley.；Nicene and Post-Nicene Fathers, Second Series；Vol. 6.；Edited by Philip Schaff and Henry Wace.；Buffalo, NY: Christian Literature Publishing Co.,；1893.；<http://www.newadvent.org/fathers/3001035.htm>.}

% structure: p0001 source=h1 target=section reason=page-title

\section{第三十六封信}

\Emph{致教宗\Person{达玛苏斯}}\par

\EditorialNote{热罗尼莫对上述信件的回复。对于第二和第四个问题，他引荐达玛苏斯参阅戴尔都良、诺瓦替安和奥力振的著作。其余三个问题他作了详细处理。}\par

\BibleRef{创 4:15}，他理解为意指\Quote{杀害加音的人将完成在他身上所要施行的七倍复仇。}\par

\BibleRef{出 13:18}，他提议与\BibleRef{创 15:16}调和，假定一处指的是肋未支派，另一处指的是犹大支派。然而，他提出，由七十贤士译本译为\Quote{在第五代}的词更可能意为\Quote{武装的}（钦定本）或\Quote{装载的}。在回答关于依撒格的问题时，他说：\Quote{除了那为了我们的救恩而屈尊穿上血肉的那一位外，没有人拥有完全的知识和对真理的全面把握。\Person{保禄}、\Person{撒慕尔}、\Person{达味}、\Person{厄里叟}都犯过错误，圣人只知晓天主向他们启示的事。}接着，他继续对这段经文给出殉道者\Person{依玻里多}所建议的神秘诠释。写于前一封信的次日。\par

\SourceRecord{Translated by W.H. Fremantle, G. Lewis and W.G. Martley.；Nicene and Post-Nicene Fathers, Second Series；Vol. 6.；Edited by Philip Schaff and Henry Wace.；Buffalo, NY: Christian Literature Publishing Co.,；1893.；<http://www.newadvent.org/fathers/3001036.htm>.}

% structure: p0001 source=h1 target=section reason=page-title

\section{第三十七封书信}

\Emph{致玛尔切拉}\par

\EditorialNote{玛尔切拉曾请求热罗尼莫借给她一部勒提提乌斯（奥古斯托杜努姆［即奥顿］主教）所著\WorkTitle{雅歌}注释的抄本。热罗尼莫如今拒绝这样做，理由是这部著作充满错误，其中两个例子如下：（一）勒提提乌斯将塔尔史士等同于塔尔索；（二）他以为乌法兹（在\Quote{乌法兹的金子}一语中）与刻法相同。写于罗马，公元三八四年。}\par

\SourceRecord{Translated by W.H. Fremantle, G. Lewis and W.G. Martley.；Nicene and Post-Nicene Fathers, Second Series；Vol. 6.；Edited by Philip Schaff and Henry Wace.；Buffalo, NY: Christian Literature Publishing Co.,；1893.；<http://www.newadvent.org/fathers/3001037.htm>.}

% structure: p0001 source=h1 target=section reason=page-title

\section{第三十八封信}

\Emph{致玛尔塞拉}\par

\EditorialNote{\Person{布拉埃西拉}是\Person{保拉}之女、\Person{欧斯托基}之妹，结婚七个月后丧夫。一场危险疾病导致她皈依，如今她因苦修之严苛而闻名\Place{罗马}。许多人指责她，认为她是狂热分子，\Person{热罗尼莫}作为她的神师，也受到一些责难。在本信中，他为她的行为辩护，并宣称那些指责像她这样的人的生活的人，无权自称为基督徒。写于公元385年的\Place{罗马}。}\par

1. 当\Person{亚巴郎}受命献祭他的儿子时，这试探反而坚固了他的信德。\EditorialNote{创世纪第二十二章}当\Person{若瑟}被卖到埃及，他在那里的寄居使他能供养父亲和兄弟们。当\Person{希则克雅}因死亡临近而惊慌时，他的眼泪和祈祷为他赢得了十五年的缓期。如果宗徒\Person{伯多禄}的信德因主受难而动摇，那是为了让他痛哭之后，能听到安慰的话：\Quote{喂养我的羊群。}如果\Person{保禄}——那只贪婪的狼\BibleRef{创 49:27}、那个小\Person{本雅明}——在神魂超拔中被昏盲，那是为了让他复明，并被周围黑暗的突然恐怖所引导，称那位他先前当作人来迫害的为\Quote{主}。\BibleRef{宗 9:3}\par

2. 如今，我亲爱的\Person{玛尔塞拉}，我们心爱的\Person{布拉埃西拉}也是如此。我们见她几乎连续三十天遭受高烧折磨，这热病是为教导她放弃对那不久必被虫噬的肉身过分关注。主耶稣在她病中来到她身边，牵着她的手，看，她起来伺候祂。\BibleRef{谷 1:30}从前她的生活有些漫不经心，被财富的锁链紧紧捆绑，如同死人躺在世界的坟墓里。但耶稣愤慨不已，心中忧伤，大声呼喊说：\Quote{\Person{布拉埃西拉}，出来！}\BibleRef{若 11:38}她一听到呼唤就起来并走出来，与主一同坐席。\BibleRef{若 12:2}犹太人若愿意，可以发怒威胁她；他们可以设法杀害她，因为基督使她复活。\BibleRef{若 12:10}宗徒们将荣耀归于天主就够了。布拉埃西拉知道她的生命来自那位重新赐予她生命的主。她知道现在她能拥抱那位不久之前她还畏惧作为审判者的主的脚。\BibleRef{路 7:38}那时生命几乎离开了她的身体，死亡的临近使她喘息颤抖。那时她从亲戚那里得到了什么帮助？他们轻如烟的话语有何安慰？你们这些不友善的亲属，如今她对世界已死，对基督却活着，她不欠你们什么。\BibleRef{罗 6:11}基督徒应当为此欢喜，而恼怒的人必须承认他无权称为基督徒。\par

3. 一位\Quote{脱离了丈夫法律的}寡妇，她的唯一责任是恒守寡居。\BibleRef{罗 7:2}但你会说，素色的衣服令世界不悦。那么\Person{洗者若翰}也会令世界不悦；然而，在妇女所生者中，没有一位比若翰更大。\BibleRef{路 7:28}他被称作天使，他为上主施洗，但他身穿骆驼毛的衣服，腰束皮带。\BibleRef{玛 3:4}世界是否因为寡妇的饮食粗陋而不悦？没有比蝗虫更粗陋的，然而这些正是若翰的食物。应当令基督徒反感的妇女，是那些用胭脂和化妆品涂眼抹唇的；那些脸上涂得像偶像一样白得反常的；每一滴偶然的眼泪都会在她们脸颊上留下痕迹的；她们意识不到岁月使她们衰老的；她们用别人的头发堆在头上的；她们抚平面孔，擦去皱纹的；她们在孙辈面前表现得像发抖的女学生的。基督徒妇女应以对自然施加暴力为耻，或以关心肉体来刺激欲望为耻。宗徒告诉我们：\Quote{随从肉性的人，不能使天主喜悦。}\BibleRef{罗 8:8}\par

4. 昔日，我们亲爱的寡妇在着装极其挑剔，整日站在镜前纠正缺陷。如今她勇敢地说：\Quote{我们众人脸面揭开，如同镜子反映主的荣光，就变成同样的形象，光荣上加光荣，正如从主的神来的。}那时侍女为她梳头，她那无罪的脑袋被硬梳成波浪状。如今她任头发自然，唯一的头饰是面纱。那时最柔软的羽绒床对她都嫌硬，她在一堆床垫上也难以安眠。如今她热切起身祈祷，尖声高喊\Quote{阿肋路亚}先于他人，她是第一个赞美主的人。她跪在光秃的地上，以频频的泪水洗刷曾被铅粉玷污的脸庞。祈祷后是咏唱圣咏，只有当她的脖子酸痛、膝盖颤抖、眼睛因疲倦欲闭时，才不情愿地允许自己休息。因为她的衣服是深色的，躺在地上也不会弄脏。便宜的鞋子使她能把镀金鞋的钱施舍给穷人。她的腰带没有金子和宝石装饰，是羊毛做的，朴素而极其干净。它旨在保持衣服整齐，而不是把腰勒成两段。因此，如果那蝎子斜眼看她的决心，用甜言蜜语诱惑她再次吃禁树之果，她必须咒骂着把它踩在脚下，并说，当它躺在自己被定的尘土中死去时：\BibleRef{创 3:14}\Quote{退到我后面去，\Person{撒殚}！}\BibleRef{玛 16:23}撒殚意思是敌对者，\BibleRef{伯前 5:8}而一个不喜欢基督诫命的人，不仅是基督的敌对者，他是假基督。\par

5. 但我请问你，我们究竟做了什么使人对我们反感？我们曾效法过宗徒吗？我们知道，首批门徒舍弃了他们的船和网，甚至年迈的父亲。\BibleRef{玛 4:18} 收税人从税关站起来，永远跟随了救主。\BibleRef{玛 9:9} 当一个门徒想回家告别亲属时，师傅的声音不允许。\BibleRef{路 9:61} 一个儿子甚至被禁止埋葬他的父亲，\BibleRef{玛 8:21} 仿佛表明为了主的缘故有时不尽孝道也成了一种宗教义务。\BibleRef{路 14:26} 我们却不然。如果我们拒绝穿丝绸，就被认为是隐修士；如果我们保持清醒并避免过度欢笑，就被称为酸腐而严厉。如果我们的短衣干净整洁，群众就称我们为希腊人和骗子。如果他们乐意，他们还可以开更刻薄的玩笑；他们可以到处炫耀他们能找到的每一个肥肚皮来嘲弄我们。我们的\Person{布拉埃西拉}会对他们的努力付之一笑，并耐心忍受所有这些呱呱叫的青蛙的嘲讽，因为她记得人们称她的主为贝耳则步。\BibleRef{玛 10:25}\par

\SourceRecord{Translated by W.H. Fremantle, G. Lewis and W.G. Martley.；Nicene and Post-Nicene Fathers, Second Series；Vol. 6.；Edited by Philip Schaff and Henry Wace.；Buffalo, NY: Christian Literature Publishing Co.,；1893.；<http://www.newadvent.org/fathers/3001038.htm>.}

% structure: p0001 source=h1 target=section reason=page-title

\section{书信三十九}

\Emph{致保拉}\par

\EditorialNote{布拉西拉皈依不到三个月就去世了，热罗尼莫现在写信给保拉，向她表示同情，并尽可能减轻她的悲伤。他请她记住，布拉西拉现在已在天堂，因此要控制自己，以免信仰的敌人对她的行为吹毛求疵。然后他以预言（后来已应验有余）结束：在他的著作中，布拉西拉的名字将永垂不朽。写于公元389年的罗马。}\par

1. \Quote{但愿我的头是水，我的眼是泪泉，好叫我哭泣，}不像耶肋米亚所说的，\BibleQuote{为我百姓的伤损，}\BibleRef{耶 9:1}也不像耶稣，为耶路撒冷的悲惨命运，\BibleRef{路 19:41}而是为圣德、仁慈、无邪、贞洁，以及一切美德，因为布拉西拉一死，这一切都消失了。我并不为她悲伤，而是必须为自己悲伤；我的损失太大了，无法以顺服承受。谁能带着干眼回忆那炽热的信德——它使一个二十岁的少女举起十字架的旗帜，哀悼失贞胜于哀悼丈夫的死？谁能不叹息地回忆她祈祷的热忱、谈吐的光彩、记忆的牢固和才智的敏捷？你若听她说希腊语，会以为她不懂拉丁语；然而当她用罗马语言时，她的言语毫无外国口音。她在那些为\Person{奥利振}赢得希腊钦佩的造诣上，甚至可与伟大的\Person{奥利振}媲美。因为短短几个月，更确切说几天之内，她便完全掌握了希伯来语的难点，以致在学习与歌唱圣咏上可与母亲的热情匹敌。她的衣着朴素，但这朴素并非通常那样是骄傲的标志。事实上，她的谦卑如此完美，穿着不优于婢女，唯一区别在于她步履更加自在。她的步履因虚弱而蹒跚，面容苍白颤抖，纤弱的脖颈几乎支撑不住头。但她手中始终拿着一本先知书或福音书。一想到她，我的眼就充满泪水，哽咽阻断了我的声音，我的情绪如此激动，以致舌头顶住上颚。当她躺在那儿奄奄一息，可怜的躯体因高烧而枯焦，亲属围在床边，她临终的话是：\Quote{请向主耶稣祈祷，求祂宽恕我，因为我本想做的事，却没有能力做到。}亲爱的布拉西拉，请安心，确信你的衣裳常是洁白的。\BibleRef{训 9:8}因为你有永恒童贞的纯洁。我坚信我的话是真的：皈依永远不嫌晚。耶稣对临终强盗的话就是保证：\BibleQuote{我实在告诉你，今天你就要同我在乐园里。}\BibleRef{路 23:43}最后，当她的灵魂从肉体负担中解脱，回到赐予它的天主那里；\BibleRef{训 12:7}还有，当她结束漫长的朝圣之旅，升入她古老的家园时，她的葬礼按惯例隆重举行。显贵人物走在队列前列，一块金线织成的棺罩覆盖着她的灵柩。但我似乎听到来自天上的声音说：\Quote{我不认得这些装饰；这不是我惯常穿的衣裳；这种豪华对我来说是陌生的。}\par

2. 但这是什么呢？我想阻止一位母亲的哭泣，而我自己也呻吟。我不掩饰自己的情感；这整封信是流泪写成的。连\Person{耶稣}也为\Person{拉匝禄}流泪，因为他爱他。\BibleRef{若 11:35}但他是一个可怜的安慰者，他自己也被叹息所战胜，从他的忧苦之心既挤出眼泪也挤出言语。亲爱的\Person{保拉}，我的痛苦和你的一样大。\Person{耶稣}知道此事，\Person{布拉西拉}现在跟随的正是他；圣天使也知道此事，她现在与他们为伴。我在精神上是她的父亲，在情感上是她的养父。有时我说：\BibleQuote{愿我生的那日灭亡，}又说：\BibleQuote{我有祸了，我的母亲，你生了我，使我成为全地纷争和相争的人。}\BibleRef{耶 15:10}我呼喊：\BibleQuote{上主，你是公义的……但我愿与你讨论你的审判：为何恶人的道路亨通？}\BibleRef{耶 12:1}\BibleQuote{至于我，我的脚几乎滑倒，我的脚步险些滑脱。因为我嫉妒愚昧人，看见恶人享平安，便说：天主如何知道？至高者岂有知识？看，这些就是世上亨通的恶人，他们积蓄财富。}但我想起另一些话：\BibleQuote{若我说要这样讲，看，我就得罪了你儿女的世代。}难道疑惑的巨浪不也像冲击你一样冲击我的灵魂吗？我问：为何不虔敬的人长寿享受世上的财富？为何未受教训的青少年和无辜的孩童在蓓蕾时就被剪除？为何三岁或两岁的儿童，甚至未断奶的婴孩，被魔鬼附身，长满癞病，被黄疸吞噬，而不虔敬和亵渎的人、奸夫和凶手却有健康和力量来亵渎天主？难道我们没有被告诉：父亲的罪不落在儿子身上\BibleRef{则 18:20}，以及\BibleQuote{犯罪的人必死亡}\BibleRef{则 18:4}吗？如果旧的教训仍然有效，即父亲的罪要报应在儿女身上\BibleRef{出 20:5}，那么一个老人无数的罪不能公平地报复在一个无辜的婴儿身上。我说：\BibleQuote{我真是徒然洁净了我的心，徒然洗手表示无辜，因为我终日遭灾。}但当我思想这些事时，我像先知一样学会说：\BibleQuote{我思索要明白这事，眼看实是为难；直到我进了天主的圣所，才明白他们的终局。}的确，上主的判断如同深渊。\BibleQuote{啊，天主的智慧和知识的财富何其深！他的判断何其不可测，他的道路何其难寻！}\BibleRef{罗 11:33}天主是良善的，他所做的一切也必是良善的。他命定我必须失去我的丈夫吗？我哀悼我的损失，但因为这是他的旨意，我顺从地忍受。我独生的儿子被夺去吗？打击是沉重的，却仍可承受，因为取走他的就是那赐予他的。\BibleRef{约 1:21}如果我失明，朋友的诵读将安慰我。如果我失聪，我将避开罪恶的话语，我的思想将唯独思念天主。如果除了这些试炼之外，还有贫穷、寒冷、疾病和赤身压迫我，我将等待死亡，视它们为暂时的痛苦，很快将让位于更好的结局。让我们思量智慧诗篇的话：\BibleQuote{上主，你是公义的，你的判断是正直的。}唯有那在一切困苦中赞美主，并将自己的痛苦归咎于自己的罪，感谢天主的仁慈的人，才能这样说。\par

犹大女儿们，我们被告知，因上主的一切判断而欢乐。因此，既然犹大意谓宣认，既然每个信靠的灵魂宣认其信德\BibleRef{罗 10:10}，凡声称信基督的人，必须在基督的一切判断中喜乐。我健康吗？我感谢我的造物主。我生病吗？在此情形下，我也赞美天主的旨意。因为\BibleQuote{当我软弱时，我正是刚强的；}而精神的力量在肉身的软弱中得以成全。连宗徒也必须忍受他所不喜欢的，那个他三次求主使它离开的病痛。天主的回答是：\BibleQuote{我的恩宠为你够了，因为我的能力在软弱中得以成全。}唯恐他因启示而过于自高，就赐给他一个关于人性软弱的提醒，正如在得胜将军的凯旋战车上，总有一个奴隶在民众的欢呼中不断低语：\Quote{记住，你不过是人。}\par

3. 但是，为什么我们必须自己有一天要忍受的事，现在却觉得难以忍受呢？又为何为死者悲伤？我们生来并非永生不死。\Person{亚巴郎}、\Person{梅瑟}、\Person{依撒意亚}、\Person{伯多禄}、\Person{雅各伯}、\Person{若望}、\Person{保禄}——那\BibleQuote{特选的器皿}\BibleRef{宗 9:15}——甚至连天主圣子自己都死了；那么当一个灵魂离开地上的居所时，我们岂能烦恼呢？或许他被接走，是\Quote{免得邪恶改变他的心意……因为他的灵魂蒙主喜悦，所以主急忙将他从众人中接去}——免得在人生漫长旅途中迷失于无路迷途。我们确实应为死者哀悼，但只为那被\OriginalTerm{革哈拿}{Gehenna}收纳、被\OriginalTerm{塔耳塔鲁斯}{Tartarus}吞噬、且永恒之火为其燃烧的人哀悼。但我们这些在离世时有天使护送、并被基督亲自迎接的人，反而应哀伤自己必须在这死亡的帐幕中停留更久。\BibleRef{格后 5:4}因为\BibleQuote{我们住在肉身内，便离主远}\BibleRef{格后 5:6}。我们唯一的渴望应如圣咏作者所言：\Quote{我哀叹我的漂泊岁月延长，我曾与\OriginalTerm{刻达尔}{Kedar}的居民同住，我的灵魂作了远行。}刻达尔意为黑暗，黑暗代表现世（因为经上告诉我们：\BibleQuote{光在黑暗中照耀，黑暗却没有胜过它}\BibleRef{若 1:5}）。因此，我们应为亲爱的\Person{布莱西拉}祝贺，她已从黑暗进入光明\BibleRef{弗 5:8}，并在她初露晨曦的信德之初，就获得了完成工作的华冠。若她是在满心世俗欲望与短暂享乐之时被剪除（我祈求无人如此），那么哀悼对她确是该有的，为她流多少泪也不为过。实则，因基督的仁慈，她四个月前在守寡誓愿中更新了她的圣洗，在余下的日子里摒弃了世界，只思虑修道生活。那么，你岂不害怕，救主会对你说：\Quote{你愤怒吗，\Person{帕乌拉}，因为你的女儿成了我的女儿？你恼恨我的命令，以悖逆的泪水嫉妒我拥有\Person{布莱西拉}吗？你应当知道我对你和你家人旨意。你禁绝食物，不是为了守斋，而是为了满足忧伤；这样的节制不悦于我。这样的斋戒是我的仇敌。我不收纳任何违抗我意志而离开肉身的灵魂。愚昧的哲学或可炫耀这样的殉道者，或可炫耀\Person{芝诺}、\Person{克莱欧布洛托斯}或\Person{加图}。我的神安息在\InnerQuote{他是贫苦而心灵痛悔、对我的话战栗的人}\BibleRef{依 66:2}。这就是你向我许愿要过修道生活的意义吗？你为此穿得与其他妇人不同，以修女装束打扮自己吗？哀悼属于那些穿丝绸衣服的人。在你泪水中，召叫将临，你也要死；然而你逃避我，如同逃避残酷的审判官，幻想你能避免落在我手中。那个刚愎的先知\Person{约纳}曾逃避我，但在深海他仍属我。\BibleRef{纳 2:2}若你真相信你的女儿活着，你就不会为她去往更美好的世界而忧伤。这是我藉我宗徒给你们命令，叫你们不要为安眠的人忧伤，像没有望德的外邦人一样。\BibleRef{得前 4:13}羞愧吧，因你被外邦人的榜样所羞辱。魔鬼的婢女比我的更好。因为她想象她的不信丈夫被提入天堂，而你或是不能或是不愿相信你的女儿安息在我这里。}\par

4. 你说：我为何不该哀哭？\Person{雅各伯}为\Person{若瑟}披上麻衣，当他全家聚集时，不肯受安慰。他说：\BibleQuote{我要哀哭着下到阴间到我儿子那里去。}\BibleRef{创 37:35}\Person{达味}也为\Person{阿贝沙隆}哀哭，掩面喊道：\BibleQuote{我儿\Person{阿贝沙隆}……我儿\Person{阿贝沙隆}！但愿我替你死了，\Person{阿贝沙隆}我儿！}\BibleRef{撒下 18:33}\Person{梅瑟}\BibleRef{申 34:8}和\Person{亚郎}\BibleRef{户 20:29}及其他圣人也以隆重的哀悼被悼念。对你的推理，答案很简单。的确，\Person{雅各伯}哀悼\Person{若瑟}，以为他被杀了，以为只能在阴间相遇（他说：\BibleQuote{我要哀哭着下到阴间到我儿子那里去}），但他这样做，只是因为基督尚未打开天堂之门，未以祂的血熄灭闪烁的剑和守卫\OriginalTerm{革鲁宾}{cherubim}的旋转（因此在富翁与拉匝禄的故事中，\Person{亚巴郎}和乞丐虽实在安息之所，却被描述为在阴间）。而\Person{达味}，在无效地为婴孩生命代祷后，拒绝为他哭泣，因为知道他没有犯罪，却为一个在意志（若非行为）上弑父的儿子哭泣，这做得对。\BibleRef{撒下 17:1}当我们读到\Person{梅瑟}和\Person{亚郎}按照古例受哀悼时，这不应使我们惊讶，因为即使在宗徒大事录、在福音的完全光明中，我们看见耶路撒冷的弟兄们为\Person{斯德望}大大哀哭。\BibleRef{宗 8:2}但这极大的哀哭并非指哀悼者，而是指送葬行列和随行人群。这就是圣经关于\Person{雅各伯}所说的：\BibleQuote{\Person{若瑟}上去埋葬他的父亲，与他同去的有法郎的一切臣仆、他家里的长老、埃及国的一切长老、\Person{若瑟}全家和他的兄弟们}；数行之后：\BibleQuote{又有战车和骑兵同他上去：这是一大队人。}最后：\BibleQuote{他们大大哀悼，非常悲痛。}这隆重的哀悼并没有把长期的哭泣强加给埃及人，而只是描述葬礼仪式。同样，当我们读到为\Person{梅瑟}和\Person{亚郎}哭泣时，所指即是如此。\par

我不能充分颂扬圣经的奥秘，也无法足够钦佩其最朴素言辞所传达的属灵意义。例如，我们被告知，曾为\Person{梅瑟}哀哭；然而，当描述\Person{若苏厄}的葬礼时，\BibleRef{苏 24:30} 却未提及哭泣。原因当然是，在\Person{梅瑟}之下——即在旧约法律之下——所有的人都受制于对\Person{亚当}之罪的判决，当他们进入阴间时，理所当然伴随着眼泪。因为，如宗徒所说：\BibleQuote{从\Person{亚当}到\Person{梅瑟}，死亡统治了那些没有犯罪的人。} \BibleRef{罗 5:14} 但在\Person{耶稣}之下，即在\Person{基督}的福音之下，祂为我们开启了乐园之门，死亡伴随的不是悲伤，而是喜乐。犹太人至今仍在哭泣；他们赤脚，披麻蒙灰，蜷缩在麻布中，在灰中打滚。为使他们的迷信完整，他们遵循法利塞人的愚昧习俗，吃扁豆，似乎是为了显示他们因何等微薄的食物而失去了长子名分。\BibleRef{创 25:34} 当然，他们哭泣是正确的，因为他们不信主的复活，正在为假基督的来临做准备。但我们这些已穿上\Person{基督}的，并如宗徒所说，是王家的司祭民族，\BibleRef{迦 3:27} \BibleRef{伯前 2:9} 我们不应当为死者悲伤。圣经告诉我们：\BibleQuote{\Person{梅瑟}对\Person{亚郎}及\Person{厄肋阿匝尔}和\Person{依塔玛尔}，他剩下的儿子们说：\InnerQuote{你们不可披头散发，不可撕裂衣服，免得你们死亡，也免得向全会众发怒。}} \BibleRef{肋 10:6} 他说：\Quote{不可撕裂衣服，也不可像外邦人那样哀悼，免得你们死亡。} 因为，对我们来说，罪就是死亡。在同一部书\WorkTitle{肋未纪}中，有一项条文，或许有人会视为残忍，但对信德却是必要的：大司祭被禁止接近他父亲、母亲、兄弟和儿女的尸体；\BibleRef{肋 21:10} 为的是，任何悲伤都不可分散那向上主献祭、全心奉献于神圣奥迹的灵魂。我们在福音中是否也以其他言语学到同样的教训？不是禁止门徒与家告别或埋葬亡父吗？\BibleRef{路 9:59} 关于大司祭，又说道：\BibleQuote{他不可走出圣所，不可玷污天主的祝圣，因为天主的傅油临于他。} 确实，如今我们相信了\Person{基督}，并将祂怀在我们内，借着我们领受的傅油之油，\BibleRef{若一 2:27} 我们不应当离开祂的殿宇——即我们的基督徒信仰——不应当出去与不信的外邦人混杂，而应常留在殿内，作为顺服主旨的仆人。\par

我已经直言不讳，免得你无知地以为圣经认可你的悲伤；免得你若是错了，还自以为有理。至今，我的话是针对一般基督徒妇女说的。但现在不同了。因为在你身上，据我所知，弃绝世界是彻底的；你拒绝并践踏了生活的享乐，每天致力于禁食、阅读和祈祷。像\Person{亚巴郎}一样，\BibleRef{创 12:1} 你渴望离开你的本乡和家族，舍弃美索不达米亚和加色丁人，进入应许之地。你在死亡之前已向世界死去，你将所有世俗财富施舍给了穷人，或赐予了你的子女。因此，我更加惊讶，你竟会做出在他人身上本应受谴责的行为。你怀念\Person{布拉色拉}的陪伴、她的谈吐和她的可爱之处；你无法忍受自己已失去这一切的想法。我宽恕你母亲的眼泪，但我请求你克制你的悲伤。当我想到作为母亲时，我不怪你哭泣；但当我想到作为基督徒和隐修女时，母亲就从我眼前消失了。你的伤口尚新，无论我多么轻柔的碰触，都更可能激化而非愈合。然而，你为何不尝试用理性来克服那种时间终必缓解的悲伤呢？\Person{纳敖米}因饥荒逃往\Place{摩阿布}地，在那里失去了丈夫和儿子。然而，当她失去这些天然保护者时，一个外邦女子\Person{卢德}从未离开她身边。\EditorialNote{卢德传第一章} 你看，安慰一个孤独的妇女是何等伟大的事！\Person{卢德}作为报偿，成了\Person{基督}的祖先。\BibleRef{玛 1:5} 思考\Person{约伯}所承受的大试炼，你就会明白自己过于娇弱。在他的家宅倒塌、疮痍痛苦、无数丧亲、以及最后妻子的引诱中，他仍然举目向天，保持了耐心的坚忍。我知道你要说什么：\Quote{这一切临到他，是因为他是个义人，为要考验他的义德。} 好吧，你任选其一。要么你是圣洁的，那么天主是在证明你的圣德；要么你是个罪人，那么你无权抱怨。因为若是如此，你所受的远不及你应得的。\par

我何必重复老故事？且听一个现代的实例。圣\Person{默拉尼}，她在基督徒中因真正的高贵而卓著（愿主使你和我在祂的日子里有分于她！），在她丈夫的尸体尚未埋葬、尚温时，又不幸在同一时刻失去了两个儿子。后续令人难以置信，但\Person{基督}是我的证人，我所言为真。你会不会以为她在疯狂中会解开头发、撕裂衣服、捶打胸膛？然而，她眼中没有落下一滴眼泪。她一动不动地站在那里；然后俯伏在\Person{基督}脚下，微笑着，仿佛用双手抓住了祂。她说：\Quote{主啊，从今以后，我要更乐意地事奉祢，因为祢已将我从一个重担中释放出来。} 但也许她剩下的孩子战胜了她的决心。不，实在，她那样不看重他们，以至于将她所有的一切给了她唯一的儿子，然后，不顾即将到来的冬天，乘船去了\Place{耶路撒冷}。\par

6. 请祢宽恕自己，我恳求祢，宽恕\Person{布拉埃西拉}，她如今与\Person{基督}一同为王；至少宽恕\Person{欧斯托基雅}，她的年幼和缺乏经验依赖于祢的指导和教诲。如今魔鬼咆哮并抱怨他被轻视，因为他看到祢的一个孩子被高举在凯旋中。他在已离去的那位身上未能取得的胜利，他希望从仍存留的那位身上获得。对子女过度的爱，就是对\Person{天主}的冷淡。\Person{亚巴郎}乐意准备杀死他的独生子，而祢却抱怨若几个孩子中的一个获得了她的冠冕？我不能不叹息地说出我要说的话。当祢被抬出殡仪队伍时晕厥过去，人群中可听到这样的窃窃私语：\Quote{这岂不是我们常说的吗？她为她的女儿哭泣，因禁食而致死。她想要她再婚，好能有孙子孙女。我们还要忍多久才不驱逐这些可憎的隐修士出\Place{罗马}？我们为什么不拿石头砸他们或把他们扔进\Place{台伯河}？他们误导了这位不幸的女士；她不是出于选择而做修女，这是显然的。没有一个异教母亲为她的子女哭泣像她为\Person{布拉埃西拉}那样。} 你想，当\Person{基督}听到这样的话，祂会承受多大的悲伤？\Person{撒殚}又是多么得意洋洋，他急于夺取祢的灵魂！他以看似自然和合理的悲伤的诉求引诱祢，并始终将\Person{布拉埃西拉}的形象放在祢眼前，他的目的是杀死这位得胜者的母亲，然后袭击她被遗弃的妹妹。我这样说不是要吓唬祢。\Person{主}是我的见证，我现在对祢说话，如同我站在祂的审判座前。无意义的眼泪是可憎的。祢的眼泪是可憎的，亵圣的眼泪，不信的眼泪；因为它们不知止境，将祢带到死亡的边缘。祢尖叫呼喊，如同内心着火，尽力结束自己。但\Person{耶稣}以祂的慈悲来到祢和像祢一样的人面前，说：\BibleQuote{为什么哭泣？女孩子没有死，只是睡着了。} \BibleRef{谷 5:39} 旁观者可能嘲笑祂；这样的不信是\Person{犹太人}当得的。如果祢因悲伤俯伏在祢女儿的坟墓前，祢也将听到天使的责备：\BibleQuote{为什么在死人中寻找活人呢？} \BibleRef{路 24:5} 因为\Person{玛利亚玛达肋纳}正是这样做了，当她认出\Person{主}的声音呼唤她并俯伏在祂脚前时，祂对她说：\BibleQuote{不要摸我，因为我还没有升到我的父那里去；} \BibleRef{若 20:17} 这就是说，祢不配触摸那位复活者，因为祢还以为她仍在坟墓中。\par

7. 你想，我们的\Person{布拉埃西拉}要承受多大的十字架和折磨，当她看到\Person{天主}对祢发怒——即使只是一点点！此刻当祢哭泣时，她向祢呼喊：\Quote{若祢曾爱过我，母亲，若我曾吃祢的奶，若我曾受祢教导，请不要嫉妒我的高升，不要这样做使我们永远分离。祢以为我孤身一人吗？代替祢，我如今有了\Person{玛利亚}，\Person{主}的母亲。这里我看到许多从前不认识的人。我的同伴远比在地上的一切好。这里我有福音中的女先知\Person{亚纳}的陪伴；\BibleRef{路 2:36} 而且——什么能燃起祢更热烈的喜乐呢——我在短短三个月内获得了她多年劳苦才赢得的。我们两人确实都是寡妇，都得到了贞洁的棕榈枝。祢可怜我离开了世界？该是我可怜祢，也确实可怜祢，祢仍被囚在世界的监狱中，天天与愤怒、贪财、情欲以及种种引诱灵魂走向毁灭的试探斗争。若祢真想作我的母亲，必须中悦\Person{基督}。不中悦我\Person{主}的，就不是我的母亲。} 她还说了很多别的话，此处略过；她也为祢向天主祈祷。我也确信，她为我转求，求天主赦免我的罪，以报答我为她得救而冒犯她家人的警告和劝告。\par

8. 因此，只要气息仍在我身体内活动，只要我仍在享受生命，我承诺、宣告并保证，\Person{布拉埃西拉}的名字将永远在我口中，我的劳苦将奉献给她的光荣，我的才华将用于赞美她。我写的每一页都将有\Person{布拉埃西拉}的名字。我的言论所到之处，她也将随我拙劣的著作同行。童贞女、寡妇、隐修士和司铎在阅读时，将看到她在我心中的形象多么深刻。永恒的怀念将补足她短暂的生命。她与\Person{基督}一同活在天堂，也将在人的口舌上活着。现在很快会过去，让位给未来；那未来将不偏不倚、毫无偏见地判断她。作为一个无子女的寡妇，她将居于有子女的母亲\Person{保拉}和童贞女\Person{欧斯托基雅}之间。在我的著作中，她永不死亡。她将听见我常与她的妹妹或母亲谈论她。\par

\SourceRecord{Translated by W.H. Fremantle, G. Lewis and W.G. Martley.；Nicene and Post-Nicene Fathers, Second Series；Vol. 6.；Edited by Philip Schaff and Henry Wace.；Buffalo, NY: Christian Literature Publishing Co.,；1893.；<http://www.newadvent.org/fathers/3001039.htm>.}

% structure: p0001 source=h1 target=section reason=page-title

\section{\WorkTitle{第四十封信}}

\Emph{致\Person{玛尔塞拉}}\par

\EditorialNote{塞杰斯塔的奥纳苏斯，即本信所述之人，乃是耶柔米在罗马的反对者之一。此处他被以一种有损作者文雅的方式加以嘲讽。信件的日期为公元385年。}\par

1. 被称为外科医生的医学人士被认为是残忍的，但其实值得同情。因为用无情的刀割去他人坏死的肉而不为他的痛苦所动，这难道不可悲吗？治疗病人的医生对他人的痛苦麻木不仁，不得不显得像他的敌人，这难道不可悲吗？然而自然秩序正是如此。真理总是苦涩的，而欢愉等待行恶。\Person{依撒意亚}赤身露体而毫不羞愧，作为将来被掳的预兆。\BibleRef{依 20:2} \Person{耶肋米亚}从\Place{耶路撒冷}被派往\Place{幼发拉底河}（\Place{美索不达米亚}的一条河流），将他的腰带留在\Place{加色丁}营中，在与他的人民为敌的\Place{亚述人}中间任其朽坏。\BibleRef{耶 13:6} \Person{厄则克耳}奉命吃由杂粮做成、撒上人和牲畜粪的饼。\BibleRef{则 4:9} 他不得不看着自己的妻子死去而不得流泪。\BibleRef{则 24:15} \Person{亚毛斯}被赶出\Place{撒玛黎雅}。\BibleRef{亚 7:12} 他为何被赶出？无疑，此事与其他事例一样，因为他是一位属灵的外科医生，切除被罪感染的部位，敦促人悔改。宗徒\Person{保禄}说：\BibleQuote{因为我给你们说实话，就成了你们的仇敌吗？}\BibleRef{迦 4:16} \Person{救主}自己也经历了同样的事，许多门徒因祂的话难懂而退去了。\BibleRef{若 6:60}\par

2. 因此，我因揭发他们的过错而冒犯了许多人，这并不奇怪。我打算动手切除一个癌变的鼻子；让患肿块的人战栗吧。我想斥责一只聒噪的寒鸦；让乌鸦明白她是令人讨厌的。然而，在\Place{罗马}难道只有一个人\par

谁是鼻孔被疤痕毁容的？\par

难道只有\Place{塞杰斯塔}的\Person{奥纳苏斯}一人，将两颊鼓得像气囊，在舌上平衡空洞的言辞？\par

我说某些人凭借犯罪、伪证和虚假托辞，达到了这样或那样的高位。这与你何干？你知道这指控并不涉及你。我嘲笑一个没有委托人的辩护人，讥讽一个卖文为生的油滑口才。这与你何干？你是一个如此文雅的演说家。我突发奇想要抨击唯利是图的司铎。你已经富有，何必生气？我想把\Person{伏尔甘}关起来，用他自己的火焰烧死他。你是他的客人还是邻居，竟想从火中救出偶像的神龛？我选择嘲笑幽灵、猫头鹰和\Place{尼罗河}怪物；而无论我说什么，你都当作是针对你。无论我指笔于何项过错，你都大叫你被点中。你揪住我，把我拖进法庭，荒谬地指控我写讽刺诗，而我写的只是平实的散文！\par

所以你真的以为自己是个漂亮人物，只因为有一个幸运的名字！但事实根本不是这样。破片被称作破片，正是因为光线不能穿透它。命运女神被称作\Quote{宽恕者}，恰恰是因为她们从不宽恕。复仇女神被称作\Quote{仁慈的}，正因为她们不显仁慈。而在日常言谈中，埃塞俄比亚人被称为\Quote{银人}。不过，如果揭发缺点总让你生气，我现在用\Person{佩尔西乌斯}的话来抚慰你：\Quote{愿你成为我主和小姐女儿的猎物！愿漂亮女子们争抢你！愿你脚下的土地变成玫瑰床！}\par

3. 尽管如此，我还是给你一个提示，如果你想显得最好看，该隐藏哪些特征。脸上不要露出鼻子，嘴巴闭上。这样你才有希望被认为既英俊又雄辩。\par

\SourceRecord{Translated by W.H. Fremantle, G. Lewis and W.G. Martley.；Nicene and Post-Nicene Fathers, Second Series；Vol. 6.；Edited by Philip Schaff and Henry Wace.；Buffalo, NY: Christian Literature Publishing Co.,；1893.；<http://www.newadvent.org/fathers/3001040.htm>.}

% structure: p0001 source=h1 target=section reason=page-title

\section{\WorkTitle{书信四十一}}

\Emph{致\Person{玛尔克拉}}\par

\EditorialNote{曾有人试图劝玛尔克拉皈依蒙丹派，热罗尼莫在此为她总结蒙丹派的主要教义，并与教会的教义相对照。写于公元385年，罗马。}\par

1. 关于有人从\BibleBook{若望}{福音}中搜集、用以攻击你的那些段落——即我们的救主应许祂将往父那里去，并要派遣护慰者——关于这些，\BibleBook{宗徒}{大事录}告诉我们这些应许是在何时作出，又是在何时实现在了。我们得知，从主升天算起过了十天，从祂复活算起过了五十天，那时圣神降临，信徒们的舌头分开了，以致各人都说各种语言。那时，一些尚未相信的人声称门徒们被新酒灌醉了，伯多禄却站在宗徒和众人当中说：\BibleQuote{犹太人和一切住在耶路撒冷的人哪，这件事你们应当知道，也应当侧耳听我的话。这些人不是醉了，像你们所想的那样，因为现在只是白天第三时辰。但这正是\Person{岳厄尔}先知所说的：\InnerQuote{在末日，天主说：我要将我的神倾注在一切有血肉的人身上，你们的儿子和女儿要说预言，你们的青年要见异象，你们的老人要作梦；在那些日子里，我要将我的神倾注在我的仆人和我的婢女身上。}}\BibleRef{宗 2:14}\par

2. 既然如此，宗徒\Person{伯多禄}——主在他身上建立了教会\BibleRef{玛 16:18}——已明确说明主的预言和应许在当时当地已得到应验，我们岂能再为自己要求另一种应验呢？如果蒙丹派回答说，后来\Person{斐理伯}的四个女儿说了预言\BibleRef{宗 21:9}，又提到一位名叫\Person{阿加波}的先知，并且在对神恩的分配中，除了宗徒、教师等之外，也提到了先知，而且\Person{保禄}本人也曾预言过许多关于未来的异端和世界末日的事；那么我们告诉他们，我们并非完全拒绝预言——因为主的苦难已经证实了预言——而是拒绝接受那些言论与新旧圣经不相符合的先知。\par

3. 首先，我们在信德规矩上与蒙丹派不同。我们区分圣父、圣子和圣神为三个位格，但将他们联为一体。他们却追随\Person{撒伯里乌}的学说，将天主圣三压缩到单独一个位格的狭窄界限内。我们虽然不鼓励第二次婚姻，但予以允许，因为\Person{保禄}吩咐年轻的寡妇再嫁。\BibleRef{弟前 5:14} 他们却认为重复婚姻是极其严重的罪，犯此罪者应被视为犯奸淫者。我们按照宗徒传统（全世界都与我们一致）每年守一个四旬期的斋戒；而他们每年守三个，仿佛三位救主受了苦。当然，我并非说在一年中其他时间（始终排除五旬节期）守斋是不合法的，只是说在四旬期内是义务，在其他时间则是选择。此外，在我们这里，主教位居宗徒之职；在他们那里，主教不是第一，而是第三。因为他们将\Place{夫黎基雅}的\Place{佩普撒}的宗主教放在首位，其次是将那些被称为管家的执事，主教则被降至第三或几乎最低的地位。无疑，他们的目的是使他们的宗教显得更自命不凡，将我们置诸首位者置于末位。此外，他们几乎对一切过犯都关闭天主的教会之门，而我们则天天读到：\BibleQuote{我更喜欢罪人的悔改，而不是他的死亡}\BibleRef{则 18:23}，以及\BibleQuote{他们跌倒了，难道不能再起来吗？——上主说}\BibleRef{耶 8:4}，又一处说：\BibleQuote{归来的背逆之子啊，我要医治你们的背逆。}\BibleRef{耶 3:22} 他们的严厉并不能阻止他们自己犯大罪，远非如此；但我们与他们之间有这样的区别：他们因自以为是而羞于承认自己的过犯，而我们对我们的过犯行补赎，因此更容易获得赦免。\par

4. 我略过他们那据说由受凯旋式殉道的婴儿构成之罪之圣事。我宁愿说，我不相信这些；流血的指控很可能是假的。但我必须驳斥那些人公开的亵渎，他们说天主先在旧约中决意藉\Person{梅瑟}和先知拯救世界，但发现自己无法实现这一目的，便从童贞女取了身体，以基督的圣子身份宣讲，并为我们的救恩而受死。又说，当祂通过这两个步骤仍不能拯救世界时，祂最后藉圣神降临到\Person{蒙丹}和那两个癫狂的女人\Person{普黎斯卡}和\Person{玛克西米拉}身上；因此，这个残缺不全、绝了嗣的\Person{蒙丹}拥有了\Person{保禄}从未自称的那样圆满的知识；因为\Person{保禄}满足于说：\BibleQuote{我们如今所知道的，只是局部的；我们作先知的，也只是局部的}，又说：\BibleQuote{我们现在是藉着镜子观看，模糊不清。}\BibleRef{格前 13:9}\par

这些话无需反驳。揭露蒙丹派的不信就是战胜它。在这样一封短信中，我也无需逐一驳斥他们所提出的种种荒谬。你精通圣经；而且我认为，你写信并非因为你被他们的诡辩所困扰，而只是想了解我对他们的看法。\par

\SourceRecord{Translated by W.H. Fremantle, G. Lewis and W.G. Martley.；Nicene and Post-Nicene Fathers, Second Series；Vol. 6.；Edited by Philip Schaff and Henry Wace.；Buffalo, NY: Christian Literature Publishing Co.,；1893.；<http://www.newadvent.org/fathers/3001041.htm>.}

% structure: p0001 source=h1 target=section reason=page-title

\section{\WorkTitle{第四十二封书信}}

\Emph{\Person{玛尔大辣}}\par

\EditorialNote{应玛尔大辣之请，热罗尼莫向她解释基督所说的亵渎圣神之罪，并指出诺瓦齐安的解释站不住脚。写于罗马，公元385年。}\par

1. 你提出的问题简短，答案也显明。福音中有这样一段话：\Quote{凡说话干犯人子的，可得赦免；但谁若说话干犯圣神，不论今世或来世，都不得赦免。} \BibleRef{玛 12:32} 现在，如果诺瓦齐安断言只有背教的基督徒才能干犯圣神，那么显然，亵渎基督的犹太人并未犯此罪。然而他们是邪恶的佃户，曾杀害先知，又正谋害主；\BibleRef{玛 21:33} 且沉沦如此之深，以致天主子告诉他们，祂来正是为了拯救他们。\BibleRef{玛 18:11} 因此，必须向诺瓦齐安证明，那永不得赦免的罪，并非那些受酷刑而肠断的人在痛苦中否认其主时的亵渎，而是那些人的吹毛求疵的喧嚷：他们看见天主的作为是德性的果实，却将德性归之于魔鬼，并宣称所行的奇迹不属于神圣的卓越，而属于魔鬼。这正是我们救主论证的全部要旨，当祂教导说，撒殚不能驱逐撒殚，祂的国不能自相分裂。\BibleRef{玛 12:25} 如果魔鬼的目的是损害天主的受造物，祂如何会愿意医治病人，并将自己从祂所附的身体中驱除呢？让诺瓦齐安证明，在被迫于审判官面前献祭的人中，有谁曾宣称福音所记载的事迹不是由天主子所行，而是由魔鬼之王贝尔则步所行；\BibleRef{玛 12:24} 然后他才能坚持他的论点，即这是那永不得赦免的亵渎圣神的罪。\par

2. 但再提出一个更深入的问题：让诺瓦齐安告诉我们，他如何区分说话干犯人子与亵渎圣神。因为我认为，按照他的原则，那些在迫害中否认基督的人，只是说话干犯人子，并没有亵渎圣神。因为当一个人被问是否基督徒，而他声明自己不是时，显然在否认基督即人子时，他并未侮辱圣神。但如果他否认基督也涉及否认圣神，这位异端者也许能告诉我们，人子如何能被否认而不干犯圣神？如果他以为这里的\Quote{圣神}一词意指圣父，那么否认者在否认中并未提及圣父。当宗徒伯多禄被一个使女的问话吓住而否认主时，他是干犯人子还是干犯圣神？如果诺瓦齐安荒诞地曲解伯多禄的话：\Quote{我不认识那个人，} \BibleRef{玛 26:74} 认为否认的不是基督的默西亚身份，而是祂的人性，那么他就使救主成为说谎者，因为救主曾预言，祂自己，即祂的天主性，必须被否认。现在，当伯多禄否认天主子时，他痛哭流涕，并以三次宣认抹去了三次否认。\BibleRef{若 21:15} 因此，他的罪并非那永不得赦免的亵渎圣神的罪。显然，此罪涉及亵渎，即因某人的作为而称他为贝尔则步，而他的德行证明他是天主。如果诺瓦齐安能举出一个背教者称基督为贝尔则步的例子，我立刻放弃我的立场，承认这样的跌倒者无法获得赦免。在折磨下屈服并否认自己是基督徒是一回事，说基督是魔鬼是另一回事。这一点你若仔细阅读那段经文，你自己就会看到。\par

3. 我本应更充分地讨论此事，但几位朋友来访寒舍，我无法拒绝把自己交给他们。但若不能立即答复你，又似乎傲慢，于是我将一个广泛的主题压缩成寥寥数语，寄给你的不是一封信，而是一封释疑短札。\par

\SourceRecord{Translated by W.H. Fremantle, G. Lewis and W.G. Martley.；Nicene and Post-Nicene Fathers, Second Series；Vol. 6.；Edited by Philip Schaff and Henry Wace.；Buffalo, NY: Christian Literature Publishing Co.,；1893.；<http://www.newadvent.org/fathers/3001042.htm>.}

% structure: p0001 source=h1 target=section reason=page-title

\section{第四十三书}

\Emph{致\Person{玛尔塞拉}}\par

\EditorialNote{热罗尼莫将自己与奥力振的日常生活相对照，悲哀地承认自己的缺失。然后他向玛尔塞拉指出乡村生活相对于城市生活的优点，并暗示自己有意尝试一下。写于公元385年的罗马。}\par

1. \Person{盎博罗削}——这位向\Person{奥力振}（真正的金刚石与铜铁般的人）提供金钱、材料与抄写员以出版其无数著作之人——在从\Place{雅典}寄给朋友的一封信中写道，他们从未在没有诵读经文的情况下共同进餐，也从未就寝，除非有一位兄弟的声音将某段圣经读给他们听。事实上，日夜的安排是这样的：祈祷只为阅读让路，阅读又为祈祷让路。\par

2. 我们这些畜牲般的人，何曾做过类似的事？我们若诵读超过一小时，便打哈欠；我们揉搓额头，徒劳地试图抑制倦怠。然后，在这番伟大功绩之后，我们为了缓解，再次投身于世务之中。\par

我不谈那些使我们的心智迟钝的宴饮，也不愿估量我们花费在互相拜访上的时间。接着我们陷入交谈；我们浪费言语，在背后攻击他人，细数他们的生活方式，我们挑剔他们，也反过来被他们挑剔。这就是我们晚餐时及餐后所关注的\Quote{佳肴}。然后，当宾客离去，我们结算账目，这些账目定然引起我们的愤怒或焦虑。前者使我们如咆哮的狮子，后者则徒然试图为未来多年做准备。我们不记得福音的话：\BibleQuote{你这糊涂人，今夜就要索回你的灵魂；你所预备的一切，将归谁呢？}\BibleRef{路 12:20}我们所买的衣服，不仅为实用，也为炫耀。一有机会省钱，我们便加快脚步，言辞敏捷，竖起耳朵。若听到家中的损失——此类事常发生——我们的面容便沮丧阴郁。一文钱的收益令我们欢喜；半文钱的损失使我们忧伤。一个人有如此多副面孔，以致先知祈祷说：\Quote{主啊，在你的城中消散他们的形象。}因为我们本是按天主的肖像和模样受造的，\BibleRef{创 1:26}却是我们的邪恶使我们戴上面具。正如在舞台上，同一个演员时而扮演强壮的赫拉克勒斯，时而柔化为温柔的维纳斯，时而又因扮演库柏勒而颤抖；同样，我们——若不属世界，就会被世界所恨，\BibleRef{若 15:19}——每犯一罪，便有一副相应的面具。\par

3. 因此，既然我们一生中大部分时间在汹涌的大海上航行，我们的船只时而因狂风巨浪而摇晃，时而在险恶的礁石上撞碎，那么，就让我们尽快驶向乡村宁静的港湾吧。在那里，像牛奶、家常面包和我们亲手浇水栽培的蔬菜这类乡间美味，将为我们供应粗糙却无害的食物。如此生活，睡眠不会将我们召离祈祷，饱足不会使我们放弃阅读。夏季，树荫将为我们提供幽静；秋季，清新的空气和脚下散落的树叶会邀我们驻足休息；春季，田野鲜花盛开，我们的圣咏因鸟儿的啁啾而更加动听。当冬季带着霜雪来临，我不必购买燃料，无论睡眠还是守夜，我都比城里更暖和。至少，就我所知，我将以更少的开支抵御寒冷。让罗马保留其喧嚣与忙乱，让竞技场的残酷表演继续，让群众在赛马场喧闹，让观众在剧院纵情，并且——我绝不能完全忽略我们的基督徒朋友——让贵妇之家举行其日常会议。我们紧紧依靠主，并将我们的希望寄托于上主天主，是有益的，这样，当我们以现今的贫穷换取天国时，我们便能高呼：\BibleQuote{在天上除你以外，我还有谁？在地上除你以外，我也没有所爱慕的。}确实，若我们在天上能寻得如此福乐，我们便应悲叹曾在世上追求那贫乏而短暂的享乐。再会。\par

\SourceRecord{Translated by W.H. Fremantle, G. Lewis and W.G. Martley.；Nicene and Post-Nicene Fathers, Second Series；Vol. 6.；Edited by Philip Schaff and Henry Wace.；Buffalo, NY: Christian Literature Publishing Co.,；1893.；<http://www.newadvent.org/fathers/3001043.htm>.}

% structure: p0001 source=h1 target=section reason=page-title

\section{第四十四封书信}

\Emph{致玛尔塞拉}\par

\EditorialNote{玛尔塞拉曾送一些小礼物（可能是给保拉和欧斯托钦的），热罗尼莫现在以她们的名义写信感谢她。他注意到这些礼物不仅适合那两位女士，也适合他自己。写于公元385年的罗马。}\par

当我们身体分离时，我们惯常以心神彼此交谈。\BibleRef{哥 2:5}我们各尽其能。你送我们礼物，我们回你感谢信。既然我们是已蒙面纱的贞女，我们有责任表明你精美的礼物下隐藏的意义。那么，苦衣乃是祈祷和斋戒的标记，椅子提醒我们贞女永不外出，蜡烛则以灯光催促我们迎接新郎的来临。\BibleRef{玛 25:1}杯也告诫我们要克制肉身，常备殉道。圣咏作者说：\Quote{何其明亮，是主的杯爵，使饮者沉醉！}此外，当你向已婚妇女赠送驱蚊蝇的小拂尘时，这是巧妙的暗示，要她们立即抑制淫欲，因为\Quote{死蝇}，我们得知，\Quote{使香膏变坏}。那么，在这样的礼物中，贞女能找到模范，已婚妇女能找到榜样。对我而言，你的礼物也传达了一个教训，尽管是相反的一种。因为椅子适合闲人，苦衣适合补赎者，杯是为口渴者准备的。我将乐于点亮你的蜡烛，只要它能驱散黑夜的恐怖和良心有愧的恐惧。\par

\SourceRecord{Translated by W.H. Fremantle, G. Lewis and W.G. Martley.；Nicene and Post-Nicene Fathers, Second Series；Vol. 6.；Edited by Philip Schaff and Henry Wace.；Buffalo, NY: Christian Literature Publishing Co.,；1893.；<http://www.newadvent.org/fathers/3001044.htm>.}

% structure: p0001 source=h1 target=section reason=page-title

\section{《书信第四十五》}

致阿塞拉\par

\EditorialNote{杰罗姆在离开罗马前往东方后，致信阿塞拉，驳斥他所遭受的诽谤，特别是关于他与保拉和欧斯托基姆的亲密关系。写于公元385年8月，在奥斯蒂亚的船上。}\par

1. 倘若我以为自己能够报答你的善意，那我就太愚蠢了。只有天主才能代替我酬报那奉献给祂的灵魂。我实在是如此不配得到你的眷顾，以致我从未敢衡量其价值，甚至不敢奢求它会为了基督的缘故而赐予我。有些人视我为恶人，满身罪孽；这样的说法比起我实际的罪过，尚属宽恕。然而，你对待恶人，却善待他们，视他们为善人。判断别人的仆人，是危险的；\BibleRef{罗 14:4} 而出言毁谤义人，是一种不易被宽恕的罪。那日子必然到来，那时你我要为他人哀哭；因为有不少人将陷于火中。\par

2. 有人说我是个臭名昭著的叛徒，一个滑头无赖，一个用撒殚的伎俩说谎欺骗他人的人。我倒想知道，哪一种做法更安全？是捏造或相信这些针对无辜者的指控，还是即便对罪人也拒绝相信？有些人吻过我的手，却用毒蛇的舌头攻击我；同情在他们唇上，而恶意的喜悦却在他们的心里。主看见了他们，并嗤笑他们，将我这卑微之身和他们一同留待将来的审判。有人攻击我的步态或笑声；另一人在我的外貌上挑剔；还有人怀疑我态度的真诚。这就是我生活了近三年所处的同伴。\par

我常发现自己被贞女们环绕，我尽力为其中一些人阐释神圣的经卷。我们的学习带来了频繁的交往，这很快发展为亲密，进而产生了相互信任。如果她们曾在我身上看到任何有失基督徒体统的行为，就让她们说出来吧。我拿过谁的钱吗？我不是鄙视所有礼物，无论大小吗？我手中有过钱币的叮当声吗？\BibleRef{撒上 12:3} 我的言语暧昧不清，或是我的眼神轻浮吗？不；我的性别是我唯一的罪行，而即使这一点，也只有当有人谈论保拉要去耶路撒冷时，才受到攻击。那么，很好。他们相信了我的控告者撒谎时的话；为何他们不相信他收回前言时的话呢？他现在和那时是同一个人，然而他先前声称我有罪，如今却承认我清白。诚然，一个人在受刑时的言语比在欢乐的时刻更可信，除非人们倾向于相信那些旨在取悦他们耳朵的谎言，或者更糟的是，那些之前尚未被编造、他们却怂恿别人编造的故事。\par

3. 在我认识圣善的保拉的家庭之前，整个罗马都回荡着对我的赞美。几乎所有人都一致认为我配得上主教之职。可敬的达马苏斯只谈论我的话语。人们称我为圣洁、谦逊、有口才。\par

我曾踏入轻浮女子的门槛吗？我曾被丝绸衣服、闪亮的宝石、涂抹的脸庞或黄金的炫耀所迷惑吗？在罗马所有的贵妇中，只有一个能征服我，那就是保拉。她哀恸并禁食，浑身污秽，双眼因哭泣而模糊。她整夜祈求主的怜悯，常常在晨光中发现她仍在祈祷。圣咏是她唯一的歌谣，福音是她全部的话语，贞洁是她唯一的放纵，禁食是她生活的主调。唯一吸引我的女人，是我甚至从未在餐桌上见过的。然而，当我开始因她显著的贞洁而尊敬、敬重、崇敬她时，我先前的一切德行立刻就离我而去了。\par

4. 啊！嫉妒，你竟从撕裂自己开始！啊！撒殚那阴险的恶意，你总是迫害圣洁的事物！在罗马所有的贵妇中，唯一引起丑闻的是保拉和梅拉尼娅，她们藐视财富、抛弃儿女，高举主的十字架作为宗教的旗帜。倘若她们常去浴场，或选择使用香水，或利用她们的财富和寡妇身份享受生活、不受约束，她们就会被尊为高贵而圣洁的贵妇。而如今，当然，她们穿上苦衣和灰烬就是为了显得美丽，她们忍受禁食和污秽只是为了进入火焚的地狱！仿佛她们不能与那受大众喝彩的人群一同灭亡！如果攻击她们生活方式的是外邦人或犹太人，她们至少会得到安慰，因为她们只是令基督自己都未能取悦的人不悦。但可耻的是，攻击她们的是基督徒，他们忽视自己家庭的照管，无视自己眼中的大梁，却去寻找邻人眼中的木屑。\BibleRef{玛 7:3} 他们拆毁每一种虔敬的宣认，以为只要能否定他人的圣德、诽谤每一个人、证明灭亡的人很多、罪人是广大的一群，就能为自己的厄运找到补救。\par

5. 你每日沐浴；另一个人却视这种过度讲究为污秽。你饱食野禽并以吃鲟鱼为傲；我相反，以豆子填饱肚腹。你从成群欢笑的女郎中获得快乐；我更喜欢哭泣的保拉和梅拉尼娅。你贪图别人的东西；她们鄙视自己的所有。你喜欢蜂蜜调味的酒；她们喝冷水，更觉甜美。你认为那些你不能立刻拥有、吃光、吞下的东西都是损失；她们相信圣经，期待将来的美物。而如果她们错了，她们所依赖的身体复活是愚蠢的妄想，那与你有何相干？我们这一边，也不赞成你那样的生活。你可以尽情用你的好东西把自己养肥；我则宁愿选择苍白和消瘦。你这样的人认为我不幸；我们觉得你更为不幸。这样，我们互换了想法，彼此看来都是疯狂的。\par

6. 我匆忙写下此信，亲爱的\Person{亚塞拉}女士，当我上船时，满怀忧伤与眼泪，但我感谢我的天主，因我堪当承受世人的仇恨。\BibleRef{若 15:18}请为我祈祷，使我在巴比伦之后，能再见耶路撒冷；愿约匝达克之子\Person{耶叔亚}统治我，\BibleRef{盖 1:1}而非\Person{拿步高}；愿那位名字意为\Quote{助佑}的\Person{厄斯德拉}前来，将我带回本国。我曾是愚昧的，竟愿在异乡唱上主的歌，并离开西乃山，寻求埃及的帮助。我忘了福音警告我们\BibleRef{路 10:30}：那从耶路撒冷下去的人，立刻落在强盗手中，被剥去衣服，被打伤，半死半活被丢下。但是，虽然司铎和肋未人可能不顾我，仍有那好心的\Person{撒玛黎雅人}，当人对他说：\Quote{你是撒玛黎雅人，并且有魔鬼，}\BibleRef{若 8:48}他否认有魔鬼，但没有否认他是撒玛黎雅人，\BibleRef{若 8:49}这词在希伯来语中相当于我们的\Quote{守护者}。人们称我为惹祸者，而我接受这称号作为对我信德的承认。因为我不过是个仆人，而犹太人仍称我主人为魔术师。同样，宗徒也被称为欺骗者。我所受的试探，无非是人所常受的。\BibleRef{格前 10:13}我作为十字架的士兵，所经历的苦难何其少！人们将一项我并未犯的罪归咎于我；但我知道，我必须通过恶名与美名，进入天国。\BibleRef{格后 6:8}\par

7. 请代我问候\Person{保拉}和\Person{欧斯托钦}，她们无论世人如何想，在基督内永远是属于我的。也问候你的母亲\Person{阿尔比纳}、你的姊妹\Person{玛尔策拉}，以及\Person{玛尔策利纳}和圣\Person{斐利琪达}；并对她们所有人说：\Quote{我们众人都要站在基督的审判台前，\BibleRef{罗 14:10}那时各人生活的原则必要显明出来。}\par

现在，贞洁与童贞的光辉典范，我恳求你在祈祷中记念我，并借你的转求平息海中的风浪。\par

\SourceRecord{Translated by W.H. Fremantle, G. Lewis and W.G. Martley.；Nicene and Post-Nicene Fathers, Second Series；Vol. 6.；Edited by Philip Schaff and Henry Wace.；Buffalo, NY: Christian Literature Publishing Co.,；1893.；<http://www.newadvent.org/fathers/3001045.htm>.}

% structure: p0001 source=h1 target=section reason=page-title

\section{第四十六書信}

\Signature{\Person{葆拉}和\Person{欧斯托基乌姆}致\Person{玛尔塞拉}}\par

\EditorialNote{热罗尼莫以葆拉和欧斯托基乌姆的名义写信给玛尔塞拉，描述圣地的魅力，敦促她离开罗马，到白冷与旧日同伴相聚。信的大部分篇幅用于反驳一种异议，即基督受难后圣地已受诅咒。此信写于公元386年，寄自白冷，白冷从此成为热罗尼莫终生居住之地。}\par

1. 爱情无法衡量，急躁没有界限，渴望不容迟缓。因此，我们不顾自身的软弱，更多依赖意志而非能力，竟想——虽为学生——教导我们的师长。我们就像谚语中的母猪，妄想教训发明女神。是你首先点燃了我们的火绒；是你首先以言传身教，激励我们采纳如今的生活。正如母鸡聚集小鸡，你曾将我们收于翼下。如今你难道要让我们没有母亲在身边，随意乱飞吗？你难道要让我们惧怕猛禽的俯冲和每一只飞鸟的影子吗？与你分离，我们只能尽其所能：发出哀怨的悲叹，用呜咽多于眼泪恳求你，将我们所爱的玛尔塞拉归还我们。她温和，她柔美，她比最甜的蜜更甜。因此，她不可对我们严厉冷酷，因为正是她的魅力激励我们效法她的生活。\par

2. 假设我们所求是最好的，那么急于得到它并不值得羞耻。如果所有圣经都同意我们的观点，那么我们敦促你采取你常敦促我们的做法，并不显得过于大胆。\par

天主对亚巴郎所说的最初话语是什么？\BibleQuote{你离开你的家乡、你的家族，往我将指示你的地方去。}\BibleRef{创 12:1} 这位圣祖——第一个接受基督应许的人——被吩咐离开加色丁，离开那座混乱之城及其\OriginalTerm{勒霍博特}{Rehoboth}（即宽阔之处）\BibleRef{创 10:11}；也要离开史纳尔平原，那里曾建起高耸入云的高傲之塔。他必须渡过这世界的波涛，涉过其河流——那些圣徒们坐下流泪、思念熙雍的河流，以及革巴尔河，厄则克耳曾被他头顶的头发提到耶路撒冷。\BibleRef{则 8:3} 亚巴郎经历这一切，为要住在一片从上而得雨水浇灌的应许之地，不像埃及那样从下而得，\BibleRef{申 11:10} 它不生产给软弱患病者的青菜，\BibleRef{罗 14:2} 一片仰望天上早雨和晚雨的土地。那地有山有谷，\BibleRef{申 11:11} 高高在海面之上。它完全缺乏世俗的吸引力，但精神的吸引力却因此更大。主之母玛利亚离开平原，往山区去，当她接到天使的消息后，意识到自己怀有天主之子。古代培肋舍特人被征服、他们魔鬼般的狂妄被击溃、他们的勇士仆倒在地时，\BibleRef{撒上 17:49} 正是从这城出发，有一队欢腾的灵魂，一个和谐的合唱团，歌颂我们的达味战胜了千万敌军。同样，天使握剑，在摧毁整座不虔敬之城时，在耶步斯人敖尔难的禾场上标记了主的圣殿。这里早早显明，基督的教会将不是出自以色列，而是出自外邦人。回头翻看创世纪，\BibleRef{创 14:18} 你会发现这城曾是默基瑟德的治所，这位撒冷王作为基督的预像，向亚巴郎献上饼和酒，甚至在那时就祝圣了基督徒以救主的体和血祝圣的奥迹。\par

3. 也许你会默默责备我们背离了圣经的顺序，让我们混乱的叙述随意漫游，想到什么提什么。如果这样，我们再说一遍开头的话：爱情没有逻辑，急躁没有规则。在雅歌中有一个诫命，说得很艰难：\BibleQuote{请把你的爱情规范在我身上。}所以我们恳求，若我们犯错，那不是出于无知，而是出于情感。\par

那么，为了再提出一些更不合时宜的事，我们必须回溯更遥远的时代。传统认为，在这城——甚至就在这地方——亚当曾生活并死去。我们主被钉十字架的地方被称为\OriginalTerm{髑髅地}{Calvary}，因为元祖的头骨埋在那里。于是，第二亚当，即基督的血，从十字架上滴下，洗净了被埋葬的元祖——第一亚当的罪，如此就应验了宗徒的话：\BibleQuote{你这睡着的人，醒来吧，从死者中起来，基督必要光照你。}\BibleRef{弗 5:14}\par

要一一列举从这个地方派出的所有先知和圣者，将很冗长。凡对我们奇怪而神秘的事，对这城这地却是熟悉而自然的。它有三个名称，正是借此证实天主圣三的道理。因为它首先称为耶布斯，然后称为撒冷，再后称为耶路撒冷：第一个意思是\Quote{被践踏}，第二个是\Quote{和平}，第三个是\Quote{和平的异象}。因为我们是逐步达到目标；只有在被践踏之后，才能被提升而见到和平的异象。因着这和平，和平之人撒罗满在那里诞生，\Quote{他的居所在于和平}。万王之王、万主之主，他的名与城名共同显明他是基督的预像。我们还需要提及达味及其后裔，他们都在这里为王吗？正如犹太地高于所有省份，这城也高于整个犹太地。更简洁地说，省份的荣耀来自其首都；所有肢体的声望无不出于头部。\par

4. 你长久以来渴望发言；我们写下的书信已察觉这点，而纸张也已明白你将提出的问题。你会回答我们说：昔日如此，当\BibleQuote{主爱慕熙雍的城门，超过雅各伯所有的帐幕}，且她的根基在圣山上时。然而，这些诗句也可作更深的理解。但自那时起，情况已改变。复活的主以雷霆之声宣告：\BibleQuote{你们的房屋将留给你们荒凉}。祂含泪预言其倾覆：\BibleQuote{耶路撒冷，耶路撒冷！你杀死先知，又用石头砸死那些被派遣到你这里来的人；我多少次愿意聚集你的子女，如同母鸡聚集小鸡在翅膀底下，你却不愿意！看，你们的房屋将留给你们荒凉。}\BibleRef{玛 23:37} 圣所的幔子裂开了；\BibleRef{玛 27:51} 一支军队包围了耶路撒冷，它已被主的血玷污。因此，它的护守天使已离弃它，基督的恩宠也已收回。犹太史家\Person{约瑟夫斯}本人证实，在主被钉十字架时，从圣殿中传出天上权势的声音说：\Quote{我们离开这里吧。} 这些及其他考虑表明，哪里恩宠增多了，那里罪恶也格外增多。\BibleRef{罗 5:20} 再者，当宗徒们领受命令：\BibleQuote{你们要去使万民成为门徒，} \BibleRef{玛 28:19} 当他们自己说：\BibleQuote{天主的圣言本来先应讲给你们听，但因你们拒绝接受，……看，我们转向外邦人。} \BibleRef{宗 13:46} 那时，犹太地的一切属灵重要性及其与天主的古老亲密关系，都藉宗徒们转移给了万民。\par

5. 这个难题被有力地陈述出来，甚至可能使精通圣经的人也感到困惑；但尽管如此，它仍可轻易解决。主曾为耶路撒冷的沦陷而哭泣，\BibleRef{路 19:41} 祂若不爱它，就不会这样做。祂因爱\Person{拉匝禄}而为他哭泣。\BibleRef{若 11:35} 实情是：犯罪的是人民，而非地方。一座城市的陷落与其居民的遭戮紧密相连。耶路撒冷若被毁灭，那是为使它的百姓受惩罚；圣殿若被推翻，那是为使其中象征性的祭献被废除。至于其地点，时光的流逝反而赋予它新的庄严。古代的犹太人尊敬至圣所，是因其中所含之物——革鲁宾、赎罪盖、约柜、玛纳、\Person{亚郎}的棍杖以及金香坛。\BibleRef{希 9:3} 主的坟墓难道不值得更深的尊敬吗？每当我们进入其中，就看见身着殓布的救主；若我们停留片刻，又看见天使坐在祂脚边，以及卷好的汗巾放在祂头边。早在\Person{若瑟}凿出这坟墓之前，它的荣耀就已在\Person{依撒意亚}的预言中预示：\BibleQuote{他的安息之所将是光荣的}，\BibleRef{依 11:10} 意指主埋葬的地方应在普世受到尊崇。\par

6. 那么，你会说，我们如何在\Person{若望}所写的默示录中读到：\BibleQuote{那从深渊中上来的巨兽……要杀死他们（显然是指两位先知），他们的尸体要躺卧在大城的街上，这城按灵意称为索多玛和埃及，也就是他们的主被钉在十字架上的地方？} 如果主被钉十字架的大城是耶路撒冷，而祂被钉十字架的地方按灵意被称为索多玛和埃及，那么既然主是在耶路撒冷被钉十字架，耶路撒冷就必是索多玛和埃及。首先，圣经不能自相矛盾。一本书不能推翻全书的主旨。一节经文不能废止一卷书的意义。在默示录中早十行写道：\BibleQuote{起来，量量天主的殿和祭台，并量量在殿中朝拜的人。但圣殿外院要留下，不必量，因为它已交给了外邦人；他们将要践踏圣城四十二个月。} \BibleRef{默 11:2} 默示录是\Person{若望}在主受难后很久才写成的，但其中仍称耶路撒冷为圣城。但若是如此，他怎能按灵意称它为索多玛和埃及呢？不能回答说，被称为圣的耶路撒冷是将来的天上之城，而被称为索多玛的是行将倾覆的尘世之城。因为关于巨兽的描写所指的正是未来的耶路撒冷：\BibleQuote{那从深渊中上来的巨兽……要与两位先知交战，并要战胜他们，杀死他们；他们的尸体要躺卧在大城的街上。} \BibleRef{默 11:7} 在书的末尾，对此又有更远的描述：\BibleQuote{城是四方的，长宽都一样；天使用金芦苇量了城，共有一万二千\InnerQuote{斯塔狄}，长、宽、高都一样；又量了城墙，按照人的尺寸，即天使的尺寸，共有一百四十四肘。城墙是用碧玉造的，城是纯金的，} \BibleRef{默 21:16} 等等。然而，既是四方形的，何来长与宽？长与宽等于高又是何种量法？墙用碧玉、全城用纯金、地基和街道用宝石、十二个门各用珍珠发亮，这又如何可能？\par

7. 显然，这段描述不能按字面理解（事实上，假设一座城的长度、宽度和高度都等于一万二千斯塔狄亚是荒谬的），因此其中的细节必须神秘地理解。加音最初建造并以他儿子的名字命名的这座大城\BibleRef{创 4:17}必须被视为代表这个世界，即魔鬼——那位控告弟兄者、注定灭亡的杀害兄弟者——用罪恶（以犯罪为黏合剂）建造并充满不义的世界。因此它被灵意称为索多玛和埃及。经上记载：\BibleQuote{索多玛必回复先前的状态}，\BibleRef{则 16:55}也就是说，世界必须恢复原来的样子。因为我们不能相信索多玛、哈摩辣、阿德玛和责颇殷\BibleRef{申 29:23}将被重建：它们必须永远化为灰烬。我们从未读到埃及被当作耶路撒冷：它总是代表这个世界。要从圣经中收集无数证据证明这一点会令人厌烦：我仅举一处，这一处清楚地将这个世界称为埃及。宗徒犹达，雅各伯的兄弟，在他的公函中写道：\BibleQuote{我愿意提醒你们，虽然你们曾一次知道这事：耶稣既然从埃及地救出了百姓，后来却把那些不信的人除灭了。}\BibleRef{犹 1:5}并且，免得你以为指的是农的儿子若苏厄，经文继续写道：\BibleQuote{至于那些没有守住所居尊位、离弃自己住处的天使，主用永远的锁链将他们拘留在黑暗中，直到大日的审判。}\BibleRef{犹 1:6}此外，为使你确信在任何提到埃及、索多玛和哈摩辣的地方，所指的不是这些地方，而是现今的世界，他随即按这意思提到它们。他写道：\BibleQuote{同样，索多玛、哈摩辣和周围的城市，也照样放纵情欲，随从逆性的肉身，就受了永火的刑罚，作为鉴戒。}\BibleRef{犹 1:7}但是，当主的受难与复活之后，玛窦告诉我们：\BibleQuote{磐石崩裂，坟墓开了；许多长眠的圣者的身体也复活了，在耶稣复活后走出坟墓，进入圣城，显现给许多人。}\BibleRef{玛 27:51}我们不可像许多人荒谬地那样径直将此解释为天上的耶路撒冷：在那里圣者们身体的显现不能作为主复活的记号给人类。因此，既然众福音作者和全部圣经都称耶路撒冷为圣城，而圣咏作者也吩咐我们\Quote{在他脚凳前}敬拜主，就不要让任何人称它为索多玛和埃及，因为主禁止人指着它起誓，因为它\BibleQuote{是大君王的城。}\BibleRef{玛 5:35}\par

8. 你说，这片土地因汲取了主的血而受诅咒。那么，人们凭什么认为伯多禄和保禄——基督军队的领袖——为基督流血的那些地方是有福的呢？如果人的和仆人的宣信是光荣的，那么他们的主和天主的宣信岂不是更加光荣？我们在各处尊敬殉道者的坟墓，将他们的圣灰敷在眼睛上，甚至可能的话用唇舌触碰它们。然而有些人认为我们应当忽略主自己被埋葬的坟墓。如果我们拒绝相信人的见证，至少让我们相信魔鬼和他的天使。\BibleRef{玛 25:41}因为当他们在圣墓前被赶出被附的身体时，他们呻吟战栗，如同站在基督的审判座前，并为他们曾经钉死那位现在他们在他面前畏缩的主而懊悔太迟。如果——正如一种邪恶的理论所坚持的——自从主受难以来，这个圣地已经变得可憎，为什么保禄如此急忙赶到耶路撒冷在那里过五旬节？\BibleRef{宗 20:16}然而对那些劝阻他的人，他说：\BibleQuote{你们为什么哭泣，使我心碎？我为主耶稣的名不但准备被捆绑，也准备死在耶路撒冷。}\BibleRef{宗 21:13}我还需要提及那些其他圣洁而卓越的人吗？他们在基督宣讲之后，向耶路撒冷的弟兄们奉献了他们的感恩礼物和供物。\par

9. 时间不允许我回顾自主升天以来的时期，也不能一一列举那些因感到自己的虔诚和知识不够完美、德行缺乏最后润色，除非在福音首次从刑架上闪亮的地方朝拜基督，而来到耶路撒冷的主教、殉道者和神学家。如果一位著名的演说家责备一个人因为去了利利贝学希腊文而不是雅典，在西西里学拉丁文而不是罗马（显然是因为每个省份都有自己的特性），那么，我们能认为一个基督徒的教育是完整的，如果他没有参观基督徒的雅典吗？\par

10. 这样说，我们并非否认天主的国就在我们内\BibleRef{路 17:21}，也不是说其他地方没有圣人；我们只是最强烈地断言，那些在普世中居首位的人，在这里比肩而立。我们自己是末后，而非首先；然而我们来到此地，是为了见到万民中的首位。在教会的一切饰物中，我们的隐修士和童贞女团体是最精美的；它如同美丽的花朵或无价的宝石。高卢的每位显要人物都急忙赶来。不列颠人，\Quote{与世隔绝}，一旦在信仰上有所进步，就离开落日，去寻找一个他仅从圣经和传闻中得知的地方。我们还需要提及亚美尼亚人、波斯人、印度和阿拉伯的民众吗？或者我们的邻邦埃及，那里隐修士如此众多；本都、卡帕多细亚；以及科厄勒-叙利亚、美索不达米亚和丰饶的东方？为应验救主的话：\BibleQuote{无论身体在哪里，鹰也必聚在那里}\BibleRef{路 17:37}，他们全都聚集于此，在这同一座城市中展现出最繁多的美德。他们语言不同，却在信仰上合一，几乎每个民族都有自己的唱诗班。然而在这庞大的集会中，没有傲慢，没有对节制的轻蔑；所有人都追求谦逊，这是基督徒最大的德行。谁若是最末后，在此被视为首先。\BibleRef{玛 19:30}他们的衣着既不引人注意，也不招人羡慕。人无论以何种面貌出现，既不受责备，也不受奉承。长时间的禁食在这里对谁也无益。挨饿得不到尊敬，适度进食也不受谴责。\BibleQuote{各人} \BibleQuote{或站住或跌倒，}\BibleRef{罗 14:4}谁也不论断别人，免得被主论断。\BibleRef{玛 7:1}在别处常见的诽谤，在这里完全不存在。感官享乐和纵欲远离我们。在这城里，祈祷的场所如此之多，一天的时间都不够走遍它们。\par

11. 但是，既然人人都赞美自己所能及之物，让我们现在转到那曾庇护基督和玛利亚的客栈小屋。\BibleRef{路 2:7} 用什么样的言辞和语言，才能将救主的洞穴呈现在你们面前呢？那婴孩啼哭的马槽，最好用沉默来尊崇，因为言语不足以赞美它。宽阔的门廊在哪里？镀金的天花板在哪里？由注定灭亡之人悲惨劳作营建的宅邸在哪里？由无名富豪为人的卑微躯体行走而建的昂贵厅堂在哪里？拦截天空的屋顶在哪里——好像还有什么比穹苍更美好？看，在这大地的可怜缝隙中，诸天的创造者诞生了；在这里他被裹在襁褓里；在这里他被牧羊人看见；在这里他被星指出来；在这里他被贤士朝拜。我想，这个地方比那塔佩亚岩石更神圣，那岩石因屡遭雷击而显明为天主所不悦。\par

12. 请阅读\BibleBook{若望}{默示录}，思考其中所歌唱的：那身穿紫红袍的妇人，她额上写着亵渎的话，七座山，众水，以及巴比伦的结局。\Quote{我的百姓，从那城出来，}主说，\Quote{免得你们分担她的罪恶，也免得你们遭受她的灾祸。}\BibleRef{默 18:4} 也请回头阅读\BibleBook{}{耶肋米亚}，留意他所写的类似的话：\Quote{你们要从巴比伦中逃奔，各救自己的性命。}\BibleRef{耶 51:6} 因为\Quote{巴比伦大城已经倾倒了，成了魔鬼的住所，各种污秽之灵的巢穴。}\BibleRef{默 18:2} 诚然，罗马有一个圣教会，宗徒和殉道者的战利品，对基督的真实宣认。信仰曾由一位宗徒在那里宣讲，异教已被践踏，基督徒的名号日益高举。然而，这座城市的排场、权势和规模，被人看见和看见别人，拜访和接受拜访，交替的奉承和诋毁，谈论和倾听，以及即使最不愿意时也必须面对如此众多的人群——所有这些都同样与修道生活的原则相悖，并对修道生活的宁静是致命的。因为当人们出现在我们的路上时，我们要么看见他们来而被迫说话，要么没看见他们而招致傲慢的指责。有时，在回访中，我们不得不穿行于傲慢的门廊和镀金的大门，面对挑剔仆役的喧闹。但是，正如我们上面所说，在基督的小屋里，一切都是简单而质朴的：除了咏唱圣咏外，一片寂静。无论转向何处，耕地的农夫唱着阿肋路亚，辛劳的割草人用圣咏鼓舞自己，修剪葡萄园的葡萄园工唱着一首达味的诗歌。这些是乡间的歌曲；这些，用流行的话说，是它的情歌；这些是牧羊人吹的口哨；这些是农夫用来助其劳作。\par

13. 但我们这是在做什么？我们忘记了自己所应尽的职责，一心只想自己所愿的。难道那时刻永远不会到来吗？那时，气喘吁吁的使者将带来消息，说我们亲爱的\Person{玛尔策拉}已经抵达巴勒斯坦的海岸，每一队隐修士和每一群贞女将齐声唱起欢迎之歌？我们的兴奋之情已使我们急忙前去迎接你：不等车辆，我们立刻徒步赶路。我们将握住你的手，凝视你的面容；久候之后，我们终于拥抱你，将难以分离。难道那日子永远不会到来吗？那时我们将一同进入救主的洞穴，与祂的姐妹和母亲一同在主墓前哭泣？\BibleRef{若 19:25} 然后我们将以嘴唇亲吻十字架的木头，并在橄榄山上与上升的主一同祈祷、立志。\BibleRef{宗 1:9} 我们将看见拉匝禄裹着殓布走出来，\BibleRef{若 11:43} 我们将看见约旦河的水因主的洗礼而得洁净。\BibleRef{玛 3:13} 然后我们将前往牧羊人的草场，\BibleRef{路 2:8} 我们将一同在达味的陵墓祈祷。\BibleRef{列上 2:10} 我们将看见亚毛斯先知在他的岩石上吹牧羊人的号角。我们将匆匆前往，即使不是到帐篷，也到亚巴郎、依撒格、雅各伯以及他们三位卓越妻子的纪念碑。我们将看见那水泉，太监在那里受斐理伯的浸没。\BibleRef{宗 8:36} 我们将朝圣撒玛黎雅，并排敬礼洗者若翰、厄里叟和敖巴狄雅的遗骸。\BibleRef{列下 13:21} 我们将进入那些山洞，在迫害和饥荒时期，先知团体曾在那里得食。\BibleRef{列上 18:3} 只要你来，我们将前往纳匝肋，正如其名所示，那是加里肋亚的花朵。不远之处，将可见到加纳，那里水变成了酒。\BibleRef{若 2:1} 我们将前往大博尔山，\BibleRef{玛 17:1} 并看到那里的帐棚，救主在那里与圣父和圣神同在，而非如伯多禄曾愿的那样与梅瑟和厄里亚同在。然后我们来到革乃撒勒湖，在那里我们将看见五千人被五个饼吃饱，以及四千人被七个饼吃饱的地方。纳因城将映入我们眼帘，在那城门口，寡妇的儿子复活了。赫尔孟山也将可见，还有恩多尔溪流，息色辣在那里被击败。我们的目光也将落在葛法翁，那是我们主如此多奇迹的地点——是的，还有整个加里肋亚。当我们由基督陪伴，途经史罗和贝特耳，以及那些如同旌旗般设立以纪念主胜利的教会，回到我们的洞穴时，我们将衷心歌唱，尽情哭泣，不住祈祷。被救主的箭所伤，我们将彼此说：\Quote{我的灵魂所爱的，我找到了祂；我必抱住祂，决不放开。}\par

\SourceRecord{Translated by W.H. Fremantle, G. Lewis and W.G. Martley.；Nicene and Post-Nicene Fathers, Second Series；Vol. 6.；Edited by Philip Schaff and Henry Wace.；Buffalo, NY: Christian Literature Publishing Co.,；1893.；<http://www.newadvent.org/fathers/3001046.htm>.}

% structure: p0001 source=h1 target=section reason=page-title

\section{第四十七封信}

\Emph{\Person{德西德里乌斯}}\par

\EditorialNote{热罗尼莫邀请他在罗马的两位老朋友，\Person{德西德里乌斯}和他的姐妹（或妻子）\Person{塞雷尼拉}，到\Place{白冷}来与他相聚。这位\Person{德西德里乌斯}有可能（但可能性不大）即是\Person{阿基坦的德西德里乌斯}，此人后来促使热罗尼莫撰写反对\Person{维吉兰提乌斯}的文章。}\par

这封信（写于公元393年）与前一封相隔七年，在此期间所写的所有信件均已完全失传。\par

我深感惊讶，我卓越的朋友，您的好意促使您对我所言如此，我喜乐自己得到了一位既雄辩又真诚之人的见证；但当我从您转向我自己时，我感到懊恼，因为由于我的不配，您的赞美和称颂之言反使我沉重而非提升。您当然知道，我的原则是高举谦卑的标杆，并准备通过暂时行走在最低处来攀登高峰。因为我是什么，我有什么意义，竟值得学问之声为我作证，或让一位文风如此迷人、几乎使我打消写信念头的人，将雄辩的棕榈枝放在我的脚前？但我必须尝试，好使那不求自己益处、但常求他人利益的\BibleRef{格前 13:5}爱德，至少能回报一份恭维，既然它无法传递教训。\par

我向您和您圣洁而可敬的姐妹\Person{塞雷尼拉}表示祝贺，她名副其实，已践踏了世间汹涌的波涛，到达了基督平静的港湾；这一幸福——如果我们可信赖您名字的预兆——也在等待着您。因为我们读到，圣\Person{达尼尔}被称为\Quote{渴望之人}，是天主的朋友，因为他渴望知道天主的奥秘。因此，我乐于做可敬的\Person{保拉}托付给我的事。我以主的爱德敦促并恳求你们二位，务请光临我们这里，并希望访问圣地能促使你们以此大礼丰富我们。即使你们不关心我们的团体，作为信徒，你们仍有责任在主的脚曾站立的地方朝拜，并亲眼目睹祂诞生、十字架和受难尚存的痕迹。\par

我的几篇小文章已飞出巢窠，轻率地自行寻求出版的荣耀。我没有寄给您，以免寄去您已经有的作品。但如果您想借阅副本，可以向我们的圣姐妹\Person{玛尔塞拉}（她住在\Place{阿文蒂诺山}上）或那位圣人\Person{多姆尼奥}（他是我们时代的\Person{罗特}）借用\BibleRef{伯后 2:7}。与此同时，我期待您的到来，一旦您来，我将把我所有的一切都给您；或者，如果有任何阻碍使您不能与我们相聚，我将乐意为您送去您所期望的论著。效法\Person{特兰奎卢斯}和\Person{希腊人阿颇罗尼乌斯}的榜样，我写了一本关于从宗徒时代到我们这个时代著名人物的书；在列举了许多人之后，我在最后一页写下了自己，如同一个未及时出生的人，是所有基督徒中最小的\BibleRef{格前 15:8}。在此，我发现有必要简要叙述一下我的著作，直至\Person{德奥多西}皇帝在位的第十四年。如果您从上述人士那里得到这本论著，发现其中提到您尚未拥有的一些作品，如果您愿意，我将逐步为您抄录。\par

\SourceRecord{Translated by W.H. Fremantle, G. Lewis and W.G. Martley.；Nicene and Post-Nicene Fathers, Second Series；Vol. 6.；Edited by Philip Schaff and Henry Wace.；Buffalo, NY: Christian Literature Publishing Co.,；1893.；<http://www.newadvent.org/fathers/3001047.htm>.}

% structure: p0001 source=h1 target=section reason=page-title

\section{第四十八封信}

\Emph{致帕玛丘斯}\par

\EditorialNote{这是对热罗尼莫之前不久所写两卷驳若维尼安著作的辩护；他将这些著作的副本寄到了罗马。帕玛丘斯和他的其他朋友没有发表这些著作，因为他们认为热罗尼莫过分抬高童贞而贬低了婚姻。他现在写信来证明自己的立场，为此他从那本惹人非议的著作中大量摘录。这封信的日期是公元393或394年。}\par

1. 你一直保持沉默，这就是我迄今没有写信的原因。因为我担心，如果我先写信给你而尚未收到你的来信，你会认为我不是一个尽责的，而是一个令人厌烦的通信者。但现在，我被你那封极其令人愉快的信所激发，这封信要求我诉诸基本原则来辩护我的观点，我以张开双臂的方式——如俗语所说——迎接我昔日的老同学、同伴和朋友；我期待你成为我贫乏著作的捍卫者；也就是说，我首先要争取你的判断，使你能做出对我有利的裁决，并让我这位辩护人了解我受到攻击的所有要点。因为你所喜爱的西塞罗，以及在他之前——在他那篇短论中——安东尼乌斯都写道，胜利的首要条件是要仔细熟悉自己要辩护的案情。\par

2. 某些人指责我，因为在我反对若维尼安的书中，我（如他们所说）过分赞美了童贞，贬低了婚姻；他们断言，宣扬贞洁直到不留下妻子与贞女的比较，就等于谴责婚姻。如果我没有记错争论的要点，我与若维尼安之间的分歧在于：他将婚姻置于与童贞同等的地位，而我则认为它较低；他宣称两者之间几乎没有或根本没有差别，而我则断言有很大差别。最后——在天主安排下通过你的中介达到的结果——他已被定罪，因为他竟敢将婚姻与永久的贞洁等同起来。或者，如果贞女和妻子应被视为相同，罗马怎么拒绝听从这个不敬神的教义？贞女来自男人，但男人并非来自贞女。没有折中之道。要么我的观点被采纳，要么若维尼安的观点被采纳。如果我因将婚姻置于童贞之下而受责备，那么他因将两者置于同等地位而应受称赞。反之，如果因假设他们相等而被他定罪，他的定罪必须被视为对我论著的认可。如果世俗的人因感到他们处于比贞女低下的地位而恼怒，我奇怪神职人员和隐修士——两者都过独身生活——为何不称赞他们一贯实践的东西？他们离开妻子以模仿贞女的贞洁，然而却坚持已婚妇女与她们一样好。他们要么应该重新回到他们所放弃的妻子身边，要么，如果他们坚持与妻子分开生活，他们将不得不——用生活而非用言语——承认，在将童贞置于婚姻之上时，他们选择了更好的道路。那么，我只是圣经的新手，第一次阅读圣书吗？童贞与婚姻之间的界限如此细微，以至于我未能注意到？我当然不会不知道这句话：\Quote{不要过于正义！}因此，当我试图保护自己的一面时，另一面却受了伤；更明白地说，当我与若维尼安进行短兵相接的交战时，摩尼教在背后刺伤了我。我问道，在我的著作的最前面，我是否没有写下如下前言：\Quote{我们不是马克西昂或摩尼教的门徒，来贬低婚姻。我们也没有被塔提安——恩克拉提特派的领袖——的错误所欺骗，以致认为所有同居都是不洁的。因为他不仅谴责和拒绝婚姻，也谴责和拒绝天主所造供我们享用的食物。\BibleRef{弟前 4:3}我们知道在大户人家，不仅有金银的器皿，也有木器和瓦器。\BibleRef{弟后 2:20}我们也知道，在匠师保禄所立的基督根基上，有人用金银、宝石建造，有人却用草、木、禾秸建造。\BibleRef{格前 3:10}我们并非不知道\BibleQuote{婚姻应受尊重……床笫是无玷污的。}\BibleRef{希 13:4}我们读过天主的第一个命令：\BibleQuote{你们要生育繁殖，充满大地。}\BibleRef{创 1:28}但我们虽然允许婚姻，却更偏好从婚姻产生的童贞。金子比银子更贵重，但银子因此就不是银子了吗？从一棵树，偏好它的果实胜过它的根与叶，这是对树的侮辱吗？从谷物，将穗子置于秆和叶之上，这是对谷物的伤害吗？如同果实来自树，谷粒来自稻草，同样童贞来自婚姻。一百倍的、六十倍的、三十倍的出产\BibleRef{玛 13:8}可能都来自同一土壤、同一播种，但它们会在数量上大不相同。三十倍出产代表婚姻，因为表达该数字时手指的并拢，暗示着温柔亲切的亲吻或拥抱，恰当地代表了夫妻关系。六十倍出产指寡妇，她们处于困苦和磨难的地位。因此，它们由另一手指放在下方以表达六十的数字所预表；因为一旦尝过快乐又能抵制其诱惑是极其困难的，所以这样做的回报也相应地大。此外，一百倍——我请求读者给予最大的注意——需要从左手换到右手；但虽然手不同，手指却与左手表示已婚妇女和寡妇的相同；只是在这个例子中，它们形成的圆圈代表童贞的冠冕。}\par

3. 请问：一个这样说的人，是否在谴责婚姻？若我曾称童贞为金，则我也曾称婚姻为银。我曾说明那结一百倍、六十倍、三十倍的果实——全都出自同一土壤、同一播种，虽然数量大不相同。我的读者中会有人如此不公，不以我的话，而以自己的意见来评判我吗？无论如何，我对待婚姻比大多数拉丁和希腊作者温和得多；他们把那百倍果实归于殉道者，六十倍归于贞女，三十倍归于寡妇，从而表明在他们看来，已婚者被排除在好地和大父的种子之外。但惟恐有人以为我开头谨慎，而著作其余部分却轻率，我岂不在标出其划分之后，一到实际问题时立即引入了以下内容：\Quote{我请求你们所有男女，既包括贞女和节欲者，也包括已婚或再婚者，以你们的祈祷来协助我的努力。}若维尼安是所有人的公敌，但我岂能把我需要其祈祷、恳求其帮助我工作的人斥为摩尼教异端？\par

4. 由于信件的简短不容我们在一点上耽搁太久，现在让我们转向其余部分。在解释宗徒的证言时：\Quote{妻子没有权柄主张自己的身体，乃是丈夫；照样，丈夫也没有权柄主张自己的身体，乃是妻子。}\BibleRef{格前 7:4}我们附加了以下内容：\Quote{整个问题关乎那些生活在婚姻中的人，他们是否可以休妻——这也是主在福音中所禁止的事。\BibleRef{玛 19:9}因此，宗徒也说：\InnerQuote{男人不碰妻子或女人，是好的。}\BibleRef{格前 7:1}仿佛接触有危险，接触者无法逃脱。因此，当埃及妇人想要触摸若瑟时，他丢下外衣，从她手中逃走了。\BibleRef{创 39:12}但是，那人一旦娶了妻子，除非双方同意，不能与她断绝关系或休她——只要她没有犯罪——他必须向妻子尽本分，因为他自愿承担了强制履行的义务。}一个宣称主有命令说妻子不可被休，且天主所联合的，人不可未经同意分开\BibleRef{玛 19:6}的人——这样的人能说是谴责婚姻吗？同样，在接下来的经文中，宗徒说：\Quote{可是每人都有天主各自己的恩赐，一个这样，一个那样。}\BibleRef{格前 7:7}在解释这话时，我们做了如下评论：\Quote{他自己愿意什么，是清楚的。但由于在教会中有不同的恩赐，\BibleRef{格前 12:4}我也允许婚姻，免得我好像谴责本性。也请思考：童贞的恩赐是一回事，婚姻的恩赐是另一回事。因为若已婚妇女和贞女有同一赏报，他绝不会在给出节制的劝告后接着说：\InnerQuote{可是每人都有天主各自己的恩赐，一个这样，一个那样。}既然每一类有其各自的恩赐，那么各类之间必定有分别。我承认婚姻和童贞一样，都是天主的恩赐，但恩赐与恩赐之间有很大差别。最后，宗徒本人谈到一个曾生活在乱伦中后来悔改的人说：\InnerQuote{你们倒不如饶恕他，安慰他，}\BibleRef{格后 2:7}并且\InnerQuote{你们饶恕谁什么，我也饶恕。}\BibleRef{格后 2:10}并且，免得我们以为人的恩赐是小事，他又加上：\InnerQuote{因为我若饶恕了什么，我所饶恕的，是为你们的缘故在基督面前饶恕的。}\BibleRef{格后 2:10}基督的恩赐各不相同。因此，作为祂预像的若瑟有一件彩色外衣。\BibleRef{创 37:23}同样，在第四十四篇圣咏中，我们读到有关教会的话：\InnerQuote{王后身穿金绣彩衣，侍立在祢右边。}宗徒伯多禄也（谈到丈夫和妻子）\InnerQuote{同为多种恩宠的承继人。}在希腊文中，表达更加显著，所用词汇是\OriginalTerm{彩色}{\GreekText{ποικίλη}}，即\InnerQuote{多种颜色的。}}\par

5. 那么我问，人们顽固地闭眼拒绝看那明如白昼之事，是什么意思？我曾说过，教会中有不同的恩赐，童贞是一种恩赐，婚姻是另一种。不久之后，我用了以下话语：\Quote{我也允许婚姻是天主的恩赐，但恩赐与恩赐之间有很大差别。}能说我谴责那最明确地宣布为天主恩赐的事物吗？再者，若若瑟被视为主预像，他的彩色外衣是贞女、寡妇、独身者和已婚者的预像。任何有分于基督长袍的人，能被视为外人吗？我们岂不曾说到王后本人——即救主的教会——穿着金色绣衣，有各种颜色？此外，当我在讨论与以下经文相关的婚姻时，\BibleRef{格前 7:8}我仍坚持同一观点。我说：\Quote{这段经文确实与当前的争议无关；因为，跟随主的决定，宗徒教导说，妻子除非因奸淫，不可被休；如果她被休，她在丈夫有生之年不能再嫁，或者至少她应尽可能与丈夫和好。在另一节中，他表达了同样的意思：\InnerQuote{妻子…丈夫活着的时候，是受束缚的；丈夫若死了，她就脱离了丈夫的法律，\BibleRef{罗 7:2}她可以随意再嫁，只要在主内\BibleRef{格前 7:39}}，意即嫁给基督徒。因此，宗徒虽然允许在主内第二次或第三次婚姻，却禁止甚至第一次与外教人结婚。}\par

6. 我要求我的批评者打开耳朵，意识到一个事实：我允许了\Quote{在主内}的第二次和第三次婚姻。那么，如果我没有谴责第二次和第三次婚姻，我又怎能禁止第一次呢？此外，在我解释宗徒的话的那段经文里：\BibleQuote{你原是受割损而蒙召的吗？不要变成未受割损的。你原是未受割损而蒙召的吗？不要受割损。}\BibleRef{格前 7:18}（确实，有些最谨慎的圣经注释者认为这段经文是指法律上的割损和奴役），我不是以最清楚的措辞为婚姻纽带辩护吗？我的原话是这样的：\Quote{\InnerQuote{如果你是在未受割损时蒙召，就不要受割损。}宗徒说，你信的时候已经有了妻子。不要以为你因信基督就有理由离开她。因为\InnerQuote{天主召叫了我们，是为了和平。}\InnerQuote{受割损算不得什么，不受割损也算不得什么，只在于遵守天主的诫命。}\BibleRef{格前 7:19}没有行为，无论独身还是婚姻都没有丝毫用处，因为连基督徒特征的信德，如果没有行为，也被说成是死的，\BibleRef{雅 2:17}照此条件，\Person{维斯塔}或\Person{朱诺}的贞女（她们始终忠于一个丈夫）也可以要求被列入圣人之列。稍后他还说：\InnerQuote{你是作奴隶蒙召的吗？不要为此烦恼；但如果你能获得自由，就使用它更好；}\BibleRef{格前 7:21}这就是说，如果你有妻子，并且依附于她，要尽你对她应尽的义务，你的身体不由自己作主——或者更明白地说——如果你是一个妻子的奴隶，不要为这事难过，不要为失去童贞叹息。即使你能找到离开她的借口，享受洁德自由，也不要为了自己的利益而毁掉别人。再忍耐妻子一段时间，不要太急于克服她的不愿。等待她效法你的榜样。只要你耐心，你的妻子有一天会成为你的姐妹。}\par

7. 在另一段经文中，我们讨论了\Person{保禄}这样说的原因：\BibleQuote{关于童贞者，我没有主的命令，但我蒙主仁慈，忠信地提出我的意见。}\BibleRef{格前 7:25} 在这里，我们虽然赞扬了童贞，但也谨慎地给了婚姻应有的地位。\Quote{如果主命令了童贞，}我们说过，\Quote{祂就会似乎谴责婚姻，并取消人性那个种子基地，童贞本身正是从那里产生的。如果祂切断了根，怎能期待果实呢？除非祂首先打好地基，又怎能建造大厦，或者用覆盖整个范围的屋顶给它加冕呢？} 如果我们将婚姻说成是根，童贞是其果实，并且我们将婚姻作为基础，永久洁德的大厦或屋顶建立在其上，那么我的指责者中谁能如此吹毛求疵或盲目，以至于在眼前既有大厦和屋顶本身时，却忽略大厦及其屋顶所建立的基础呢？再次，在另一处，我们引用了宗徒的见证：\BibleQuote{你束缚于妻子吗？不要寻求解脱。你从妻子解脱了吗？不要寻求妻子。}\BibleRef{格前 7:21} 对此我们附加了以下评论：\Quote{我们每个人都被分配了自己的领域。让我拥有我的，你也保持你的。如果你束缚于妻子，不要离开她。如果我从妻子解脱，就不要让我寻求妻子。正如我不解开已建立的婚姻纽带，你也应避免将目前解脱于这类纽带的人捆绑在一起。} 还有另一段经文对我们就童贞和婚姻所持的观点提供了明确无误的见证：\Quote{宗徒没有给我们设下圈套，\BibleRef{格前 7:35} 也没有强迫我们成为我们所不愿的。他只敦促我们去做光荣和合宜的事，恳切地激励我们事奉主，常常渴望取悦祂，并寻求祂为我们预备的旨意。我们应该像警觉且武装的士兵，立即执行给他们的命令，不假思索地执行，没有训道者所说给予世人\InnerQuote{用以受锻炼}的那种心智劳苦。} 最后，在我们将贞女与已婚女性进行比较的结尾处，我们这样总结了讨论：\Quote{当事物一好而另一物更好时；当好的事物有与更好的事物不同的酬报；并且当酬报不止一种时，那么，显然，存在着恩赐的多样性。婚姻与童贞之间的区别，正如不做恶与行善之间的区别一样大——或者，更中肯地说，正如善与更善之间的区别一样大。}\par

8. 随后，我们继续这样发言：\Quote{宗徒在结束他对婚姻与守贞的讨论时，小心地保持中间路线以区分二者，既不偏右也不偏左，而是行走于王道，\BibleRef{户 20:17}，从而履行了\BibleQuote{不要极端正义}的训诫。\BibleRef{训 7:16} 此外，当他继续比较一夫一妻制与再婚时，他将再婚置于一夫一妻之后，正如他先前将婚姻置于守贞之下。} 我们岂非藉此言辞清楚表明了圣经中以\Quote{右}和\Quote{左}术语所预表的内容，以及我们如何理解\InnerQuote{不要极端正义}这句话的含义？我们若随从犹太人和外邦人的情欲，对性交感到焦渴，便是向左转；我们若随从摩尼教徒的错误，以贞洁为借口陷入不洁的网罗，便是向右转。但我们若向往守贞且不谴责婚姻，便是行走于王道。此外，任何人若读到我关于第二次婚姻的意见如下，怎能如此不公地批评我的拙著，指责我谴责第一次婚姻呢？\Quote{诚然，宗徒允许第二次婚姻，但仅限于那些执意如此、不能自制的妇女，\BibleRef{格前 7:9} 免得她们\InnerQuote{开始纵欲违背基督而嫁人，因此受到谴责，因为她们抛弃了起初的信德}，他作出这一让步是因为许多人\BibleQuote{转去随从了撒殚}\BibleRef{弟前 5:15}。然而她们若能守寡，就更有福。对此他立即加上自己的权威说：\BibleQuote{这是我的意见。} 此外，免得有人认为那只是人的权威分量不足，他接着说：\BibleQuote{我也以为自己有天主的神。}\BibleRef{格前 7:40} 因此，当他劝人节欲时，他诉诸的不是人的权威，而是天主的神；但当他允许人结婚时，他并未提及天主的神，而是允许审慎的考虑来平衡，为了各人的需要而放宽他律例的严格。} 既然已经提出证据表明宗徒允许第二次婚姻，我们随即附加了以下评论：\Quote{正如因奸淫的危险而允许贞女结婚，且本身并非值得向往之事因此可被谅解，同样，因同一危险而允许寡妇再婚。因为一个女人认识一个男人（尽管他是第二个或第三个丈夫），好过她认识多个男人。换言之，她委身于一人而非多人，更可取。} 诽谤尽可恣意。我们在此谈的不是第一次婚姻，而是第二次、第三次，或（如果你愿意）第四次。但唯恐有人将我的话（即女人委身于一人而非多人更好）应用于第一次婚姻，而我的全部论述涉及的是再婚与三婚，我便用以下话表明我对这些行为的看法：\Quote{\BibleQuote{一切都可行，但不全有益处。}\BibleRef{格前 6:12} 我不谴责再婚者，也不谴责三婚者，甚至（说得极端些）八婚者。我还要作更大的让步：我甚至愿意接受一个悔改的奸淫者。在任何可能公平的地方，都必须给予公平的考量。}\par

9. 我的诽谤者应因宣称我谴责第一次婚姻而羞惭，当他读到我刚才引用的话：\Quote{我不谴责再婚者或三婚者，甚至（说得极端些）八婚者。} 不谴责是一回事，称赞是另一回事。我可以允许一种做法为可行，而不称赞它为有功。但如果我显得严厉地说：\Quote{在任何可能公平的地方，都必须给予公平的考量，} 我想，没有人会认为我残忍或严苛，如果他读到为贞女和已婚者所预备的地方，与为三婚者、八婚者和忏悔者所预备的地方不同。基督自己，虽在肉体上是贞洁的，在精神上却是一夫一妻者，只有一位妻子，即教会，\BibleRef{弗 5:23} 我已在论述的后半部分指出。然而，我竟然被认为谴责婚姻！我虽用了以下这些话，却仍被说成谴责婚姻：\Quote{无可置疑的事实是，肋未司祭出自亚郎、厄肋阿匝尔和丕乃哈斯的后裔；既然他们都是已婚者，倘若我们被禁欲派的错误所迷惑，主张婚姻本身应受谴责，那么我们大可以他们为例来反驳。} 在此我责怪禁欲派的首领塔提安拒绝婚姻，然而我自己却被说成谴责婚姻！此外，当我将贞女与寡妇对比时，我自己的话表明我对婚姻的看法，并提出了我提议的三重等级：贞女、寡妇（不论是实际或事实上）和已婚妻子。\Quote{我不否认}——这是我的原话——\Quote{那些受洗后继续守寡者的幸福，也不低估与丈夫在贞洁中生活的妻子的功绩；但是，正如寡妇比顺从丈夫的妻子从天主那里获得更大的赏报，她们也必须满足于看到贞女被置于自己之上。}\par

10. 再者，在解释宗徒致迦拉达人所作的证言时，\Quote{没有人能因法律的行为而成义，}我曾发言如下：\Quote{婚姻也是法律的行为。因此，不生育子嗣者是有诅咒的。婚姻即使是在福音之下也是许可的；但一件事是让步于软弱，另一件事是应许德行的赏报。}看我的明确声明：婚姻在福音中是允许的，然而已婚者彼此尽夫妻之责时，不能接受贞洁的赏报。如果已婚者对此感到愤慨，就让他们不要向我发怒，而要向圣经发怒；不仅如此，更要向所有的主教、司铎、执事、以及全体司祭和肋未人发怒，因为他们知道，如果他们履行婚姻义务，就不能献祭。再者，当我从默示录引用证据时，我关于贞女、寡妇和已婚妇女的观点难道不清楚吗？\Quote{这些人是唱新歌的人\BibleRef{默 14:3}，除了贞女外没有人能唱这歌。这些人就是\InnerQuote{首先献给天主和羔羊的初熟之果}\BibleRef{默 14:4}，他们是没有瑕疵的。如果贞女是天主的初熟之果，那么度节欲生活的寡妇和已婚妇女就必须排在初熟之果之后——也就是说，在第二位和第三位。}因此，我们把寡妇和已婚妇女放在第二位和第三位，为此，我们被一个异端者的疯狂指控为完全谴责婚姻。\par

11. 全书我以极为温和的语气对童贞生活、寡居生活和婚姻生活发表了许多言论。但为简洁起见，我在此只引一段，且这段是如此性质，我想没有人会反对它，除非有人想证明自己是恶意的或疯狂的。在描述我们的主造访加里肋亚加纳的婚宴\BibleRef{若 2:1}时，在另外一些评论之后，我加了以下的话：\Quote{他只去过一次婚宴，这教导我们女子应只结婚一次；如果我们未能给予婚姻其应有的位置——即在童贞和贞洁的寡居之后——这一事实可能对童贞不利。但是，既然只有异端者才谴责婚姻并践踏天主的诫命，我们欣然聆听我们的主所说的一切赞美婚姻的话。因为教会并不谴责婚姻，只是使之处于从属地位。它并不完全拒绝婚姻，只是加以规范，知道（正如我前面所说）\InnerQuote{在大户人家，不但有金器银器，也有木器瓦器；有的作贵重的用途，有的作卑贱的用途。所以，人若自洁……他必成为贵重的器皿……预备行各样的善事。}\BibleRef{弟后 2:20}}我在这里说，我欣然聆听宗徒所说的一切赞美婚姻的话。我是否欣然聆听对婚姻的赞美，却又谴责婚姻？我说，教会并不谴责婚姻，而是使之处于从属地位。无论你喜欢与否，婚姻都从属于童贞和寡居。即使婚姻继续履行其功能，教会也不谴责它，只是使之处于从属地位；它并不拒绝它，只是加以规范。如果你愿意，你有能力登上贞洁的第二步。为什么你站在第三级、最低级，却不愿赶快向上攀登呢？\par

12. 既然我如此频繁地提醒读者我的观点；既然我像谨慎的旅行者一样步步谨慎，反复声明，虽然我承认婚姻本身是可接受的，但我更喜爱节欲、寡居和童贞，那么明智而宽厚的读者本应根据我的总体方向来判断那些看似难懂的话，而不应指责我在同一本书中提出不一致的观点。因为谁如此迟钝或写作经验不足，竟会称赞又谴责同一件事物，毁坏自己所建立的，又建立自己所毁坏的，并在击败对手之后最后将剑转向自己？如果我的指责者是乡下出身，或不熟悉修辞学或逻辑学的技巧，我会原谅他们缺乏洞察力；如果我看到他们的无知而非恶意是过失所在，我也不会指责他们控告我。然而，那些受过通识教育的智者，其目的与其说是理解我，不如说是伤害我，对于这些人，我有这样一个简短的回答：他们应当纠正我的错误，而不仅仅是谴责我的错误。竞技场已开放，我喊道；你的敌人已列好阵势，他的位置明显，并且（如果我可以引用维吉尔的话）——\par

敌人正在召唤你：与他正面交锋。\par

这些人应当回应他们的对手。他们应当保持在辩论的范围内，而不应挥舞教师的教鞭。他们的书应旨在指出我的陈述在哪些方面偏离了真理，在哪些方面超过了它。因为，尽管我不听吹毛求疵者的话，但我会听从导师的建议。当你自己占据城墙上的有利位置时，却指导战士如何作战，这是一种不值得推荐的教学方式；当你自己沉浸在香水中时，指控流血的士兵懦弱是不光彩的。我说这话并不是让自己被指责为自夸，说当他人沉睡时，只有我进入了竞技场。我的意思只是，那些在这场战斗中看到我受伤的人，可能在战斗方法上过于谨慎。我不希望你们参与一场除了自卫外无事可做的遭遇战，你的右手保持不动，而左手操纵盾牌。你必须要么出击，要么倒下。除非我看到你的对手被剑杀死，否则我不认为你是胜利者。\par

13. 你们无疑是学识渊博的人；但我们也曾在学府中学习，并且像你们一样，从\Person{亚里士多德}——或者更确切地说，从他得自\Person{高尔吉亚}的——教诲中学到了不同的说话方式；我们也知道，其中有一件事：为展示而写作的人采用一种风格，为说服而写作的人采用另一种。在前一种情况下，辩论是散漫的；为驳斥对手，时而引用这个论证，时而引用那个。人们随心所欲地论证，说着一件事而意指另一件。用一句谚语来说：\Quote{一只手给你面包，另一只手里却拿着石头。}在后一种情况下，需要某种坦诚和面容的开放。因为提出问题和阐述已被证明的事是一回事。前者需要辩手，后者需要教师。我身处战火之中，生命时刻危险；而你，自称教导我的人，却是个书斋之人。\Quote{不要，}你说，\Quote{出其不意地攻击或从侧面刺伤。直击你的对手。你应当以诡计而非力量为耻。}好像战斗的最高境界不是虚张声势而攻其不备。我恳请你，读一读\Person{狄摩西尼}或\Person{西塞罗}，或者（如果你不关心那些以巧言而非真话为目的的演说家）读一读\Person{柏拉图}、\Person{泰奥弗拉斯}、\Person{色诺芬}、\Person{亚里士多德}以及其他从\Person{苏格拉底}这一泉源汲取各自智慧之流的人。他们有任何坦诚吗？他们毫无技巧吗？他们所说的每一句话不都充满意义吗？而且这意义不总是为了胜利吗？\Person{奥利振}、\Person{美多丢}、\Person{欧瑟伯}和\Person{阿波利拿里}写了大量著作反对\Person{凯尔苏斯}和\Person{波菲利}。请想，他们用以推翻魔鬼之灵所造之物的论证是何等精妙，策略何等隐秘。有时，的确，他们不得不说的不是他们所想的，而是必需的；为此他们甚至用外邦人自己的主张来对付对手。我不提拉丁作家，如\Person{特土良}、\Person{西彼廉}、\Person{米努提乌斯}、\Person{维克托里努斯}、\Person{拉克坦提乌斯}、\Person{依拉利}，免得我似乎不是在为自己辩护，而是在攻击别人。我只提宗徒\Person{保禄}，他的话语我每听到一次，就觉得不是话语，而是雷霆。读他的书信，尤其是写给罗马人、迦拉达人和厄弗所人的，在其中他都身处激战，你将看到他在从旧约中引证时是多么巧妙和谨慎，他是多么小心地遮掩他的目的。他的话语看似极其单纯：一个天真无邪、不谙世事的人——既无技巧设计两难困境，也无技巧逃避它——的表达。然而，无论你往哪边看，它们都是霹雳。他的辩论踌躇，却凡所着手无不获胜。他背对敌人只是为了战胜他；他假装逃跑，只是为了击杀。因此，若有人应受诽谤，那就是他；我们应该这样对他说：\Quote{你用来反对犹太人或其它异端者的那些引证，在其自身语境中的含义与你书信中所给的含义不同。我们看到，这些段落被你的笔所掳掠，被迫为你服务以赢得胜利，而在它们出自的著作中，它们毫无争论意义。}难道他不会用救主的话回答我们：\Quote{对于外人我有一种说话方式，对于内人有另一种；群众听到我的比喻，但解释只给我的门徒。} \BibleRef{玛 13:10} 主向法利塞人提出问题，但不加以阐明。教导门徒是一回事，战胜对手是另一回事。\Quote{我的奥秘属于我，}先知说，\Quote{我的奥秘属于我和那些属于我的人。}\par

14. 你对我感到义愤，因为我只是使\Person{约维尼安}静默，而没有教导他。我说你？不，是他们，那些听到他被绝罚而悲伤的人，他们通过赞扬他人所持有、而他们自己也持有的异端，来指控自己所谓的正统。我本应请他和平投降！我无权无视他的挣扎，违背他的意愿将他拖入真理的捆绑！如果胜利的欲望使我说出任何违反圣经规则的话，如果我也像辩论中强者常做的那样，以结果证明手段的正确，那么我或许会使用这样的语言。然而，事实上，我宣讲宗徒的教导，而不是凭我自己建立教义；我的职责只是一个注释者。因此，任何看起来难懂的言论都应归咎于我所诠释的著者，而不是我，诠释者；除非，他所说的与别人所代表的不一致，并且我通过不公正的解释曲解了他话语的明显含义。如果有人以此指控我不诚实，让他从圣经本身来证明他的指控。\par

我曾在我的书中说过：\Quote{如果\InnerQuote{男人不亲近女人是好的}，那么亲近女人对他就是坏的，因为只有坏才是好的反面。但若它虽是坏的，却被宽免为小罪，那么便允许它，以防止比坏更糟的事，}如此直到下一章开头。上述是我对宗徒的话的评注：\Quote{男人不亲近女人是好的。但为了避免淫乱，男人当各有自己的妻子，女人当各有自己的丈夫。}\BibleRef{格前 7:1}我的意思与宗徒的原意有何不同？只在于他肯定地说，而我犹豫不定。他确定信理，我冒昧探究。他公开说：\Quote{男人不亲近女人是好的。}我怯生生地询问，男人不亲近女人是否\Emph{是好的}。若如此犹豫，便不能说是肯定地说。他说：\Quote{不亲近是好的。}我则加上了\Quote{好}的可能反义。紧接着我这样说：\Quote{请注意宗徒的谨慎。他并没有说：\InnerQuote{男人不要妻子是好的}，而是说\InnerQuote{男人不亲近女人是好的}；仿佛亲近女人本身就包含危险——亲近者无法逃脱的危险。}因此，你看，我并非在解释关于夫妻的法律，而只是讨论性关系的普遍问题——相较于贞洁、童贞及天使的生活，\Quote{男人不亲近女人是好的}如何如何。\par

\Quote{\OriginalTerm{虚而又虚}{Vanity of vanities}}，训道者说，\Quote{万事皆虚}。\BibleRef{训 1:2}但如果一切受造物都是善的，因为出自善的造物主之手，怎么万事皆是虚呢？如果大地是虚，那么诸天也是虚吗？——以及天使、座天使、宰制天使、掌权天使以及其他各品德能天使？不；如果那些本身为善的事物，因出自善的造物主之手，也被称为虚，那是因为它们与更美的事物相比较。例如，与灯相比，灯笼便毫无用处；与星相比，灯便全然不发光；最亮的星在月亮前也黯然失色；将月亮置于太阳旁，它便不再明亮；将太阳与基督相比，它便是黑暗。\Quote{我是自有者，}天主说；\BibleRef{出 3:14}如果你将一切受造物与他相比，它们便不存在。\Quote{不要把您的权杖，}艾斯德尔说，\Quote{交给那虚无者}，\BibleRef{艾 14:11}——就是说，交给偶像和魔鬼。而她祈求不要将自己和自己的人交给他们，他们无疑就是偶像和魔鬼。在约伯传中我们也读到彼耳达得论及恶人时如何说：\Quote{他的信赖必从帐幕中被拔除，毁灭如君王般践踏他。那不存在者的同伴也要住在他的帐幕中。}这显然是指魔鬼，他必定存在，否则便不能说他有同伴。然而，因为他失落于天主，所以被称为不存在者。\par

我在类似的意义上宣称触摸女人是坏事——我不是说妻子——因为不触摸是好事。我又加上：\Quote{我称童贞为细粮，婚姻为大麦，奸淫为牛粪。}诚然，粮食和大麦都是天主的受造物。但在福音中奇迹般供应的两群人之中，较大的一群是靠大麦饼喂养，较小的一群是靠细粮面包。\Quote{祢，上主，}圣咏作者说，\Quote{祢要拯救人和牲畜。}我自己也用其他话说过同样的事，我曾称童贞为金子，婚姻为银子。此外，在讨论那十四万四千没有沾染妇女的盖印童贞者时，我曾努力表明，所有没有保持童贞的人，与天使及我们的主耶稣基督的完美贞洁相比，都被视为污秽。但如果有人认为我将童贞与婚姻之间的距离等同于细粮与大麦之间的距离是苛刻或应受谴责的，就请他阅读圣盎博罗削的著作\WorkTitle{论寡妇}，他会在其中找到关于童贞和婚姻的其他论述，如下：\Quote{宗徒并没有如此毫无保留地表达他对婚姻的偏好，以致熄灭人们对童贞的渴望；他以推荐节制作开始，只是随后才屈尊提及相反之事的补救措施。虽然对强壮者，他指出了他们崇高召叫的奖赏，\BibleRef{斐 3:14} 但他不让任何人在路上昏倒；\BibleRef{玛 15:32} 他称赞那些领路的人，却不轻视那些殿后的人。因为他自己已经学会，主耶稣给了某些人大麦饼，免得他们在路上昏倒，却将自己的身体给了另一些人，使他们努力获得他的国度；}紧接其后：\Quote{那么，婚姻的束缚不应作为罪行来避免，而应作为沉重的负担来拒绝。因为法律束缚妻子在辛苦和忧伤中生儿育女。她的欲望要归于丈夫，使他管辖她。\BibleRef{创 3:16} 因此，不是寡妇，而是新娘被交托给生育的辛苦和忧伤。不是童贞女，而是已婚妇女受到丈夫的支配。}在另一处，宗徒说：\Quote{你们是重价买来的，所以不要作人的奴仆。}\BibleRef{弗 6:6} 你看他多么清楚地定义了伴随婚姻状态的奴役。稍后又说：\Quote{那么，如果连好的婚姻都是奴役，坏的婚姻又当如何？在其中丈夫和妻子不能圣化彼此，而只能互相毁灭。}我关于童贞和婚姻所详述的，盎博罗削却简洁而一针见血地陈述，将许多意义压缩在几句话中。他描述童贞是推荐节欲的方式，婚姻是不节制的补救。当他从宽泛的原则下降到具体细节时，他意味深长地向童贞者展示崇高召叫的奖赏，同时安慰已婚者，免得他们在路上昏倒。他赞扬一类人，却不鄙视另一类人。他将婚姻比作摆在一群人面前的大麦饼，将童贞比作赐给门徒的基督身体。在我看来，大麦和细粮之间的差异比大麦和基督身体之间的差异小得多。最后，他称婚姻为沉重的负担，应尽可能避免，也是再明确不过的奴役标记。他还做了许多其他陈述，他在他的三卷著作\WorkTitle{论童贞}中详细阐述了这些。\par

15. 从所有这些考虑可以看出，我关于童贞和婚姻所说的并无新意，而是在所有方面都遵循了更早的作家的判断——也就是说，盎博罗削及其他讨论过教会教义的人。\Quote{我宁愿跟随他们的错误，也不愿抄袭今日作家的乏味迂腐。}让已婚者，如果他们愿意，因我说过\Quote{我问你，那禁止人祈祷、阻止人领受基督身体的是什么样的好东西？}而勃然大怒吧。当我履行丈夫的本分时，我不能满足节欲的要求。同一宗徒在另一处命令我们常常祈祷。\BibleRef{得前 5:17} \Quote{但如果我们应当常常祈祷，我们就绝不能屈从于婚姻的要求，因为每次我向妻子履行义务，就使自己无法祈祷。}我这样说时，显然依赖宗徒的话：\Quote{你们不要彼此亏负，除非两相情愿，暂时分房，为专务祈祷。}\BibleRef{格前 7:5} 保禄宗徒告诉我们，当我们与妻子交合时，就不能祈祷。那么，如果性交阻止了较不重要的事——即祈祷——那么它更会阻止更重要的事——即领受基督身体。伯多禄也劝勉我们节欲，以使我们的\Quote{祈祷不受阻碍。}\BibleRef{伯前 3:7} 我想知道，我在这所有事上犯了什么罪？我做了什么？我有什么错？如果溪流的水浑浊，错的不是河床，而是源头。我受到攻击，是因为我胆敢在宗徒的话上加上了我自己的这句话：\Quote{那阻止人领受基督身体的是什么样的好东西？}如果是这样，我将简短回答：哪个更重要，祈祷还是领受基督身体？当然是领受基督身体。那么，如果性交阻碍较不重要的事，它就更阻碍那更重要的事。\par

我在同一篇论著中曾说过，达味和与他在一起的人不能合法地吃供饼，除非他们回答说，三天来他们没有与女人同房——当然不是与妓女，因为法律禁止与妓女行淫，而是与她们自己的妻子，他们合法地结合了。再者，当百姓即将在西乃山上接受法律时，他们被命令三天远离自己的妻子。\BibleRef{出 19:15} 我知道在罗马，信友们习惯经常领受基督的身体，对此我既不谴责也不认可。\BibleQuote{各人应在自己心中充分确信}\BibleRef{罗 14:5} 但我呼吁那些在当日行房事后领受共融者的良心——他们先是如佩尔西乌斯所说，\Quote{在流动的溪流中洗去夜晚}——我问他们，为什么不敢走近殉道者或进入圣堂？基督在外是一心，在家却是另一心吗？在圣堂内不合法的，在家里也不合法。没有什么能瞒过天主。在祂面前，\Quote{黑夜亮如白昼}。各人应省察自己，然后才可接近基督的身体。\BibleRef{格前 11:28} 当然，将共融推迟一天或两天，并不会使基督徒更加圣洁，也不会让我今天不配得的，明天或后天就配得。但如果我因未能分享基督的身体而忧伤，这确实帮助我暂时避免妻子的拥抱，并让对基督的爱胜过对婚姻的爱。你会说，这是一条严苛的诫命，难以承受。世俗之人谁能承受呢？我回答说：\BibleQuote{能领受的，就领受吧}\BibleRef{玛 19:12}；不能承受的，就自己看着办。我的职责不是说每个人能做什么或愿意做什么，而是说圣经教导什么。\par

16. 再者，有人对我关于宗徒的以下评论提出了反对：但恐怕有人会从前面所引的话（即\BibleQuote{为能专心祈祷，再结合一起}的上下文中，以为宗徒希望这种圆满，而不仅仅是为避免更糟糕的堕落而让步，他立刻加上\BibleQuote{免得撒殚因你们的不节制而诱惑你们}。\BibleRef{格前 7:5} \BibleQuote{再结合一起} 这句话传达了何等高贵的宽容！他羞于用更直白的话说出，只将它置于撒殚的诱惑之上，并根源于不节制。我们还需要费力解释这句难懂的话吗，作者自己已经解释了意思？他说：\BibleQuote{我说这话，是出于许可，而不是出于命令} 我们还在犹豫，是否将婚姻说成是一种许可而不是一种命令？或者我们害怕这种许可会排除第二次或第三次婚姻或其他情况？我在这里说了什么宗徒没有说过的话？我想，那句\Quote{他羞于用更直白的话说出}。当我想到他说\Quote{再结合}，却没有提及为何事，他是以一种含蓄的方式指出他不愿公开说出的事情——即性交。或者反对的是后面的话——\Quote{只将它置于撒殚的诱惑之上，并根源于不节制}？这不正是宗徒自己的话，只是重新排列了次序——\BibleQuote{免得撒殚因你们的不节制而诱惑你们}？或者人们吹毛求疵是因为我说了：\Quote{我们还在犹豫，是否将婚姻说成是一种许可而不是一种命令？} 如果这话听起来刺耳，那应归咎于宗徒，他说：\BibleQuote{我说这话，是出于许可，而不是出于命令}，而不是归咎于我，因为我除了重新排列顺序外，既没有改变词句，也没有改变意思。\par

17. 书信简短，迫使我匆忙前行。因此，我转向剩余的问题。宗徒説：\BibleQuote{我説，对未婚者和寡妇，他们若像我一样保持下去，为他们便好。但若他们不能自制，就让他们结婚；因为结婚比燃烧更好。} \BibleRef{格前 7:8} 我如此解释这一段：\Quote{当他准许已婚者使用婚姻，并显明自己的意愿和让步之后，他转向未婚者或寡妇，将自身的榜样摆在他们面前。他称他们为有福的，若他们像他一样保持下去，\BibleRef{格前 7:8} 但他接着说：\BibleQuote{他们若不能自制，就让他们结婚。}他以此重复了先前的话：\BibleQuote{但只是为了避免奸淫}，以及\BibleQuote{免得撒但因你们的不能自制诱惑你们。}当他说：\BibleQuote{他们若不能自制，就让他们结婚}，他给出这话的理由是：\BibleQuote{结婚比燃烧更好。}结婚之所以好，只因为燃烧是坏的。但若除去情欲之火，他就不会说\BibleQuote{结婚更好}。因为一件事被称为更好，是与更坏的事相对，而不是与公认的好事对比。这就如同他说：\InnerQuote{有一只眼睛比没有眼睛好。}} 不久之后，我向宗徒致辞说：\Quote{如果婚姻本身是好的，就不要将其与燃烧相比，而只需说：\InnerQuote{结婚是好的。}我必定怀疑一件事物的善，它只是在更大的恶面前才成为较小的恶。就我而言，我希望它不是较轻的恶，而是彻头彻尾的善。} 宗徒希望未婚女子和寡妇戒绝房事，激励她们效法他的榜样，并称她们有福，若她们像他一样保持下去。但若她们不能自制，且受到诱惑欲以奸淫而非节欲来扑灭情欲之火，他告诉她们，结婚比燃烧更好。对此训诫，我作了如下评论：\Quote{结婚之所以好，只因为燃烧是坏的，} 不是提出我自己的见解，而只是解释宗徒的训诫：\BibleQuote{结婚比燃烧更好；} 即，嫁人比犯奸淫更好。如果你们教导燃烧或奸淫是好的，那么善仍会被更善所超越。但如果婚姻只是比它所优先的恶稍好一点，它就不能具有那无瑕疵的完美与福乐，以至与天使的生命相比。假设我说：\Quote{作贞女比作已婚妇人更好；} 在此，我将更好的事物优先于好的事物。但若我再进一步说：\Quote{结婚比犯奸淫更好；} 在此，我不是将更好的事优先于好事，而是将好事优先于坏事。这两种情况有极大的差别；因为贞女相对于婚姻，是好中之更好，而婚姻相对于奸淫，是坏中之好。请问，我在这一解释中犯了什么罪？我坚定的目的不是要曲解圣经以适应自己的意愿，而只是说出我所认为的含义。注释者不宜大谈自己的见解；他的职责是阐明他所声称要解释的作者的含义。因为，如果他反驳他试图解释的作者，他将证明自己是他的对手而非解释者。当我自由地表达自己的意见，而不是注释圣经时，那么任何乐意的人都可以指控我苛刻地谈论婚姻。但是，如果他找不到这样指控的根据，就应将我注释中显得严厉或苛刻的段落归咎于被注释的作者，而非我，我只是他的解释者。\par

18. 另一个对我的指控实在令人难以忍受！有人提出，我在解释宗徒关于夫妻的话时，\BibleQuote{这样的人肉身必有苦难，} 我説：\Quote{我们由于无知，原以为婚姻至少在肉身方面会有喜乐。但如果已婚者在肉身方面要有苦难——这似乎是他们唯一可能获得快感的事——那么还有什么动机能使女子结婚呢？因为，除了在心神和灵魂上受苦，她们甚至在肉身上也要受苦。} \BibleRef{得前 5:23} 如果我列举婚姻的苦难，如婴儿啼哭、子女夭折、流产的几率、家庭损失等等，难道我就谴责了婚姻吗？当圣德达玛穌仍在世时，我曾写了一本书驳斥厄尔维丢，题为\WorkTitle{论荣福玛利亚的永久童贞}，在其中，为适当颂扬贞女的福乐，我不得不说了许多婚姻的苦难。那位卓越的人——精通圣经，又是童贞教会的童贞博士——可曾在我的论述中发现任何可指责之处？此外，在我写给欧斯托基的书信中，我用了更严厉的言辞谈论婚姻，却没有人因此感到冒犯。不，每一位热爱贞洁的人都侧耳聆听我对节欲的赞颂。请读戴尔都良、读西彼廉、读盎博罗削，要么连同他们一起指控我，要么连同他们一起赦免我。我的批评者类似于普劳图斯笔下的人物。他们唯一的才智在于诽谤；他们通过轮番攻击各方，试图使自己显得博学。因此，他们不加区别地将指责加于我和我的对手，并声称我们二人都被击败，尽管两者中必有一人获胜。\par

此外，当我讨论\OriginalTerm{再婚}{digamy}和\OriginalTerm{三婚}{trigamy}时，曾说过：\Quote{一个女人认识一个男人——即使他是第二个丈夫或第三个——比认识几个男人更好；她向一个人卖淫比向许多人更容易容忍}，我难道没有立刻附上我这样说的理由吗？\Quote{福音中的撒玛黎雅妇人，当她宣告她现在的丈夫是她的第六个时，主责备她，因为那并不是她的丈夫。}就我而言，我现在再次自由地宣告：再婚在教会中并未被定罪——不，三婚也未被定罪——一个女人可以合法地嫁给第五个丈夫，或第六个，或更多，就像她可以嫁给第二个一样；但这样的婚姻虽未被定罪，却也不受赞美。它们是对不幸命运的缓解，绝不有助于守贞的荣耀。我在别处也说过类似的话。\Quote{当一个女人结婚不止一次——无论她结婚两次还是三次，区别不大——她就不再是一夫一妻者了。\InnerQuote{凡事都可行……但不都有益处。}\BibleRef{格前 6:12}我不谴责再婚者或三婚者，甚或，举一个不可能的例子，八婚者。让一个女人如果有必要就嫁第八个丈夫吧；只要她停止卖淫。}\par

19. 我现在要谈那段我被指责说过的话——至少根据真实的希伯来文本——即在创造的第二天后，并没有像第一天、第三天及之后几天那样插入\Quote{天主看了认为好}\BibleRef{创 1:10}这句话，并且我立即加了以下评论：\Quote{我们被暗示在数字二中有某种不善，它将我们与合一分开，并预表了婚姻的纽带。正如在诺厄方舟的记载中，所有成对进入的动物都是不洁的，而那些取单数的则是洁净的。}\BibleRef{创 7:2}在此陈述中，有人对我关于第二天所说的话提出了暂时的反对，要么因为上述经文确实出现，而我说它们不出现；要么因为假设它们出现，我对它们的理解与上下文显然要求的不同。至于所问的经文（即\Quote{天主看了认为好}）不出现，让他们不要相信我的证据，而是相信所有犹太和其他译者——即阿奎拉（\Person{Aquila}）、息玛古（\Person{Symmachus}）和特敖多提昂（\Person{Theodotion}）的证据。但如果这些经文在其他天的记述中出现，而在这一天的记述中不出现，要么他们给出一个比我所给的更合理的理由来解释它们的不出现，要么，如果没有这样的理由，就让他们无论喜欢与否，接受我提出的解释。此外，如果在诺厄方舟中，所有成对进入的动物都是不洁的，而取单数的则是洁净的，并且关于文本的准确性没有争议，那么请他们解释，如果可能，为何这样写。但如果他们无法解释，那么，无论他们愿意与否，他们必须接受我对这件事的解释。要么提供更好的饭食，邀请我做客，要么就满足于我提供给你的餐食，无论它多么粗陋。\par

我现在必须提及处理过这个奇数问题的教会作家。他们中，希腊方面有克肋孟（\Person{Clement}）、依玻里托（\Person{Hippolytus}）、奥力振（\Person{Origen}）、狄约尼削（\Person{Dionysius}）、欧瑟伯（\Person{Eusebius}）、狄迪莫（\Person{Didymus}）；拉丁方面有戴尔都良（\Person{Tertullian}）、西彼廉（\Person{Cyprian}）、维多利诺（\Person{Victorinus}）、拉克坦提（\Person{Lactantius}）、依拉利（\Person{Hilary}）。西彼廉写给福图纳图（\Person{Fortunatus}）关于数字七的话，从他寄给那人的信中看得很清楚。或者我应当提出毕达哥拉斯（\Person{Pythagoras}）、塔兰托的阿尔基塔（\Person{Archytas of Tarentum}）和普布利乌斯·西庇阿（\Person{Publius Scipio}）在（西塞罗的）第六卷《论共和国》（\WorkTitle{论共和国}）中的推理。如果我的指责者不愿听这些，我就让文法学校向他们高声朗诵维吉尔（\Person{Virgil}）的诗句：\par

奇数乃天主所喜。\par

20. 像我曾说的，童贞比婚姻更洁净，偶数必须让位于奇数，旧约的预像证实了福音的真理：这似乎是一项大罪，颠覆教会，为世人所不容。我认为，在我的论文中受到指责的其他各点较不重要，或者归结为这一点。因此，我未予回答，既为了不超出我可用的篇幅，也为了不显得不信任你的才智，因为我知道你甚至在我请求之前就愿意为我辩护。那么，在我最后一息，我声明：无论是现在还是从前，我从未谴责过婚姻。我仅仅是回答了一个对手，毫无担心自己一方的人会为我设陷阱。我将童贞高举到天上，不是因为我自己拥有它，而是因为，由于没有拥有它，我更加钦佩它。诚然，谦逊而坦率地承认：赞美别人身上你所缺乏的，这是一种美德。我身体的重量使我贴在地面，但我难道不钦佩飞翔的鸟，不赞美鸽子吗？——正如维吉尔所说的，它\par

在流动的路径上滑翔，以不动的迅捷翅膀？\par

没有人可以欺骗自己，没有人可以听从谄媚的声音而奔向毁灭。第一次童贞来自出生，第二次来自重生。这话不是我的，而是一句古老的格言：\Quote{没有人能服侍两个主人；}\BibleRef{玛 6:24}即肉体和精神。因为\Quote{肉体的私欲相反圣神的引导，圣神的引导相反肉体的私欲；这两者互相敌对，}\BibleRef{迦 5:17}致使你们不能行你们所愿意的。因此，当我的小作品中有什么看起来严厉时，不要看我所说的话，而要看它们所取自的圣经。\par

21. 基督亲自是童贞者，祂的母亲也是童贞者；是的，虽为祂的母亲，她仍是贞女。因为耶稣曾穿过关闭的门户进入，\BibleRef{若 20:19}，在祂的坟墓——一个从最坚硬磐石中凿出的新墓穴——中，前后都没有人躺卧。\BibleRef{若 19:41}玛利亚是\Quote{封闭的园子……密封的泉源}，\BibleRef{歌 4:12}从这泉源流出，按约厄尔所说，那灌溉绳索或荆棘溪谷的河流；绳索指我们从前被捆绑的罪恶之绳，\BibleRef{箴 5:22}荆棘指那些窒息家主所播种子的荆棘。\BibleRef{玛 13:7}她就是厄则克耳先知所说的东门，永远关闭、永远发光，既隐藏又揭示至圣所；通过她，\Quote{正义的太阳}，\BibleRef{拉 4:2}我们的\Quote{按照默基瑟德品位的大司祭}，\BibleRef{希 5:10}出入往来。让我的批评者向我解释：耶稣如何能穿过关闭的门户进入，同时祂又允许人触摸祂的手和肋旁，显示祂有骨有肉，证明祂是真实的身体而非幻影；然后我将解释圣玛利亚如何同时是母亲和贞女。她在婚配前已是母亲，生育儿子后仍保持童贞。因此，如我正要说的，童贞的基督和童贞的玛利亚为两性奉献了童贞的初果。宗徒们或是童贞者，或是已婚却过着独身生活。被选为主教、司铎和执事的人，或是童贞者，或是鳏夫；至少当他们领受圣职后，便誓许终身贞洁。我们为何自欺，并感到烦恼？如果我们不断追求性欲的满足，却发现贞洁的棕榈枝不赐予我们，我们便希望过着奢华的生活，享受妻子的怀抱，同时却渴望与基督在贞女和寡妇中为王。难道饥饿与饱足、污秽与华丽、麻衣与丝绸，将有同一个赏报吗？拉匝禄\BibleRef{路 16:19}在世时受了苦，那穿紫红袍、肥胖光滑的富人活着时享受了肉体的美好；但现在他们死了，处境却不同了。痛苦变成了满足，满足变成了痛苦。而我们要么跟随拉匝禄，要么跟随那富人，这取决于我们自己。\par

\SourceRecord{Translated by W.H. Fremantle, G. Lewis and W.G. Martley.；Nicene and Post-Nicene Fathers, Second Series；Vol. 6.；Edited by Philip Schaff and Henry Wace.；Buffalo, NY: Christian Literature Publishing Co.,；1893.；<http://www.newadvent.org/fathers/3001048.htm>.}

% structure: p0001 source=h1 target=section reason=page-title

\section{第四十九封书信}

\Emph{致\Person{帕玛邱}}\par

\EditorialNote{热罗尼莫附上前一封信，感谢帕玛邱为制止他反对约维尼安的论文所作的努力，但声言这些努力已属徒然，并劝告他若心中仍存有丝毫犹疑，当转向圣经以及奥力振等对圣经所作的注释。此信与前信同时写成。}\par

1. 基督的谦逊有时要求我们甚至对朋友也保持沉默，并在安静中滋养我们的谦卑，因为旧友情的重续会招致自私之嫌。因此，当你保持沉默时，我也保持了沉默，并且不愿与你理论，以免被人认为我急于讨好有权势之人而不是培养一位朋友。但现在，由于回复你的来信已成为一项义务，我将努力总是先于你，与其说是回答你的询问，不如说是独立于它们而写作。这样，如果我先前以沉默显示了谦逊，今后我将更以发言显示它。\par

2. 我十分认可你的善意和远见，你因此撤回了一些我的反对\Person{约维尼安}的论文的抄本。然而，你的努力是徒劳的，因为从城里来的几个人多次向我大声朗读他们在罗马遇到的一些段落。在这个省区，这些书也已经流传开了；正如你自己在\Person{贺拉斯}作品中读到的：\Quote{言出难收}。我不像当今大多数作家那样幸运——也就是说，能够随时修改我的小作品。一旦我写了什么，无论是我的仰慕者还是我的恶意者——出于不同的动机，但以同样的热情——都会把我的作品广泛传播给公众；他们的评价，无论是赞扬还是批评，往往走向极端。他们不是根据作品的优点，而是根据自己愤怒的情绪来评判。因此，我做了我能做的。我已将一篇为所述作品辩护的文字献给了你，坚信当你阅读后，你自会替我解答别人的疑问；或者，如果你也对此嗤之以鼻，那么你就必须以某种新的方式来解释宗徒的那段（\EditorialNote{格前第七章}）论述贞洁与婚姻的章节。\par

3. 我这样说并非为了激你亲自就此主题写作——虽然我知道你在研读圣经方面的热忱比我更大——而是为了让你迫使我的折磨者去写。他们有学识，在他们自己眼中不是平庸的学者；他们不仅有能力责备我，也有能力教导我。如果他们就此主题写作，我的观点与他们的比较起来就会更快被忽视。我请求你阅读并仔细思考宗徒的话，然后你会看到——为了避免误解——我对已婚者的态度比他所表露的温和得多。\Person{奥力振}、\Person{狄约尼削}、\Person{皮埃里乌}、\Person{凯撒勒雅的欧瑟比}、\Person{狄迪莫}、\Person{阿波里纳里}，在对这封书信的解释上使用了很大的自由度。当\Person{皮埃里乌}筛选并阐述宗徒的意思时，遇到这句话：\BibleQuote{我愿众人如同我自己一样}\BibleRef{格前 7:7}，他对此评论道：\Quote{说这话的保禄明显在宣讲禁绝婚姻。}这里的过错是我的吗？或者我应该为严厉负责？与\Person{皮埃里乌}的这句话相比，我所写的一切都显得温和。请查阅上述作者们的注释，并利用教会的图书馆；那么你将更快地以你希望的方式完成你已经很好开始的事业。\par

4. 我听说全城的希望都集中在你身上，主教和人民都一致希望你的高升。做主教已是大事，配得做主教则更为重大。\par

如果你阅读了我从希伯来文译成拉丁文的十六位先知的书，并且读完表示满意我的劳动，这个消息将鼓励我从书桌中取出一些现在锁在其中的其他作品。我最近将\WorkTitle{约伯传}译成了我们的母语：你可以从你的同族，圣善的\Person{玛尔塞拉}那里借到一份副本。请读希腊文和拉丁文两种版本，并将旧译本与我的翻译相比较。你将清楚看到它们之间的差别如同真理与谬误。我将一些对十二位先知的注释送给了可敬的司铎\Person{多姆尼欧}，还有\WorkTitle{列王纪}四卷——即两卷\WorkTitle{撒慕尔纪}和两卷\WorkTitle{玛拉基姆}。如果你愿意读这些，你自己便会明白理解圣经，特别是先知书是多么困难；以及由于译者的错误，那些对犹太人来说清楚流畅的段落，对我们却错漏百出。最后，你千万不要在我这些小作品中寻找那种你为基督的缘故而不屑于在\Person{西塞罗}作品中寻找的雄辩。为教会使用所作的译本，即使可能具有文学魅力，也应尽可能加以掩饰和避免；以便它不向哲学家的闲散学校和少数门徒说话，而是面向整个人类。\par

\SourceRecord{Translated by W.H. Fremantle, G. Lewis and W.G. Martley.；Nicene and Post-Nicene Fathers, Second Series；Vol. 6.；Edited by Philip Schaff and Henry Wace.；Buffalo, NY: Christian Literature Publishing Co.,；1893.；<http://www.newadvent.org/fathers/3001049.htm>.}

% structure: p0001 source=h1 target=section reason=page-title

\section{第五十封书信}

\Emph{致\Person{多米诺}}\par

\EditorialNote{多米诺，一位罗马人（在\WorkTitle{书信第四十五篇}中被称为\Quote{我们时代的命运}），曾写信告诉热罗尼莫，一位无知的隐修士在诽谤他反对约维尼安的著作。热罗尼莫在回信中严厉斥责了批评者的愚昧，并评论了他行为的不正直。他在这封信的结尾以强调的语气重申了他最初的立场。写于公元394年。}\par

1. 你的来信既充满了情意，也充满了抱怨。情意是你自己的，它促使你不停地警告我迫在眉睫的危险，并且使你为我\par

\Quote{对最安全的事物也充满疑虑和恐惧。}\par

抱怨是针对那些不喜爱我的人，他们在我的罪中寻找攻击我的机会。他们议论自己的弟兄，毁谤自己母亲的儿子。你写信告诉我这些事，不，是对一个特别的人：一个在街头、十字路口和公共场所都能见到的游荡者；一个隐修士，喧哗的多嘴者，唯一擅长的是诽谤，并且急于，不顾自己眼中的大梁，去掉邻人眼中的木屑。\BibleRef{玛 7:3}你还告诉我，他公开向我讲道，用牙齿啃噬、撕扯、撕碎我反对约维尼安所写的著作。你还告诉我，这个土生土长的辩证家，这个普劳提恩一派的支柱，既未读过亚里士多德的\WorkTitle{范畴篇}，也未读过他的\WorkTitle{解释篇}，也未读过他的\WorkTitle{分析篇}，更没有读过西塞罗的\WorkTitle{论题篇}；但是，由于他只在不学无术的圈子中活动，除了与软弱的女子交往之外不常去别处，他竟敢构造不合逻辑的三段论，并用微妙的论证来拆解他所谓的我的诡辩。我是多么愚蠢，竟以为没有哲学就不可能了解这些主题；并以为写作中更重要的部分是删除而不是书写！我徒然翻阅了亚历山大的注释；一位熟练的老师用波菲利的\WorkTitle{导论}教导我逻辑，也毫无用处；而且——轻看人类的学问——我竟因拥有纳西昂的额我略和狄迪莫斯作为我圣经的教理讲授者而一无所获。我学习希伯来文是徒劳的努力；我自幼每日研读法律与先知、福音与宗徒的辛劳也是同样徒然。\par

2. 此人无师自通，已达到完美，成为圣神的器皿和自学成才的天才。他超越西塞罗在口才上，亚里士多德在辩论上，柏拉图在谨慎上，阿里斯塔库斯在学问上，以及那个黄铜般的狄迪莫斯在著作数量上；不仅狄迪莫斯，而且在他那个时代所有作家在圣经知识上都比不上他。据说你只要给他一个题目，他总像卡尔内亚德斯一样，随时准备为正义或反对正义，为此方或彼方辩论。世界逃脱了一场大危险，民事诉讼和继承诉讼也从张开的深渊中被拯救出来，就在他轻视法庭、转向教会的那一天。因为，若他不愿意，谁能被证明无罪呢？若他一旦开始用手指计算案件的要点，铺开他的三段论之网，什么罪犯他的辩护会救不了呢？他只要跺跺脚，定定眼神，皱皱眉头，动动手，捋捋胡子，他立刻就会在陪审团眼中撒上灰尘。难怪这样一位完全的拉丁语学者和如此高深的口才家能战胜可怜的我，我——由于离开罗马已有一段时间，没有说拉丁语的机会——一半是希腊人，如果不是完全的蛮族的话。我说，难怪他能战胜我，因为他的口才已经压倒了约维尼安本人。善哉耶稣！什么！连约维尼安那个伟大而聪明的人！确实如此聪明，以至于没人能理解他的著作，当他歌唱时，只为自己——也为缪斯女神！\par

3. 我亲爱的父亲，请警告此人不要说出与他的誓言相悖的话，不要用言语破坏他服装所宣称的贞洁。无论他选择守童身还是已婚的独身——选择权在于他自己——他不可将妻子与童贞女相比较，因为那就等于白费功夫反对约维尼安的口才。我听说，他喜欢拜访寡妇和童贞女的小屋，并皱起眉头向她们讲授神圣文学。他在自己房间里究竟教这些可怜的女子些什么？是让她们确信童贞女不比妻子更好？是让她们充分利用青春年华，吃喝、常去浴室、生活奢侈、不嫌弃使用香水？还是向他们宣讲贞洁、斋戒和忽略自己的身体？无疑，他灌输的戒律充满德性。但若是这样，他应在公开场合承认私下所说的话。或者，若他的私下教导与公开一致，他就应当完全远离少女的社交。他是个年轻人——一个隐修士，在他自己眼中是个雄辩家（难道珍珠不是从他口中落下，他优雅的语句不是洒满了喜剧的盐和幽默吗？）——因此，我感到惊讶，他能毫不脸红地常去贵族府邸，频繁拜访已婚女士，使我们的宗教成为争论的话题，用曲解语言的方式歪曲基督的信仰，此外，还毁谤他在主内的弟兄。然而，他或许认为我有错误（因为\BibleQuote{我们在许多事上都有过失}，并且\BibleQuote{若有人在言语上没有过失，他就是完全人}\BibleRef{雅 3:2}）。那样的话，他应当写信来责备我或质疑我，就像潘马丘所采取的步骤，他是一位学识和地位都很高的人。对于后者，我尽己所能辩护，并在一封长信中解释了我的确切意思。他至少应当效仿你的谨慎，你曾提取并整理看似引起不快的内容，向我询问修正或解释，而不要以为我疯狂到在同一本书中既支持婚姻又反对它。\par

4. 愿他珍惜自己，愿他珍惜我，愿他珍惜基督徒的名号。愿他认清自己身为隐修士的身份，不是靠着空谈和争辩，而是藉着缄默与安静。愿他诵读耶肋米亚的话：\Quote{人最好在青年时背负上主，孤独地坐着默静无言，因为这是他加给他的担子。}\BibleRef{哀 3:27} 或者，如果他真有权利对一切作者施用审查之杖，并自以为是博学之士，因为只有他了解若委尼安（你知道这句谚语：巴尔布斯最懂巴尔布斯的意思）；然而，正如阿提利乌斯提醒我们的：\Quote{我们并非人人都是作者。} 若委尼安本人——一个前所未有、不学无术的文人——将最公正地向他指出这一点。他说：\Quote{主教们谴责我，}这不是理由，而是背叛。我不想要无名小卒的回答，他们虽有压制我的权力，却无教导我的才智。愿那与我作对的人，其语言我能理解，战胜他即战胜一切。\par

\Quote{我深知：请相信我，我曾感受到\newline{}英雄的力量，当他从盾牌上跃起，\newline{}掷出呼啸的投枪。}\par

他在辩论中坚强有力，缜密而固执，是那种埋头奋战的人。他常常在街头从深夜到凌晨向我叫喊。他体格健壮，肌肉如运动员。我私下相信他是我的教导的追随者。他从不脸红，也不斟酌言辞：唯一的目标就是尽可能大声说话。他以口才闻名，以致他的言语被当作模范教给那些卷发的青年。多少次，当我在集会中遇见他时，他激起我的愤怒，使我动怒！多少次，他朝我吐唾沫，然后唾沫满身地离去！但这些是粗俗的方法，我的任何追随者都能做到。我诉诸书籍，诉诸那些必须传于后世的纪念。让我们以著作说话，让沉默的读者在我们之间评判；而且，正如我有一群门徒，愿他也有一群——谄媚者和依附者，配得上他们的主人格纳托和福尔米奥。\par

5. 我亲爱的栋尼奥，在街角或药铺里喋喋不休，并对世界作出判断，并非难事。\Quote{某人演讲精彩，某人演讲拙劣；此人通晓圣经，那人疯狂；此人能言善道，那人一言不发。} 但谁认为他配如此评判每一个人？在每条街上对一个人大声疾呼，将不确定的指控而不是明确的罪名堆在他头上，这算不了什么。任何小丑或好讼之徒都能做到。让他伸出手，拿起笔，行动起来；让他写书，并在其中证明他所能证明的一切。让他给我机会回应他的雄辩。如果我愿意，我可以以牙还牙；当我受伤时，我可以用牙齿咬住对手。我也受过良好的教育。正如尤维纳尔所说：\Quote{我也常将手从戒尺下缩回。} 关于我，也可以用贺拉斯的话说：\Quote{逃离他；他角上有干草。} 但我宁愿作那位的门徒，祂说：\BibleQuote{我把背转给打击者……我没有遮掩我的脸，免受侮辱和唾污。}\BibleRef{依 50:6} 祂受辱骂，却不还口。\BibleRef{伯前 2:23} 在掌掴、十字架、鞭打、辱骂之后，在最后时刻，祂还为钉祂的人祈祷，说：\BibleQuote{父啊，宽赦他们罢！因为他们不知道他们做的是什么。}\BibleRef{路 23:34} 我也原谅一位弟兄的错误。我确信，他是被魔鬼的诡计欺骗了。在妇女中间，他被视为聪明而有口才；但当我的拙作到达罗马时，他畏惧我作为竞争对手，便试图夺走我的桂冠。他下定决心，世上没有任何人应该取悦他这位雄辩家，除非那些令人敬畏而非寻求喜爱的人，那些显得更令人畏惧而非受宠爱的人。作为一个手段高超的人，他渴望像老战士那样，一剑同时击倒他的两个敌人，并向每个人表明，无论他采取什么观点，圣经总是支持他。那么，他必须屈尊俯就，寄给我他的说明，并纠正我轻率的言辞，不是通过责难而是通过教导。如果他试图这样做，他就会发现，在沙发上看似有力的东西，在法庭上并不同样有力；在妇女的纺锤和针线篮之间讨论神圣律法的教义，与在有教养的人中间讨论它们是两回事。而事实上，他毫不迟疑、不知羞耻地一再高喊：\Quote{热罗尼莫谴责婚姻，}并且，当他经常穿梭于怀孕的妇女、啼哭的婴儿和婚床之间时，他压制了宗徒的话，只是为了给我——可怜的我——增添恶名。然而，当他日后动笔写书，与我短兵相接时，那时他就会感受到，那时他就会陷入困境；那时厄壁鸠鲁和亚里斯提卜不会在他身边；猪倌们不会来帮助他；多产的母猪连哼哼都不会。因为我也可以和屠尔努斯一起说：\par

\Quote{父亲，我也能投出有力的投枪，\newline{}当我击中时，伤口会流出鲜血。}\par

但如果他拒绝写作，并以为辱骂和批评同样有效，那么，尽管有陆地和海洋以及民族横亘在我们之间，他至少必须听到我呼喊的回声：\Quote{我不谴责婚姻，}\Quote{我不谴责婚配。} 的确——我这样说，是为了让他清楚明白我的意思——我希望每个人都娶妻，只要他们因为夜间害怕而无法独自入睡。\par

\SourceRecord{Translated by W.H. Fremantle, G. Lewis and W.G. Martley.；Nicene and Post-Nicene Fathers, Second Series；Vol. 6.；Edited by Philip Schaff and Henry Wace.；Buffalo, NY: Christian Literature Publishing Co.,；1893.；<http://www.newadvent.org/fathers/3001050.htm>.}

% structure: p0001 source=h1 target=section reason=page-title

\section{第五十一封书信}

来自塞浦路斯\Place{萨拉米斯}主教\Person{埃皮法尼乌斯}，致耶路撒冷主教\Person{若望}。\par

这两位主教之间曾因奥利振主义\OriginalTerm{奥利振主义}{Origenism}的争论产生冷淡，当时该争论正处高峰。\Person{埃皮法尼乌斯}公开指责\Person{若望}是奥利振主义者，并且违反教会法规授予\Person{热罗尼莫}的兄弟\Person{保利尼安}司铎圣职，为的是让白冷的隐修院从此完全独立于\Person{若望}。自然，\Person{若望}对此表示愤慨并流露不满。本信是\Person{埃皮法尼乌斯}对其所作所为的一种半道歉性质的信件，但如同所有此类信件一样，似乎只是使事情更糟。该争论在本书的\Quote{\WorkTitle{对耶路撒冷若望的驳斥}}中有详细记述，尤其见§11-14。\par

有一个有趣的段落（§9）叙述了\Person{埃皮法尼乌斯}在阿纳布拉塔撕毁了一幅教堂幔帐，上面绘有\Quote{基督或某位圣人的肖像}——这是早期圣像破坏精神\OriginalTerm{圣像破坏精神}{iconoclastic spirit}的实例。\par

本信原以希腊文写成，由\Person{热罗尼莫}应作者之请译为拉丁文。日期为公元394年。\par

致主主教及亲爱的兄弟\Person{若望}，\Person{埃皮法尼乌斯}致候。\par

1. 亲爱的弟兄，我们切不可滥用我们作为神职人员的身份，使之成为骄傲的借口，而应通过殷勤遵守天主的诫命，在实际上成为我们名义上所自称的人。因为，如果圣经说：\Quote{\BibleQuote{他们的阄分必不使他们得利}}，那么，我们这些不仅在思想感情上，而且在言语上犯罪的人，又能靠我们神职身份的骄傲得到什么益处呢？我听说你当然对我发怒、生气，并且威胁要写信说我——不仅仅寄往个别地方和省份，而是直到天涯海角。我们对天主的敬畏在哪里？那敬畏本应使我们因主所说的话而战栗——\Quote{\BibleQuote{凡向弟兄发怒而无理由的，应受审判}}\BibleRef{玛 5:22}。我并非特别在意你随意写什么。因为依撒意亚告诉我们，写在纸莎草上又抛在水里的信函——很快就会被时间与潮水冲走。我没有做过对你不利的事，没有加给你任何伤害，也没有用暴力从你那里强求过什么。我的行动涉及一座隐修院，其成员是外方人，不受你的教省管辖。此外，他们对我这卑微之人的敬意，以及我经常写给他们的信，已开始使他们不愿与你共融。因此，我觉得我若过于严格或谨小慎微，可能会使他们疏远拥有古老信仰的教会，于是我将其中一位弟兄祝圣为执事，等他履行了执事职务后，又收纳他进入司铎行列。我想你本应对此心怀感激，因为你一定知道，是天主的敬畏迫使我有此行动，尤其是当你记起天主的司祭职在任何地方都是相同的，我只是为教会的需要作了准备。诚然，教会的每一位主教都有他管辖下属的教会，没有人能超越自己的范围\BibleRef{格后 10:14}，但基督的爱——没有伪善的爱\BibleRef{罗 12:9}——给我们立了榜样；我们不仅要考虑所做的事，还要考虑做事的时间、地点、方式和动机。我看到隐修院中有众多可敬的弟兄，而可敬的司铎\Person{热罗尼莫}和\Person{味增爵}因谦逊和卑微，不愿行他们的职分所允许的祭献，也不愿在那一部分比任何其他部分更有益于基督徒得救的使命上劳苦。此外，我知道你无法找到或按手在这位天主的仆人身上，他曾多次从你那里逃走，只因他不愿意承担司铎的繁重职务，而且别的主教也不容易找到他。因此，我相当惊讶，当按照天主的安排，他与隐修院的执事及其他弟兄一起来到我这里，为我曾对他们有的某种抱怨或别的什么事作赔补。当集祷经\OriginalTerm{集祷经}{Collect}在毗邻我们隐修院的别墅教堂中举行时——他完全不知情，也毫不怀疑我的意图——我吩咐几位执事抓住他并堵住他的嘴，免得他急于挣脱而奉基督之名向我发誓。首先，我祝圣他为执事，将天主的敬畏摆在他面前，强迫他服务；他极力反抗，大叫自己不堪当，并抗议说这种重担超过了我的能力。我好不容易才克服了他的抗拒，用圣经段落尽力说服他，并将天主的诫命摆在他面前。当他在奉献神圣祭献时服务之后，我又一次费了很大力气堵住他的嘴，祝圣他为司铎。然后，我用先前同样的论据说服他坐在为司铎保留的位置上。此后，我写信给隐修院的可敬司铎及其他弟兄，责备他们没有就他的事写信给我。因为在一年前，我曾听到他们中许多人抱怨没有人能为他们举行主的圣事。于是全体一致同意请他承担这一职责，指出他对隐修院团体该有多么大的用处。我责备他们没有写信给我并要求我祝圣他，尽管他们本有机会这样做。\par

2. 这一切，如我刚才所说，我是依靠着那份基督徒的爱德而做的，我相信您对我这卑微之人怀有这份爱德；更不用说，我是在一座隐修院内举行祝圣礼，而不是在您管辖的范围内。塞浦路斯主教们的温和与谦逊，以及他们的纯朴，是多么有福啊！尽管在您的精致与辨别力看来，这似乎只值得天主的怜悯！因为许多与我共融的主教，在我的教省内祝圣了我无法捕获的司铎，并将执事和副执事送到我这里，我很乐意接纳。我自己也曾催促已故的、可纪念的斐洛主教，以及可敬的特奥普雷普斯主教，为基督的教会作出安排，在那些塞浦路斯的教会中祝圣司铎——这些教会虽然被认为属于我的教区，却恰好靠近他们——原因是我教省辽阔而分散。但就我而言，我从未祝圣女执事，也未将她们送入其他人的教省，更没有做过任何撕裂教会的事。那么，您为何认为有理由对我所做的这项天主工程如此愤怒和恼恨呢？这工程是为了建设弟兄们，而非破坏他们。\BibleRef{格后 10:8} 此外，我十分惊讶您对我的圣职人员所说的断言：您曾通过那位可敬的司铎、修道院院长额我略，给我捎来口信，要我不得祝圣任何人，而我答应服从，说：\Quote{我是少年人吗？还是我不懂教规呢？} 我以上主的话向您说实话：我对这一切一无所知，也从未听说过，完全不记得说过任何这类话。然而，由于我担心我不过是人，可能在这么多事情中忘记此事，我已询问了可敬的额我略和他的同伴则诺司铎。其中，院长额我略回答说他对此事毫不知情，而则诺说，鲁菲努斯司铎在散漫的谈话中说了这些话：\Quote{您认为这位可敬的主教敢祝圣任何人吗？}但谈话没有进一步。然而我，埃皮法尼乌斯，从未收过这个口信，也从未回复过。所以，亲爱的，不要让您的愤怒压倒您，也不要让您的恼怒左右您，以免您徒劳地困扰自己；以免您被认为是提出这个委屈，只是为了给另一种倾向的发作留出空间，从而为自己寻找犯罪的机会。正是为了避免这个，先知向主祈祷说：\Quote{不要使我的心偏向邪恶的言语，去为我的罪过找借口。}\par

3. 我还惊讶地听说，某些惯于搬弄是非、总是在听到的事情上添油加醋，以挑起兄弟间的委屈和纷争的人，成功使您不安，说我向天主献祭时，习惯为您这样祈祷：\Quote{主啊，赐予若望恩宠，使他正确相信。}请不要以为我如此无知，竟能如此公开地说这话。向您坦白实情，我最亲爱的兄弟，虽然我在心里不断使用这个祈祷，但我从未把它透露给别人的耳朵，以免似乎羞辱了您。但当我按照礼仪诵念圣事所需祈祷时，我就为众人，也为您和他人说：\Quote{守护他，使他宣讲真理，}或者至少说：\Quote{主啊，求祢赐予他祢的助佑，守护他，使他宣讲真理之言，}视话语的场合和特定祈祷的次序而定。因此，我恳求您，亲爱的，并俯伏在您脚前，求您应允我和您自己这一个祈祷：使您得救，正如经上记载，\BibleQuote{脱离这邪恶的世代}。\BibleRef{宗 2:40} 亲爱的，请远离奥利金的异端和所有异端。因为我看到，您所有的怒气都是针对我而发，仅仅因为我告诉您，您不应赞美那位是亚略的精神父亲、一切异端的根源和源头的人。当我恳求您不要误入歧途，并警告您后果时，您曲解了我的话，使我流泪悲伤；不仅是我，还有在场的许多其他公教徒。我认为这就是您此刻愤怒和激情的根源。为此，您威胁要发出反对我的信函，并将您的说法四处传播；如此，您为了捍卫您的异端而点燃人们反对我的激情，破坏了我向您显示的爱德，行事如此不慎，以至于让我后悔曾与您共融，并因此支持了奥利金的错误观点。\par

4. 我直言不讳。用圣经的话说，我不惜剜出使我失足的眼睛，砍下使我跌倒的手和脚。\BibleRef{玛 18:8} 你们不管是我的眼睛、我的手还是我的脚，也必须同样对待。因为哪个天主教徒，哪个以善行装饰自己信德的基督徒，能平静地听闻\Person{奥利振}的教导和劝诫，或相信他那奇异的讲道呢？他告诉我们：\Quote{圣子不能看见圣父，圣神不能看见圣子。}这些话出现在他的\WorkTitle{论首要原理}一书中；我们这样读到，\Person{奥利振}也这样说了。\Quote{因为正如说圣子能看见圣父不合宜，同样，认为圣神能看见圣子也不合宜。}再者，谁能容忍\Person{奥利振}的主张：人的灵魂曾经是天上的天使，因在上界犯罪而被投到这个世界，被禁锢在身体里如同在坟墓或冢中，为前世的罪过受罚；并且信徒的身体不是基督的殿，而是被定罪者的监牢？此外，他滥用寓意法曲解叙述的真义，无休止地增添空话；用五花八门的论证动摇质朴之人的信德。他时而主张灵魂（希腊语中称为\Quote{凉的东西}，源于一个意为\Quote{变凉}的词）之所以得名，是因为从天上降到下界时失去了原先的热力；时而又说我们的身体在希腊语中被称为\Quote{锁链}（源于一个意为\Quote{链条}的词），或（类比我们拉丁语中的词）\Quote{堕落的东西}，因为我们的灵魂从天上坠落；而希腊语丰富的习语所提供的另一个表示身体的词，被许多人理解为\Quote{坟墓的纪念碑}，因为灵魂被封闭在其中，如同死者的尸体被封闭在坟墓和冢中。如果这样的教义是真的，我们的信仰何在？复活的宣讲何在？宗徒的教导何在？这教导至今仍在基督的教会中存留。对亚当及其后裔、对诺厄及其儿子的祝福何在？\BibleQuote{你们要生育繁殖，充满大地。}按照\Person{奥利振}的说法，这些话必定是咒诅而非祝福；因为他把天使变成人的灵魂，强迫他们离开至高的处所，降至更低之处，仿佛天主仅凭祝福之工就不能赐予人类灵魂，除非天使犯罪；仿佛地上每次生育都必有一次天上的堕落。那么，我们就要放弃宗徒和先知、法律以及我们主救主自己的教导，尽管祂在福音中如雷贯耳地大声说话。而\Person{奥利振}却命令并敦促——且不说强制——他的门徒不要祈求升入天上，免得再次犯罪，比在地上罪更重，又被掷回世界。这种愚蠢狂乱的见解，他通常通过曲解圣经的意义，使之完全不是它们本来的意思来证实。他引用圣咏的这段经文：\BibleQuote{因我的邪恶，在祢谦卑我之前，我走错了路}；还有：\BibleQuote{我的灵魂，回到你的安息吧}；又有：\BibleQuote{领我的灵魂出离监牢}；以及\BibleQuote{我要在活人之地向上主认罪}，尽管神圣经文的意思与他的解释——他强行扭曲以支持自己异端——无疑不同。这种做法是摩尼派、诺斯替派、厄俾翁派、玛尔基翁派以及其他八十种异端信徒所共有的，他们都是从圣经的纯洁泉源中汲取证据，却不是按照写下来的意思解释，而是试图使教会著作的朴素语言符合自己的愿望。\par

5. 对于他所竭力主张的一个观点，我几乎不知道是应为之流泪还是大笑。这位奇妙的博士竟妄自教导说，魔鬼将再次恢复其昔日的状态，将返回他原有的尊严，并重新升入天国。哦，可怕！一个人竟如此疯狂愚昧，竟以为\Person{洗者若翰}、\Person{伯多禄}、宗徒兼福音作者\Person{若望}、\Person{依撒意亚}、\Person{耶肋米亚}以及其他先知，都将与魔鬼一同成为天国的共同继承人！我跳过他对\OriginalTerm{皮衣}{coats of skins}的空洞解释，也不提他为使我们相信这些皮衣代表人的肉身而付出的努力和论据。在许多其他事情中，他这样说：\Quote{天主难道是皮匠或马具匠，要为动物预备皮革，并为\Person{亚当}和\Person{厄娃}缝制皮衣吗？}他接着说：\Quote{显然，他指的是人的肉身。}如果是这样，那么在皮衣、不服从和从乐园堕落之前，\Person{亚当}如何并非以寓意而是以字义说：\BibleQuote{这才是我的骨中之骨，肉中之肉}；\BibleRef{创 2:23}。或者神圣叙述的依据是什么：\BibleQuote{上主天主使\Person{亚当}陷入沉睡，他就睡了；然后取了他的一根肋骨，并用肉填补了原处；上主天主用从人身上取来的肋骨，造了一个女人}\BibleRef{创 2:21}给他？或者\Person{亚当}和\Person{厄娃}在吃了禁树上的果子后，能用无花果树叶遮盖什么肉身？\BibleRef{创 3:7}。谁能耐心地听\Person{奥利振}危险的论点，当他否认这肉身的复活时（正如他在解释第一篇\BibleBook{}{圣咏}的书和其他许多地方所明确表明的）？或者谁能容忍他，当他把乐园放在第三层天上，并将圣经所提及的从地上转移到天上之处，并以寓意的方式解释\BibleBook{}{创世纪}中提到的所有树木，实际上说这些树是天使的\OriginalTerm{潜能}{potencies}，而这段经文的真正含义并不允许这种解释？因为神圣的圣经并没有说：\Quote{天主将\Person{亚当}和\Person{厄娃}安置在地上}，而是说：\Quote{祂将他们逐出乐园，并使他们住在乐园对面。}祂没有说\Quote{在乐园之下。}\BibleQuote{祂安置了\OriginalTerm{革鲁宾}{cherubims}和一把旋转的火焰剑，以看守生命树的道路。}\BibleRef{创 3:24}祂没有说上升（到那里）。\BibleQuote{有一条河从伊甸流出。}\BibleRef{创 2:10}祂没有说\Quote{从伊甸流下。}\BibleQuote{它分开并成了四条支流。第一条名叫\Place{丕雄}……第二条名叫\Place{基红}。}我亲眼见过\Place{基红}的水，用我肉体的眼睛见过。\Person{耶肋米亚}所指的正是这条\Place{基红}，他说：\Quote{你要在往\Place{埃及}的路上做什么？去喝\Place{基红}的混水吗？}我也喝过\Place{幼发拉底}大河的水，不是属灵的，而是实际的水，能用你的手触摸、用你的口饮用的。但既然有可以看见和饮用的河流，那里也必然有无花果树和其他树木；正是关于这些，主说：\BibleQuote{你可以随意吃园中所有树上的果子。}\BibleRef{创 2:16}它们与其他树木和木材一样，正如河流与其他河流和水一样。但若水是可见而真实的，那么无花果树和其他木材也必然是真实的，\Person{亚当}和\Person{厄娃}起初必然是以真实而非幻影的肉身被造的，并非如\Person{奥利振}要我们相信的，后来因他们的罪才接受了肉身。但你会说：\Quote{我们读到\Person{圣保禄}被提到第三层天，进入乐园。}你正确地解释这话：\Quote{当他提到第三层天，然后加上乐园一词时，他表明天在一个地方，乐园在另一个地方。}难道不应该人人拒绝并鄙视\Person{奥利振}那样的诡辩吗？他说天上的水——即穹苍以上的水——不是水，而是具有天使能力的英雄存有，而地上的水——即穹苍以下的水——则是相反种类的\OriginalTerm{潜能}{potencies}（即魔鬼）。如果是这样，为什么我们在洪水记载中读到天上的窗户开了，洪水的水泛滥？因此深渊的源泉开了，整个大地都被水覆盖。\BibleRef{创 7:11}\par

6. 噢！那些离弃\BibleBook{}{箴言}教训的人是何等的疯狂和愚昧，书中说：\BibleQuote{我儿，应遵守你父亲的命令，不要放弃你母亲的教训。}\BibleRef{箴 6:20}他们却转向错谬，对愚昧人说他要作他们的首领，也不鄙视愚昧人所说的愚昧话，正如圣经所证：\Quote{愚昧人说话愚昧，他的心思领悟虚妄。}我恳求你们，亲爱的诸位，并凭我对你们所怀的爱，我乞求你们——如同怜惜我自己的肢体一样——以言语和书信来成就经上所记：\Quote{上主，憎恨祢的人，我岂不是恨他们吗？起来反抗祢的，我岂不是忧愁吗？}\Person{奥力振}的话是仇敌的话，可憎可恶，相反天主及其圣徒；不仅是我引用的那些，还有无数其他。因为我现在无意反驳他所有的意见。\Person{奥力振}并没有活在我的时代，也没有抢夺我的东西。我并非因遗产或任何世俗之事而憎恶他或与他争吵；但——说实话——我悲伤，极度悲伤，看到我的许多弟兄，尤其是那些最有前途、已达到圣职最高品级的人，被他的雄辩所欺骗，被他最邪曲的教训所吞噬，成为魔鬼的食物，从而使这句预言应验了：\Quote{他嘲笑一切堡垒，他的食物是精选的，他聚集俘虏如同沙粒。}但愿天主释放你，我的弟兄，以及托付给你的基督的圣民，还有你身边所有的弟兄，特别是司铎\Person{鲁菲诺}，脱离\Person{奥力振}的异端及其他异端，以及它们所导致的灭亡。因为，若因一两个反对信德的字句，许多异端已被教会所弃绝，那么，制造如此邪恶的诠释和如此有害的教义以摧毁信德，实则自称为教会之敌的人，又怎能不被视为异端者呢！因为，在他其他的恶行中，他还妄言\Person{亚当}失去了天主的肖像，尽管圣经从未如此说过。若果真如此，世上所有的受造物就永远不会服从\Person{亚当}的后裔——即整个人类；然而，宗徒的话说：\BibleQuote{一切……已被人类制伏了。}\BibleRef{雅 3:7}因为，若人没有——连同他们对万物的权威——天主的肖像，万物就绝不会被人制伏。但神圣的圣经将这一点与赐给\Person{亚当}及其后裔的祝福之恩宠联系在一起。没有人能通过曲解词义妄言这恩宠只赐给了一个人，只有他（和他的妻子，即，当他由泥土形成时，她由他的一根肋骨造成）是按照天主的肖像受造的，而那些随后在母胎中受孕、并非如\Person{亚当}那样受造的人，并不拥有天主的肖像；因为圣经紧接着附上了这样的陈述：\BibleQuote{亚当活到二百三十岁，认识了自己的妻子厄娃，她生了一个儿子，按自己的肖像和模样给他取名舍特。}再者，在第十代，即二千二百四十二年之后，天主为维护祂自己的肖像，并表明祂赐给人类的恩宠仍存于他们内，颁布了这条诫命：\BibleQuote{凡有生命的肉……你们不可吃，流血的肉也不可吃。凡流人血的，我必向人追讨血债；因为我是按天主的肖像造了人。}从\Person{诺厄}到\Person{亚巴郎}，经过了十代，\BibleRef{创 11:10}从\Person{亚巴郎}到\Person{达味}，又过了十四代，\BibleRef{玛 1:17}这二十四代加起来共二千一百一十七年。然而，圣神在第三十九篇圣咏中，哀叹众人行走于虚妄之中，且屈服于罪恶，如此说：\BibleQuote{因为人人所行走的，无非是肖像。}同样，在\Person{达味}之后，他的儿子\Person{撒罗满}为王时，我们读到对天主的模样的类似提及。因为在那部以他名字命名的\BibleBook{}{智慧篇}中，\Person{撒罗满}说：\BibleQuote{天主创造了人，使他成为不朽的，并使他成为祂永恒性的肖像。}\BibleRef{智 2:23}此后大约一千一百一十一年，我们在\WorkTitle{新约}中读到，人并没有失去天主的肖像。因为宗徒兼主的兄弟\Person{雅各伯}——我前面已提到过他——为使我们不陷入\Person{奥力振}的圈套，教导我们人确实拥有天主的肖像和模样。他在对舌头的略微散漫的论述之后，接着说：\BibleQuote{它是一个不受管束的邪恶……我们用它赞颂天主父，也用它诅咒按照天主的模样受造的人。}\BibleRef{雅 3:8}此外，\Person{保禄}，那位\BibleQuote{被选器皿}\BibleRef{宗 9:15}，在他的宣讲中完全保持了福音的道理，教导我们人是按照天主的肖像和模样受造的。他说：\Quote{男人}，他说，\Quote{不应留长发，因为他是天主的肖像和光荣。}\BibleRef{格前 11:7}他仅提及\Quote{肖像}一词，但通过\Quote{光荣}一词说明了模样的本质。\par

7. 代替你所说若我能提出三个圣经证据就满意的，看我已给了你七个。那么，谁还能容忍奥力振的愚妄？我不愿用更严厉的措辞，免得自己变得像他或他的追随者一样——他们冒着灵魂的危险，擅自将他们头脑中首先出现的任何东西当作教条来断言，并向天主发号施令，而他们本应向祂祈祷或从祂学习真理。因为他们中有些人说，亚当先前所接受的天主的肖像在他犯罪时失掉了。另一些人猜测，天主圣子预定要从玛利亚取的身体就是造物主的肖像。有些人将这种肖像等同于灵魂，另一些人等同于知觉，还有的等同于德性。这些人把它当作圣洗，那些人断言人正是藉着天主的肖像而行使普世统御。他们像醉酒的人，时而吐出这个，时而吐出那个，而他们本应避免如此严重的风险，并藉着单纯的信仰获得救恩，不否认天主的话。他们本应将祂恩赐的确切知识以及祂按自己的肖像和模样创造人的特定方式留给天主。他们放弃了这条道路，使自己陷入许多微妙的问题中，并藉此陷入了罪的泥沼。但我们，亲爱的诸位，相信主的话，并知道天主的肖像留在所有人身上，我们将认识人在哪方面按祂的肖像受造留给了祂。不要让任何人被若望书信中的那段话所迷惑，有些读者不明白，那里说：\BibleQuote{现在我们是天主的子女，将来如何，还没有显明；但我们知道，当祂显现时，我们将要像祂，因为我们将看见祂实在怎样。} \BibleRef{若一 3:2} 因为这指的是那时将要显耀于祂圣徒的光荣\BibleRef{伯前 5:1}；正如在另一处我们读到\BibleQuote{从光荣到光荣}\BibleRef{格后 3:18}的话，这光荣圣徒们甚至在此世已经领受了凭据和一小部分。他们的首领是梅瑟，他的脸极其发光，明亮如太阳的光辉。其次是厄里亚，他被火马车接升了天\BibleRef{列下 2:11}，并未感到火焰的效力。斯德望被石头砸时，他的脸像天使的脸，众人都看见了。\BibleRef{宗 6:15} 我们在少数案例中验证的这一点，应理解为适用于所有人，为叫所写的得以应验。\Quote{凡自洁的，必被列在蒙福者中。} 因为，\BibleQuote{心里洁净的人是有福的，因为他们要看见天主。} \BibleRef{玛 5:8}\par

8. 既然如此，亲爱的诸位，要守护你自己的灵魂，并停止向我发怨言。因为神圣的圣经说：\BibleQuote{你们也不要彼此抱怨，像他们中有些人抱怨，就被蛇所灭。} \BibleRef{格前 10:10} 宁可向真理屈服，并爱我——我爱你们也爱真理。愿平安的天主，按祂的仁慈，赐给我们使撒殚被践踏在基督徒的脚下，\BibleRef{罗 16:20} 并避开一切邪恶的机会，使爱与平安的纽带不在我们之间断裂，也不使正确信仰的宣讲受到任何阻碍。\par

9. 此外，我听说某些人对我有这怨言：当我陪同你们到那称为贝特耳的圣地，与你们一起按照教会的惯例举行集祷时，我来到一个叫阿纳布拉塔的村庄，经过时看见那里有一盏灯点着。我问那是什么地方，得知是一座教堂，便进去祈祷，发现教堂门上挂着一块帘子，染色绣花。上面有一个人像，或是基督，或是某位圣人的像——我记不清是谁的像。我见此情景，不愿一个人像挂在基督的教堂里，违反圣经的教导，就把它撕了下来，并劝告该处的看守者把它用作某个穷人的殓布。他们却发怨言说，若我决意撕掉它，我就应当给他们另一块帘子作为替换。我一听这话，就答应给一块，并说会立即送去。此后有一些耽搁，因为我一直在寻找一块质量最好的帘子来代替原先的那块，并认为应当派人到塞浦路斯去买。现在我已送去所能找到的最好的一块，恳请你吩咐该处的长老从读经员手中收下我送去的帘子，然后下令，不得在任何基督的教会中挂起其它这类——与我们的宗教相悖的帘子。像你这样正直的人，应当谨慎除去一件不适合基督教会以及那些受托照管的基督徒的冒犯机会。当心加拉提亚的帕拉迪乌斯——此人曾是我所亲爱的，但现在亟需天主的怜悯——因为他宣讲和教导奥力振的异端；要留意他不要引诱任何受托于你照管的人误入他错谬教义的歧途。愿你在主内安康。\par

\SourceRecord{Translated by W.H. Fremantle, G. Lewis and W.G. Martley.；Nicene and Post-Nicene Fathers, Second Series；Vol. 6.；Edited by Philip Schaff and Henry Wace.；Buffalo, NY: Christian Literature Publishing Co.,；1893.；<http://www.newadvent.org/fathers/3001051.htm>.}

% structure: p0001 source=h1 target=section reason=page-title

\section{书信第五十二篇}

\Emph{致\Person{Nepotian}}\par

\EditorialNote{\Person{Nepotian} 是 \Person{Heliodorus}（见第十四封书信）的侄子。他与其叔父一样，弃军职而从圣职，当时是 \Place{Altinum} 的司铎，\Person{Heliodorus} 在那里任主教。此信是一篇系统论述司铎职责及其应遵循的生活准则的文章。它广为流传，并引起了许多人对 \Person{Jerome} 的愤慨。写于公元 394 年。}\par

1. 我亲爱的\Person{Nepotian}，你在海外来信中一再请求我为你勾勒几条生活准则，说明一个弃绝世俗服务而成为隐修士或神职人员的人，如何能持守基督的正道，不被引入邪恶的巢穴。我年轻时——甚至还是个孩子——当我以旷野的艰苦生活来遏制青春激情的最初冲击时，我曾向你可敬的叔父\Person{Heliodorus}写过一封规劝信，信中充满眼泪与怨言，向他展示了他所弃绝的朋友的感情。在那封信中，我扮演了适合我年龄的角色，并因我仍沉浸于修辞学家的方法与格言中，我用许多学苑的华丽辞藻来装饰它。然而现在，我的头灰白了，额头皱纹密布，下巴垂着像牛一样的肉垂，如\Person{维吉尔}所说：\par

\Quote{冰冷的血凝滞在我心周围。}\par

他在另一处唱道：\par

\Quote{老年带走一切，甚至智慧也带走。}\par

再往下一点：\par

\Quote{我如此多的歌已离我而去，\newline{}甚至我的声音也已弃我而逝。}\par

2. 但为了不显得我只引用世俗作品，请听圣经的神秘教导。\Person{大卫}曾是一个战士，但到七十岁时，年岁使他发冷，以致任何东西都不能使他温暖。于是从以色列全境寻找一个少女——\Place{叔南}女子\Person{阿彼沙格}——与王同寝，温暖他苍老的身体。\BibleRef{列上 1:1} 你若拘泥于那使人死的字句\BibleRef{格后 3:6}，这难道不像是某种闹剧故事或阿特拉笑剧中的粗俗笑话吗？一个发冷的老人裹在毯子里，却只在少女的怀抱中才暖和起来。\Person{巴特舍巴}仍然活着，\Person{阿彼盖耳}还在，圣经提到名字的其他妻子和妾还剩下。然而她们都被视作冰凉而弃绝，只有在这年轻女孩的拥抱中老人才变得温暖。\Person{亚巴郎}远比大卫年长，但只要\Person{撒辣}还在世，他未寻求别的妻子。\Person{依撒格}的岁数比大卫多一倍，却从未因与年老的\Person{黎贝加}同住而感觉寒冷。我不提洪水前的年代的人，他们虽活到九百岁，肢体不仅衰老且腐朽，却也未求助于少女的拥抱。以色列的领袖\Person{梅瑟}活到一百二十岁，却未从\Person{漆颇辣}身边更换。\par

3. 那么，这位叔南女子是谁？这位妻子和少女，如此炽热足以温暖寒冷，却又如此圣洁，不致在她所温暖的人身上激起情欲？\BibleRef{列上 1:4} 愿最智慧的\Person{撒罗满}告诉我们他父亲所爱的女子；愿和平之人向我们讲述战争之人的拥抱。\BibleRef{编上 28:3} \BibleQuote{你要获取智慧，获取明达：不要忘记，也不要离弃我口中的言语。不要离弃她，她就必保护你；喜爱她，她就必保守你。智慧是首要的，因此你要获取智慧，用你的一切所得获取明达。高举她，她就必使你升扬；你拥抱她，她就必给你带来尊荣。她必将华冠加在你头上，将荣冕赐给你。}\BibleRef{箴 4:5}\par

几乎一切身体的优点都随年龄改变，唯有智慧增长，而其他一切衰败。禁食、守夜、施舍变得更困难。睡在地上、四处迁徙、款待旅人、为穷人辩护、祈祷的恳切与恒心、探访病人、为施舍而做的体力劳动——总之，一切以身体为媒介的行为都随身体的衰退而减少。\par

现在，有一些青年仍充满活力，凭着辛劳与炽热的热情，以及圣善的生活和向主耶稣的不断祈祷，获得了知识。我不是指这些人，也不是说在他们身上智慧之爱是冷淡的——因为这在许多老年人身上因年龄而枯萎。我意思是，青春本身必须应对情欲的攻击，在邪恶的诱惑和肉身的刺痛中，如同火在青枝中被窒息，无法发展其应有的明亮。但当人们在青年时代从事可嘉的事业，并昼夜默想天主的法律时，他们随着时间学习，新的经验和智慧随年岁而来，因此从过去的追求中，他们的晚年收获喜乐之果。因此，那位希腊智者\Person{地米斯托克力}，在活了一百零七年后，意识到自己行将就木，据说他极为后悔，正当他开始变得智慧时却不得不离开生命。\Person{柏拉图}在八十一岁去世，笔仍握在手中。\Person{伊索克拉底}在文学与教学工作中活到九十九岁。我不提其他哲学家，如\Person{毕达哥拉斯}、\Person{德谟克利特}、\Person{色诺克拉底}、\Person{芝诺}和\Person{克雷安特}，他们在极其年迈时仍显示追求智慧的青年活力。我转向诗人：\Person{荷马}、\Person{赫西奥德}、\Person{西摩尼得斯}、\Person{斯特西克鲁斯}，他们都活到高龄，临终时却各唱一首比往常更甜美的天鹅之歌。\Person{索福克勒斯}因年迈及疏忽财产而被儿子指控为老糊涂时，他向法官宣读了他新近创作的《俄狄浦斯》剧，展现出如此大的智慧——尽管时光侵蚀——以致将法庭肃静的沉默变成了剧场的掌声。无怪乎监察官\Person{加图}，那位最雄辩的罗马人，在老年既不羞于学习希腊语，也未绝望于成功。至于荷马，他叙述从\Person{涅斯托尔}的舌头上——尽管已相当年老体弱——流出比蜜更甜的话语。\par

甚至阿彼沙格这个名字的神秘含义也指向老人更大的智慧。因为它的翻译是：\Quote{我的父亲是超越的}，或\Quote{我父亲的咆哮}。\Quote{超越的}一词是隐晦的，但在本段中指示卓越，暗示老人有更大的智慧储备，甚至因丰富而溢出。在另一段中，\Quote{超越的}与\Quote{必需的}形成对照。此外，阿彼沙格，即\Quote{咆哮}，恰当地用于波浪的声音，以及我们从大海听到的汩汩声。由此显而易见，神圣声音的隆隆声以超越人间声音的响亮回荡在老人的耳中。再者，按我们的语言，叔能的意思是\Quote{朱红}，暗示智慧之爱通过宗教研习变得温暖而炽热。虽然这颜色指向主宝血的奥秘，但也展现了智慧的温暖光芒。因此，在\BibleBook{}{创世纪}中，收生婆将朱红线拴在培勒兹手上——培勒兹\Quote{分开者}，如此称谓因为他分开了此前分隔两族的壁垒。\BibleRef{创 38:28}同样，以流血的奥秘为参照，朱红绳索是妓女辣哈布（教会的预表）在耶里哥毁灭时挂在窗户上以保全她房屋的。\BibleRef{苏 2:18}因此，在另一处，\WorkTitle{圣经}论及圣徒说：\BibleQuote{这些是从勒卡布父亲家的温暖中出来的。}在\WorkTitle{福音}中，主说：\BibleQuote{我来把火投在地上，我多么切望它燃烧起来！}\BibleRef{路 12:49}这火当在门徒心中燃起时，迫使他们说：\BibleQuote{当祂在路上与我们谈话，给我们讲解圣经的时候，我们的心不是火热的吗？}\BibleRef{路 24:32}\par

你要问，这些深奥的引证有何目的？为向你表明，你不必期待从我这里得到男孩式的夸夸其谈、华丽的情感、矫揉造作的风格，以及每个段落末尾那简短尖锐的格言——它们引起听众的喝彩。让智慧独独拥抱我；让她依偎在我怀中，我的永不衰老的阿彼沙格。她确实是纯洁无瑕的，永远是一位童贞女，因为虽然她日日受孕，不断地生产，却像玛利亚一样保持童贞。当宗徒说\BibleQuote{在圣神内要火热}\BibleRef{罗 12:11}，他是说\Quote{忠于智慧}。当我们的主在\WorkTitle{福音}中宣告，在世界穷尽时——按\Person{匝加利亚}的预言，牧人将变成愚拙\BibleRef{匝 11:15}——\BibleQuote{许多人的爱情必要冷淡}\BibleRef{玛 24:12}，祂的意思是智慧将衰败。因此，请听——引用\Person{圣西彼廉}的话——\Quote{有力而非优雅的言语}。请听这个人，他虽然与你同为圣职弟兄，但在年岁上是你的父亲；他能将你从信德的摇篮引向属灵的成年；当他逐步建立圣洁生活的规则时，也能通过教导你来教导他人。我当然知道，从你尊敬的舅父\Person{赫利奥多罗斯}（现今是基督的主教）那里，你已学到并每日学习一切神圣之事；在你身上，你拥有生活的准则和德行的模范。那么，请接受我的建议，不值什么，并将我的训诫与他相比。他将教给你隐修士的成全，而我将向你展示神职人员的全部职责。\par

那么，一个神职人员，当他服侍基督的教会时，必须首先明白他的名字意味着什么；然后，当他认识到这一点时，必须努力成为他被称的那一位。因为希腊词\OriginalTerm{份额}{\GreekText{κλῆρος}}意思是\Quote{份额}或\Quote{产业}，所以神职人员被称为\Quote{神职人员}，要么因为他们是主的份额，要么因为主自己是他们的份额和分得的部分。现在，在他自己身上是主份额的人，或有主作他份额的人，必须如此行为，以致拥有主并被主拥有。他拥有主，并与先知同说：\BibleQuote{上主是我的产业}，他不能有主以外的任何东西。因为如果他拥有主以外的东西，主将不是他的份额。假设，例如，他拥有金银、产业或镶嵌家具；有了这样的份额，主将不屑于成为他的份额。我，如果我是主的份额，祂产业的界绳，就不在其余支派中获得份额；但像司祭和肋未人一样，我靠什一之物生活，\BibleRef{户 18:24}并服务祭台，靠其祭品维持。\BibleRef{格前 9:13}只要有食粮和衣服，我就以此满足，\BibleRef{弟前 6:8}作为十字架的门徒，我将分担它的贫穷。我恳求你，因此，并且\par

我一再地，并且再三地劝诫你；\par

不要以你的军事经验为标准来衡量神职的义务。在基督的旗帜下，不要寻求世俗的利益，免得你拥有了比你最初成为神职人员时更多的东西，听到人们羞耻地说：\BibleQuote{他们的产业不会使他们受益。}欢迎穷人和陌生人在你简陋的饭桌旁，好让基督与你同席。一个从事商业、从贫穷变为富有、从卑微升到高位的神职人员，你要躲避他如同瘟疫。因为\BibleQuote{恶友会败坏善行。}\BibleRef{格前 15:33}你轻视金子；他喜爱金子。你鄙弃财富；他急切追求财富。你喜爱寂静、温良、隐居；他乐于闲谈、厚颜、广场、街道和药铺。在如此悬殊的品行差异中，哪里会有和谐的情感呢？\par

女人的脚几乎不应跨越你家的门槛。对所有基督的贞女，你当一视同仁或一概漠视。不要与她们同住一屋檐下，也不要倚靠你过去的节欲。你不能比达味更圣洁，也不能比撒罗满更有智慧。常记住：是女人将耕种乐园者逐出了他的产业。你若患病，可由弟兄中一人照料；或你的姐妹、母亲，或某位信德已蒙认可的妇人。但若你没有这样亲近的人，或行为贞洁的人，教会赡养着许多年长妇女，她们的服务可令你受惠，自己也获益，这样连你的疾病也能以施舍行为结出果实。我见过身体康复却预伏灵魂疾病的例子。你若时常注视一个人的面容，为她服务便会有危险。若在圣职工作中需要探访寡妇或贞女，切勿独自进屋。让你的同伴是那些与为伍不会使你蒙羞的人。若你带一名读经者、辅祭或咏唱者同行，让他们的品德而非衣着成为装饰；他们不可用卷发钳卷发，而应以举止彰显贞洁。你不得与女子独坐，也不可在无见证的情况下见她。若她有秘密要透露，她必会有某位保姆、管家、贞女、寡妇或有夫之妇。她不可能孤苦得只有你敢听她倾诉秘密。要提防一切引发猜疑的事物；为避免丑闻，当避开所有可能赋予它口实的行为。频繁赠送手帕、袜带、面巾、先尝过的菜肴——更不必说充满爱意的字条——这些事神圣的爱一无所知。那些诸如\Quote{我的蜜糖}\Quote{我的宝贝}\Quote{你是我一切的魅力和喜悦}的亲昵爱语，恋人们可笑的殷勤和愚蠢行为，我们在舞台上面红耳赤，在世俗人中深恶痛绝。在隐修士和圣职人员身上，他们以誓愿装饰司铎职，而司铎职又装饰了誓愿，我们更当憎恶这些。我说这些，并非担心你或圣洁生活者会遭遇这类恶事，而是因为在各种阶层、职业、男女中，都有善有恶；谴责恶的同时，我赞扬善。\par

6. 说来可耻，偶像司祭、戏子、骑师和娼妓可以继承财产；唯独圣职人员和隐修士受到法律限制——这限制并非由迫害者，而是由基督徒皇帝所立。我不抱怨法律，但我痛心我们竟应受如此严厉的律令。烙铁无疑是好的，但我为何有伤口需要它？法律是严明而有远见的，即便如此，贪婪仍不受约束。我们凭信托的虚构蔑视律法；仿佛皇帝谕令重过基督的诫命，我们畏惧法律而轻视福音。若必须有继承人，母亲对子女有首要权利，教会对其羊群——她所生、所养、所哺育的成员——也有权利。我们为何要闯入母亲与子女之间？\par

主教的光荣在于为穷人所需做准备；但所有司铎的耻辱是积聚私产。我，假设生在穷人家，乡间茅屋，几乎无法以普通小米和家常面包填饱空腹，如今却嫌弃精白面粉和蜂蜜。我熟记各种鱼类的名目。我能无误地说出贻贝在哪个海岸被拾起。我能凭味道区分飞鸟来自哪个行省。美味佳肴令我愉悦，因其原料珍稀，最终我却在它们昂贵的代价中找到乐趣。\par

我也听说有些人向无子女的老人和妇女献奴性的关怀。他们端盆、围床，亲手做最令人作呕的事。他们焦急等待医生到来，颤抖着嘴唇询问病人是否好转。若老人短暂显出恢复活力的迹象，他们便痛苦万分。他们假装高兴，但贪婪的心却暗受折磨。他们害怕努力白费，将留恋生命的老者比作默突舍拉始祖。若他们不将心放在暂时奖赏上，在天主前可得多大的赏报！他们以何等巨大的劳苦追逐空虚遗产！少些劳力本可为他们购得基督的珍珠。\par

7. 时常阅读神圣圣经；勿让圣卷离手。学习你所要教导的。\Quote{坚守那合乎道理、确实的圣言，正如你所受的教导，好能以健全的道理劝勉并驳斥反对者。继续你所学和所确信的事，因为你知道你是从谁学来的；}并\Quote{时常准备好，向询问你们心中盼望和信德的理由的人，提出答辩。}\BibleRef{伯前 3:15} 不要让你的行为否定你的言语；免得你在教会讲道时，有人心里回应说：\Quote{你为何不实践你所宣讲的？这是一位贪爱美食者成了监察者！他腹中饱满，却向我们宣讲禁食。好比强盗指责他人贪婪。}在基督的司铎身上，口、心和手应一致。\par

要服从你的主教，视他为你灵魂的慈父。儿子爱父亲，奴隶畏惧主人。祂说：\BibleQuote{我若是父亲，我的尊荣在哪里呢？我若是主人，我的敬畏在哪里呢？}\BibleRef{拉 1:6} 在你的情形中，主教集多种应受尊敬的身份于一身：他同时是隐修士、长上、叔父，而且从前已教导你一切神圣之事。我也要说：主教应当知道自己不是主人而是司铎。他们应给圣职人员应得的尊敬，好让圣职人员也向主教献上该有的尊重。演说家多米提乌斯有一句妙语可引用于此：\Quote{为何我应认你为元老院领袖，而你不肯认我作普通成员的权利呢？}我们应认识到，主教和他的长老们如同亚郎和他的儿子们。既然只有一位主、一座圣殿，也应只有一项职务。让我们常记宗徒伯多禄给司铎的训诫：\BibleQuote{务要牧养你们中间的天主的羊群，照管他们，不是出于勉强，而是出于甘心，正如天主所愿；不是为卑鄙的财利，而是出于热忱；也不是作天主的羊群的主人，而是作羊群的榜样，}并且要甘心乐意，好使\BibleQuote{当总司牧显现时，你们可以领受那永不衰残的荣耀的冠冕。}\BibleRef{伯前 5:4} 某些教会有一种不良风俗，即主教在场时长者保持沉默，仿佛主教会是妒忌或不耐烦的听众。宗徒保禄写道：\BibleQuote{若旁边坐着的得了启示，那先说话的就当闭口不言。因为你们都可以一个一个地说预言，叫众人学习，叫众人得劝勉；且先知的灵是顺服先知的，因为天主不是叫人混乱，而是叫人平安。}\BibleRef{格前 14:30} \BibleQuote{智慧之子使父亲欢乐；}\BibleRef{箴 10:1} 主教应因自己的辨别力而喜悦，因他拣选了这样的人作基督的司铎。\newline{}\par

8. 在教会教导时，要力求唤起叹息而非喝彩。让听众的眼泪成为你的荣耀。长老的话语应由其诵读圣经来调味。不要作一个朗诵者或夸夸其谈者，一个毫无章法地瞎扯的人；而要显出你精通深奥之事，熟悉天主的奥秘。信口开河、语速飞快以致令无知的群众惊异，这是无知的标志。自信常常解释它一无所知的事；而当它说服了别人时，也欺骗了自己。我的老师纳齐安的额我略，当我曾请他解释路加所说的\OriginalTerm{第二个第一的安息日}{\GreekText{σάββατον} \GreekText{δευτερόπρωτον}}时，他开玩笑地回避了我的请求说：\Quote{我将在教会中告诉你，在那里，当全体民众为我喝彩时，你将被迫不得不承认你完全不知道的事。因为，若你独自保持沉默，人人都会把你当作傻瓜。}没有比单凭口若悬河来欺骗普通群众或未受教育的会众更容易的事了：他们最钦佩自己不懂的东西。请听马尔库斯·图利乌斯，那位享有崇高颂词的主题：\Quote{你会成为第一流的演说家，若不是有德摩斯梯尼；而他也会成为唯一的演说家，若不是有你。}请听他在为昆图斯·加里乌斯的辩护词中关于不熟练的演讲者和大众喝彩的言论，那样你就不会被这些幻觉所愚弄。他说：\Quote{我告诉你们的是我最近的一次亲身经历。一位有着诗人和文化人之名的人写了一本书，题为\WorkTitle{诗人和哲学家的对话}。书中他将欧里庇得斯与米南德、苏格拉底与伊壁鸠鲁置于对话中——这些人的生活我们知道不仅相隔多年，而且相隔若干世纪。然而他引发了无限的喝彩和不断的欢呼。因为剧院中有许多与他同派别的人：他们和他一样从未学过文字。}\newline{}\par

9. 在衣着上，既要避免暗沉的颜色，也要避免亮丽的颜色。炫耀和邋遢同样应当避免；因为前者虚浮，后者骄傲。不戴亚麻围巾算不得什么；值得称赞的是没有钱买它。既没有餐巾也没有手帕却炫耀，却带着装满的钱袋，这是可耻和荒谬的。\newline{}\par

有些人给穷人一点施舍，为的是自己获得更多，在施舍的伪装下只是寻求财富。这样的人是施舍的追逐者而非施舍者。他们的方法如同捕鸟、野兽和鱼的方法：在钩上放一小块饵——为了钓取一位已婚妇人的钱包！教会已托付给主教；让他留心指派谁作他的施赈员。对我来说，宁可没有钱可施舍，也不要不以为耻地乞讨我打算囤积的东西。而希望显得比基督的主教更慷慨，这也是傲慢。\Quote{并非所有事对我们所有人都是公开的。}在教会中，一个是眼睛，另一个是舌头，另一个是手，另一个是脚，还有耳朵、肚子等等。请读保禄致格林多人书，学习身体如何由不同肢体组成。\BibleRef{格前 12:12} 粗野愚钝的弟兄不可因无知而自以为圣人；有学问、有口才的人也不可仅凭口才来衡量自己的圣德。在两样不完美的品性中，神圣的质朴要好于有罪的雄辩。\newline{}\par

10. 如今许多人建造教堂；他们用闪亮的大理石筑墙和柱，天花板以金装饰，祭台镶嵌宝石。然而对基督仆役的遴选却毫不关心。但愿无人以犹大殿宇的财富、其桌子、灯台、香炉、盘碟、杯爵、调羹以及其余金器来反驳我。如果这些都蒙主悦纳，那是在司祭必须献祭、羊血曾赎罪的时代。它们是预表未来之事的预像，并且是\BibleQuote{为劝诫我们这些末世的人而写的}\BibleRef{格前10:11}。但如今，我们的主以其贫穷圣化了他居所的贫穷。因此，让我们思念他的十字架，视财富如粪土。我们为何钦佩基督所称的\BibleQuote{不义的钱财}\BibleRef{路16:9}？我们为何珍爱并喜欢那些伯多禄所夸耀不拥有的东西\BibleRef{宗3:6}？或者，如果我们坚持拘泥字句，且觉得提及金银财富如此愉快，就让我们一同持守金银和其他一切吧。让基督的主教们必须娶妻，且是童贞女\BibleRef{肋21:14}。让最善意的司祭若带有伤痕或容貌毁损就被剥夺职务\BibleRef{肋21:17}。让身体的癞病被认为比灵魂的斑点更严重。让我们生育繁殖，充满大地\BibleRef{创1:28}，但让我们不宰杀羔羊，不庆祝逾越节\BibleRef{申16:5}，法律禁止这些行为。让我们在七月搭帐篷\BibleRef{肋23:40}，并用号角声传扬严肃的斋戒\BibleRef{岳2:15}。但是，如果我们以属神的事与属神的事相比\BibleRef{格前2:13}，并与保禄一同承认\BibleQuote{法律是属神的}\BibleRef{罗7:14}，并记起达味的话：\BibleQuote{求你开启我的眼睛，使我得以看出你法律的奇事}；如果在此基础上，我们如主所解释的那样解释它——他这样解释了安息日\BibleRef{玛12:1}——那么，弃绝犹太人的迷信时，我们也必须弃绝金子；或者，若赞同金子，我们也须赞同犹太人。因为我们要么连同金子一同接受他们，要么连同金子一同谴责他们。\par

11. 避免招待世俗之人，尤其是那些因官位而趾高气扬者。你是基督的司祭——贫穷而受钉者，靠陌生人的面包为生。如果执政官的侍从或兵士在你门前守望，如果行省的总督在你家吃得比在他府上更好，这是一件耻辱。如果你以愿意为不幸者和受压迫者说情为借口，我回答：一个世俗的法官会更尊重一个克己的教士，而不是一个富有的教士；他会更看重你的圣德而非你的财富。或者，如果他是一个不在杯觥之前就不愿听教士为困苦者陈情的人，我宁愿放弃他的帮助，而向基督求助，他能比任何法官更有效且更快地帮助人。诚然，\BibleQuote{倚靠主胜过倚靠世人；倚靠主胜过倚靠权贵}。\par

愿你的呼吸从不带有酒味，以免别人对你说：\Quote{你给我的与其说是亲吻，不如说是酒的滋味。}贪酒的司祭被宗徒谴责\BibleRef{弟前3:3}，也被旧律禁止。据说，供职于祭台的人不可喝酒和\OriginalTerm{烈酒}{shechar}。在希伯来语中，所有致醉的饮品都称为\OriginalTerm{烈酒}{shechar}，不论是用谷物或苹果汁酿造的，不论你是从蜂巢蒸馏出一种粗糙的蜜酒，还是用椰枣榨出液汁，或是从谷物煎汁中滗出浓糖浆。任何致醉并扰乱心智平衡的东西，你都要像酒一样回避。我并不是说我们要谴责天主的受造之物。主自己曾被称作\Quote{贪酒的人}，并且出于胃弱的缘故，酒被适度允许给弟茂德。我只要求饮酒者遵守其年龄、健康或体质所要求的限度。但是，如果我完全不喝酒，却因青春而热血沸腾，因血气旺盛而体质强壮且精力充沛，我宁愿放弃那杯，因为我不能不怀疑其中有毒。希腊人有一句精妙的谚语，或许可以翻译如下，\par

\Quote{肥腹无柔情。}\par

12. 只给自己安排你能承受的斋戒，让你的斋戒纯洁、贞洁、简单、适度，而非迷信。如果不使用油，却追求最麻烦和最珍稀的食物——无花果干、胡椒、坚果、椰枣、精面粉、蜂蜜、开心果——这有什么益处呢？菜园的一切资源被尽力利用，以免我们吃家常面包；追求美味，我们背向天国。我听说有些人颠倒自然和人类的规律：他们既不吃面包也不喝水，只喝磨碎草药的稀汤和甜菜汁——不是用杯，而是用贝壳。真羞耻，我们对这样的愚行毫不脸红，对这样的迷信毫不厌恶！最可耻的是，我们在美味中间还寻求节制的美名。最严格的斋戒是面包和水。但因为这不带来荣耀，且我们所有人都靠面包和水生活，这就不被当作斋戒，而仅是日常普通的事。\par

13. 不要追求赞词，免得你赢得众人喝采的同时，却蔑视了天主。宗徒说：\BibleQuote{如果我还求人的欢心，我就不是基督的仆人了。}\BibleRef{迦 1:10} 当他成为基督的仆人时，他便不再求人的欢心。基督的士兵在美名与恶名中前进，\BibleRef{格后 6:8} 一左一右。没有赞扬能使他得意，没有侮辱能将他压垮。他不因财富而骄傲，也不因贫穷而畏缩。他同样鄙视喜乐与忧愁。白日太阳不伤他，夜里月亮也不害他。不要在街角祈祷，\BibleRef{玛 6:5} 免得人的喝采打断你祈祷的正直。不要放宽你的衣穗，为炫耀而佩戴\OriginalTerm{经匣}{phylacteries}，\BibleRef{玛 23:5} 也不要违背良心，像法利塞人那样追求私利。你愿知道主要求你穿什么服装吗？要有明智、正义、节制、刚毅。让这些成为你四方的地平线，让它们成为驮你的四马战车，使你——基督的御者——全速奔向目标。没有比这些更宝贵的项链；没有宝石能形成更灿烂的星群。藉着它们，你被装饰，被环绕，被四面保护。它们既是你的防御，也是你的光荣；因为每一颗宝石都变成了盾牌。\par

14. 也要当心饶舌的舌头和发痒的耳朵。既不要诋毁他人，也不要听信诋毁者。圣咏作者说：\Quote{你坐着，毁谤你的兄弟，诬蔑你母亲的儿子。你做了这些事，我默不作声；你便以为我像你一样，但我要责备你，将这一切摆在你眼前。} 要管住你的舌头，避免挑剔，谨守你的言语。要知道，你在判断别人时，便是定自己的罪；你所指责的过错，自己也有份。不要借口说：\Quote{如果别人告诉我，我不能无礼对待他们。} 没有人愿意对不情愿的听者说话。箭从不射入石头，常常反弹到射手身上。让诋毁者从你不愿听的态度中学会不轻易诋毁。\Person{撒罗满}说：\Quote{不要接近好诽谤的人，因为他们的灾祸必突然来临；谁知道两者的毁灭呢？} ——即诋毁者和听信诋毁者的人的毁灭。\par

15. 你有责任探望病人，了解已婚妇女的家庭和子女，守护显要人物的秘密。因此，要致力于保持你的舌头和眼睛一样纯洁。绝不谈论女人的身材，也不要让一个家庭知道另一个家中的事。\Person{希波克拉底}在教导学生之前，先让他们宣誓，迫使他们向他效忠。他以沉默约束他们，并规定他们的言语、步态、服装和举止。何况我们这些受托照管灵魂医药的人，更应当爱所有基督徒的家如同自己的家。让他们知道我们是悲伤中的安慰者，而非欢乐时的客人。那位常被邀请赴宴从不拒绝的司铎，很快便会受人轻视。\par

16. 我们不要寻求礼物，即使被请求接受也极少接受。因为\BibleQuote{施予比领受更为有福。}\BibleRef{宗 20:35} 不知何故，那个请求你接受礼物的人，往往因你接受而轻视你；而如果你拒绝，他会奇妙地更加尊重你。宣讲节德的人不应做婚姻的撮合者。为什么读宗徒的话\BibleQuote{今后有妻子的，要像没有一样}\BibleRef{格前 7:29}的人，却催促童女结婚？为什么身为\OriginalTerm{一夫一妻者}{monogamist}的司铎\BibleRef{弟前 3:2}，却怂恿寡妇再嫁？司铎们被吩咐要不顾自己的利益，又怎能管理他人的家务？从朋友那里抢夺东西是偷窃，而欺骗教会则是亵圣。当你收到要施舍给穷人的钱时，如果犹豫不决，而群众正在挨饿，你就比强盗更坏；从中克扣一部分，更是犯下最严重的罪行。我因饥饿而痛苦，你凭什么判断什么能满足我的渴望？要么立刻分发你所收到的，要么你作为胆怯的\OriginalTerm{救济品分发者}{almoner}，就让捐赠者自己去分配他的礼物。在你口袋满满的时候，我不应还有所需。没有人比我更能照顾我的需要。最好的救济者是那个不留任何东西给自己的人。\par

17. 我亲爱的\Person{内波提安}，尽管我关于\WorkTitle{童贞}的论著——我在\Place{罗马}为圣善的\Person{欧斯托钦}所写的那一篇——曾遭受抨击，但你仍迫使我在十年之后，再次在\Place{白冷}开口，并暴露在众舌的刺伤之下。因为只有什么都不写才能逃避批评，但你的请求使这不可能；我提笔时就知道所有反对者的箭都会射向我。我恳求他们保持沉默，停止反对；因为我写信给他们，不是作为对手，而是作为朋友。我没有责骂犯罪的人，只是警告他们不要再犯罪。我对自己的判断和对他们的判断一样严格。当我希望除去邻人眼中的木屑时，我先拔出了自己眼中的大梁。\BibleRef{玛 7:3} 我没有诽谤任何人，没有暗示任何名字。我的话不是针对个人，我对缺点的批评完全是普遍的。如果有人想对我发怒，他首先得承认自己符合我的描述。\par

\SourceRecord{Translated by W.H. Fremantle, G. Lewis and W.G. Martley.；Nicene and Post-Nicene Fathers, Second Series；Vol. 6.；Edited by Philip Schaff and Henry Wace.；Buffalo, NY: Christian Literature Publishing Co.,；1893.；<http://www.newadvent.org/fathers/3001052.htm>.}

% structure: p0001 source=h1 target=section reason=page-title

\section{Letter 53}

\Emph{To \Person{Paulinus}}\par

\EditorialNote{热罗尼莫敦促诺拉的保利诺主教（参见第五十八封书信）勤奋研读圣经，为此提醒他不仅当效法外邦智者，更当效法宗徒保禄的求学热忱。随后，他详细叙述了两个约中的各卷书籍及其内容与教训。最后，他恳请保利诺完全舍弃世俗财富，全心奉献于天主。此信写于公元三九四年。}\par

1. 我们的兄弟\Person{安波罗削}随同您的小礼物给我带来了一封极为令人愉快的信函；这信虽在友谊之初到来，却保证了历经考验的忠诚和长久的依恋。一种由基督亲自缔结的真正友谊，不依赖于物质利益，不依赖于对方的在场，也不依赖于虚伪和夸张的奉承；而是像我们这样的友谊，由对天主的共同敬畏和对神圣圣经的共同研习所铸成。\par

我们在古老传说中读到，人们穿越行省、横渡海洋、访问异族，只为亲眼见到那些仅从书本上认识的人。\Person{毕达哥拉斯}拜访了\Place{孟斐斯}的先知们；\Person{柏拉图}除了访问\Place{埃及}和\Place{塔兰托}的\Person{阿尔基塔斯}，还极为仔细地探索了\Place{意大利}沿岸那片古称\Place{大希腊}的区域。就这样，那位其学说响彻\Place{学园}殿堂的富有影响力的\Place{雅典}大师，同时成了朝圣者和学生，谦逊地选择学习他人所教，而不是过于自信地提出自己的观点。事实上，他对学问的追求——似乎满世界地在他面前飞逃——最终导致他被海盗俘虏，卖身为奴与一位暴君。于是，他成了囚犯、奴仆和奴隶；然而，由于他始终是一位哲学家，他仍比购买他的人更伟大。我们又读到，某些贵族从\Place{西班牙}和\Place{高卢}最遥远的地方旅行前去拜访\Person{蒂托·李维}，聆听他那如奶泉般涌流的雄辩。就这样，一个人的名声比城市本身的魅力更能吸引人们前往\Place{罗马}；那个时代出现了一个前所未闻的、值得注意的奇观：人们进入大城，却把注意力不是放在城市上，而是放在别的东西上。\Person{阿波罗尼奥斯}也是一个旅行者——我指的是那个被普通人称为术士、而被追随毕达哥拉斯的人称为哲学家的那位。他进入\Place{波斯}，穿越\Place{高加索}，途经\Place{阿尔巴尼人}、\Place{斯基泰人}、\Place{马萨格泰人}和\Place{印度}最富饶的地区，最后渡过那条宽阔的河流\Place{丕宋河}\BibleRef{创 2:11}，来到\Place{婆罗门人}那里。在那里，他看见\Person{希阿尔卡斯}坐在金宝座上，喝着\Place{坦塔罗斯}的泉水，并听他给少数门徒讲授天体的性质、运行和轨道。此后，他旅行途经\Place{厄兰人}、\Place{巴比伦人}、\Place{迦勒底人}、\Place{玛代人}、\Place{亚述人}、\Place{帕提亚人}、\Place{叙利亚人}、\Place{腓尼基人}、\Place{阿拉伯人}和\Place{培肋舍特人}；然后回到\Place{亚历山大里亚}，又前往\Place{埃塞俄比亚}，去探望\Place{赤裸智者}和那片著名的在沙漠沙土中铺展的\Place{太阳之桌}。所到之处，他都学到一些东西，并且由于他不断前往新地方，他变得越来越有智慧和更好。\Person{斐洛斯特拉托斯}用了八卷书详细记录了他的一生。\par

2. 但我何必只限于提及世间人物呢？宗徒\Person{保禄}，蒙选的器皿\BibleRef{宗 9:15}，外邦人的导师，能大胆地说：\BibleQuote{你们寻求基督在我内说话的凭证吗？}\BibleRef{格后 13:3}，他知道在他内确实拥有那位最伟大的客人——即使在他访问\Place{大马士革}和\Place{阿拉伯}之后，\BibleQuote{他上\Place{耶路撒冷}去见\Person{伯多禄}，同他住了十五天}\BibleRef{迦 1:17}……又过了十四年，他带着\Person{巴尔纳伯}和\Person{弟铎}，把他的福音传给宗徒们，免得他白白地奔跑或曾经奔跑\BibleRef{迦 2:1}。口传的话语具有一种难以言说的隐藏力量，直接从说话者之口传入听者耳中的教导比任何其他形式都更令人印象深刻。当\Person{德摩斯梯尼}反对\Person{埃斯基涅斯}的演说在后者流放\Place{罗得岛}期间向他诵读时，在一片赞赏和掌声中，他叹道：\Quote{要是你能亲耳听到那头野兽自己发表他的句子该多好！}\par

3. 我提出这些例子并非因为我有什么可以让你学到功课，而是要向你表明，你的热忱和好学——即使你无法依靠我的帮助——本身就是值得称赞的。一个愿意学习的心灵，即使没有老师，也值得赞赏。对我来说重要的不是你找到了什么，而是你寻求什么。蜡是柔软的，即使没有工匠和模压师的手去塑造它，它也很容易成型。它已经潜在地成为它所能成为的一切。\Person{宗徒保禄}曾坐在\Person{加玛里耳}脚前学习\Person{梅瑟}法律和先知，并且欣然如此，因为披上这属灵的铠甲，他能勇敢地说：\BibleQuote{我们争战的武器不是属血肉的，而是靠着天主有能力攻破坚固的营垒}；我们以此争战，\BibleQuote{攻破诡辩，以及一切高举起来反对天主知识的傲慢之物，并将一切思想掳去，使之服从基督；并且预备好惩罚一切不顺服的事}。\BibleRef{格后 10:4}他写信给\Person{弟茂德}，\Person{弟茂德}自幼受教于圣经，他劝勉他殷勤研读，\BibleRef{弟后 3:14}不要忽略因长老团按手所赐的恩典。\BibleRef{弟前 4:14}他又命令\Person{弟铎}，在主教的其他德行（他简要描述了）中，应谨慎寻求圣经的知识：他说，主教必须持守\BibleQuote{那按照教训所教导的忠信之言，以便能藉着纯正的道理劝勉人，并驳斥那反对的人}。\BibleRef{铎 1:9}事实上，神职人员缺乏学识，除了自己外不能帮助任何人；尽管他生活的德性也许能建立基督的教会，但他因无力抵抗那些试图拆毁教会的人，所造成的伤害也同样巨大。先知\Person{哈盖}说——更确切地说，主藉\Person{哈盖}的口说——\BibleQuote{现在你们询问司铎关于法律的事}。\BibleRef{哈 2:11}因为司铎的重要职务就是回答那些询问法律的人。在\WorkTitle{申命记}中我们读到：\BibleQuote{询问你的父亲，他必指示你；你的长老，他们必告诉你}。\BibleRef{申 32:7}也在\WorkTitle{圣咏第一一九篇}中：\BibleQuote{你的诫命成了我在临时居所中的诗歌}。\Person{达味}在描述义人时，将他比作乐园中的生命树，在其他优点中也提到：\BibleQuote{他喜爱上主的法律，昼夜默想祂的法律}。\Person{达尼尔}在最庄严的异象结束时宣告：\BibleQuote{义人必发光如同星辰；那些智慧者，即有学问的人，必发光如同穹苍}。\BibleRef{达 12:3}因此，你可以看到，无知的义人与有学识的义人之间有多大的区别。拥有前者的被比作星辰，拥有后者的被比作诸天。然而，按照希伯来原文的精确含义，这两种说法都可以理解为指有学识的人，因为原文是这样读的：\BibleQuote{那些智慧的人必发光如同穹苍的光辉；那些引导多人归义的人必如同星辰直到永远}。为什么\Person{宗徒保禄}被称为特选的器皿？\BibleRef{宗 9:15}无疑是因为他是法律和圣经的宝库。我们主的渊博教导使法利塞人惊得哑口无言，他们惊讶地发现\Person{伯多禄}和\Person{若望}虽未学习文字，却通晓法律。因为对于他们，圣神立即启示了那些别人通过每日学习默想才能得到的东西；正如经上所记，\BibleRef{得前 4:9}他们\BibleQuote{是受天主教导的}。救主刚满十二岁时，发生了圣殿中的事；\BibleRef{路 2:46}但当祂询问长老们关于法律的问题时，祂智慧的提问所传达的信息多于所寻求的。\par

4. 但或许我们应当称\Person{伯多禄}和\Person{若望}为无知的人，他们二人都能自称：\BibleQuote{我虽然在言语上粗鲁，但在知识上却不粗鲁}。\BibleRef{格后 11:6}难道\Person{若望}只是一个粗鲁未受教育的渔夫吗？若是如此，他从哪里得到这些话：\BibleQuote{在起初已有圣言，圣言与天主同在，圣言就是天主}。\BibleRef{若 1:1}希腊文中的\OriginalTerm{圣言}{Logos}有多种含义。它意指言语、理性、计算以及个体事物的原因，事物因它而存在。所有这些我们都正确地归于基督。这个真理，所有学问的\Person{柏拉图}并不知晓，所有雄辩的\Person{德摩斯梯尼}也不明白。\BibleQuote{我要摧毁智者智慧，废除贤者聪明}，经上说。\BibleRef{格前 1:19}真正的智慧必须摧毁虚假的智慧，尽管宣讲的愚妄\BibleRef{格前 1:21}与十字架不可分离，\Person{保禄}却也讲\BibleQuote{在成全者中的智慧，但不是这世代的智慧，也不是这世代将要消灭的掌权者的智慧}，而是讲\BibleQuote{天主奥妙的智慧，即隐藏的智慧，是天主在万世之前所预定}。\BibleRef{格前 2:6}天主的智慧就是基督，因为基督，正如我们被告知的，是\BibleQuote{天主的德能和天主的智慧}。\BibleRef{格前 1:24}祂就是那隐藏在奥秘中的智慧，我们在第九篇圣咏的标题中也读到：\BibleQuote{为圣子的隐秘之事}。在祂内蕴藏着一切智慧和知识的宝库。那隐藏在奥秘中的同一位就是在万世之前所预定的。现在，正是在法律和先知中，祂被预定和预表。为此原因，先知被称为先见者，\BibleRef{撒上 9:9}因为他们看见了别人所未曾看见的祂。\Person{亚巴郎}看见了祂的日子，就欢喜快乐。\BibleRef{若 8:56}向叛逆的百姓封闭的天，向\Person{厄则克耳}敞开了。\Person{达味}说：\BibleQuote{求你开启我的眼睛，使我得以看见你法律中的奇妙}。因为\BibleQuote{法律是属神的}，\BibleRef{罗 7:14}需要启示才能使我们理解它，并且当天主揭开祂的面容时，得以看见祂的荣耀。\par

5. 在\WorkTitle{默示录}中，有一本书卷展开，用七封印封着，\BibleRef{默 5:1}你若把它交给一个学者，说：\Quote{读这卷书}，他必回答你：\BibleQuote{我不能，因为它是封着的。}\BibleRef{依 29:11}今天，有多少人自以为有学问，但圣经对他们来说是一本封着的书，除非藉着那拥有达味钥匙的一位，他们不能打开它，\BibleQuote{他开了，没有人能关；他关了，没有人能开。}\BibleRef{默 3:7}在\WorkTitle{宗徒大事录}中，那位神圣的太监（或者不如说是\BibleQuote{人}，因为圣经如此称呼他\BibleRef{宗 8:27}），当他读依撒意亚先知书时，斐理伯问他：\BibleQuote{你明白你所读的吗？}他回答说：\BibleQuote{若没有人指教我，怎能明白呢？}\BibleRef{宗 8:30}现在暂且转向我自己，我既不比这位太监更圣洁，也不比他更勤恳。他从厄提约丕雅来，也就是从地极来到圣殿，撇下了王后的宫殿；他是法律和神圣知识的热爱者，即使在战车上也阅读圣经。然而，尽管他手中有书卷，心中记住主的话，甚至唇舌念出这些话，他仍然不认识那在书中他虽不认识却敬拜的那一位。于是斐理伯来了，向他展示了隐藏在文字下的耶稣。真是奇妙的导师！顷刻间，太监相信了并接受了圣洗；他成了信者和圣人。他不再是个学生，而是个老师；他在旷野中教会的水泉里找到了比在会堂的金碧辉煌的圣殿中更丰盈的恩宠。\par

6. 我仅仅略举了这些例子（由于信函的篇幅限制，不能更详尽地论述），以使你深信，在圣经中，除非有一位向导为你指明道路，否则你无法进步。我暂且不提文法家、修辞学家、哲学家、几何学家、逻辑学家、音乐家、天文学家、占星家、医生的知识，这些技艺对人类最为有用，可归纳为教导、方法和精通三个方面。我将转向不那么重要的技艺，这些技艺更需要手的灵巧而非脑力。农夫、泥瓦匠、木匠、木工和金工、羊毛整理工和漂洗工，以及那些制造家具和廉价器皿的工匠，若没有合格人士的指导，便无法达到他们追求的目标。正如\Person{贺拉斯}所说：\par

\Quote{惟独医生自称精于治疗之术，\newline{}除了木匠，无人试图接合木材。}\par

7. 诠释圣经的艺术是唯一一门普天之下人人都自称精通的艺术。再次引用\Person{贺拉斯}：\par

\Quote{无论受过教导或未受过教导，我们都在写诗。}\par

多嘴的老妇人、昏聩的老翁、花言巧语的诡辩家，每一个人都拿起圣经，将它们撕成碎片，在还未学习之前就教导别人。有些人皱着眉头，夸夸其谈，把互相矛盾的话拼凑起来，在软弱的妇女中间对圣著作哲学讨论。其他人——我惭愧地说——从妇女那里学习他们本应教导男人的东西；仿佛这还不够，他们大胆地向别人解释他们自己毫不理解的东西。我不提那些像我一样，在开始学习圣经之前就熟悉世俗文学的人。这些人用文辞的优美取悦大众的耳朵，就以为他们所说的每一句话都是天主的法律。他们不屑于注意到先知和宗徒的本意，却将矛盾的经文调整得符合自己的意思，仿佛这是一种伟大的教学方式——实则是最糟糕的——歪曲作者的观点，强迫圣经勉强服从他们的意愿。他们忘了我们读过\Person{荷马}和\Person{维吉尔}的集句诗；但我们绝不会因为维吉尔（没有基督的那位）有这样的诗句就称他为基督徒：——\par

\Quote{现在贞女归来，萨图尔努斯统治再现，\newline{}现在从高天降生了一个新生婴儿。}\par

另一行诗句可能是圣父对圣子说的：——\par

\Quote{万岁，独生子，我的能力和威严。}\par

还有另一行可能追随救主在十字架上说的话：——\par

\Quote{他说了这样的话，就钉在那里不动了。}\par

但这一切都是幼稚的，如同江湖骗子的戏法。试图教导你所不知道的东西是徒劳的，并且——如果我可以稍带感情地说——更糟糕的是对自己的无知浑然不知。\par

8. 有人会对我们说，\WorkTitle{创世纪}无需解释；它的主题太简单——世界的诞生、人类的起源、大地的分裂、口音的混乱，以及希伯来人下到埃及！\WorkTitle{出谷纪}无疑同样浅显，因为它仅仅包含十灾、十诫和一些神秘而神圣的诫命！\WorkTitle{肋未纪}的意义当然不言自明，尽管它所描述的每一项祭献，甚而至于它所包含的每一句话，对亚郎衣饰的描述，以及所有与肋未人相关的规条，都是天上之事的预像！\WorkTitle{户籍纪}也是如此——难道它的数字、巴郎的预言，以及旷野中的四十二个安营之地，不都是许多奥秘吗？\WorkTitle{申命纪}也是如此，即第二法律或福音法律的预像——它岂非在展示已知之事的同时，以新的光照亮古老的真理？这就是宗徒所夸耀的，他愿意在教会中说那\Quote{五句话}。\BibleRef{格前 14:19}再论约伯，那忍耐的典范，在他的论述中包含了何等的奥秘！该书以散文开始，很快转入诗歌，最后又回到散文。它通过提出命题、假设前提、引证证明、得出结论的方式，阐明了逻辑学的一切法则。书中出现的单个词语也充满意义。且不说其他主题，它预言人身体的复活，比任何人所展示的都更清晰、更谨慎。约伯说：\BibleQuote{我知道我的救赎主活着，末了我要从大地兴起；我必再以我的皮囊包裹自己，我将在我的肉身内看见天主。我必亲眼看见，我的眼睛要看见，并非别人。这希望在我怀里存留。}我要继续论\Person{农的儿子若苏厄}——他在名称和行为上都是主的预像——他渡过约但河，征服敌国，为得胜的民族分地，并在他所处理的每一座城市、村庄、山岭、河流、山溪和边界上，刻划了天上耶路撒冷，即教会的属灵疆域。\BibleRef{迦 4:26}在\WorkTitle{民长纪}中，每一位民间领袖都是一个预像。摩阿布女子卢德实现了依撒意亚的预言：\BibleQuote{主啊，请从旷野的岩石，向熙雍女儿的山岭，遣发一只羔羊，作大地的统治者。}在厄里之死和撒乌耳被杀的形象中，撒慕尔显示了旧法律的废除。同样，在匝多克和达味身上，他见证了新司铎职和新王国的奥秘。希伯来文称为\OriginalTerm{玛拉基姆}{Malâchim}的\WorkTitle{列王纪}第三、四卷，记载了犹大王国从撒罗满到耶苛尼雅的历史，以及以色列王国从乃巴特之子雅洛贝罕到被掳往亚述的曷舍亚的历史。如果你只看叙事，文字很简单，但如果你透过表面看隐藏的意义，你会发现对教会人数稀少以及异端者向她发动的战争之描述。十二位小先知的作品压缩在一卷窄小的篇幅内，其预像意义远不同于字面意义。欧瑟亚多次提及厄弗辣因、撒玛黎雅、若瑟、依次勒耳、淫妇之妻和淫妇之子女\BibleRef{欧 1:2}，以及一个困在丈夫内室中的淫妇，长久地守寡并穿着丧服，等待丈夫回到她身边。贝特尔之子岳厄尔描述了十二支派之地如何被蝗虫、蝻子、蚂蚱和霉病蹂躏\BibleRef{岳 1:4}，并预言在先前的人民覆亡之后，圣神要浇灌在祂的仆人和婢女身上\BibleRef{岳 2:29}；这同一位圣神，即要在熙雍楼上浇灌在一百二十位信徒身上的那一位。这些信徒从一渐升至十五，构成阶梯，在\OriginalTerm{登阶圣咏}{psalms of degrees}中有神秘地暗示。亚毛斯，虽然不过是乡间的\Quote{牧人}和\Quote{采集桑葚的人}\BibleRef{亚 7:14}，无法用几句话解释。因为谁能充分讲述大马士革、迦萨、提洛、厄东、摩阿布、阿孟子民的三次和四次过犯，以及第七和第八位的犹大和以色列？他向撒玛黎雅山上的肥牛说话\BibleRef{亚 4:1}，见证大屋和小屋都要倒塌\BibleRef{亚 6:11}。他时而看见那造蝗虫的\BibleRef{亚 7:1}，时而看见主站在一座涂抹或由金钢石制成的墙上\BibleRef{亚 7:7}，时而看见一篮将要带来刑罚的熟果子\BibleRef{亚 8:1}，时而看见地上的饥荒，\BibleQuote{不是食物的饥荒，也不是渴水的饥渴，而是听不到主的话}\BibleRef{亚 8:11}。亚北底亚，其名意为天主的仆人，向身染红血的厄东和土生之物发出雷霆；因它与兄弟雅各伯不断争竞，他以圣神之矛击打它。约纳，最美丽的鸽子，他的船难预表主的受难，唤回世界悔改，当他向尼尼微宣讲时，向所有外邦人宣告救恩。摩拉舍提人米该亚，与基督同为承继者，宣告抢夺者女儿的毁灭，并围攻她，因为她击打了以色列判官的腮骨。纳鸿，世界的安慰者，斥责\Quote{血城}，当它倾覆时喊道：\BibleQuote{看啊，在山上那传报佳音者的脚步，多么美丽！}\BibleRef{鸿 1:15}哈巴谷，如同坚固不屈的角力者，站在守望台上，将脚踏在堡垒上\BibleRef{哈 2:1}，以便默观被钉的基督，说：\BibleQuote{祂的光荣遮盖了诸天，大地充满了祂的赞美。祂的光明如同日光，从祂手中射出双角；在那里隐藏着祂的力量。}\BibleRef{哈 3:3}索福尼亚，即主的近卫和奥秘知悉者，听见\BibleQuote{从鱼门有哀号，从第二门有嚎啕，从山丘有巨大的崩裂声}\BibleRef{索 1:10}。他宣告\BibleQuote{向居住在锅坑的居民哀号，因为迦南所有的子民都灭亡了，所有满载银子的人都被剪除了。}哈盖，即欢喜快乐的人，他含泪播种以便欢乐收割，专注于重建圣殿。他代表主（即圣父）说：\BibleQuote{不久，再片刻，我将震动诸天、大地、海洋和旱地；我必震动万民，万民所渴望的必然来到。}\BibleRef{盖 2:6}匝加利亚，记念其主的人，给了我们许多预言。他看见若苏厄\BibleQuote{穿着污秽的衣服}\BibleRef{匝 3:3}，一块有七只眼睛的石头\BibleRef{匝 3:9}，一座纯金的灯台和七盏灯，以及灯盏左右两边的两棵橄榄树\BibleRef{匝 4:2}。在描述马匹（红、黑、白和有斑点的）之后\BibleRef{匝 6:1}，以及从厄弗辣因斩断战车、从耶路撒冷斩断战马之后\BibleRef{匝 9:10}，他预言一位贫穷的君王，他要\BibleQuote{骑着一匹驴驹}\BibleRef{匝 9:9}。玛拉基亚，所有先知中最后一位，公开谈论以色列的弃绝和外邦人的蒙召。\BibleQuote{我不喜欢你们，万军的上主说，我也不接受你们手中的祭品。因为从日出到日落，我的名在外邦中伟大，在各处都向我的名献上馨香和纯洁的祭品。}至于依撒意亚、耶肋米亚、厄则克耳和达尼尔，谁能完全理解或充分解释他们？他们中的第一位似乎不是在撰写预言，而是在撰写福音。第二位谈到一根杏树枝\BibleRef{耶 1:11}和一个朝北的沸腾锅\BibleRef{耶 1:13}，以及一只改变斑点的豹子\BibleRef{耶 13:23}。他还用不同的韵律写了四次字母诗。四人中第三位厄则克耳的开头和结尾极其晦涩，如同\WorkTitle{创世纪}的开头，希伯来人在三十岁以前都不学习。四位先知中最后一位达尼尔，通晓时节，关心全世界，用清晰的语言宣告那块从山上非人手凿出的石头，它要推翻所有王国\BibleRef{达 2:45}。达味，他集我们的\Person{西蒙尼得}、\Person{品达}、\Person{阿尔卡俄斯}、\Person{贺拉斯}、\Person{卡图卢斯}和\Person{塞雷努斯}于一身，用竖琴歌唱基督；用十弦瑟呼唤他从地下界复活。撒罗满，爱慕和平及主的人，纠正道德，教导本性，联合基督与教会，并唱一首甜美的婚姻歌曲来庆祝那神圣的婚约。艾斯德尔，教会的预像，拯救她的子民脱离危险，在杀死名意为邪恶的哈曼之后，为后代留下一个值得纪念的日子和一个盛大节日。\WorkTitle{编年纪}，即旧约中被遗漏之事或概略，是如此重要和有价值，以至于没有它，任何自诩通晓圣经的人都将使自己在自己眼中成为笑柄。其中使用的每一个名字，甚至词语的连接，都用来阐明\WorkTitle{列王纪}中略过的事件以及福音所提出的问题。厄斯德拉和乃赫米雅，即主的助佑者和安慰者，合为一卷。他们重建圣殿，修筑城墙。在他们的篇章中，我们看到以色列人众回归故土，我们读到司铎、肋未人、以色列人和归化者的名单，我们甚至得知负责建造城墙和塔楼的不同家族。这些记述表面是一种意义，但内里却是另一种意义。\par

9. [In Migne, 8.] 你看，我被对圣经的热爱所冲昏，已经超越了一封信的界限，却仍未完全达到目的。我们只听到了我们应当知道和渴望之事，以便我们也能与圣咏作者一起说：\Quote{我的灵魂因极其热切地渴望你的典章而迸发。}但苏格拉底关于自己的那句话——\Quote{我只知道我一无所知}——也在我们身上应验了。关于新约，我将简要处理。玛窦、马尔谷、路加和若望是主的四活物，是真正的革鲁宾，即知识之库。在他们身上，整个身体布满眼睛，他们闪烁如火花，\BibleRef{则 1:7}他们疾行回转如闪电，\BibleRef{则 1:14}他们的脚是直脚，\BibleRef{则 1:7}且向上举起，他们的背也有翅膀，准备向四方飞翔。他们彼此相连，相互交织：\BibleRef{则 1:11}像轮中套轮滚动前进，\BibleRef{则 1:16}随圣神的气息所向而行。\BibleRef{则 1:20}宗徒保禄写信给七个教会（第八封书信——致希伯来人书——通常未被计入）。他训导弟茂德和弟铎；他为费肋孟的逃跑奴隶求情。关于他，我认为与其写得不当，不如闭口不言。\WorkTitle{宗徒大事录}似乎只是朴实无华地叙述了新诞教会的幼年；但一旦我们意识到其作者是福音中所称赞的路加医生，就会发现他所有的话语都是患病灵魂的良药。宗徒雅各伯、伯多禄、若望和犹达发布了七封书信，既属灵又切题，有短有长，言辞虽短而内容深远，以至于读它们的人几乎没有不感到茫然的。\WorkTitle{若望默示录}有多少词语，就有多少奥秘。我这样说还不及该书所应得的。一切对它的赞美都不足够；每一个词中都隐藏着多重含义。\par

10. [In Migne, 9.] 我恳求你，我亲爱的弟兄，以这些书为伴，默想它们，不知别的，不寻别的。这样的生活对你来说岂非天堂在世上的预尝？不要因圣经的简朴或其词汇的贫乏而反感；这或是翻译者的错误所致，或是故意为之：因为如此，它更适合教导不识字的会众，使有学问的人和无学问的人能从同一句话中各得其所。我不至于迟钝或冒昧到自称明白它，或能在地上采摘根在天上的果实，但我承认我渴望如此。我比那闲坐的人更进一步，我虽不敢自称老师，却乐意自荐为同窗。\Quote{凡求的，就必得到；找的，就必找到；敲门的，就必给他开门。}\BibleRef{玛 7:8}让我们在地上学习那将随我们进入天堂的知识。\par

11. [In Migne, 10.] 我将张开双臂欢迎你，并且——如果我可以像赫尔马戈拉斯那样自夸和说傻话——我会力图与你一同学习任何你渴望研究的内容。在此地的欧瑟伯以兄弟之情待你；他告诉我你品格的真诚、你对世界的轻蔑、你在友谊中的坚定和你对基督的爱，使你的信加倍珍贵。这封信本身就（无需他帮忙）显露出你的谨慎和文采的魅力。那么，我恳求你，赶快砍断而不是解开那阻止你船只在海上航行的缆绳。因轻贱自己的财物而变卖它们并意图弃绝世界的人，不会想要高价出售。你必须花费在行装上的钱可视为赚得的。有一句古话说：吝啬鬼既缺乏他已拥有的，也缺乏他没有的。信徒拥有全世界的财富；非信徒连一文钱也没有。让我们常常活\Quote{似乎一无所有，却拥有一切。}\BibleRef{格后 6:10}食物和衣服，这些是基督徒的财富。\BibleRef{弟前 6:8}如果你的财产在你的权下，\BibleRef{宗 5:4}就变卖它；若不然，就把它丢弃。\Quote{若有人拿走你的外衣，连斗篷也让他拿去。}\BibleRef{玛 5:40}你一味拖延，希望推迟行动：除非——如你争辩——除非我零星谨慎地变卖货物，否则基督将无法喂养他的穷人。不，凡把自己奉献给天主的人，已一次而永远地将一切献给了他。宗徒们只不过舍弃了船和网。\BibleRef{玛 4:18}那寡妇只向银库投了两个铜钱，\BibleRef{谷 12:41}然而她将胜过拥有全部财富的克罗伊斯。那常常默想自己必死的人，就能轻易轻视一切。\par

\SourceRecord{Translated by W.H. Fremantle, G. Lewis and W.G. Martley.；Nicene and Post-Nicene Fathers, Second Series；Vol. 6.；Edited by Philip Schaff and Henry Wace.；Buffalo, NY: Christian Literature Publishing Co.,；1893.；<http://www.newadvent.org/fathers/3001053.htm>.}

% structure: p0001 source=h1 target=section reason=page-title

\section{第五十四封信}

\Emph{致\Person{富丽雅}}\par

\EditorialNote{一封指导寡妇如何最好地保持寡居的信（根据热罗尼莫的说法，是\Quote{三种贞洁中的第二种}）。富丽雅曾一度想再婚，但最终放弃了这一意图，专心照顾她年幼的孩子和年迈的父亲。热罗尼莫生动地描绘了她在罗马所面临的危险，为她制定了行为准则，并将她托付给司铎艾克苏佩里（后来的图卢兹主教）照料。这封信写于公元394年。}\par

你恳求并哀求我在信中写信给你——或者说是回信给你——你应当采取什么样的生活方式来保持寡妇的冠冕并保持你贞洁的名声无玷。我的心灵欢喜，我的肾腑雀跃，我的心快乐，因为你渴望像你神圣记忆中的母亲蒂蒂亚娜那样，在婚后长久地保持贞洁。她的祈祷和恳求蒙应允。她成功地在她唯一的女儿身上重新获得了她生前所拥有的。你们家族有一个崇高的特权：从\Person{卡米路斯}时代起，你们家很少有人或没有人被描述为缔结第二次婚姻。因此，如果你继续保持寡妇的身份，这不会像你作为基督徒，未能守住外邦妇女许多世纪以来所嫉妒地守护的东西那样，给你带来更多的赞美，而是耻辱。\par

我对\Person{保拉}和\Person{欧斯托基}——你们家族最美丽的花朵——只字不提；因为我的目的是劝勉你，我不愿显得是在赞扬她们。\Person{布拉西拉}我也略过不提，她紧随她的丈夫——你的兄弟——进入坟墓，在短暂的生命中完成了长期的德行。\BibleRef{智 4:13} 但愿男人们会效仿女人们可赞美的榜样，而皱纹满面的老年最终会付出青春一开始就乐意付出的代价！这样说，我是明知故犯地将手伸进火里。男人们将皱起眉头，向我挥舞紧握的拳头；\par

愤怒的\Person{克瑞墨斯}将用洪亮的声音咆哮。\par

领袖们将团结起来反对我的信；贵族们将对我发出雷鸣般的吼叫。他们会喊叫说我是个巫师和诱惑者，应当被流放到天涯海角。如果他们愿意，还可以加上撒玛黎雅人的绰号；因为其中我只认出这是对我主的一个称呼。但有一件事是确定的：我不将女儿与母亲分离，我不用福音的话：\BibleQuote{任凭死人埋葬他们的死人}。\BibleRef{玛 8:22} 因为凡信基督的人都活着；而且凡信祂的，\BibleQuote{他自己也应当照样行事，如同祂行事一样}。\BibleRef{若一 2:6}\par

暂且停止那些诽谤者持续加诸于所有自称为基督徒的人身上的恶意中伤，以免他们因羞耻的恐惧而不敢追求德行。除了书信往来，我们彼此并不相识；既然没有肉性上的相识，就不可能有不基于宗教动机的交往。\BibleQuote{尊敬你的父亲，}诫命说，\BibleRef{出 20:12}但前提是他没有使你与你的真父亲分离。承认血缘关系，但仅限于你的父母认识他的造物主。如果他未能做到，\Person{达味}将向你歌唱：\Quote{女儿，请听，请看，请侧耳；忘记你的民族和你的父家。君王就会恋慕你的美貌，因为他是你的主。}伟大的奖品是为忘记父母而设的：\Quote{君王将恋慕你的美貌。}你已经听见，已经看见，已经侧耳，忘记了你的民族和你的父家；因此君王将恋慕你的美貌，并对你说：\BibleQuote{我的爱，你全然美丽，毫无瑕疵}。\BibleRef{歌 4:7} 有什么比一个被称为天主的女儿、不为自己寻求外在装饰\BibleRef{伯前 3:3}的灵魂更美丽呢？她信基督，因着这伟大希望的嫁妆，走向她的新郎；因为基督同时是她的新郎和她的主。\par

婚姻涉及多少麻烦，你已在婚姻状态本身中学到了；你已饱食鹌鹑肉直至厌烦；你的嘴里充满了苦胆；你排出了不消化和不健康的食物；你缓解了翻腾的胃。为什么你要再次吞下对你不适的东西？\BibleQuote{狗回到自己的呕吐物，洗净的猪回到泥潭中打滚}。\BibleRef{伯前 2:22} 即使是野兽和飞鸟也不会两次落入同一个陷阱。你担心如果\Person{卡米路斯}的血脉灭绝，你不给你父亲提供一个可以爬在他胸前并流口水在他脖子上的小家伙？好像所有结婚的人都有孩子！好像即使有孩子，他们总是像他们的祖先！\Person{西塞罗}的儿子展示了他父亲的雄辩吗？你自己的\Person{科尔内利娅}，贞洁和多产的典范，有理由为她是她的\Person{格拉古兄弟}的母亲而喜乐吗？指望确定的后代是荒谬的，因为你可以看到很多人没有孩子，而另一些人有了又失去了。那么你将把你的巨大财富留给谁？留给不会死的基督。你立谁为你的继承人？那位已经是你的主的。你的父亲会难过，但基督会喜乐；你的家人会悲伤，但天使会与你一同喜乐。让你父亲随他所愿处理他自己的东西。你不属于你所出生的那个他，而属于你借以重生、并用他自己的宝血以巨大代价赎买了你的那一位。\BibleRef{宗 20:28}\par

小心乳母、侍女和类似的毒物，她们试图通过吸你的血来满足她们的贪婪。她们建议你做的不是对你最好的，而是对他们最好的。她们不断地在你耳边重复\Person{维吉尔}的诗句：——\par

你将把整个青春浪费在孤独的悲哀中，\newline{}并放弃可爱的孩子——爱的礼物吗？\par

凡有神圣贞洁之处，也有简朴的生活；凡有简朴的生活，仆人们便受损。他们得不到的，在他们看来是被剥夺了。他们只看实际收到的金额，而不看相对的比例。他们一见到基督徒，立即重复那句陈腔滥调：\Quote{希腊人！骗子！}他们散播最可耻的谣言，而当这样的谣言出自他们自己时，他们假装是从别人那里听来的，这样既制造了故事又夸大了它。一个虚假的传闻传开了；当它传到已婚妇女那里，被她们的舌头煽动，便传遍各省。你可以看到许多这样的妇人——脸上涂脂抹粉，眼睛像毒蛇，牙齿用浮石磨光——疯狂地咆哮着、指责基督徒。其中一位妇人，\par

一件紫罗兰色的披肩披在她肩上，\newline{}她拖长声音说着一些矫揉造作的话，用鼻子说话，\newline{}用结结巴巴的舌头把一半的话说得很琐碎。\par

于是其余的人齐声附和，每条长凳都发出嘶哑的赞同声。他们得到我同会弟兄的支持，这些人发现自己受到攻击，便反过来攻击别人。他们总是善于攻击我，却在自己辩护时哑口无言；就好像他们自己不是隐修士，就好像每一句反对隐修士的话也不会危及他们的神性祖先——神职人员。对羊群的伤害给牧人带来耻辱。另一方面，我们不得不赞美一位隐修士的生活，他尊崇基督的司铎，拒绝诽谤他赖以成为基督徒的那个会秩。\par

6. 我这样说话，我的基督内的女儿，并非因为我怀疑你会忠于你的誓愿（若你对独身的福分不确定，你就绝不会请求一封劝告信），而是为了让你认识到仆人们仅仅想为了自己的利益出卖你的恶意，亲戚们可能为你设下的陷阱，以及一位父亲善意却错误的建议。我承认他对你怀有爱，但我不能承认这是按知识的爱。我必须与宗徒一同说：\BibleQuote{我可以为他们作证，他们对天主有热诚，但不是按着知识。}\BibleRef{罗 10:2} 你当效法——我不能说得太多——你那位圣洁的母亲，每逢我记起她时，她对基督的热诚便涌上我心头；不仅是她的热诚，还有因禁食而显得苍白的面容、她给穷人的施舍、她对天主仆人的礼貌、她的衣着和心灵的谦卑，以及她言语中一贯的节制。至于你的父亲，我带着敬意说话，并非因为他是贵族并担任执政官，而是因为他是基督徒。愿他忠于自己的信仰。愿他因生了一个为基督而非为世界的女儿而喜乐。反而，他应当悲伤，因为你白白失去了童贞，因为婚姻的果实不属于你。他给了你的丈夫在哪里？即使他可爱而善良，死亡仍会夺走一切，他的去世结束了你们之间的血肉纽带。我恳求你，抓住机会，化必要为德行。在基督徒的生命中，我们不看开端，而看结局。保禄开始得不好，但结局很好。犹大的开端值得称赞，但他的结局因背信弃义而受谴责。读厄则克耳：\BibleQuote{义人的正义在他犯罪之日不能救他；恶人的邪恶在他转离邪恶之日也不致使他跌倒。}\BibleRef{则 33:12} 基督徒的生命是真正的雅各伯的梯子，天使在上面上去下来，\BibleRef{创 28:12} 而上主站在梯子顶端，向那些滑倒的人伸出手，并藉着对祂的默观扶持那些攀登者疲惫的脚步。但是，祂不愿恶人丧亡，却愿他归化并生活，祂却厌恶那些温吞的人，\BibleRef{默 3:16} 他们很快就令祂厌烦。那得赦免多的，爱得也多。\BibleRef{路 7:47}\par

7. 在福音中，一个妓女获得了救恩。如何？她以泪水受洗，并用那先前欺骗了许多人的头发擦干主的脚。她不戴飘动的头巾或吱嘎作响的靴子，她不用锑黑化眼睛。然而在她污秽中，她比任何时候都更可爱。在一位基督徒妇女的脸上，胭脂和白铅有什么地位？前者模拟脸颊和嘴唇的自然红润；后者模拟脸和脖子的白净。它们只用于点燃青年的情欲，刺激淫欲，并表明不贞洁的心。一个女人怎能为自己的罪哭泣呢？她的泪水暴露了她的真实肤色，并在脸颊上留下沟痕。这样的装饰不是来自主的；这种面具属于假基督。一个女人怎能自信地将面孔举向天堂，而她的造物主必定认不出它来？用少女和青春虚荣来为这类行为辩解是徒劳的。一位寡妇，既已不再有丈夫取悦，按宗徒的话说是一位真正的寡妇，\BibleRef{弟前 5:5} 除了坚忍外什么也不需要。她记得过去的享乐，她知道什么曾给她快乐，她如今失去了什么。她必须以严格的禁食和守夜来熄灭魔鬼燃烧的箭。\BibleRef{弗 6:16} 如果我们守寡，我们必须要么穿得如我们所说的，要么说得如我们穿着的。为什么我们所说的是一回事，所做的又是另一回事？舌头谈论贞洁，但身体的其他部分却显示不节制。\par

8. 关于服饰和装饰就说到这里。但一个寡妇\BibleQuote{生活奢华}——这不是我的话，而是宗徒的话——\BibleQuote{虽生犹死}。\BibleRef{弟前 5:6} 这是什么意思——\BibleQuote{虽生犹死}？在不明真相的人看来，她似乎活着，而不像她实际那样死于罪恶；是的，从另一种意义上说，她对基督是死的，因为基督无所不知。\BibleQuote{犯罪的人，他必死亡。}\BibleRef{则 18:20}\BibleQuote{有些人的罪过是明显的，即先受审判；有些人的罪过是随后跟随的。同样，善工也是明显的，即使不明显，也不能隐藏。}\BibleRef{弟前 5:24} 这些话的意思是：某些人犯罪如此故意和明目张胆，以至于你一看见他们就知道他们是罪人。但另一些人的缺陷被狡猾地掩盖，我们只能通过后来的信息才得知。同样，有些人的善行是众所周知的，而另一些人的善行只有通过长期交往才能了解。那么，我们何必夸耀自己的贞洁呢？贞洁若没有它的同伴和随从——节欲和简朴生活——就无法证明其真实性。宗徒\BibleQuote{磨练自己的身体，使身体服从灵魂，免得我给别人报捷，自己反而落选}；\BibleRef{格前 9:27} 而一个仅仅因食物丰盛而情欲炽烈的少女，她能自信于自己的贞洁吗？\par

9. 我说这些，当然不是谴责天主所造、应感恩享用的食物，\BibleRef{弟前 4:4} 而是为了消除青少年中引发感官享乐的诱因。无论是燃烧的埃特纳火山、武尔坎的领地、维苏威火山，还是奥林匹斯山，都不及那些因酒足饭饱而血气方刚的青年的骨髓燃烧得猛烈。许多人将贪婪踩在脚下，如同放下钱袋一般轻易。强制性的沉默可以为咒骂的舌头做补赎。外貌和着装方式可以在一小时内改变。所有其他罪恶都是外在的，外在的容易抛弃。唯独欲望是天主植入人内、为引导人生育子女的，它是内在的；一旦越过界限，便成为罪恶，并因自然法则而呼喊性交。因此，这是一项伟大的功绩，需要不懈的努力去克服你内在的本性：活在肉身中却不随从肉身，日复一日地与己争斗，以传说中的百眼阿耳戈斯那样的一百只眼睛注视你内中关闭的仇敌。这就是宗徒用另一句话所说的：\BibleQuote{凡人所犯的一切罪，都在身体以外；唯独犯邪淫的，是得罪自己的身体。}\BibleRef{格前 6:18} 医生和那些撰写人体本性著作的人，尤其是加仑在他的\WorkTitle{论健康}一书中说，男孩、青年、壮年男女的身体因内在的热力而炽热，因此对这些年龄的人而言，所有助长这种热力的食物都有害；相反，在饮食中摄取冷的、清凉的东西则极有利于健康。相反，他们告诉我们，温热的食物和陈酒对因体液和寒冷而受苦的老人有益。因此，救主说：\BibleQuote{你们要谨慎，免得你们的心为宴饮沉醉及人生的挂虑所压住。}\BibleRef{路 21:34} 宗徒也说：\BibleQuote{不要醉酒，醉酒使人放荡。}\BibleRef{弗 5:18} 难怪陶匠这样谈论他所造的器皿，连那位只求了解和描述人之习性的喜剧诗人也告诉我们：\par

谷神和酒神不在，爱神便低垂。\par

10. 那么，首先，在你度过青春年少之前，只饮清水，因为清水天然是最具凉爽性的饮料。如果你的胃足够强壮能够承受的话，就这样做；但如果你的消化力弱，请听宗徒对弟茂德说的话：\BibleQuote{因为你的胃口和屡次患病，稍微用点酒。}\BibleRef{弟前 5:23} 至于你的食物，必须避免所有燥热的菜肴。我不只是指肉食（虽然关于肉食，蒙拣选的器皿这样表明心迹：\BibleQuote{最好是不吃肉，不喝酒}\BibleRef{罗 14:21}），也包括蔬菜。一切刺激性的或难以消化的食物都要拒绝。请相信，对于年轻的基督徒来说，没有比吃蔬菜更好的了。因此他在另一处说：\BibleQuote{那软弱的人，只吃蔬菜。}\BibleRef{罗 14:2} 这样，身体的热力必须用凉性食物来调节。达尼尔和三个孩子以豆类为生。\BibleRef{达 1:16} 他们还是少年，尚未进入巴比伦王油炸长老们的那口锅。而且，由于天主特别的恩宠，他们的身体状况因饮食而改善。我们并不期望我们也会如此，但我们期待灵魂的力量增强，因为肉躯越弱，灵魂越强。有些渴望贞洁生活的人，以为只需戒绝肉食，便在途中半途而废。他们以蔬菜填满胃，而蔬菜只有少量、适度食用时才无害。如果要我说出我的想法，没有什么比未消化的食物和打嗝不畅的呼吸更能使身体发热、点燃情欲的了。至于你，我的女儿，我宁愿伤害你的谦逊，也不愿因轻描淡写而危及我的论点。凡带有感官享乐种子的一切，都视为毒药。清淡的饮食，让食欲永远不满足，胜过三天的禁食。每天吃一点，比有些日子完全禁食、另一些日子暴饮暴食要好得多。那种慢慢降入土地的雨水最好。突然猛烈倾泻的阵雨会冲走土壤。\par

11. 用餐时，要想到之后你立即得祈祷和诵读。定下圣经的固定行数，当作你向主所尽的义务。决不可让自己入睡，除非你用这纺织的纬线充满你胸中的篮子。在圣经之后，你应阅读博学之士的著作；至少是那些信德广为人知者的著作。你无需到泥沼中寻金；你已有许多珍珠，用这些买那颗珍珠\BibleRef{玛 13:45}。要像耶肋米亚所说的，站在多条路上，这样你才能找到通向圣父的正路。将你对项链、宝石和丝绸衣服的喜爱，换成研读圣经的热忱。进入流奶流蜜的应许之地\BibleRef{出 33:3}。吃细面和油。让你的衣服，像若瑟的一样，是多种颜色的\BibleRef{创 37:23}。让你的耳朵，如同耶路撒冷的耳朵一样\BibleRef{则 16:12}，被天主的话穿透，让新谷的珍贵颗粒挂在它们上面。在可敬的\Person{埃克苏佩里乌斯}身上，你有一位多年考验、信德坚定的人，随时准备用他的忠告给你持续的支持。\par

12. 要用不义的钱财结交朋友，为在你们被接进永恒的帐幕时，他们能迎接你们\BibleRef{路 16:9}。施舍你的财富，不要给那些吃野鸡的人，而要给那些只有普通面包可吃、只解饥而不引欲的人。顾念贫穷和匮乏的人。凡求你的，就给他\BibleRef{玛 5:42}，但尤其应行善给有同样信德的人\BibleRef{迦 6:10}。赤身露体的，要给他衣服穿；饥饿的，要给他食物吃；患病的，要去看顾他\BibleRef{玛 25:35}。每次你伸出手时，要想到基督。注意，当你主上主向你讨求施舍时，不要增加不属于祂的财富。\par

13. 要避开年轻人的交往。绝不要让那些长发、装扮华丽、轻浮的年轻人出现在你的屋檐下。躲避歌手如同躲避祸患。将那些靠弹唱谋生的女人赶出你的家，她们是魔鬼的唱诗班，其歌声如塞壬一般致命。不要自以为是寡妇而放纵，在一群太监的簇拥下招摇过市。对于因性别和年龄而软弱的人来说，滥用己见、以为悦乐的事就是合法的，这是极为有害的。\BibleQuote{凡事都可行，但不全有益处}\BibleRef{格前 6:12}。没有一个卷发的管家、美丽的乳兄弟或红润的脚夫可以跟在你脚后。有时，女主人的性情可从女仆的衣着推断出来。要寻求圣洁的童贞女和寡妇的陪伴；若有必要与男子交谈，不要避开见证人，且你的言谈应充满自信，使第三者的进入既不令你吃惊，也不使你脸红。面孔是心灵的镜子，女人的眼睛无需言语就透露出心中的秘密。最近我见到一件极其可悲的丑闻传遍整个东方。那女子的年龄、风度、衣着、表情、所交往的杂众、精美的餐桌和如王的排场，都表明她是尼禄或萨尔达纳帕鲁斯的新娘。他人的伤痕应使我们谨慎。\Quote{引起麻烦的人受鞭打时，愚昧者也会变聪明。}圣洁的爱不知急躁。虚浮的谣言很快被平息，后来的生活对以往做出判断。的确，没有人能在不受诽谤的情况下跑完人生的赛程；恶人以此为慰藉来挑剔良善，以为犯罪的人越多，罪本身的严重性就越小。然而，稻草之火很快熄灭，蔓延的火焰若缺乏燃料也会迅速熄灭。一年前盛传的传闻是真是假，一旦罪过停止，丑闻也会停止。我说这些，不是因为我担心你身上有什么不对，而是因为我深爱着你，任何安全我都不能不担忧。哦！愿你见到你的姊妹，愿你能听到她圣洁唇舌的雄辩，看到那令她弱小身躯充满活力的伟大精神。你会听到新旧约的全部内容从她心中涌出。禁食是她的运动，祈祷是她的消遣。如同米黎盎在法老淹死后拿起鼓，向童贞女唱诗班唱道：\BibleQuote{让我们歌颂上主，因他大获全胜；将马和骑士投在海中。}\BibleRef{出 15:21}她教导同伴成为乐女，却是为基督的乐女；成为琴师，却是为救主的琴师。她日夜以此忙碌，备好灯油，等候新郎的到来\BibleRef{玛 25:4}。因此，你当效法你的女亲属。愿罗马在你身上拥有比罗马更伟大的城市——我是指白冷——在她身上所拥有的。\par

14. 你富有，因此容易供给那些有需要的人食物。愿德行消耗那为放纵所预备的；一个决心蔑视婚姻的人，无需惧怕任何贫困。要拥有成群结队的童贞女，好带领她们进入君王的内室。支持寡妇，好使她们如同紫罗兰，与童贞女的百合花和殉道者的玫瑰花交织。这就是你应为基督编织的花环，以代替祂头戴的荆棘冠冕\BibleRef{玛 27:29}，祂在其中担负了世界的罪过。愿你最尊贵的父亲在你身上找到他的喜乐和支持；愿他从女儿那里学到从前从妻子那里学到的教训。他的头发已灰白，双膝颤抖，牙齿脱落，岁月在他的额头上刻下皱纹；死亡已近在门口，火葬堆几乎就在旁边准备好了。无论我们喜欢与否，我们都在变老。让他为自己准备长途旅行所需的干粮。让他带上那些否则他将不情愿留下的东西，不，让他事先送到天上去；若他拒绝，这些东西将被大地占有。\par

15. 年轻的寡妇中，有些人\Quote{已经转离跟随撒殚，当他们开始背叛基督而纵欲}并希望再婚时，通常提出这样的借口：\Quote{我那微薄的祖业日渐减少，我继承的财产被浪费，仆人羞辱我，使女忽视我的命令。谁将在官长面前为我出头？谁来负责我田产的地租？谁来照管我子女的教育和我奴隶的抚养？}说来可耻，她们竟将此当作再婚的理由，而这本应阻止她们再婚才对。因为再婚的母亲，为她的儿子们带来的不是保护者，而是仇敌；不是父亲，而是暴君。她被情欲点燃，忘记了腹中的果实，在毫不知情的孩子们中间，这位刚刚还在痛哭的寡妇又装扮成新娘。为何要拿你的财产和仆人的无礼作借口？承认那可耻的真相吧。没有女人是为了避免与丈夫同居而结婚的。至少，如果情欲不是你的动机，那么仅仅为了增加财富而行淫，简直是无稽之谈。你只不过是以珍贵的永恒恩宠为代价，换取微薄而短暂的收益！如果你已有子女，为何要再婚？如果你无子女，为何不惧怕重蹈先前的不孕？为何将不确定的收益置于确定的自尊损失之前？\par

今天为你订立婚约，但不久你就要被迫立遗嘱。你的丈夫会假装生病，并为你做他想要你为他做的事。然而，他一定会活下去，而你一定会死去。或者，如果你和第二任丈夫生了儿子，家庭争斗和内部争吵是必然的。你将不被允许爱你的头生孩子，也不被允许善待你亲生的子女。你将不得不秘密地给他们食物；即便如此，你现任丈夫也会对你前任心怀怨恨，除非你恨你的儿子们，否则他会认为你仍然爱他们的父亲。但你的丈夫可能已与前妻有子女。若是如此，当他带你回家时，即使你是世上最仁慈的人，所有修辞学家的陈词滥调、喜剧诗人和滑稽剧作家的台词都会朝你掷来，说你是残忍的继母。如果你的继子生病或头痛，你会被诬蔑为下毒者。如果你拒绝给他食物，你就是残忍；而如果你给他，你就会被认为用巫术迷惑了他。我问你，第二次婚姻能带来什么好处，足以补偿这些恶事？\par

16. 我们想知道寡妇应该怎样吗？让我们读一读\BibleBook{路加}{福音}。他说：\Quote{有一位先知\Person{亚纳}，}\Quote{是\Person{阿协尔}支派\Person{法奴耳}的女儿。}\BibleRef{路 2:36} \Person{亚纳}这个名字的意思是恩宠。\Person{法奴耳}在我们的话中是\Quote{天主的面}。\Person{阿协尔}可以翻译为\Quote{有福}或\Quote{财富}。她从年轻时直到八十四岁，一直肩负着寡妇的重担，不离圣殿，昼夜以禁食和祈祷事奉；因此她获得了属天主的恩宠，被称为\Quote{天主的面}的女儿，并分享了她祖先的\Quote{有福和财富}。让我们回想\Place{匝尔法特}的寡妇（\EditorialNote{\BibleBook{列王纪}{上}第十七章}），她更关心满足\Person{厄里亚}的饥饿，而不是保全自己和儿子的性命。尽管她相信如果当晚没有食物，她和儿子必死无疑，但她决定让她的客人存活。她宁愿牺牲自己的生命，也不愿忽视施舍的义务。在她的一把面里，她找到了种子，由此收获了主所赐的丰收。她撒下她的面，看！油从其中流出。在犹大地，谷物稀缺，因为麦子已经死了；\BibleRef{若 12:24} 但在一个外邦寡妇的家里，油却如泉涌。在\BibleBook{友弟德}{传}——如果有人接受它为正经的话——我们读到一位因禁食而憔悴、穿着丧服的寡妇，她外在的褴褛与其说是为亡夫悲伤，不如说是她期待新郎来临的心境。我看见她手执刀剑，沾满血迹。我认出了她从敌营带回来的\Person{敖罗斐乃}的头颅。在这里，一个女人战胜了男人，贞洁斩首了情欲。她迅速更换衣服，在胜利的时刻重新穿上她自己卑贱的衣裳，这衣裳比世间一切荣华更为华美。\EditorialNote{\BibleBook{友弟德}{传}第十三章}\par

17. 有些人因误解将\Person{德波辣}列入寡妇之中，并以为军队的首领\Person{巴辣克}是她的儿子，但圣经所言不同。我在此提到她，是因为她是一位女先知，且被列入民长之中，又因为她可能曾如圣咏作者所说：\Quote{你的言语在我的上颚多么甘甜！在我的口中比蜂蜜更甜。} 她被称为蜜蜂，确实恰当，因为她以圣经的花朵为食，充满圣神的馨香，并以先知之唇将花蜜的甘甜收集为一。然后是\Person{纳敖米}，希腊文作\OriginalTerm{受安慰者}{\GreekText{παρακεκλημένη}}，即受安慰者。当她的丈夫和子女客死异乡时，她带着贞洁返回家乡，并在旅途中以此支持自己，留住了她的\Place{摩阿布}儿媳，为要使\Person{依撒意亚}的预言应验：\Quote{上主，从旷野的磐石到\Place{熙雍}女儿的山，请派出羔羊来统治大地。} 我进而论及福音中的那位寡妇，她虽贫穷，却比所有\Place{以色列}人更富足。\BibleRef{谷 12:43} 她只有一粒芥菜子，却将酵母放在三斗面中；她将父与子的宣认与圣神的恩宠结合，将两个小钱投入银库。她所有的全部财产，她的一切所有，她都献在她信德的两个盟约中。这两位色辣芬以三重咏唱光荣天主圣三，并被存放在教会的宝库中。它们也是火钳的腿，藉此取来火炭洁净罪人的嘴唇。\BibleRef{依 6:6}\par

18. 但是，我何必从历史中回忆实例，从经书中列举圣善妇女的典范，何况在你自己的城市里，你眼前就有许多妇女，她们的榜样你大可以效法？我在此不谈她们的功绩，以免有奉承之嫌。只需提及圣\Person{玛尔塞拉}便够了：她忠于自己出身和地位的要求，为我们树立了合乎福音的生活。\Person{亚纳}\Quote{从她童贞时与丈夫同居了七年}；\BibleRef{路 2:36} \Person{玛尔塞拉}与丈夫同居了七个月。\Person{亚纳}期待基督的来临；\Person{玛尔塞拉}紧握着\Person{亚纳}曾抱在怀中的主。当基督仍是啼哭的婴孩时，\Person{亚纳}歌颂祂；如今基督已得胜，\Person{玛尔塞拉}宣扬祂的荣耀。\Person{亚纳}向所有期待\Place{以色列}得救赎的人谈论祂；\Person{玛尔塞拉}与得救赎的万民同声呼喊：\Quote{兄弟不能赎，但人将赎，} 又借用另一篇圣咏：\Quote{有一位生于她，至高者亲自建立了她。}\par

大约两年前，我清楚记得，我曾出版了一本反对\Person{约维尼安}的书，在其中凭借圣经的权威粉碎了对方就宗徒允许再婚一事所提出的异议。我无需在此重复我的论证，因为你可以在那篇论文中找到全部内容。为不超过一封书信的篇幅，我只想给你最后一条忠告：每天想着你必将死去，你就绝不会再想结婚。\par

\SourceRecord{Translated by W.H. Fremantle, G. Lewis and W.G. Martley.；Nicene and Post-Nicene Fathers, Second Series；Vol. 6.；Edited by Philip Schaff and Henry Wace.；Buffalo, NY: Christian Literature Publishing Co.,；1893.；<http://www.newadvent.org/fathers/3001054.htm>.}

% structure: p0001 source=h1 target=section reason=page-title

\section{书信五十五}

\Emph{\Person{阿曼德}}\par

\EditorialNote{一封非常有趣的书信。\Person{阿曼德}，\Place{波尔多}的一位司铎，曾写信给\Person{热罗尼莫}，请求解释三处圣经经文，即：\BibleRef{玛 6:34}、\BibleRef{格前 6:18}、\BibleRef{格前 15:25}、\BibleRef{格前 15:26}，并在同一封信中以一位\Quote{姐妹}（\Person{蒂埃里}认为可能是\Person{法比奥拉}）的名义提出了以下问题：\Quote{一个因丈夫的恶行而与其离婚的女子，在其丈夫尚在世时被迫再婚，她是否可以在不先做补赎的情况下与教会共融？}热罗尼莫在回信中给出了所要求的解释，但对关于那位\Quote{姐妹}的进一步问题，他给予了断然的否定回答。此信写于公元394年左右。}\par

1. 短函不容长篇解释；将大量内容压缩于狭小空间，只能对引发诸多思绪的主题略提数语。你问我\BibleBook{玛窦}{福音}中这段话的含义：\Quote{不要为明天忧虑，因为明天自有明天的忧虑；一天的难处一天当就够了。}\BibleRef{玛 6:34} 在圣经中，\Quote{明天}意指将来的时间。如\BibleBook{}{创世纪}中\Person{雅各伯}说：\Quote{我的义就在明日为我作证。}又如，当\Person{勒乌本}、\Person{加得}二支派及\Person{默纳协}半支派建筑了一座祭坛，全以色列派使者前往时，他们对大司祭\Person{丕乃哈斯}回答说，他们建筑祭坛是免得\Quote{明日}有人对他们的子孙说：\Quote{你们与主无分。}在旧约中你还能找到许多类似经文。因此，基督禁止我们为将来的事忧虑，却允许我们为当下的事思虑，因祂深知我们必死性情的软弱。祂接下来的话\Quote{一天的难处一天当就够了}应理解为：我们只需思虑今生当下的烦恼。何必将思绪延伸至未知之事，至那些我们或无法获得、或获得后不久必失去之物呢？希腊词\OriginalTerm{邪恶}{\GreekText{κακία}}在拉丁译本中译为\Quote{wickedness}，实际有两种截然不同的含义：邪恶与患难，后者在希腊语中称为\OriginalTerm{患难}{\GreekText{κακωσίν}}。在这段经文中，\Quote{患难}比\Quote{邪恶}是更好的译法。但若有人反对，坚持\OriginalTerm{邪恶}{\GreekText{κακία}}必须意为\Quote{wickedness}而非\Quote{患难}或\Quote{烦恼}，那么其含义应与\Quote{整个世界都卧在邪恶中}以及主祷文中的\Quote{救我们脱离凶恶}相同；这段话的用意便是，我们今生与这世界的邪恶争战，已经足够了。\par

2. 其次，你问我关于蒙福宗徒\Person{保禄}的\BibleBook{格林多}{前书}中这段话：\Quote{你们要逃避淫乱。人所犯的，无论什么罪，都在身体以外；惟有行淫的，是得罪自己的身体。}\BibleRef{格前 6:18} 让我们稍往前追溯，继续读直到这些话语，因为我们不应只从经文的后部，或可说是从章节的末尾，来寻求整个段落的意义。\Quote{身体不是为淫乱，乃是为主；主也是为身体。并且天主已经使主复活，也要用祂的大能使我们复活。岂不知你们的身子是基督的肢体吗？我可以将基督的肢体作为娼妓的肢体吗？绝对不可！岂不知与娼妓联合的，便是与她成为一体吗？因为祂说：二人要成为一体。但与主联合的，便是与主成为一灵。你们要逃避淫乱。人所犯的，无论什么罪，都在身体以外；惟有行淫的，是得罪自己的身体。}\BibleRef{格前 6:13} 等等。这位神圣的宗徒一直在论证节制，在此之前刚说过：\Quote{食物是为肚腹，肚腹是为食物；但天主必要叫这两样都废坏。}\BibleRef{格前 6:13} 现在他转而论及淫乱。因为饮食过度是情欲之母；饱食涨腹、酒醉浸透的肚腹必然导致感官的激情。正如别处所说：\Quote{人器官的排列暗示着他恶习的进程。}因此，诸如偷盗、杀人、抢劫、伪证等罪，在犯下之后都可以悔改；无论利益如何引诱，良心总会打击犯罪者。惟有淫欲和感官快乐，在悔改的那一刻又再次经历过去的诱惑、肉体的瘙痒和罪的引诱；以至于我们为纠正这些过犯所付出的思考，本身就成了新的罪恶之源。或者换个角度来看：其他罪恶都在我们之外；无论我们做什么，都是针对他人。但淫乱既玷污了行淫者的良心，也玷污了他的身体；并且按照主的话：\Quote{因此，人要离开父母，与妻子连合，二人成为一体。}行淫者同样与娼妓成为一体，藉着使那本是基督的殿成为娼妓的身体，而得罪了自己的身体。为不遗漏希腊注释者的任何解释，我再给你们一个说明。他们说，用身体犯罪和在身体内犯罪是两回事。偷盗、杀人及其他一切罪（淫乱除外），我们是用手在我们之外犯下的。惟有淫乱是我们在自身之内、在我们的身体内犯下的，而不是用身体针对他人。介词\Quote{用}表示犯罪所用的工具，而介词\Quote{在……内}表示情欲的领域是我们自身。另一些人则这样解释：按圣经，人的身体就是他的妻子；当人行淫时，他被说是得罪自己的身体，即得罪他的妻子，因为他藉着自己的淫乱玷污了她，并导致她虽自身无罪，却因与他结合而成了罪人。\par

3. 我发现你的询问书信附有一张短笺，其中包含以下话语：\Quote{问他（也就是我），一个女人因丈夫是奸夫和\OriginalTerm{鸡奸者}{sodomite}而离开了他，并发现自己被迫另嫁他人，那么在最初所离之夫在世时，她能否在不为自己过错做补赎的情况下与教会共融？}当我阅读所提出的案例时，我想起了一节经文：\Quote{他们为自己的罪找藉口。}我们都是人，都迁就自己的过错；凡我们自己的意志促使我们做的事，我们都归因于自然之必然。就好像一个年轻人会说：\Quote{我被身体所压倒，天性的炽热点燃了我的情欲，我身体的结构及其生殖器官要求性交。}或者一个杀人犯会说：\Quote{我缺衣少食，一无所有。我流别人的血，是为了救自己免于冻饿而死。}因此，告诉那位就这样向我询问她处境的姊妹，不是我的裁决，而是宗徒的裁决。\Quote{弟兄们（我向那些知道法律的人说话），你们岂不知道法律只管辖人活着的时期？因为有丈夫的女人，在法律上被束缚于她的丈夫，只要丈夫活着；但如果丈夫死了，她就从丈夫的法律中解脱了。所以，如果丈夫还活着，她嫁给了另一个男人，她将被称为奸妇。}\BibleRef{罗 7:1}在另一处：\Quote{妻子被法律束缚，只要她的丈夫活着；但如果丈夫死了，她就自由了，可以随意嫁人，只要在主内。}\BibleRef{格前 7:39}因此，宗徒剪除了每一种借口，并清楚宣告，如果女人在丈夫活着时再嫁，她就是奸妇。你不要向我谈论强暴者的暴力、母亲的恳求、父亲的命令、亲属的影响、仆人的傲慢和阴谋、家中的损失。丈夫可能是个奸夫或\OriginalTerm{鸡奸者}{sodomite}，可能染上各样罪行，并因他的罪被妻子离弃；但他仍然是她的丈夫，只要他活着，她就不可再嫁。宗徒不是凭自己的权柄颁布此项法令，而是凭那在他内说话的基督的权柄。因为他遵循了基督在福音中的话：\Quote{凡休妻的，除非因为奸淫，就是使妻子犯奸淫；凡娶那被休的妇女的，也犯奸淫。}\BibleRef{玛 5:32}请注意他所说的：\Quote{凡娶那被休的妇女的，就犯奸淫。}无论是她离弃丈夫，还是丈夫离弃她，娶她的男人仍是奸夫。因此，宗徒们看到婚姻的轭如此沉重，就对祂说：\Quote{如果人和妻子的情形是这样，倒不如不结婚，}主回答说：\Quote{能领受的，就领受吧。}并且立即以三个阉人的例子，他显示了不受肉体羁绊的守贞之福。\BibleRef{玛 19:10}\par

4. 我一直不太能确定她所说的\Quote{发现自己被迫}再嫁到底是什么意思。她所说的这个逼迫是什么？难道她被人群压倒并违背意愿被强暴？如果是这样，为什么她这个受害者事后没有离开那个强暴她的人？让她读梅瑟的书，她就会发现，如果在城里有人对已许配的处女施暴，而她没有喊叫，她就被当作奸妇受罚；但如果她是在田野被强迫，她就无罪，唯有那强暴者应受法律制裁。\BibleRef{申 22:23}因此，如果你的姊妹（正如她所说，是被迫进入第二次结合）愿意领受基督的身体，并且不被视为奸妇，就让她做补赎；至少从她开始悔改之时起，不再与那个第二个丈夫——他应当被称作奸夫而非丈夫——有进一步的来往。如果她觉得这太难，如果她不能离开她曾经爱过的人，并且不愿将主置于感官享乐之上，就让她听宗徒的宣告：\Quote{你们不能喝主的杯，又喝魔鬼的杯；你们不能分尝主的筵席，又分尝魔鬼的筵席。}\BibleRef{格前 10:21}在另一处：\Quote{光明与黑暗有什么相通？基督与\OriginalTerm{贝里雅耳}{Belial}有什么相和？}\BibleRef{格后 6:14}我接下来要说的话可能听起来新颖，但其实并不新，而是旧的，因为它有旧约的见证支持。如果她离开第二个丈夫，并希望与第一个丈夫和好，她现在不能这样做；因为申命记中记载：\Quote{若有人娶了妻，成了她的丈夫，后来发现她在他眼中不蒙喜悦，因为他在她身上找到了什么不洁的事，他就可以写一张休书，交在她手中，打发她离开夫家。她离开夫家之后，可以去嫁别人。如果后夫恨她，也写休书交在她手中，打发她离开夫家；或者娶她为妻的后夫死了；那么，打发她离去的原夫不可再娶她为妻，因为她已被玷污；这在主面前是可憎恶的；你不可使主你的天主赐给你为业之地陷于罪恶。}\BibleRef{申 24:1}所以我恳求你，尽力安慰她，敦促她寻求救恩。有病的肉体需要刀和烙铁。要怪的是伤口，而非医术；如果医生以真正仁慈的残忍，不忍则不忍，且仁慈是残忍的。\par

5. 你第三个也是最后一个问题涉及同一封书信中的一段，其中宗徒讨论复活时，来到这些话：\BibleQuote{因为祂必须为王，直到把一切仇敌放在祂脚下。最后被毁灭的仇敌就是死亡。因为祂已使万有服从在祂脚下。但说万有都服从祂时，显然那使万有服从于祂的除外。当万有都服从了祂时，子自己也要服从那使万有服从于祂的，好使天主成为万有中的万有。}\BibleRef{格前 15:25} 我感到惊讶，你竟决定就这一段来问我，因为那位可敬之人，\Place{普瓦捷}主教\Person{希拉利}，在他反对亚略派的论著的第十一卷中已对此作了全面的探讨和解释。不过，我至少可以说几句话。这段经文的主要绊脚石在于子被说成是服从于父。那么，哪个更可耻、更屈辱：是服从于父（这常是虔敬之爱的标记，如圣咏所说\BibleQuote{我的灵魂确实服从于天主}），还是被钉在十字架上并成为十字架的诅咒？因为\BibleQuote{凡挂在树上的，都受诅咒。}\BibleRef{迦 3:13} 如果基督为了我们成为诅咒，为把我们从法律的诅咒中解救出来，那么，祂也为了我们而服从于父，以便也使我们服从于祂，正如祂在福音中所说的：\BibleQuote{除非经过我，谁也不能到父那里去，}\BibleRef{若 14:6} 以及\BibleQuote{当我从地上被举起时，我要吸引众人归向我。}\BibleRef{若 12:32} 所以，基督在信徒中服从于父；因为所有信徒，乃至全人类，被视为祂身体的肢体。但在不信者中，即在犹太人、外邦人和异端者中，祂被说成是不服从的；因为这些祂身体的肢体并不服从于信德。但在世界末日，当所有祂的肢体看见基督，即他们自己的身体，为王时，他们也将服从于基督，即服从于他们的身体，好使整个基督的身体服从于天主父，并使天主成为万有中的万有。祂不是说\Quote{叫父成为万有中的万有}，而是说\Quote{叫天主成为}，这个称号本来属于天主圣三，不仅可以指父，也可以指子和圣神。所以，祂的意思是说\Quote{叫人性服从于天主性。} 这里我们所说的\Quote{人性}，不是指希腊人称为\OriginalTerm{仁爱}{philanthropy}的那种温和与慈善，而是指全人类。此外，当祂说\Quote{叫天主成为万有中的万有}时，应如此理解：目前，我们的主和救主并非万有中的万有，而只是在每人心中的一部分。例如，祂在\Person{撒罗满}那里是智慧，在\Person{达味}那里是慷慨，在\Person{约伯}那里是忍耐，在\Person{达尼尔}那里是对未来之事的认识，在\Person{伯多禄}那里是信德，在\Person{丕乃哈斯}和\Person{保禄}那里是热忱，在\Person{若望}那里是童贞，在其他人那里是其他德行。但当万物的终结来临时，祂将成为万有中的万有，因为那时圣人们将各自拥有所有德行，并且所有人都将完全拥有基督。\par

\SourceRecord{Translated by W.H. Fremantle, G. Lewis and W.G. Martley.；Nicene and Post-Nicene Fathers, Second Series；Vol. 6.；Edited by Philip Schaff and Henry Wace.；Buffalo, NY: Christian Literature Publishing Co.,；1893.；<http://www.newadvent.org/fathers/3001055.htm>.}

% structure: p0001 source=h1 target=section reason=page-title

\section{第五十七封书信}

\Emph{致\Person{帕玛丘}：论翻译的最佳方法}\par

\EditorialNote{写于公元395年，致\Person{帕玛丘}（参见\WorkTitle{第六十六封书信}）。前一年，\Person{热罗尼莫}在此描述的背景下（§2）将\WorkTitle{第五十一封书信}（从\Person{埃皮法尼乌斯}致\Place{耶路撒冷}的\Person{若望}）译为拉丁文。他的译本很快公开并遭到某人严厉批评，\Person{热罗尼莫}未指名，但怀疑受\Person{鲁菲努斯}煽动（§12）。被指控篡改原文后，他驳斥指控并为其翻译方法辩护（重意不重字，§5），援引古典作家（§5）、教会作家（§6）和新约作者（§7-10）的做法。}\par

此后，当\Person{鲁菲努斯}将他翻译的\Person{奥力振}\OriginalTerm{《第一原理》}{\GreekText{περί} \GreekText{᾿Αεχῶν}}版本发表时（\Person{热罗尼莫}认为这是误导人的），他援引这封信作为他行为的充分依据。参见\WorkTitle{第八十封书信}和\WorkTitle{第八十一封书信}，以及\Person{鲁菲努斯}为本系列第三卷所写的\OriginalTerm{《第一原理》}{\GreekText{περί} \GreekText{᾿Αεχῶν}}序言。\par

1. 宗徒\Person{保禄}在\Person{阿格黎帕}王面前为所受的指控辩护时，为了让听众明白并自信必获成功，他开口先自幸说：\BibleQuote{我认为自己是幸运的，\Person{阿格黎帕}王，因为今日我在你面前，要为犹太人控告我的一切事进行辩护；尤其因为你精通犹太人的一切风俗和争议。}\BibleRef{宗 26:2}。他读过耶稣的话：\BibleQuote{那说给愿意听的人听的人是有福的。}\BibleRef{德 25:9}。\par

2. 大约两年前，上述教父\Person{埃皮法尼乌斯}给\Person{若望}主教写了一封信，先批评他的一些观点，然后温和地召叫他悔改。由于作者的名望或信函的优雅，全\Place{巴勒斯坦}争相索要副本。现在，我们隐修院里有一个人，在他的国家颇有名望，即\Person{克雷莫纳的欧瑟伯}。他发现这封信广为流传，无知者和有识者都因其教导和文笔纯正而赞赏它，于是开始请求我为他将它译成拉丁文，同时简化论证，以便他更容易理解；因为他本人完全不懂希腊文。我同意了他的请求，叫来一名秘书，迅速口述了我的译本，并在页边简要标出各章的内容。事实上，他请我这样做只是为了他自己，我也要求他替我将抄本保密，不要轻易传播。一年半过去了，上述译本竟通过一种新奇策略从他的书桌传到了\Place{耶路撒冷}。因为一个假隐修士——很可能被收买（这是很有理由相信的），或出于他自己的恶意（正如诱惑者徒劳地试图说服我们的那样）——成了一个第二个\Person{犹达斯}，偷走了\Person{欧瑟伯}的文稿财产，给了仇敌一个毁谤我的机会\BibleRef{犹 1:9}。他们告诉无知的人说我篡改了原文，没有逐字翻译，把\Quote{亲爱的朋友}代替了\Quote{尊敬的大人}，更可耻的是！我删节了译文，省略了\OriginalTerm{尊敬的教父}{\GreekText{αἰδεσιμῶτατε} \GreekText{Πάππα}}这些词。这些以及类似的琐碎小事构成了对我的指控。\par

3. 首先，在我为我的译本辩护之前，我想问那些将智慧与狡猾混淆的人几个问题。你们从哪里得到那封信的抄本？谁给你们的？你们怎么有脸拿出通过欺诈得来的东西？如果连我们自己的墙壁和书桌内都不能隐藏秘密，还有什么安全的地方？如果我在法庭上对你们提起这样的指控，我就可以依法追究你们；即使在财政案件中，法律也会惩罚多管闲事的告密者，并在接受叛徒的背叛时谴责叛徒。因为他们虽然欢迎告发带来的利益，却谴责告发者的动机。不久前，一位名叫\Person{赫西基乌斯}的执政官级人物（宗主教\Person{加玛里耳}对他发动了一场无情的战争）被\Person{德奥多西}皇帝判处死刑，仅仅因为他通过被收买的秘书获取了皇家的文件。我们也读到古代历史中，背叛\Place{法里斯契}儿童的教师被送回他的学生那里，被捆绑交给他们，因为罗马人民拒绝接受不光彩的胜利。当\Place{伊庇鲁斯}国王\Person{皮洛士}因伤躺在营中时，他的医生提议毒死他；但\Person{法布里基乌斯}认为国王死于背叛是可耻的，便将叛徒捆绑送回他的主人那里，拒绝认可犯罪，即使受害者是敌人。法律维护的原则，敌人也维持，连战争和刀剑都不能侵犯的原则，在基督的隐修士和司铎中至今无人质疑。他们当中现在有人敢皱起眉头、打着响指，挥霍他的气息说：\Quote{就算他用了贿赂或其他手段！他做了合乎他目的的事。} 这真是为欺诈辩护的奇怪托词，好像强盗、小偷和海盗不这样做似的。当然，当\Person{亚纳斯}和\Person{盖法}引诱不幸的\Person{犹达斯}时，他们也只是做了他们认为对自己有利的事。\par

4. 假设我想在笔记本里写下这样或那样无聊的琐事，或对圣经作些评注，反驳诽谤我的人，平息我的愤怒，练习使用套话，并为战斗之日储备锋利的箭矢。只要我不公开这些想法，它们只是不友善的言辞，不足以构成诽谤指控；事实上，它们甚至称不上不友善的言辞，因为公众从未听闻。你可以贿赂我的奴隶，收买我的门客。你可以像寓言所说的那样，用你的黄金潜入\Person{达那厄}的闺房；然后，掩饰你所做的一切，你可以指责我造假；但如果你这样做，你将不得不承认自己犯了比你能加诸于我的更严重的罪。一个人斥责你是异端者，另一个人斥责你是教义败坏者。你自己保持沉默；你不敢回答；你攻击翻译者；你挑剔音节，并以为如果你的诽谤没有引起反驳，你的辩护就完备了。假设我在翻译中犯了错误或遗漏。你的整个案件都依赖于此；这就是你向控告者提供的辩护。难道因为我是个糟糕的翻译者，你就不是异端者了吗？请注意，我并没有说我知道你是异端者；我把这种认识留给你的控告者，写给那封信的人；我所要说的是，当一个人控告你时，你却攻击另一个人，这是极端的愚蠢；当你自己满身伤痕时，却通过伤害一个仍然平静、不进攻的人来寻求安慰。\par

5. 在上述评论中，我假设我对那封信作了改动，并且一个简单的翻译可能包含错误，尽管不是故意的。然而，既然信本身表明在意义上没有作出任何改变，没有增加任何内容，也没有将任何教义强行塞入其中，\Quote{显然他们的目的是理解却什么也不理解}；当他们想要指责别人的技巧不足时，反而暴露了自己的不足。因为我自己不仅承认，而且自由地宣告：在从希腊文翻译时（除了圣经的情况，其中连词序也是一个奥秘），我采用意义对意义而非字对字的方法。对于这种做法，我有\Person{图利乌斯}的权威，他翻译了\Person{柏拉图}的\WorkTitle{普罗泰戈拉}、\Person{色诺芬}的\WorkTitle{家政论}，以及\Person{埃斯基涅斯}和\Person{德摩斯梯尼}相互辩论的两篇优美演说。他作了哪些省略、增添和改动，用自己的语言习惯取代了另一种语言的习惯，这里不是细说的时候。我满足于引用那位翻译者在他的演说序言中所说的话：\Quote{我认为有责任承担一项劳动，虽非必要对我自己，却将证明对学习者有益。我翻译了两位最雄辩的阿提卡演说家的最高贵演说，即\Person{埃斯基涅斯}和\Person{德摩斯梯尼}相互辩论的那些演说；但我翻译时不是作为翻译者，而是作为演说家，保留意义但改变形式，通过调整隐喻和词汇以适合我们自己的习语。我不认为必须字对字翻译，而是再现了整体风格和重点。我不认为自己有义务逐字支付给读者，而是只给他等值的东西。} 在他工作的结尾，他又说：\Quote{如果我的译文被发现符合这个标准——我相信它会如此——我将非常满意。在译文中，我利用了原著的一切优点，即思想、表达方式和主题安排，而只在不违反我们审美观念的前提下遵循实际的措辞。如果我写的一切在希腊文中找不到，至少我努力使它与之相符。} \Person{贺拉斯}，一位敏锐而有学识的作家，在他的\WorkTitle{诗艺}中也给有技巧的翻译者提出了同样的建议：——\par

且勿怀过分焦虑的心思\newline{}逐字逐句去翻译。\par

\Person{泰伦提乌斯}翻译了\Person{米南德}；\Person{普劳图斯}和\Person{凯基利乌斯}翻译了旧喜剧诗人。他们何时拘泥于字词？难道他们在翻译中不是首先考虑保持原著的优美和魅力吗？像你这样的人所谓的忠实誊写，有学识的人称之为有害的琐细。大约二十年前，我的老师们就是这样；即使在那时，我也曾成为现在加诸于我的类似错误的受害者，尽管我从未想象过\Emph{你}会指责我。在将\Place{凯撒勒雅}的\Person{欧瑟伯}的\WorkTitle{编年史}翻译成拉丁文时，我在序言中作了如下评论：\Quote{遵循他人所划定的路线，有时不偏离它们是很困难的，并且在翻译中保持另一种语言中最恰当的表述的魅力也是很难的。每一个词都有它自己的含义，而我也许没有等同的词来翻译它；如果我用迂回的方式达到目标，我就得走许多里路来覆盖短短的距离。这些困难之外，还要加上倒装的曲折、格的用法的差异、隐喻的分歧；最后还有这种语言特有的、如果可以这样称呼的话，固有的特性。如果我逐字翻译，结果将显得笨拙；如果因必要而改变任何词序或措辞，我又似乎背离了翻译者的职能。} 在进行了长时间的讨论后——这里详述会很冗长——我又补充了以下内容：\Quote{如果有人以为翻译不会损害风格的魅力，就让他将\Person{荷马}逐字译成拉丁文——不，我还要更进一步说，让他译成拉丁文散文——结果将是词序显得可笑，而最雄辩的诗人几乎语无伦次。}\par

6. 我引用自己的著作，唯一目的是要证明：至少从年轻时起，我一直致力于传达意义而非字词；但如果这些著作提供的权威被认为不足，请阅读并思考那本叙述圣\Person{安当}生平的著作中关于此事的一篇简短序言。\Quote{直译从一种语言到另一种语言会遮蔽意义；生长的茂盛会减少收成。因为当措辞被格和比喻所束缚时，它不得不通过冗长的迂回解释那些本可用几个词就能说明白的事。我已在应你要求所译的\WorkTitle{圣\Person{安当}事迹}中尽量避免这个错误。我的译本始终保存意义，虽然不总是保留原文字词。让其他人去捕捉音节和字母吧，你则当寻求意义。} 若要我详述所有按意义翻译者的见证，时间也不够。目前只需提及宣信者\Person{依拉略}即可，他将一些关于\Person{约伯}的讲道和几篇关于\BibleBook{}{圣咏}的论述从希腊文译成了拉丁文；但他并没有受制于字句的昏睡，也没有被贫乏文化的陈腐字面主义所束缚。他如同征服者一般，将原文的意义掳获到自己的语言中。\par

7. 世俗作家与教会作家若采用此路线，本不足为奇，当我们想到\WorkTitle{七十贤士译本}的译者、福音作者和宗徒们在处理圣经时亦如此作为时。我们在马尔谷福音读到\BibleRef{谷 5:41}主说\OriginalTerm{塔里塔，古木}{Talitha cumi}，并且立刻加上\Quote{这被解释为：闺女，我吩咐妳，起来。}福音作者可能被指控说谎，因为他加了\Quote{我吩咐妳}这句话，因为希伯来文只有\Quote{闺女，起来。}为强调这一点，并造成呼唤和命令的印象，他添加了\Quote{我吩咐妳}。再次，在玛窦福音\BibleRef{玛 27:9}，当叛徒犹大退回三十块银钱，并用它们买了陶匠的田地时，记载说：\Quote{这就应验了耶肋米亚先知所说的话：\InnerQuote{他们拿了那三十块银钱，即以色列子民为被估定者所估定的价钱，给了它们为陶匠的田地，照主所吩咐我的。}}这段经文在耶肋米亚书中根本找不到，而是在匝加利亚书中，用相当不同的词句和完全不同的顺序。事实上，\WorkTitle{拉丁通行本}翻译如下：\Quote{我要对他们说：如果你们看为美，就给我价钱，否则就拒绝。他们就称了三十块银钱作为我的价钱。主对我说：把它们投入熔炉，并察看它是否像我被他们所试炼的那样。我就拿了那三十块银钱，抛在主殿中。}显然，\WorkTitle{七十贤士译本}的译文与福音作者的引文大相径庭。在希伯来文中，尽管意思相同，词句却完全不同，且排列不同。它说：\Quote{我对他们说：你们若以为美，就给我价钱；若不，就作罢。他们便称了三十块银钱作为我的价钱。主对我说：把它丢给陶匠；那是我被他们所估定的美好价钱。我就拿了那三十块银钱，在主殿中丢给陶匠。}他们可能会指控宗徒篡改了他的译本，因为它既不符合希伯来文，也不符合\WorkTitle{七十贤士译本}的译者；更糟糕的是，他们可能会说他弄错了作者的名字，写成了耶肋米亚，而应该是匝加利亚。我们绝不如此谈论一位基督的门徒，他致力制定教义，而非追逐字句音节。再举一个出自匝加利亚的例子，福音作者若望从希伯来文引用\Quote{他们要仰望他们所刺透的}，而我们在\WorkTitle{七十贤士译本}中读到\Quote{他们要仰望我，因为他们嘲弄了我}，在拉丁文译本中则是\Quote{他们要因他们所嘲弄或侮辱的事而仰望我。}这里，福音作者、\WorkTitle{七十贤士译本}和我们自己的译本都不同；然而，语言的分歧由精神的一致所弥补。在玛窦福音中，我们又读到主向宗徒们宣讲逃跑，并用一段匝加利亚的经文坚固他的劝告。他说：\Quote{经上记载：\InnerQuote{我要打击牧人，羊群中的羊就要四散。}}但在\WorkTitle{七十贤士译本}和希伯来文中，读法不同，因为不是天主在说话——如福音作者所认为的——而是先知向天主圣父恳求说：\Quote{打击牧人，羊群就要四散。}按我的判断——且我有一些谨慎的批评者与我同在——福音作者犯了一个错误，竟将先知的话归于天主。同一福音作者又写道，因天使的警告，若瑟带着婴孩和他的母亲逃往埃及，住在那里直到黑落德去世；\Quote{这就应验了主藉先知所说的话：\InnerQuote{我从埃及召回了我的儿子。}}\BibleRef{玛 2:13}拉丁文抄本并不如此记载，但在欧瑟亚书中\BibleRef{欧 11:1}，真正的希伯来文如下：\Quote{以色列年幼时，我就爱了他，从埃及召回了我的儿子。}\WorkTitle{七十贤士译本}则译作：\Quote{以色列年幼时，我就爱了他，从埃及召回了他的儿子们。}难道因为译者对一段主要指向基督奥秘的经文做了另一种解释，他们就要被全然拒绝吗？或者，我们岂不更应该因他们人性的软弱而宽恕译者的缺点，正如雅各伯所说的：\Quote{我们众人在许多事上都有过失。若有人在言语上没有过失，他就是完人，也能勒住全身。}\BibleRef{雅 3:2}同一福音作者的书页中又写道：\Quote{他来到并住在一座名叫纳匝肋的城：为应验先知所说的话：\InnerQuote{他将被称为纳匝肋人。}}\BibleRef{玛 2:23}让这些喜好辞藻和苛求文笔的批评家告诉我们，他们在哪里读过这些话；如果他们不能，就让我告诉他们，这些话在依撒意亚书中。\BibleRef{依 11:1}因为在我们在某处读到并翻译为\Quote{从叶瑟的树干将生出一根权杖，从他的根上将长出一根枝条}的地方，按希伯来成语是这样写的：\Quote{从叶瑟的根将生出一根权杖，从他的根上将长出一个纳匝肋人。}如果用一个词替换另一个词是非法的，那么\WorkTitle{七十贤士译本}怎能省略\Quote{纳匝肋人}这个词呢？无论隐藏奥秘还是轻视奥秘，都是亵渎。\par

8. 让我们转到其他段落，因为信函篇幅有限，不容我们过于详细地讨论任何一点。同一\Person{玛窦}记载说：—\Quote{这一切事的发生，是为应验主借先知所说的话：\InnerQuote{看，一位童贞女将怀孕，要生一个儿子，人要称祂的名字为厄玛奴耳。}} \OriginalTerm{七十贤士译本}{Septuagint}的译文是：\Quote{看，一位童贞女将受孕，要生一个儿子，你要称祂的名字为厄玛奴耳。} 如果人们吹毛求疵，显然\Emph{受孕}并不完全等同于\Emph{怀孕}，而\Quote{你要称}也与\Quote{人要称}不同。此外，在\OriginalTerm{希伯来文}{Hebrew}中我们读到这样：\Quote{看，一位童贞女将怀胎生子，要称祂的名字为依玛奴耳。} 阿哈次不会这样称祂，因为他被指为不信；犹太人也不会，因为他们注定要否认祂；而是那将要怀胎生子的，那童贞女本人。在同一福音书中我们读到，黑落德因贤士的到来而惊慌，就召集了经师和司祭，问他们基督当生在何处，他们回答说：\Quote{在犹大的白冷，因为先知这样记载：\InnerQuote{你犹大地啊的白冷，你在犹大的首领中绝不是最小的，因为将由你出来一位领袖，祂将治理我的百姓以色列。}}\BibleRef{玛 2:5}在\OriginalTerm{拉丁通行本}{Vulgate}中，这段经文如下：\Quote{你白冷，厄弗辣大的家，你在犹大的千邑中虽小，却有一位将为我从你出来，作以色列的首领。} 如果你查看希伯来文如下，就会对玛窦和七十贤士译本之间措辞和顺序的差异更加惊讶：\Quote{但是你白冷厄弗辣大，你在犹大的千邑中虽小，却有一位将为我从你出来，作以色列的统治者。}\BibleRef{米 5:2}逐一考虑福音书作者的话：\Quote{你犹大地啊的白冷。} 因为\Quote{犹大地}希伯来文作\Quote{厄弗辣大}，而七十贤士译本作\Quote{厄弗辣大的家}。福音书作者写道：\Quote{你在犹大的首领中绝不是最小的。} 在七十贤士译本中这是：\Quote{你在犹大的千邑中虽小，} 而希伯来文给出：\Quote{你在犹大的千邑中虽小。} 这里有一个矛盾——不仅是字面上的——在福音书作者和先知之间；因为至少在此处，七十贤士译本和希伯来文是一致的。福音书作者说他（弥赛亚）在犹大的首领中并不小，而他所引用的经文恰恰说的是相反的意思：\Quote{你确实小而且微小；但由你——虽然又小又微——却要为我出来一位以色列的领袖}，这个思想与宗徒的话语一致：\BibleQuote{天主拣选了世上软弱的，来羞辱那强壮的。}\BibleRef{格前 1:27}此外最后一句\Quote{治理}或\Quote{牧放我的百姓以色列}，显然在原文中也有不同表达。\par

9. 我提及这些段落，不是要指控福音作者篡改——这指控只配像策尔苏斯、波斐利和尤利安这样的不敬之人提出——而是要让我的批评者认识到自己的无知，并求得这样的体谅，使他们在我一封简单书信的事上可以允许我，无论他们喜不喜欢，他们在神圣圣经的事上都必须允许宗徒们。马尔谷，伯多禄的门徒，这样开始他的福音：\BibleQuote{耶稣基督福音的开始，正如在\BibleBook{依撒意亚}{书}上所记载的：看，我派遣我的使者在你面前，他要在你前面预备你的道路。旷野中有呼号者的声音：你们当预备上主的道路，修直他的途径。} 这引文由两位先知的话组成，即玛拉基亚和依撒意亚。因为第一部分：\BibleQuote{看，我派遣我的使者在你面前，他要在你前面预备你的道路。} 出现在\BibleBook{玛拉基亚}{书}的结尾。\BibleRef{拉 3:1} 但第二部分：\BibleQuote{旷野中有呼号者的声音，等等。} 我们在\BibleBook{依撒意亚}{书}中读到。\BibleRef{依 40:3} 那么，马尔谷凭什么在他书的开头就写下\BibleQuote{正如在\BibleBook{依撒意亚}{书}上所记载的：看，我派遣我的使者}，而正如我们所说，这话根本不在\BibleBook{依撒意亚}{书}中，而是在十二位先知中最后一位\BibleBook{玛拉基亚}{书}中？让无知的自负者若能做到就来解答这个难题吧，然后我愿为自己的错误请求原谅。同一马尔谷向我们呈现救主这样对法利塞人说：\BibleQuote{你们从来没有读过，达味在需要时和饥饿时，他和他的人所做的事吗？他怎样进了天主的殿，在阿彼雅塔尔大司祭的时候，吃了陈设饼，这饼除了司祭外，别人是不许吃的？} \BibleRef{谷 2:25} 现在让我们翻阅\BibleBook{撒慕尔}{纪}，或通常所称的\BibleBook{列王}{纪}，我们会发现那里大司祭的名字不是阿彼雅塔尔，而是阿希默肋客，\BibleRef{撒上 21:1} 后来他与其余司祭一同被多厄格奉撒乌耳之命处死。现在让我们继续看宗徒保禄，他这样写信给格林多人：\BibleQuote{因为他们若知道，就不会把光荣的主钉在十字架上了。但是，正如经上所载：眼所未见，耳所未闻，人心所未想到的，就是天主为爱祂的人所准备的。} \BibleRef{格前 2:8} 有些作者在这段经文上转向伪经的狂言，断言这引文出自厄里亚的默示录；而事实上，它出现在\BibleBook{依撒意亚}{书}中，按希伯来文原文是：\BibleQuote{从太初以来，人未曾听见，未曾耳闻，未曾眼见，天主啊，除你以外，你为那等候你的人所预备的是什么。} 七十贤士译本则完全不同地翻译了这些话：\BibleQuote{从太初以来，我们未曾听见，我们的眼也未曾看见除你以外有任何神，以及你的真实作为，你必向那等候你的人显示仁慈。} 因此我们看明了这引文出自何处，但宗徒并没有逐字直译他的原文，而是使用了释义，用不同的话表达了意义。在同一宗徒致罗马人的书信中，他引用了\BibleBook{依撒意亚}{书}的这些话：\BibleQuote{看，我在熙雍安放了一块绊脚石和使人失足的磐石。} \BibleRef{罗 9:33} 这个译文与希腊文译本不同，却与希伯来文原文一致。七十贤士译本给出了相反的意义：\BibleQuote{免得你碰到绊脚石或失足的磐石上。} 宗徒伯多禄同意保禄和希伯来文，写道：\BibleQuote{但对那些不信的人，是一块绊脚石和使人失足的磐石。} 从所有这些段落可以清楚看出，宗徒和福音作者在翻译旧约圣经时，力求传达意义而非字句，他们并不十分在意保留形式或结构，只要能将主题清晰地向理解力阐述即可。\par

10. 路加，福音作者和宗徒的同伴，描述基督的首位殉道者斯德望在犹太会众中讲述如下。\BibleQuote{雅各伯带着七十五个人下到埃及，他自己死了，我们的祖先被带到舍根，安放在亚伯拉罕用银钱从舍根的父亲哈摩尔的儿子们那里买来的坟墓里。} \BibleRef{宗 7:15} 在\BibleBook{创世}{纪}中，这段记载完全不同，因为是亚伯拉罕从赫特人厄斐龙，祚哈尔的儿子，在赫贝龙附近，用四百银币买了一块双洞田地和周围的田地，并把他妻子撒辣埋葬在那里。而在同一书中我们读到，雅各伯从美索不达米亚带着妻子儿子回来后，在舍根城，即舍根地方，搭帐棚，并住在那里，从哈摩尔，舍根的父亲，用一百只羊羔买了一块搭帐棚的田地，在那里筑了一座祭台，呼求了以色列的天主。而亚伯拉罕不是从哈摩尔，舍根的父亲那里买洞，而是从厄斐龙，祚哈尔的儿子那里买；他也不是埋葬在舍根，而是在赫贝龙，即讹传的阿贝尔。而十二位圣祖并非埋葬在阿贝尔，而是在舍根，在雅各伯所买（不是亚伯拉罕所买）的田地里。我把这个难题的解答推迟一下，好让那些挑剔我的人去考察，并明白在对待圣经时，我们应当看重的是意义而不是字句。在希伯来文中，\BibleBook{圣咏}{第二十二篇}以主在十字架上所说的话开始：\Emph{Eli Eli lama azabthani}，意思是 \OriginalTerm{我的天主，我的天主，你为什么舍弃了我}{Eli Eli lama azabthani}？让我的批评者告诉我，为什么七十贤士译本在这里加上了\Quote{求你眷顾我}这几个字。因为它的译文如下：\BibleQuote{我的天主，我的天主，求你眷顾我，你为什么舍弃了我？} 他们无疑会回答说，增加几个字无损于意义。那么让他们承认，如果我在口授的仓促中遗漏了几个字，我并没有因此危及教会所处的地位。\par

11. 现在要详细列举七十贤士译本作了多少增添和删节，以及教会抄本中用剑号和星号标记的所有经文，未免冗长。犹太人听到我们翻译的\Person{依撒意亚}这段经文时，常会嘲笑：\BibleQuote{有种子在\Place{熙雍}、有仆人在\Place{耶路撒冷}的人，是有福的。}在\BibleBook{亚毛斯}{书}中，在描述放纵情欲\BibleRef{亚 6:4}之后，又有这些话：\BibleQuote{他们视这些事为跛脚的，似乎不会飞翔。}这是一句极富修辞的句子，完全不逊于\Person{图利}。然而，我们该如何对待希伯来原文呢？这些经文和其他类似经文在希伯来原文中被省略了，数量如此之多，若要全部重现，需要数不清的书籍。上述星号以及我的译本与旧译本（经细心读者比较）共同显示了省略的数量。不过，七十贤士译本在教会中仍保有其正当地位，或许因为它是所有译本中最早的，在\Person{基督}来临之前就已出现，又或许因为它曾被宗徒们使用（但仅在不与希伯来原文相悖之处）。另一方面，我们理应拒绝\Person{阿奎拉}——这位改宗者兼争议性的译者，他不仅试图翻译词语，还要翻译词源。谁能接受他用\OriginalTerm{倾注、收成、闪耀}{\GreekText{χεῖμα} \GreekText{ὀπωρισμός} \GreekText{στιλπνότης}}来翻译\BibleQuote{五谷、新酒和油}\BibleRef{申 7:13}呢？或者，由于希伯来文除了冠词之外还有其他前缀，他竟以令人不悦的学究气一个音节一个音节、一个字母一个字母地翻译，如此写道：\OriginalTerm{连同天和地}{\GreekText{σὺν} \GreekText{τὸν} \GreekText{ὀυρανὸν} \GreekText{καὶ} \GreekText{σὺν} \GreekText{τὴν} \GreekText{γήν}}，这种结构希腊语和拉丁语都无法接受，我们自己的作者中有不少例子可以证明。有多少在希腊语中优美的短语，如果逐字翻译，在拉丁语中听来并不悦耳；同样，有多少我们拉丁语中悦耳的短语——假设词语顺序不变——在希腊语中却不动听。\par

12. 但为了绕过这无边的讨论领域，并向您——最基督徒的贵族、最高贵的基督徒——表明我的译本被指责的是何种篡改，我将把书信开头几行以希腊原文和我的译文呈现在您面前，以便您从一项控诉中对全部作出判断。书信开头写着希腊文原文：\par

\SourceRecord{Translated by W.H. Fremantle, G. Lewis and W.G. Martley.；Nicene and Post-Nicene Fathers, Second Series；Vol. 6.；Edited by Philip Schaff and Henry Wace.；Buffalo, NY: Christian Literature Publishing Co.,；1893.；<http://www.newadvent.org/fathers/3001057.htm>.}

% structure: p0001 source=h1 target=section reason=page-title

\section{第五十八书信}

\Emph{致保利诺}\par

\EditorialNote{在这封致诺拉保利诺的第二封信中，热罗尼莫劝阻他不要前往圣地朝圣，并描述了耶路撒冷本来的面目而非应有的样子。他然后给朋友提供了类似之前给内波齐亚诺的建议，称赞保利诺为狄奥多西皇帝所写的颂词（现已失传），并比较了他的文风与拉丁教会伟大作家，最后推荐了他的信使，即那个维吉朗提乌斯，他很快就成为他最轻蔑的对象。写于公元395年左右。}\par

1. \BibleQuote{善人从心里的善库发出善来，} \BibleRef{玛 12:35} 并且 \BibleQuote{每棵树凭它的果实就能被认识。} \BibleRef{路 6:44} 你以你自己的德行尺度来衡量我，并因你自己的伟大而夸大我的渺小。你在宴席上取了末座，好让家主叫你再上座。\BibleRef{路 14:10} 我身上有什么，或我有什么品质，配得博学之人的称赞？我这个渺小卑微的人，竟能得到曾为我们最敬虔的君主辩护的嘴唇的颂扬？我最亲爱的弟兄，不要以我的年岁来估量我的价值。白发不是智慧；智慧才与白发等同。至少撒罗满这样说：\BibleQuote{智慧是人的白发。} \BibleRef{智 4:9} 梅瑟在拣选七十长老时，也被吩咐要取那些他确认为长者的人，不是按年岁而是按明智来拣选。\BibleRef{户 11:16} 并且，孩童达尼尔审判老年人，在青春年华时谴责年老的不节制。我再说，不要以年岁来秤量信德，也不要因我比你更早加入基督旗下而自以为胜过于你。宗徒保禄，那个从迫害者中选出的器皿，\BibleRef{宗 9:15} 虽然在宗徒次序中最后，但在功劳上却是第一。因为虽然后来，他却比众人更劳苦。\BibleRef{格前 15:10} 对犹达斯曾说过：你是与我同吃甜食的人，我的向导和我的熟人；我们曾在人群中同行于天主的殿中；但救主却指责他出卖朋友和老师。维吉尔的一行诗很好地描述了他的结局：\par

\Quote{他从高梁上系起可怕的死亡。}\par

相反，垂死的强盗以十字架换取了天堂，并将谋杀的刑罚转变为殉道。如今有多少人活得如此之久，以至于他们承载的是尸体而非身体，就像粉饰的坟墓，里面充满死人的骨头！\BibleRef{玛 23:27} 新燃起的炽热比长期持续的微温更为有效。\par

2. 至于你，当你听到救主的劝言：\BibleQuote{你若愿意成为成全的，去，变卖你所有的，施舍给穷人，然后来跟随我，} \BibleRef{玛 19:21} 你就将他的话付诸行动；并赤身跟随赤裸的十字架，你更容易攀登雅各伯的梯子，因为不带任何重担。你的衣着随着你的信念改变，你并不追求显眼的褴褛，以免你的钱袋依然饱满。不，以洁净的手和清白的良心，你以在神贫上和实际上都贫穷为荣耀。穿着悲伤或毁容的面容，假装并炫耀斋戒，或在你保留大笔收入时穿着廉价的外衣，并没有什么伟大。昔日的泰贝富翁克拉特斯，在前往雅典学习哲学的路上，抛弃了数不清的金子，相信财富与美德不能相容。他给我们，贫穷的基督的门徒，一个多么大的对比！我们却往口袋里塞满金子，并以施舍为借口坚守旧有的财富。当我们如此怯懦地扣留自己的东西时，我们怎能忠实地分配属于别人的东西？\BibleRef{路 16:12} 当肚子饱足时，谈论斋戒很容易。值得称赞的不是曾到过耶路撒冷，而是在那里度过了良好的生活。我们应当赞美和寻求的城，不是那杀害先知、\BibleRef{玛 23:37} 流基督血的城，而是那被河流的溪流所喜乐的城，是建在山上的城，不能隐藏，\BibleRef{玛 5:14} 是宗徒称为圣徒之母的城，\BibleRef{迦 4:26} 他乐于在其中与义人一同拥有公民权。\par

3. 我这样说，并非使自己陷于矛盾之嫌，也非谴责自己曾走过的路。我相信，我像\Person{亚巴郎}一样离开家乡和亲族，并非徒然。但我并不敢限制天主的全能，也不将连天都不能容纳的那一位拘禁在一小块土地上。每个信徒受审判，不是因他住在这里或那里，而是按他信德应得的赏报。真正的朝拜者既不在这\Place{耶路撒冷}，也不在\Place{革黎斤山}朝拜圣父；因为\BibleQuote{天主是神，朝拜祂的人必须以心神和真理朝拜祂。}\BibleRef{若 4:24}\Quote{风随意向哪里吹，}\Quote{大地和其中的万物属于上主。}\BibleRef{民 6:36}且有许多人从\Place{东方}和\Place{西方}而来，\BibleRef{路 13:29}坐在\Person{亚巴郎}的怀中；\BibleRef{路 16:22}那时，天主就不再仅在\Place{犹大}被认识，祂的名不再仅在\Place{以色列}被尊崇；宗徒的呼声传遍全地，他们的话语传到地极。救主亲自在圣殿中对门徒说：\BibleQuote{起来，我们离开这里，}\BibleRef{若 14:31}又对犹太人说：\BibleQuote{你们的房屋将留给你们荒芜。}\BibleRef{玛 23:38}如果天地都要过去，\BibleRef{路 21:33}显然所有尘世之物也都将过去。因此，目睹十字架和复活的地点，仅对那些背负自己十字架、每日与基督一同复活、并以此显示自己配得如此神圣居所的人才有益处。那些说\BibleQuote{上主的圣殿，上主的圣殿}\BibleRef{耶 7:4}的人，应听从宗徒的话：\BibleQuote{你们就是上主的圣殿，}\BibleRef{格后 6:16}\Person{圣神}\BibleQuote{住在你们内。}\BibleRef{罗 8:11}从进入天堂的庭院而言，从\Place{不列颠}与从\Place{耶路撒冷}同样容易；因为\BibleQuote{天主的国就在你们内。}\BibleRef{路 17:21}\Person{安东尼}和\Place{埃及}、\Place{美索不达米亚}、\Place{本都}、\Place{卡帕多细亚}、\Place{亚美尼亚}的众多隐修士从未见过\Place{耶路撒冷}，但乐园的门却为他们远在异地而开。真福\Person{依拉略}虽是\Place{巴勒斯坦}本地人和居民，却只在\Place{耶路撒冷}停留了一天，一方面不愿在近在咫尺时忽略圣地，另一方面也不愿将天主局限在地域之内。从\Person{哈德良}时代到\Person{君士坦丁}时代——约一百八十年——那曾目睹复活的地点被\OriginalTerm{朱庇特}{Jupiter}像占据；而在竖立十字架的磐石上，异教徒设立了一尊\OriginalTerm{维纳斯}{Venus}的云石雕像，成为朝拜的对象。最初的迫害者确实以为，玷污我们的圣地会剥夺我们对苦难和复活的信仰。就连我自己的\Place{白冷}，如今这个全世界最可敬的地方，圣咏作者曾歌颂：\Quote{真理从地上发出，}却被\OriginalTerm{塔慕次}{Tammuz}的丛林遮蔽，\BibleRef{则 8:14}即\OriginalTerm{阿多尼}{Adonis}的丛林；正是在那婴儿基督发出最初啼哭的洞穴里，人竟为\OriginalTerm{维纳斯}{Venus}的淫夫哀悼。\par

4. 你或许会说，我为何要做这些遥远的暗示？是为使你确信：你虽未见过\Place{耶路撒冷}，你的信德却一无所缺，而我住在\Place{耶路撒冷}也并不因此而更好。请放心，无论你住在这里或别处，你的善工与我们的主必得同样的赏报。的确，若我坦诚表达自己的感受，当我考虑你的誓愿以及你弃绝世界的热忱时，我认为只要你在乡下生活，一处地方与另一处同样好。离开城市和它们的拥挤，住在一小块土地上，在独处中寻求基督，独自与耶稣一起在山上祈祷，\EditorialNote{路六}常接近圣地：远离城市，我再说，你将永不失去你的圣召。我的劝告不涉及主教、长老或圣职人员，因为他们有不同职责。我只对一位隐修士说话：他原是世间的名人，将财产的代价放在宗徒脚前，\BibleQuote{将财产的代价放在宗徒脚前}\BibleRef{宗 4:37}以告诉人们必须践踏钱财，并且决定过孤独隐退的生活，永远继续拒绝曾经拒绝过的一切。倘若苦难与复活的场景不是在一个人口稠密的城市——有宫廷、驻军、妓女、戏子、小丑以及通常这类中心聚集的各色人等——倘若拥挤的人群是由隐修士组成，那么城市对那些选择了隐修生活的人就是可向往的居所。但正如现实情况，那真是愚蠢至极：先是弃绝世界、否定祖国、离开城市、自称为隐修士，然后却生活在比在自己本国更庞大的人群中，过着同样的生活。人们从世界各地涌向这里，城市充满了各种种族的人，男女拥挤如此之多，以至于你将在这里充分忍受你曾在别处部分发现的罪恶——而你原本想逃避它。\par

5. 既然你以弟兄的身份问我应走哪一条路，我就坦白地对你说。如果你愿意尽司铎的职务，并向往属于主教的职务或尊荣，就住在城市和城垣之内，这样借着使他人得救而使自己的灵魂获益。但如果你希望在实际生活中成为你名义上所是的人——即隐修士，也就是独居者，那么你和城市有什么关系呢？城市不是独居者的家，而是人群的家。每一种生活方式都有其代表人物。例如，让罗马的将军们效法像卡米卢斯、法布里基乌斯、雷古鲁斯和西庇阿这样的人。让哲学家们以毕达哥拉斯、苏格拉底、柏拉图和亚里士多德为楷模。让诗人们努力与荷马、维吉尔、米南德和泰伦提乌斯争胜。让历史学家们追随修昔底德、萨卢斯特、希罗多德和李维。让演说家们以吕西亚斯、格拉古兄弟、狄摩西尼和图利乌斯为师。至于我们自己的情况，让主教和司铎们以宗徒或他们的同伴为榜样；既然他们拥有这些人曾拥有的职位，就应努力展现同样的美德。最后，让我们隐修士以保禄、安东尼、犹利安、希拉利翁和玛加略们的生平为我们要效法的楷模。回到圣经的权威，我们有厄里亚和厄里叟为师傅，有先知弟子们为领袖；他们住在田野和僻静之处，在约旦河边支搭帐棚。\BibleRef{列下 6:1} 勒加布的儿子们也属于这类人，他们不喝酒和烈性饮料，住在帐棚中；天主的声音通过耶肋米亚赞扬他们，\EditorialNote{耶肋米亚第三十五章}并许给他们：他们的后代中总不断有人侍立在天主面前。\BibleRef{耶 35:19} 第七十一篇圣咏的标题很可能就是指这个：\Quote{属于约纳达布的儿子们和那些首先被掳去的人。} 这里所指的人是勒加布的儿子约纳达布，在列王纪中记载他上了耶胡的战车。\BibleRef{列下 10:15} 他的儿子们一直住在帐棚中，直到最后（由于加色丁军队的入侵）被迫进入耶路撒冷，他们被描述为\BibleRef{耶 35:11}首先遭受被掳的人；因为他们离开孤独生活的自由后，觉得城中的拘禁如同监禁一样。\par

6. 既然你并非完全自主，而是与一位在主内是你的姊妹的妻子结合，我恳求你——无论在这里或那里——远离人群、正式和礼貌的拜访以及社交聚会，所有这些都会束缚灵魂。让你的食物粗糙——比如卷心菜和豆类——并且直到傍晚才吃。有时作为一种佳肴，你可以吃一点小鱼。渴慕基督并以真粮为食的人，不太在意那些终将化为粪土的美味。食物一旦经过咽喉你就无法品尝其味，那么对你来说，面包和豆类也就够了。你有我反对约维尼安的著作，其中更详尽地谈论了轻视食欲和口味。让某本圣书常在你手中。不断祈祷，俯身使你的心灵上达于主。经常守夜，常常空腹入睡。避免闲言碎语和一切自我夸耀。逃避花言巧语的奉承者如同逃避公开的敌人。亲手分发物资以减轻穷人和弟兄们的痛苦。我再说，亲手，因为人间的诚信是稀有的。你不信我的话吗？想想茹达斯和他的钱囊。不要为膨胀的灵魂寻求卑贱的服装。避免与世俗之人交往，尤其是那些有权势的人。为何再看那些因轻视而成为隐修士的事物？最重要的是，让你的姊妹远离已婚妇女。如果她周围的妇女穿着丝绸和宝石，而她的衣着朴素，她既不要烦恼也不要自得。因为这样做，她要么后悔自己的决定，要么播下骄傲的种子。如果你已经以忠信地管理自己的财产而闻名，就不要拿别人的钱来施舍。你明白我的意思，因为主在一切事上给了你智慧。要如同鸽子般纯朴，不设圈套害人；又要如同蛇般机警，不让人设圈套害你。\BibleRef{玛 10:16}因为一个容许自己被欺骗的基督徒，其过错几乎与试图欺骗别人的人一样大。如果有人总是或几乎总是与你谈论金钱（除非是施舍，这是人人都可以谈论的话题），就要把他当作商贾而非隐修士。除了食物、衣服和明显必要之物外，不要给任何人任何东西；因为狗不能吃孩子的饼。\BibleRef{玛 15:26}\par

7. 基督的真正圣殿是信徒的灵魂；装饰它，给它穿衣，向它奉献礼物，在它里面欢迎基督。当基督在祂贫穷的人身上\BibleRef{玛 25:40}面临饿死的危险时，墙壁上镶满宝石有何用？你的财产不再属于你，而是委托你管理。要记住阿纳尼雅和撒斐辣，他们因惧怕未来而保留了自己的财产；你要小心，不要鲁莽地挥霍属于基督的。也就是说，不要因判断错误，将穷人的财物给予那些并非穷人的人；免得，如一位智者所说，爱德成了爱德的死亡。不要看\par

华丽的衣饰或加图空洞的虚名。\newline{}按佩尔西乌斯的话，天主说：\newline{}\Quote{我知道你的思想，能读懂你内心的深处。}\par

成为基督徒是大事，而不仅仅是看起来像。不知怎的，那些最不讨基督喜悦的人往往最讨世界的喜欢。我说这些，并不像母猪教训密涅瓦；而是如同朋友劝诫朋友，在你开始新道路之前劝告你。我宁愿能力不足，也不愿缺乏服侍你的意愿；因为我的愿望是：在我跌倒的地方，你能站稳。\par

8. 我非常愉快地阅读了你送给我的那本书，书中你以智慧和雄辩为皇帝 \Person{Theodosius} 进行了辩护；你对主题的安排尤其令我满意。在前几章中，你超越了他人，而在后几章中，你超越了你自己。你的文风简洁而优雅，具有 \Person{Tully} 的全部纯洁性，却又蕴含深意。因为，正如有人所说，那种只被人称赞辞藻的演讲是失败的。你成功地保持了主题的连续性和逻辑的关联。无论取哪一句，它都是前文的结论或后文的引言。\Person{Theodosius} 是幸运的，因为他有你这样的基督教演说家为他辩护。你使他的紫袍光彩夺目，并将他有益的律法奉献给后世。继续前进并取得成就吧，因为如果这就是你在战场上的初次尝试，那么当你成为一位训练有素的战士时，你还有什么不能做到呢？啊！我多么希望能引导像你这样的天才，不是（如诗人所唱的）穿过 \Place{Aonian} 群山和 \Place{Helicon} 之巅，而是经过 \Place{Zion}、\Place{Tabor} 和 \Place{Sinai} 的高处。如果我能教导你我所学到的，并传给你先知们的神秘卷轴，那么我们可能会产生一些希腊人以其全部学识都无法展示的东西。\par

9. 因此，请听我说，我的同仆、我的朋友、我的兄弟；请稍作倾听，让我告诉你应当在圣经中如何行走。我们在神圣书卷中所读的一切，虽然外表闪耀发光，但内在却更为甘甜。\Quote{想要吃果仁的人必须先破开坚果。} \BibleQuote{求你开启我的眼睛，} \Person{David} 说，\BibleQuote{使我得以明白你法律中的奇事。} 如果如此伟大的先知都承认自己处于无知的黑夜中，那么你认为，我们这些不过是婴孩和尚未断奶的婴儿，所被笼罩的误解之夜该有多么深沉！这帕子不仅蒙在 \Person{Moses} 的脸上，\BibleRef{格后 3:14} 也蒙在福音书作者和宗徒们身上。救主对群众只讲比喻，并且为表明他的话有属灵意义，就说：\Quote{有耳可听的，就应当听。} \BibleRef{路 8:8} 除非一切所写的都由那\Quote{拿着 \Person{David} 的钥匙，开了就没有人能关，关了就没有人能开的} \BibleRef{默 3:7} 那位开启，否则没有人能解开锁链，将之呈现在你面前。只要你有了惟独他能赐予的根基；不，只要他的手指甚至掠过你的作品，你的著作就会无比美好、无比博学、无比迷人、无比拉丁化。\par

10. \Person{Tertullian} 含义丰富，但风格粗犷生硬。有福的 \Person{Cyprian} 像一股纯净的泉水，柔和甜美地流淌，但由于他专注于劝勉德行和迫害带来的烦恼，他从未探讨过圣经。\Person{Victorinus} 尽管拥有殉道者的冠冕，却无法表达他所知道的。\Person{Lactantius} 拥有堪与 \Person{Tully} 媲美的雄辩才华：但愿他教导我们教义的准备能像他推翻他人教义那样充分！\Person{Arnobius} 冗长而不均衡，常常因对主题划分不当而混乱。那位可敬的 \Person{Hilary} 因其高卢悲剧风格而显得崇高；然而，虽然他饰有希腊修辞的花朵，有时却陷入冗长的句子，为学识较浅的兄弟提供绝不轻松的阅读。对其他已故或仍在世的作者，我什么也不说；其他人将根据善恶来评判他们。\par

11. 我要来对你说，我的同奥秘者、我的同伴、我的朋友；我的朋友，我这样说，虽然尚未亲自认识：请你不要怀疑一个如此亲密的人会是谄媚者。你宁可认为我失误或受情感误导，也不愿认为我会对朋友阿谀奉承。你拥有伟大的智力和取之不尽的词汇，你的措辞流畅纯洁，你的流畅和纯洁中蕴含着智慧。你的头脑清晰，所有感官敏锐。如果你在这智慧和口才之上，再加上对圣经的认真研究与知识，我很快就会看到你能够据守我们的堡垒，抵挡一切来犯者；你将与 \Person{Joab} 一同登上 \Place{Zion} 的屋顶，\BibleRef{编上 11:5} 并在房顶上唱出你在内室中所学到的。\BibleRef{路 12:3} 请你束紧腰，束紧你的腰。正如 \Person{Horace} 所说：\par

\Quote{生命不给予人任何礼物，除非他们辛劳。}\par

在教会中显出你曾在元老院中那样的名人风范。为你自己预备财富，你可以每天花费却不会耗尽。趁你仍然强壮，头上尚无白发时预备它们：在 \Person{Virgil} 的话到来之前：\par

\Quote{疾病向你袭来，以及阴郁的老年，\newline{} 和痛苦，以及残酷死亡的严峻无情。}\par

我对你的平庸并不满足：我渴望你所做的一切都是最卓越的。\par

我是多么由衷地欢迎了可敬的司铎 \Person{Vigilantius}，他自己的口比这封信更能告诉你。他为何这么快离开我们又重新出发，我说不上来；而且，我实在不愿伤害任何人的感情。不过，尽管他只是一个过路人，急于继续他的旅程，我还是设法留住他，直到我让他尝到了我对你的友谊。这样你就可以从他那里了解你想知道的关于我的事。请代我问候你可敬的姊妹和同仆，她与你一同在主内打那美好的仗。\par

\SourceRecord{Translated by W.H. Fremantle, G. Lewis and W.G. Martley.；Nicene and Post-Nicene Fathers, Second Series；Vol. 6.；Edited by Philip Schaff and Henry Wace.；Buffalo, NY: Christian Literature Publishing Co.,；1893.；<http://www.newadvent.org/fathers/3001058.htm>.}

% structure: p0001 source=h1 target=section reason=page-title

\section{第五十九封书信}

\Emph{致 \Person{玛尔切拉}}\par

对\Person{玛尔切拉}在一封未保存的信中向\Person{热罗尼莫}提出的五个问题的回答。问题如下。\par

（1）眼所未见，耳所未闻\BibleRef{格前 2:9}是什么？\Person{热罗尼莫}回答说，它们是属灵之事，因此只能以属灵的方式辨识。\par

（2）将\Person{基督}比喻\BibleRef{玛 25:31}中的绵羊和山羊等同于基督徒和异教徒，难道不是错误吗？它们难道不是善人和恶人吗？对于这个问题，\Person{热罗尼莫}引导\Person{玛尔切拉}参考他反对\Person{约维尼安}的著作（第二卷，第18-23节）。\par

（3）\Person{保禄}说，有些人要\Quote{活着存留直到主来临}，并且他们要\Quote{被提到空中与主相遇}\BibleRef{得前 4:15}。我们是否应认为这个提升是肉身的，并且那些被提升的人将免于死亡？是的，\Person{热罗尼莫}回答说，但他们的身体将被荣耀。\par

（4）\BibleRef{若 20:17}的\Quote{不要摸我}如何与\BibleRef{玛 28:9}的\Quote{她们上前抱住他的脚}相协调？\Person{热罗尼莫}回答说，第一种情况是，\Person{玛利亚玛达肋纳}没有认出\Person{耶稣}的天主性；而第二种情况中，妇女们认出了。因此，她们得到了被拒绝给予她的特权。\par

（5）复活的\Person{基督}在升天之前是只与门徒同在，还是同时也与天使同在，在天堂并其他各处？根据\Person{热罗尼莫}，后者是正确的教义。他写道：\Quote{天主性完全存在于各处。因此，\Person{基督}同时与宗徒和天使同在，在\Person{圣父}内，也在海洋的极深处。此后，他与\Person{多默}同在\Place{印度}，与\Person{伯多禄}同在\Place{罗马}，与\Person{保禄}同在\Place{伊利里亚}，与\Person{弟铎}同在\Place{克里特}，与\Person{安德肋}同在\Place{阿哈雅}。}这封信的日期是公元395年或396年。\par

\SourceRecord{Translated by W.H. Fremantle, G. Lewis and W.G. Martley.；Nicene and Post-Nicene Fathers, Second Series；Vol. 6.；Edited by Philip Schaff and Henry Wace.；Buffalo, NY: Christian Literature Publishing Co.,；1893.；<http://www.newadvent.org/fathers/3001059.htm>.}

% structure: p0001 source=h1 target=section reason=page-title

\section{第六〇书}

\Emph{致\Person{赫利奥多鲁斯}}\par

\EditorialNote{这是热罗尼莫最优美的书信之一，写给他的老友、现为阿尔蒂努姆主教的赫利奥多鲁斯，安慰他失去侄儿内波提安的悲痛，内波提安不久前死于热病。热罗尼莫试图缓解朋友的悲伤：（1）将异教徒的绝望或顺从与基督徒的希望对比；（2）称颂死者作为人和司铎；（3）回顾当时困扰帝国的邪恶，并主张内波提安已从中被移除。整封信充满深沉而真诚的情感。其年代为公元396年。}\par

1. 才智浅薄者无法驾驭宏大题旨，若不自量力，必在尝试中失败；且题材愈伟大，无力表述其壮丽者便愈被彻底压倒。属于我、属于你、也属于我们所有人——更确切地说，属于基督，且因属于基督而更属于我们——的\Person{内波提安}，已离弃我们这些长辈，以致我们痛悔交加，被难以承受的悲伤压倒。我们原以为他是我们的继承人，如今我们所拥有的只有他的遗体。我的才智将为何人辛劳？我拙劣的文字将渴望取悦何人？他在何处？我那工作的推动者，他的声音比天鹅绝唱更甜美？我心茫然，手在颤抖，双目蒙雾，口舌结巴。无论我说什么，既然他再也听不见，便如同未说。我的笔似乎也感受到他的失去，我的蜡板显得黯淡而哀伤；前者生锈，后者发霉。每当我试图用言语表达，并把这\OriginalTerm{颂词}{encomium}的花朵撒在他的坟上时，我的眼中便充满泪水，悲伤复返，除了他的死，我什么也想不到。古时有习俗，子女在父母遗体前公开赞美死者，并（如同唱悲歌时）使听者眼中流泪、胸中叹息。但在我们这里，看，秩序颠倒了：为给予我们这一打击，自然已丧失其权利。因为年轻人本应对长辈表示的尊敬，如今我们这些长辈却向他表示。\par

2. 那我该怎么办？我要与你们一同流泪吗？宗徒禁止我，因为他称已死的基督徒为\BibleQuote{那些睡了的人}\BibleRef{得前 4:13}。同样，主在福音中说：\BibleQuote{女孩子不是死了，是睡着了}\BibleRef{谷 5:39}；拉匝禄复活时，也被说成是睡着了。\BibleRef{若 11:11}不，我要欢喜快乐，因为\Quote{他早早被接去，免得邪恶改变他的心意}，因为\Quote{他的灵魂使主喜悦}。然而，虽然我不愿屈服，与情感斗争，泪水仍流下双颊，尽管有德行的教导和复活的希望，一股遗憾的激情压倒了我过于柔顺的心。哦，死亡！你分开了因爱而结合在一起的兄弟，你多么残酷、无情，竟如此拆散他们！\Quote{主从旷野吹来一阵热风，枯干了你的脉络，使你的泉源荒凉。}你吞吃了我们的\Person{约纳}，但他仍在你腹中活着；你将他如死人一般带走，好让世界风暴平息，我们的\Place{尼尼微}因他的宣讲而得救。是他，是的，他征服了你，杀死了你——那逃走的先知，离开了他的家，放弃了继承权，将宝贵的生命交在寻找它的人手中。就是他，古时藉欧瑟亚威胁你：\BibleQuote{死亡啊！我要做你的灾难；阴府啊！我要做你的毁灭}\BibleRef{欧 13:14}。藉他的死，你死了；藉他的死，我们活着。你吞吃了，又被吞吃。当你被贪欲吸引，渴望他取的身体，以为是你贪婪口中所掠夺之物时，你的内脏已被钩子咬伤了。\par

3. 救主基督，你所造的我们感谢你，当你被杀时，你击杀了我们强大的仇敌。在你来临之前，有谁比人更悲惨？他们惧怕永死，活着却只为灭亡！因为\BibleQuote{从亚当到梅瑟，死亡为王了，连那些没有像亚当一样犯罪的人，也隶属其权下}\BibleRef{罗 5:14}。如果\Person{亚巴辣罕}、\Person{依撒格}和\Person{雅各伯}都在地狱，谁能进入天国？如果你的朋友们——甚至连那些本身未犯罪的人——都因另一人的罪而受罚，并且触犯了亚当的惩罚，那么对于那些心里说\Quote{没有天主}，自甘\Quote{败坏可憎}，且被称说\Quote{他们偏离了正道，变成无用的；没有行善的，连一个也没有}的人，我们该如何看待？\BibleRef{罗 3:12}即使\Person{拉匝禄}被看见在\Person{亚巴辣罕}怀中，在安息之处，地狱仍无法与天国相比。基督来临前，\Person{亚巴辣罕}在地狱；基督来临后，强盗在乐园里。因此，当他复活时，\BibleQuote{许多长眠圣人的身体复活了，并出现在天耶路撒冷}\BibleRef{玛 27:52}。那时应验了这句话：\BibleQuote{你这睡眠的，醒来吧！从死者中起来，基督必要光照你}\BibleRef{弗 5:14}。\Person{洗者若翰}在旷野呼喊：\BibleQuote{你们悔改吧！因为天国临近了}\BibleRef{玛 3:2}。因为\BibleQuote{从洗者若翰的日子起，天国是强取夺来的，强力的人就夺取了它}\BibleRef{玛 11:12}。那把守乐园路径的旋转火剑和守在门口的\OriginalTerm{革鲁宾}{cherubim}\BibleRef{创 3:24}，都因基督的圣血而熄灭并被解除。我们应许在复活中得到这一点，并不奇怪：因为我们这些活在肉身却不随从肉身的人，\BibleRef{格后 10:3}已拥有天上的国籍，当我们还在世上时，我们就被告知\BibleQuote{天国在我们中间}\BibleRef{路 17:21}。\par

4. 此外，在基督复活前，天主仅在\Quote{犹大为人所知}，\Quote{在以色列中他的名为大}。而那些认识他的人，虽认识他，仍被拖入地狱。那时，从印度到不列颠、从北极冰带到炎热的太洋之滨，地球上的居民在哪里？世上无数的百姓在哪里？大军在哪里？\par

言语不同，衣着和武器不同？\par

他们像鱼和蝗虫、像苍蝇和蚊子一样被压碎。因为除了认识他的造主，每个人不过是一头牲畜。但现在，所有民族的言论和著作都在宣告基督的受难与复活。我不提犹太人、希腊人和罗马人，这些民族曾被主用祂十字架上所写的称号\BibleRef{路 23:38}归于祂的信德。灵魂的不朽及其在肉体解体后的持续——这些毕达哥拉斯梦想过、德谟克利特拒绝相信、苏格拉底在狱中为安慰自己面对判决而讨论的真理——如今已成为印度人、波斯人、哥特人和埃及人熟悉的主题。凶猛的贝西人和一群穿着兽皮的野蛮人，他们曾为纪念死人而献人祭，如今已从他们刺耳的不和谐中爆发出来，进入十字架的甜美旋律，基督成了全世界的唯一呼声。\par

5. 我们能做什么，我的灵魂？我们该转向何处？我们该先拿起什么？我们该跳过什么？你忘记修辞学家的准则了吗？你被悲伤占据，被眼泪淹没，被啜泣阻碍，以致忘记了所有逻辑顺序？你自幼研习的学问在哪里？你一直称赞的\Person{阿那克萨戈拉}和\Person{特拉蒙}的那句话在哪里：\Quote{我知道我生了一个必死的人}？我读了\Person{克兰托}为了缓解悲伤而写的书，\Person{西塞罗}模仿了它。我读过\Person{柏拉图}、\Person{第欧根尼}、\Person{克利托马库斯}、\Person{卡涅阿德斯}、\Person{波西多尼乌斯}的安慰性著作，他们曾以书籍或书信在不同时期努力减轻各种人的悲伤。因此，即使我自己的才智枯竭，也能从他们开启的泉源中重新得到滋润。他们向我们展示了无数的榜样，特别是\Person{伯里克利}和\Person{苏格拉底}的学生\Person{色诺芬}。前者在失去两个儿子后戴上花冠并发表演说；而后者在献祭时听说儿子战死沙场，据说他放下了花冠，但得知儿子英勇倒下的消息后，他又重新戴上。至于那些罗马将领，他们的英雄德性像星星一样闪耀在拉丁史册上，我该说什么呢？\Person{普尔维卢斯}在奉献卡皮托利欧山时接到儿子猝死的消息，他下令葬礼在他不在场的情况下进行。\Person{卢基乌斯·保卢斯}在两个儿子葬礼之间的一周内凯旋进城。我略过\Person{马克西穆}、\Person{加图}、\Person{加卢斯}、\Person{皮索}、\Person{布鲁图}、\Person{斯凯沃拉}、\Person{梅特卢斯}、\Person{斯考鲁斯}、\Person{马略}、\Person{克拉苏}、\Person{马塞卢斯}、\Person{奥菲狄乌斯}这些人，他们在悲伤和战争中表现出同样的坚韧，\Person{图利乌斯}在他的书\WorkTitle{论安慰}中陈述了他们的丧亲之痛。我略过他们，以免显得我宁愿选用别人的言辞和哀伤而非自己的。然而，如果我们的信德未能超越不信的成就，这些例子也足以确保我们感到羞愧。\par

6. 那么让我转向我真正的主题。我不会与雅各伯和达味一同为在律法下死去的儿子捶胸哀悼，但我要在福音中与基督一同迎接他们复活。犹太人的哀悼是基督徒的喜乐。\Quote{哭泣可能会持续一夜，但喜乐在早晨来临。}\BibleQuote{黑夜已深，白昼将近}\BibleRef{罗 13:12}。因此，当梅瑟去世时，人们为他哀悼\BibleRef{申 34:8}，但当若苏厄下葬时，既无泪水也无葬礼的排场\BibleRef{苏 24:30}。所有能从圣经中关于哀悼的内容，我已经在写给罗马的保拉的安慰信中简要陈述了。现在我必须走另一条路到达同一目标。否则我似乎是在一条曾经踏过但久已废弃的小径上重新行走。\par

7. 我们确实知道我们的\Person{尼波提亚努斯}与基督同在，并已加入圣人的歌咏团。他在世上与我们一起遥远地摸索、尽力寻求的一切，如今在近处看见，以至他能说：\Quote{我们曾听说，如今在万军上主的城中，在我们天主的城中，亲眼看见了。} 但我们仍无法忍受失去他的感觉，若不是为他，也是为我们自己悲伤。他所享有的幸福越大，失去如此巨大福佑带给我们的悲伤就越深。\Person{拉匝禄}的姐妹们忍不住为他哭泣，尽管知道他会复活。救主自己——为显示祂拥有真实的人性——也为祂即将复活的人哀哭\BibleRef{若 11:35}。祂的宗徒尽管说：\Quote{我渴望离世与基督同在}，以及别处\Quote{对我来说，活着就是基督，死了就是获益}，却感谢天主将\Person{厄帕夫辣}（他曾\Quote{病重垂危}）归还给他，免得他忧上加忧。这些话并非源于不信的恐惧，而是出于真爱的深切惋惜。你既是\Person{尼波提亚努斯}的叔父又是主教（即肉身和灵魂上的父亲），失去如此亲爱的人，你会多么更深地哀痛，仿佛心被撕碎。我求你为你的悲伤设一个界限，记住那句格言：\Quote{凡事不可过分。} 暂时包扎你的伤口，聆听你一直喜爱其美德之人的赞词。不要因为你失去了如此完美的模范而悲伤：倒要为他曾属于你而喜乐。如同人在小板上描绘大地的轮廓，在这小小的卷轴中，你也能看到他的美德——即使不是完全描绘，至少是勾勒了轮廓。我恳求你接受我的心意，而非实际成就。\par

8.  修辞学家在此类情形下的建议是：你应首先追溯受颂扬者的远祖，历数其功绩，然后逐步接近你的英雄；如此，藉先祖的德能使他更显光辉，并表明他或是有德之人的相称继承者，或是为原本卑微的门第增添了光彩。然而，因我的职责是歌颂灵魂，故不拘泥于那些肉体上的优势，而\Person{乃波提安}生前始终轻视这些。我也不夸耀他的家族——那不属于他，而属于他人；因为连那些圣人\Person{亚巴郎}和\Person{依撒格}都生有罪人\Person{依市玛耳}和\Person{厄撒乌}为子。另一方面，被宗徒列于义人行列\BibleRef{希 11:32}的\Person{依弗大}，却是妓女之子。\BibleRef{民 11:1}经上记载：\BibleQuote{犯罪的人，必遭死亡。}因此，没有犯罪的人必要生存。父母的德行与罪过都不归算于儿女。天主只从我们在基督内重生的那一刻起计算我们。\Person{保禄}，教会的迫害者，早晨是\Person{本雅明}中掠夺的狼\BibleRef{创 49:27}，到了傍晚却\Quote{赐予食物}，即将自己交托给绵羊\Person{阿纳尼雅}。同样，让我们视我们的\Person{乃波提安}为初生的婴孩和未受教的孩子，他在一瞬间从\Place{约旦河}的水中重生，为我们而生。\par

9.  另一个人或许会描述你如何为了他的得救离开了东方和旷野，以及你如何以返回的希望安慰我最亲爱的同伴：这一切都是为了尽可能地拯救你的妹妹——当时她是寡妇，带着一个幼子——或者，若她拒绝你的劝诫，至少拯救你的可爱的小侄子。关于他，我曾预言说：\Quote{即使你的小侄子紧紧搂着你的脖子。}另一个人，我说，会叙述\Person{乃波提安}尚在宫廷服务时，如何在戎装和洁白耀眼的亚麻衣下，皮肤被苦衣磨损；如何在这世界的权贵面前，嘴唇因禁食而变色；如何仍在一主之戎装中服侍另一位主；以及他如何佩带剑带，只为扶助寡妇和孤儿、困苦和不幸之人。至于我，我不喜欢人们拖延对天主的完全奉献。当我读到百夫长\Person{科尔乃略}\EditorialNote{宗徒大事录第十章}是一个义人时，我立即听到了他的圣洗。\par

10. 我们仍可赞许这些事，如同幼小信德的襁褓。一个曾在异旗下忠心服役的士兵，当他来事奉自己的君王时，必定配得桂冠。当\Person{Nepotian}放下佩带、改变装束时，他将所得的全部薪饷施舍给了穷人。因为他读过这话：\BibleQuote{你若愿意成全，变卖你所有的，施舍给穷人，然后来跟随我。}\BibleRef{玛 19:21} 又说：\BibleQuote{你们不能事奉两个主人：天主和玛门。}\BibleRef{玛 6:24} 他除了为遮盖自己、抵御寒冷的一件普通长衣和斗篷外，什么也没有留下。他的衣着按本省款式裁剪，既不特别雅致，也不显得寒酸。他每日急切渴望前往\Place{埃及}的隐修院，或拜访\Place{美索不达米亚}的团体，至少也要独自生活在\Place{达尔马提亚}群岛——这些岛屿仅以\Place{阿尔提努姆}海峡与大陆相隔。但他不忍心撇下他的主教叔父，在叔父身上他看到了许多德行的模范，无需远行便可向他学习。在同一人身上，他既找到了可效法的隐修士，又找到了可敬重的主教。通常发生的事没有发生：长期的亲密并未导致轻慢，轻慢也未生出藐视。他敬重叔父如父亲，每日都为某种新美德而赞叹不已。简言之，他成为了一名圣职人员，历经通常的阶段后，被祝圣为司铎。好\Person{耶稣}！他是何等叹息、呻吟！何等斋戒、躲避众人目光！他生平第一次也是唯一一次向叔父发怒，抱怨加于他的重担过于沉重，他的青春不适合司铎职。但他越推辞，越赢得众人的心；他的推辞恰恰证明了他配得这他不愿接受的职分，正因他自认不配，所以越发配得。我们也有我们的\Person{弟茂德}；我们也看见了那与白发同等的智慧；\BibleRef{智 4:9} 我们的\Person{梅瑟}拣选了一位他确知是长者的长者。\Person{Nepotian}视圣职为负担而非荣誉。他的首要之事是以谦逊平息嫉妒，其次是不给人留下把柄，使那些攻击他青春的人能惊叹他的贞洁。他帮助穷人，探望病人，激励人好客，以温柔的话语安抚人，与喜乐的人同乐，与哭泣的人同哭。\BibleRef{罗 12:15} 他是盲人的杖，饥饿者的食物，沮丧者的希望，丧亲者的安慰。每一种德行在他身上都如此显著，仿佛他只有这一种德行。在同僚司铎中，他在工作上总是最杰出，却永远满足于最低的位置。他所行的任何善事都归于叔父；但若结果不如预期，他便说叔父并不知情，是自己的错误。在公共场合，他视叔父为主教；在家中，他视他为父亲。他性情的严肃因愉快的表情而缓和；他的欢笑是快乐的，却从不放声大笑。基督的贞女和寡妇，他敬如母亲，劝如姐妹，\BibleQuote{保持着完全的纯洁}。\BibleRef{弟前 5:2} 回家后，他常把圣职身份留在门外，全心投入隐修士的严格规矩。他祈祷频繁，守夜不倦，将泪水献给天主，而非世人。他调节斋戒——如同车夫调节马匹的步伐——根据身体的疲劳或精力。在叔父的餐桌旁，他浅尝食物，既避免迷信，又保持节制。宴席交谈时，他习惯从圣经中提出话题，谦逊聆听，迟疑作答，支持正确，驳斥错误，但两者都不带苦毒；宁可教导他的对手，而不打败他。如此纯朴的谦逊装点着他的青春，使他坦承各个论点来自何处；这样，他虽不追求学识之名，却被视为最有学识者。他会说：这是\Person{Tertullian}的意见，那是\Person{Cyprian}的；这是\Person{Lactantius}的，那是\Person{Hilary}的；\Person{Minucius Felix}如此说，\Person{Victorinus}这样讲，\Person{Arnobius}那样说。有时他也会引用我，因为他因我与叔父的亲密而爱我。的确，通过不断的阅读和长久的默想，他已使自己的胸中成为一座基督的图书馆。\par

11. 他多少次从海外来信催我为他写点什么！他多少次提醒我福音中夜间求帮的人\BibleRef{路 11:5}，以及那向残忍法官不断恳求的寡妇！当我默默忽略他的请求，因羞于回答而让求情者脸红时，他推出他的叔父来加强请求，知道既然是为别人求情，叔父会更愿意开口，而且我因尊重他的主教职位，更易获得恩准。于是我便做了他所愿的事，在一篇简短的文章中把我们彼此的友谊献给永恒的纪念。\Person{尼波提安}收到后，夸口说他比\Person{克雷兹}更富有，比\Person{大流士}更富裕。他把它捧在手里，用眼吞食，藏在怀里，挂在嘴边。常常当他躺在床上展开它时，他枕着那宝贵的书页入睡。无论是陌生人还是朋友来访，他都乐于让他们知道我对他的见证。他通过细心的标点和多样的强调，弥补了我小书的不足，以至于当它被朗读时，似乎总是他而不是我在使人愉悦或不满。如此的热忱从何而来，若不是来自对天主的爱？如此不倦地研习基督的律法从何而来，若不是来自对赐律者的渴望？让别人钱上加钱，直到钱袋鼓满；让别人降低身份去依赖已婚妇女；让别人作为隐修士比作为世俗之人更富有；让别人在侍奉贫穷的基督时拥有他们在侍奉富有的魔鬼时从未拥有的财富；让教会因那些在世上曾是乞丐之人的富有而喘不过气来。我们的\Person{尼波提安}鄙视金钱，只求写下的书籍。然而，尽管他在肉体上轻视自己，在贫穷中却行走得更加辉煌，他仍然寻求一切可以装饰教会的东西。\par

12. 与前面所说的相比，我下面要说的可能显得琐碎，但即使在琐事中，同样的精神也显明出来。因为正如我们赞美造物主，不仅因为他是天地、太阳和海洋、大象、骆驼、马、牛、豹、熊和狮子的创造者，而且也是极微小生物的创造者：蚂蚁、蚊蝇、蠕虫等等，我们对它们的形状比对它们名字更熟悉，而在所有这些受造物中，我们都同样敬畏那创造技巧。同样，献给基督的心灵，在小事上表现出与大事情上同样的认真，知道它必须为每一句废话交账\BibleRef{玛 12:36}。\Person{尼波提安}尽心保持祭台明亮，教堂墙壁无烟灰，地面适当清扫。他确保守门人常在其位，门帘挂在适当位置，圣所清洁，器皿闪光。他对每一项礼仪所显示的小心恭敬，使他忽略了任何大小职责。每当你去教堂找他，你总会在那里找到他。\par

在\Person{昆图斯·法比乌斯}身上，古人欣赏一位贵族和罗马史的作家，然而他的绘画比他的写作更出名。我们自己的\Person{贝匝肋耳}\BibleRef{出 31:2}和\Person{希兰}，一个\Place{提尔}妇人的儿子，在圣经中被描述为充满智慧和天主的神，因为他们分别制作了会幕和圣殿的器具。因为，正如肥沃的田地和富饶的耕地有时会过度生长出茎或穗，杰出的才华和充满美德的心灵也必然会溢出优雅多样的成就。因此，在希腊人中，我们听说有一位哲学家曾夸口说，他穿的一切，直到斗篷和戒指，都是自己做的。我们可以对我们的朋友说同样的赞美，因为他用花卉、叶子和葡萄藤的素描装饰了教堂的长廊和殉道者的厅堂，以至于教堂中一切吸引人的东西，无论是因其位置还是外观，都见证了管理它的司铎的辛劳和热忱。\par

13. 你蒙福前行，在善良中！这样的开始之后，我们该期待怎样的结局！唉！必死之人的不幸命运，以及所有不活在基督里的生命的虚空！我的话语啊，你为何退缩？为何躲闪？我害怕来到结局，仿佛我能推迟他的死亡或延长他的生命。\BibleQuote{凡有血肉的都像草，凡人的荣耀都像草上的花。}\BibleRef{伯前 1:24}如今那俊美的面容和庄严的身影，仿佛一件美丽的衣裳披在他可爱的灵魂上，在哪里呢？百合开始枯萎，唉！当南风吹起时，紫色的紫罗兰渐渐褪成苍白。然而，当他因发烧而燃烧，当疾病的火焰干涸他血管的泉源时，他喘着气，疲倦地仍然试图安慰他悲伤的叔父。他的面容闪耀着喜乐，周围的人都哭泣，只有他一个人微笑。他甩开斗篷，伸出手，看见别人看不到的，甚至试图起身迎接新来的人。你会以为他是在开始旅行，而不是在死去，而且他不是离开所有的朋友，只是从一些人走向另一些人。泪水顺着我的脸颊流下，无论我如何坚忍心志，我无法掩饰我感到的悲伤。谁能想到在这样的时刻他会记起与我的亲密，在他挣扎求生时会回想起学习的甜美？然而，他抓住叔父的手对他说：\Quote{请把这件我在服侍基督时所穿的短衣送给我的挚友，我在年龄上的父亲、但在职位上的兄弟，并将你一向给予你侄子的爱转移给那位和你我一样亲爱的人。}说完这些话，他握着叔父的手，口中念着我的名字，离世了。\par

14. 我知道你是何等不愿意以如此代价证明你百姓的挚爱，并且你宁愿在保有幸福的同时赢得同胞的爱。这种情感的表达，在万事顺遂时令人愉快，在悲哀时刻更是加倍的受欢迎。整个阿尔蒂诺、整个意大利都为内波提安哀恸。大地接收了他的身体；他的灵魂已交还给基督。你失去了一位侄子，教会失去了一位司铎。本该跟随你的人却走在了你前面。在你所持有的职务上，按众人判断，他理应继任。因此，一个家庭有幸产生两位主教：第一位值得庆贺，因为他已担任此职；第二位值得哀悼，因为他被过早带走而未能就任。柏拉图认为智者的整个生活应当是默想死亡；哲学家们赞美这一见解并将其推崇备至。但宗徒的话更有力量：\Quote{我因你们的荣耀而天天冒死。}因为拥有理想是一回事，实现它又是另一回事。活着为了死是一回事，死了为了活是另一回事。圣贤和基督徒都必须死：但前者总是因他的荣耀而死，后者却是进入荣耀而死。因此，我们也应当预先思考那终有一天要临到我们、且无论我们愿意与否都不能远去的结局。因为即使我们活到九百岁或更久，如洪水之前的人那样，甚至玛土撒拉的日子也加给我们\BibleRef{创 5:27}，那漫长的时光一旦过去、到了尽头，也等于虚无。因为对活过十年的人和活过千年的人，一旦生命的终结和死亡无情的审判来临，所有过去无论长短都是一样的；除非一个人越年长，他必须带走的罪担也越沉重。\par

首先不幸的凡人从生命中失去最美好的日子；\newline{}疾病和衰老接踵而至；\newline{}最后残酷的死亡索取他的猎物。\par

诗人奈维乌斯也说：\par

凡人必须忍受许多痛苦。\par

据此，古人虚构说尼俄柏因哭泣过多而变成石头，其他妇人则变形为野兽。赫西俄德也哀悼人的生日而欢喜他们的死亡，恩尼乌斯明智地说：\par

众人比国王多一个优势：\newline{}他们可以哭泣，而国王流泪则是耻辱。\par

如果国王不可哭泣，主教也不可哭泣；实际上主教比国王更少特权。因为国王统治不愿服从的臣民，主教统治甘愿服从的臣民。国王以恐怖强迫服从；主教以成为仆人来行使主权。国王看守人的身体直到他们死去；主教拯救人的灵魂得永生。众人的目光都转向你。你的房屋建在瞭望台上；你的生活为他人设定了克己的界限。你无论做什么，众人都以为他们可以照样做。不要如此行事，以致那些寻求挑剔理由的人以为他们理应攻击你，或者那些渴望效法你的人被迫作恶。尽最大努力——甚至超过你的能力——克服你心神的敏感，节制你涌流的眼泪。否则，你对你侄子的深情可能被不信者理解为对天主的绝望。你该惋惜他不是当作死了，而是当作不在眼前。你该显得在寻找他，而非失去了他。\par

15. 然而我为何要试图医治一个我想已经被时间和理性缓解的悲伤呢？我何不向你揭示——它们并不难寻——我们统治者的苦难和我们时代的灾祸呢？那失去生命之光的人并不那么值得怜悯，那逃脱如此巨大灾祸的人才值得庆贺。亚流异端的庇护者君士坦提乌斯在匆忙与敌人交战时，死于莫普苏村，大为懊恼地将帝国留给了他的敌人。出卖自己灵魂的朱利安、杀害基督徒军队的凶手，在米底感受到了他先前在高卢否认的基督的手。渴望为罗马增添新领土的他，反而失去了先前获取的领土。约维安刚刚尝到统治的甜美，便被煤火闷死：这是人类权力短暂的一个好例证。瓦伦提尼安因血管破裂而死，他的出生地荒芜，国家未得复仇。他的兄弟瓦伦斯在色雷斯被哥特人击败，葬身于战死之地。格拉提安被军队出卖，被沿途城市拒绝入城，成为敌人的笑柄；里昂啊，你的城墙仍留有那血腥之手的痕迹。小瓦伦提尼安尚是青年——我可以说，只是个孩子——在逃亡、流放和以流血重获权力之后，在他兄弟丧生的城市不远处被处死。不仅如此，他的尸体还被悬挂示众以羞辱他。我还能说什么呢？普罗科皮乌斯、马克西穆斯、尤金尼乌斯，他们掌握主权时令万国恐惧，却一个个成为他们征服者面前的俘虏，并且——对位高权重者最残酷的创伤——在死于刀剑之前，感受了卑贱奴役的刺痛。\par

16.\par

有人会说：这就是君王的命运；\newline{}闪电总是击打山顶。\par

因此，我要谈到私人地位的人，而谈论他们时，我不会追溯超过最近两年。事实上，我将满足于——省略其他所有人——只叙述最近三位执政官的各自命运。阿邦丹提乌斯在皮提乌斯沦为乞丐般的流亡者。鲁菲努斯的头颅被长矛挑着带到君士坦丁堡，他被砍下的手挨家挨户乞求施舍，以羞辱他那贪得无厌的欲望。提马西乌斯，从最高职位上突然被抛下，认为能在阿萨过着默默无闻的生活是一种逃脱。我正在描述的并非少数不幸者的灾难，而是整个人类命运所依赖的线索。当我想到我们这个时代的灾祸时，我不寒而栗。二十多年来，在君士坦丁堡与尤利安阿尔卑斯山之间，罗马人的鲜血每日流淌。西徐亚、色雷斯、马其顿、达尔达尼亚、达契亚、色萨利、阿哈伊亚、伊庇鲁斯、达尔马提亚、潘诺尼亚诸省——所有这些都被哥特人、萨尔马特人、夸德人、阿兰人、匈奴人、汪达尔人和马尔克人洗劫、掠夺和劫掠。多少天主的已婚妇女和贞女，贤德而高贵的女士，成了这些野蛮人的玩物！主教被俘，司铎和低级圣职人员被杀。教堂被推翻，基督的祭台旁拴着马匹，殉道者的圣髑被挖出。\par

哀伤与恐惧四处弥漫，\newline{}死亡呈现出无数形态。\par

罗马世界正在崩溃：我们却昂首而不低头。你想，科林斯人、雅典人、拉刻代蒙人、阿卡迪亚人，以及任何受蛮族统治的希腊人，现在还有何等勇气？我只提了几个城市，但这些城市曾经是重要城邦的首都。东方似乎确实能免受所有这些灾祸；如果说这里的人们惊慌失措，那只是因为来自其他地区的坏消息。但是看！就在刚过去的一年里，狼群（不再是阿拉伯的，而是整个北方的）从高加索最偏远的堡垒中被释放出来，并在短时间内席卷了这些广大的省份。他们占领了多少隐修院！使多少河流染成了红色！他们围攻了安提约基雅，并包围了哈吕斯河、库德努斯河、奥龙特斯河和幼发拉底河上的其他城市。他们俘虏了大批人。阿拉伯、腓尼基、巴勒斯坦和埃及，因恐惧而幻想自己已经沦为奴隶。\par

纵有百口百舌，\newline{}铁喉铜胸，\newline{}亦不能数尽世人的无尽苦痛。\par

其实，我无意撰写历史；我只想为你的和我的悲伤洒下几滴眼泪。至于其他，若要以应有的方式处理这样的主题，\Person{修昔底德}和\Person{撒路斯提乌斯}也会哑口无言。\par

17. 内波提安是有福的，他既看不见这些事，也听不见这些事。我们是不幸的，因为我们或自己受苦，或看到我们的弟兄受苦。然而我们竟渴望活着，并把那些远离这些灾祸的人视为可怜而非有福。我们早就感到天主在发怒，却不去平息祂。正是我们的罪使蛮族强大，正是我们的恶行击败了罗马的士兵；而且，仿佛这里供屠杀的材料还不够，内战在我们中间造成的毁灭几乎比外敌的刀剑还多。那些以色列人必定是可怜的，与他们相比，拿步高被称为天主的仆人。\BibleRef{耶 27:6}我们也是不幸的，我们如此令天主不悦，以致祂利用蛮族的怒火来执行祂对我们的愤怒。然而，当希则克雅悔改时，一夜之间一个天使就消灭了十八万五千亚述人。\BibleRef{列下 19:35}当约沙法特赞美上主时，上主使祂的敬拜者得胜。\BibleRef{编下 20:5}同样，当梅瑟与阿玛肋克作战时，他得胜不是靠刀剑，而是靠祈祷。\BibleRef{出 17:11}因此，如果我们希望被高举，就必须先俯伏。唉！我们的羞辱和愚昧甚至达到了不信的程度！罗马的军队，曾是世界的征服者和主人，如今一见到敌人就恐惧战栗，并败给那些不能行走、一旦下马就以为死了的人。我们不明白先知的话：\Quote{一人叱责，千人逃遁。}\BibleRef{依 30:17}我们没有切除疾病的根源，正如我们必须做的那样才能除去疾病本身。否则，我们很快就会看到敌人的箭让位于我们的标枪，他们的帽子让位于我们的头盔，他们的坐骑让位于我们的战马。\par

18. 然而，我已超越了安慰者的职责，在禁止你为一个死人哭泣的同时，我自己却为全世界的死者哀悼。伟大的国王薛西斯，他曾削平山岳、填平海洋，从高处望着他面前那数不清的军队，据说想到一百年后他所看见的人没有一人还活着，便哭泣起来。哦！如果我们能登上一个如此高的瞭望塔，从那里可以看见整个大地展开在我们脚下，那么我将向你展示一个世界的废墟，民族与民族交战，王国与王国碰撞；有些人受折磨，有些人被刀剑杀死，有些人被波浪吞没，有些人被掳为奴；这里有婚礼，那里有葬礼；有人出生，有人死亡；有人生活富裕，有人乞讨为生；不仅是薛西斯的大军，尽管那也很庞大，而是世上所有现在活着但注定很快逝去的居民。语言不足以表达如此宏大的主题，我所能说的一切都必然不及现实。\par

19. 那么，让我们返回自身，从天上下来，片刻注视那更贴近我们的事。请问，你是否意识到自己成长的阶段？你能确定你何时成为婴儿、儿童、青年、成年、老人吗？我们每天都在改变，每天都在死去，却自以为永恒。我用于口授、书写、审阅并修改的每一刻，都是从我的生命中夺去的时光。我的秘书所划的每一点，都是从我有限的岁月中消失的。我们写信并回复他人的信，我们的书信跨越海洋，当船只在层层波浪中犁出沟槽时，我们必须活着的时刻逐一消逝。我们唯一的收获是我们在基督的爱中如此紧密相连。\BibleQuote{爱德含忍，慈祥；爱德不嫉妒，不自夸，不傲慢；凡事包容，凡事相信，凡事盼望，凡事忍耐。爱德永不止息。}它永远活在心中，因此我们的\Person{内波提安}虽不在眼前却仍临在，尽管我们相隔遥远，他仍向每人伸出手。是的，在他身上，我们有了彼此爱德的保证。那么，让我们在精神上结合，让我们以情感彼此联结，让我们这些失去儿子的人，展现出真福主教\Person{克罗马提乌斯}失去兄弟时所怀有的同样坚韧。让我们书写的每一页都回响他的名字，让我们所有的书信都充满他的名字。如果我们不能再将他拥抱在怀中，就让我们将他紧握在记忆里；如果我们不能再与他交谈，就让我们永不停止谈论他。\par

\SourceRecord{Translated by W.H. Fremantle, G. Lewis and W.G. Martley.；Nicene and Post-Nicene Fathers, Second Series；Vol. 6.；Edited by Philip Schaff and Henry Wace.；Buffalo, NY: Christian Literature Publishing Co.,；1893.；<http://www.newadvent.org/fathers/3001060.htm>.}

% structure: p0001 source=h1 target=section reason=page-title

\section{第六十一封信}

\Emph{致 \Person{维吉兰修斯}}\par

\EditorialNote{\Person{维吉兰修斯}在访问\Place{耶路撒冷}（他作为\Place{诺拉}的\Person{保利努斯}信件的携带者前往——见\WorkTitle{第五十八封信}，第11节）后返回西方时，曾公开指责\Person{热罗尼莫}倾向于\Person{奥力振}的异端。\Person{热罗尼莫}现在以最严厉的语调写信给他，否认奥力振主义的指控，并将其对手的指控归咎于无知和亵渎。他特别谴责\Person{维吉兰修斯}对\BibleBook{达尼尔}{书}中\Quote{非人手凿出的石头}的解释，并敦促他为自己的罪过忏悔，在这种情况下，他将有和\Person{奥力振}所认为的魔鬼一样多的宽恕机会！这封信常被引用以显示\Person{热罗尼莫}处理\Person{奥力振}著作的方式。\Person{热罗尼莫}随后写了一篇驳斥\Person{维吉兰修斯}著作的文章，这是他所有争论性著作中最激烈和最不合理的。见本卷中其译文。亦见\WorkTitle{第一百零九封信}。此信日期为公元396年。}\par

1. 既然你不肯相信自己耳朵所听见的，我或许可以公正地拒绝以书信满足你；因为，如果你不相信活生生的声音，你也不太可能向书写的纸张让步。但是，既然基督在祂自己身上给我们展示了完美谦卑的典范，亲吻了背叛祂的人，并在十字架上接受了强盗的悔改，我现在在你缺席时告诉你，正如我先前在你面前时已告诉过你的那样：我阅读并且阅读了\Person{奥力振}，只是如同我阅读\Person{阿波利纳尔}或其他作家的作品一样，这些作家的书籍在某些方面教会并不接纳。我绝非说所有这些书中的内容都应被谴责，但我承认其中有值得责备之处。然而，我的任务和研究是通过阅读许多作者，从尽可能多的作品中采集不同的花朵，与其说是以证明一切为目标，不如说是选择那些好的。我拿起许多作家，以便从许多作品中学到许多东西；根据那所写的：\BibleQuote{考查一切，持守那美好的。}因此，我非常惊讶你试图将\Person{奥力振}的教义强加于我，而你对他在许多点上的错误教导至今全然不知。我是异端者吗？那么请问，为什么异端者如此不喜欢我？而你是正统的，你或是违背自己的信念和自己口中的话勉强签字，因而成为一个言行不一者，或是故意签字，因而成为一个异端者？你忽略了\Place{埃及}；你放弃了所有那些你的教派可以自由公开争论的省份；你单单选中了我来攻击，我不仅谴责而且公开谴责所有相反教会的教义。\par

2. \Person{奥力振}是异端者，确实；但这对我有什么可剥夺的呢？我并不否认他在许多点上确实是异端。他在身体的复活上犯了错误，在灵魂的状态上犯了错误，他错误地以为魔鬼可能悔改，并且——比这些更严重的错误——他在关于\BibleBook{依撒意亚}{先知书}的注释中宣称先知提到的色辣芬是天主圣子和圣神。如果我不承认他犯了错误，或者如果我不每天咒骂他的错误，我就成了他过错的同伙。因为当我们接受他著作中好的东西时，我们绝不能绑定自己也要接受其中邪恶的东西。然而，在许多段落中，他很好地解释了圣经，解释了先知书中晦涩的地方，并阐明了极大的奥秘，无论是在旧约还是在新约中。那么，如果我接受了他好的东西，并切除或改变或忽略了他邪恶的东西，我是否应当被认为有罪，因为通过我，那些阅读拉丁文的人接受了他著作中好的部分而对其坏的部分一无所知？如果这是罪行，那么宣信者\Person{依拉略}必须被定罪；因为他将\Person{奥力振}的\WorkTitle{\BibleBook{}{圣咏}释义}和\WorkTitle{\BibleBook{约伯}{书}讲道}从希腊文译成了拉丁文。\Place{韦尔切利}的\Person{欧瑟伯}，一位见证过同样宣信的人，也必须被责备；因为他将那位同名的异端者关于所有圣咏的\WorkTitle{所有圣咏注释}译成了我们的语言，然而却省略了不健全的部分，只翻译了那些有益的部分。我不提\Place{佩塔维}的\Person{维克多利诺}和其他人，他们仅仅跟随并扩展了\Person{奥力振}对圣经的解释。我若这样做，或许显得不太急于为自己辩护，而更像是为自己寻找同犯。我要论到你自己：你为什么保留他关于\BibleBook{约伯}{书}的论著的副本？在这些论著中，当他与魔鬼以及关于星辰和天进行争论时，他说了一些教会不接纳的话。难道只有你，以你那非常聪明的头脑，有权对所有希腊和拉丁作家作出判决，用你的审查官之杖一挥，把一些从我们的图书馆中驱逐出去，又接纳另一些，并随心所欲地宣称我是天主教徒或异端者吗？而我就可以被禁止拒绝错误的东西，并谴责我经常已经谴责过的东西吗？读一读我关于\BibleBook{厄弗所}{书}所写的，读一读我的其他著作，特别是关于\BibleBook{}{训道篇}的注释，你就会清楚地看到，从年轻时候起，我从未被任何人的影响所吓倒而赞同异端的邪说。\par

3. 认识自己的无知并非小益。认识自己的分寸，乃是人的智慧，免得他在魔鬼的煽动下被引走，使全世界都见证他的无能。我想，你一心想抬高自己，在你的家乡夸耀说，我无力回应你的口才，我惧怕你如\Person{克吕西波}般的尖锐讽刺。基督徒的谦逊使我克制，我不愿以尖刻的言辞暴露我那可怜斗室的隐居。否则，我很快就能揭露你所有的勇敢和凯旋游行。但我把这些留给别人去谈论或嘲笑；而我本人，作为一个基督徒对另一个基督徒说话，我恳求你，我的兄弟，不要假装你知道得比实际更多，免得你的笔会宣告你的天真和单纯——至少我对此缄默不语的那些品质，虽然你自己看不见，别人却在你身上看见。因为那时你将给所有人讥笑你愚昧的理由。从你最早的童年起，你就接受了其他的教导，习惯了不同种类的学校教育。同一个人很难既是柜台上金币的检验者，又是圣经的检验者；或既是酒类鉴赏家，又是阐释先知或宗徒的专家。至于我，你将我肢解，我们可敬的兄弟\Person{奥克阿诺斯}你指控为异端，你厌恶司铎\Person{味增爵}和\Person{保利尼安}的判断，我们的兄弟\Person{欧瑟比}也不讨你喜欢。只有你是我们的\Person{加图}，罗马民族中最雄辩的人，你希望我们接受你所说的话，如同智慧本身的话语。请回想那一天，当我讲道论及复活及复活身体的真实性时，你跳到我旁边，拍手跺脚，为我的正统欢呼。但现在，自从你开始航海旅行，舱底水的臭气影响了你的头脑，你只把我当作异端者来记念。我能为你做什么呢？我相信了可敬的司铎\Person{保利诺}的书信，并未想到他对你的判断可能有误。而且，当你把信交给我时，我立刻注意到你言辞中的某种不连贯，但我以为这是由于你缺乏文化和知识，而非头脑混乱。我并不责备那位可敬的作者，他写信给我时无疑更愿意隐瞒他所知道的，而不愿在他的信中指控一个既受他庇护又托付他送信的人；但我责怪自己，因为我依赖他人的判断而非自己的判断，而且，当我的眼睛看到一件事时，我却根据一张纸片上的证据相信了与所见不同的事。\par

4. 因此，不要再以你的卷轴烦扰我、压垮我。至少节省你的钱财，你雇用秘书和抄写员，让同一个人既为你抄写又为你喝彩。也许他们的赞美是因为他们从为你抄写中获利。如果你要锻炼心智，就把自己交给文法修辞教师，学习逻辑，在哲学家的学校中受教；等你学会了这一切，你或许会开始沉默。然而，我这样做是愚蠢的，竟为一位能教导众人的人寻找教师，并试图限制一个不知如何说话却又无法缄默的人的言词。那句古希腊谚语千真万确：\Quote{对驴子而言，竖琴毫无用处。}就我而言，我甚至认为你的名字也是矛盾的。因为你的整个心智都在沉睡，你甚至打鼾，你所陷入的睡眠——不如说是昏睡——是如此深沉。事实上，在你以亵渎的嘴唇说出的其他亵渎话中，你竟敢说\BibleBook{达尼尔}{书}\BibleRef{达 2:34}中的那座山，从那里有一块石头非人手所凿而出，就是魔鬼，而那块石头是\Person{基督}，祂从\Person{亚当}（其罪恶此前已使他与魔鬼相连）取了肉身，由童贞女所生，为要将人类从那座山，即魔鬼那里分离出来。你的舌头应当被割下，撕成碎片。有哪个真正的基督徒会把这形象解释为魔鬼，而不是将其归于全能的天主圣父，或以如此可怕的罪恶玷污全世界的耳朵？如果你的解释曾被任何人——我不说天主教徒，而是异端者或外教人——接受过，那就让你的话被视为虔诚。但另一方面，如果\Person{基督}的教会从未听说过这样的不敬，而你是第一个通过其口，使那曾说过\BibleQuote{我要像至高者一样}\BibleRef{依 14:14}的那一位宣告祂就是\Person{达尼尔}所说的那座山的人，那么你就悔改，披上麻衣和灰，以急流的眼泪洗去你可怕的罪责；但愿这桩不敬得以赦免，并且，假设\Person{奥力振}的异端是真的，愿你在魔鬼本人获得赦免时也获得宽恕——这魔鬼从未被定罪过比通过你所说的更大的亵渎话。你对我本人的侮辱，我耐心忍受；你对天主的不敬，我不能忍受。因此，在这封信的后半部分，我可能显得比我一开始所说的要更尖刻一些。然而，既然你曾一度悔改并向我请求宽恕，那么你又犯下一桩需要重新做补赎的罪，是极其愚蠢的。愿\Person{基督}赐你恩宠，让你倾听并缄默，让你明了，然后才发言。\par

\SourceRecord{Translated by W.H. Fremantle, G. Lewis and W.G. Martley.；Nicene and Post-Nicene Fathers, Second Series；Vol. 6.；Edited by Philip Schaff and Henry Wace.；Buffalo, NY: Christian Literature Publishing Co.,；1893.；<http://www.newadvent.org/fathers/3001061.htm>.}

% structure: p0001 source=h1 target=section reason=page-title

\section{\newline{}第六十二封书信}

\newline{}\Emph{致\Person{特兰奎利努斯}}\par

\newline{}\EditorialNote{特兰奎利努斯是热罗尼莫的一位罗马朋友，他曾写信（1）告诉热罗尼莫欧坎努斯在罗马反奥利振派的立场，以及（2）询问奥利振的著作是否有某些部分可以安全有益地研读。热罗尼莫对欧坎努斯的消息表示欢迎，并肯定地回答了特兰奎利努斯的问题。他将奥利振与戴尔都良、阿波利纳里等人并列，这些人的著作虽然有异端观点，但仍被继续阅读。写于公元396或397年。}\par

\newline{}1．尽管我以前对此有所怀疑，但现已证实，联系精神与精神的纽带比任何物质联系都要坚固。因为你，我尊敬的朋友，以你全部的灵魂紧贴着我，而我也因基督的爱与你联合。我坦率而真诚地向你无瑕的心灵说：你写信的纸张，你所写下的文字——尽管无声——却让我感受到你的深情。\par

\newline{}2．你告诉我，有许多人被奥利振的错误教导所迷惑，而那位圣洁的人，我儿欧坎努斯，正与他们的疯狂搏斗。我痛心于那些淳朴的人失去了平衡，但也欣喜于像欧坎努斯这样博学的人正尽力使他们回归正轨。再者，你向我这个微不足道的人询问意见，关于阅读奥利振著作的适宜性。你是说，我们是否应与我们的兄弟福斯提努斯一起完全拒绝他，还是如其他人所说，部分地阅读他？我的意见是，我们有时可以为了学识而阅读他，正如我们阅读戴尔都良、诺瓦图斯、阿尔诺比乌、阿波利纳里以及其他一些希腊和拉丁教会作家的作品一样，并且应当根据宗徒的话来选择其著作中的善而避免恶：\BibleQuote{凡事考验，好的，要持守。}\BibleRef{得前 5:21}然而，那些因性情乖僻而对他过分偏爱或过分厌恶的人，在我看来，正应验了先知的话：\BibleQuote{祸哉，那些称恶为善，称善为恶，以苦为甜，以甜为苦的人！}\BibleRef{依 5:20}因为，他的教导能力不应使我们接受他的错误观点，但他的错误观点也不应使我们完全拒绝他关于圣经所作的有益注释。但如果他的崇拜者和贬损者执意彼此较劲，并且不讲中庸、不守节制，非要一概赞成或反对他的著作，那么我宁愿做一个虔诚的粗人，也不做一个有学问的亵渎者。我们可敬的兄弟，执事塔提安，衷心问候你。\par

\SourceRecord{Translated by W.H. Fremantle, G. Lewis and W.G. Martley.；Nicene and Post-Nicene Fathers, Second Series；Vol. 6.；Edited by Philip Schaff and Henry Wace.；Buffalo, NY: Christian Literature Publishing Co.,；1893.；<http://www.newadvent.org/fathers/3001062.htm>.}

% structure: p0001 source=h1 target=section reason=page-title

\section{\WorkTitle{第六十三封书信}}

\Emph{致\Person{德斐斐洛}}\par

\EditorialNote{当\Person{耶柔米}与\Person{厄彼法尼}一方和\Person{鲁菲努}与\Person{若翰·耶路撒冷}另一方之间发生争论时（参见第五十一封书信），\Place{亚历山大}的主教\Person{德斐斐洛}受到后一方的呼吁，派遣了司铎\Person{依西多尔}向他报告此事。\Person{依西多尔}的报告不利于\Person{耶柔米}，因此\Person{德斐斐洛}拒绝回复\Person{耶柔米}的几封信。最后他写信劝告\Person{耶柔米}遵守教会的教规。\Person{耶柔米}在回信中表示，遵守教规一直是他首要的目标。然后他劝诫\Person{德斐斐洛}对奥利振主义者过于宽容，并宣称这会产生最坏的结果。这封信的日期大概是公元397年。}\par

\Person{耶柔米}致最蒙祝福的教宗\Person{德斐斐洛}。\par

1. 您的圣德会记得，在您对我保持沉默的那段时间，我从未停止尽我写信的职责，与其说考虑到您凭智慧所做的事，不如说考虑到我应当做的事。如今我收到了您的恩宠的来信，我看出我对福音的阅读并非没有果实。因为如果一个妇人的多次祈祷改变了一位不妥协的法官的决心\BibleRef{路 18:2}，那么我的不断恳求岂不更会软化您如慈父般的心吗？\par

2. 我感谢您对教会教规的提醒。的确，\BibleQuote{主所爱的，他必管教，又鞭打凡所收纳的儿子。}\BibleRef{希 12:6}然而我向您保证，没有什么比维护基督的权利、遵循教父们所划定的路线、并永远牢记\Place{罗马}的信德更符合我的目标；那信德是宗徒所赞美的\BibleRef{罗 1:8}，\Place{亚历山大}的教会也自夸是分享者。\par

3. 许多虔诚的人对您在对待那骇人的异端上的长期容忍感到不悦，并认为您以为能凭这样的温和来改正那些攻击教会要害的人。他们相信，当您等待少数人的悔改时，您的行动却在助长被遗弃之人的胆量，并使他们的党派壮大。在基督内再见。\par

\SourceRecord{Translated by W.H. Fremantle, G. Lewis and W.G. Martley.；Nicene and Post-Nicene Fathers, Second Series；Vol. 6.；Edited by Philip Schaff and Henry Wace.；Buffalo, NY: Christian Literature Publishing Co.,；1893.；<http://www.newadvent.org/fathers/3001063.htm>.}

% structure: p0001 source=h1 target=section reason=page-title

\section{第六十四封信}

\Emph{致\Person{法比奥拉}}\par

\EditorialNote{法比奥拉对\Place{白冷}的访问因匈奴入侵的威胁而缩短，这迫使\Person{热罗尼莫}和他的朋友们暂时避难于巴勒斯坦海岸。法比奥拉在此与同伴告别，启航前往意大利，但直到\Person{热罗尼莫}完成了这封供她使用的信函（§22）。信函中包含对大司祭礼服的神秘阐述，体现了\Person{热罗尼莫}一贯的机敏与学识。类似论著被认为出自\Person{戴尔都良}和\Person{科尔多瓦主教奥西乌}，但已失传。其年代为公元396或397年。}\par

\SourceRecord{Translated by W.H. Fremantle, G. Lewis and W.G. Martley.；Nicene and Post-Nicene Fathers, Second Series；Vol. 6.；Edited by Philip Schaff and Henry Wace.；Buffalo, NY: Christian Literature Publishing Co.,；1893.；<http://www.newadvent.org/fathers/3001064.htm>.}

% structure: p0001 source=h1 target=section reason=page-title

\section{第六十五封信}

\Emph{致\Person{普林奇皮亚}}\par

\EditorialNote{这是一篇对《圣咏集》第四十五篇的注释，致\Person{玛尔塞拉}的朋友和同伴\Person{普林奇皮亚}（见第一二七封信）。\Person{热罗尼莫}在开始他的论述之前，先为自己为妇女写作的做法辩护，这种做法使他受到许多愚蠢的嘲笑。他在\Person{索弗罗尼乌斯}注释的献词中也讨论了同一主题。这封信的日期是公元397年。}\par

\SourceRecord{Translated by W.H. Fremantle, G. Lewis and W.G. Martley.；Nicene and Post-Nicene Fathers, Second Series；Vol. 6.；Edited by Philip Schaff and Henry Wace.；Buffalo, NY: Christian Literature Publishing Co.,；1893.；<http://www.newadvent.org/fathers/3001065.htm>.}

% structure: p0001 source=h1 target=section reason=page-title

\section{书信第六十六}

\Emph{致帕玛基乌斯}\par

\EditorialNote{帕玛基乌斯，一位罗马元老，失去了他的妻子保利纳，她是保拉的一个女儿，当时她尚在青春之花绽放之际。直到两年后，热罗尼莫才敢写信给他；他这样做时，对保利纳的生平和美德着墨不多。或许是因为可说的很少。信的大部分内容是对帕玛基乌斯本人的称赞，尽管他身居高位，却成了隐修士，过着严苛的克己生活。热罗尼莫赞许他与法比奥拉在波尔图共同建立的旅客收容所，并描述了自己在伯利恒类似的机构。他还提到了保拉、欧斯托基乌姆和已故的布拉埃西拉，均以最高的赞誉。信的日期是公元397年。}\par

1. 假如一处伤口已经愈合，皮肤上已结疤，那么任何旨在去除疤痕的治疗，在试图改善外观时，必定会重新引发原伤的刺痛。经过两年不合时宜的沉默，我的慰问现在来得太迟了；然而即便如此，我仍担心此刻的言语可能更加不合时宜。我害怕触动您心中的伤痛之处，我的话会再次点燃那已被时间和反省治愈的伤口。因为谁会有如此迟钝的耳朵或如此铁石般的心肠，以至于听到您的保利纳的名字而不哭泣？即使是喝希尔卡尼亚母虎的奶长大的人，他们也必然流泪。谁能以干涸的眼睛看着一朵初绽的玫瑰、一朵未发育的花蕾被如此过早地砍伐和枯萎，它尚未形成杯状，也未展开其深红花瓣的骄傲展示？在她身上，一颗无价的珍珠破碎了。在她身上，一块鲜亮的翡翠碎裂了。只有疾病才让我们知道健康的幸福。我们只有失去时才更意识到我们曾经拥有。\par

2. 我们在比喻中读到的那块好地结出果实，有的一百倍，有的六十倍，有的三十倍。\BibleRef{玛 13:8} 在这三种收成中，我认出了一种象征，即基督的不同奖赏，落到了三个血缘和道德美德紧密相连的女人身上。欧斯托基乌姆采摘童贞的花朵。保拉清扫辛劳的寡居打谷场。保利纳保持婚姻的床榻无玷。一位有这样的母亲，在地上就赢得了基督应许在天上赐予的一切。然后，为了完备这所谓的\Quote{四圣}——由同一家庭产生——并以男人的美德匹配女人的美德，三个女人有了一个合适的同伴：帕玛基乌斯，他像厄则克耳描述的革鲁宾，\BibleRef{则 10:8} 是第一个的姻兄，第二个的女婿，第三个的丈夫。我说丈夫吗？不，不如说是一位最忠诚的兄弟；因为婚姻的语言不足以描述圣神的神圣纽带。这四马战车由耶稣掌握缰绳，正是这样的骏马，哈巴谷歌唱道：\Quote{骑上你的马，让你的骑行成为救恩。}怀着同样的决心，尽管速度不同，他们竭力追逐胜利者的棕榈枝。他们的颜色不同，目标却一致。他们同负一轭，服从同一个御者，不等鞭子催促，而是以新的努力回应他声音的呼唤。\par

3. 让我暂时使用哲学的语言。按照斯多葛派的看法，有四种美德如此紧密相连且彼此一致，以至于缺少一种便缺少全部。它们是明智、正义、勇敢和节制。\BibleRef{智 8:7} 尽管你们所有人都拥有这四种，但每人各以一种出类拔萃。您拥有明智，您的母亲拥有正义，您的童贞姐妹拥有勇敢，您的已婚妻子拥有节制。我说您是明智的，因为谁能比蔑视世俗的愚妄、追随\Quote{天主的德能和天主的智慧}基督的人更明智呢？\BibleRef{格前 1:24} 或者，有什么比您的母亲更好的正义实例呢？她将自己的财产分给子女，以自己对财富的藐视教导他们应当将情感寄托于真正的目标。谁比欧斯托基乌姆立下了更好的勇敢榜样？她决心成为童贞女，从而冲破了贵族的大门，推翻了执政官家族的骄傲。作为罗马贵妇中的第一人，她使罗马的第一家族屈服。有比保利纳更大的节制吗？她读了宗徒的话：\Quote{婚姻应受尊重，床第应是无玷的}，\BibleRef{希 13:4} 不敢奢望她童贞姐妹的幸福或寡母的节制，宁愿选择守住较低路径的安全轨道，而不是踏着空气迷失在云中。一旦她进入婚姻状态，她日夜所思就是，一旦她的结合蒙福得子，从此她将活在第二个等级的贞洁中，并且\par

虽是女人，在这崇高事业中领先，\par

她会引导她的丈夫走上同样的道路。她不离开他，而是期待他成为救恩中的同伴。由于几次小产发现自己的子宫并非不育，她不能完全放弃生育的希望，不得不让自己的不情愿屈服于婆婆的急切和丈夫的烦恼。因此她遭受了许多苦，如同辣黑耳所受的，\BibleRef{创 35:16} 尽管她没有像她那样生下一个痛苦之子、右手之子，她所渴望的继承者正是她的丈夫。我据可靠消息得知，她顺服于丈夫的意愿，并非要利用天主的原始命令\Quote{你们要忠信，要生育繁殖，充满大地}，\BibleRef{创 1:28} 而是她只渴望有子女，以便为基督生养童贞女。\par

4. 我们读到，司铎丕尼哈斯的妻子，一听见上主的约柜被掳去，突然临产阵痛，生下了一个儿子，起名叫依加波特，并在接生婆手中死去。辣黑耳的儿子被称为本雅明，即\InnerQuote{优异之子}或\InnerQuote{右手之子}；但另一位（后来成为天主的一位杰出司铎）的儿子却从约柜得名。同样的事在我们今日也发生了：因为保利娜安息后，教会产下了隐修士帕玛基乌斯，他出身望族，联姻显贵，富有施舍，谦卑自下。宗徒致格林多人书信说：\Quote{弟兄们！你们看看你们是怎样蒙召的：按肉眼来看，你们中有智慧的人并不多，有权力的人也不多，有显贵身世的人也不多。}\BibleRef{格前 1:26}初生的教会需要如此，好使芥菜籽渐渐长成大树，\BibleRef{玛 13:31}并使福音的酵母逐渐发起来，使整个教会面团发酵。\BibleRef{玛 13:33}在我们今日，罗马拥有昔日世界所不知的事物。那时，智慧人、有权势者或贵族中很少有基督徒；如今，许多智慧人、有权势者和贵族不仅是基督徒，甚至是隐修士。在他们中间，我的帕玛基乌斯是最智慧、最有权势、最高贵的；他在大人物中为最大，在领袖中为领袖，他是所有隐修士的总司令。他和类似他的人，就是保利娜在世时所渴望的、在她死后赐予我们的后裔。\Quote{不生育的不妊者，欢乐吧！未经产痛者，欢呼高唱吧！}\BibleRef{依 54:1}因为你顷刻间生出了有如罗马穷人数目的儿子。\par

5. 昔日装饰保利娜颈项和面容的闪闪发光的宝石，如今为穷人购买食物。她的丝绸衣裳和金线锦缎换成了柔软的羊毛衣服，旨在御寒，而非暴露身体以招致虚浮的赞美。从前服务于奢华的一切，如今都服务于德行。那个伸着手、常无人听见而大声呼喊的盲人，是保利娜的继承人，是与帕玛基乌斯一同继承的人。那个几乎拖不动自己的可怜跛子，仰赖一位温柔女孩的帮助得以维生。那些从前涌出成群访客的大门，如今只有可怜的人围聚。一人患水肿，大腹便便，濒临死亡；另一人又聋又哑，无法乞讨，却因不能乞讨而更令人同情；另一人自幼残疾，乞求施舍，却可能自己无法享受。还有一人因黄疸而肢体腐烂，身体已如尸体却还活着。用维吉尔的诗句来说：\par

即使我有百口、百舌，\newline{}我也无法道尽人们无数的苦难。\par

这就是帕玛基乌斯无论走到哪里都伴随的护卫；在这些人的身上，他服侍基督自己；他们的肮脏反使他的灵魂洁白。如此，他向着天国急速前进，如竞赛的施主般仁慈，如给穷人的表演提供者般善良。别的丈夫在妻子的坟墓上撒布紫罗兰、玫瑰、百合和紫色的花朵，以完成这项柔情义务来抚慰心中的悲伤。我们亲爱的帕玛基乌斯也浇灌保利娜的圣灰和敬爱的遗骨，但用的是施舍的香膏。这些就是他珍爱亡妻余烬的香料和香水，因为他知道经上记载：\Quote{水能熄灭烈火，施舍能补赎罪过。}\BibleRef{德 3:30}怜悯有多大能力，它将赢得何等崇高的赏报，真福西彼廉在一部宏大的著作中加以阐述。但以理也证明这一点，他劝告最不虔敬的君王——若君王愿意听从——藉着怜悯穷人而得救。\BibleRef{达 4:27}保利娜的母亲完全可以为保利娜的继承人而高兴。当她看到女儿的财富仍用于她所关心的事情时，她不会后悔财富已转入他人之手。相反，她必定庆幸，自己毫不费力，愿望就得以实现。可供分配的总数仍和从前一样：只是分配者换了。\par

6. 谁能相信，一位是弗里家族的光荣、祖父和曾祖父都是执政官的人，竟然身穿暗色衣服在紫袍的元老中穿行，而且，不但不因遇见同伴的目光而脸红，反而嘲笑那些嘲笑他的人！\Quote{有一种羞耻带来死亡，有一种羞耻带来生命。}隐修士的首要德行就是轻视人言，并常记住宗徒的话：\Quote{如果我还求人的欢心，我就不是基督的仆人了。}\BibleRef{迦 1:10}同样，主对先知说，他使他们的脸像铜城、金钢石和铁柱，好使他们不惧怕人民的侮辱，反而用严厉的目光使嘲笑者无地自容。一颗紧绷的心，更容易为凌辱所胜，而非为恐怖所胜。那些任何酷刑都不能吓倒的人，有时却被怕羞耻的心所征服。的确，一个有出身、有口才、有财富的人，在大街上避开有权势者的交往，混在人群中，贴近穷人，与无知者平等交往，不再作领袖而成为民众的一员，这绝非小事。他越谦卑，就越被高举。\BibleRef{路 14:11}\par

7. 一颗珍珠会在污秽中闪耀，一颗上等宝石会在泥沼中发光。这正是主所应许的，祂说：\BibleQuote{尊重我的，我必尊重他。} \BibleRef{撒上 2:30} 其他人可能将此理解为将来，那时忧愁变为喜乐，世界过去，圣人们将领受永不朽坏的花冠。但就我而言，我看到给圣人的应许甚至在此生就已实现。在他全心服事基督之前，\Person{帕玛丘}已是元老院中的知名人物。然而还有许多其他元老佩戴着总督等级的徽章。整个世界充满了类似的装饰。他确实位居第一等，但除了他还有别人。当他在一些人之上时，另一些人在他之上。最显赫的特权若赐予众人便失去威望，荣誉本身在好人眼中也会因不配得的人拥有而减色。因此\Person{图利乌斯}精妙地评论\Person{凯撒}，当他想要提拔一些追随者时，\Quote{他与其说荣耀他们，不如说羞辱了他所安置他们的荣耀职位。} 如今基督的所有教会都在谈论\Person{帕玛丘}。全世界都赞赏他，视他为穷人，而此前却因富有而忽视他。有什么比执政官更辉煌呢？然而这荣誉只持续一年，当另一个人接任后，前任便让位。每个人的桂冠都消失在人群中，有时凯旋本身也因庆祝者的缺陷而受损。这个职位一度由贵族传给贵族，只有出身高贵者才能担任；执政官\Person{马略}——尽管他是战胜努米底亚、条顿人和辛布里人的胜利者——因家世卑微而被认为不配；\Person{西庇阿}则因勇武而提前赢得此职位——如今这伟大职位只需从军便可获得；胜利的闪亮袍子现在包裹着不久前还是乡巴佬的人。因此我们所得的比所给的更多。我们所舍弃的微小，所拥有的伟大。基督所应许的一切都如期实现，我们因所放弃的而领受了百倍。\BibleRef{玛 19:29} 这就是依撒格撒种的土地\BibleRef{创 26:12}，依撒格，他因随时准备赴死\EditorialNote{创世纪第二十二章}而在福音来到之前就背负了福音的十字架。\par

8. \BibleQuote{你若愿意成为成全的，}主说，\BibleQuote{就去变卖你所有的，分给穷人....然后来跟随我。} \BibleRef{玛 19:21} 你若愿意成为成全的。伟大的事业总是留给听见之人的自由选择。因此宗徒避免把贞洁定为必行义务，因为主谈到那些为天国自阉的人时最后说：\BibleQuote{能领受的，就领受吧。} \BibleRef{玛 19:12} 因为，如宗徒所说：\BibleQuote{不在乎那愿意的，也不在乎那奔跑的，只在乎那施恩的天主。} \BibleRef{罗 9:16} 你若愿意成为成全的。没有强迫加于你：若要赢得奖赏，必须凭你的自由意志。因此，如果你愿意成为成全的，并渴望像先知、像宗徒、像基督自己，就不要卖掉你财产的一部分（免得对匮乏的恐惧成为不忠的契机，而你和\Person{阿纳尼雅}及\Person{撒斐辣}一同丧亡\EditorialNote{宗徒大事录第五章}），而要卖掉你所有的一切。当你卖掉一切后，不要将所得给予富人或者高傲的人，而要给予穷人。给每人足够当下所需，但不要给钱去增加他已有的财富。\BibleQuote{牛在场上踹谷的时候，不可笼住牠的嘴。} \BibleRef{格前 9:9} 并且\BibleQuote{工人配得到他的报酬。} \BibleRef{弟前 5:18} 还有\BibleQuote{那些在祭坛上服务的，分享祭坛上的物品。} \BibleRef{格前 9:13} 也要记住这些话：\BibleQuote{只要有衣有食，就当知足。} \BibleRef{弟前 6:8} 凡你见到冒着热气的美味、煮熟的野鸡、沉重的银盘、烈性的骏马、长发男仆、昂贵的衣物、绣花挂毯，在那里什么都不要给。因为你要给的人比你这位施予者还富足。而且，把属于穷人的东西给那些不穷的人，也是一种亵渎。然而，要成为一个成全的、完备的基督徒，仅仅轻视财富，或者挥霍扔掉自己的金钱（那东西可以一瞬间失去又找到）是不够的。底比斯的\Person{克拉特斯}这样做过，\Person{安提斯泰尼}和其他一些人也这样做过，他们的生活显示出他们有许多缺点。为了获得属灵的荣耀，基督的门徒必须比世俗的哲学家，比那追逐谣言和民众呼声的卑鄙奴仆做得更多。你轻视财富还不够，除非你也跟随基督。只有那离弃自己的罪并携手德行的人，才是跟随基督。我们知道基督是智慧。祂是那人在田地中发现的宝藏。\BibleRef{玛 13:44} 祂是那颗用买卖许多珍珠买来的无价宝石。\BibleRef{玛 13:45} 但如果你爱一个被掳的女子，即世俗的智慧，如果没有别的美貌吸引你，就要剃光她的头发，剪去她那诱人的发丝，也就是文采，修剪她的死指甲。\BibleRef{申 21:11} 用先知所说的碱水\BibleRef{耶 2:22} 洗净她，然后与她同享安逸，说：\BibleQuote{她的左手在我头下，她的右手环绕我。} 然后那被掳的女子会给你带来许多儿女；从摩阿布女子，她将成为以色列女子。基督是那圣化，没有它没有人能见天主的圣容。基督是我们的救赎，因为祂同时是我们的救赎者和赎价。基督是万有，以至于为基督舍弃一切的人可以找到一位代替万有，并能自由地宣告：\BibleQuote{主是我的产业。}\par

9. 我清楚看出你对神圣学问怀有热烈的爱慕，你远不像某些轻率之人那样，试图教导自己一无所知的事物，而是把学习你将要教导的内容作为首要目标。你的书信纯朴自然，散发着先知的气息，洋溢着宗徒的意味。你不追求做作的雄辩，也不像孩童般以工整的对句平衡浅薄的句子。迅速泛起的泡沫也同样迅速地消失；肿胀虽增大身体的体积，却有害健康。况且，加图这句格言十分精辟：\Quote{\OriginalTerm{足够快则足够好}{Fast enough if well enough}}。诚然，很久以前，在我们年轻的日子里，当那位圆满的演说家在开场白中使用这句格言时，我们曾对此放声大笑。我想你还记得我们雅典讲堂中那位讲者所犯的错误，以及整个房间如何因学生们齐声喊出的呼声而回荡：\Quote{\InnerQuote{足够快则足够好}}。据法比乌斯所言，如果只有工匠才能评判工匠，那么技艺必会繁荣。没有人能充分评价一位诗人，除非他自己也能写诗。没有人能理解哲学家，除非他熟悉他们所持的各种理论。物质可见的产品最能由制造它们的人来评估。我们文人面临着多么残酷的命运，你可以从以下事实看出：我们不得不依赖公众的判断；而许多人在众人面前是令人生畏的对手，但若单独看他，他一定会被轻视。我顺带提及此点，是为了若有可能，使你满足于有学识之人的耳朵。不要理会无知之人对你能力的评论；但日复一日地汲取先知的精髓，好使你认识基督的奥秘，并与圣祖们分享这奥秘。\par

10. 无论你读或写，无论你醒或睡，愿亚毛斯的牧人号角常在你耳中回响。愿喇叭声唤醒你的灵魂，愿天主的爱将你带出自我；然后你在床上寻找你灵魂所爱的那一位，\BibleRef{歌 3:1}，并勇敢地说：\Quote{我身睡卧，我心却醒。} \BibleRef{歌 5:2} 当你找到了他，并抓住了他，不要让他走。若你片刻沉睡，他从你手中逃脱，不要立即绝望。出去到街上，吩咐耶路撒冷的众女子；那时你必找到他，正午时分因激情而困倦沉醉，或在同伴的羊群旁被夜露沾湿，或被多种香料熏香，在园中的苹果树间。在那里将你的双乳给他，让他吮吸你博学的胸怀，让他安息在他的产业中，他的羽毛如镀银的鸽子，他的内里闪耀着金光。这个婴孩，这个小儿，以奶油和蜂蜜为食，\BibleRef{依 7:14}，并在崎岖的山间长大，很快长成壮士，迅速掠取你内一切反对他的，时机成熟时劫掠（属灵的）大马士革，并将（属灵的）亚述王锁上锁链。\par

11. 我听说你在波尔图建造了一座接待外客的旅舍，并在奥索尼亚海岸上栽种了一棵亚巴郎树上的枝条。你像埃涅阿斯一样，正在勾勒一个新营地的轮廓；只是，当他抵达台伯河畔时，因急需而不得不吃那形成神谕所指的桌子的方形扁饼，而你却能建造一座面包之屋，堪比我所在的白冷小村；并且，在这里，经过长久的匮乏，你打算以突然的丰盛来满足旅客。做得好！你超越了我卑微的开端。你达到了最高点。你从树根走到了树顶。你是世界首城中隐修士之首；因此，你追随圣祖之首是正确的。让罗特（其名意为\OriginalTerm{转离者}{one who turns aside}）选择平原，\BibleRef{创 13:5}，让他走著名毕达哥拉斯字母表中左边容易的分岔路。但你要像撒辣那样，\BibleRef{创 23:19}，在陡峭多石的高处为自己预备一座纪念碑。让书籍之城近在咫尺；当你摧毁了巨人，阿纳克的后裔，\BibleRef{苏 15:14}，就将你的产业转变成喜乐和欢愉。亚巴郎在金银、牲畜、财物和衣物上都很富足；他的家室如此庞大，以致在紧急情况下他能带领一支精选的年轻人队伍上战场，能追赶至丹，并击杀那已使五王溃败的四王。\BibleRef{创 14:13} 他频繁地施行款待，从不将任何人拒之门外，最终堪当亲身接待天主。他不满足于吩咐仆人和婢女侍奉客人，也不因让别人照料而减少他所施的恩惠；反而像得了宝贝一样，他和妻子撒辣亲自投身于款待的义务。他亲手给客人洗脚，亲自从牛群中背回一只肥牛犊。客人用餐时，他站在一旁伺候，并将撒辣亲手烹制的菜肴摆在面前——尽管他自己打算禁食。\par

12. 我对你所怀的敬意，亲爱的弟兄，使我想起这些事；因为你不但要向基督献上你的钱财，也要献上你自己，作为\BibleQuote{活的祭献，圣洁的，中悦天主的，这是你们合理的敬礼} \BibleRef{罗 12:1}，并且你要效法人子，祂\BibleQuote{来不是为受服事，而是为服事人}。\BibleRef{玛 20:28} 族长对客旅所做的，我们的主和师傅也为祂的仆人和门徒做了。\BibleQuote{以皮换皮，凡人所有的，他都会为性命舍弃。但}，魔鬼说，\BibleQuote{你触碰他的肉，他必当面诅咒你。} \BibleRef{约 2:4} 那古老的仇敌知道，与不洁的战争比与贪婪的战争更难。轻易丢弃那从外附着于我们的事物是容易的，但内部边境的战争则危险得多。我们必须拆解那连结在一起的东西，必须割裂那紧密结合的东西。匝凯富有，而宗徒们贫穷。他归还了所夺取的四倍，并将剩下的一半施舍给穷人。他迎接基督为客，救恩临到了他的家。\BibleRef{路 19:2} 然而，因为他身材矮小，够不上宗徒的高度，他没有被列入十二宗徒之中。至于财富，宗徒们根本放弃了什么，但就意愿而言，他们全都放弃了整个世界。如果我们向基督献上我们的灵魂和财富，祂必欣然接受我们的祭献。但如果我们只将外在的东西献给天主，而将内在的东西献给魔鬼，这分法便不公平，且天主的声音说：\Quote{你献得正当，却没有分得正当，岂不犯了罪？}\par

13. 你身为贵族阶层的领袖，首先树立了隐修的表率，这不应成为你自夸的事，而应成为谦卑的事，因为你知道天主子成了人子。无论你多么自谦，都不能比基督更谦卑。即使你赤足行走，穿着朴素，与穷人为伍，屈尊进入贫寒的住所，成为盲人的眼睛、弱者的手、跛子的脚，你担水、劈柴、生火——即使你做了这一切，主所受的锁链、掌掴、唾弃、鞭打、十字架和死亡，又在哪里呢？即使你做了我所说的一切，你仍被你的姐妹欧斯托基和保拉所超越：因为考虑到她们性别的软弱，她们在相对上做得更多，尽管在绝对上更少。当你的岳父托克索提乌斯还在世，他的女儿们仍沉迷于世时，我自己不在罗马，而在旷野——我但愿我仍在那里。但我听说她们嫌泥土路难行，由阉人抱着走，不喜欢走不平的路，觉得丝绸衣服太重，阳光太灼热。但现在，她们衣着褴褛朴素，比起从前，堪称女英雄。她们修剪灯芯、生火、扫地、洗菜、把白菜放进锅里煮、摆桌、递杯、上菜、来回服侍别人。然而，她们同屋的贞女并不缺乏。难道她们没有可以承担这些职责的仆人吗？当然不是。她们不愿让别人在体力劳作上超过她们，而她们自己在心灵刻苦上超过别人。我说这一切，并非怀疑你心灵的热忱，而是为了加快你奔跑的步伐，在战火炽热时为你增添温度。\par

14. 至于我，正在这个省建造一座隐修院，旁边还有一个招待所；这样，若瑟和玛利亚碰巧来到白冷，他们不会找不到栖身之所和欢迎。事实上，从世界各地蜂拥而来的隐修士人数如此庞大，我既不能放弃我的事业，也承受不了如此重担。福音的警告几乎在我身上应验了，因为我当初没有充分计算将要建的塔的代价；\BibleRef{路 14:28} 因此，我不得不派遣我的弟弟保利尼安到意大利，去卖掉一些逃过蛮族之手的破旧别墅，以及从我们共同父母继承的财产。因为我不愿，既然我已经开始了，就放弃服事圣徒，免得招来吹毛求疵和嫉妒之人的嘲笑。\par

15. 在写这封信的结尾时，我回想起我那个四马车的比喻，并且记起，若布拉西拉得以存留分享你的决心，她会成为第五匹。我几乎忘了提及她，你们中第一个去迎接主的人。你们曾经是五人，我现在看到是两人和三人。布拉西拉和她的姊妹保利纳安眠于甜蜜的睡眠中：你和两旁的另外两人将更轻快地飞向基督。\par

\SourceRecord{Translated by W.H. Fremantle, G. Lewis and W.G. Martley.；Nicene and Post-Nicene Fathers, Second Series；Vol. 6.；Edited by Philip Schaff and Henry Wace.；Buffalo, NY: Christian Literature Publishing Co.,；1893.；<http://www.newadvent.org/fathers/3001066.htm>.}

% structure: p0001 source=h1 target=section reason=page-title

\section{第六十八封书信}

\Emph{致卡斯特鲁提乌}\par

\BibleRef{若 9:3} 以及（二）通过告诉他\Person{安东尼}和\Person{狄迪默}的故事。这封信的日期是公元397年。\par

1. 我尊敬的执事之子\Person{赫拉克利乌}向我报告，说你因渴望见我，竟来到了\Place{基萨}；而且，你虽是\Place{潘诺尼亚}人，因而本是陆地动物，却不畏惧\Place{亚得里亚海}的波涛以及\Place{爱琴海}和\Place{爱奥尼亚海}的危险。他告诉我，若不是我们的弟兄出于亲切的关怀拦阻了你，你本已完成你的目的。但我仍然感谢你，并视之为一份善意。因为在朋友之间，必须以意愿代替行动。敌人往往给我们行动，但只有真诚的依恋才能给予我们意愿。如今我写信给你，恳求你不要认为临到你身上的身体疾病是出于罪。当宗徒们推测那生来瞎眼的人，并向我们的主和救主问道：\BibleQuote{谁犯了罪，是他还是他的父母，使他生来瞎眼呢？}他们被告知：\BibleQuote{不是他犯了罪，也不是他的父母，而是为叫天主的作为在他身上显扬出来。}\BibleRef{若 9:2} 难道我们没有看见许多外邦人、犹太人、异端者和各样观点的人，沉溺于情欲的泥沼，浴血，比豺狼更凶猛，比鸢鸟更贪婪，而这一切灾祸却没有临近他们的居所吗？他们不像别人那样受打击，因此他们向天主狂傲，甚至将脸抬起向着天。另一方面，我们知道圣人们遭受疾病、苦难和匮乏，也许他们被引诱说：\Quote{我真的白白洁净了我的心，在无辜中洗了我的手。}但立刻他们又责备自己：\Quote{若我说：我要这样讲，看哪，我就得罪了你儿女的一代。}如果你以为你的瞎眼是出于罪，而且医生常能治愈的疾病是天主愤怒的证据，那么你会认为\Person{依撒格}是个罪人，因为他完全失明，以致受骗祝福了一个他本不想祝福的人。\EditorialNote{创世纪第二十七章} 你也会指责\Person{雅各伯}有罪，他的视力变得如此模糊，以致看不见\Person{厄弗辣因}和\Person{默纳协}，\BibleRef{创 48:10} 尽管以内在的眼和先知的精神，他能预见遥远的未来和那要出自他王室的基督。\BibleRef{创 49:10} 难道有哪个君王比\Person{约史雅}更圣洁吗？然而他却死于埃及人的刀下。\BibleRef{列下 23:29} 难道有比\Person{伯多禄}和\Person{保禄}更崇高的圣人吗？然而他们的血染红了\Person{尼禄}的刀剑。且不说人，难道天主子没有承受十字架的耻辱吗？而你竟认为那些在今世享有幸福和快乐的人是蒙福的？天主对罪人最强烈的愤怒，是当他不发怒时。因此，在\Person{厄则克耳}中他对耶路撒冷说：\Quote{我的妒火必离开你，我将安静，不再发怒。}因为\BibleQuote{主所爱的，祂必管教，又鞭打凡所收纳的儿子。}\BibleRef{希 12:6} 父亲若不爱儿子，就不会管教他。师傅若不看弟子有希望的迹象，就不会纠正他。一旦医生放弃照顾病人，那就是他绝望的征兆。你应当这样回答：\BibleQuote{正如拉匝禄在他的一生中\BibleRef{路 16:25} 受了苦，我也愿现在甘心忍受痛苦，好使将来的荣耀为我存留。}因为\Quote{患难不再兴起第二次。}如果\Person{约伯}，一个在他那个时代圣洁、无玷、正义的人，遭受了可怕的苦难，他自己的书解释了原因。\par

2. 为避免冗长或超出书信的适当限度而重复旧事，我将简单向你叙述一件我童年发生的事。\Person{圣亚大纳削}，\Place{亚历山大}的主教，曾召请\Person{真福安东尼}到该城去驳斥那里的异端者。于是\Person{狄迪默}，一位博学但失明的人，来拜访这位隐修士；在谈到圣经时，\Person{安东尼}不禁赞赏他的才能，并称赞他的洞见。最后他说：\Quote{对于眼睛的失明，你难道不后悔吗？}起初\Person{狄迪默}羞于回答，但当问题被重复第二次和第三次时，他坦承他的失明对他是一个极大的悲伤。于是\Person{安东尼}说：\Quote{我感到惊讶，一个智者竟会为失去一个他与蚂蚁、苍蝇和蚊子共有的官能而悲伤，而不因拥有一个只有圣人和宗徒才被堪当享有的官能而欢喜。}从这个故事，你可以看出拥有属灵的视觉比属肉体的视觉好得多，并拥有罪恶的尘埃不能落入的眼睛。\BibleRef{路 6:42}\par

虽然你今年未能前来，但我并未对你前来失望。如果这位可敬的执事——这封信的携带者——再次被你情感的网络所困，如果你与他同行来到这里，我将很高兴欢迎你，并欣然承认，拖延的偿付因利息的丰厚而得以弥补。\par

\SourceRecord{Translated by W.H. Fremantle, G. Lewis and W.G. Martley.；Nicene and Post-Nicene Fathers, Second Series；Vol. 6.；Edited by Philip Schaff and Henry Wace.；Buffalo, NY: Christian Literature Publishing Co.,；1893.；<http://www.newadvent.org/fathers/3001068.htm>.}

% structure: p0001 source=h1 target=section reason=page-title

\section{第六十九封书信}

\Emph{\Person{奥克阿诺}}\par

\EditorialNote{\Person{奥克阿诺}是一位热忱于信德的\Place{罗马}贵族，他请求\Person{热罗尼莫}支持他抗议\Place{西班牙}主教\Person{卡特里乌斯}，因为\Person{卡特里乌斯}违反了宗徒关于主教应是一妻之夫的规则，第二次结婚。\Person{热罗尼莫}拒绝按所建议的方式行事，理由是\Person{卡特里乌斯}的第一次婚姻是在他受洗之前，因此不能考虑。所以他建议\Person{奥克阿诺}放弃此事。这封信的日期是公元397年。}\par

一、我从未想到，儿子\Person{奥克阿诺}，罪犯会攻击皇帝的仁慈，或者刚刚从监狱释放的人，在亲身经历了那里的污秽和镣铐之后，会抱怨允许给别人的宽免。在福音书中，嫉妒他人得救的人如此被说：\BibleQuote{朋友，因为我好，你就眼红吗？}\BibleRef{玛 20:15} \BibleQuote{天主把众人都禁锢在不信从之中，是为了要怜悯众人。}\BibleRef{罗 11:32} \BibleQuote{罪恶在哪里越多，恩宠在那里也越格外丰富。}\BibleRef{罗 5:20} 埃及的长子被击杀，连以色列的一头牲畜也没有留在埃及。\OriginalTerm{卡因派}{Cainites}的异端出现在我面前，那曾被击杀的毒蛇抬起它破碎的头，不是部分地，而是像往常一样完全地摧毁了\Person{基督}的奥秘。这异端宣称，有一些罪是\Person{基督}不能用自己的血洗净的，旧过犯留在身体和灵魂上的伤痕有时如此之深，以致于祂所提供的良药无法抹去。这不是说\Person{基督}白死了吗？如果有一些人是祂不能使之复活的，祂确实白死了。当\Person{洗者若翰}指着\Person{基督}说：\BibleQuote{看，天主的羔羊，除免世罪者！}\BibleRef{若 1:29} 如果真有一些人活着而\Person{基督}没有除去他们的罪，他就是在说谎。因为或者必须证明这些人不属于这个世界，以致\Person{基督}的恩宠无视他们；或者如果承认他们属于这个世界，我们就面临两难选择：要么他们已从罪恶中被解救，这证明了\Person{基督}拯救所有人的能力；要么他们仍未得救，仿佛仍处于罪行的指控之下，这证明了\Person{基督}的无能。但是我们决不敢对全能者认为祂在任何事上无能。因为\BibleQuote{凡父所作的，子也照样作。}\BibleRef{若 5:19} 将软弱归于圣子，就等于将软弱归于圣父。牧人背负整只羊，而不仅仅是这只或那只的一部分；宗徒的所有书信都不断谈论\Person{基督}的恩宠。而且，唯恐这个恩宠的一次宣告似乎微不足道，\Person{伯多禄}说：\BibleQuote{愿恩宠与平安丰富地赐予你们。}\BibleRef{伯前 1:2} 圣经应许丰盛，我们却声称匮乏。\par

2. 你问这一切的目的是什么？我回答：你记起你所提的问题。问题是：一位名叫\Person{卡特里乌斯}的西班牙主教，年事已高且在司铎职上多年，娶过两位妻子：一位在他受洗之前，她去世后，另一位在他经过洗礼池之后；你认为他违反了宗徒的训诫，宗徒在主教资格清单中命令主教应是\Quote{只作过一个妻子的丈夫}。\BibleRef{弟前 3:2}我很惊讶你单独攻击一个人，而整个世界都充满了在类似情况下被祝圣的人；我不是指司铎或更低级的神职人员，而只谈主教们，如果我逐一列举他们，人数足以超过参加\Place{阿里米诺}会议的人群。不过，我不应该通过指控多数人来为一个人辩护；也不应该，如果理性谴责一项罪，就以其犯罪人数多作为借口。在\Place{罗马}，一位能言善辩的辩论家用一句俗语所说，把我夹在了两难困境的犄角之间：无论我转向哪边，都被紧紧抓住。他问：\Quote{娶妻是有罪还是无罪？}我出于单纯，不够谨慎以避免为我设下的陷阱，回答说无罪。然后他提出另一个问题：\Quote{在洗礼中，被除去的是善行还是恶行？}这里，我的单纯又促使我说，是罪被赦免。这时，正当我开始自以为安全时，两难困境的犄角从这边和那边开始向我合拢，它们先前隐藏的尖端开始显露。他说：\Quote{如果娶妻不是罪，如果洗礼赦免罪，那么没有被除去的一切都保留下来。}瞬间，一片黑暗迷雾在我眼前升起，仿佛被一个强壮的拳击手击中。然而，回想起归于\Person{克吕西普}的那个诡辩：\Quote{无论你说谎还是说实话，你都在说谎}，我再次镇定下来，用我自己的两难困境回击对手。我说：\Quote{请问我，洗礼使人成为新人，还是不？}他勉强承认是的。我追问：\Quote{是使他完全成为新人，还是部分？}他回答：\Quote{完全。}于是我问道：\Quote{那么在洗礼中，旧人没有任何残留吗？}他表示同意。至此，我提出论证：如果洗礼使人成为新人，并创造了一个全新的存在，并且在新人中没有残留任何旧人，那么曾经在旧人中的东西就不能归咎于新人。起初，我那多刺的朋友沉默了；但后来，犯了\Person{皮索}的错误，尽管无话可说，却无法保持沉默。汗水出现在他的额头，脸颊变得苍白，嘴唇颤抖，舌头贴在上颚，喉咙干渴；恐惧（而非年龄）使他畏缩。最后，他脱口而出这些话：\Quote{难道你没有读过宗徒只允许被祝圣为司铎者是一个妻子的丈夫，而他强调的是婚姻的事实而非缔结的时间吗？}既然这家伙用三段论向我挑战，并且我看出他正试探一些复杂棘手的问题，我便开始用他自己的武器反击他。我说：\Quote{那么，宗徒挑选谁为主教？是受过洗的人，还是望教者？}他拒绝回答。但我再次猛攻，重复我的问题第二次、第三次。你会以为他是因过度哭泣而变成石头的\Person{尼俄柏}。我转向听众说：\Quote{善良的人们，我唤醒对手还是在睡梦中绑住他，对我来说都一样；但给一个不抵抗的人上镣铐更容易。}如果宗徒接纳进入神职队伍的不是望教者，而是信徒，并且被祝圣为主教的始终是信徒，那么作为信徒，就不能将望教者的过错归咎于他。这就是我向瘫痪的对手投掷的标枪。这就是我向他扔去的颤抖的长矛。最后，他张开嘴，吐出了他心中的内容。他脱口而出：\Quote{当然，这是宗徒\Person{保禄}的教义。}\par

3. 为此，我引出宗徒的两封书信：第一封致弟茂德，第二封致弟铎。在第一封中有以下段落：\BibleQuote{若有人渴望主教的职分，就是渴望善工。所以主教必须无可指摘，只做一个妻子的丈夫，有节制、谨慎、端庄、好客、善于教导；不酗酒、不打人……但要温和、不争斗、不贪财；善于管理自己的家，使儿女顺服，凡事庄重。（因为人若不知道管理自己的家，怎能照管天主的教会呢？）不可初信，恐怕他被骄傲冲昏，就落在魔鬼的刑罚里。而且他必须在外人那里有美名，免得他落在辱骂和魔鬼的网罗里。}\BibleRef{弟前 3:1} 而在致弟铎的书信开头，立即定下以下诫命：\BibleQuote{我从前留你在克里特，是要你将那未办完的事办妥，并照我所吩咐你的，在各城设立长老：若有人无可指摘，只做一个妻子的丈夫，儿女信主，没人告他们是放荡不服约束的，就可以设立。因为主教既是天主的管家，必须无可指摘，不任性、不酗酒、不打人、不贪无耻之财；却要好客、喜爱善人、有节制、公正、圣洁、有分寸；坚守所教真实的道理，就能以纯正的教训劝勉人，并且驳倒争辩的人。}\BibleRef{铎 1:5} 在两封书信中都命令：只有一夫一妻的人应被选入圣职，无论是作主教还是作长老。的确，在古人那里，这两个名称是同义词，一个指职务，另一个指圣职人员的年岁。无论如何，没有人能怀疑宗徒是仅指受了洗的人说的。因此，如果一个人在受洗前并不具备所有必需的资格，这决不影响他被祝圣为主教（因为问的是他是什么，而不是他过去是什么），那么，为什么以前的婚姻——这事本身不是罪过——却成为他接受圣职的障碍呢？你辩称：既然他的婚姻不是罪，所以在受洗时并没有被除去。这对我来说真是新闻，竟然要把本身不是罪的事当作罪。一切淫乱和公开恶行的沾污、对天主的亵渎、杀父和乱伦、改变两性的自然用途为逆性之用\BibleRef{罗 1:26}以及一切反常的淫欲，都在基督的泉源中洗去了。难道婚姻的污点竟然是无法消除的，而卖淫比体面的婚姻得到更宽大的判断吗？\Person{卡特里乌斯}可能会说：我不责备你拥有众多情妇和成群的宠妾；我不指控你的杀人流血和像猪在污秽泥沼中的打滚；但你却要为我多年的亡妻而搅扰我，我与她结婚本是为了避免陷入你所犯的那些罪。把这些话告诉那些作为教会庄稼、教会将其收进仓廪的外邦人吧；告诉那些寻求加入信友行列的望教者吧；告诉他们，我说，不要在受洗前缔结婚姻，不要进入体面的婚姻，而要与苏格兰人和阿塔科蒂人以及\Person{柏拉图}的\WorkTitle{理想国}中的民众一样，实行妻子共有和子女不区分，不仅如此，还要提防婚姻的丝毫表象；免得他们信了基督以后，基督对他们说，他们以前所拥有的不是姘妇或情妇，而是合法的妻子。\par

4. 每人當審查自己的\OriginalTerm{良心}{conscience}，為自己一生中對良心所施的暴力而哀慟；當他對自己從前的過犯作出真實判斷時，當聽耶穌的責備：\BibleQuote{假善人，先從你眼中取出木柴，然後你才能看得清楚，從你弟兄眼中取出木屑。}\BibleRef{瑪 7:5} 誠然，我們如同經師和法利塞人一樣，濾出蚊蚋，吞下駱駝；我們將薄荷與茴香獻上十分之一，卻忽略了天主所要求的正義判斷。妻子與妓女之間有何可比之處？將已因死亡而解除的婚姻作為指控的理由，而放蕩的生活卻為自己贏得讚美的花環，這公平嗎？他若前妻還在，絕不會另娶；至於你，如何為你隨意進行的獸性結合辯護？也許你會說，你害怕結婚，以免因此失去領受\OriginalTerm{聖秩}{ordination}的資格。他娶妻是為了生子；你因娶娼婦而失去了得子的希望。他在順從本性、求主祝福時，退入內室：\BibleQuote{你們要生育繁殖，充滿大地。}\BibleRef{創 1:28} 你卻在慾火中公然違反公共道德。他以羞怯的帷幕遮掩合法的享樂；你卻在眾人眼前無恥地追求不合法的享樂。為他寫著：\BibleQuote{婚姻應受尊重，床第應是無玷的；}對你卻讀出：\BibleQuote{因為淫亂和姦淫的人，天主必要審判}\BibleRef{希 13:4} 以及\BibleQuote{誰若毀壞天主的宮殿，天主必要毀滅他。}我們被告知，在\OriginalTerm{聖洗}{Baptism}中，一切罪過都獲得赦免；一旦我們蒙受了天主的仁慈，以後便不必畏懼祂作為法官的嚴厲。\OriginalTerm{宗徒}{Apostle}說：\BibleQuote{你們中有些人從前也是這樣，但是你們已經洗淨了，你們已經聖化了，你們已經因主耶穌的名並因我們天主的聖神成義了。}\BibleRef{格前 6:11} 因此所有罪過都獲赦免，這話是確實而可信的。但我問你，為何你的污穢被洗淨，而我的潔淨卻成了不潔？你回答說：\Quote{不，它沒有成為不潔，它仍是原樣。如果它是污穢，它就會像我的一樣被洗淨。}我想知道你這迴避之詞是何意。你的話語似乎像杵的圓頭一樣沒有尖銳之點。一件事不是罪，因此它就是罪嗎？或者一件事不是不潔，因此它就是不潔嗎？你說主沒有赦免，因為祂沒有什麼可赦免；但正因為祂沒有赦免，那未被赦免的仍舊存留。\par

5. 圣洗的真正效果为何，以及在基督内被祝圣的水所传递的真实恩宠是什么，我将在稍后告诉你；现在我先按照其应得的方式来处理这个论点。谚语常说：\InnerQuote{一个坏结只需一个坏楔子来劈开。}反对者所引用的经文，\BibleQuote{一个主教必须是一个妻子的丈夫}，可以有完全不同的解释。宗徒来自犹太人，而初期基督徒教会是从以色列的遗民中聚集起来的。保禄知道，法律允许男人通过多位妻子生育子女，\BibleRef{出 21:10}，并且他也知道圣祖们的榜样使一夫多妻制在人民中成为熟悉的事。甚至司祭们亦可自行决定享有同样的自由。\BibleRef{肋 21:7} 因此他命令教会的司铎不应主张这种自由，不应同时娶两三个妻子，而是每人一次只应有一个妻子。也许你会说，我所提出的这个解释是有争议的；既然如此，请听另一个解释。你不可垄断地扭曲法律以迎合你的意愿，而不是扭曲你的意愿以符合法律。有些人通过牵强的解释说，这段经文中的妻子应理解为教会，丈夫应理解为主教。尼西亚大公会议的神长们曾颁布法令，任何主教不得从一教会调任至另一教会，以免轻视贫穷而童贞的教区，去寻求富有而奸淫的教区的拥抱。因为正如\OriginalTerm{争论}{\GreekText{λογισμόι}}一词，即\Quote{争辩}，是指信仰中的儿子们的过犯和恶行，而关于管理家庭的诫命是指灵魂和身体的正确引导，\BibleRef{弟前 3:4}，同样，主教们的妻子应理解为他们的教会。关于这些，依撒意亚书上写道：\BibleQuote{你们这些妇女，赶快从炫耀中出来，因为这是无知的百姓。} 又说：\BibleQuote{富有的妇女啊，起来，听我的声音。} \BibleRef{依 32:9} 在箴言书中：\BibleQuote{谁能找到贤德的妇女？她的价值远胜珍珠。她丈夫的心信赖她。} \BibleRef{箴 31:10} 同一本书中也写道：\BibleQuote{智慧的妇女建造房屋；愚昧的妇女亲手拆毁。} \BibleRef{箴 14:1} 他们说，这并不贬损主教职位的尊严；因为同一比喻也用于天主。耶肋米亚写道：\BibleQuote{如同妻子背弃丈夫那样背弃我，以色列家！} \BibleRef{耶 3:20} 宗徒也用了同样的比喻：\BibleQuote{我曾把你们许配给一个丈夫，} 他对他的皈依者说，\BibleQuote{要把你们当作贞洁的童女献给基督。} \BibleRef{格后 11:2} 女人一词在希腊文中是歧义的，在这些地方都应理解为妻子。你会说这种解释牵强，且对原意有所扭曲。既然如此，就恢复圣经的简单含义，免得我不得不在你的地盘上与你争战。我要问你以下问题：一个人，在领洗前曾与一个妾室同居，在她死后领了洗并娶了妻，他能成为神职人员吗？你会回答说能，因为他的第一个伴侣是妾而不是妻。那么，宗徒所谴责的，似乎不是单纯的性关系，而是婚姻契约和夫妻权利。我们看到，许多人由于经济拮据而拒绝承担婚姻的重担。他们不娶妻，而与女仆同居，并将她们所生的孩子当作自己的孩子抚养。这样，如果藉着皇帝的恩赐，他们的情妇获得了穿戴主妇袍服的权利，他们就会立刻处于宗徒的轭下，极不情愿地不得不把她们的伴侣当作合法妻子来接受。但如果他们的贫困使他们无法得到我所说的那样一道皇帝谕旨，那么教会的法令就会随罗马法律的变化而变化。因此，你要注意，不要将\BibleQuote{一个妻子的丈夫}，即\Quote{一个女人的丈夫}，解释为认可无序的性关系而只谴责婚姻契约。\par

我提出所有这些解释，并不是为了抵制所讨论经文真实而简单的含义，而是为了向你表明，你必须按圣经所写的方式来理解圣经，并且你不可使救主所订立的圣洗仪式失去其效能，也不可使整个十字架的奥迹归于无效。\par

6. 现在让我履行刚才所作的承诺，以修辞学者的全部技巧，歌颂水和圣洗圣事。起初，地是混沌空虚，没有耀眼的太阳或苍白的月亮，也没有闪亮的星辰。只有无机、无形的物质，而且这物质还隐没在深渊中，被扭曲的黑暗笼罩。天主的神像驾车者一样运行在水面上\BibleRef{创 1:2}，从水中产生了初生的世界，这是从圣洗池中得生的基督徒孩子的预像。在天与地之间造成了穹苍，称为天——按希伯来语就是\Emph{Shamayim}，或\Quote{出自水中的}——天上的水与其他水分开，以赞美天主。因此在厄则克耳先知的异象中，革鲁宾之上显现出一块铺开的水晶\BibleRef{则 1:22}，即被凝聚而更浓密的水。最初的生物从水中出现；信徒则从圣洗池中长出翅膀，飞向天堂。人是由泥土形成的\BibleRef{创 2:7}，天主在他手心里握着奥秘的水。在\Place{厄登}园中种植了一个乐园，园中央有一道泉水分为四支\BibleRef{创 2:8}。这就是\Person{厄则克耳}后来所描述的同一道泉水，它从圣殿流出，流向日出之处，直到治愈苦水并使死者复活。当世界陷入罪恶时，只有水的洪流才能再次洁净它。但一旦邪恶的污秽之鸟被赶走，圣神的鸽子就来到\Person{诺厄}那里\BibleRef{创 8:8}，正如随后它来到约但河内的\Person{基督}那里\BibleRef{玛 3:16}，它用喙衔着一根象征恢复和光明的枝子，给全世界带来和平的喜讯。\Person{法郎}和他的军队不愿让天主的子民离开\Place{埃及}，被淹没在\Place{红海}中，这预表了我们的圣洗圣事。他的毁灭在\WorkTitle{圣咏集}中描述如下：\Quote{你用你的大能赋予海洋以德能：你在水中打碎了巨龙的头：你把\OriginalTerm{里外雅堂}{Leviathan}的头打得粉碎。}因此，蝰蛇和蝎子出没于干旱之地\BibleRef{申 8:15}，每当它们接近水时，就表现得像疯了一样或失去理智。正如木头使\Place{玛辣}变甜，使七十棵棕榈树得以享受它的溪流，同样，十字架使律法的水成为赋予生命的，给基督的七十位宗徒。\Person{亚巴郎}和\Person{依撒格}挖掘水井，\Person{培肋舍特人}却试图阻止他们。\Place{贝尔舍巴}，即盟约之城\BibleRef{创 21:31}，以及\Place{基红}，\Person{撒罗满}加冕之地，其名字都源于泉水。\Person{厄里厄则尔}在井旁找到了\Person{黎贝加}\BibleRef{创 24:15}。\Person{辣黑耳}也是一个打水之人，并因此从篡夺者\BibleRef{创 27:36}\Person{雅各伯}那里赢得了亲吻\BibleRef{创 29:10}。当\Place{米德扬}司铎的女儿们难以接近水井时，\Person{梅瑟}为她们开了一条路，将她们从暴行中解救出来\BibleRef{出 2:16}。主的前驱在\Place{撒冷}（这名意为平安或圆满）以泉水为基督预备百姓\BibleRef{若 3:23}。救主自己并未宣讲天国，直到祂以圣洗圣事的浸入洁净了\Place{约但河}。水是祂第一个奇迹的材料，并且\Place{撒玛黎雅}妇人被邀请从井中解渴\BibleRef{若 4:13}。祂暗中对\Person{尼苛德摩}说：\BibleQuote{人除非由水和圣神而生，不能进天主的国。}\BibleRef{若 3:5}正如祂的地上之旅始于水，也终于水。祂的肋旁被长矛刺透，流出了血和水，这是圣洗圣事和殉道的双重象征。此外，在祂复活后，当派遣宗徒们到外邦人那里时，祂命令他们以天主圣三的奥迹为这些人付洗\BibleRef{玛 28:19}。犹太人民为自己的恶行忏悔后，立即被\Person{伯多禄}遣去受洗\BibleRef{宗 2:38}。\Place{熙雍}尚未受产痛就生了儿女，一个民族霎时产生\BibleRef{依 66:7}。\Person{保禄}——教会的迫害者，那只来自\Person{本雅明}的贪狼\BibleRef{创 49:27}——在基督的一只羊\Person{阿纳尼雅}面前低头，只有在应用圣洗的救治时才恢复视力。通过先知的宣读，\Place{埃塞俄比亚}女王\Person{甘达刻}的太监预备好接受基督的圣洗\BibleRef{宗 8:27}。虽然违反自然，但\Place{埃塞俄比亚}人还是改变了肤色，豹子改变了斑点\BibleRef{耶 13:23}。那些只接受了\Person{若翰}的圣洗而不认识圣神的人，重新受了洗\BibleRef{宗 19:1}，以免有人以为未受圣神圣化的水能足以拯救犹太人或外邦人。\Quote{上主的声音在水上……上主在大水之上……上主使洪水居住其中。}祂的\BibleQuote{牙齿如同一群剪了毛的母羊，从洗浴中上来；每只都怀双胎，其中没有不生育的。}\BibleRef{歌 4:2}如果其中没有不生育的，那么它们所有的乳房都充满乳汁，并能随宗徒说：\BibleQuote{我的孩子们！我愿为你们再受产痛，直到基督在你们内形成；}\BibleRef{迦 4:19}和\BibleQuote{我喂给你们的是奶，不是硬食物。}\BibleRef{格前 3:2}正是圣洗的恩宠，\Person{米该亚}先知所指的：\BibleQuote{他必再怜悯我们，将我们的罪孽踏在脚下，将我们的一切罪恶投入海洋深处。}\BibleRef{米 7:19}\par

7. 那么，你怎能说一切的罪都在圣洗礼池中被淹没，而一个人的妻子却仍浮在水面作为反对他的证据呢？圣咏作者说：\Quote{过犯得赦免、罪恶得遮盖的人，是有福的。上主不归咎于他罪愆的人，是有福的。}看来我们还要在这首歌里加上一句，说：\Quote{上主不归咎于他有妻子的人，是有福的。}让我们也听听厄则克耳——被称为\Quote{人子}\BibleRef{则 2:1}——关于那真正人子——基督徒——德行的宣告：\Quote{我要从外邦人中领你们出来……那时我要在你们身上洒清水，洁净你们脱离一切不洁……我还要赐给你们一颗新心，在你们五内注入一种新精神。}\Quote{脱离一切不洁，}他说，\Quote{我要洁净你们。}若一切都挪去了，就什么也不能留下。若不洁被洁净了，那么清洁岂不更应免受玷污？\Quote{我还要赐给你们一颗新心，在你们五内注入一种新精神。}是的，因为\Quote{在耶稣基督内，割损或不割损都不算甚么，惟有新的受造物。}因此我们唱的这首歌也是一首新歌，\BibleRef{默 14:3}且脱去旧人，\BibleRef{弗 4:22}我们不按文字的旧样行走，而按心神的新样行走。\BibleRef{罗 7:6}这就是那块新石头，其上写着新名，\Quote{除了那领受的以外，没有人认识。}\BibleRef{默 2:17}\Quote{难道你们不知道，}宗徒说，\Quote{我们这受过圣洗归于耶稣基督的人，就是受圣洗归于祂的死吗？所以我们借着圣洗已与祂同葬于死中，为的是基督怎样由父的光荣中从死者复活，我们也怎样在新生活中行走。}\BibleRef{罗 6:3}我们屡屡读到\Quote{新样}与\Quote{更新}，却依然没有更新能擦去\Quote{妻子}一词带来的污点吗？我们借着圣洗与基督同葬，并因信那使祂从死者复活的天主的作为而复活了。\Quote{我们纵然死在过犯中，和未受割损的肉身中，天主却使我们与基督一同生活；赦免了我们的一切过犯，涂抹了那相反我们、攻击我们的律法规条，把它从我们中间除掉，钉在十字架上。}\BibleRef{哥 2:13}难道当我们的整个存在与基督同死，当旧\Quote{规条}中所记载的一切罪过都被涂抹时，唯独\Quote{妻子}一词仍存留吗？倘若我试图按次序向你们陈明圣经中所有与圣洗功效有关的段落，或解释那第二次出生——虽为第二次，在基督内却是我们的第一次——的奥秘教义，时间就不够用了。\par

8. 在我结束口述之前（因为我已经觉察到这封信已超过了信函的适当限度），我想简要解释一下宗徒在前几节中描述将被选为主教者的生活的经文。这样我们便会认出他是万民的导师，不仅因为他称赞一夫一妻，而且因为他的一切训诲。同时我恳请没有人会认为我在写作中所图的是抹黑如今的司铎。我唯一的目标是促进教会的利益。正如演说家和哲学家在论述完美的演说家和完美的哲学家时，并非贬低\Person{德摩斯梯尼}和\Person{柏拉图}，而只是提出抽象的理想；同样，当我描述一位主教并解释为主教职所定下的资格时，我只是为司铎们提供一面镜子。每个人的良心都会告诉他，他看到镜中映出的影像取决于他自己，是因其畸形而令他悲伤，还是因其美丽而令他喜悦。我现在转向所讨论的经文。\BibleRef{弟前 3:1} \BibleQuote{若有人想要监督的职分，他愿做善工。} 你看，是工作，不是地位；劳苦，不是享乐；工作是为了在谦卑中增长，而不是因高位而骄傲。\BibleQuote{那么，监督必须是无可指摘的。} 同样的，他对弟铎说：\BibleQuote{若有人无可指摘。} \BibleRef{铎 1:6} 所有的德行都包含在这一词中；这样，他似乎要求一种不可能的完美。因为如果每一个罪，甚至每一句闲话，都该受责备，世上谁能无罪无可指摘呢？然而，那被选为教会牧者的人，必须是一个与其他人相比，其他人理当被视为羊群的人。修辞学家将演说家定义为能言善道的善人。要配得上如此崇高的荣誉，他必须在生活与言语上无可指摘。因为一位教师若言行不一，便失去了一切影响力。\BibleQuote{只作一个妻子的丈夫。} 关于这一要求，我已在前面谈过。现在我只提醒你，如果一夫一妻是在圣洗前被强调的，那么其他条件也必须在圣洗前被强调。因为不可能认为其余义务只对领圣洗者具有约束力，而这一条却也对未领圣洗者具有约束力。警醒（或\Quote{节制}，因为希腊文\OriginalTerm{警醒（或节制）}{\GreekText{νηφαλιος}}两种含义都有）、明智、行为端庄、好客、善于教导。在天主的圣殿中服务的司铎被禁止饮酒、喝烈酒，\BibleRef{肋 10:9} 以免他们的头脑被醉酒麻痹，并使他们的理智能在天主的事奉中尽职。藉着\Quote{明智}一词，那些以纯朴为借口为司铎的愚昧辩护的人被排除。因为如果头脑不健全，所有肢体都会出错。\Quote{行为端庄}这一短语是对前面\Quote{无可指摘}一词的延伸。没有缺点的人被称为\Quote{无可指摘}；富有德行的人被称为\Quote{行为端庄}。或者，按照\Person{图利乌斯}的格言\Quote{最重要的是，你做事要做得优雅}可以有不同的解释。因为有些人如此不了解自己的分寸，\BibleRef{格后 10:14} 如此愚钝愚蠢，以致于因其姿态、步态、衣着或谈吐而成为旁观者的笑柄。他们幻想自己懂得什么是雅致，于是用华丽服饰和身体装饰打扮自己，并举办自诩优雅的宴会：但所有这类衣着和炫耀的尝试比乞丐的破布更令人厌恶。关于司铎必须作教导的责任，我们在旧约中有诫命，\BibleRef{申 17:9} 以及对铎更详细的指示。\BibleRef{铎 1:9} 因为一个纯良而不张扬的品行，其沉默造成的危害与其榜样带来的益处一样大。若是要吓退贪婪的狼，就必须靠狗吠和牧人的棍杖。\BibleQuote{不沉溺于酒，不打人。} 与所要追求的德行相对，他列出了要避免的恶习。\par

9. 我们已学会了我们应当是什么，现在让我们学习司铎不应当是什么。沉湎于酒是宴饮和狂欢者的过错。身体因饮酒发热，很快就沸腾于情欲。饮酒意味着放纵，放纵意味着感官享受，感官享受意味着破坏贞洁。凡耽于享乐的人，活着也是死的\BibleRef{弟前 5:6}，喝得酩酊大醉的人不仅是死的，而且是埋葬了的。一小时的放荡使\Person{诺厄}暴露了他那在六十年清醒中一直遮盖着的裸体\BibleRef{创 9:20}。\Person{罗特}在一次酩酊恍惚中，无意中在放纵之上又加了乱伦，酒征服了那个\Place{索多玛}未能征服的人\BibleRef{创 19:30}。一个好斗的主教被那将自己的背转给鞭打者\BibleRef{依 50:6}、被辱骂却不还口\BibleRef{伯前 2:23}的那位所定罪。\Quote{但是要节制}；一样好事摆在两样坏事面前。酗酒和暴躁必须由节制来抑制。\Quote{不好辩，不贪婪}。没有什么比无知者的自以为是更傲慢，他们以为喋喋不休就能说服人，总是准备争辩，用浮夸的辞藻向托付给他们的羊群大发雷霆。司铎必须避免贪婪，连\Person{撒慕尔}也教导了这点，当他在全体人民面前证明他没有拿过任何人的东西。同样的教训也由宗徒们的贫穷所表明，他们惯常从弟兄们那里接受供养和慰藉，并夸耀自己除了衣食之外，既没有也不想要任何东西\BibleRef{弟前 6:8}。\BibleBook{弟茂德前}{书}称为贪婪的，\BibleBook{弟铎}{书}则公开谴责为贪图不义之财\BibleRef{铎 1:7}。\Quote{好好管理自己的家}。不是靠增加财富，不是靠提供王侯般的宴席，不是靠拥有一堆精工制作的盘子，不是靠慢蒸野鸡使热量穿透骨头而不熔化上面的肉；不，而是首先要自己的家庭表现出他必须灌输给别人的品行。\Quote{使自己的子女完全端庄地服从}。他们不可效法\Person{厄里}的儿子们——他们与圣殿门廊中的女人同寝，以为宗教在于掠夺，便将祭品中最好的部分转用来满足自己的口腹之欲。\Quote{不是初皈依者，免得他骄傲自大，陷入魔鬼所受的判决}。我无法充分表达我对如此巨大盲目的惊愕：人们竟讨论像洗礼前结婚这类问题，并指责人们在洗礼中已死去、甚至与基督一同活于新生命中的事情，而没有人遵守这条关于主教不得是初皈依者的如此清晰无误的命令。昨日还是望教者，今日就成了主教；昨日还在斗兽场，今日就在教会中；昨晚还在竞技场，今早却站在祭台前；片刻前还是演员的赞助人，如今却成为童贞女的奉献者。宗徒难道不知道我们的推托和遁词吗？他对我们的愚蠢论调一概不知吗？他不仅说主教必须是一妻之夫，还命令他必须无瑕可指、警觉、清醒、品行端正、好客、善于教导、节制、不嗜酒、不斗殴、不好辩、不贪婪、不是初皈依者。然而我们对所有这些要求视而不见，只注意到候选人的妻子。谁能不举例说明这警告的必要性：\Quote{免得他骄傲自大，陷入魔鬼所受的判决？}片刻之间被立为司铎的人，不知道标志着最卑微信徒的谦卑和温良，不知道基督徒的礼貌，不够智慧以轻视自己。他从一个尊位升到另一个尊位，却未曾斋戒，未曾哭泣，未曾自责自省，未曾通过持续默想来修正生活，未曾将财产施舍给穷人。然而他从一个教区调往另一个教区，也就是说，他从骄傲走向骄傲。毫无疑问，宗徒所指的魔鬼的审判和堕落就是傲慢。所有片刻之间成为师傅（在成为门徒之前）的人都陷入其中。\Quote{并且他必须在教外人中也有好名声}。最后的要求与第一个相似。真正\Quote{无瑕可指}的人不仅得到自己家人的一致认可，也得到外人的认可。所谓教外之人是指犹太人、异端者和外邦人。因此，一位基督徒主教必须如此，以致那些挑剔他宗教信仰的人不敢挑剔他的生活。但目前我们却看到太多主教愿意像赛马中的车夫一样，用金钱换取民众的掌声；而有些主教如此普遍被憎恨，以致无法从他们的人民那里榨取钱财，而小丑却能用几个手势做到这一点。\par

10. 这些就是，吾子\Person{奥恰努}，教会的导师们应以忧惧之心要求他人并自己遵守的条件。这些也是他们在选择司铎时应遵循的教规；因为他们不可按私人的仇恨与敌对来诠释基督的法律，也不可迎合必然反噬心存怨恨者之恶感。试想\Person{卡特里乌}的品德何等无可指摘，在怀有恶意者眼中，他的生活除了在领洗前缔结的一桩婚姻外，别无苛责之处。\BibleQuote{那说不可奸淫的，也说不可杀人的；我们虽未犯奸淫，但若杀人，便成了犯法者。}\BibleRef{雅 2:11} \BibleQuote{凡遵守全部法律而只在一点上跌倒的，他就成了违反全部法律的。}\BibleRef{雅 2:10} 因此，当他们将洗礼前缔结的婚姻拿来指责我们时，我们必须要求他们遵守一切给予受洗者的诫命。因为他们放过许多不允许的事，却指责许多允许的事。\par

\SourceRecord{Translated by W.H. Fremantle, G. Lewis and W.G. Martley.；Nicene and Post-Nicene Fathers, Second Series；Vol. 6.；Edited by Philip Schaff and Henry Wace.；Buffalo, NY: Christian Literature Publishing Co.,；1893.；<http://www.newadvent.org/fathers/3001069.htm>.}

% structure: p0001 source=h1 target=section reason=page-title

\section{第七十书}

第七十书。致\Place{罗马}的演说家\Person{玛格努斯}。\par

\EditorialNote{热罗尼莫感谢\Person{玛格努斯}，一位\Place{罗马}的演说家，因为他帮助一名名叫\Person{塞贝修斯}的青年人为自己所做的大错事来向热罗尼莫道歉。然后他回应了玛格努斯对他喜爱引用世俗作家作品之风的批评，并用教会教父和圣经作者的榜样来辩护这一做法。最后他暗示，这一反对实际上并非来自\Person{玛格努斯}本人，而是来自\Person{鲁菲诺}（此处绰号卡尔普尔尼乌斯·拉纳里乌斯）。这封信写于公元397年。}\par

1. 我们的朋友\Person{塞贝修斯}从你的劝诫中获益，我从他的忏悔中得知，这甚于从你的信函中得知的。说来奇怪，他受责备后给我的快乐，大于他先前任性时给我的痛苦。确实在父亲这里有一种宽容，在儿子那里有一种孝爱之间的冲突：前者急于忘记过去，后者渴望承诺未来的顺服。因此，你我都应当同样喜乐：你因为成功考验了一位学生，我因为再次接纳了一个儿子。\par

2. 你在信末问我，为什么我有时在著作中引用世俗文学的例句，从而以异教的污秽玷污了教会的洁白。我现在简要回答你的问题。若不是你的心思完全被\Person{图利}占据，你绝不会问这个问题；若不是你习惯于研读\Person{沃尔卡提乌斯}而不是圣经及其注释者，你也绝不会问这个问题。因为谁不知道在\Person{梅瑟}和先知书中，有引自从外邦人书籍的段落，而且\Person{撒罗满}曾向\Place{提洛}的哲学家提出问题，并回答了他们向他提出的问题？在\BibleBook{}{箴言}的开头，他嘱咐我们要理解谨慎的格言和精明的谚语、比喻和隐晦的言谈、智者的言语和他们深奥的谜语；\BibleRef{箴 1:1}所有这些正属于辩证家和哲学家的领域。宗徒\Person{保禄}在写给\Person{弟铎}的书信中，也引用了诗人\Person{厄庇默尼德}的一句话：\BibleQuote{克里特人常说谎，是恶兽，是懒惰的肚腹。}\BibleRef{铎 1:12}这句话的后半部分后来被\Person{卡利马科斯}采纳。将这句话逐字译成拉丁文未能保留韵律并不奇怪，因为\Person{荷马}译成同一语言时，即使译成散文也几乎无法理解。在另一封书信中，\Person{保禄}引用了\Person{墨南德}的一句话：\BibleQuote{坏友毁善德。}当他在\Place{阿勒约帕哥}与\Place{雅典}人辩论时，他引用\Person{阿拉托}为证人，引用他的话：\BibleQuote{我们也是祂的子孙}\BibleRef{宗 17:28}，希腊文作\OriginalTerm{我们也是祂的子孙}{\GreekText{τοῦ} \GreekText{γὰρ} \GreekText{καὶ} \GreekText{γένος} \GreekText{ἐσμεν}}，这是一行英雄诗的结尾。似乎这还不够，这位基督军队的领袖、基督事业的不可战胜的辩护者，巧妙地将一块偶然的铭文转化为信仰的证明。\BibleRef{宗 17:22}因为他从真正的\Person{达味}那里学会了从敌人手中夺过刀剑，并用敌人的刀刃砍下傲慢的\Person{哥肋雅}的头。他曾在\BibleBook{}{申命纪}中读到上主的命令：当一个被俘的女子剃了头、剃了眉毛和所有头发、剪了指甲后，才可以娶她为妻。\BibleRef{申 21:10}那么，我仰慕她形体的美丽和她的雄辩风姿，难道不也希冀使这被我俘虏并为我婢女的世俗智慧成为真正\Person{以色列}的主妇吗？或者，剃去并剪除她身上一切死亡的东西，无论是偶像崇拜、享乐、错误还是色欲，我将她洁净而纯正地娶来，并为\Person{万军的上主}通过她生养仆人？我的努力促进了基督家人的利益，我与外邦人所谓的玷污增加了我同仆的人数。\Person{欧瑟亚}娶了一个淫妇，\Person{哥默尔}，\Person{狄布辣因}的女儿，这淫妇为他生了一个儿子，名叫\Person{依次勒耳}，即天主的子民。\BibleRef{欧 1:2}\Person{依撒意亚}说到一把锋利的剃刀，它剃去\BibleQuote{罪人的头和脚上的毛}\BibleRef{依 7:20}；\Person{厄则克耳}剃自己的头，作为那座曾为淫妇的\Place{耶路撒冷}的象征，\BibleRef{则 5:1}标志着她身上一切没有知觉和生命的东西都必须被除去。\par

3. \Person{西彼廉}，一位因其雄辩和殉道而闻名的人，据\Person{菲尔米安}告诉我们，曾因在他的驳\Person{德默特琉}的论文中使用了先知和宗徒的段落而受到攻击，后者宣称这些段落是捏造和杜撰的，而不是使用了哲学家和诗人的段落——这些作者的权威，作为外邦人的德默特琉无法轻易反驳。\Person{策尔苏}和\Person{波斐利}曾写文章攻击我们，并得到了有力的回应：前者被\Person{奥利振}驳斥，后者被\Person{美多德}、\Person{欧瑟伯}和\Person{阿波利纳利}驳斥。\Person{奥利振}写了一部八卷的论著，\Person{美多德}的作品长达万行，而\Person{欧瑟伯}和\Person{阿波利纳利}分别写了二十五卷和三十卷。读了这些，你会发现与他们相比，我在学识上只不过是一个初学者；而且，由于我的才智长期荒废，我几乎只能像在梦中一样回忆起我少年时学得的东西。\Person{尤利安}皇帝在\Place{帕提亚}战役期间抽出时间写了七卷反对基督的书，但正如诗篇传说中经常发生的那样，他只是用自己的剑伤了自己。如果我试图用哲学家和斯多葛派的教义来反驳他，你无疑会禁止我用\Person{海格立斯}的棍棒去打一条疯狗。的确，他很快就在战场上感受到了我们\Person{纳匝肋人}（或如他常说的\Person{加里肋亚人}）的手，一枪刺入要害，为他恶毒的诽谤付出了应有的代价。为了证明犹太民族的古老，\Person{约瑟夫}写了两卷书驳斥\Place{亚历山大}的语法学家\Person{阿彼翁}；在这些书中，他引用了如此多的世俗作家的段落，以至于我惊讶一个从小就接受圣经教育的希伯来人，竟也能通读整个希腊图书馆。我还要提到\Person{斐洛}吗？批评家称他为第二个或犹太人的柏拉图。\par

4. 现在让我逐一列举我们自己的作家。难道\Person{库阿德拉图斯}——一位宗徒的门徒兼\Place{雅典}教会的主教——不是曾在\Person{哈德良}皇帝（趁他参加\Place{厄琉息斯}秘仪之机）向他呈上一部为我们的宗教辩护的论著吗？他的卓越才华引起众人极大的钦佩，竟平息了一场最激烈的迫害。哲学家\Person{亚里斯提德}，一位极富口才的人，向同一位皇帝呈上一部由哲学作家的摘录编成的为基督徒辩护的著作。其后，另一位哲学家\Person{犹斯定}仿效他的榜样，向\Person{安敦尼·庇护}及其儿子们并元老院呈上一部论著《驳斥外邦人》，其中他卫护了十字架的耻辱，并大胆宣讲基督的复活。我是否还需提及\Place{撒尔底斯}的\Person{默利托}、\Place{希拉波里斯}教会大司祭\Person{阿波里纳瑞}、\Place{格林多}的主教\Person{狄约尼削}、\Person{塔提安}、\Person{巴尔德撒内}、继承殉道者\Person{波提诺}的\Person{依肋内}？他们都在多卷著作中解释了若干异端的兴起，并把它们一一追溯至其流出的哲学源头。\Person{潘泰诺}，一位斯多葛学派的哲学家，因为其学识的卓著声誉，被\Place{亚历山大里亚}主教\Person{德默特琉}派往\Place{印度}，向那里的婆罗门和哲学家宣讲基督。\Person{克莱孟}，\Place{亚历山大里亚}的一位司铎，按我的判断是众人中最博学的，撰写了八卷《杂集》和同样卷数的《大纲》，一部驳斥外邦人的论著，以及三卷称为《导师》的著作。难道这些著作缺乏学识，或者它们不是从哲学的核心中汲取的吗？效法他的榜样，\Person{奥力振}撰写了十卷《杂集》，其中他将基督徒和哲学家各自持有的意见相互比较，并通过引用\Person{柏拉图}、\Person{亚里士多德}、\Person{努墨尼乌}和\Person{科尔努图斯}的话，证实了我们宗教的一切教义。\Person{弥尔提亚德}也撰写了一部极佳的驳斥外邦人的论著。此外，\Person{希玻里图}和一位名叫\Person{阿颇罗尼}的罗马元老院议员各自编纂了护教著作。\Person{犹利乌·阿非利加努斯}——他撰写了自己时代的史书——的著作至今犹存，同样存在的还有\Person{塞奥多洛}（后来被称为\Person{额我略}）的著作，他不仅具有宗徒般的德行，还拥有宗徒般的奇迹。我们仍拥有\Place{亚历山大里亚}主教\Person{狄约尼削}、\Place{劳狄刻雅}教会大司祭\Person{阿纳多利}、司铎\Person{庞斐禄}、\Person{皮埃利乌}、\Person{路济安}、\Person{马尔基翁}、\Place{凯撒勒雅}主教\Person{欧瑟伯}、\Place{安提约基雅}的\Person{欧斯塔提}、\Place{亚历山大里亚}的\Person{亚大纳削}、\Place{厄梅撒}的\Person{欧瑟伯}、\Place{塞浦路斯}的\Person{特黎斐利}、\Place{西托波利}的\Person{阿斯忒里}、宣信者\Person{塞拉比雍}、\Place{波斯特拉}主教\Person{提托}，以及\Place{卡帕多西亚}的\Person{巴西略}、\Person{额我略}和\Person{安斐洛基}的著作。所有这些作家在他们的著作中如此频繁地穿插哲学家的学说和格言，以至于你可能很容易难以分辨更应赞赏哪一种：是他们的世俗学识，还是他们对圣经的了解。\par

5. 我将转向拉丁作家。有什么比\Person{戴尔都良}的风格更博学、更犀利呢？他的《护教篇》和《驳斥外邦人》包含了世间的一切智慧。\Person{米努齐乌·费利克斯}，一位罗马法庭的律师，搜罗了所有异教文献来润饰他的《屋大维》以及他的《驳斥占星家》的篇章（除非这后一部作品是假托其名的）。\Person{阿尔诺比乌}出版了七卷驳斥外邦人的著作，他的学生\Person{拉克坦提乌斯}出版了同样卷数，此外还有两卷：一卷《论愤怒》，另一卷《论天主的创造活动》。如果你阅读其中任何一部，你就会发现它们是\Person{西塞罗}对话的缩影。殉道者\Person{维多里诺}，虽然作为作家缺乏学识，但并非缺乏运用他所拥有学识的意愿。然后是\Person{西彼廉}。他的笔法何等简洁，何等通晓一切历史，何等雄辩的修辞和论证，他涉及了\Quote{偶像不是神}这一主题！\Person{希拉利奥}，我时代的一位宣信者和主教，在卷数和风格上都模仿了\Person{昆体良}的十二卷书，并且在他针对医生\Person{狄奥斯库鲁}的短论中也展现了他作为作家的能力。在\Person{君士坦丁}统治时期，司铎\Person{尤文库}用诗体叙述了我们的主和救主的故事，并且毫不犹豫地将福音的庄严语句强行纳入韵文。关于其他已故和在世的作家，我就不说了。他们的目的和能力对所有读到他们的人都是显而易见的。\par

6. 你切不可抱持错误的观点，以为在与外邦人打交道时可以诉诸他们的文献，而在其他所有讨论中则应忽略文献；因为几乎所有这些作家的著作——除了那些像\Person{伊壁鸠鲁}一样没有学问的人——都充满了学识和哲学。我实际上倾向于猜测——这个想法在我口授时涌上心头——你自己完全清楚在这一问题上学者们一贯的做法。我相信你向我提出这个问题，只是作为另一个人的代言人，此人因对\Person{撒路斯提乌斯}史书的喜爱而可能被称为\Person{卡尔普尔尼乌·拉纳里乌}。请恳求他不要因为自己没有牙齿而嫉妒食客的牙齿，也不要因为自己是鼹鼠而轻忽瞪羚的眼睛。这里如你所见，有丰富的讨论素材，但我已填满了我所能支配的篇幅。\par

\SourceRecord{Translated by W.H. Fremantle, G. Lewis and W.G. Martley.；Nicene and Post-Nicene Fathers, Second Series；Vol. 6.；Edited by Philip Schaff and Henry Wace.；Buffalo, NY: Christian Literature Publishing Co.,；1893.；<http://www.newadvent.org/fathers/3001070.htm>.}

% structure: p0001 source=h1 target=section reason=page-title

\section{\WorkTitle{第七十一封书信}}

\Emph{致路奇纽斯}\par

\EditorialNote{路奇纽斯是巴提卡（Bætica）的一位富有的西班牙人，他遵照当时的苦修理念，与妻子特奥多拉（Theodora）一起立了节欲誓愿。他对圣经研究很感兴趣，提议访问白冷，并在公元397年派遣了几位抄写员到那里，为他抄写热罗尼莫的主要著作。热罗尼莫现在将这些抄写员在返回家乡时托付了以下这封信。在信中，他鼓励路奇纽斯实现前往白冷的计划，描述了他寄给他的书籍，并回答了有关教会习惯的两个问题。他还送给他一些小礼物。}\par

不久之后（该信写于公元398年），路奇纽斯去世了，热罗尼莫写信给特奥多拉以安慰她的丧夫之痛（第七十五封书信）。\par

1. 你的信突然到来，是我未曾料到的；它以这意想不到的方式，藉着它所传达的喜讯，帮助我从麻木中苏醒过来。我急忙以爱臂拥抱这位我从未见过的人，并在心中默默自语：——\Quote{我恨不得有鸽子般的翅膀！好飞去安息！}那么，我就要找到\BibleQuote{我所爱的}\BibleRef{歌 3:1}。如今，主的话在你身上真正应验了：\BibleQuote{许多人要从东方和西方来，同亚巴郎一起坐席}\BibleRef{玛 8:11}。那些日子里，我的路奇纽斯的信德已在科尔乃略——\BibleQuote{所谓意大利营的百夫长}——身上预表了。\BibleRef{宗 10:1} 当宗徒保禄写给罗马人说：\BibleQuote{每逢我往西班牙去的时候，我要到你们那里去；因为我盼望在路过时见到你们，并蒙你们送我往那里去}\BibleRef{罗 15:24}，他藉着先前成功的叙述，表明了他期望从那个省份得到什么。短短时间内，他奠定了福音的基础，\BibleQuote{从耶路撒冷直到依里利谷}\BibleRef{罗 15:19}，他带着锁链进入罗马，为要释放那些陷于错误和迷信锁链中的人。他在自己租住的房屋里住了两年，\BibleRef{宗 28:30}，为要将我们在两约中所说的永恒之屋赐给我们。这位宗徒，捕人的渔夫，\BibleRef{玛 4:19} 撒下了他的网，在无数种类的鱼中，将你如同一条华丽的黄金鲷鱼捞了上来。你已抛下了苦海、咸潮和山裂；你已藐视了那在水中掌权的里外堂（Leviathan）。你的目标是同耶稣一起寻求旷野，并唱出先知的诗歌：\Quote{我的灵魂渴慕你，我的肉身切望你，在干旱疲乏无水之地；好得见你的能力和荣耀，正如我在圣所中见过你一样。}或如他在另一处唱道：\Quote{看，我就远走，住在旷野；我急急逃避暴风暴雨。}既然你已离开索多玛，正赶往山中，我以慈父的心恳求你，不要回头看。你的手已抓住犁柄，\BibleRef{路 9:62} 救主的衣繸，\BibleRef{玛 9:20} 以及祂那被夜露沾湿的卷发；\BibleRef{歌 5:2} 不要松开。不要从德行的屋顶下来，去寻找你昔日所穿的衣服，也不要从田间返回家里。\BibleRef{玛 24:17} 不要像罗特那样倾心于平原或那宜人的花园；\BibleRef{创 13:10} 因为那些地方不是像圣地那样由天降雨，而是由那被死海浸咸的约但河的浑水浇灌的。\par

2. 许多人开始，但很少人坚持到底。\Quote{赛跑的都跑，但只有一个得奖赏。}但对我们而言，经上却说：\BibleQuote{你们也该这样跑，好能得奖赏}\BibleRef{格前 9:24}。我们的竞赛大师并不吝啬；祂不把棕榈枝给一人而羞辱另一人。祂希望祂的所有运动员都能同样赢得花冠。我的灵魂喜乐，然而这极大的喜乐却使我悲伤。像卢德一样，\BibleRef{卢 1:14} 当我试图说话时，泪如雨下。匝凯，即刻悔改的人，被当作堪当接待救主为客的人。\BibleRef{路 19:5} 玛尔大和玛利亚预备了宴席，然后迎接主进去。\BibleRef{若 12:2} 一个妓女用眼泪洗祂的脚，并用善功的香膏膏抹祂的身体，为祂的安葬预备。\BibleRef{谷 14:8} 西满癞病人邀请主和祂的门徒，并没有被拒绝。\BibleRef{玛 26:6} 对亚巴郎说：\BibleQuote{你离开你的故乡、你的家族和父家，往我所要指示你的地方去}\BibleRef{创 12:1}。他离开加色丁，离开美索不达米亚；他寻求所不知道的，为不失去他所找到的。他认为不可能既保存自己的国家又保存主；即使在那个早期，他就已在实践达味先知的话：\Quote{我在你面前是客旅，是寄居的，像我列祖一样。}他被称作\Quote{希伯来人}，在希腊文中是\OriginalTerm{经过者}{\GreekText{περάτής}}，一个经过者，因为他不满足于现有的美善，而是忘记背后，向着前面奔跑。\BibleRef{斐 3:13} 他使圣咏作者的话成为自己的：\Quote{他们行走，力上加力。}因此，他的名字具有奥秘意义，并为你开辟了一条道路，使你寻求的不是自己的事，而是别人的事。你也要像他一样离开你的家，并以那些在基督里与你相连的人为父母、弟兄和亲戚。祂说：\BibleQuote{凡遵行我天父旨意的人，就是我的弟兄、姊妹和母亲}\BibleRef{玛 12:50}\par

3. 你身边有一位曾是你肉身上的伴侣、如今是你精神上的伴侣；曾是妻子、如今是姊妹；曾是女人、如今是男人；曾是较低者、如今是平等者。她与你同负一轭，一同向着那天国奔驰。\par

过分精细地管理自己的收入、过于细算自己的花费——这些习惯不易放下。然而，若瑟为逃脱埃及妇人，只得将外衣留在她手中。\BibleRef{创 39:12} 那跟随耶稣的青年，身上披着一块麻布，当仆役们袭击他时，他只得抛掉地上的遮盖，赤身逃走。\BibleRef{谷 14:51} 厄里亚乘着火马车升天时，也将他的羊皮外衣留在地上。厄里叟用他曾耕作的牛和犁具来祭献。\BibleRef{列上 19:21} 我们在《德训篇》中读到：\Quote{凡摸沥青的，必被它玷污。}\BibleRef{德 13:1} 只要我们仍忙于世俗之事，只要我们的灵魂被财产和收入束缚，我们就不能自由地思念天主。\Quote{因为正义与不法有什么相通？光明与黑暗有什么联系？基督与贝里雅耳有什么和谐？信者与不信者有什么相干？}\BibleRef{格后 6:14} \Quote{你们不能，}主说，\Quote{事奉天主而又事奉玛门。}\BibleRef{玛 6:24} 如今，舍弃金钱是为那些走在初阶路上的人，而非为已臻成全的人。异教徒如安提斯泰尼和底比斯的克拉特斯，此前早已做到。但将自己奉献于天主，这才是基督徒和宗徒的标志。他们如同那穷寡妇，将两个小钱投入银库，并将一切所有献给主，堪当听主的话：\Quote{你们也要坐在十二个宝座上，审判以色列十二支派。}\BibleRef{玛 19:28}\par

4. 你自能明白我为何提及这些事；我不明说，却是在邀请你定居于圣地。你的富足曾补助许多人的缺乏，好使他们的富足将来也补助你的缺乏；\BibleRef{格后 8:14} 你曾用\Quote{不义的钱财结交朋友，为使他们将来接你们到永久的帐幕里。}\BibleRef{路 16:9} 这样的行为值得称赞，堪与宗徒时代的德行相比。那时，如你所知，信徒们变卖产业，将价银带来放在宗徒脚前；\BibleRef{宗 4:34} 这象征性的行动旨在表明人必须践踏贪财。但主更渴望信徒的灵魂，而非他们的财富。我们在《箴言》中读到：\Quote{人的赎金是他自己的财富。}的确，我们可以将\Quote{自己的财富}理解为非来自他人或掠夺之物，正如诫命所说：\Quote{要以你正义的劳动恭敬天主。}但更好的理解是，将人\Quote{自己的财富}视为那些隐藏的宝藏，没有贼能偷，没有强盗能夺。\BibleRef{玛 6:20}\par

5. 至于我那卑微的著作——你说你渴望得到它们，这并非因为它们本身有什么价值，而纯粹出于你的善意——我已交给你的仆从们抄写，我看过他们抄录的纸本，并多次命令他们对照原文仔细校订。因为来往的朝圣者太多，我无法阅读那么多卷书。他们离开我时，我正被一场久病困扰，今年四旬期我正慢慢康复。因此，你若发现影响文意的错误或遗漏，那不可归咎于我，而要归咎于你的仆从们；那是抄写者的无知或疏忽所致，他们写下的是自己以为的意思，而非他们找到的文字，在尝试纠正他人错误时，反而暴露了自己的错误。有一个谣言传到你说我翻译了若瑟夫的著作和圣者帕皮亚及波利卡普的卷册。我既无闲暇也无能力将那些杰作的魅力用另一种语言保存起来。至于奥力振和狄迪慕斯，我翻译了一些，以向同乡展示希腊教义的若干样本。\OriginalTerm{希伯来真本}{Hebraica veritas}——除了\OriginalTerm{八部正典}{Octateuchus}（我目前尚在手头）——我已供你的仆从和抄写员使用。你无疑已拥有\OriginalTerm{七十贤士译本}{Septuaginta}的译本，那是多年前我为学生认真修订过的。至于新约，我已将其恢复至希腊原文的权威形式。因为旧约的真本只能通过参照希伯来文来检验，同样，新约的真本则需诉诸希腊文来判断。\par

6. 你问我是否应在安息日守斋，并每日领受圣体圣事——据传闻，这是罗马及西班牙教会的习惯。这两点都已被雄辩家希波利特讨论过，几位作者也收集了不同作者的相关段落。我能给你的最好建议是：教会传统——尤其是那些不与信德相悖的——应按照前代所传的形式遵守；一个教会的行为不应因与另一个教会相反而被废除。关于守斋，我真希望我们能不间断地守斋——如《宗徒大事录》所载，保禄和与他同信的人甚至在五旬节期和主日也如此行。他们不应被指控为\OriginalTerm{摩尼教}{Manichaeismus}，因为肉身的食物不应优先于属灵的食物。关于圣体圣事，你可以在任何时候毫无良心不安地领受，我不会反对。你可以聆听圣咏作者的话：\Quote{请体验并请观看，上主是何等美善；}你也可以如他那样歌唱：\Quote{我的心涌出优美的言辞。}但不要误解我的意思。你不可在庆节日守斋，也不可在五旬节期间平日禁食。在这些事上，每个省份可随其意愿，而所传的传统应被视为宗徒的法律。\par

7. 你从衣橱中送我两件小斗篷和一领羊皮外衣，并要求我自己穿着或施予穷人。作为回报，我向你和你在主内的姊妹送去四件适合你们修道生活和日常所需的粗毛布小苦衣，因为它们象征着贫穷，也是持续补赎的外在见证。我还附上了一部手稿，内含依撒意亚十个最隐晦的异象，我最近以评注加以阐明。当你们看到这些微物时，请思念你所喜爱的朋友，并加快你一度推迟的旅程。因为\BibleQuote{人的道路不在自己手中}，而是主\BibleQuote{指引他的脚步；}\BibleRef{耶 10:23}若有何阻碍——但愿没有——妨碍你前来，我祈愿距离不会割断那在爱中相连的人，并愿我能通过书信往还，使我在缺席中仍能感到我的\Person{卢奇尼乌斯}临在。\par

\SourceRecord{Translated by W.H. Fremantle, G. Lewis and W.G. Martley.；Nicene and Post-Nicene Fathers, Second Series；Vol. 6.；Edited by Philip Schaff and Henry Wace.；Buffalo, NY: Christian Literature Publishing Co.,；1893.；<http://www.newadvent.org/fathers/3001071.htm>.}

% structure: p0001 source=h1 target=section reason=page-title

\section{\WorkTitle{第七十二封书信}}

\Emph{致\Person{维塔利斯}}\par

\EditorialNote{维塔利斯曾问热罗尼莫：圣经告诉我们撒罗满和阿哈次在十一岁时成为父亲，圣经是否可信？这个难题之前也曾困扰过热罗尼莫本人（见第36封信第10节，或许维塔利斯是从那里得到的），在这封信中他提出了几种可能应对的方式。如有必要，他很愿意接受这一声称的事实，理由是圣经中有许多听起来难以置信却是真实的事情，并且自然不能抗拒自然之主（§2）。然而，他倾向于认为这个问题无关紧要且不重要。这封信的日期是公元398年。}\par

\SourceRecord{Translated by W.H. Fremantle, G. Lewis and W.G. Martley.；Nicene and Post-Nicene Fathers, Second Series；Vol. 6.；Edited by Philip Schaff and Henry Wace.；Buffalo, NY: Christian Literature Publishing Co.,；1893.；<http://www.newadvent.org/fathers/3001072.htm>.}

% structure: p0001 source=h1 target=section reason=page-title

\section{第七十三封信}

\Emph{致\Person{埃瓦格留斯}}\par

\EditorialNote{埃瓦格留斯曾送给热罗尼莫一份匿名论文，其中将默基瑟德等同于圣神，并询问他对此理论的看法。热罗尼莫在回信中驳斥该观念为荒谬，并坚称默基瑟德是一个真实的人，可能如犹太人所说，是诺厄的长子闪。这封信的日期是公元398年。}\par

\SourceRecord{Translated by W.H. Fremantle, G. Lewis and W.G. Martley.；Nicene and Post-Nicene Fathers, Second Series；Vol. 6.；Edited by Philip Schaff and Henry Wace.；Buffalo, NY: Christian Literature Publishing Co.,；1893.；<http://www.newadvent.org/fathers/3001073.htm>.}

% structure: p0001 source=h1 target=section reason=page-title

\section{书信七十四}

\Signature{致罗马的鲁菲努斯}\par

\BibleRef{列上 3:16}。热罗尼莫详细叙述了这段经文，将此叙事视为一个比喻，将假母亲和真母亲作为会堂与教会的预像。此信的日期是公元398年。\par

\SourceRecord{Translated by W.H. Fremantle, G. Lewis and W.G. Martley.；Nicene and Post-Nicene Fathers, Second Series；Vol. 6.；Edited by Philip Schaff and Henry Wace.；Buffalo, NY: Christian Literature Publishing Co.,；1893.；<http://www.newadvent.org/fathers/3001074.htm>.}

% structure: p0001 source=h1 target=section reason=page-title

\section{第七十五封信}

\Emph{致德奥多拉}\par

\EditorialNote{德奥多拉是有学问的西班牙人卢西尼乌斯（关于他，见第七十一封信）的妻子，最近失去了丈夫，这一丧事促成了本信的写作。在此信中，热罗尼莫讲述了卢西尼乌斯的许多美德，尤其是他抵抗马库斯异端（诺斯替派）的热忱，该异端在他生前盛行于西班牙。本信的写作年份是公元399年。}\par

1. 我被圣善且为我深为敬重的卢西尼乌斯安眠的悲伤消息所压倒，几乎无法口述哪怕一封短信。确实，我并不哀叹他的命运，因为我知道他已归于更好的境界：如同梅瑟，他可以说：\BibleQuote{我要转向观看这大异象，}\BibleRef{出 3:3}，但我心中充满遗憾，因为我未能得见他的面容——我原相信他不久就会来到此处。关于死亡命运的预言确实真实：它分离兄弟，用严酷无情的手拆散那些以爱德之链相连的名字。但我们有这样的安慰：死亡被主的话所击杀。因为经上说：\BibleQuote{死亡啊！我要作你的灾殃；阴府啊！我要作你的毁灭，}以及下一节：\BibleQuote{东风要来，主的风要从旷野吹来，他的泉源必干涸，他的水泉必枯竭。}\BibleRef{欧 13:14} 因为，正如依撒意亚所说：\Quote{从叶瑟的树干将生出一根枝条，从它的根上将长出一根嫩芽}；而祂自己在雅歌中说：\BibleQuote{我是沙仑的玫瑰，谷中的百合。}\BibleRef{歌 2:1} 我们的玫瑰就是死亡的毁灭，祂死了，好让死亡在祂的死亡中死亡。但是，当经上说祂要从\Quote{旷野}被带来时，是指童贞女的子宫，这子宫没有性交或受孕，却以婴孩的形象将天主赐给了我们，能借圣神的炽热熄灭情欲的泉源，并吟唱圣咏的话：\Quote{我好像在干涸无路无水之地，这样在圣所中向你显现。}因此，当我们面对死亡严酷无情的必然性时，我们被这安慰所支撑：我们将很快再次见到我们如今哀悼其缺失的人。因为他们的终结不称为死亡，而是睡眠和安眠。因此，蒙福的宗徒也禁止我们为那些安眠的人忧伤\BibleRef{得前 4:13}，告诉我们相信那些我们知道现在安眠的人，将来必从睡眠中被唤醒，当他们的睡眠结束时，他们将再次与圣人们一同警醒，并与天使一同歌唱：\Quote{天主受享光荣于高天，良善的人受享平安于地。}在天上，没有罪恶，只有光荣、永久的赞美和不倦的歌唱；但在世上，叛乱、战争和不和掌权时，平安必须借祈祷获得，并且它并非在所有人中，而只在良善的人中，他们留意宗徒的问候：\BibleQuote{愿恩宠与平安由天主我们的父和主耶稣基督赐予你们。}\BibleRef{罗 1:7} 因为\Quote{他的居所在平安中，他的住所在于熙雍，}即在于瞭望台，在于教义和德性的高处，在于信者的灵魂中；因为后者的天使每日得见天主的面\BibleRef{玛 18:10}，并以揭开的面容默观天主的荣耀。\par

2. 因此，虽然你已在路上奔跑，但我仍如俗语所说，鞭策一匹自愿的马，恳求你，在你为你的卢西尼乌斯（一个真正的兄弟）而惋惜的同时，也当因他现在与基督一同为王而喜乐。因为，正如智慧书上所记，他被\BibleQuote{取去，免得邪恶改变他的心意……因他的灵魂蒙主喜悦……他在短时间内完成了长久的工作。}\BibleRef{智 4:11} 我们更有理由为自己哭泣，因为我们每日与自己的罪恶争战，被恶习玷污，受到创伤，并且必须为每一句闲话交账。\BibleRef{玛 12:36} 他现在已胜利，无忧无虑，从高处俯视你，支持你的奋斗，不仅如此，他还为你在他身边预备了位置；因为他对你的爱与感情仍与从前一样——那时他忽略作为丈夫对你的权利，决心在世上待你如同姊妹，或者我可以说如同兄弟，因为性别的差异对婚姻是必要的，但对贞洁的联系则不然。既然即使在肉身中，如果我们重生在基督内，就不再是希腊人、外邦人、奴隶、自由人、男人、女人，而是众人在祂内合而为一，\BibleRef{迦 3:28} 那么当这必朽坏的穿上不朽坏，这必死的穿上不死时，岂不更是如此？\BibleRef{格前 15:53} 主告诉我们：\Quote{在复活时，}\BibleQuote{他们也不娶也不嫁，而是像天上的天使一样。}\BibleRef{玛 22:30} 如今，当经上说他们不娶不嫁，而是像天上的天使时，并非否定一个自然真实的肉身，而只是表明将要来临的荣耀的伟大。因为话不是说\Quote{他们将成为天使}而是\Quote{他们将像天使}：因此，虽然应许了与天使相似，却否认了与他们同一。基督告诉我们：\Quote{他们将要，}\Quote{像天使，}即像天使；因此他们不会不再是人。他们确实将是荣耀的，并享有天使的光辉，但他们仍然是人；宗徒保禄仍是保禄，玛利亚仍是玛利亚。那时，那异端将陷入混乱，它提出伟大但模糊的应许，只是为了夺走那些既谦逊又确定的希望。\par

3. 如今我一旦提到了\Quote{异端}这个词，何处能找到足够响亮的号角来宣告我们亲爱的卢西尼乌斯的雄辩？当巴西里德的污秽异端在西班牙猖獗，并如瘟疫般蹂躏比利牛斯山与海洋之间的省份时，他完整地持守了教会的信仰，全然拒绝接纳阿玛基尔、巴贝隆、阿布拉克萨斯、巴尔萨姆以及荒谬的勒乌西博拉。这些怪异的名字，为了激发无知之人和软弱妇女的心灵，他们假称源自希伯来来源，用野蛮的组合恐吓单纯之人，而这些组合越是难以理解，人们反而越加崇拜。这异端的增长，由里昂教会的依勒内主教为我们描述，他是宗徒时代的人，曾是帕皮亚斯的门徒，而帕皮亚斯是福音作者若望的听众。他告诉我们，有一个名叫马尔谷的人，出身于诺斯替派的巴西里德家族，起初来到高卢，用他的教导污染了隆河和加龙河所灌溉的地区，并且特别以他的错误误导了出身高贵的妇女；他向她们承诺某些秘密的奥迹，并通过魔法技艺和暗中放纵非法交合来赢得她们的情感。依勒内接着说，后来马尔谷越过比利牛斯山，占据了西班牙，他的目标是寻找富人的房屋，尤其是其中的妇女，关于这些人我们被告知：\BibleQuote{被各种私欲牵引，常常学习，终久不能明白真理。} \BibleRef{弟后 3:6} 所有这一切，他大约在三百年前写在他的极其博学且雄辩的著作中，这部著作题为 \WorkTitle{驳斥一切异端}。\par

4. 根据这些事实，你凭智慧将认识到，我们亲爱的卢西尼乌斯多么值得赞美：他塞住耳朵，以免听到对流血者所作的判决，并且散尽所有财产，分施给穷人，使他的义德永远常存。他不满足于将自己的慷慨施予本国，还向耶路撒冷和亚历山大的教会送去足够多的黄金，以缓解大量人群的匮乏。但许多人将钦佩和颂扬他的这种慷慨，而我则更愿赞扬他对圣经研究的热情和勤奋。他多么热切地请求我那些可怜的作品！他实际上派了六名抄写员（因为在这个省份，懂拉丁文的书记员匮乏）来抄录我从青年时期直到现在所口述的一切。这样的荣誉自然不是给我这微不足道的人——我是所有基督徒中最小的，住在伯利恒附近的岩石中，因为我知道自己是个罪人——而是给基督，祂在祂的仆人中受到尊敬，并如此应许他们：\BibleQuote{谁接待你们，就是接待我；谁接待我，就是接待那派遣我来的。} \BibleRef{路 9:48}\par

5. 因此，我亲爱的女儿，请把此信视为爱情促使我为你的丈夫所写的墓志铭。如果你认为我能胜任任何属神的工作，请大胆命令我承担它：这样，未来的世代将会知道，那位在\BibleBook{依撒意亚}{先知书}中论到自己说\BibleQuote{祂使我成为磨亮的箭；将祂藏匿在箭袋中} \BibleRef{依 49:2} 的那一位，用祂的利箭如此创伤了两个被极广阔的海域和陆地隔开的人，以至于尽管他们在肉体上彼此不认识，在精神上却在爱中紧密相连。\par

愿你在身体和灵魂上，被那位撒玛黎雅人——即救主和守护者——所保守，关于祂在\BibleBook{圣咏}{}中说：\BibleQuote{保护以色列的那一位，既不瞌睡，也不睡觉。} 愿那位守护者和圣者——曾降临到达尼尔那里的那位——也降临于你，使你也能说：\BibleQuote{我睡眠，但我的心警醒。} \BibleRef{歌 5:2}\par

\SourceRecord{Translated by W.H. Fremantle, G. Lewis and W.G. Martley.；Nicene and Post-Nicene Fathers, Second Series；Vol. 6.；Edited by Philip Schaff and Henry Wace.；Buffalo, NY: Christian Literature Publishing Co.,；1893.；<http://www.newadvent.org/fathers/3001075.htm>.}

% structure: p0001 source=h1 target=section reason=page-title

\section{第七十六封书信}

\Emph{致\Person{阿彼伽乌斯}}\par

\EditorialNote{阿彼伽乌斯是本信的收信人，他是西班牙贝提卡地区的一位盲人司铎。他曾请求热罗尼莫以祈祷帮助他对抗邪恶，热罗尼莫如今写信鼓励并安慰他。他在结尾处特别将寡妇特奥多拉托付给他照顾。此信应与致卡斯特鲁提乌斯的信（第六十八封）对照阅读。它写于前一封信的同时。}\par

1. 虽然我自知多罪，每日屈膝祈祷：求你不要记念我幼年的罪愆和我的过犯；但因我知道宗徒曾说过：\Quote{免得人自高自大，落在魔鬼所受的审判里}，又另一处记载：\BibleQuote{天主拒绝骄傲的人，却赐恩宠给谦逊的人}\BibleRef{雅 4:6}，因此，我从幼年至今，最极力避免的就是心高气傲与颈项僵硬——这些事常招致天主的义怒。因为我知道我的主、上主、天主在祂肉身的卑微中曾说：\BibleQuote{你们跟我学吧，因为我是良善心谦的}\BibleRef{玛 11:29}，在这之前祂也藉达味的口歌咏：\BibleQuote{上主，求你纪念达味和他所有的谦逊}。又在另一处读到：\BibleQuote{心高气傲必先跌倒，谦下尊前必见光荣}\BibleRef{箴 18:12}。所以，我恳求你不要以为我收到了你的信却默不作声。我恳求你不要把别人的不诚实和疏忽归咎于我。因为当我被召唤回应你的善意时，我为何要继续沉默，用我的沉默拒绝你所提供的友谊呢？我一向积极寻求与善人亲密交往，甚至主动争取他们的情谊。我们读到：\Quote{两个人总比一个人好……若跌倒，这人可以扶起他的同伴……三股合成的绳子不易折断，弟兄帮助弟兄必得高举。}因此请大胆写信给我，藉频繁的交谈克服别离的困扰。\par

2. 你不可因缺少了那属于蚂蚁、苍蝇和爬虫也有的肉眼而悲伤；倒应欢喜，因为你拥有那在\WorkTitle{雅歌}中所说的眼睛：\BibleQuote{我妹子，我新娘，你夺了我的心；你以你的一只眼夺了我的心}\BibleRef{歌 4:9}。这眼睛正是用以看见天主的眼睛，梅瑟指着它说：\BibleQuote{我要过去看这大异象}\BibleRef{出 3:3}。我们甚至读到世上有些哲学家挖出自己的眼睛，为将全部思想转向纯洁深邃的理智。一位先知曾说：\BibleQuote{死亡从你的窗户进来}。我们的主也对宗徒们说：\BibleQuote{凡注视妇女而好色者，心里已经与她犯了奸淫}\BibleRef{玛 5:28}。因此他们受命举目观看田地，因为田地已经发白，可以收割了。\BibleRef{若 4:35}\par

3. 你请求我以劝言为你杀死你内的拿步高、辣布沙刻、乃步匝辣当和敖罗斐乃。如果他们活在你内，你就绝不会求我帮助。不，他们已死在你内，你已开始藉则鲁巴贝耳、约叔亚（约匝达克之子大司祭）、厄斯德拉和乃赫米雅重建耶路撒冷的废墟。你不把工钱放入破漏的口袋\BibleRef{盖 1:6}，却为自己积蓄财宝在天上\BibleRef{玛 6:20}。你若寻求我的友谊，是因为你相信我乃是基督的仆人。\par

我把我的神女\Person{特奥多拉}托付给你——尽管她无需他人推荐，她自身就是推荐。她曾是已故的\Person{卢奇尼乌斯}（愿主赐福）的妻子，更可以说是他的姊妹。请告诉她，不要厌烦所踏上之路；她惟有辛苦穿越旷野，才能到达圣地。提醒她，不要以为离开埃及便完成了德行的功夫；要让她知道，要到达乃波山和约旦河\BibleRef{户 33:47}，她必须经过无数陷阱；她必须在基耳加耳重新受割礼；耶里哥必须在她面前倾覆，为司祭号角之声所摧毁\BibleRef{苏 6:20}；阿多尼责德克必须被杀；哈依和哈祚尔——昔日最美丽的城邑——都必须倾覆。\par

与我一同在隐修院的弟兄们问候你；我藉你诚恳地问候那些惠然顾念我的可敬人士。\par

\SourceRecord{Translated by W.H. Fremantle, G. Lewis and W.G. Martley.；Nicene and Post-Nicene Fathers, Second Series；Vol. 6.；Edited by Philip Schaff and Henry Wace.；Buffalo, NY: Christian Literature Publishing Co.,；1893.；<http://www.newadvent.org/fathers/3001076.htm>.}

% structure: p0001 source=h1 target=section reason=page-title

\section{第七十七封书信}

\Emph{致\Person{奥瑟努斯}}\par

\EditorialNote{这篇关于法比奥拉的颂辞是在她不安的一生于公元399年结束后写成的。热罗尼莫讲述了她犯罪和忏悔的故事（见第五十五封书信），她在波尔图建立的医院，她前往白冷的访问，以及她研读圣经的热忱。他提到自己曾为她撰写关于大司祭祭服的论述（第六十四封书信），并在她去世时，应她的请求正从事关于以色列人在旷野中四十二站注释的工作（第七十九封书信）。这最后一篇注释现在随此信一同送给奥瑟努斯。热罗尼莫也赞扬帕玛丘作为法比奥拉一切劳作的伴侣。此信日期为公元399年。}\par

1. 几年以前，我安慰了可敬的\Person{葆拉}，当时她因\Person{布莱西拉}的安眠而刚遭丧女之痛。四年前的夏天，我为主教\Person{赫略多若}撰写了\Person{内波提安}的墓志铭，并尽我所能表达了对失去他的悲痛。仅两年过去，我寄了一封简短的书信给我亲爱的\Person{帕玛丘}，关于他\Person{宝利纳}的突然离世。面对如此有学识的人，我羞于多说，也不愿将他自己已有的思想还给他：恐怕自己不像一位安慰朋友的人，而更像一个多管闲事、在完美者面前指导的人。但现在，我儿\Person{奥瑟努斯}，你交付我的这项义务，我欣然接受，即使无人请求我也愿主动承担。因为当需要论述新的美德时，旧的主题本身也会变新。在先前的事例中，我不得不缓和并克制一位母亲的慈爱、一位叔父的悲痛和一位丈夫的思念；根据各自的不同需要，我必须从圣经中应用不同的药方。\par

2. 今天你给我的主题是\Person{法比奥拉}，基督徒的赞美，外邦人的惊奇，穷人的哀痛，隐修士的安慰。在我选择首先论述她品格的哪一点时，都会因随后出现的那些而黯然失色。我该赞扬她的斋戒吗？她的施舍更大。我该推荐她的谦卑吗？她信仰的光辉更明亮。我该提及她刻意朴素的衣着，她自愿选择平民的装束和奴隶的外衣以羞辱丝绸长袍吗？改变性情比改变衣着是更大的成就。对我们而言，舍弃傲慢比舍弃黄金和宝石更难。因为即使我们抛弃这些，有时我们还会以其实虚荣的卑贱自满，并以一种可出售的贫穷来博取大众的掌声。但一种寻求隐藏并在内心意识中珍视的美德，只诉诸于天主的判断。因此，我要赐予法比奥拉的颂辞将是全新的：我必须忽略修辞学家的顺序，只从她悔改和忏悔的摇篮开始说起。另一位作家，心存学派，或许会提出\Quote{因延迟而拯救了罗马的那个人}\Person{昆图斯·马克西穆斯}，以及整个\Person{法比安}家族；他会描述他们的奋斗和战争，并会庆幸法比奥拉通过如此高贵的血脉来到我们中间，显示枝干中不明显的品质仍存在于根中。但作为白冷客栈和上主马槽的爱好者，在那里童贞女分娩并产下了婴孩天主，我将推出基督婢女的家谱，不是来自历史上著名的家族，而是来自教会的卑微。\par

3. 因为在最初就有礁石挡路，她被谴责的风暴所淹没，因她离弃了第一位丈夫而另嫁了第二位，所以在她悔改之前，我必先为她辩白。她前夫被指控的罪过如此可怕，连娼妓或普通奴隶都无法忍受。如果我一一列举，就会破坏那位妻子的英勇，她宁愿承担分离的责备，也不愿诋毁那与她成为一体的丈夫的品格并暴露他的污点。我只提出这一理由，足以使一位贞洁的已婚妇女和基督徒女人获得宽免。主曾命令，妻子不可被离弃\Quote{除非为奸淫的事，并且若被离弃，她必须保持未婚}。现在，给与男人的命令在逻辑上也适用于女人。因为不可能，既然一个犯奸淫的妻子要被离弃，而一个不节制的丈夫却被保留。宗徒说：\Quote{与娼妓结合的，成为一体。} \BibleRef{格前 6:16} 因此，与嫖客和不洁之人结合的女人也与他成为一体。凯撒的法律确实与基督的法律不同：\Person{帕比尼安}命令一件事，我们的\Person{保禄}命令另一件事。地上的法律对男人的不贞给予自由，只谴责诱奸和通奸；情欲在妓院和女奴中不受约束地放纵，仿佛罪过由受害者的身份构成，而非由攻击者的意图。但在我们基督徒看来，对女人不合法的，对男人同样不合法，既然两者事奉同一天主，两者都受同样义务的约束。因此法比奥拉离弃了一个犯罪的丈夫——他们说得完全正确——他犯了这种那种罪行，几乎所有的邻舍都传扬这些罪，但妻子独自拒绝揭露。但如果有人指责她休夫后没有保持未婚，我承认这是一个过错，但同时声明这可能是出于必要。宗徒告诉我们：\Quote{结婚比燃烧更好。} \BibleRef{格前 7:9} 她当时相当年轻，无法保持寡居。用宗徒的话说，她看到肢体的律法对抗理性的律法，\BibleRef{罗 7:23} 她感到自己被锁链拖曳着，作为俘虏走向婚姻的放纵。因此，她认为公开承认自己的软弱并接受不幸婚姻的表象，比以独身的名义操持娼妓的行当更好。同一位宗徒愿意年轻的寡妇结婚，生育子女，不给仇敌毁谤的机会。\BibleRef{弟前 5:14} 他立刻解释他的意愿：\Quote{因为有些人已经转去随从撒殚。} \BibleRef{弟前 5:15} 因此法比奥拉内心确信：她认为自己休夫是合法的，休夫后可以自由再婚。她不知道福音的严格性剥夺了女人在丈夫还在世时再婚的一切借口；不知道这一点，尽管她设法躲避魔鬼的其他攻击，她在此处无意中暴露自己，受他伤害。\par

4. 然而，我为何还要纠缠于这些陈旧且已被遗忘的事，试图为\Person{法比奥拉}已亲自忏悔过的过错辩护呢？谁能相信，在她的第二任丈夫去世后——此时大多数寡妇，一旦摆脱了奴役的轭，便变得漫不经心，放纵自己，比以往更自由，频繁出入浴场，在街上溜达，到处显露她们娼妓般的面孔——就在这个时候，\Person{法比奥拉}清醒过来了呢？然而，正是在那时，她穿上了苦衣，公开承认自己的错误。正是在那时，在全体罗马人面前（在那座原先属于那位死于\Person{凯撒}剑下的\Person{拉特兰}的大殿里），她站在忏悔者的行列中，在主教、司铎和众人面前，露出她散乱的头发、苍白的容颜、污秽的双手和未洗的脖颈——众人见她哭泣，也都流泪。这样的补赎，还有什么罪过不能涤除？这样的眼泪，还有什么根深蒂固的污点不能洗去？\Person{伯多禄}用三次宣认，抹去了他的三次否认。如果\Person{亚郎}犯下了渎神之罪，将熔化的金子铸成牛犊的头，那么他兄弟的祈祷就弥补了他的过犯。\BibleRef{出 32:30} 如果圣洁的\Person{达味}，最温良的人，犯下了谋杀和奸淫的双重罪，他用了七天的禁食来赎罪。他躺在地上，在灰中打滚，忘记了他的王权，在黑暗中寻求光明。 \BibleRef{撒下 12:16} 然后，他将目光转向那被他深深冒犯的天主，用哀痛的声音呼号说：\Quote{我得罪了你，惟独得罪了你，在你眼前行了这恶，} 以及 \Quote{求你使我重获你救恩的喜乐，以慷慨的精神扶持我。} 那位以他的德行教导我如何站立而不跌倒的人，也以他的忏悔教导我，若跌倒，如何重新站起。在列王中，我们读到有谁像\Person{阿哈布}那样邪恶？经上论他说：\Quote{从来没有像阿哈布那样出卖自己，在上主眼前行恶的人。} 为了纳波特的血，\Person{厄里亚}斥责了他，先知向他宣布了天主的愤怒：\Quote{你杀了人，还要霸占产业吗？……看，我要降祸于你，消除你的后裔} 等等。然而，当\Person{阿哈布}听见这些话时，\Quote{就撕裂了自己的衣服，身穿苦衣，禁食……穿着苦衣，缓缓而行。} \BibleRef{列上 21:27} 然后，上主的话传到了提市贝人\Person{厄里亚}那里，说：\Quote{你看见\Person{阿哈布}在我面前怎样自卑吗？因为他在我面前自卑，我在他在世的日子不降这祸。} \BibleRef{列上 21:28} 啊，幸福的忏悔！它引来了天主的垂顾，并通过承认错误，改变了天主愤怒的判决！同样的事迹在\BibleBook{}{编年纪} \BibleRef{编下 33:12} 中归于\Person{默纳舍}，在\BibleBook{约纳}{书} \BibleRef{纳 3:5} 中归于\Place{尼尼微}，在福音中归于税吏。\BibleRef{路 18:13} 前者不仅获准得到宽恕，而且恢复了王位；次者阻止了天主即将降下的愤怒；而第三者，以手捶胸，\Quote{连举目望天也不敢。} 然而，尽管如此，税吏以谦卑地承认自己的过错回去，成了义人，远胜过法利塞人傲慢地夸耀自己的德行。不过，这里不是宣讲忏悔的地方，我也没有在写文章反对\Person{孟他努}和\Person{诺瓦坦}。否则，我会说忏悔是\Quote{悦乐天主的祭献}，我会引用圣咏作者的话：\Quote{天主的祭献就是破碎的精神}，以及\Person{厄则克耳}的话：\Quote{我更喜欢罪人的悔改，而非他的死亡} \BibleRef{则 18:23}，还有\Person{巴路克}的话：\Quote{起来，起来，耶路撒冷！} 以及许多其他由先知号角发出的宣告。\par

5. 但是我要说一件事，它既对我的读者有益，也与当前的主题相关。既然\Person{法比奥拉}在地上不以主为耻，那么主在天上也必不以她为耻。\BibleRef{路 9:26} 她将自己的伤口袒露在众人眼前，\Place{罗马}含泪看着那毁损她美貌的伤疤。她裸露四肢，披头散发，紧闭双唇。她不再进入天主的教堂，而是像\Person{梅瑟}的姊妹\Person{米黎盎}一样，\BibleRef{户 12:14} 独自坐在营外，直到那驱逐她的司铎亲自召她回来。她像\Place{巴比伦}的女儿从她娇嫩的宝座上下来，拿起磨石磨面，赤脚走过泪水的河流。\BibleRef{依 47:1} 她坐在炭火上，炭火成了她的帮助。那张曾取悦她第二任丈夫的脸，如今她用拳头击打；她憎恨珠宝，回避装饰，不能忍受观看细麻衣。事实上，她为自己所犯的罪而哀哭，如同犯了奸淫一般，并花费许多药费，急切地想治愈她唯一的伤口。\par

6. 在\Person{法比欧拉}的罪的浅滩上搁浅之后，我之所以如此长久地谈论她的补赎，是为了能为描述她的赞美打开一个更大且毫无阻碍的空间。当她在全教会的眼前恢复共融时，她做了什么？在顺遂之日，她没有忘记患难；\BibleRef{德 11:25} 既已遭遇一次船难，她不愿再次面对海洋的风险。因此，她没有重蹈旧日生活的覆辙，而是变卖并处理了她所有能拿到的财产（那财产庞大且与她的地位相称），将其换成金钱，并以此施惠于穷人。她是第一位建立医院的人，在那里她可以聚集街上的患者，照顾那些不幸患病和匮乏的人。我现在需要一一列举人类的种种疾病吗？需要提及被割掉的鼻子、被挖出的眼睛、被烧焦一半的脚、布满疮痍的手吗？或是浮肿和萎缩的肢体？或是长满蛆虫的病肉？她常常亲自背负患黄疸病或污秽不堪的人。她也常常清洗伤口中流出的脓液，那些即使是男人也不忍目睹的东西。她亲手喂食病人，用点滴的液体湿润垂死者几乎不动的嘴唇。我知道许多富有且虔诚的人，他们无法克服对这些景象的自然厌恶，便通过他人之手施行这项慈悲工作，给予金钱而非亲自援助。我并不责备他们，也绝不将他们意志的软弱理解为信德的缺失。然而，我虽原谅这种娇气，却无限赞颂那超越它的心灵的炽热热忱。伟大的信德视这些琐事为无物。但我知道那位身穿紫袍的富人因没有帮助拉匝禄而遭受了何等可怕的惩罚。\BibleRef{路 16:19} 我们所鄙视的可怜虫，我们连看都不愿看，甚至看到就作呕的，其实与我们一样是人，由同样的泥土造成，由同样的元素形成。他所受的一切，我们也可能受。那么，让我们把他的伤口当作我们自己的，这样，我们对他人痛苦的一切麻木都会因自怜而消退。\par

我即便有一百张嘴、铜铸的喉咙，\newline{}也无法说尽那些可怕疾病的形式。\par

\Person{法比欧拉}在受苦的穷人身上所奇妙缓解的（疾病）如此之多，以至于许多健康的人都嫉妒起病人来。然而，她对神职人员、隐修士和贞女们同样慷慨。有哪一座隐修院不是靠法比欧拉的财富维持的？有哪一个赤裸或卧床不起的人没有穿上她提供的衣服？曾有谁需要帮助而她没有迅速而毫不犹豫地给予？甚至罗马本身也不足以容纳她的慈悲。她或亲身，或通过可敬且可靠的人，从一个岛屿到另一个岛屿，不仅将她的慷慨施予\Place{伊特鲁里亚海}沿岸，也施予沃尔西人的地区，那里沿着那些偏僻蜿蜒的海岸，有隐修士们的团体。\par

7. 突然间，她不顾所有朋友的劝告，决心乘船前往耶路撒冷。在那里，她受到一大群人的欢迎，并短暂地接受了我的款待。事实上，当我回想起我们的会面时，我仿佛现在就看到她在这里，而不只是在过去。蒙福的耶稣，她对神圣经卷是何等的热忱、何等的认真！为了满足那种真正的渴望，她会迅速翻阅先知书、福音书和\BibleBook{圣咏}{集}；她会提出问题，并将答案珍藏在她内心的宝库中。然而，这种聆听的热忱并没有带来任何满足感：她在增进知识的同时，也增加了忧伤，\BibleRef{训 1:18} 并且火上浇油，只不过为更炽烈的热忱提供了燃料。有一天，我们面前有\Person{梅瑟}所写的\WorkTitle{户籍纪}一书，她谦逊地询问我其中众多名字的含义。她问道，为什么同一个支派在这段经文和那段经文中以不同方式被联系？何以巫者\Person{巴郎}在预言基督的未来奥迹时\BibleRef{户 24:15}，比几乎任何其他先知都更明确地谈论祂？我尽我所能回答并试图满足她的疑问。然后，她将书卷进一步展开，看到了人民离开埃及后前往约但河水的所有驻停地列表。当她问我每一个地点的含义和缘由时，我对一些表示不确定，对另一些以确信的语气处理，并在几个情况下坦承无知。于是她开始更激烈地追问，像是责备我不应该不知道我所不知道的事，同时却声称自己不配理解如此深奥的奥迹。总之，我羞于拒绝她的请求，允许她从我这里强索一个承诺，即我会专门为她写一部关于这个主题的著作。直到现在，我一直推迟履行我的诺言：如我现今所见，是出于天主的旨意，好使这部著作奉献给她的纪念。如同在先前的一部作品中，我用司铎的祭衣覆盖了她，同样，在当前的页面中，她可以欣喜，因为她已穿越这个世界的旷野，最终来到了预许之地。\par

8. 但我继续已开始的工作。当我为如此伟大的贵妇寻找合适的住所时——她对独修生活的唯一概念包括一处像\Person{玛利亚}的客栈那样的休憩之所——忽然，使者往来飞报，整个东方惊恐万状。消息传来：匈奴大军从\Place{麦奥提斯}倾巢而出（他们盘踞在冰封的\Place{塔纳伊斯河}与粗野的马萨格泰人之间，那里有亚历山大之门阻挡\Place{高加索山}后的蛮族）；他们骑着快马四处奔袭，使全世界充满恐慌和流血。当时罗马军队因内战滞留在\Place{意大利}，不在东方。关于这些匈奴人，\Person{希罗多德}告诉我们，在\Person{玛待王大流士}治下，他们曾奴役东方二十年，并向\Place{埃及}人和\Place{埃塞俄比亚}人逐年索取贡赋。愿\Person{耶稣}使这些野兽的进一步侵袭远离罗马世界！他们的到来处处出人意料，速度胜过传言；他们来到时，既不宽恕宗教，也不顾身份年龄，甚至对啼哭的婴儿也毫无怜悯。孩子们被迫在尚不能说已开始生活之前死去；幼小的孩子尚未意识到自己的悲惨命运，却可见他们在敌人的手中和武器前微笑。人们普遍认为侵略者的目标是\Place{耶路撒冷}，而且他们急于前往这座特殊城市是由于对黄金的过度贪求。因此，平时无人照管的城墙得到修缮。\Place{安提约基雅}处于被围状态。\Place{提洛}渴望与陆地隔绝，重新寻回它古代的岛屿。我们也被迫装备船只，停泊在海岸附近，以防敌人到来。无论海风多么猛烈，我们不得不更惧怕蛮族而非海难。然而，我们担忧的与其说是自身安全，不如说是与我们同行的贞女的贞洁。就在那时，我们中间也有纷争，内部的冲突使我们与蛮族的战斗相形见绌。我自己坚持留在东方久已定居的住所，并顺从对圣地的深厚热爱。\Person{法比奥拉}习惯从一个城市迁往另一个城市，除了行李中的物品外别无财产，于是返回了她的故乡；在曾经富裕的地方过贫穷的生活，寄居在别人的房屋里——她从前曾在自己家中接待过许多客人——并且，不过分延长我的叙述，在\Place{罗马}的眼皮底下，将她卖掉那笔财产所得的收益施舍给穷人，\Place{罗马}知道她卖掉了它。\par

9. 我只痛惜这一点：在她身上，圣地失去了一串最美丽的项链。\Place{罗马}收回了先前失去的，外教人的轻率毁谤之舌被他们自己的眼睛所见证驳倒。他人或许称赞她的怜悯、谦卑、信德；我更愿赞扬她灵魂的热忱。我年轻时曾写信劝勉\Person{赫略多洛}度隐修生活的信，她烂熟于心；每当她看到\Place{罗马}的城墙，她就抱怨自己身陷囹圄。她忘记自己的性别，不顾自己的软弱，只渴望独处，事实上她的灵魂所萦绕之处正是她身在之处。朋友的劝告无法留住她；她急切地想挣脱城市，如同摆脱束缚之地。她也没有将施舍的分配交给别人；她亲自分配。她的愿望是，在公允地将钱财分给穷人之后，她本人为了基督的缘故得到他人的支持。她是如此匆忙、如此不耐迟延，你会以为她即将启程。由于她总是准备好的，死亡无法发现她毫无准备。\par

10. 当我执笔称赞她时，我亲爱的\Person{帕玛丘}似乎突然出现在我面前。他的妻子\Person{保利纳}安眠，使他得以警醒；她先于丈夫离去，好让他存留下来作基督的仆人。虽然他继承妻子的遗产，但如今占有他所继承的却是别人——我指的是穷人。他和\Person{法比奥拉}竞相在\Place{波尔图}设立像\Person{亚巴郎}那样的帐篷。他们之间兴起的竞赛在于显示仁慈的卓绝。双方都战胜了，双方也都战败了。两人都承认自己同时是胜利者和失败者，因为每人都想独自完成的事，两人一起完成了。他们联合资源，合并计划，好使和谐推进那些竞争本会使之化为乌有的事。计划一经提出，便立即执行。买了一所房子作为避难所，人群蜂拥而入。\Quote{在雅各伯中没有劳苦，在以色列中没有患难。}海洋载着航海者，在此登陆时受到欢迎。旅行者匆忙离开\Place{罗马}，在启航前利用这温和的海岸。\Person{普布利乌斯}曾在\Place{马尔他岛}上为一位宗徒——且不容反驳——为一艘船上的船员所做的，\BibleRef{宗 28:7}法比奥拉和帕玛丘为大批人反复做了；他们不仅供给贫困者的需要，而且他们的慷慨如此普遍，以至于为已经拥有一些的人提供了额外的资源。全世界都知道在\Place{波尔图}建立了一家旅客之家；\Place{不列颠}在夏天获悉，而\Place{埃及}和\Place{帕提亚}在春天就已经知道。\par

11. 在这位高贵妇人的死亡中，我们看到了宗徒话语的应验：\Quote{一切事都协助那些敬畏天主的人获得益处。}她预感到将要发生的事，曾写信给几位隐修士，请他们来解除她所承受的重担；因为她愿藉不义的钱财结交朋友，为叫这些朋友在她死后接她到永恒的帐幕里。\BibleRef{路 16:9}他们来到她那里，她使他们成为她的朋友；她照自己所愿的方式安眠，终于卸下重担，更轻盈地飞升天堂。\Person{法比奥拉}生前对\Place{罗马}是何等伟大的奇观，这在她死后百姓的行为中显露出来。她刚呼出最后一口气，刚把灵魂交还给属于她的基督，\par

飞翔的传闻宣告着灾祸\par

聚集了全城来参加她的葬礼。圣咏被咏唱，殿宇的金色天花板因高举的\Quote{阿肋路亚}呼声而震动。\par

老少歌咏团颂扬她的功绩，\newline{}并歌唱她圣洁灵魂的赞美。\par

她的胜利远比\Person{弗里乌斯}战胜高卢人、\Person{帕皮里乌斯}战胜萨谟奈人、\Person{西庇阿}战胜\Place{努曼提亚}、\Person{庞培}战胜\Place{本都}更为光荣。他们征服了肉体的力量，而她却战胜了属神的邪恶。\BibleRef{弗6:12}我似乎现在还听见引领队伍前行的骑兵队，以及数千拥挤前来参加葬礼的人群的脚步声。街道、门廊和屋顶上，但凡能看见的地方，都不足以容纳观众。在那一天，罗马看见她所有的百姓聚集一处，在场的每一个人都自以为在她悔改的光荣中有一份。难怪人们要因一个人的得救而如此欢腾，因为她的归化，使天上的天使也喜乐。\BibleRef{路15:7}\par

我给你这礼物，\Person{Fabiola}，这是我年老力衰之时的最好赠品，权作葬礼的献仪。我多次赞美过那些始终保持衣裳洁白、跟随羔羊无论祂往哪里去的贞女、寡妇和已婚妇女。\BibleRef{默14:4}那一生不受玷污而受称颂的，真是有福。但让嫉妒离去，让责难保持沉默。如果家主是良善的，为什么我们的眼要邪恶呢？\BibleRef{玛20:15}那落在强盗中的灵魂已被基督背在肩上带回家。在我父的家里有许多住处。\BibleRef{若14:2}罪恶在那里越多，恩宠在那里也格外丰富。\BibleRef{罗5:20}那得赦免多的，爱得也多。\BibleRef{路7:47}\par

\SourceRecord{Translated by W.H. Fremantle, G. Lewis and W.G. Martley.；Nicene and Post-Nicene Fathers, Second Series；Vol. 6.；Edited by Philip Schaff and Henry Wace.；Buffalo, NY: Christian Literature Publishing Co.,；1893.；<http://www.newadvent.org/fathers/3001077.htm>.}

% structure: p0001 source=h1 target=section reason=page-title

\section{第七十八封信}

\Emph{致\Person{法比奥拉}}\par

\EditorialNote{一篇论以色列子民四十二站或安营之处的论文，原计划为\Person{法比奥拉}所写，但直到她去世后才完成。与前一封信一起寄给\Person{奥西安努斯}。这些站象征基督徒的朝圣之旅，真正的希伯来人急于从地上赶往天堂。}\par

\SourceRecord{Translated by W.H. Fremantle, G. Lewis and W.G. Martley.；Nicene and Post-Nicene Fathers, Second Series；Vol. 6.；Edited by Philip Schaff and Henry Wace.；Buffalo, NY: Christian Literature Publishing Co.,；1893.；<http://www.newadvent.org/fathers/3001078.htm>.}

% structure: p0001 source=h1 target=section reason=page-title

\section{第七十九封书信}

\Emph{致萨尔维娜}\par

\EditorialNote{这是热罗尼莫写给萨尔维娜（一位宫廷贵妇）的慰问信，为其丈夫内布里迪乌斯去世而作。在为自己贸然致信一位完全陌生的人致歉后，热罗尼莫赞扬了内布里迪乌斯的美德，尤其是他的贞洁和对穷人的慷慨。接着，他用毫不阿谀的言辞警告萨尔维娜作为寡妇将面临的危险，并建议她全力以赴悉心教养她当前主要责任所在的儿子和女儿。这封信的语气有些傲慢，很难说是热罗尼莫最成功的作品之一。然而，萨尔维娜后来终身从事虔诚事业，成为金口若望的一位女执事。此信写于公元400年。}\par

1. 我恐怕我的尽责意愿会让人指责我自私；尽管我只是效法那说\BibleQuote{跟我学习，因为我是良善心谦的}，\BibleRef{玛 11:29}的一位，但我所采取的做法可能被归因于对名声的渴望。人们可能会说，我与其说是在安慰一个受苦的寡妇，不如说是在试图潜入宫廷；我以提供安慰为借口，实际上是在寻求与权贵的友谊。显然，凡知道这条诫命的人都不会这样认为：\BibleQuote{不可尊重穷人的情面}，\BibleRef{肋 19:15}，这一吩咐是为了避免我们在怜悯的借口下作出不公的判决。因为每个人应根据其案件的是非曲直来判断，而不是根据其地位。富人的财富若善用，不会成为障碍；穷人的贫困若在污秽和匮乏中未能远离恶行，也不能成为推荐。这些事物的证据在圣经时代和我们的时代都不缺乏；因为亚巴郎，尽管拥有巨额财富，却是\Quote{天主的朋友}，而穷人每天因犯罪而被法律逮捕和惩罚。我此刻致信的人既富足又贫穷，以至于她无法说出她实际拥有什么。因为我说的不是她的钱包，而是她灵魂的纯洁。我不认识她的面容，但我熟知她的美德；因为名声对她赞不绝口，她的年轻使她的贞洁更加值得称赞。她为年轻丈夫的悲伤为所有妻子树立了榜样；她的顺从证明她相信他并未失去，而是先她而去。她丧亲之痛的重大彰显了她信仰的真实。因为当她忘记失去的内布里迪乌斯时，她知道在基督内他仍与她同在。\par

但我为何写信给一个对我来说是陌生人的人？有三个理由。首先，因为（如同司铎当做的）我爱所有的基督徒如同我的子女，并以促进他们的福祉为光荣。其次，因为内布里迪乌斯的父亲与我有着最紧密的联系。最后——这是比其它理由更强的一个——因为我未能拒绝我的儿子阿维图斯。他以一种胜过寡妇对不义法官的纠缠\BibleRef{路 18:1}的固执，屡次写信给我，并在我面前提出众多我之前处理过类似主题的例子，以至于他克服了我谦逊的犹豫，使我决心做不是最合适于我的事，而是最符合他愿望的事。\par

2. 由于内布里迪乌斯的母亲是皇后的妹妹，他是在姑母的怀抱中长大的，另一个人也许会称赞他如此深受不败的皇帝德敖多西的喜爱。实际上，德敖多西从非洲为他寻得一位最高等级的配偶，这位配偶因她的祖国当时内乱频繁，而成为其忠诚的一种人质。我应当在开头就说，内布里迪乌斯似乎预感自己会早逝。因为置身于皇宫的辉煌和其身份而非年资所赋予的高位中，他始终生活得像一个相信很快就要去与基督相遇的人。关于意大利营的百夫长科乃略，神圣的叙述告诉我们天主如此完全接纳他，以至于派遣了一位天使到他那里；这位天使告诉他，他的功绩促成了那奥秘，即伯多禄从割礼的狭窄界限被带到未受割礼的广阔领域。他是第一个由宗徒受洗的外邦人，在外邦人中他被分别出来以获得救恩。关于这人，经上记载：\BibleQuote{在凯撒勒雅有一个人，名叫科乃略，是意大利营的百夫长，他是个虔诚而敬畏天主的人，全家都敬畏天主，他慷慨施舍给百姓，时常向天主祈祷。}\BibleRef{宗 10:1} 所有关于他说的话，我宣称——只是换了名字——属于我亲爱的内布里迪乌斯。他是如此\Quote{虔诚}，如此热爱贞洁，以至于在结婚时仍是纯洁的。他如此真实地\Quote{敬畏天主，全家都敬畏天主}，以至于忘记了自己的高位，将所有时间与隐修士和神职人员一起度过。他给百姓的施舍如此丰厚，以至于他的门口不断挤满了成群的患者和穷人。确实，他\Quote{时常向天主祈祷}，为使他得到最好的结果。因此\Quote{他被迅速接去，免得邪恶改变他的心意……因为他的灵魂悦乐了主。}因此，我可以如实将宗徒的话应用在他身上：\BibleQuote{我确实看出天主是不看情面的；但在每个民族中，凡敬畏祂而实行正义的，都是祂所悦纳的。}\BibleRef{宗 10:34} 作为士兵，内布里迪乌斯没有因他的披风、剑带和成队的勤务兵而受害；因为当他穿着皇帝的制服时，他已被征召为天主服务。另一方面，那些披着粗布外套、暗淡束腰外衣、肮脏和贫困的外表，却以行为否定其崇高宣称的人，一无所获。我们在福音中找到主对另一位百夫长的见证：\BibleQuote{我在以色列从未见过这样大的信德。}\BibleRef{玛 8:10} 再追溯更早的时代，我们读到若瑟，他在贫乏和富足时都证明了自己的正直，他作为奴隶和主人均教导灵魂的自由。他被提升到法郎身边，佩戴了王权的标志；\BibleRef{创 41:42} 然而他如此为天主所爱，以至于在所有圣祖中，他独自成为两个支派的父亲。\BibleRef{创 41:50} 达尼尔和三个孩子被派管理巴比伦的事务，被列为国家的王子；尽管他们穿着拿步高的服饰，心中却事奉天主。摩尔德开和艾斯德尔也身穿紫红袍、丝绸和珠宝，以谦卑克服了骄傲；虽然他们是俘虏，却如此受尊敬，以至于能够向征服者发号施令。\par

3. 这些论述旨在表明，我所提到的这位青年，利用他与王室的亲属关系、富足的财富和外在的权力标记，作为修德的助益。因为，正如训道者所言：\BibleQuote{智慧是保障，钱财也是保障} \BibleRef{训 7:12}。我们不可仓促断定这句话与主的话相冲突：\BibleQuote{我实在告诉你们：富人难进天国；我再告诉你们：骆驼穿过针眼，比富人进天国还容易。} \BibleRef{玛 19:23}。若是如此，圣经中描述为巨富的税吏\Person{匝凯}的得救，就与主的宣言矛盾了。但\BibleQuote{在人不可能，在天主却可能} \BibleRef{谷 10:27}，宗徒的劝诫教导我们，他这样写信给\Person{弟茂德}：— \BibleQuote{你要嘱咐今世富足的人，不要心高气傲，也不要寄望于无常的财富，而要寄望于赐予我们百物享用的永生天主；又要他们行善，在善事上富足，甘心施舍，乐意与人分享，为自己积存美好的根基，以备将来，好叫他们把握真正的生命。} 我们学会了骆驼如何穿过针眼：那背上有峰的牲畜，卸下驮子后，可以生出鸽子的翅膀，栖息在由芥菜籽长成的树枝上。\BibleRef{玛 13:31} 在\WorkTitle{依撒意亚书}中，我们读到骆驼，即\Place{米德扬}、\Place{厄法}和\Place{舍巴}的独峰驼，它们驮着黄金和乳香来到上主的城邑。\BibleRef{依 60:6} 同样，作为预象的骆驼，\Person{依市玛耳}商人\BibleRef{创 37:25}将香料、乳香和没药（就是在\Place{基肋阿得}出产、有益于医治创伤的那种\BibleRef{耶 8:22}）运到\Place{埃及}；他们何其幸运，在买卖\Person{若瑟}的过程中，竟以世界的救主为商品。\Person{伊索}寓言也告诉我们，一只老鼠饱食后，不能再像先前那样爬出去。\par

4. 我亲爱的\Person{内布里狄乌斯}日复一日思考着这些话：\BibleQuote{那些想要致富的人，就陷于诱惑和罗网}，魔鬼的\BibleQuote{和许多私欲} \BibleRef{弟前 6:9}。皇帝恩赐给他的所有钱财，或他的官职徽章为他带来的一切，他都用来周济穷人。因为他知道主的诫命：\BibleQuote{你若愿意成为成全的，去，变卖你所有的，施舍给穷人，然后来跟随我。} \BibleRef{玛 19:21} 由于他不能字面遵行这些指示（因他有妻子、幼小的儿女和众多的家仆），他便藉着不义之财结交朋友，使他们将来接他到永恒的帐幕里\BibleRef{路 16:9}。他没有像宗徒们那样抛弃父亲、渔网和船\BibleRef{玛 4:18}，一次性舍弃自己的弟兄，而是用自己多余的供给他人的缺乏，以便将来他们的富余可以补足他的缺乏\BibleRef{格后 8:14}。收信的那位女士知道，我所叙述的只是我耳闻之事，但她也知道，我并非希腊作家，以谄媚回报所受的恩惠。愿基督徒远离这种指控！我们有衣有食，就当知足。\BibleRef{弟前 6:8}\par

5. 然而，不可认为我只赞扬\Person{内布里狄乌斯}的慷慨施舍，尽管我们从这句话中学到施舍的重要：\BibleQuote{水能消灭烈火，施舍能补赎罪过。} \BibleRef{德 3:30} 现在我要继续谈他的其他美德，其中每一项都只在少数人身上找到。谁曾进入\Place{巴比伦}王的火窑而不被烧伤？\BibleRef{达 3:25} 可有青年人的衣服不被他的\Place{埃及}女主人抓住？可有阉人的妻子满足于无子女的婚床？可有任何人不对宗徒所说的争战感到惊骇：\BibleQuote{我发觉在我的肢体内，另有一条法律，与我理智的法律交战，并把我掳去，叫我隶属于那在我肢体内的罪的法律？} \BibleRef{罗 7:23} 但令人惊奇的是，\Person{内布里狄乌斯}虽在宫中长大，作为奥古斯都们的同伴和同学（他们的餐桌由全世界供应，由陆地和海洋服侍）；\Person{内布里狄乌斯}，我再说，虽处丰裕之中，又在青春年华，却比少女更谦逊，从未给丑闻流言以任何最微小的口实。此外，他虽是王公的朋友、同伴和表亲，并与他们一同受教育——这甚至使陌生人变得亲密——但他没有让骄傲膨胀自己，或对不如他幸运的人皱眉蔑视：不，他对所有人都友善，既爱王公如兄弟，又尊敬他们如君王。他常宣称自己的健康和安全有赖于他们。对于他们的侍从以及所有那些以数量增添宫廷威严的官员，他都通过表达关怀而如此妥当地赢得了他们的好感，以至于实际上不如他的人，因他的关注而相信自己是他的同侪。以德行遮盖自己的身份，并赢得那些被迫让你居先之人的爱戴，并非易事。有哪个寡妇不靠他的帮助？有哪个孤儿不从他身上找到父亲？整个东方的众主教常向他呈上不幸者的祈祷和苦难者的恳求。每当他向皇帝求恩，要么是为穷人求周济，要么是为俘虏求赎金，要么是为受难者求宽恕。因此，王公们也很乐意应允他的请求，因为他们深知，他们的慷慨不仅惠及一人，而且惠及多人。\par

6. 为什么我还要拖延结局？\Quote{凡有血肉的尽都如草，它的一切荣华都像田野的花。}\BibleRef{依 40:6} 尘土已归于尘土。\BibleRef{创 3:19} 他已在主内安眠，与他的祖先同眠，享满寿数、光明，并在美好的老年中受滋养。因为\Quote{智慧是人的白发。}\BibleRef{智 4:9} \Quote{他在短时间内}已\Quote{完成了长久的工作。}\BibleRef{智 4:13} 如今在他的位置上，我们有他可爱的孩子。他的妻子是他贞洁的继承人。对于那些想念他父亲的人，小\Person{内布里迪乌斯}再次向他展示，因为\par

他生前的眼睛、双手和容貌就是如此。\par

父母卓越性情的火花闪烁在儿子身上：孩子的面容如同镜子，透露出性格上的相似。\par

那窄小的身躯里蕴藏着一颗英雄的心。\par

与他在一起的还有他的妹妹，一篮玫瑰与百合，象牙与紫色的混合。她的面容虽像父亲，却更为迷人；她的肤色则来自母亲，与双亲如此相似，以致两人的轮廓都映照在她的样貌中。她甜美如蜜，是所有亲人的骄傲。皇帝不嫌弃将她抱在怀中，皇后最喜欢的是把她放在膝上抚育。人人都跑着争相第一个接住她。时而她依偎在这个人的颈项上，时而她在另一个人的臂弯中被宠爱。她咿呀学语，结结巴巴，而这笨拙的舌头使她更显可爱。\par

7. 因此，\Person{沙尔维纳}，你有了那些足以在你眼前代表你缺席丈夫的婴儿：\BibleQuote{看，子女是天主赐予的产业；胎中的果实是他的赏报。}在一位丈夫的位置上，你得到了两个孩子，因此你的情感有了比从前更多的对象。你所能给予他的一切，都可以给他们。以爱调和悲伤，因为他虽已离去，他们却仍与你同在。在天主眼中，善加抚养子女并非微末的功绩。请听宗徒的劝告：\BibleQuote{寡妇被登记在册，不可少于六十岁，只做过一个丈夫的妻子，因善行而有美名：如果她抚养了子女，如果她款待了旅客，如果她洗过圣徒的脚，如果她救济过受苦的人，如果她殷勤地追随一切善工。}\BibleRef{弟前 5:9}在此你得知天主要求于你的德行清单，你应如何活出你所持守的寡妇之名，以及借由哪些善行你可达至那仍然向你敞开的第二级贞洁。不要因宗徒不准六十岁以下的寡妇被选录而困扰，也不要以为他意在拒绝那些尚年轻的人。要相信你确实被他所选，他曾对门徒说：\BibleQuote{不要让人轻视你的青年，}\BibleRef{弟前 4:12}也就是说，不要因你年龄不足而轻视你，而非因你缺乏节制。若他的意思并非如此，所有六十岁以下成为寡妇者都将不得不嫁人。他是在训练一个在基督内仍未受教的教会，并为各阶层的人作准备，尤其是穷人，因为对他们的照管已交托给他和\Person{巴尔纳伯}。\BibleRef{迦 2:9}因此，他只愿那些不能亲手劳动、且是真正的寡妇，\BibleRef{弟前 5:3}经年岁和生命证明的人，由教会的劳苦来供养。\Person{厄里}的孩子们的行径使身为司铎的他成为天主的憎恶。反之，天主因那些\BibleQuote{在信德、爱德、圣德与贞洁上持守}之人的德行而息怒。\BibleQuote{\Person{弟茂德}啊，}宗徒喊道，\BibleQuote{你要保守自己纯洁。}\BibleRef{弟前 5:22}我绝不敢怀疑你能做出任何错事；然而，提醒一位青春与财富引她入诱惑的人，只是一种善意。你要将我将要说的话当作不是对你说的，而是对你的青春年华说的。一位\BibleQuote{生活奢侈的寡妇，虽生犹死。}\BibleRef{弟前 5:6}这是那位\BibleQuote{蒙选的器皿}\BibleRef{宗 9:15}所说的，这言语从他的宝库中取出，他曾大胆地说：\Quote{你们要寻求基督在我内说话的凭据吗？}然而，这些话语出自一位亲身承认人性身体软弱的人，他说：\BibleQuote{我愿意行的善，我反而不行；我不愿意作的恶，我倒去作。}\BibleRef{罗 7:19}又说：因此，\BibleQuote{我痛击我的身体，使它成为奴仆，免得我给别人传报了，自己反而落选。}\BibleRef{格前 9:27}如果\Person{保禄}害怕，我们中谁还敢自信？如果\Person{达味}——天主的朋友——和那爱天主的\Person{撒罗满}，也像别人一样被战胜，如果他们的跌倒是为警告我们，他们的忏悔是为引导我们得救，在这滑溜的人生中，谁能保证自己不跌倒呢？决不要让野鸡出现在你的餐桌上，也不要肥美的斑鸠或\Place{爱奥尼亚}的黑鸡，或任何那些昂贵到能飞走最大财产的鸟类。不要以为当你拒绝猪肉、兔肉、鹿肉和其他四足动物的美味时，你就在戒除肉食。区别不在于脚的数目，而在于味道的精致。我知道宗徒曾说过：\BibleQuote{凡天主所造的样样都好，若以感恩的心领受，没有一样是可弃绝的。}\BibleRef{弟前 4:4}但同一宗徒又说：\BibleQuote{不吃肉、不喝酒，才是好的，}\BibleRef{罗 14:21}并在另一处说：\BibleQuote{不要醉酒，醉酒使人放荡。}\BibleRef{弗 5:18}\BibleQuote{凡天主所造的样样都好}——这命令是为那些操心如何取悦自己丈夫的人所设的。\BibleRef{格前 7:34}让那些服侍肉身、身体因情欲沸腾、与丈夫结合、一心想要生育的人吃肉吧。让那些子宫沉重的女人用肉塞满自己的胃。但你已将一切放纵葬在你丈夫的坟墓里：在他的棺架旁，你用泪水洗净了涂满胭脂和铅粉的脸庞；你已将白袍和金边便鞋换成了深色长袍和黑鞋；如今只缺少一件事：在禁食上持之以恒。让苍白和脏污从此成为你的珠宝。不要用羽绒床褥娇养你年轻的肢体，也不要用热水澡点燃你青春的热血。请听一位外教诗人在贞洁寡妇口中说过的话：\par

\Quote{他，我的首任丈夫，夺走了我的所爱。\newline{}就这样吧：让他把它们留在坟墓中。}\par

如果普通的玻璃如此珍贵，那么一颗珍贵的珍珠又价值几何？如果一位外教寡妇因顺从天主眷顾的自然法则而能谴责一切感官放纵，那么我们对一位基督徒寡妇，她对自己的贞洁所欠的不是一位已死的人，而是那位她在天上将与之一同为王的那一位，又该期待什么呢？\par

8. 我恳求你，不要以为这些普遍性的评论——它们适用于所有年轻女子——是要侮辱你或责备你。我怀着焦虑的心情写信，却祈祷你永远不必知道我所担忧的是什么。女人的名誉是柔弱的植物；它像美丽的花朵，在一丝微风中凋零，在一息气息中枯萎。尤其当她处于容易陷入诱惑的年龄，又缺乏丈夫的权威时更是如此。因为连丈夫的影子也是妻子的保障。寡妇要那么大的宅第和大群随从做什么？作为仆人，她当然不可轻视他们，但作为男人，她应在他们面前感到羞惭。如果大宅需要这样的奴仆，至少应派一个品德无瑕的老人管理他们，他的尊严可以保护主母的荣誉。我知道许多寡妇，虽然闭门独居，仍未能免于与仆人过于亲近的嫌疑。这些仆人当他们衣着超过身份，或身体肥胖光滑，或处于易受情欲影响的年龄，或因为知晓主人私下对他们的宠爱而在举止上过于自信时，就成了猜疑的对象。因为这种骄傲，无论怎样小心隐藏，必定会以轻看其他同是仆人的形式显露出来。我之所以说这些看似多余的话，是要你以全勤谨守护你的心 \BibleRef{箴 4:23}，并防备一切可能针对你的丑闻。\par

9. 不要带着卷发的管家与你同行，不要有柔弱的戏子，不要有唱着毒蜜般甜美的魔鬼歌手，不要有漂亮而剃得光滑的青年。不要让任何戏场上的恭维、任何谄媚的阿谀与你有关。让你身边总有寡妇和贞女；让你的安慰者是你同性的人。主母的品格由女仆的品格来判断。只要你身边有一位圣洁的母亲，只要你有一位誓守贞洁的姑母陪伴，你就不该忽略他们，而冒着风险寻求陌生人的陪伴。让你手中常有圣经，并如此频繁地祈祷，使那些常侵袭年轻人的邪恶思想之箭，因此找到盾牌来击退它们。要逃过那些内在动因的开端是困难的，不，甚至是不可能的；希腊人很有深意地称之为\OriginalTerm{预激情}{\GreekText{προπάθειαι}}，即\Quote{激情的预向}。事实上，罪的暗示撩拨我们所有人的心思，而决定在于我们自己的心，是接纳还是拒绝到来的思想。自然界的主自己在福音中说：\BibleQuote{由心里发出来的是恶念、凶杀、奸淫、邪淫、盗窃、妄证、亵渎。} \BibleRef{玛 15:19} 从另一本书的见证清楚看到：\BibleQuote{人从小心里就怀着恶念。} \BibleRef{创 8:21}，而且灵魂在宗徒所列举的本性私欲之作为与圣神之作为之间摇摆，\BibleRef{迦 5:19}，时而渴望前者，时而渴望后者。因为\par

没有人完全无过；\newline{}最好的人只是过错最少。\par

再者，引用同一诗人的诗句：\par

人们在白皙的皮肤上挑剔瑕疵。\par

先知用不同的词句表达了同样的意思：他说\Quote{我如此困苦，甚至不能说话}，在同一卷书中又说\Quote{你们要动怒，却不可犯罪。}\Place{塔伦图姆}的\Person{阿尔基塔斯}曾对一个粗心的管家说：\Quote{要不是我在生气，我早已把你鞭打至死。}因为\BibleQuote{人的忿怒并不成就天主的正义。} \BibleRef{雅 1:20} 这里关于一种情绪波动的说法，可以适用于所有情绪。正如动怒是人的本性，而抑制怒气是基督徒的本分，其他情欲也是如此。肉体总是贪恋肉体的事，并以它的诱惑吸引灵魂去享受致命的快乐；但我们基督徒应当借着对基督更强烈的爱来克制对感官放纵的欲望。我们应当借着斋戒来驯服我们内心那桀骜不驯的野兽，使它所渴望的不是淫欲而是食物，并且平稳安详地前行，以圣神为它的骑手。\par

10. 我为何这样写？为要向你表明，你不过是人，除非谨慎提防，否则也容易陷入人的情欲。我们所有人都是由相同的泥土造成，由同样的元素形成。无论我们穿丝绸还是破衣，我们都受制于相同的欲望。它不惧怕王袍，也不鄙弃乞丐的污秽。因此，宁可让肚腹受苦，也不让灵魂受损；宁可管束身体，也不受它奴役；宁可失去自持，也不失去贞洁。我们不要自欺，以为总能靠补赎回头。因为这最多不过是苦难的药方。让我们避免造成一个需要痛苦治疗的伤口。因为安然驶入救恩港湾，船货无损，是一回事；而赤身紧抱一块木板，被波浪抛来抛去，一次次撞在锋利的礁石上，是另一回事。一个寡妇应当不知道再婚是被允许的；她应当对宗徒的话一无所知：\Quote{与其欲火焚身，不如结婚。}\BibleRef{格前 7:9}除去那被认为是更坏的东西——燃烧的危险，婚姻就不再被看作好事。当然，我驳斥异端者的诽谤，我知道\Quote{婚姻应受尊重，床第也是无玷的。}\BibleRef{希 13:4}然而，\Person{亚当}即使在被逐出乐园后也只有一个妻子。那受诅咒、沾满血迹的\Person{拉默客}，是\Person{加音}的后裔，他是第一个从一根肋骨造出两个妻子的人；然后，那时种下的再婚的幼苗，被洪水的惩罚完全毁灭了。的确，宗徒在致\Person{弟茂德}书时，因害怕淫乱而被迫认可再婚。他的话是这样：\Quote{所以我愿意年轻的寡妇再嫁，生育儿女，治理家务，不给敌人以辱骂的借口。}但他立即加上这让步的理由：\Quote{因为有些人已转去随从了\Person{撒旦}。}\BibleRef{弟前 5:14}可见，他所提供的不是给站立者的冠冕，而是给跌倒者的援手。如果再婚只是被看作妓院的替代品，它又算什么呢！\Quote{因为有些人，}他写道，\Quote{已转去随从了\Person{撒旦}。}整个事情的要旨是，如果一位年轻的寡妇不能或不愿自制，她宁愿让一个丈夫进入她的床，也不要让魔鬼进入。\par

真是个高尚的选择，只有在优先于\Person{撒旦}时才被采纳！在古代，连\Place{耶路撒冷}也行淫，向每个过路的人张开双脚。\BibleRef{则 16:25}她是在\Place{埃及}初次失身，在那里她的乳房被揉弄。\BibleRef{则 23:3}后来，当她来到旷野，因不耐她的领袖\Person{梅瑟}的迟延，并在情欲之刺的疯狂驱使下说道：\Quote{这就是领你们出埃及地的神。}\BibleRef{出 32:4}她就领受了不良的条例和有害的诫命，藉此她不能存活，\BibleRef{则 20:25}却要因此受罚。那么，当宗徒在另一处论到年轻的寡妇说：\Quote{她们开始对基督放纵情欲时，便想嫁人，她们应受惩罚，因为抛弃了起初的信德，}\BibleRef{弟前 5:11}之后，他对那些将要放纵情欲的人，赐予再婚的不良条例和有害的诫命，难道不奇怪吗？因为他允许再婚的理由，也会证明一个女人嫁给第三个、甚至——如果她愿意——第二十个丈夫是合理的。他显然想向他们表明，他并非那么急于让她们嫁人，而是让她们避免情夫。最亲爱的基督内的女儿，我把这些事记在你心上，并常常重复，为使你忘记背后的，努力面前的。\BibleRef{斐 3:13}你有像你自己一样的寡妇可作典范：在希伯来史中著名的\Person{友弟德}，和福音中著名的\Person{法奴耳}的女儿\Person{亚纳}。她们俩日夜在圣殿里生活，以祈祷和禁食保存了贞洁的宝藏。一位是教会的预象，斩断魔鬼的头；另一位首先在怀中接抱了世界的救主，并向她启示了将要来的神圣奥秘。\BibleRef{路 2:36}最后，我恳求你，将我书信的简短归因于并非缺乏言辞或题材，而是出于深切的谦逊，使我害怕过分长久地强加于陌生人的耳朵，并使我畏惧那些读我文字之人的隐秘评判。\par

\SourceRecord{Translated by W.H. Fremantle, G. Lewis and W.G. Martley.；Nicene and Post-Nicene Fathers, Second Series；Vol. 6.；Edited by Philip Schaff and Henry Wace.；Buffalo, NY: Christian Literature Publishing Co.,；1893.；<http://www.newadvent.org/fathers/3001079.htm>.}

% structure: p0001 source=h1 target=section reason=page-title

\section{第八十封信}

\Emph{\Person{鲁菲努斯}致\Person{玛加略}}\par

\EditorialNote{\Person{鲁菲努斯}在从\Place{白冷}返回罗马后，发表了\Person{奥利振}著作\WorkTitle{论原理}的拉丁文译本。}，\WorkTitle{论原理}。他为此写了序言，即在此处刊印于\Person{热罗尼莫}书信集中。他声称以\Person{热罗尼莫}自己所译并大为赞扬的\Person{奥利振}注释为典范，声明遵循其榜样，意译了该论文晦涩的段落，并意译了该论文晦涩的段落，并删除了那些看似异端而应归咎于窜改者的部分。这篇序言对\Person{热罗尼莫}（虽未提其名）的虚伪赞誉以及对\Person{奥利振}原文的公开操纵，在罗马引起极大困惑（参看第81、83、84封信），并导致了\Person{鲁菲努斯}与\Person{热罗尼莫}之间的争论，该争论在本书绪论中有所描述，并在本系列第三卷中详述。时值公元398年。\par

1. 我知道，为数众多的弟兄出于对圣经知识的热情，曾恳请精通希腊文学的学者将\Person{奥利振}变成罗马人，把他的教导传入拉丁耳朵。其中一位学者，一位亲爱的弟兄和同伴，应\Person{达玛苏}主教之请，从希腊文译出了\Person{奥利振}关于\BibleBook{}{雅歌}的两篇讲道，并为该著作写了一篇雄辩而赞美的序言，这序言不能不激起所有人热切阅读和研究\Person{奥利振}的愿望。他宣称，对这位义人的灵魂而言，雅歌中的话是适用的：\Quote{君王引我进入了内室}；他接着这样说：\Quote{在\Person{奥利振}的其他著作中，他超越了所有前代作者；在处理\BibleBook{}{雅歌}时，他超越了自己。}在他的序言中，他承诺将这些\Person{奥利振}的讲道以及尽可能多的其他作品给予罗马耳朵。他的风格无疑是吸引人的，但我看得出他的目标比单纯的翻译者更为远大。他不满足于传达\Person{奥利振}的话语，而是渴望自己成为老师。至于我，我只是追随他所认可和开创的事业，但我缺乏他那雄辩的言辞来修饰这位伟人的言论。我甚至担心我的不足和拉丁文掌握不够会严重损害这个人的声誉，而这位作者曾恰当地称他为教会中仅次于宗徒的知识与智慧导师。\par

2. 我常反复思量此事，保持沉默，拒绝听从弟兄们的请求——他们时常催促我承担这项工作。但是，最忠实的弟兄\Person{玛加略}，你的坚持如此之大，以至于连无能也无法抵抗。为躲避你不断施加的纠缠，我违背初衷做出了让步；但条件如下：尽可能按照我先辈所订的翻译规则进行翻译，特别是遵循我刚才提到的那位作者的做法。他已经将\Person{奥利振}的七十多篇讲道以及几部关于宗徒的注释译成了拉丁文；在这些译文中，但凡希腊原文有绊脚石之处，他都在译文中予以平息，并如此修改了用语，以至于拉丁文作者找不到任何与我们的信仰相冲突的词。因此，我打算追随他的脚步，即便不展现同样有力的雄辩，也至少遵循相同的规则。我不会重现\Person{奥利振}著作中与其其他声明不一致或矛盾的段落。这些矛盾的原因，我在\Person{庞菲利乌斯}所撰的奥利振著作辩护中已更为详细地陈述于你，我并附加了一篇自己的短文。我已经给出我认为清楚的证据，证明他的著作在许多地方被异端者和心怀恶意之人窜改，尤其是你现下敦促我翻译的这些著作。\OriginalTerm{论原理}{\GreekText{περὶ} \GreekText{᾿Αρχῶν}}，即原理或权能之书，实际上在其他方面极其晦涩难懂。因为它们所讨论的主题是哲学家们穷尽一生也无法发现任何东西的。在处理这些主题时，\Person{奥利振}已尽力使对造物主的信仰和对受造物的理性说明服务于宗教，而非像哲学家那样服务于不敬神。因此，在他的著作中，但凡我遇到关于天主圣三的陈述与他处所忠实陈述的观点相矛盾之处，我均删除了这些段落，视为被窜改和误导的，或者代之以我见他屡次坚持的观点。再者，凡他因急于推进主题而写得简略或晦涩、依赖读者的技能和智慧之处，我为了使行文清晰，已试图通过添加我在他其他著作中读到的任何更明确的陈述来加以解释。但我未曾添加任何自己的东西。所用词语可见于他其他著作：是他的，而非我的。我在此提及这一点，是为了打消吹毛求疵者再次挑剔的借口。但这些乖谬好争议之人最好小心自己的所作所为。\par

3. 同时，我已承担起这项伟大任务——如果天主允诺你的祈祷——不是为了堵住诽谤者之口（除了天主也许无人能做到），而是为了给那些渴望者提供进步的知识。\par

我在天主圣父、圣子、圣神面前，恳求并责令每一位将要阅读或抄写我这书的人，凭他对将来国度的信仰，凭从死者中复活的奥迹，凭那\BibleQuote{为魔鬼和他的使者预备的}\BibleRef{玛 25:41}永火，正如他盼望不永远承受那\BibleQuote{有哀号和切齿}\BibleRef{玛 22:13}的地方，以及那\BibleQuote{他们的虫子不死，火也不灭}\BibleRef{谷 9:44}的地方一样，他不得增添所写的，不得删减，不得插入，不得更改，而应将他的抄本与据以制作的副本核对，逐字校正，并正确标点。凡未经适当校正和标点的手稿，他必须弃绝；因为否则，因缺乏标点而导致的文本困难，将使原本晦涩的论证对阅读者来说更加晦涩。\par

\SourceRecord{Translated by W.H. Fremantle, G. Lewis and W.G. Martley.；Nicene and Post-Nicene Fathers, Second Series；Vol. 6.；Edited by Philip Schaff and Henry Wace.；Buffalo, NY: Christian Literature Publishing Co.,；1893.；<http://www.newadvent.org/fathers/3001080.htm>.}

% structure: p0001 source=h1 target=section reason=page-title

\section{第八十一封书信}

致\Person{鲁菲努斯}\par

\EditorialNote{一封由热罗尼莫写给鲁菲努斯的友好劝诫信，以回应他收到的\WorkTitle{论原理}的译本。}（参见前一封信）。它首先寄给\Person{潘玛丘}，后者背信弃义地将其扣押，从而使两位朋友和解的一切希望都化为泡影。这封信的日期是公元399年。\par

1. 你在罗马逗留了一段时间，这一点从你本人的文字中可以看出来。但我深信，若不是因你母亲的哀伤使你踌躇，渴望见到灵性双亲的愿望必会引你返回故乡——惟恐你远在他方尚难忍受的悲痛，在家乡会更不堪承受。\par

至于你抱怨说，人们只听从激情的驱使，不肯赞同你我的判断；主是我的良心的见证：自我们和好以来，我心中从未怀有怨恨去伤害任何人；恰恰相反，我竭尽全力防止任何偶然事件被归咎于恶意。但我又能做什么呢？只要每个人都以为他有权按自己的方式行事，并认为散播诽谤是在报复而非制造中伤？真正的友谊绝不应当隐瞒其想法。\par

关于\OriginalTerm{论原理}{\GreekText{περὶ} \GreekText{᾿Αρχῶν}}诸书的简短序言已寄给我，我从文风认出那是你的手笔。你写作的意图你最清楚；但即使是愚人也能看出它必然会被如何理解。隐晦地，更确切地说，公开地，我便是针对的对象。我自己也时常虚构争议来进行演说练习。因此我现在可以重拾这一陈旧的技巧，用你的方式称赞你。但我绝不愿模仿我所责难你的东西。事实上，我已克制自己的感情，以致我不对你提出任何控诉；尽管受了伤害，我也不愿反过来伤害朋友。但下次，若你想跟随某人，请满足于你自己的判断。我们所追求的对象或是善，或是恶。若是善，便无需他人帮助；若是恶，多人同罪并非借口。我宁愿以朋友的方式向你进行劝诫，而不愿公开发泄我因所遭受的委屈而产生的愤怒。我希望你明白，当我与人和好时，我便成为他真诚的朋友，而不会——借用\Person{普劳图斯}的比喻——一只手递面包，另一只手握着石头。\par

2. 我的兄弟\Person{保利尼安}尚未从家中返回，我想你会在\Place{阿奎莱亚}\Person{克罗玛修}可敬的主教家中见到他。我也差遣可敬的司铎\Person{鲁菲努斯}因公务经罗马前往\Place{米兰}，并请他向你传达我的情意和敬意。我也向其余朋友传达同样的信息；免得如宗徒所说：\BibleQuote{你们若彼此相咬相吞，你们要小心，免得彼此消灭。}\BibleRef{迦 5:15}剩下的只是你和你朋友们，因不愿冒犯那些不容忍的人，而显示你们的节制。因为你很难找到像我这样的人。没有多少人会对虚假的颂词感到满意。\par

\SourceRecord{Translated by W.H. Fremantle, G. Lewis and W.G. Martley.；Nicene and Post-Nicene Fathers, Second Series；Vol. 6.；Edited by Philip Schaff and Henry Wace.；Buffalo, NY: Christian Literature Publishing Co.,；1893.；<http://www.newadvent.org/fathers/3001081.htm>.}

% structure: p0001 source=h1 target=section reason=page-title

\section{第八十二封信}

\Emph{致\Place{亚历山大}主教\Person{德奥斐洛}}\par

\EditorialNote{在他之前尝试（见第六十三封信）两年后，德奥斐洛再次写信给热罗尼莫，敦促他与若望·耶路撒冷和解。热罗尼莫回答说他最渴望的就是和平，但这必须是真正的和平，而不是空洞的休战。他非常痛苦地谈到若望，指控若望策划将他从巴勒斯坦流放。他还谈到他兄弟保利尼安被祝圣一事（见第五十一封信），并引用普瓦捷的依拉略的例子为自己翻译奥利振的注释辩护。这封信应与本卷中的\WorkTitle{驳若望·耶路撒冷}一文对照阅读。日期为公元399年。}\par

你的来信显示出你拥有主的遗产，当他即将回归圣父时，对宗徒们说：\BibleQuote{我把平安留给你们，我将我的平安赐给你们。}\BibleRef{若 14:27}并拥有那话语所描述的福分：\BibleQuote{缔造和平的人是有福的。}\BibleRef{玛 5:9}你如同父亲般劝勉，如同师傅般教导，如同主教般训诫。你来到我这里，不是带着棍棒和严厉，而是本着仁慈、温和与良善的精神\BibleRef{格前 4:21}。你的开场白呼应了基督的谦卑，祂拯救人不是藉着雷霆和闪电\BibleRef{希 12:18}，而是作为马槽中啼哭的婴孩和十字架上静默受苦者。你曾读到那预言论及一位祂的预像者：\Quote{上主，求祢记念达味和他的所有温良。}你深知这预言后来如何在祂身上应验。祂说：\BibleQuote{你们跟我学习，因为我是良善心谦的。}\BibleRef{玛 11:29}你从圣经中引用了许多赞美和平的段落，你像蜜蜂一样飞过圣经的花丛，你用巧妙的雄辩采集了所有甘甜且有利于和谐的话语。我本来已在追求和平，但你使我加快了步伐；我的帆已经为航行张起，但你的劝勉为它们吹来了更强的风。我畅饮和平的甘泉，不是勉强和厌恶，而是热切地张开口。\par

2. 但我能做什么呢？我只能渴望和平，却没有能力实现它。即使这份渴望能在天主面前获得赏报，它的徒劳也必然使渴望者悲伤。当宗徒说：\BibleQuote{尽你们所能，与众人和平相处。}\BibleRef{罗 12:18}他很清楚和平的实现取决于双方的同意。先知确实喊道：\Quote{他们说平安，平安，却没有平安。}用行动颠覆和平却用言辞宣称和平并不难。指出和平的好处是一回事，努力争取和平是另一回事。人们的话语可能都是为了合一，但他们的行动却可能实施束缚。我和别人一样渴望和平；我不仅渴望，而且祈求。但我想要的和平是基督的和平；真正的和平，没有怨恨的和平，不涉及战争的和平，不减少对手而能团结朋友的和平。我怎能将控制称为和平？我必须按其正确名称称呼事物。哪里有仇恨，就让人谈论仇隙；哪里有相互尊重，那里才可谈论和平。至于我，我既不撕裂教会，也不脱离教父们的共融。可以说，从我幼年起，我就被公教的乳汁喂养；没有人能比一个从未做过异端者的人更善尽教会成员之责。但我不知道没有爱的和平，也不知道没有和平的共融。在福音中我读到：\BibleQuote{若你带礼物到祭台前，在那里想起你的弟兄有什么对你不满之处，就把礼物留在祭台前，先去与你的弟兄和好，然后再来献礼物。}\BibleRef{玛 5:23}既然我们连自己的礼物都不能献，除非与弟兄和好；那么，如果我们心中怀有敌意，又怎能领受基督的身体呢？如果我怀疑主祭者的爱德，又怎能心安理得地走近基督的圣体圣事并回应阿们呢？\par

3. 请耐心听我说，不要把真话当作谄媚。是否有人不愿与你共融？是否有人在你伸出手时转过脸去？是否有人在神圣的宴席上给你\Person{犹达斯}之吻？你接近时，隐修士们不是战栗而是欢喜。他们跑来迎接你，离开沙漠中的洞穴，情愿以谦卑来征服你。是什么迫使他们出来？岂不是对你的爱？是什么聚集了沙漠中散居的居民？岂不是他们对你的尊敬？一位父亲应当爱他的儿女；不仅是父亲，一位主教也应当被他的儿女所爱。两者都不应被畏惧。有一句古老的谚语：\Quote{人所畏惧的，他必憎恨；所憎恨的，他必欲其死。}因此，尽管圣经为青年人以敬畏为智慧的开端，\BibleRef{箴 1:7} 它又告诉我们：\BibleQuote{完美的爱驱除恐惧。}\BibleRef{若一 4:18} 你不向他们索取服从，因此隐修士们服从你。你献上亲吻，因此他们低头。你表示自己是一名普通士兵，因此他们立你为将军。这样，你从众人之一变为众人之上。自由一旦被压制，就容易被激起。没有人能从自由人那里得到比不强迫他为奴者更多的东西。我知道教会的规章；我知道她的圣职人员属于何等品级；我从人和书籍中，日日至今，学习并收集了许多事情。温和的\Person{达味}的王国，很快就被一个用蝎子鞭责打百姓、自以为他的手指比父亲的腰还粗的人所分裂。\BibleRef{列上 12:10} 罗马人民甚至连国王的傲慢都不能容忍。\Person{梅瑟}是以色列军队的领袖；他对\Place{埃及}降下十灾；天空、大地、海洋都听从他的命令：但他被描述为\BibleQuote{极为谦和，超过所有当时}那时\BibleQuote{在地上的人。}\BibleRef{户 12:3} 他保持了四十年的权柄，因为他以温和与谦和缓和了职务的傲慢。当百姓用石头砸他时，他为他们祈求；\BibleRef{出 17:4} 不仅如此，他宁愿自己的名字从天主的册上被涂抹，也不愿托付给他的羊群灭亡。\BibleRef{出 32:31} 他试图效仿那位他知道会把迷羊扛在肩上的牧者。\BibleQuote{善牧}——这是主自己的话语——\BibleQuote{为羊舍命。}他的一位门徒，为了他的弟兄、按血统是他的同胞以色列人的缘故，情愿被弃绝于基督。如果\Person{保禄}可以宁愿自己失丧，好使那迷失的不致失丧，那么善良的父母岂不更不应惹儿女发怒，\BibleRef{弗 6:4} 或是因过于严厉而苦待那些本性温和的人呢？\par

4. 一封信的篇幅迫使我自己节制；否则，愤慨会使我冗长。在一封其作者视为调解性的、但在我看来充满恶意的信函中，我的对手承认我从未诽谤他或指责他异端。那么，他为何散布谣言，说我染上了那可怕的疾病并背叛教会，从而诽谤我？为何他如此轻易地放过真正的攻击者，却如此热切地伤害我这个从未伤害他的人？在我兄弟被祝圣之前，他对自己与\Person{埃皮法尼乌斯主教}之间的教义分歧只字未提。那么，是什么\Quote{迫使}他——我用他自己的话——公开争论一个尚未有人提出的论点？像你这样充满智慧的人深知此类讨论的危险，在这种情况下沉默是最安全的做法；除非，确实，在某些场合，必须处理重大事件。要将最博学的作家们在多卷著作中详细讨论的所有主题压缩成一篇讲道——如他所夸耀的——需要何等的才能与口才！但这与我无关：那应由讲道的听众和信函的作者来注意和意识到。但就我而言，他应主动承认我并未对他提出指控。我不在场，也没有听到讲道。我只是众人中的一个，甚至几乎不算其中之一；因为在别人呼喊时，我保持了沉默。让我们让被告与原告对质，并将信任赐予那侍奉、生活和教义看来最好的人。\par

5. 你看不出吗？我对许多事闭目不见，对其他的事也只是以最粗略的方式触及，暗示我所猜想的，而不是说出我所想的。\par

我理解并赞赏你的策略；你为了教会的和平，在听到\Place{塞壬}的歌声时捂住耳朵。而且，自幼在圣学中受训的你，确切知道你使用的每个词句的含义。你故意使用模棱两可的措辞和精心平衡的句子，以免定他人的罪，或否认我们。但是，纯洁的信德和坦白的宣认并非寻求诡辩或迂回。简单相信的，必须以同样的简单来宣认。就我而言，我愿呼喊——即使是在\Place{巴比伦}的刀剑与烈火中，\Quote{为何答案回避问题？为何没有坦率、直接的宣告？}自始至终，都是退缩、妥协、含糊：仿佛他在谷粒尖上行走。他渴望和平的热血沸腾；却不肯给出一个直接的答案！别人可以自由地侮辱他；因为当他受辱时，他不敢报复。与此同时，我保持沉默：目前，我愿让人认为我太忙，或无知，或害怕；因为我若指控他，他将会如何对待我——如果我称赞他（正如他本人承认我所做的）时，他却暗中诽谤我？\par

6. 他的整封信与其说是对他信德的阐述，不如说是一堆针对我本人的诽谤。没有任何那种人与人之间可以互用而不显谄媚的礼节，他一再提起我的名字，嘲弄它，把它丢来丢去，仿佛我被从生命册上抹去。他以为他写的信把我打得鼻青脸肿；而我，自幼便被关在隐修小室中，素以成为一个好人而非看似好人为目标，竟为他所对准的这些琐事而活。不错，我们中有些人他提及并表尊重，但只是为了以后更痛地伤害我们。好像我们就没有可供揭露的秘密似的！他的一项指控是，我们允许了一个奴隶被祝圣。而他自己也有同一阶层的圣职人员，他一定读过关于\Person{敖乃息摩}的事，他因\Person{保禄}在狱中重生，从奴隶变成了执事。然后他又抛出，这个奴隶是一个告密者；而为了不必证明指控，他声称只是听说的！唉，如果我选择重复众人的闲话并听信搬弄是非者，他早就该知道，我也知道人尽皆知的事，也和别人一样听过同样的故事。他还宣称，祝圣给这个奴隶是作为他散布诽谤的奖赏。这样的狡猾和机巧难道不令人惊恐吗？对于如此压倒性的雄辩，还有什么回答？散布诽谤和遭受诽谤，哪个更好？指控一个你之后可能希望得到其爱的人，还是宽恕一个罪人？一个告密者被任命为执政官，比被任命为市政官，更容易容忍吗？他知道我略过什么、说什么；我自己听到什么，以及因对基督的恐惧，我或许拒绝相信什么。\par

7. 他指控我将\Person{奥力振}翻译成了拉丁文。在此我并非独行，因为精修者\Person{依拉略}也做了同样的事，我们两人都一致：我们翻译了所有有用的部分，而删除了所有有害的部分。让他自己读一读我们的译本，如果他有能力（既然他常与意大利人交谈、每日交往，我想他不会不懂拉丁文）；或者，如果他不能完全理解，就让他用他的译员，那样他就会知道，我基于他指控我的工作，应得的只有赞扬。因为，我一直承认\Person{奥力振}作为圣经诠释者和批评家的伟大功绩，但我一贯否认他学说的真实性。那么，是我把他释放到人群中的吗？我是其他像他那样的讲道者的担保人吗？不，因为我知道宗徒和所有其他讲道者之间的区别。前者总是说真理；后者作为人，有时会犯错。为\Person{奥力振}辩护，承认他的过失然后用别人也犯过类似错误来原谅他，这无疑是奇怪的辩护！好像当你不敢公开为一个人辩护时，你可以通过将他的错误推给许多别人来保护他！至于他所说的\Person{奥力振}的六千卷书，任何人都不可能读过从未写成的书；而对我来说，我更倾向于认为这个谎言出自那个声称听到的人，而非据说讲述的人。\par

8. 他又断言，我兄弟是所发生分歧的原因，他满足于待在隐修小室，把圣职视为繁重而非荣耀。尽管直到今日，他以不诚恳的和平誓言敷衍我们，他却通过告诉西方的主教们，说一个几乎还是个孩子的青年在他的教区中被祝圣为\Place{白冷}的司铎，从而在他们的心中引起了骚动。如果这是事实，所有\Place{巴勒斯坦}的主教都应该知道。因为可敬的教宗\Person{厄比法尼}的隐修院——称为老隐修院——我兄弟在那里被祝圣为司铎，位于\Place{厄肋乌特洛波里}地区，而不是\Place{厄里亚}地区。再者，他的年龄为您的圣座所熟知；既然他现在已三十岁，我认为在此点上无可指责。事实上，这个特定的年龄因基督所取人性的奥迹而标志着圆满和成熟。让他记起旧约法律，他会看到，二十五岁之后肋未人便可被选为司祭；或者，如果他更愿意在此段按希伯来文，他会发现司祭候选人必须三十岁。他竟敢说\BibleQuote{旧事已过；看，一切都成了新的。}\BibleRef{格后 5:17}，让他听听宗徒对\Person{弟茂德}说的话：\BibleQuote{不要让人小看你的年轻。}\BibleRef{弟前 4:12}。当然，当我的对手自己祝圣为主教时，他也不比我兄弟现在大多少。而如果他认为年轻对主教无妨碍，但对司铎却有妨碍，因为年轻的长老是自相矛盾，那么我问这个问题：为什么他自己祝圣了这个年龄或更年轻的司铎，而且还在别人的教会中服务？但如果除非我兄弟同意屈服并放弃祝圣他的主教，他就不能与我兄弟和平相处，那么他清楚表明，他的目标不是和平而是报复，除非他能施加他现在威胁的所有惩罚，否则他不会满足于安宁平静。即使他自己祝圣了我兄弟，对后者也毫无区别。他如此热爱隐居，以至在那时也会继续安静生活，不履行其职务。而如果主教认为应当因此撕裂教会，那么除了对所有献祭者应得的尊敬之外，他对他便无所亏欠。\par

9. 关于他那冗长的自我辩护，我更应称之为对我的攻击，就说到这里。在这封信里，我只简短粗略地回应他，好让他从我所说的话中看出我没有说出的话，并知道作为一个人，我是理性的动物，能够理解他的机敏；我并非如此迟钝或愚钝，以至于只抓住他话语的声音而不领会其含义。我现在请求您原谅我的恼怒，并允许我这样认为：如果回击是傲慢的，那么提出无根据的指控更加傲慢。然而，我的回答只是暗示了我本可以说出口的话，而非实际说出口。人们为何要在远方寻求和平？又为何希望凭借命令来强制执行和平？让他们显露出缔造和平者的身份，和平就会立刻随之而来。他们为何用您的圣座之名来恐吓我们？而您的信——与他们的严厉和威胁性言辞形成奇怪对比——只散发着和平与温和？因为我知道，\Person{伊西多尔}司铎带给我的您的这封信，确实是为了和平与和谐，而这些虚伪宣称渴望和平的人却拒绝把它交给我。让他们选择他们喜欢的任何一种说法吧。要么我是一个好人，要么我是一个坏人。如果我是好人，就让我安静；如果我是坏人，他们为何渴望与坏人为伍？的确，我的对手已经通过经验学到了谦卑的价值。他如今撕裂那些他以前自愿将分离开的东西结合在一起的事物，这证明他现在将曾经联合的东西分离，是受他人唆使而做的。\par

10. 最近，他寻求并得到了一项针对我的流放令，我但愿他能执行它，这样，正如意愿被归算为行为一样，我也不仅在意愿上，而且在事实上戴上流放的荣冠。基督教会是通过倾洒自己的血而非他人的血，通过忍受凌辱而非施加凌辱而建立的。迫害使它成长；殉道给它加冕。或者，如果与我共处的基督徒们在好施严厉方面是独一无二的，只知道迫害而不知忍受迫害，那么这里有犹太人，有持各种虚假学说的异端者，尤其是最污秽的摩尼教。为什么他们不敢对他们说一句话？为什么唯独他们想把我流放？难道只有我这个与教会共融的人，是唯一可以被说成撕裂教会的人吗？我问您，要么他们像我一样驱逐这些人，要么如果他们容纳这些人，也应当容纳我，这难道不是一个公平的要求吗？尽管如此，他们通过流放来尊荣人，因为这样做使他们与异端者的团体分离。可耻的是，一位隐修士竟威胁隐修士们，并获取针对他们的流放令；而且是一位自称占据宗座的隐修士。但隐修团体并不屈服于恐怖手段：它宁愿将脖子暴露在即将落下的刀剑之下，也不愿让双手被束缚。难道每一位隐修士不都是背井离乡的人吗？难道他不都是从整个世界流放的人吗？哪里有需要公共权力、敕令的花费、为了得到它而往返于大地的旅程呢？只消他用小指碰我一下，我就会自动去流放。\BibleQuote{大地和其中的万物属于上主。}基督不被局限于任何一处。\par

11. 此外，当他写道，虽然我似乎与他断绝了共融，但实际上我藉着你们和罗马教会与他共融时，他不必舍近求远，因为我在这里巴勒斯坦也以同样的方式与他相连。唯恐这仍然显得遥远，就在这白冷村，我尽可能与他的长老们共融。由此可见，个人的懊恼不可被当作教会的原因，一个人的怒气，甚或被他煽动的几个人的怒气，也不应被称作教会的不悦。因此，我现在重述我在信首所说的话：我本人渴望基督的和平，我祈求和谐，我请求你们劝告他，不要强求和平，而要立志和平。愿他因过去对我施加的侮辱所导致的痛苦而满足。愿他藉一点新的爱德抹去旧的创伤。愿他显出他从前自愿尊重我时的本色。愿他的话语不再染上从他人心中流出的苦涩。愿他做他自己所愿的事，而非他人强迫他所愿的事。或者，作为教宗，他对所有人一视同仁地行使权威；或者，作为宗徒的追随者，他为所有人的救恩而服侍众人。\BibleRef{格前 9:19} 若他显出这般模样，我情愿甘心让步，张开双臂；他必会视我为朋友和亲人，并会察觉在基督内我顺服他如同顺服诸位圣人一样。宗徒写道：\Quote{爱德含忍慈祥；爱德不嫉妒……不矜夸……凡事包容，凡事相信。} \BibleRef{格前 13:4} 爱德是诸德之母，宗徒论及信、望、爱的话，\BibleRef{格前 13:13} 如同那不易折断的三股绳。\BibleRef{训 4:12} 我们信，我们望，并藉着信德和望德，在爱德的联系中彼此结合。\BibleRef{哥 3:14} 正是为这些德行，我和其他人才离弃家园，也正是为这些，我们愿在田野间独自和平相处，毫无争执；向基督的主教们——只要他们宣讲正确的信仰——致以应有的敬意，非因畏惧他们如主人，乃因尊敬他们如父亲，同时对主教们也以主教之礼相待，但不愿在主教权威的荫蔽下被迫服侍那些我们不愿服从的人。我并未如此心高气傲，以致不晓得对基督的司铎们当尽的义务。因为凡接待他们的，不是接待他们，乃是接待那差遣他们的主，他们就是祂的主教。\BibleRef{若 13:20} 但愿他们满足于自己所拥有的尊荣。愿他们知道自己是父亲而非主人，尤其对于那些蔑视世俗野心、视平安与安息为至上者而言。愿全能的基督天主因你们的祈祷，使我和我的对手在真诚而坚实的互相尊重中合一，而非在虚假空洞的和平中，以免我们彼此相咬相吞，同归于尽。\BibleRef{迦 5:15}\par

\SourceRecord{Translated by W.H. Fremantle, G. Lewis and W.G. Martley.；Nicene and Post-Nicene Fathers, Second Series；Vol. 6.；Edited by Philip Schaff and Henry Wace.；Buffalo, NY: Christian Literature Publishing Co.,；1893.；<http://www.newadvent.org/fathers/3001082.htm>.}

% structure: p0001 source=h1 target=section reason=page-title

\section{\WorkTitle{第八十三封书信}}

\Emph{来自\Person{帕玛丘斯}与\Person{奥克亚努斯}}\par

\EditorialNote{一封来自\Person{帕玛丘斯}和\Person{奥克亚努斯}的信，他们在信中表达了对\Person{鲁菲努斯}所译\Person{奥力振}论著《论首要原理》（见第八十封书信）的困惑，并请求\Person{热罗尼莫}为他们逐字翻译这部作品。写于公元399年或400年。}\par

1. \Person{帕玛丘斯}与\Person{奥克亚努斯}致司铎\Person{热罗尼莫}，安康。\par

一位可敬的弟兄给我们带来了一些书页，其中包含某人对\Person{奥力振}一篇题为\WorkTitle{\OriginalTerm{论首要原理}{\GreekText{περὶ} \GreekText{ἀρχῶν}}}的论著的拉丁文翻译。这些书页包含许多使我们可怜的理智感到困惑、在我们看来不合乎公教信仰的内容。我们还怀疑，为了替作者辩护，他著作中的许多段落已被删除；这些段落若保留下来，本可清楚证明其教导的不虔敬性质。因此，我们恳请您惠予关注此事，这不仅有益于我们自己，也有益于城中所有居民；我们请求您用自己的语言出版上述\Person{奥力振}的著作，完全按照作者本人写就的样子；我们还希望您揭露其辩护者所引入的增改之处。您还要驳斥并推翻我们寄给您的书页中所有出于无知或违背公教信仰的言论。该书的作者在其著作序言中，以极其巧妙的方式，虽未提及您，却暗示他不过是完成了您自己曾许诺的工作，从而间接表明您与他意见一致。请消除人们无法避免的疑虑，并驳倒您的攻击者；因为如果您忽视他的暗示，人们会认为您承认其真实性。\par

\SourceRecord{Translated by W.H. Fremantle, G. Lewis and W.G. Martley.；Nicene and Post-Nicene Fathers, Second Series；Vol. 6.；Edited by Philip Schaff and Henry Wace.；Buffalo, NY: Christian Literature Publishing Co.,；1893.；<http://www.newadvent.org/fathers/3001083.htm>.}

% structure: p0001 source=h1 target=section reason=page-title

\section{第八十四封信}

致\Person{帕玛丘}与\Person{奥西亚努斯}\par

\EditorialNote{一封平静的信，热罗尼莫在其中界定并辩护自己对奥力振的态度，但过度淡化了他早期对奥力振的热情。他以西彼廉钦佩戴尔都良的方式钦佩奥力振，但绝不采纳他的错误。然后他描述了自己的学习，并叙述了他对阿波利纳里斯、狄迪穆斯和一位名叫巴-阿尼纳的犹太人的亏欠。信的其余部分涉及奥力振的错误、其著作的文本状况，以及殉道者庞菲鲁斯所撰写的颂词（热罗尼莫毫无充分理由地攻击其真实性）。信的日期是公元400年。}\par

1. 你寄来的篇幅令我既受赞誉又感惭愧；因为它们虽称赞我的能力，却质疑我信仰上的真诚。然而，在\Place{亚历山大港}和\Place{罗马}，甚至可以说在整个世界，善人们都已习惯以同样的方式提及我的名字，他们对我评价止于这一点：他们不能容忍成为异端者却未将我列入其数。我将搁置人身攻击，只回应具体指控。因为以辱骂还辱骂、以攻击回击对手，对案情毫无助益。我们受命不以恶报恶\BibleRef{得前 5:15}，反要以善胜恶\BibleRef{罗 12:21}，饱受侮辱，并转脸颊给打者\BibleRef{玛 5:39}。\par

2. 有人指控我曾有时赞扬\Person{奥力振}。如果我没有弄错，我仅在两处这样做过：在（致\Person{达玛苏}的）\BibleBook{}{雅歌}讲道词的简短序言中，以及在我的\WorkTitle{希伯来名字}一书的前言中。在这些段落中，是否涉及教会的教义？是否谈及圣父、圣子及圣神？或谈论肉身的复活？或灵魂的状况与质料？我仅赞扬了他译文的简洁和注释，丝毫不涉及教会的信仰或教义。只涉及伦理，藉清晰的解释驱散寓言的迷雾。我赞扬了注释者而非神学家，有才智的人而非信徒，哲学家而非宗徒。但如果人们想知道我对\Person{奥力振}的真实判断，让他们阅读我关于\BibleBook{训道}{篇}的注释，让他们通读我关于\BibleBook{厄弗所}{书}的三卷书：他们将会看到我一贯反对他的教义。赞扬一种体系以至于赞同其亵渎，是何等愚蠢！有福的\Person{西彼廉}以\Person{戴尔都良}为师，其著作可证；然而他虽欣喜欢迎这位博学且热忱的作家的才能，却并未跟随他追随\Person{孟他努}和\Person{马克西米拉}。\Person{阿波利纳里斯}撰写了一部极有分量的反对\Person{波菲利}的著作，\Person{欧瑟比}则编纂了一部出色的教会史；然而前者割裂了基督降生成人的人性，后者则是亚略不敬神明的公开拥护者。\Person{依撒意亚}说：\Quote{祸哉，那些称恶为善、称善为恶，以苦为甜、以甜为苦的人。}\BibleRef{依 5:20}我们不得贬低反对者的美德——如果他们有任何值得称赞之处——但也不得赞扬朋友的缺陷。每个案件必须根据其本身的是非来判断，而不是参考相关的人。\Person{路奇里乌斯}因诗句不匀称而受到\Person{贺拉斯}正当的抨击，但他也因其机智和迷人的风格而同样受到正当地赞扬。\par

3. 在我年少时，我被极大的求知欲所驱使，但我并不像某些人那样狂妄到自学。在安提约基雅，我常听劳狄刻雅的阿波利纳里讲学，并出席他的讲座；然而，虽然他教导我圣经，但我从未接受他那些有争议的关于经文含义的学说。最后，我的头发斑白，看起来更像一位老师而非门徒。但我仍前往亚历山大，听了狄迪莫斯的课。我对他深表感激：因为我不知道的事，我向他学习；而我已经知道的，也没有忘记。他的教导如此出色。人们以为我现在已终止了学习。然而，我再次来到耶路撒冷和白冷。我费了多少周折和花费，才请到巴兰尼纳斯在夜晚掩护下教导我。因为他害怕犹太人，在我面前活像个尼苛德摩的翻版。\BibleRef{若 3:2} 所有这些，我常在著作中提到。阿波利纳里和狄迪莫斯的学说是相互矛盾的。两位领袖的阵营必定将我从不同方向拉扯，因为我承认两人都是我的老师。如果憎恶某些人、厌恶某个种族是合宜的，那么我奇怪地对那些受割礼者怀有反感。因为直到今日，他们仍在撒殚的会堂里迫害我们的主耶稣基督。\BibleRef{默 2:9} 然而，难道有人能因我有一位犹太老师而责怪我吗？有人竟敢用我写给狄迪莫斯并称他为老师的信来反对我吗？这似乎是一桩大罪，我一个门徒竟称一位年高博学的人为老师。然而，当我请求看看那封长久以来被扣留以诋毁我的信时，信中除了礼貌的措辞和一些问候语外，别无他物。这样的指控是愚蠢且轻浮的。更切题的做法是展示一段我为异端辩护或赞扬奥力振某些邪恶学说的文字。在以撒意亚描写两个色辣芬呼喊的部分，\BibleRef{依 6:2} 他将这两个解释为圣子和圣神；但我不曾将这可憎的解释改为指两个盟约吗？我手头有那本二十年前出版的书。在我的许多著作中，尤其是在注释中，我抓住机会攻击了这位异教大师。如果我的对手声称我比任何人更致力于收集奥力振的著作，我回答说，我只希望我能拥有所有神学家的著作，以便通过勤奋研究它们来弥补我自己才智的迟钝。我承认我收集了他的著作；但因为我知道他所有的著作，我并不追随他的错误。我作为基督徒向基督徒说话：请相信一个考验过他的人。他的学说是毒害人的，不为圣经所知，甚至是对圣经的暴力。我读过奥力振，我再重复，我读过他；如果读他是罪行，我认罪；确实，这些亚历山大的著作耗尽了我的钱财。如果你们相信我，我从未是奥力振主义者；如果你们不相信我，我现在已经不再是一个了。但即便如此仍不能令你们信服，你们将迫使我在自卫中写作反对你们所喜爱的人，以致如果你们不相信我否认他，你们就会在我攻击他时相信我。但我发现，当我犯错时比当我改过时更容易得到信任。这并不奇怪，因为那些自以为是朋友的人以为我与他们在其体系的奥秘中是同门。他们猜想我不愿在那些他们看来如畜生和泥土般的人面前宣扬隐秘的教义。因为他们认为珍珠不应轻易投在猪前，圣物不应给狗。\BibleRef{玛 7:6} 他们赞同达味的话：\Quote{我将你的话藏在我心里，免得我得罪你；} 并在另一处称义人是\Quote{与邻舍说实话}，即与那些\Quote{信仰之家}的人。\BibleRef{迦 6:10} 从这些经文他们得出结论，我们这些尚未入门的人应被告知谎言，免得仍是吃奶的婴孩，被过于坚固的食物噎住。现在，伪誓和谎言如何进入他们的奥秘，并成为他们之间的联结，从奥力振的《杂集》第六卷中看得最清楚，在那卷中他将基督教教义与柏拉图的观念调和起来。\par

4. 那么我该做什么？否认我赞同\Person{奥力振}的看法？他们不会相信我。发誓说我不是？他们会嘲笑我，说我撒谎。我将做一件他们害怕的事。我将公开他们的神圣仪式和奥秘，揭露他们用以欺骗像我这样单纯之人的诡计。也许，尽管他们在我否认时不信任我的声音，但他们可能会在我指控时相信我的笔。他们特别担心一件事，那就是他们的著作有一天会被用作反对他们导师的证据。他们随时准备宣誓作出声明，然后又以同样虚假的誓言否认。当被要求签名时，他们耍花招，寻找借口。一个说：\Quote{我不能谴责别人没有谴责过的东西。}另一个说：\Quote{教父们对此没有作出决定。}他们就这样诉诸世界的判断，以推迟必须同意一项谴责。另一个人更肯定地说：\Quote{我怎能谴责\Place{尼西亚}大公会议没有触及的人呢？因为谴责了\Person{亚略}的大公会议，如果它不赞同\Person{奥力振}的教义，也一定会谴责他。}换句话说，他们有责任一次性用同一剂药方治愈教会所有的疾病；按同样的推理，我们必须否认圣神的尊威，因为那次大公会议没有提到祂的本性。但问题涉及的是\Person{亚略}，不是\Person{奥力振}；是圣子，不是圣神。大公会议的主教们宣布他们坚持当时被否认的信理；对于无人提出的难题，他们未置一词。然而，他们暗中打击了\Person{奥力振}，视其为亚略异端的根源：因为，在谴责那些否认圣子与圣父同体的人时，他们既谴责了\Person{亚略}，也谴责了\Person{奥力振}。按照这些人的立场，我们无权谴责\Person{瓦伦丁}、\Person{马尔西翁}、弗里吉亚派或\Person{摩尼}，这些人没有一个被\Place{尼西亚}大公会议点名，然而无疑他们在时间上早于那次大公会议。但当他们发现自己被迫要么签字认同，要么离开教会时，你会看到一些奇怪的扭曲。他们修饰言辞，重新排列，使用含混的表达；以便可能的话，既持守我们的宣认，也持守反对者的宣认，被毫无区别地称为异端者和天主教徒。仿佛不是在同样的精神下，\Person{德尔斐的阿波罗}（有时被称为\Person{洛西亚斯}）向\Person{克罗伊斯}和\Person{皮洛士}发出神谕；用相似的手法欺骗了两个时间上相隔甚远的人。为了清楚说明我的意思，我将举几个例子。\par

5. 他们说，我们相信身体的复活。这一宣认，只要真诚，就无可指摘。但因为有天上的形体，也有地上的形体\BibleRef{格前 15:40}，且稀薄的空气和以太按其本性都称为形体，所以他们使用\Quote{形体}一词代替\Quote{血肉}，以便正统之人听到他们说\Quote{形体}时以为是指血肉，而异端者会明白他们是指精神。这是他们的第一件诡计，如果被识破，他们便想出新的伎俩，假装无辜，反而指控我们心怀恶意。仿佛他们是坦率的信徒，他们说：\Quote{我们相信血肉的复活。}现在当他们这样说时，无知的群众认为应该满意了，特别是因为这些确切的话出现在信经中。如果你继续深究他们，场中就会响起不赞成的嗡嗡声，他们的支持者会喊道：\Quote{你已听到他们说他们相信血肉的复活；你还想要什么？}公众的偏爱就从我们这边转到他们那边，他们被称为诚实的人，而我们被看作诬告者。但如果你坚定脸色，紧紧抓住他们关于血肉的承认，进而追问他们是否主张那可见可触、能行走能说话的血肉的复活，他们先是笑，然后表示同意。当我们询问复活是否会重新展现头发、牙齿、胸脯、胃、手、脚以及身体的所有其他部分时，他们再也忍不住大笑起来，告诉我们那样我们就需要理发师、蛋糕、医生和鞋匠了。他们反过来问我们，难道我们相信复活后男人的脸颊仍然粗糙，女人的光滑，性别像现在一样区分他们的身体吗？如果我们承认这一点，他们立即从我们的承认中推导出最粗俗唯物主义的结论。因此，他们虽然主张整个身体的复活，却否认其各个部分的复活。\par

6. 现今不是以修辞言辞攻击邪僻道理的时候。无论是\Person{西塞罗}的丰富辞藻，还是\Person{德摩斯梯尼}的激昂雄辩，都不能充分表达我的情感之热烈，若试图揭露这些异端者所用的诡辩：他们口头上声称相信复活，心中却加以否认。因为他们的妇人抚弄自己的乳房，拍打胸脯，捏弄腿臂，说：\Quote{若这些脆弱的身体还要复活，复活于我们有何益处？不！我们若要像天使一样\BibleRef{玛 22:30}，我们就该有天使的身体。}这就是说，他们不屑于带着基督复活时所具有的肉和骨来复活\BibleRef{路 24:39}。现在暂且假定，我年轻时误入歧途，受教于外教哲学的学府，在信仰之初对基督宗教的教义一无所知，以为在\Person{毕达哥拉斯}、\Person{柏拉图}和\Person{恩培多克勒}那里所读到的，也包含在宗徒的著作中：我假定，我说，我相信这一切，那么你们为何还要追随一个在基督里尚是婴孩、尚是吃奶者的错误呢？你们为何从一位尚不认识宗教的人那里学习不虔敬呢？船难之后仍有一块木板可攀附；人亦可借着坦诚的告解来补赎罪过。你们曾在我误入歧途时追随我；现今我已回头，也请追随我。青年时我们曾迷途；既已年老，让我们改过迁善。让我们汇合眼泪与呻吟，一同哭泣，回归于创造我们的主。不要等待魔鬼的悔改；因为那是虚妄的期待，会将我们拖入地狱的深渊。生命必须在此处寻求或丧失。若我从未追随\Person{奥力振}，你们试图诋毁我就是徒劳的；若我曾是他的门徒，就请效法我的补赎。你们已相信我的告解，也要相信我的否认。\par

7. 但会有人说：\Quote{如果你知道这些事，为何还在你的著作中赞扬他？} 我今天仍会赞扬他，若非你和像你这样的人赞扬他的错谬。我仍会觉得他的才华有吸引力，若非有些人被他的不虔敬所吸引。宗徒说：\Quote{凡事要查验，好的要持守。}\BibleRef{得前 5:21}\Person{拉克坦提乌斯}在他的著作中，尤其是致\Person{德梅特里安}的书信中，完全否认圣神的实体存在，并随从犹太人的错误，说那些提及圣神的章节是指圣父或圣子，而且\Quote{圣神}一词只是证明天主性中这两位是圣的。但谁能禁止他阅读他的\WorkTitle{神圣教理}——他在其中以极大的才能攻击外邦人——仅仅因为他的这个观点应受憎恶？\Person{阿波利纳里斯}写了杰出的论文反对\Person{波菲利}，我称赞他的劳苦，尽管在许多点上我因他的愚蠢而鄙视他的学说。如果你们也愿意承认\Person{奥力振}在某些事上有错误，我就不会再说什么。承认他对圣子的看法有误，对圣神更是错谬，指出他在论及人的灵魂从天坠落时所犯的不虔敬，并证明他虽然口头上断言肉身的复活，却以其他说法破坏了这句话的力量。例如，在许多世代和一\Quote{万物的恢复}之后\BibleRef{宗 3:21}，\Person{加俾额尔}和魔鬼一样，\Person{保禄}和\Person{盖法}一样，童贞女和娼妓一样。一旦你们摒弃了这些错误说法，并用审查之杖将它们与教会的信仰分开，我就可以安全地阅读剩下的部分，先服了解药，就不再惧怕毒药了。例如，我说过这样的话：\Quote{奥力振在其他著作中超越了所有作家，而在注释\WorkTitle{雅歌}时，他超越了自己}，这对我无害；我也不怕面对我早年年轻时所说过的话，我曾称他为教会圣师。难道有人会假装说，我受特别请求翻译一个人的著作，就必须指责他吗？我必须在序言中说：\Quote{我翻译其著作的这位作者是异端者；读者，提防他，不要读，快逃离这条毒蛇；或者，如果你执意要读，要知道我所翻译的这些论文已被异端者和恶人篡改；但你无需害怕，因为我已纠正了他们败坏的所有地方}，换句话说，我应该说：\Quote{我翻译的作者是异端者，但我作为他的译者，是一位天主教徒。} 事实是，你和你们一派急于想表现得直率、真诚和诚实，却对修辞学的戒律和演讲术的技巧关注太少。因为你们承认他的\WorkTitle{论首要原理}是异端，并试图将责任推给别人，这给读者造成了困难；你们诱导他们审查作者的全部生平，并从其余著作中对问题作出判断。另一方面，我则足够明智地默默修正了我希望修正的地方：因此，通过无视罪行，我避免了偏见对待罪犯。医生告诉我们，严重的疾病不应施以治疗，而应顺其自然，否则治疗会加剧病情。奥力振在\Place{提洛}去世至今已将近一百五十年。然而，有哪个拉丁作家曾冒险翻译他的\WorkTitle{论复活}和\WorkTitle{论首要原理}、他的\WorkTitle{杂集}和\WorkTitle{注释集}或如他自己所称的\WorkTitle{卷册}？有谁曾因如此臭名昭著的著作而使自己蒙羞？我不比\Person{希拉利}更有口才，也不比\Person{维克多理}更忠于信仰，他俩都曾翻译他的\WorkTitle{讲道集}，但并非精确翻译，而是独立的意译。最近，\Person{安波罗削}也借用了他的\WorkTitle{六日创世}，但表达的是\Person{希玻里托}和\Person{巴西略}的观点，而非奥力振的。你们声称以我为榜样，对其他人如鼹鼠般盲目，却以瞪羚的眼睛审视我。好吧，如果我曾对奥力振心怀恶意，我本可以翻译这些书，让拉丁读者知道他的最坏作品；但我从未这样做，尽管多人要求，我总是拒绝。因为我一向不习惯对我欣赏其才华之人的错误幸灾乐祸。奥力振本人若仍活着，很快就会与你们这些自称保护他的人争执，并会像\Person{雅各伯}那样说：\BibleQuote{你搅乱了我，使我在这地的居民中成了可憎的。}\BibleRef{创 34:30}\par

8. 有人想要赞美\Person{奥力振}吗？让他像我一样赞美他。他自幼就伟大，确实是殉道者的儿子。在\Place{亚历山大里亚}，他主持了教会学校，继承了大有学问的司铎\Person{克莱孟}。他如此憎恶肉欲，以至于出于对天主的热情，但并非按着真知\BibleRef{罗 10:2}，他用刀自宫。他践踏贪婪。他熟记圣经，日夜努力阐释其含义。他在教会中发表了超过一千篇讲道，出版了无数他称为\OriginalTerm{卷册}{tomes}的注释。这些我现在略过，因为我的目的不是列出他的著作。我们中谁能读完他所写的一切？谁又能不钦佩他对圣经的热忱？如果有人以\Person{热诚者犹达斯}的精神向我提起他的错误，他将在\Person{贺拉斯}的诗句中得到回答：\par

\Quote{诚然，\Person{荷马}有时也会打盹，\newline{}但他并非没有理由：\newline{}这个过失是小罪，因为他的作品很长。}\par

我们不要效仿那人的过失，因为他的德行我们无法企及。其他人在信德上曾犯过错误，不论希腊人还是拉丁人，但我不可提他们的名字，以免人们认为我是在以他人的错误、而非奥力振自己的优点来为他辩护。\Quote{你会说，这是指责他们，而非为他开脱。}如果我说他没有犯错，或者我承认宗徒保禄或来自天上的天使，如\BibleRef{迦 1:8}所说，在败坏信德的事上也该听从，那你的说法就对了。但既然我坦率承认他错了，我就可以像阅读其他人一样阅读他，因为如果他错了，他们也同样错了。但你可能会说，如果错误是许多人的通病，为什么你只攻击他一人？我回答，因为只有他被你们称颂为宗徒。除去你们对他过分的爱，我也愿除去我极度的厌恶。你们收集别人书中的错误言论，只是为了捍卫奥力振的错误；你们却将他捧上天，根本不容许他犯过任何错误。不论你是谁，既然在宣讲新道理，我恳求你，请怜悯罗马人的耳朵，请怜悯一个曾受宗徒称赞的教会的信德。\BibleRef{罗 1:8}为何在四百年后，你们却试图教导我们罗马人这些至今我们一无所知的教义？为何你们公开宣扬伯多禄和保禄拒绝承认的观点？至今从未听说过这样的教导，然而世界已经成了基督徒的世界。至于我，我将在晚年坚持我少年时重生所接受的信仰。他们称我们为泥土城人，由尘土所造，粗鲁而属血肉，因为他们说，我们拒绝接受属神的事物；但他们自己当然是耶路撒冷的公民，他们的母亲在天上。\BibleRef{迦 4:26}我并不轻视基督在其内降生并复活的肉身，也不蔑视那经过烘烤成为洁净器皿、在天上掌权的泥土。然而我奇怪，为何那些贬低肉身的人却随从肉身生活，并珍惜且娇养那个是他们的仇敌的东西。也许他们确实想要应验圣经的话：\BibleQuote{爱你们的仇敌，祝福那迫害你们的。}我爱肉身，但只爱那贞洁的、童贞的、以守斋克制的肉身；我爱不是它的作为，而是它本身，那知道必要受审判、因而为基督殉道而死、被鞭打、被撕裂、被火烧的肉身。\par

他们争辩说某些异端者和心怀恶意之人篡改了奥力振的著作，这种愚昧也可如下证明。有谁能比奥力振的支持者欧瑟伯和狄迪莫斯更聪明、更有学问、更能言善辩呢？前者在其六卷本的\WorkTitle{护教书}中断言奥力振与他自己意见相同；后者虽试图为他的错误辩解，却承认他确实犯了错误。他无法否认他所看到的文字，便试图加以曲解。说异端者做了添加是一回事，但主张异端言论值得称赞是另一回事。如果奥力振的著作在全球都被篡改，并且在一天之内，如同米特里达梯的法令一样，从他的书中剪除了所有真理，那他的情况将是独特的。即使假设他的某一篇论文被篡改，难道他所有在不同时间地点发表的作品都可能被窜改吗？奥力振本人在写给罗马主教法比盎的信中，对自己发表了错误言论表示忏悔，并责备安博罗削过于仓促地将本意仅为私下流传的内容公之于众。然而直到今天，他的门徒仍在寻找托辞，以证明他著作中所有引起非议的内容并非出自他本人，而是出自他人。\par

10. 再者，当他们提到庞斐禄曾赞扬奥利振时，我个人非常感激他们，认为我配与那位殉道者一同被毁谤。因为，诸位先生，如果你们告诉我，奥利振的书被他的敌人篡改以败坏其声誉，那我为何不能反过来指出，他的朋友和追随者自己编撰了一卷书，假托庞斐禄之名，借一位殉道者的见证来为其师长辩护，以洗刷恶名？看哪，你们自己就在奥利振的著作中修改那些（照你们说）他从未写过的段落；然而，当有人说某人出版了一本他实际上并未出版的书时，你们却感到惊讶。但你们的说法很容易通过查阅奥利振已出版的著作来检验；而庞斐禄除此之外别无出版，因此毁谤更容易将一本书归在他名下。因为请给我看庞斐禄的其他任何作品，你们在任何地方都找不到，这本是他的唯一著作。那么我如何知道这书是庞斐禄所写？你们会说，文风和语调应当告知我。但我决不相信，一位如此博学的人会将他的才华初果用于辩护那些可疑且声名狼藉的立场。这篇论文冠以\Quote{护教}之名，本身就暗示了先前有指控；因为未经攻击则无需辩护。现在我仅提出一个论点，然而，这个论点的力量只有愚蠢和厚颜无耻才能否认。这部归于庞斐禄的论著，几乎包含了欧瑟伯第六卷为奥利振辩护的前一千行。然而，在该作品的其余部分，作者引用了段落试图证明奥利振是大公的。但欧瑟伯和庞斐禄彼此如此和谐，以至于他们似乎共享同一个灵魂，甚至其中一人更进而采用另一人的名字。那么他们怎么会在这一点上如此根本地分歧呢？欧瑟伯在他的所有著作中都证明奥利振是亚略派，而庞斐禄却将他描述为尚未举行的尼西亚大公会议的支持者。由此可见，这书不属于庞斐禄，而属于狄迪默或其他某人，他砍掉了欧瑟伯第六卷的开头部分，自己补上了其余部分。但我愿意慷慨，并承认这书是庞斐禄所写，只是庞斐禄尚未成为殉道者。因为他必定是在殉道之前写了这书。而你们会说，他为何配得殉道？当然是为了让他以殉道者的死来消除自己的错误，用他流出的血洗去他唯一的过失。普天之下有多少殉道者，他们在死前曾是罪的奴隶！难道我们要因那些犯罪者后来成了殉道者，就掩饰他们的罪吗？\par

11. 我最亲爱的弟兄们，我急迫地口述了这封回复你们信件的信；并且克服了我的顾虑，拿起笔来攻击那个我曾赞扬其才能的人。诚然，我宁愿冒声誉之险，也不愿冒信德之险。我的朋友们使我陷入如此尴尬的两难境地：我若沉默不言，就会被视为有罪；我若辩护，就会被视为仇敌。两者皆难，但我将选择较轻的一个。争吵可以和解，但亵渎得不到宽恕。我将判断留给你们，让你们发现我在翻译\WorkTitle{论首要原理}时耗费了多少辛劳；因为一方面，若从希腊文改动任何内容，作品便不再是翻译，而是曲解；另一方面，逐字忠实于原文，绝无助于保留其雄辩的魅力。\par

\SourceRecord{Translated by W.H. Fremantle, G. Lewis and W.G. Martley.；Nicene and Post-Nicene Fathers, Second Series；Vol. 6.；Edited by Philip Schaff and Henry Wace.；Buffalo, NY: Christian Literature Publishing Co.,；1893.；<http://www.newadvent.org/fathers/3001084.htm>.}

% structure: p0001 source=h1 target=section reason=page-title

\section{第八十五封书信}

致\Person{波肋诺}\par

\BibleRef{出 7:13} 如何与自由意志调和？以及（2）为何信徒的子女被称为圣的\BibleRef{格前 7:14}，而不依靠圣洗的恩宠？关于第一个问题，热罗尼莫请波肋诺参考他新译的奥力振的著作《\WorkTitle{论首要原理}》。关于第二个问题，他引用戴尔都良的解释。写于公元400年。\par

1. 你的话语敦促我写信给你，但你的口才又阻止我这样做。因为作为写信者，你几乎与图利不相上下。你抱怨我的信简短且不精致：这并非出于马虎，而是出于对你的敬畏，唯恐写得更长反而送你更多可挑剔的句子。此外，向你这样的好人坦白说吧：每当船只驶向西方时，同时有那么多信件要求我回复，若我回复所有写信人，我将无法完成我的任务。因此，我忽略了文辞的讲究，也不修改秘书的文稿，随口说出首先想到的话。这样，我写信给你时，是把你当作朋友，而不是批评家。\par

2. 你的信提出了两个问题：第一，为什么天主使法郎的心刚硬，以及为什么宗徒说：\Quote{这样看来，不在乎那愿意的，也不在乎那奔跑的，只在乎那显示仁慈的天主。}\BibleRef{罗 9:16}以及其他似乎否定自由意志的话；第二，如何解释那些由信徒——即由已领洗的父母——所生的子女，\BibleRef{格前 7:14} 若没有后来领受并持守的恩宠之赐予，他们便不能得救。\par

3. 你的第一个问题，奥力振在他的著作《\WorkTitle{论首要原理}》中回答得最出色，此书应我的朋友帕玛丘的请求，我最近翻译了。这项工作使我如此忙碌，以致无法履行对你的承诺，再次推迟寄送我的《\WorkTitle{达尼尔注释}》。的确，帕玛丘虽杰出而深爱我，若只有他一人请求，我也许会推迟到别的时间；但事实上，罗马几乎所有我们的弟兄都同样要求，声称许多人处于危险之中，有些人甚至接受了奥力振的异端教导。因此我被迫翻译了一本书，其中坏处多于好处，并恪守这一原则：既不增添也不删减，而保持希腊文真实含义的拉丁文全貌。你可以从上述弟兄那里借阅我的译本；但对你而言，希腊文同样适用；你既然能饮自源头，就不必转向我贫乏心智提供的浑浊溪流。\par

4. 此外，由于我是对一位有学识、精通圣经与世俗文学的人说话，我愿向阁下提醒这一点：不要以为我是个粗鲁的小丑，谴责奥力振所写的一切——如同他轻率的朋友们错误断言的那样——也不要以为我像哲学家狄约尼削那样突然改变立场。事实上，我只是拒绝他那令人反感的教义。因为我知道，那称恶为善、称善为恶，以苦为甜、以甜为苦的人，\BibleRef{依 5:20} 同受诅咒。谁会过分赞扬别人的教导以至默认亵渎呢？\par

5. 你的第二个问题，戴尔都良在他的著作《\WorkTitle{论独婚}》中讨论过，他宣称信徒的子女被称为圣洁，是因为他们犹如信仰的候选人，且未受偶像崇拜的玷污。也请思考：我们在会幕中读到的器皿以及礼仪崇拜所需的一切，都被称为圣洁；尽管严格来说，除了认识并敬拜天主的受造物外，没有什么是圣洁的。但圣经的用法有时也称那些洁净的、被净化的、或做过补赎的人为圣洁。例如，关于巴特舍巴，经上记载她由不洁中成为圣洁\BibleRef{撒下 11:4}，圣殿本身也被称为圣地。\par

6. 我恳求你内心不要默默指责我虚荣或虚伪。天主在我的良心中为我作证，上述无法避免的情形使我正在着手撰写注释时退却了；而且你知道，心不在焉时做的事绝不会做得好。我欣然接受你送来的帽子，虽是小礼物，却情意非小，正适合温暖老人的头部。礼物与馈赠者同样令我欢喜。\par

\SourceRecord{Translated by W.H. Fremantle, G. Lewis and W.G. Martley.；Nicene and Post-Nicene Fathers, Second Series；Vol. 6.；Edited by Philip Schaff and Henry Wace.；Buffalo, NY: Christian Literature Publishing Co.,；1893.；<http://www.newadvent.org/fathers/3001085.htm>.}

% structure: p0001 source=h1 target=section reason=page-title

\section{第八十六书}

\Emph{致泰奥菲罗}\par

\EditorialNote{热罗尼莫祝贺泰奥菲罗在反对奥利振主义的运动中取得胜利，并谈及他的使者普里斯库和欧布鲁斯在巴勒斯坦所做的善工。然后（以其态度上的奇特转变）请求泰奥菲罗宽恕耶路撒冷的若望，因为他无意中收留了一位被开除教籍的埃及人。此信写于公元400年。}\par

热罗尼莫致最蒙福的教宗泰奥菲罗：我最近收到您圣座的公文，结束了您长期的沉默，并召我返回职守。因此，尽管可敬的弟兄普里斯库和欧布鲁斯迟迟才将您的信件带给我，但他们如今正以信德的热忱从巴勒斯坦的一端奔走到另一端，驱赶并逼入洞穴那些异端的毒蛇，我特写数语，祝贺您的成功。普世都因您的胜利而欢腾。万民喜乐的群众凝望着您在亚历山大竖起的十字架旗帜，以及那标志着您战胜异端的闪耀战利品。愿您的勇气蒙福！愿您的热忱蒙福！您已表明，您长期的沉默是出于策略而非本意。我向您的尊威非常坦诚地说：我曾因您过于宽容而忧伤，且对舵手所定的航向一无所知，我曾渴望那些堕落之人的毁灭。但如今我看出，您的手已举起，若您延迟打击，那只是为更重地打击。至于某人所得到的欢迎，您没有理由对本城的主教感到不悦；因为您在信中并未就此下指示，他若对他一无所知的案件做决定，实属鲁莽。不过我认为，他既不愿也不敢在任何事上令您烦恼。\par

\SourceRecord{Translated by W.H. Fremantle, G. Lewis and W.G. Martley.；Nicene and Post-Nicene Fathers, Second Series；Vol. 6.；Edited by Philip Schaff and Henry Wace.；Buffalo, NY: Christian Literature Publishing Co.,；1893.；<http://www.newadvent.org/fathers/3001086.htm>.}

% structure: p0001 source=h1 target=section reason=page-title

\section{\WorkTitle{第八十七书信}}

\Emph{从\Person{德奥斐洛}致\Person{热罗尼莫}}\par

\EditorialNote{德奥斐洛告知热罗尼莫，他已将奥力振主义者逐出尼特里亚的隐修院，并敦促他藉写作反对流行异端来显示其对信德的热情。信函日期为公元400年。}\par

德奥斐洛主教致最亲爱、最敬爱的兄弟、司铎热罗尼莫。可敬的阿加托主教与敬爱的亚大纳削执事被派往你处，传达有关教会的信息。你得知其内容后，我深信你将赞同我的决议，并为教会的胜利而欢欣。因为我们已用先知之镰\BibleRef{岳 3:13}砍倒了某些邪恶的狂热分子，他们急于在尼特里亚的隐修院里播种奥力振的异端。我们记起了宗徒的警告之语：\Quote{你当以全权告诫。}\BibleRef{铎 2:15}因此，你这一方，既然希望分享这份赏报，就当赶紧用圣经的言论使那些受骗的人回转。我们的愿望是，若可能，在我们有生之年，不仅护卫公教信德和教会规矩，也保护托付给我们的子民，并平息一切异端学说。\par

\SourceRecord{Translated by W.H. Fremantle, G. Lewis and W.G. Martley.；Nicene and Post-Nicene Fathers, Second Series；Vol. 6.；Edited by Philip Schaff and Henry Wace.；Buffalo, NY: Christian Literature Publishing Co.,；1893.；<http://www.newadvent.org/fathers/3001087.htm>.}

% structure: p0001 source=h1 target=section reason=page-title

\section{第八十八封书信}

\Emph{致特奥菲鲁斯}\par

\EditorialNote{回复前一封信，热罗尼莫再次祝贺特奥菲鲁斯成功镇压奥力振主义，并告知其努力已在远至意大利的西方结出果实。随后，他请求特奥菲鲁斯寄送其会议（最近在亚历山大里亚举行）的决议。此信写于公元400年。}\par

热罗尼莫致最可敬的教宗特奥菲鲁斯。您的圣座的信件给我带来了双重的喜乐，部分是因为它的携带者是那些可敬且值得尊敬的人，即阿加托主教和亚他那修执事，部分是因为它显示了您对信仰的热忱，以对抗最邪恶的异端。您的圣座的声音响彻世界，令所有基督的教会喜乐，魔鬼的毒计被平息了。古蛇 \BibleRef{默 12:9} 不再嘶嘶作响，而是扭动、被开膛，隐藏在黑暗的洞穴中，无法忍受太阳的光辉。在您写信之前，我已经向西方发送了信件，向我同语言的人指出异端者所使用的一些诡辩。我认为这是天主特别眷顾的结果，您和我同时写信给教宗阿纳斯塔修斯，从而在不知情的情况下确认了我的见证。现在您直接敦促我这样做，我将比以前更热心地使远近的单纯灵魂从他们的错误中回头。如果需要，我也不会犹豫招致某些人的憎恨，因为我们应当中悦天主而不是人：\BibleRef{宗 5:29}。尽管事实上他们更急于辩护他们的异端，而我和其他人更少攻击它。同时，我恳求如果您有任何关于此事的会议决议，请惠寄给我，以便我靠着如此伟大主教之权威，更加自由和自信地为基督开口。味增爵司铎两天前从罗马到达，并谦卑地向您致敬。他反复告诉我，罗马和几乎整个意大利在基督之后得救归功于您的信件。因此，最可爱和最可敬的教宗，请勤勉行事，无论何时有机会，写信给西方的诸位主教，不要犹豫——用您自己的话说——\Quote{用锋利的镰刀割断邪恶的幼苗}。\par

\SourceRecord{Translated by W.H. Fremantle, G. Lewis and W.G. Martley.；Nicene and Post-Nicene Fathers, Second Series；Vol. 6.；Edited by Philip Schaff and Henry Wace.；Buffalo, NY: Christian Literature Publishing Co.,；1893.；<http://www.newadvent.org/fathers/3001088.htm>.}

% structure: p0001 source=h1 target=section reason=page-title

\section{第八十九封信}

\Emph{德奥斐洛致热罗尼莫}\par

\EditorialNote{这封信（可能日期早于前三封）将隐修士德奥达多推荐给热罗尼莫；德奥达多从罗马来，宣布罗马教会对奥利金主义的谴责，曾访问现已清除异端的尼特里亚隐修院，并希望在返回西方之前也朝拜圣地。信件的日期是公元400年。}\par

德奥斐洛主教，致最敬爱的主内弟兄热罗尼莫司铎。我得知了隐修士德奥达多的计划——此事也将为你的圣德所知——我予以赞同。他不得不离开我们，启航前往罗马；他若不能先拜访并如骨肉般拥抱你及与你同在隐修院的诸位可敬弟兄，便不愿动身。我相信，你必会因他带给你的消息而欢喜——这消息是：这里的教会已重获安宁。他已视察了尼特里亚的所有隐修院，并能向你讲述其中隐修士们的节制与温良；以及奥利金主义者如何被压制驱散，教会如何恢复和平，主的纪律如何得以维护。我多么希望，那些在你附近据说正暗中损害真理的人，也能摘下伪善的面具。我不得不如此写信，因为你邻近的弟兄们对他们有所误解。因此，你们要谨慎，避开这类人；正如经上所载：\Quote{若有人不将教会的信仰带给你们，不要向他问安。}当然，如此写信给你，或许实属多余——你本能使迷途者回改；但当那些关心信仰的人劝告智者与博学者时，也无妨。请以我的名义，问候与你同在的所有弟兄。\par

\SourceRecord{Translated by W.H. Fremantle, G. Lewis and W.G. Martley.；Nicene and Post-Nicene Fathers, Second Series；Vol. 6.；Edited by Philip Schaff and Henry Wace.；Buffalo, NY: Christian Literature Publishing Co.,；1893.；<http://www.newadvent.org/fathers/3001089.htm>.}

% structure: p0001 source=h1 target=section reason=page-title

\section{第九十封信}

\Emph{自德奥斐洛致厄丕法尼}\par

\EditorialNote{德奥斐洛写信给厄丕法尼，召集在塞浦路斯召开的会议以谴责奥力振主义，并要求他通过可信的信使将会议决议连同德奥斐洛自己的公函副本送至君士坦丁堡。他对这最后一点的焦虑源于消息称某些被绝罚的隐修士已启航前往君士坦丁堡，向主教若望·金口若望陈述他们的案件。这封信的日期是公元400年。}\par

\Person{德奥斐洛}致其挚爱的主内兄弟、同僚主教\Person{厄丕法尼}。\par

主曾对祂的先知说：\BibleQuote{看，我今天派你到列邦列国，为要拔出、拆毁、毁灭、……建造、栽植。}\BibleRef{耶 1:10}历代以来，祂将同一恩宠赐予祂的教会，使祂的身体\BibleRef{弗 1:23}得以完整保存，使异端思想的毒害无处得胜。如今我们也看见这话应验了。因为基督的教会\BibleQuote{没有玷污或皱纹，或任何这类事}\BibleRef{弗 5:27}已用福音的剑砍倒了从洞穴中爬出的奥力振派的毒蛇，并从其致命的传染中解救了\Place{尼特里亚}丰硕的隐修士团体。我已将我的行动简短的记述（限于时间所允许）写成一封通函，不加区别地寄给众人。由于阁下在我之前常参与此类争战，故现在当务之急是坚固那些在战场中的人，并为此召集全岛的主教。应寄一封公函给我、\Place{君士坦丁堡}的主教，以及你认为合适的人；以便普世一致明确谴责\Person{奥力振}本人以及他所创立的可耻异端。我获知若干诽谤真信德的人，名为\Person{亚摩尼}、\Person{欧瑟比}、\Person{欧提米}，因对异端的新热忱而乘船前往\Place{君士坦丁堡}，企图以他们的诡计诱骗尽可能多的新皈依者，并与他们不虔敬的老同伴重新商议。因此，请你留意将此事告知\Place{伊扫里亚}、\Place{旁非利亚}及邻近各省的所有主教；此外，若你认为合适，可附上我的信，使我们众人同一心灵，以主耶稣基督的能力，将这些交付撒殚，以摧毁他们所怀的不虔敬。\BibleRef{格前 5:4}为确保我的文书迅速抵达\Place{君士坦丁堡}，请派遣一位勤快的信使，一位圣职人员（如同我派遣尼特里亚隐修院的神长们以及其他隐修士，有学识、节制的），以便他们到达时能亲自报告所行的事。最重要的是，我恳请你向主献上热切的祈祷，使我们在这次争战中也能得胜；因为\Place{亚力山大里亚}和全\Place{埃及}的人民心中充满了极大的喜乐，因少数人被逐出教会，使教会的身体得以保持纯洁。请问候与你同在的弟兄们。我们这里的人在主内问候你。\par

\SourceRecord{Translated by W.H. Fremantle, G. Lewis and W.G. Martley.；Nicene and Post-Nicene Fathers, Second Series；Vol. 6.；Edited by Philip Schaff and Henry Wace.；Buffalo, NY: Christian Literature Publishing Co.,；1893.；<http://www.newadvent.org/fathers/3001090.htm>.}

% structure: p0001 source=h1 target=section reason=page-title

\section{第九十一封信}

\Emph{从\Person{厄皮法尼乌斯}致\Person{热罗尼莫}}\par

\EditorialNote{一封从厄皮法尼乌斯致热罗尼莫的兴高采烈的信，他在信中描述了他会议的成功（由德奥斐洛建议召开），寄给热罗尼莫一份教区会议书信副本，并敦促他继续翻译关于奥利振争议的拉丁文献。写于公元400年。}\par

致他最亲爱的主、子及兄弟，司铎\Person{热罗尼莫}，\Person{厄皮法尼乌斯}在主内问安。写给所有公教徒的公函特别属于您；因为您，对信德满怀热诚而反对一切异端，尤其反对\Person{奥利振}和\Person{阿波利纳里}的门徒；他们的毒根和深植的不虔敬，全能的天主已将其拖到我们中间，使他们在\Place{亚歷山大}被掘出后，在全世界枯萎。因为，我所爱的子，要知道\Person{阿玛肋克}已连根带枝被消灭，十字架的凯旋已立在\Place{勒非丁}的山上。\BibleRef{出 17:8}因为正如当\Person{梅瑟}的手高举时，以色列就获胜；同样，主坚固了祂的仆人\Person{德奥斐洛}，将祂的旗帜竖立在\Place{亚歷山大}教会的祭台上以反对\Person{奥利振}；使在他身上应验了这句话：\Quote{你要将这事记下作为纪念，因为我要将\Person{奥利振}的异端连同\Person{阿玛肋克}本人从天下全然涂抹。}为了不显得我重复同样的话而使信变得冗长，我寄给您实际寄给我的信函，好叫您知道\Person{德奥斐洛}对我说了什么，以及主在我晚年赐予了何等大的福分，以这位伟大主教之证词，批准了我一向所宣讲的原则。我想此刻您也已发表了某些作品，并且，如我在前一封信中就此主题向您建议的，您已为使用您自己语言的读者撰写了一篇论文。因为我听说，有些在信仰上遭遇船破的人\BibleRef{弟前 1:19}也来到了西方，且不满足于自己的毁灭，他们还想拉别人与他们一同丧亡；仿佛他们认为罪人的众多会减轻罪过的严重，而地狱的火焰并不因加添更多木柴而扩大。藉着您并与您一起，我们向与您同在隐修院中事奉天主的可敬弟兄们致以最亲切的问候。\par

\SourceRecord{Translated by W.H. Fremantle, G. Lewis and W.G. Martley.；Nicene and Post-Nicene Fathers, Second Series；Vol. 6.；Edited by Philip Schaff and Henry Wace.；Buffalo, NY: Christian Literature Publishing Co.,；1893.；<http://www.newadvent.org/fathers/3001091.htm>.}

% structure: p0001 source=h1 target=section reason=page-title

\section{第九十二封书信}

\Emph{德奥斐禄致巴勒斯坦与塞浦路斯主教会议书}\par

这份主教会议书信是公元400年在亚历山大举行的会议的决议，旨在谴责奥力振主义。原文为希腊文，由热罗尼莫译为拉丁文。\par

这封书信以相同的措辞发送给巴勒斯坦主教们和塞浦路斯主教们。我们复制了两份副本的标题。致巴勒斯坦主教的那份开始如下：致至爱的主、弟兄及同为主教的厄乌洛吉、若望、泽比亚诺、奥森提、狄约尼削、根纳迪、则诺、德奥多西、迪克特里、颇斐理、撒图尔尼诺、阿兰、保禄、阿摩尼、赫利亚诺、欧瑟比、另一位保禄，以及所有聚集在埃利德奉献庆典的天主教主教们，德奥斐禄在主内致候。\par

致塞浦路斯主教的那份开始如下：致至爱的主、弟兄及同为主教的厄丕法尼、玛尔西安、阿加彼图、波埃修斯、厄尔比迪、恩塔西、诺尔巴诺、玛塞多尼、亚里思托、则诺、阿西亚提库、赫拉克利德、另一位则诺、齐里亚古、阿弗洛狄托，德奥斐禄在主内致候。\par

书信要点如下。\par

我们亲自访问了尼特利亚的隐修院，发现奥力振异端在其中造成了巨大破坏。伴随着一种奇异的狂热：人们甚至自残或割舌以显示他们如何蔑视身体。我发现有些这样的人从埃及到了叙利亚和其他国家，在那里攻击我们和真理。\par

奥力振的著作已在主教会议中被宣读并一致谴责。以下是其主要错误，主要见于\OriginalTerm{论本原}{\GreekText{περὶ} \GreekText{᾿Αρχῶν}}。\par

1. 圣子与我们相比是真理，但与圣父相比则是虚妄。\par

2. 基督的国度终有一天会终结。\par

3. 我们应只向圣父祈祷，而非向圣子。\par

4. 我们的身体在复活后仍将腐朽可朽、终有一死。\par

5. 天堂中亦无完美之物；天使自身也有缺陷，其中一些天使以犹太人的祭献为食。\par

6. 星辰自知其运动，魔鬼藉其运行预知未来。\par

7. 巫术若真实，并非邪恶。\par

8. 基督曾为人类受苦一次；祂将再为魔鬼受苦。\par

奥力振派曾试图胁迫我；他们甚至煽动外教人，谴责摧毁塞拉比斯神庙；并企图使两名被控犯有重罪的人脱离教会的司法管辖权。其中一人是由依西多尔错误列入寡妇名单的妇人，另一人即是依西多尔本人。他是异端派别的旗手，他的财富为他们的暴力行动提供了无尽的资源。他们试图谋害我；他们夺取了尼特利亚的隐修院教堂，一度阻止主教们进入并阻止举行礼仪。如今，他们如同则步耳（贝耳则步）在地上来回行走。\par

我未曾伤害他们；我甚至保护过他们。但我不允许旧日友谊（与依西多尔）损害我们的信德与纪律。我恳请你们，无论他们到哪里，都要抵制他们，阻止他们扰乱托付给你们的弟兄。\par

\SourceRecord{Translated by W.H. Fremantle, G. Lewis and W.G. Martley.；Nicene and Post-Nicene Fathers, Second Series；Vol. 6.；Edited by Philip Schaff and Henry Wace.；Buffalo, NY: Christian Literature Publishing Co.,；1893.；<http://www.newadvent.org/fathers/3001092.htm>.}

% structure: p0001 source=h1 target=section reason=page-title

\section{\WorkTitle{第九十三封书}}

\Emph{来自\Place{巴勒斯坦}诸位主教致\Person{德奥斐洛}}\par

耶路撒冷会议致\Person{德奥斐洛}的主教会议信函，回复前述信件。译文仍由\Person{热罗尼莫}完成。\newline{}以下为概要：\par

我们已行你所愿的一切，巴勒斯坦几乎完全摆脱了异端的污染。我们但愿不仅奥力振派，连犹太人、撒玛黎雅人和外教人也能被铲除。奥力振主义在我们中间并不存在。你所描述的那些教义在此地从未听闻。我们弃绝那些持有此类教义者，也弃绝阿波利纳里的追随者，并且不会接纳任何被你绝罚的人。\par

\SourceRecord{Translated by W.H. Fremantle, G. Lewis and W.G. Martley.；Nicene and Post-Nicene Fathers, Second Series；Vol. 6.；Edited by Philip Schaff and Henry Wace.；Buffalo, NY: Christian Literature Publishing Co.,；1893.；<http://www.newadvent.org/fathers/3001093.htm>.}

% structure: p0001 source=h1 target=section reason=page-title

\section{\WorkTitle{第九十四封书信}}

\Emph{\Person{狄约尼削}致\Person{德奥斐洛}}\par

\EditorialNote{在这封信（由\Person{热罗尼莫}翻译成拉丁文）中，\Place{里达}的主教\Person{狄约尼削}赞扬\Person{德奥斐洛}在对抗\OriginalTerm{奥力振主义}{Origenism}中取得的重大胜利，并敦促他继续努力反对这一异端。写于公元400年。}\par

\SourceRecord{Translated by W.H. Fremantle, G. Lewis and W.G. Martley.；Nicene and Post-Nicene Fathers, Second Series；Vol. 6.；Edited by Philip Schaff and Henry Wace.；Buffalo, NY: Christian Literature Publishing Co.,；1893.；<http://www.newadvent.org/fathers/3001094.htm>.}

% structure: p0001 source=h1 target=section reason=page-title

\section{第九十五封书信}

\Emph{从教宗\Person{亚纳斯塔修}致\Person{辛普利西安}}\par

\EditorialNote{应\Person{德奥斐洛}的请求，罗马主教\Person{亚纳斯塔修}致信米兰主教\Person{辛普利西安}，告知他本人如同\Person{德奥斐洛}一样谴责了\Person{奥力振}，因为克雷莫纳的\Person{欧瑟伯}已将其亵渎言论呈报给他。后者向他展示了一本\Person{鲁菲努斯}翻译的\WorkTitle{论原则}的副本。此信写于公元400年。}\par

致我主及弟兄\Person{辛普利西安}，\Signature{亚纳斯塔修}\par

1. 牧者理应对羊群倍加看顾与警醒。同样，在瞭望塔上，谨慎的守望者也日夜为城池警戒。又如暴风雨来临、海上危险时，船长焦虑不安，唯恐狂风巨浪将船撞上礁石。怀着同样的心情，我们可敬的兄弟和同僚主教\Person{德奥斐洛}不懈地关注那些关乎救恩之事，以免天主子民在不同的教会因阅读\Person{奥力振}的作品而陷入可怕的亵渎之中。\par

2. 因此，我们根据上述主教的来信，告知你的圣德：我们同样——我们被安置在\Place{罗马}城中，宗徒之长、荣耀的\Person{伯多禄}首先在此建立了教会，又以其信德坚固了它——为了任何人不得违反诫命阅读我们提及的这些书籍，也谴责了这些书；并且我们以热切的祈祷敦促\Person{天主}和\Person{基督}所启示给圣史们教导的诫命得以严格遵守。我们敦促人们记住可敬的宗徒\Person{保禄}那充满预言与警告的话语：\BibleQuote{如果任何人给你们宣讲另一个福音，与我们给你们所宣讲的福音不同，当受诅咒。}\BibleRef{迦 1:8} 因此，坚守这诫命，我们声明：凡是过去\Person{奥力振}所写、与我们的信仰相悖的一切，我们同样予以拒绝和谴责。\par

3. 我藉司铎\Person{欧瑟伯}之手将本信寄递给你的圣德，他是一位充满炽热信德并对主怀着爱戴的人。他向我展示了一些亵渎的章节，我在审断它们时不禁战栗。如果\Person{奥力振}发表了任何其他著作，你要知道，这些著作及其作者都同样被我谴责。愿主保佑你，我理应尊敬的主内弟兄。\par

\SourceRecord{Translated by W.H. Fremantle, G. Lewis and W.G. Martley.；Nicene and Post-Nicene Fathers, Second Series；Vol. 6.；Edited by Philip Schaff and Henry Wace.；Buffalo, NY: Christian Literature Publishing Co.,；1893.；<http://www.newadvent.org/fathers/3001095.htm>.}

% structure: p0001 source=h1 target=section reason=page-title

\section{\WorkTitle{第九十六封信}}

\Emph{来自 \Person{德奥斐洛}}\par

\EditorialNote{热罗尼莫翻译的德奥斐洛401年逾越节信函。其中德奥斐洛详细驳斥了阿波利纳利和奥力振的异端。}\par

\SourceRecord{Translated by W.H. Fremantle, G. Lewis and W.G. Martley.；Nicene and Post-Nicene Fathers, Second Series；Vol. 6.；Edited by Philip Schaff and Henry Wace.；Buffalo, NY: Christian Literature Publishing Co.,；1893.；<http://www.newadvent.org/fathers/3001096.htm>.}

% structure: p0001 source=h1 target=section reason=page-title

\section{第九十七封信}

\Emph{致\Person{帕玛邱斯}与\Person{玛尔塞拉}}\par

\EditorialNote{此信是热罗尼莫致帕玛邱斯与玛尔塞拉，附有泰奥斐洛于402年发出的逾越节书信的译本及希腊原文。他之所以附上原文，是因前一年有人投诉他的翻译不准确。写于402年。}\par

1. 春天再次降临，我以东方的货物丰富你们，将亚历山大里亚的珍宝送往罗马：正如经上所载，\Quote{天主自南方而来，圣者由帕兰山而来，甚至一片浓荫。}（因此，在 \BibleBook{}{雅歌} 中，新娘欢快地呼喊：\BibleQuote{我满怀喜悦坐在他的荫下，他的果实甘甜可口。} \BibleRef{歌 2:3}）如今，依撒意亚的预言确实应验了：\BibleQuote{在那日，在埃及地必将有一座祭坛献给上主。} \BibleRef{依 19:19} \BibleQuote{罪恶在哪里增多，恩宠在那里也格外丰富。} \BibleRef{罗 5:20} 那些曾哺育婴孩基督的人，如今以炽热的信德捍卫祂的成年；那些曾从黑落德手中拯救祂的人，已准备好再次从这亵渎者和异端者手中拯救祂。德默特琉将奥力振逐出亚历山大城；但如今，因着泰奥斐洛，他被全世界放逐。正如路加将\Emph{\WorkTitle{\BibleBook{宗徒}{大事录}} \BibleRef{宗 1:1}}献给他的那位一样，这位主教因其对天主的爱而得名。那蠕动的蛇如今何在？那毒蛇处于何种困境？他的\par

人的面孔连接着狼的身躯。\par

那曾嘶嘶爬行于世、并夸口说泰奥斐洛主教和我都是其错误追随者的异端，如今何在？那些无耻犬类的狂吠如今何在？他们为赢得头脑简单的人，虚假地宣称我们支持他们的事业。被泰奥斐洛的权威和口才击溃后，他们如今如同鬼魔，只能从地下喃喃自语。\BibleRef{撒上 28:13} 因为他们不认识那由上而来\BibleRef{若 8:23}、只讲论天上之事的祂。\par

2. 愿这一毒蛇的种类\BibleRef{玛 3:7}要么诚心接受我们的道理，要么始终坚持己见；这样我们就能知道谁该尊敬、谁该回避。事实上，他们发明了一种新型的补赎，将我们视为敌人仇恨，尽管他们不敢否认我们的信仰。请问，他们这种既非时间也非理智所能治愈的懊恼究是何故？当刀剑在战场上闪烁，人倒地、血流成河时，敌对的双手往往在友好中紧握，战争的狂怒被出人意料的和平取代。唯有这异端的党羽无法与教会人士达成协议；因为他们从内心否认被勉强逼出的口头同意。当他们的公开亵渎大白于众，当他们感到听众对他们群起反对时，他们就装出一副天真的样子，宣称自己是第一次听到这些道理，且先前对他们的老师所讲授的毫不知情。当你手持他们的著作时，他们用口否认亲手所写的内容。诸位先生，你们为何要困扰\Place{普罗庞提斯海}，更换住地，流窜各国，并用满口泡沫撕咬一位杰出的基督主教及其追随者？如果你们的撤回是真诚的，你们应当以同样对信仰的热忱取代先前对错误的狂热。你们为何从各处拼凑这些诅咒的破布？又为何谩骂那些你们无法抗拒其信仰之人的生活？难道因为你们认为某些人视我们为罪人，你们就不再是异端者？难道因为你们能指出我们耳朵上的疤痕，不敬神就不再玷污你们的口唇？只要你们有豹的斑点、埃塞俄比亚人的皮肤，\BibleRef{耶 13:23}知道我也被痣斑标记，对你们的背信又有何助益？看，教宗泰奥斐洛被自由地允许证明奥力振是异端者；而门徒们并不为老师的话辩护。他们只是假装这些著作已被异端者篡改和窜改，如同许多其他作家的作品一样。因此，他们试图维持他的事业，不是靠自己的信仰，而是靠他人的错误。对于这些异端者，我仅言于此；他们在对我们不公的敌意狂怒中，泄露了内心的秘密情感，并证明了胸中溃烂创伤的不可治愈性。\par

3. 但你们是基督徒，是元老院的光明：因此请接受我附上的这封信。今年我将它同时以希腊文和拉丁文寄出，以免异端者再次谎称我对原文作了许多改动和增补。我必须承认，我辛勤努力，以译文的同样优雅来保存原文措辞的魅力：并且严格遵守固定界限，绝不偏离，我尽力保持作者流畅的口才，并按其语气传达他的言论。我是否在这两点上成功，必须留待你们的判断。至于这封信本身，你们应当知道它分为四部分。在第一部分，泰奥斐洛劝勉信友庆祝主的逾越节；在第二部分，他击杀了阿波林纳留；在第三部分，他摧毁了奥力振；而在第四即最后部分，他劝勉异端者做补赎。如果你们觉得反对奥力振的论战不够充分，应当记得去年信中已充分讨论了奥力振主义；而这篇我刚翻译的，旨在简洁，无需进一步详述该主题。此外，其针对阿波林纳留的简洁清晰的信仰宣认，并不缺乏辩证的精妙。泰奥斐洛先从对手手中夺过匕首，然后刺入他的心。\par

4. 因此，恳求主，使这篇赢得希腊人喜爱的作品，也不失拉丁人的喜爱，并使全东方所赞赏称颂的，罗马也欣然接纳于心。愿宗徒伯多禄的宝座以其讲道，坚固圣史马尔谷宝座的讲道。的确，民间传言说，真福教宗阿纳斯塔修斯与泰奥斐洛有相同的热忱和精神，他甚至将异端者追赶到他们藏匿的洞穴中。此外，他自己的信函告诉我们，他在西方谴责了东方已谴责的事。愿他长命百岁，以便异端复活的嫩芽能随着时间因他的努力而枯萎死亡。\par

\SourceRecord{Translated by W.H. Fremantle, G. Lewis and W.G. Martley.；Nicene and Post-Nicene Fathers, Second Series；Vol. 6.；Edited by Philip Schaff and Henry Wace.；Buffalo, NY: Christian Literature Publishing Co.,；1893.；<http://www.newadvent.org/fathers/3001097.htm>.}

% structure: p0001 source=h1 target=section reason=page-title

\section{第九十八函}

\Emph{来自\Person{德奥斐洛}}\par

\EditorialNote{这是热罗尼莫翻译的德奥斐洛402年逾越节信函。与前一年的信函（第九十六函）一样，主要涉及阿波里纳瑞斯和奥力振的异端。}\par

\SourceRecord{Translated by W.H. Fremantle, G. Lewis and W.G. Martley.；Nicene and Post-Nicene Fathers, Second Series；Vol. 6.；Edited by Philip Schaff and Henry Wace.；Buffalo, NY: Christian Literature Publishing Co.,；1893.；<http://www.newadvent.org/fathers/3001098.htm>.}

% structure: p0001 source=h1 target=section reason=page-title

\section{第九十九封书信}

致真福德斐琳\par

\EditorialNote{热罗尼莫将德斐琳所写主后404年逾越节书信译为拉丁文，并因健康和保拉去世所致忧伤延误寄出而致歉。此信写于主后404年。}\par

至真福教宗德斐琳，热罗尼莫敬上。\par

1. 自我收到尊座的书信及逾越节论著以来，直至今日，我因忧伤、哀恸、焦虑以及从各处传来的关于教会状况的不同消息而深受困扰，几乎无法将尊著译成拉丁文。您知道那句古话的真实性：悲伤使人语塞；而当身心的疾病叠加时，更是如此。我已卧病五日，高烧不退；因此只能以极大的匆忙口述这封信。但我愿以寥寥数语向尊座表明，在翻译尊著时，我何等努力地传达原文每句中特有的辞藻魅力，并使拉丁文的风格在某种程度上与希腊文相符。\par

2. 开篇您使用了哲学的语言；虽未指名道姓，却在教导众人之际一举击倒某人。在余下的部分——这是最难完成的任务——您将哲学与修辞学结合，为我们重现了德摩斯梯尼和柏拉图。您对纵情享乐进行了何等猛烈的抨击！您对节制的德行给予了何等崇高的赞扬！您以何等隐秘的智慧论述了昼夜交替、月亮的运行、太阳的法则、我们世界的本性；始终援引圣经的权威，以免在逾越节论著中似乎借用了世俗的来源！简言之，我不敢称颂这些，以免被指控为阿谀奉承。此书在哲学部分以及在不作人身攻击的情况下为自己所支持的事业辩护的部分，均极为出色。因此，我恳求您宽恕我的迟缓：我因圣善可敬的保拉的安眠而完全崩溃，以致除了翻译此书外，迄今未写任何与神圣主题有关的作品。诚如您所知，我突然失去了我所带领的安慰者——主为我作证——并非为了我的私需，而是为了圣人们的慰藉与舒畅，她曾尽心服侍他们。您圣善而可敬的女儿欧斯托基（她因失去母亲而不肯受安慰），以及与她一起的所有弟兄，都向您谦恭致意。请惠寄您所说近著之书，以便我翻译，或至少阅读。愿主内平安。\par

\SourceRecord{Translated by W.H. Fremantle, G. Lewis and W.G. Martley.；Nicene and Post-Nicene Fathers, Second Series；Vol. 6.；Edited by Philip Schaff and Henry Wace.；Buffalo, NY: Christian Literature Publishing Co.,；1893.；<http://www.newadvent.org/fathers/3001099.htm>.}

% structure: p0001 source=h1 target=section reason=page-title

\section{第一〇〇书}

第一〇〇书。来自\Person{德奥斐禄}。\par

\EditorialNote{热罗尼莫翻译的德奥斐禄的404年逾越节信函。在信中，德奥斐禄教导罪人补赎之心，推荐守斋的实践，并谴责奥利振的谬误。}\par

\SourceRecord{Translated by W.H. Fremantle, G. Lewis and W.G. Martley.；Nicene and Post-Nicene Fathers, Second Series；Vol. 6.；Edited by Philip Schaff and Henry Wace.；Buffalo, NY: Christian Literature Publishing Co.,；1893.；<http://www.newadvent.org/fathers/3001100.htm>.}

% structure: p0001 source=h1 target=section reason=page-title

\section{\WorkTitle{第一〇六封书信}}

\Emph{致\Person{Sunnias}与\Person{Fretela}}\par

\EditorialNote{一封长信，热罗尼莫在其中回答两位侨居格提卡的\Person{Sunnias}与\Person{Fretela}向他提出的若干问题。此二人勤研圣经，却对热罗尼莫公元三八三年所译拉丁圣咏集（即所谓罗马圣咏集）与《七十贤士译本》间的常见差异感到困惑，遂寄予他一份长长的问题清单，请求解释。热罗尼莫在回信中详尽处理了所有这些问题，并向他的通信者指出，他们被其所用的《七十贤士译本》（普通版本）误导了，该译本与奥力振在《六栏圣经》中所提供的校勘本（即他自己所用的版本）差异甚大。他还表达了喜悦，发现即使在格提卡人中，圣经如今也被殷勤研读。此信日期约为公元四〇三年。}\par

\SourceRecord{Translated by W.H. Fremantle, G. Lewis and W.G. Martley.；Nicene and Post-Nicene Fathers, Second Series；Vol. 6.；Edited by Philip Schaff and Henry Wace.；Buffalo, NY: Christian Literature Publishing Co.,；1893.；<http://www.newadvent.org/fathers/3001106.htm>.}

% structure: p0001 source=h1 target=section reason=page-title

\section{第一〇七号书信}

致莱塔\par

\EditorialNote{莱塔是保拉的儿媳，她从罗马写信询问耶柔米如何教养她幼小的女儿（也称为保拉）作为献身于基督的贞女。耶柔米现在详细指示她关于孩子的训练和教育。然而，他对所提方案在罗马是否可行有些疑问，因此建议莱塔如有困难，将保拉送到白冷，在那里她将由祖母和姨妈（年长的保拉和欧斯托基）照顾。莱塔后来接受了耶柔米的建议，将孩子送到白冷，她最终接替欧斯托基成为她祖母建立的修女院院长。这封信的日期是公元403年。}\par

1. 宗徒保禄写信给格林多人，教导一个尚在基督内未受训的教会神圣的纪律，在其它诫命中也规定了这一条：\BibleQuote{女人若有不信主的丈夫，而丈夫愿意与她同住，她就不应离开他；因为不信主的丈夫因信主的妻子而成圣，不信主的妻子因信主的丈夫而成圣；不然，你们的子女就是不洁的，但如今他们是圣的。}倘若有人至今以为纪律的约束过于松弛，教师过于宽容，就让他看看你父亲的家族吧——他是极有地位和学识的人，但仍行走在黑暗之中；他便会看到，因宗徒的劝告，苦根上结出了甘甜的果实，平凡的芦苇中散发出珍贵的香脂。你自己就是混合婚姻的后代；但保拉的父母——你和我的朋友托克修——都是基督徒。谁能相信，外教大司祭阿尔比努斯竟会应一位母亲的许愿，得到一个基督徒孙女？谁能相信，欢喜的祖父会从婴儿结结巴巴的嘴唇中听到基督的阿肋路亚，并在晚年将一位天主的贞女抱在怀中？我们的期望已完全实现。那一个不信者因他圣洁信主的家庭而成圣。因为，当一个人被一群信主的儿女和孙辈围绕时，他就如同信德的候选者。依我之见，倘若他在年轻时就有这么多基督信仰的亲属，他或许早已相信基督了。即使他唾弃我的信，嘲笑它，称我为愚人或疯子，他的女婿在相信之前也这样做过。基督徒不是天生的，而是造就的。尽管镀金装饰，卡皮托利山开始显得灰暗。罗马的每一座神庙都布满烟灰和蛛网。全城深深激动，民众涌过半毁的神庙，前往殉道者的坟墓。那未曾因信服而得到的信仰，或许会因极度的羞愧而被逼得信从。\par

2. 莱塔，我在基督内最虔诚的女儿，我对你这样说，是要教你不可绝望于你父亲的得救。我希望，那为你赢得女儿的信德也能赢得你的父亲，这样你便能因整个家庭所受的祝福而欢喜。你知道主的应许：\BibleQuote{在人所不能的，在天主却可能。}\BibleRef{路 18:27} 改过从不嫌晚。强盗甚至从十字架到了乐园。\BibleRef{路 23:42} 拿步高，巴比伦王，在他身体和心都变得像野兽、被迫与荒野中的畜生一起生活后，也恢复了理智。\BibleRef{达 4:33} 且不提这些古老的故事，它们在不信者看来似乎难以置信，难道你自己的亲戚格拉古——他的名字表明其贵族出身——在几年前担任城市长官时，没有推倒、打碎、震碎密特拉的神窟及其中的所有可怖雕像吗？我指的是那些用来使崇拜者获得入门等级的东西：\OriginalTerm{乌鸦}{Raven}、\OriginalTerm{新郎}{Bridegroom}、\OriginalTerm{士兵}{Soldier}、\OriginalTerm{狮子}{Lion}、\OriginalTerm{珀耳修斯}{Perseus}、\OriginalTerm{太阳}{Sun}、\OriginalTerm{螃蟹}{Crab}和\OriginalTerm{父亲}{Father}？我重复一遍，他岂不是摧毁了这些，然后将它们作为人质送在自己前面，为自己获得了基督徒的洗礼？\par

即使在罗马本身，异教也被遗弃在孤寂中。那些曾为万邦之神者，如今与角鸮和夜鸟一起留在孤寂的屋檐下。军队的旗帜上装饰着十字架的标记。皇帝的紫袍和闪烁宝石的王冠上，装饰着那耻辱而救人的十字架形象。埃及的\OriginalTerm{塞拉皮斯}{Serapis}已经成了基督徒；而在迦萨，\OriginalTerm{玛尔纳斯}{Marnas}被囚禁哀悼，时刻期待看见自己的神庙被推翻。从印度、波斯、埃塞俄比亚，我们每日成群地欢迎隐修士。亚美尼亚的弓箭手放下了箭囊，匈奴人学习圣咏集，寒冷的西徐亚人被信德的光辉温暖。格泰人，红肤黄发，带着帐篷教堂随军而行：也许他们在与我们的战斗中得胜是因为他们信仰相同的宗教。\par

3. 我几乎误入了一个新话题，当我保持轮子转动时，我的双手在塑造一个酒壶，而我的目标却是制作一个水罐。回应你和圣玛尔塞拉的祈祷，我愿以母亲的身份对你说话，教导你如何教养我们亲爱的保拉，她出生前就已奉献给基督，受孕前就已誓愿事奉祂。这样，在我们这个时代，我们看到了先知书中关于亚纳的故事重演：她起初不生育，后来却多产。你已将一种伴有忧愁的多产，换成了永不死亡的后裔。因为我确信，既然你将首生子献给了主，你将做众子的母亲。按法律，首生子应被奉献。\BibleRef{出 13:2} 撒慕尔和参孙都是例子，洗者若翰也是，当玛利亚进来时，他欢喜跳跃。因为他听见主借着童贞女的口说话，渴望从母腹中跳出迎接祂。既然保拉是因应许而生，她的父母应当给她与出生相称的训练。你知道，撒慕尔在圣殿中长大，若翰在旷野中受训。前者作为\OriginalTerm{纳齐尔人}{Nazarite}，留长发，不喝酒和烈酒，甚至在童年时就与天主交谈。后者避开城市，穿着皮腰带，以蝗虫和野蜜为食。\BibleRef{玛 3:4} 此外，为象征他将宣讲的补赎，他穿着骆驼毛的衣服。\par

4. 必须这样教育一个灵魂，它要成为天主的圣殿。它必须学会只听、只讲属于敬畏天主的事。它不应懂得不洁的话语，不应知道世俗的歌谣。它的舌头尚在柔嫩时，就当沉浸在圣咏的甘甜中。必须让保拉远离那些有轻佻思想的男孩；就连她的使女和女仆也要与世俗的同伴隔绝。因为倘若她们学到什么恶习，就可能教唆更多。为她准备一套用黄杨木或象牙制成的字母，每个字母都有其专名。让她拿这些字母玩耍，这样就连她的游戏也能教她些东西。不仅要让她掌握字母的正确顺序，还要让她把字母的名字编成顺口溜，而且不断打乱顺序，把最后的字母放到中间，中间的放到开头，好让她既通过视觉也通过听觉认识所有字母。此外，一旦她开始用尖笔在蜡板上写字，手还不够稳时，要么你把手放在她的手上引导她柔嫩的手指，要么在板上刻好简单的字样；这样她的努力被限定在这些界线内，就能沿着为她画出的线条走，不会偏离。拼写正确时给予奖励，用小孩子喜爱的小礼物吸引她前进。让她有学习的同伴，以激发她的竞争心，这样当她看到同伴受称赞时就会受到激励。她学得慢时，你不要斥责她，而要用表扬来激发她的心智，以便她因胜过别人而欢喜，因被超越而难过。最重要的是，你必须注意不要使功课令她生厌，免得童年时形成的厌恶延续到成年。她一点一点试着拼合发音的词不应是随意的，而应是专门为此选定并堆砌起来的名称，例如先知或宗徒的名字，或者\Person{玛窦}和\Person{路加}所记载的从\Person{亚当}开始的族长名单。这样，她的舌头得到良好训练，记忆力也同时得到发展。再者，你必须为她选一位年岁、生活和学问都经得起检验的老师。一位有学问的人，我想，不会羞于为一位女亲戚或高贵的贞女做\Person{亚里士多德}为\Person{菲利普}之子所做的事：他屈尊担任启蒙教师，教他识字。不可轻视那些看似微小的事，因为少了它们，伟大的成果就无法取得。知识的初步和开端，在受过教育的人和不学无术的人口里，听起来是不同的。因此，你务必注意，不要让孩子被女人们的愚蠢哄骗所引诱，养成把长词缩短的习惯，或者用金子和紫色装饰自己。这两种习惯，一种会败坏她的谈吐，另一种会败坏她的品格。因此，她不应该在童年时学那些日后要重新学的东西。据说\Person{格拉古兄弟}的雄辩很大程度上归功于他们母亲从小对他们说话的方式。\Person{霍尔滕修斯}还在父亲膝上时就已成了演说家。早期印象很难从心中根除。羊毛一旦染成紫色，谁能再恢复它的洁白？未用过的坛子长久保留第一次盛装之物的味道和气味。希腊史告诉我们，那位不可一世的\Person{亚历山大}，全世界的君主，却无法摆脱童年时从老师\Person{莱昂尼德斯}那里学来的举止和步态的坏习惯。我们总是乐意模仿邪恶；美德显得高不可攀时，缺点很快就被复制。保拉的乳母不可酗酒、轻浮或好说闲话。她的保姆要品行端正，养父要庄重。当她看到祖父时，要跳到他怀里，搂住他的脖子，不管他愿不愿意，都要对着他唱\Quote{阿肋路亚}。她可以被祖母宠爱，对父亲微笑以显示认识他，这样让每个人都喜欢她，使全家为拥有这样一朵玫瑰花蕾而欢欣。应当立刻告诉她，她还有另一位祖母和一位姑母；她也应当知道她以新兵身份登记在哪一支军队里，以及她应召服役的统帅是谁。让她渴望与那些不在身边的人在一起，并鼓励她做出戏谑的威胁，说要离开你去找他们。\par

5. 让她身上的服装和打扮提醒她，她是许配给谁的。不要穿她的耳朵，也不要用铅粉或胭脂涂抹那已奉献给基督的脸。不要在她脖子上挂金链或珍珠，不要用宝石压她的头，也不要把她的头发染红，免得它让人想起地狱的火焰。让她的珍珠是另一种，是那种以后她能卖掉，用来购买那颗\Quote{高价珍珠}的。\BibleRef{玛 13:46} 从前有一位贵妇，名叫\Person{普雷泰克斯塔塔}，她听从丈夫\Person{许梅提乌斯}（即\Person{欧斯托基乌姆}的叔父）的命令，改变了那位贞女的衣着和容貌，按照世俗的方式梳理了她凌乱的头发，想要挫败贞女本人和她母亲公开表达的意愿。但是，就在当天夜里，一位天使在梦中向她显现。天使面貌可怖，威胁要惩罚她，并开口说：\Quote{你竟敢把你丈夫的命令置于基督的命令之上？你竟敢用亵渎的手碰触天主的贞女的头？那些手将立即枯萎，好叫你因痛苦而知道你所行的，五个月后你将被打入地狱。此外，如果你仍坚持作恶，你将被剥夺丈夫和儿女。}这些事都在适当的时候应验了：那位不幸的女子悔改得太迟，迅速的死标明了她的补赎。基督就是这样严厉惩罚那些侵犯祂圣殿的人，\BibleRef{格前 3:17} 也这样嫉妒地护卫祂珍贵的珠宝。我在此讲述这个故事，并非出于对不幸者命运的幸灾乐祸，而是要警告你，你必须怀着极大的敬畏和谨慎，持守你对天主所发的誓愿。\par

6. 我们读到司铎厄里因儿子们的罪而使天主不悦；我们也被告知，若一个人的儿子们放荡不羁，此人不可被立为主教。\BibleRef{弟前 3:4} 另一方面，关于女人经上写道：\BibleQuote{她在生育上必获救恩，倘若她们持守信德、爱德、圣德和贞洁。} 既然如此，父母对已经成年独立的子女要负责，那么对于那些尚在哺乳、软弱无力，按主的话说，\BibleQuote{不能辨别右手和左手}\BibleRef{纳 4:11}——即尚不能分辨善恶的子女，父母岂不更应负责吗？如果你采取措施保护女儿免受毒蛇咬伤，为何不也同样小心地保护她免受\Quote{全地的锤子}？不让她喝巴比伦的金杯？不让她与狄纳一同出去看异地的女儿？\EditorialNote{创世纪第三十四章} 不让她免受轻浮的舞蹈和拖地的长袍？没有人不等杯口蘸了蜜就给人服药；同样，罪恶为了更好欺骗我们，便披上美德的外貌和伪装。那么，你也许会说，我们为什么读到：\BibleQuote{儿子不承担父亲的罪过，父亲也不承担儿子的罪过，} 而是\BibleQuote{犯罪的人，必定死亡}？\BibleRef{则 18:20} 我回答，这段话是指那些有辨别能力的人，正如福音中那人的父母所说：\BibleQuote{他已经成年……他会为自己说话。}\BibleRef{若 9:21} 当儿子还年幼，思想像孩子，直到他达到能辨别毕达哥拉斯字母所指的两条道路的年龄之前，他的父母要为他的行为负责，无论是善是恶。但也许你以为，如果子女不受洗，基督徒的子女要为自己的罪负责；而父母若不让那些因年幼而无法拒绝的子女受洗，则无过犯。事实上，正如圣洗确保孩子的救恩，这反过来也给父母带来益处。你是否奉献你的孩子，原在于你的选择，但既然你已奉献了她，你若疏忽了她，便是自担风险。我泛泛地说，因为在你身上你并没有选择权——你甚至在孩子受孕之前就奉献了她。谁奉献一只瘸腿、残废或有任何缺陷的牺牲，就被视为犯了渎圣之罪。\BibleRef{申 15:21} 那么，那预备了自己身体的一部分和无玷灵魂的纯洁来迎接君王拥抱，却疏忽了她的奉献的人，岂不更要受惩罚吗？\par

7. 当保拉稍长，像她的净配一样，在智慧和身量上，并天主和人的喜爱上增长时，\BibleRef{路 2:52} 让她与父母同去她真正父亲的圣殿，但不要与父母一同从圣殿出来。让他们在世界的路上，在人群中，在亲属的拥挤中寻找她，却只在圣经的殿堂里找到她，\BibleRef{路 2:43} 询问先知和宗徒关于她所誓许的那精神婚姻的意义。让她效法玛利亚的退隐，加俾额尔发现她独自在房间，她似乎因看见一个男人而惊慌。\BibleRef{路 1:29} 让孩子效法那一位，经上记载说：\BibleQuote{君王的女儿全然荣耀在内。} 被爱之箭射中，让她对她的良人说：\BibleQuote{君王带我进了他的内室。}\BibleRef{歌 1:4} 任何时候都不要让她外出，免得巡行城内的守卫遇见她，击打、伤害她，夺去她贞洁的帕子，\BibleRef{歌 5:7} 使她赤身流血。\BibleRef{则 16:1} 反而，当有人敲她的门时，\BibleRef{歌 5:2} 让她说：\BibleQuote{我是一道墙，我的双乳像塔楼。}\newline{}\BibleRef{歌 8:10} \BibleQuote{我洗了我的脚；怎能再弄脏呢？}\BibleRef{歌 5:3}\par

8. 不要让她与别人同席，即在父母的餐桌旁；免得她看见可能渴望的菜肴。我知道有些人认为，蔑视眼前实际的快乐是一种更大的德行，但我认为，避开那必然会吸引你的事物，是一种更安全的节制。我曾是学生时，在学校读到这句话：\Quote{谴责你所容许成为习惯的事，是不当的。} 让她甚至现在学习不喝酒，\BibleQuote{其中有过度的酒。}\BibleRef{弗 5:18} 但孩子在达到强壮年龄之前，禁食对他们的柔弱身体是危险且难以承受的，如果她需要，就让她沐浴，并让她因胃的缘故喝一点酒。\BibleRef{弟前 5:23} 也要让她以肉食为养料，免得她的脚在开始奔跑之前就衰弱。但我说这话是作为让步，而非命令；因为我怕她会虚弱，而不是我要教导她放纵。此外，为什么一位基督贞女不能完全行出他人只部分行的事呢？迷信的犹太人拒绝某些动物和产品作为食物，而在印度人中，婆罗门；在埃及人中，裸身智者，全以粥、米饭和苹果为生。如果仅仅是玻璃就能值得如此多的劳作，那么一颗珍珠岂不更值得更多的劳作吗？保拉是因誓愿而生的。让她的生活像那些在同样条件下出生的人一样。如果恩宠在这两种情况下都是相同的，那么所付出的辛劳也应当相同。让她对管风琴的声音充耳不闻，甚至不知道笛子、竖琴和西塔拉琴的用途。\par

9. 愿她每日的任务，就是将从那圣经中采集的花朵献给你。让她以希腊文熟记许多诗句，但也要以拉丁文接受教导。因为，若柔嫩的嘴唇从一开始便未受此塑造，舌头就会被异国口音败坏，本乡言语也会因外来因素而堕落。你本人必须是她的老师，是她可以塑造幼年品行的典范。绝不要让她在你或她父亲身上看到任何她无法不犯罪而效仿的事物。请你们两人都记住，你们是一位奉献贞女的父母，你们的榜样将比你们的教训更能教导她。花朵易谢，恶风很快便会使紫罗兰、百合与番红花枯萎。绝不要让她在无人陪伴下出现在公共场合。绝不要让她拜访教堂或殉道者圣所，除非与她的母亲同行。不要让任何青年男子对她微笑致意；不要让任何卷发纨绔子弟向她献殷勤。若我们的小贞女要去参加守夜礼或彻夜祈祷，绝不要让她离开母亲身边半步。她不可特别挑选一个侍女作为宠儿或心腹。她对一人所说的话，所有人都应知晓。让她选择的同伴，不是一位秀丽、衣着讲究、能用甜润嗓音唱歌的姑娘，而是一位面色苍白、严肃、衣着朴素、带有忧郁色调的女子。让她以一位具有可靠信德、品格与贞洁的年长贞女为榜样，这人善于用言语和榜样教导她。她应当在夜间起身诵念祈祷文与圣咏；在早晨咏唱圣诗；在第三时辰、第六时辰、第九时辰加入队列为基督作战；最后，点亮她的灯，献上晚祭。让她在这些操练中度过白天，当夜晚来临时，让她仍沉浸其中。让祈祷之后紧接着阅读，阅读之后又回归祈祷。在如此众多且多样的工作上，时间会显得短暂。\par

10. 让她也学习如何纺羊毛，执纺竿，将篮子放在膝上，转动纺轮，用拇指捻线。让她轻蔑地拒斥丝绸织物、中国绒布和金线锦缎：她为自己做的衣服应能御寒，而非暴露它声称要遮盖的身体。让她的食物是蔬菜和麦面饼，间或有一两条小鱼。为了不浪费时间详述节制食欲的规诫（我在别处已更详细地讨论过这一主题），让她的用餐总留下几分饥饿，并能立刻开始阅读或咏唱。我强烈不赞成——尤其对年幼者——漫长而过度的禁食，其中一周又一周地加增，甚至连油和苹果也被禁食。我从经验得知，在大路上劳作的驴子疲倦时会走向客栈。我们的克己可能变成暴食，如同依西斯与库柏勒的崇拜者，他们热吞山鸡与鹧鸪，以免牙齿亵渎刻瑞斯的礼物。如果允许持续禁食，就必须加以节制，使那些长途跋涉者能够坚持到底；我们必须留心，不要刚开始良好，却中途跌倒。然而在四旬期，正如我以前所写，实行克己的人应张开每一片帆，车夫应暂时放松缰绳，加快马匹的速度。然而，对生活在世俗中的人，和对贞女与隐修士，将有不同的规则。平信徒在四旬期消耗他胃中的内壁，如同蜗牛依靠自身的汁液生活，为日后丰盛的食物和宴饮准备肚子。但对贞女和隐修士而言，情况不同；因为当他们放马驰骋时，必须记住，他们的赛程并无间歇。只限于一段时期的努力，或许可以严厉，但没有时限的努力则必须更为适度。因为在第一种情况下，赛程结束后我们可以喘息；在最后一种情况下，我们必须持续不断、不停歇地前行。\par

11. 当你到乡间小住时，不要将你的女儿留在身后。不要给她任何独立于你生活的能力或机会，让她在独自一人时感到害怕。不要让她与世俗之人交谈，或与忽视誓愿的贞女为伍。不要让她出席你仆人的婚礼，也不要让她参与家中吵闹的游戏。至于沐浴，我知道有些人满足于说，一个基督徒贞女不应与阉人或已婚妇女同浴——前者因为他们在思想上仍然是男人，后者因为怀孕的妇女令人厌恶。我自己却完全反对成年贞女沐浴。这样的人应该因想到看见自己赤身而羞愧不已。通过守夜和禁食，她折磨自己的身体，使之服从。凭借寒冷的贞洁，她试图扑灭情欲的火焰，熄灭年轻的炽热欲望。并且故意不修边幅，她急忙破坏自己的天生美貌。那么，为什么她要通过沐浴来给沉睡的火添加燃料呢？\par

12. 愿她的宝藏不是丝绸或宝石，而是圣经的手抄本；对于这些手抄本，让她不要过分看重镀金、巴比伦羊皮纸和阿拉伯式图案，而要注重文字的正确和准确的标点。让她从学习\WorkTitle{圣咏集}开始，然后从\WorkTitle{箴言篇}中收集生活规则。从\WorkTitle{训道篇}中，让她养成轻视世界及其虚幻的习惯。让她效法\WorkTitle{约伯传}中德行与忍耐的榜样。然后让她转至\WorkTitle{福音}，一旦开始便不可放下。她也应热心学习\WorkTitle{宗徒大事录}和\WorkTitle{书信}。一旦她用这些珍宝充实了心灵的宝藏，让她牢记\WorkTitle{先知书}、\OriginalTerm{七卷书}{Heptateuch}、\WorkTitle{列王纪}与\WorkTitle{编年纪}，以及\WorkTitle{厄斯德拉传}与\WorkTitle{艾斯德尔传}的卷轴。完成这一切之后，她可以安全地阅读\WorkTitle{雅歌}，但不可在此之前：因为若在开始时就读它，她将无法看出，虽然它以肉身的语言写成，却是属灵婚姻的婚歌。若不明白这一点，她将因此受害。让她避开所有伪经作品，如果她不是因这些作品所包含道理的真理，而是出于对其中所载奇迹的尊敬而被引导去阅读，让她明白它们并非真正由所署名的人所写，许多错误元素已被引入其中，并且需要无限的辨别力才能在泥土中寻找金子。让\WorkTitle{西彼廉著作}常在她手中。\WorkTitle{阿塔纳修斯书信}和\WorkTitle{依拉利论文}，她可以阅读而无跌倒之虞。让她在所有不忽视信仰之正当关注的作品和才智中获得喜乐。但如果她阅读其他人的作品，宁可评判它们，而不是追随它们。\par

13. 你会回答：\Quote{我一个世俗女子，住在罗马，身处人群之中，如何能遵守所有这些训诫呢？}若是如此，就不要承担你无法胜任的重担。当你像断奶的依撒格那样断奶了\Person{保拉}，又像穿上了衣袍的撒慕尔那样给她穿上衣服，就把她送到她的祖母和姑母那里；放弃这颗最珍贵的宝石，让它在玛利亚的内室安放，在婴孩耶稣哭泣的摇篮中安息。让她在隐修院中长大，让她在童贞女团体中成为一员，让她学习避免发誓，让她视说谎为亵渎，让她不谙世事，让她度天使般的生活，在肉身中却超越肉身，让她以为所有人类都像她自己一样。且不说其他好处，这样做将使你免于照顾她的艰巨任务和监护的责任。宁可因她不在而遗憾，也比永远为她战栗要好。因为你不得不战栗，当你观察她说什么、对谁说、向谁鞠躬、最喜欢见谁。趁她还是个婴儿，就把她交给\Person{欧斯托基}，她的每一声啼哭都为你祈祷。这样她现在就成为她的神圣同伴，日后也成为她的继承者。让她注视并爱慕，让她\Quote{从最年幼时就钦佩}一位其言语、步态和衣着都是德行教育的人。\BibleRef{德 10:30}让她坐在祖母的膝上，让祖母向她孙女重复她曾教导自己孩子的课程。长期的经验已向保拉表明如何养育、保全和教导童贞女；在她冠冕中每日编织的是象征最高贞洁的神秘世纪。啊，幸福的童贞女！幸福的保拉，\Person{托克索提乌斯}之女，她因祖母和姑母的德行，在圣德上比血统更高贵！是的，\Person{莱塔}：如果你能亲眼看到你的婆婆和你的姐妹，并认识到那激励她们小身体的伟大灵魂；你天生对贞洁的渴望是如此强烈，我不怀疑你会在你的女儿之前到她们那里去，并会将自己从法律的首道命令中解放出来，置于祂在福音中的第二道安排之下。\BibleRef{创 1:28} \BibleRef{格前 7:1}你会视对其他后裔的渴望为无物，并将自己奉献给天主的服务。但因为\BibleQuote{有时拥抱，有时不拥抱}，\BibleRef{训 3:5}又因为\BibleQuote{妻子对自己身体没有权柄}，\BibleRef{格前 7:4}又因为宗徒说\BibleQuote{各人在什么身份上蒙召，就该安于这身份} \BibleRef{格前 7:20}在主内，并且因为那负轭的人应当如此奔跑，以免让他的同伴陷入泥潭，我劝你要在你的后代身上充分偿还你暂时在自己身上推迟付出的。当\Person{亚纳}曾在会幕中献上她曾向天主誓许的儿子，她再也没有接回他；因为她认为一个将成为先知的人，与一个仍想有更多孩子的她同住一屋是不合适的。因此，在她怀了他并生下他之后，她不敢独自来到圣殿，或空手出现在主面前，而是首先向祂付清了她所欠的；然后，当她献上那伟大的祭献后，她回了家，因为她为天主生了长子，她自己获得了五个孩子。\BibleRef{撒上 2:21}你惊讶于那圣妇的幸福吗？效法她的信德。此外，只要你愿意送来保拉，我承诺我将亲自作她的导师和养父。虽已年老，我会将她扛在肩上，训练她结结巴巴的嘴唇；我的职责将比世俗哲学家的担子伟大得多；因为后者只教导了一位马其顿王，他终有一天会死于巴比伦的毒药，而我将教导基督的婢女和净配，她终有一天要在天上被献给她的主。\par

\SourceRecord{Translated by W.H. Fremantle, G. Lewis and W.G. Martley.；Nicene and Post-Nicene Fathers, Second Series；Vol. 6.；Edited by Philip Schaff and Henry Wace.；Buffalo, NY: Christian Literature Publishing Co.,；1893.；<http://www.newadvent.org/fathers/3001107.htm>.}

% structure: p0001 source=h1 target=section reason=page-title

\section{第一〇八封信}

\Emph{致欧斯托基姆}\par

\EditorialNote{这是热罗尼莫最长的一封信之一，旨在安慰欧斯托基姆因母亲最近去世而悲伤。热罗尼莫详细叙述了保拉的故事：首先谈到她高贵的出身、婚姻和在罗马的社会成功，然后叙述她的归化以及作为基督徒隐修者的一生。信中用了大量篇幅描述她东行的旅程，包括访问埃及和尼特里亚的隐修院，以及游览圣地最神圣的地点。信的其余部分描述了她在白冷的日常生活和学习，并列举了她出众的许多德行。最后以感人的笔触描述了她的去世和安葬，并给出了放置在她墓上的墓志铭。此信写于公元404年。}\par

1. 倘若我全身的肢体都化为舌头，每一条肢骸都赋有人声，我也无法完全颂扬圣洁可敬的保拉的德行。她出身高贵，但更崇高的是她的圣洁；她昔日富有今世的财物，如今却因着为基督而拥抱的贫穷而更为卓越。她是格拉古家族的后裔，西庇阿家族的苗裔，是那位她以其为名的保禄的继承人和代表，是那位曾任非洲统帅的母亲玛尔蒂亚·帕皮里亚的真确合法之女；然而，她宁选白冷而不选罗马，舍弃那金光闪闪的宫殿，住进一间泥屋。我们并非因失去了这位完美的妇人而哀伤，反而感谢天主曾拥有她，不，是仍然拥有她。因为\Quote{众人都是为天主而生活}\BibleRef{路 20:38}，那些归向主的人仍被算作祂家族的一员。诚然，我们失去了她，但天上的华厦却得到了她；因为当她在肉身内时，她远离了主\BibleRef{格后 5:6}，常常流泪抱怨说：\Quote{我寄居在默舍客，住在刻达尔的帐幕中，有祸了！我的灵魂长久做客居。} \InnerQuote{刻达尔}意为黑暗，难怪她呜咽道，连她自己也在黑暗中，因为按宗徒所说：\Quote{世界是控制于恶者。}\BibleRef{若一 5:19} 并且 \Quote{它的光明等于黑暗，} 以及 \Quote{光明在黑暗中照耀，黑暗却没有把握住它。}\BibleRef{若 1:5} 她常常呼喊：\Quote{我在你面前是客旅，是寄居的，像我所有的祖先一样。} 又说：我渴望 \Quote{离开肉身与基督同在。}\BibleRef{斐 1:23} 每当她因身体虚弱（由难以置信的节食和加倍禁食所致）而苦恼时，人们常听见她说：\Quote{我痛击并奴役我的身体，免得我给别人宣讲了，自己反而落选；}\BibleRef{格前 9:27} 以及 \Quote{最好不吃肉，不喝酒；}\BibleRef{罗 14:21} 还有 \Quote{我以禁食谦卑了我的灵魂；} 以及 \Quote{在我病苦时，你使我安卧在床；} 还有 \Quote{你的手沉重地压在我身上，我的精力在夏日的干涸中耗尽。} 当那她以奇妙的耐心忍受的痛苦刺透她时，仿佛她看见了天开了\BibleRef{宗 7:56}，她便说：\Quote{但愿我有翅膀像鸽子！我必飞去，得享安息。}\par

2. 我呼求耶稣和祂的圣人们，尤以这位可敬女性的保护者和伴侣的特别天使为证，这些绝不是谄媚和阿谀之词，而是对她品格的每一句庄严见证。的确，它们不足以颂扬一位为全世界所赞颂、为主教们所钦佩、为贞女团体所惋惜、为成群的隐修士和穷人所哭泣的德行之士。读者，你愿意一言以蔽之她的全部德行吗？她让那些依赖她的人贫穷，但不像她自己那样贫穷。她如此对待她的亲属和她小家庭中的男女（她是他们的兄弟姐妹，而非仆人），这一点并不奇怪；因为她留下了她的女儿欧斯托基姆——一位献身于基督的贞女，这篇素描正是为安慰她而作——远离她高贵的家族，只富有信德和恩宠。\par

3. 那么，让我开始叙述吧。其他人可以回溯到遥远的时光，甚至到保拉的摇篮，以及（如果我可以这样说）她的襁褓，并谈论她的母亲布莱西拉和父亲罗加图斯。前者是西庇阿和格拉古的后裔；后者出身于一个直到今天在希腊仍显赫的家族。据说他血管里流着阿伽门农的血，后者曾围攻特洛伊十年而将其摧毁。然而，我只赞美属于她自己、从她圣洁灵魂的清泉中涌出的东西。当福音中宗徒们问他们的主和救主，祂将赐给那些为祂抛弃一切的人什么时，祂告诉他们，他们要在今世得到百倍，并在来世得到永生。\BibleRef{谷 10:28} 由此我们看到，不是拥有财富值得赞美，而是为基督而舍弃财富值得赞美；我们不应以我们的特权为荣，而应视它们与天主的信实相比为渺小。诚然，救主如今已在今世对祂的仆人和使女兑现了祂的许诺。因为一个藐视单个城市荣耀的人，今日闻名全世界；一个在罗马生活时无人知晓的人，因隐藏在白冷而成为所有罗马和蛮族地区的赞叹。有哪一种族不派遣朝圣者到圣地呢？在那里谁能找到比保拉更伟大的奇观？正如在众多宝石中，最珍贵的发出最明亮的光芒，太阳以其光芒掩盖并熄灭星辰的微弱火光；同样，她以谦卑在美德和影响力上超越所有人，虽在众人中为最小，却大于一切。她越是贬抑自己，就越是基督高举她。她隐藏，却并未隐藏。她逃避荣耀，却赢得了荣耀；因为荣耀如影随形，追逐那些追求者的人，却寻找那些藐视它的人。但我不可忽略继续我的叙述，也不可因专注于一点而忘却写作的规则。\par

4. 由于拥有这样的出身，保拉嫁给了托克索提乌斯，后者血管中流淌着埃涅阿斯和尤利乌斯家族的尊贵血液。因此，她的女儿，基督的贞女欧斯托基姆，被称为尤利娅，正如他被称为尤利乌斯。\par

从伟大的尤路斯传承下来的名字。\par

我说这些事，并非因为拥有它们对那些人重要，而是因为鄙视它们的人值得称赞。世俗之人仰望在这样特权中富有的人。我们却赞美那些为救主而鄙视它们的人；且奇怪地贬低所有保留它们的人，赞扬那些不愿如此做的人。因此，出身高贵的\Person{保拉}，通过她的多产和贞洁，赢得了所有人的赞许，先是她的丈夫，然后是她的亲属，最后是整个城市。她生了五个孩子：\Person{布莱西拉}——我在罗马时曾就她的去世安慰过她；\Person{宝利纳}——她留下了可敬而令人钦佩的\Person{帕玛基乌斯}继承她的誓愿和财产，我也曾为她写有一本关于她去世的小书；\Person{欧斯托基乌姆}——她现在在圣地，是贞洁和教会的一颗宝贵明珠；\Person{鲁菲娜}——她的夭折战胜了慈母的心；以及\Person{托克索提乌斯}——之后她再无生育。由此可见，她并非想履行妻子的职责，而只是顺从她丈夫对拥有男性后代的渴望。\par

5. 当\Person{保拉}去世时，她的悲伤如此之大，以致她自己也几乎死去：然而，她随后如此完全地献身于事奉主，以至于可能让人觉得她曾渴望他的死亡。\par

我该用什么言辞来描述她那显赫、高贵、昔日富有的房屋？她将所有财富都花费在穷人身上。我怎能描述她对所有人表现出的巨大关怀，以及她甚至对素未谋面者所施的深远仁慈？哪个垂死的穷人没有裹上她赠送的毯子？哪个卧床不起的人没有得到她钱囊的支持？她在全城极其勤勉地寻找他们，并认为若有一个饥饿或患病的人由他人供养，那便是她的不幸。她的施舍如此慷慨，以至于她剥夺了自己的孩子；当亲属为此规劝她时，她宣称自己留给他们的是基督仁慈中更好的遗产。\par

6. 她也无法长久忍受她那世俗高位和显赫家族所带来的拜访和拥挤的接待。她悲伤地接受人们对她的敬意，并竭尽所能避开和逃脱那些想向她致意的人。碰巧那时，东西方的主教们受皇帝诏书召集到罗马，处理教会间的某些纷争；这样，她见到了两位最可敬的人物和基督教会的主教：安提约基雅的主教\Person{保利诺}和萨拉米斯（现称君士坦提亚）的主教\Person{埃皮法尼乌斯}。\Person{埃皮法尼乌斯}被她接待为宾客；而\Person{保利诺}虽住在别人家里，她以心中的热忱待他，仿佛他也寄宿在她家中。受他们美德的激励，她越来越时刻想着离开家园。她不顾房屋、子女、仆人、财产，总之，一切与世俗相关的事物，她渴望——独自一人，无人相伴（如果能说她是独自一人）——前往因\Person{保禄}和\Person{安当}而闻名的旷野。最后，当冬天过去，大海解冻，主教们返回各自教区时，她也以祈祷和渴望随他们航行。长话短说，她下到\Place{波尔图斯}，由她的兄弟、亲属，尤其是她自己的子女陪同——他们想用爱的表示来克服他们慈爱的母亲。终于，船帆扬起，船桨的划动将船带向深海。岸上，小\Person{托克索提乌斯}伸出双手哀求，而已经长大的\Person{鲁菲娜}默默啜泣，恳求母亲等到她出嫁。但\Person{保拉}目光望天，泪水干涸；她以对天主的爱战胜了对子女的爱。她不再将自己认作母亲，好证明自己是基督的婢女。然而她的心在撕裂，她与悲伤搏斗，仿佛被强行与自身的一部分分离。她所克服的情感之巨大，使她的胜利更令人赞叹。在战俘落入敌手的种种残酷磨难中，没有比父母与子女分离更严酷的了。虽然这违背自然法则，她以不减的信德忍受了此考验；不仅如此，她以喜乐的心寻求它：以对天主的更大爱战胜了对子女的爱，她静静地只专注于\Person{欧斯托基乌姆}——她誓愿与航行的同伴。与此同时，船继续前行，所有同船者都回望海岸。但她挪开目光，以免看见那些她无法不痛苦注视的景象。必须承认，没有母亲曾如此深爱她的孩子。出发前，她将自己所有都给了他们，在地上剥夺了自己，好能在天上找到产业。\par

7. 船在\Place{蓬提亚}岛靠岸，该岛因\Person{弗拉维亚·多米蒂拉}这位尊贵女士的流放而早已闻名；她在\Person{多米提安}皇帝时代因承认自己是基督徒而被放逐。\Person{保拉}看到那位女士度过漫长殉道岁月的囚室，便以信德的翅膀飞翔，更加渴望见到\Place{耶路撒冷}和圣地。最强风似乎微弱，最高速似乎缓慢。在穿越\Place{斯库拉}和\Place{卡律布狄斯}之间后，她进入\Place{亚得里亚海}，平安航行至\Place{墨托涅}。在此短暂停留休养疲惫之躯后，\par

她将滴着水的四肢伸在岸上：\newline{}然后驶过\Place{马勒亚}和\Place{基西拉}岛，\newline{}分散的\Place{基克拉迪群岛}，以及所有陆地，\newline{}从四面狭窄地围住海。\par

然后，她将\Place{罗得}和\Place{吕基亚}抛在身后，终于望见了\Place{塞浦路斯}。在那里，她俯伏在圣\Person{厄丕法尼}的脚下，被他挽留了十天；但这并非如他所想的那样是为了恢复体力，而是事实表明，是为了让她能从事天主的工程。因为她走遍了岛上的所有隐修院，并在她的财力范围内向那里的弟兄们提供了实质性的救济，这些弟兄是因圣人的感召从世界各地聚集到此的。接着她渡过狭海，在\Place{塞琉西亚}登陆，由此上到\Place{安提约基雅}，被可敬的宣信者\Person{保利诺}的盛情挽留了片刻。然后，凭着信德的炽热，她——一位此前一直由宦官抬着行走的贵妇——竟骑着一头驴子踏上了旅途，而且正值隆冬。\par

8. 我无需赘述她穿越\Place{柯勒叙利亚}和\Place{腓尼基}的旅程（因为我并非要给你一份她漫游的完整路线图）；我只列举那些在圣经中提及的地点。在离开罗马殖民地\Place{贝里图斯}和古城\Place{漆冬}之后，她进入了位于\Place{匝尔法特}海岸的\Place{厄里亚城}，并在那里朝拜了她的主和救主。接着经过\Place{提洛}的沙滩——\Person{保禄}曾在那里跪过\BibleRef{宗 21:5}——她来到了\Place{阿科}，即今日所谓的\Place{托勒迈依}，骑马经过\Place{默基多}平原——此地曾目睹\Person{约史雅}王之死\BibleRef{列下 23:29}——然后进入了\Place{培肋舍特}人的地境。在此，她不禁赞叹\Place{多尔}（一度最为强大的城池）的废墟；还有\Place{斯特拉顿塔}，虽然曾微不足道，但由犹太王\Person{黑落德}重建，并为尊崇\Person{奥古斯都凯撒}而命名为\Place{凯撒勒雅}。在此她见到了\Person{科尔乃略}的房子，如今已改为基督教堂；以及\Person{斐理伯}的简陋住所；和他四个女儿的寝室，这四位童女\BibleQuote{都说预言}\BibleRef{宗 21:8}。接着她到了\Place{安提帕特黎斯}，一个半毁的小镇，由\Person{黑落德}以其父\Person{安提帕特}命名；还有\Place{里达}，如今称为\Place{迪奥斯波里斯}，因复活\Person{多尔卡}\BibleRef{宗 9:36}和治愈\Person{厄乃阿}\BibleRef{宗 9:32}而闻名。不远处是\Place{阿黎玛特雅}，即埋葬主的\Person{若瑟}的村庄\BibleRef{若 19:38}，以及\Place{诺布}，曾是一座司祭城，如今却成了他们被杀尸体安息的坟墓。\Place{约培}也在附近，是\Person{约纳}逃往的港口\BibleRef{纳 1:3}；同样——若我可以引入一个诗意的传说——曾看见\Person{安德洛墨达}被绑在岩石上。再启程后，她来到了\Place{尼苛波里}，曾名\Place{厄玛乌}，主在那里于擘饼时被认出；这一行动将\Person{克罗帕}的家祝圣为教会。由此出发，她登上了上下\Place{贝特曷隆}，这是\Person{撒罗满}所建的两座城池\BibleRef{编下 8:5}，但后来在多次毁灭性战争中遭毁；在她的右边看到\Place{阿雅隆}和\Place{基贝红}，\Person{若苏厄}（\Person{农}的儿子）与五位君王作战时曾在此命令日月停止\BibleRef{苏 10:12}，并且在此他判决\Place{基贝红}人（他们凭诡计获得了盟约）为劈柴挑水的人。\EditorialNote{苏 9}在\Place{基贝亚}，如今已完全成为废墟，她停留片刻，纪念那里的罪恶，以及那个妾被肢解，并且尽管如此，\Place{本雅明}支派仍有三百人得以保存，以便日后\Person{保禄}可被称为本雅明人。\par

9. 长话短说，\Person{保拉}离开了\Place{阿狄亚贝纳}王后\Person{赫肋纳}的陵墓——她在饥荒时曾送粮食给犹太人民——向左，进入了\Place{耶路撒冷}，即\Place{耶步斯}或\Place{撒冷}，这座有三重名字的城池，在化为灰烬和荒芜之后，由\Person{艾略斯·阿德里安努斯}再次重建为\Place{艾利亚}。尽管巴勒斯坦总督是她家的亲密朋友，提前派出了差役，并下令将他的官邸供她使用，她却选择了一间简陋的小屋。不仅如此，在拜访圣地的过程中，她对每一处都表现出极大的热情和热忱，若非急切要参观其余地方，她绝不会离开一处。在十字架前，她俯伏在地朝拜，仿佛看见主悬挂其上；当她进入复活所在的坟墓时，她亲吻了天使从墓门口滚开的石头\BibleRef{玛 28:2}。她的信德如此炽热，甚至以口舔舐主的身体曾躺卧的痕迹，如同一个渴求所渴望河流的人。她在那里流了多少眼泪，发出多少叹息，倾诉多少悲伤，全耶路撒冷都知道；她祈祷的主也知道。由此出发，她登上了\Place{熙雍山}；这名字意为\Quote{堡垒}或\Quote{望楼}。这曾是大卫攻取和后来重建的城市\BibleRef{撒下 5:7}。关于它的攻取，经上写着：\BibleQuote{祸哉，阿黎耳，阿黎耳}——即天主的狮子，（确实，那时它极其坚固）——\BibleQuote{大卫攻取的城市：}关于它的重建，经上说：\BibleQuote{他的根基在圣山上；主爱熙雍的门胜过雅各伯所有的帐幕。}他不是指我们今天在尘土和灰烬中看到的门；他指的门是那些地狱不能战胜的门\BibleRef{玛 16:18}，以及众多信基督的人通过进入的门\BibleRef{默 22:14}。在那里，有人给她看了支撑教堂门廊的一根血染的柱子，据说我们的主受鞭打时曾被绑在上面。在那里，也给她看了圣神降临于一百二十位信徒灵魂的地方，从而应验了\Person{约珥}的预言\BibleRef{宗 2:16}。\par

10. 然后，在按她的财力向穷人和她的同伴们分发了钱之后，她出发前往白冷，只在路右边停留去探望辣黑耳的墓。（正是在这里，她生下了她的儿子，这儿子注定不是像他临终的母亲所称的那样是\Person{贝诺尼}，即\Quote{我疼痛之子}，而是像他属灵的父亲预言性地命名的那样是\Person{本雅明}，即\Quote{右手之子}。）\BibleRef{创 35:18} 之后，她来到白冷，进入了救主诞生的山洞。在这里，当她注视由童贞女祝圣的客栈，以及那牛认识它的主人、驴认识它主人的槽的厩棚时，\BibleRef{依 1:3} 并且那同一位先知的话语已应验之处：\BibleQuote{那在水边撒种，任凭牛和驴在脚下践踏种子的，是有福的：} 当她注视这些事，我可以说，她在我的耳旁宣告，她能用信德的眼睛看见婴孩主裹在襁褓中，在马槽里啼哭，贤士们朝拜祂，明星照耀在上，童贞女母亲，专注的养父，牧人们夜间来看\Quote{所成就的事}，因而甚至在那时就祝圣了福音作者若望的开篇之词：\BibleQuote{在起初已有圣言} 和 \BibleQuote{圣言成了血肉。} 她宣称她能看见被屠杀的无辜者，暴怒的黑落德，若瑟和玛利亚逃往埃及；她悲喜交加地喊道：\Quote{万福，白冷，食粮之家，那从天上降下来的食粮就诞生在你那里。\BibleRef{若 6:51} 万福，厄弗辣大，果实丰饶之地，你的果实就是主自己。论到你，米该亚早已预言：\InnerQuote{你犹大地的白冷厄弗辣大啊，你在犹大的千城中绝非最小，因为将由你出来一位统治以色列的统治者；祂的根源来自亘古，来自永远。因此，祂必将遗弃他们，直到那孕妇分娩的时候；那时，祂其余弟兄将归于以色列子民。} } 因为在其中诞生了那位在晓明之先受生的元首。祂由父的诞生是在万世之前；达味的族系摇篮在你那里延续，直到童贞女生下她的儿子，而相信基督的遗民归于以色列子民，并自由地向他们宣讲这样的话：\BibleQuote{天主的圣言本来应当先传给你们，但你们既然拒绝接受，并断定自己不堪当永生，看，我们就要转向外邦人。}\BibleRef{宗 13:46} 因为主曾说：\BibleQuote{我被派遣，无非是为了以色列家迷失的羊。}\BibleRef{玛 15:24} 那时，关于祂，雅各的话也应验了：\BibleQuote{权杖必不离开犹大，立法者必不离他两足之间，直到那应得者来到，祂将成为万民的期望。} 达味曾立誓，他曾许愿说：\BibleQuote{我决不进入我的帐幕，也决不上我的床铺；我不容我的眼睛睡眠，也不让我的眼睑打盹，直到我为上主寻得一个处所，为雅各的天主找到一个居所。} 他随即解释了他渴望的对象，用先知之眼看见那一位将到来，我们现在相信祂已经来了。\BibleQuote{看，我们在厄弗辣大听闻了祂，我们在田野的树林中找到了祂。} 希伯来词\OriginalTerm{祂}{Zo}，按你们从课程中学到的，并不是指\Emph{她}，即主的母亲玛利亚，而是指\Emph{祂}，即主自己。因此他大胆地说：\BibleQuote{我们要进入祂的帐幕，要在祂的脚凳前敬拜。} 至于我，虽然我是个可怜的罪人，却堪当亲吻婴孩主啼哭的马槽，并在那分娩的童贞女生下婴孩主的山洞里祈祷。\BibleQuote{这是我的安息之所，} 因为那是我主的故乡；\BibleQuote{我要住在这里，} 因为这是我的救主所选择的处所。\BibleQuote{我为我的基督预备了一盏灯。} \BibleQuote{我的灵魂要因祂而生活，我的后裔要事奉祂。}\par

之后，\Person{保拉}下山走了一小段路，来到\Place{厄达尔}塔楼，即\Quote{羊群塔}，雅各曾在此附近牧放他的羊群，而夜间守夜的牧人们得以听到这样的话：\BibleQuote{光荣归于至高天主，在地上平安，良善归于人。}\BibleRef{路 2:14} 当他们放羊时，他们找到了天主的羔羊；其洁白纯净的羊毛被天上的露水浸湿，而其他大地则干燥无水，祂的血洒在门框上时，驱走了埃及的毁灭者，\BibleRef{出 12:21} 并除去了世人的罪。\BibleRef{若 1:29}\par

11. 随即她加快脚步，沿着通往迦萨的旧路前行——迦萨即\Quote{天主的勇力}或\Quote{财富}——默默想着那外邦人的预像、厄提约丕雅太监，他虽如先知所说改变不了肤色，\BibleRef{耶 13:23}却在阅读旧约时找到了福音的源泉。\BibleRef{宗 8:27}接着向右转，她从贝特族尔行至厄市苛耳，厄市苛耳意为\Quote{一串葡萄}。探子们正是从那里带回了那串奇妙的葡萄，作为土地肥沃的凭证，\BibleRef{户 13:23}也是那自称为\Quote{我独自踹了酒醡，众民中无一人与我同在}者之预像。\BibleRef{依 63:3}不久后，她进入了撒辣的住所，目睹了依撒格的诞生地和亚巴郎橡树的遗迹——橡树下亚巴郎看见了基督的日子而欢欣。她从那里起身，上到赫贝龙，即克黎雅特阿尔巴，意为\Quote{四个人的城}。这四个人是亚巴郎、依撒格、雅各伯和伟大的亚当——按希伯来人的说法（据农的儿子若苏厄之书），他们均葬于此。但也有许多人认为加肋布是第四人，并指认一侧的纪念碑为其墓。看过这些地方后，她无心再去克黎雅特色费尔，即\Quote{书卷之村}；因为她轻视那叫人死的文字，已找到了叫人活的圣神。\BibleRef{格后 3:6}她更欣赏上泉和下泉，那是刻纳次之子耶孚乃之子敖特尼耳以南方旱地和无水产业所换得的，藉着引水灌溉了旧约的干田。原是如此，他预表了罪人在圣洗的水中为旧罪所获得的救赎。次日，日出后不久，她站在卡法巴尔加（即\Quote{降福之家}）的山脊上——亚巴郎曾为索多玛向主求情时所追到的地点。\BibleRef{创 18:23}她由此俯瞰广袤的荒野和曾属于索多玛、哈摩辣、阿德玛与责波殷的土地，看到了恩革狄的香树和左哈尔。左哈尔即\Quote{三岁的母牛}，\BibleRef{依 15:5}古时称为贝拉，\BibleRef{创 14:2}叙利亚文译作左哈尔，意为\Quote{小}。她想起罗特避难的洞穴，含泪告诫同行的童贞女们谨防\Quote{酗酒}，\BibleRef{弗 5:18}因为摩阿布人和阿孟人正是由此而来。\BibleRef{创 19:30}\par

12. 我在正午的阳光下久久徘徊，因为那时新娘找到了新郎安歇之处，\BibleRef{歌 1:7}若瑟也与兄弟们一同饮酒。\BibleRef{创 43:16}我要返回耶路撒冷，途经特科亚——亚毛斯的故乡，\BibleRef{亚 1:1}瞻仰橄榄山上闪闪发光的十字架，救主正是从此处升天到父那里。每年，一头红色母牛在此作为全燔祭献给主，其灰烬用于洁净以色列子民。\BibleRef{户 19:1}同样，按厄则克耳所言，革鲁宾离开圣殿后，在此建立了主的教会。\BibleRef{则 10:18}\par

此后，保辣参观了拉匝禄的坟墓，目睹了玛利亚和玛尔大热情款待的房舍，以及贝特法革——\Quote{司祭之颚的城镇}。在那里，一匹桀骜的驴驹——预表外邦人——接受了天主的缰绳，被宗徒们的衣服所覆盖，\BibleRef{玛 21:1}谦卑地献上自己的背供主骑乘。她从那里径直下山到耶里哥，默想福音中受伤的人、司祭和肋未人视而不见的冷酷，以及撒玛黎雅人——即护守者——的仁慈：他将半死的人扶上自己的牲口，带到教会的客栈。\BibleRef{路 10:30}她注意到那个称为阿多米宁的地方，即\Quote{血地}，因强盗频繁侵扰，那里流血甚多。她也看到了匝凯的野桑树，\BibleRef{路 19:4}这树象征悔改的善工，匝凯藉此践踏了过去流血和抢劫的罪过，犹如从美德的高峰上看见了至高者。人们还指给她看路旁盲人坐着的那个地方，他们从主那里获得视力，\BibleRef{玛 20:30}成为相信他的两个民族——犹太人与外邦人——的预像。进入耶里哥，她看到了希依耳所建的城——他献了长子阿彼兰，立门时献了幼子色古布。\BibleRef{列上 16:34}她望见了基耳加耳的营地和包皮山，\BibleRef{苏 5:3}这暗示第二次割损的奥秘；\BibleRef{罗 2:28}她凝视着从约但河床中取出的十二块石头，这些石头象征那十二根基，其上刻有十二位宗徒的名字。\BibleRef{默 21:14}她也看见了法律中最苦最荒的泉源，真正的厄里叟以其智慧使它变为甘甜丰沛的水泉。黑夜刚过，她满怀热忱赶往约但河，站在河岸上，在日出之际追忆义德的太阳升起；\BibleRef{拉 4:2}司祭们的脚如何稳立河床中央；\BibleRef{苏 3:17}后来厄里亚和厄里叟的命令如何使河水左右分开为他们让路；以及主如何以祂的洗礼洁净了曾被洪水污染、被人类毁灭玷污的水。\par

13. 若要我讲述阿苛尔山谷——即\Quote{麻烦与拥挤}——那里曾谴责偷窃与贪婪，实为冗长；\BibleRef{苏 7:24} 以及贝特耳，\Quote{天主的殿}，雅各伯贫穷一无可取，睡在光秃地面。正是在这里，他将一块在匝加利亚描述为有七眼的石头放在头下，\BibleRef{匝 3:9} 在依撒意亚被称作角石，\BibleRef{依 28:16} 他梦见一个梯子直通天上；是的，主高高在上，\BibleRef{创 28:12} 向上升的人伸手，将疏忽的人从高处抛下。此外，当她在厄弗辣因山时，她去朝圣了农的儿子若苏厄和司祭亚郎的儿子厄肋阿匝尔的坟墓，二者正好相对；若苏厄的坟墓建于提默纳色辣黑，\Quote{在加阿士山的北边}，\BibleRef{苏 24:30} 而厄肋阿匝尔的坟墓\Quote{在他儿子丕乃哈斯所有的山上}。\BibleRef{苏 24:33} 她有些惊奇，那位亲自分配土地的人竟为自己选了崎岖不平的地份。关于史罗我要说什么呢？那里至今仍可见到一座毁坏的祭台，\BibleRef{撒上 1:3} 本雅明支派在那里预先模仿了罗慕路斯抢夺萨宾妇女的事。经过舍根（不是许多人误读的息哈尔），或如今称为乃阿颇黎的地方，她进入了建在革黎斤山旁雅各伯井周围的教堂；那口井正是主又饥又渴时所坐之处，并由撒玛黎雅妇人的信德而得滋养。她放弃了五个丈夫——这五者指梅瑟五书，以及那她所称并非丈夫的第六个，即假教师多西特乌斯——找到了真正的默西亚和真正的救主。保拉由此转向，她看见了十二圣祖的坟墓，以及撒玛黎雅——黑落德为尊崇奥古斯都而改名为奥古斯都或希腊语塞巴斯特。那里安葬着先知厄里叟、敖巴狄雅和若翰洗者，凡妇人所生的，没有一个大过他的。\BibleRef{路 7:28} 在这里，她被所见奇事充满恐惧；因为她看见魔鬼在圣人们的坟墓前受各种酷刑而尖叫，人们像狼嚎叫，像狗吠叫，像狮子吼叫，像蛇嘶嘶作响，像公牛咆哮。他们扭转头并向后弯到触地；妇女也被倒悬而衣服不掉落。保拉怜悯他们全部，为他们流泪，祈求基督怜悯他们。虽然她身体虚弱，却徒步登山；因为在山上的两个洞穴中，敖巴狄雅在迫害饥荒时期曾用饼和水供养一百位先知。\BibleRef{列上 18:4} 然后她迅速经过纳匝肋——主的摇篮；加纳和葛法翁——熟悉祂所行的奇迹；提庇黎雅海——因祂在其上航行而成圣；旷野——无数外邦人用几个饼吃饱，而以色列十二支派的篮子装满了吃剩的碎块。\BibleRef{玛 14:13} 她登上了大博尔山，主在那里显圣容。在远处她看见了赫尔孟山脉；以及广阔伸展的加里肋亚平原，息色辣和他的全军曾被巴辣克击败；还有克雄河\BibleRef{民 5:21}将平地分为两部分。附近还有纳因城被指给她看，那里寡妇的儿子复活了。\BibleRef{路 7:11} 若要讲述可敬的保拉以她难以置信的信德所到过的所有地方，我的话语比时间更早耗尽。\par

14. 我现在将前往埃及，途中在索苛稍作停留，以及在三松的井——他在颚骨的洼处凿开的。在那里我将滋润干渴的嘴唇，稍事休息，然后造访摩勒舍特；昔日因米该亚先知的坟墓而闻名，如今则因它的教堂。然后绕过曷黎人和革特人的地方，玛勒沙、厄东和拉基士，穿越沙漠的荒凉旷野，旅行者的足迹消失在流动的沙中，我将来到埃及河，称为史曷尔，\BibleRef{耶 2:18} 即\Quote{泥泞的河}，并通过埃及说客纳罕话的五座城，\BibleRef{依 19:18} 以及哥笙地和左安平原，天主在那里行了奇事。我将访问诺城，即后来的亚历山大；以及尼特利亚，主的城镇，那里日复一日，众人的污秽被以纯净的德性之碱洗去。保拉一看见它，就有一位可敬可佩的主教、精修者依西多禄前来迎接，伴随无数的隐修士，其中许多是司铎或肋未品级。目睹这些，保拉欢喜地看见主在其中的荣耀彰显；但她表示自己不应受此尊荣。我何必提玛加略、阿尔色尼、色拉丕雍或其他基督的柱石呢？有哪间斗室她没有进入？有谁是她没有俯伏在他脚下的？在每一位圣人身上，她相信她看见了基督亲自显现；无论她赐予他们什么，她欣喜地感到是赐予了主。她的热忱令人惊叹，她的坚韧在一位妇女身上几乎难以置信。她忘了自己的性别和软弱，甚至渴望与随行的女孩们一同住在这成千上万的隐修士中。而且，既然他们都愿意欢迎她，她或许会请求并获准这样做；若不是她因对圣地更强烈的热爱而被吸引开。由于酷热，她从培路基雍乘海船到玛约玛，她返回如此迅速，你会以为她是一只鸟。不久之后，她决心永久居住在圣善的白冷，在一间简陋的客栈住了三年，直到她建成所需的小室和隐修院建筑，更不必说为过往旅客准备的客房，使他们在那里找到玛利亚和若瑟曾未得到的接待。至此，我结束她由欧斯托基雅和许多其他贞女陪同旅行的叙述。\par

15. 我现在可以自由地更详细地描述那构成她独特魅力的美德；在阐述这一点时，我呼求天主作证，我不是在奉承。我不增添任何东西。我不夸张任何东西。相反，我淡化了很多，以免显得在讲述难以置信的事情。我的吹毛求疵的批评者不要暗示我在凭空想象，或者像装扮\WorkTitle{伊索寓言}中的乌鸦一样，用其他鸟类的羽毛来装饰保拉。谦逊是基督徒恩宠中的首要者，而她的谦逊如此显著，以至于一个从未见过她、只因她的名声而渴望见她的人，会相信他所见的不是她，而是她最低微的使女。当她被一群贞女围绕时，她总是在衣着、言语、姿态和步态上最不起眼。从她丈夫去世直到她安息，她从未与任何男人同席吃饭，即使她知道那人居于主教职位的顶峰。她从不进浴池，除非病危。即使在最严重的高烧中，她也不睡在普通的床上，而是睡在硬地上，只铺一张山羊毛席；如果那可以称为休息的话，那使得白天和黑夜都成了几乎不间断祈祷的时间。她很好地实现了\WorkTitle{圣咏集}的话语：\BibleQuote{每夜我使我床铺游泳；我以眼泪浇灌我的卧榻}！她的眼泪如泉涌般流出，她为自己最小的过错哀恸，仿佛那是深重的罪恶。我经常警告她要爱护眼睛，留着阅读\WorkTitle{福音}，但她只说：\Quote{我必须毁损这张脸，这脸违反天主的诫命，我曾用胭脂、铅粉和锑描绘过它。我必须克苦这个曾沉溺于许多享乐的身体。我必须用不断的哭泣来弥补我长期的欢笑。我必须用粗糙的山羊毛代替我柔软的亚麻布和昂贵的丝绸。我过去曾取悦我的丈夫和世界，现在渴望取悦基督。}如果我在她伟大而显著的美德中挑选她的贞洁作为赞美的主题，我的话似乎是多余的；因为，即使她仍在世俗中时，她已为罗马所有的已婚妇女树立了榜样，她的举止如此出众，以至于最恶毒的诽谤者也不敢将丑闻与她的名字联系起来。没有比她更体贴的心，也没有比她更善待卑微者的心。她不奉承有权势的人；同时，如果骄傲和虚荣的人寻求她，她也不轻蔑地转身离开。如果她看到穷人，她就扶助他；如果她看到富人，她就督促他行善。唯独她的慷慨是无限的。事实上，她如此渴望不让任何一个需要帮助的人空手而归，以至于她借高利贷，并且经常举新债来偿还旧债。我承认我错了；但当我看到她如此挥霍地施与时，我责备了她，援引了宗徒的话：\BibleQuote{这不是要别人轻松，你们受累，而是出于均衡，即现在你们富余可以补足他们的缺乏，好使他们的富余也补足你们的缺乏。}\BibleRef{格后 8:13} 我从\WorkTitle{福音}中引用了救主的话：\BibleQuote{有两件衣服的，要分一件给那没有的}；并警告她，她可能不会总有办法随心所欲地行事。我还引用了其他同样目的的论证；但她以令人钦佩的谦逊和简洁驳斥了所有论证。\Quote{天主是我的见证，}她说，\Quote{我所做的都是为祂的缘故。我的祈祷是，我可能死时是个乞丐，不给我女儿留下一分钱，而欠陌生人我的裹尸布的钱。}然后她以这些话结束：\Quote{我，如果乞讨，会找到许多人施舍给我；但若这个乞丐不能从我这里借来给他帮助，他就会死；如果他死了，他的灵魂要向谁索取？}我希望她更谨慎地管理她的事务，但她以比我的更炽热的信德全心紧靠救主，神贫地跟随主在祂的贫穷中，将她所领受的交还给祂，并为祂的缘故成为贫穷。她终于如愿以偿，死去时留给她女儿一身的债务。这债务欧斯托基乌姆至今仍欠着，而且确实无法靠自己的努力偿还；只有基督的仁慈才能使她摆脱它。\par

16. 许多已婚妇女习惯于把礼物赠送给自己的吹鼓手，她们对少数人极其慷慨，却拒绝帮助多数人。保拉完全免于这一过失。她按照各人的需要施舍金钱，不是助长放纵，而是缓解匮乏。没有穷人空手从她面前离开。她能做到这一切，不是靠她财富的巨大，而是靠她对财富的谨慎管理。她时常将这样的话挂在嘴边：\BibleQuote{怜悯人的人是有福的，因为他们要受怜悯。}\BibleRef{玛 5:7} 以及\BibleQuote{水能扑灭烈火，施舍能补赎罪过。}\BibleRef{德 3:30} 以及\BibleQuote{要用不义的钱财结交朋友，好使他们……接你们到永久的帐幕里。}\BibleRef{路 16:9} 以及\BibleQuote{施舍……那么一切对你们就都洁净了。}\BibleRef{路 11:41} 还有达尼尔对拿步高王所说的话，在其中他劝诫王用施舍来赎罪。她希望把钱不是花在那些将要同地和世界一起消逝的石头上，而是花在那些在地上滚动的活石上；关于这些活石，\WorkTitle{若望默示录}中记载，大王之城被建造在其上；\BibleRef{默 21:14} \WorkTitle{圣经}还告诉我们，这些石头要变成蓝宝石、翡翠、碧玉和其他宝石。\BibleRef{默 21:19}\par

17. 但这些品质，她或许可以与少数人共享，而魔鬼知道，最高之德并不在于这些。因为，当\Person{约伯}失去了他的财产，他的房屋和儿女被毁灭时，\Person{撒旦}对\Person{主}说：\BibleQuote{以皮换皮，凡人所拥有的，他必愿舍去以保全性命。但如今请伸手触及他的骨与肉，他必当面咒骂你。} \BibleRef{约 2:4} 我们知道，许多人在施舍时，并未触动其身体的安逸；他们向困乏者伸出援手，自己却沉溺于感官之乐；他们粉饰外表，内里却\BibleQuote{充满死人的骨骸。} \BibleRef{玛 23:27} \Person{保辣}并不属于此类。她的自制如此之大，几乎近乎过度；她的斋戒与辛劳如此严厉，几乎削弱了她的体质。除庆节外，她几乎从不在食物中用油；由此可知她如何看待酒、酱汁、鱼、蜜、奶、蛋及其他悦口之物。有些人以为取用这些便是极端节俭；即使他们以之暴食，仍幻想自己的贞洁安全。\par

18. 嫉妒总是紧随德行的轨迹：正如贺拉斯所言，山峰永远是闪电击中的目标。我如此论及男女，并不奇怪，因为法利塞人的嫉妒竟将我们的主本人钉死在十字架上。所有圣人都曾遭人嫉恨，就连乐园也未能免于那条蛇的恶意，死亡正是因它的毒计而进入世界。\BibleRef{智 2:24} 因此，主兴起厄东人哈达德来攻击帕乌拉，为要折磨她，免得她高傲，并常用她肉体上的刺\BibleRef{格后 12:7} 警告她，不要因自己德行的伟大而自高，也不要自以为与其他女子相比已达到了完美的顶峰。就我而言，我常对她说，最好向怨恨让步，在激情面前退避。雅各就是这样对待他的兄弟厄撒乌的；达味也是这样面对撒乌耳的无情迫害的。我提醒她，前者逃往了美索不达米亚；后者则投靠了培肋舍特人，\BibleRef{撒上 21:10} 宁愿屈服于外敌，也不愿受制于家中的仇敌。然而她却这样回答：——\Quote{你的建议固然明智，但倘若魔鬼不处处与天主的仆人和婢女作战，倘若它不总是先于逃亡者到达他们选择的避难所，那倒可以听从。此外，我对圣地的热爱阻止我接受它；我在别处找不到另一个白冷。为什么我不能以忍耐战胜这恶意？为什么我不能以谦卑击破这骄傲？当有人打你的右颊，你把另一面也转给他。\BibleRef{玛 5:39} 宗徒保禄岂不是说：\InnerQuote{不可为恶所胜，反要以善胜恶。}\BibleRef{罗 12:21} 宗徒们因为主的名受辱而夸耀，难道不是吗？难道救主自己不是谦卑自下，取了奴仆的形体，服从圣父至死，且死在十字架上，\BibleRef{斐 2:7} 为以苦难拯救我们吗？若约伯没有战斗并取得胜利，他就不会获得义德的冠冕，也不会听到主说：\InnerQuote{你以为我对你说话是为了别的，无非是要你显为义吗？}\BibleRef{约 40:8} 福音中，只有那些为义而受迫害的人被称为有福的。\BibleRef{玛 5:10} 我的良心安稳，我知道自己受苦并非由于任何过错；况且，今世的苦难正是来世得赏报的缘由。}当敌人格外放肆，竟敢当面责备她时，她便吟诵圣咏中的话：\Quote{当恶人站在我面前时，我默然不语，甚至不说好话；} 又说：\Quote{我如同聋子，不听；又像哑巴，不开口；} 以及\Quote{我如同一个听不见的人，口中没有辩驳。} 当她感到受诱惑时，她默想申命记中的话：\BibleQuote{上主你们的天主试验你们，要知道你们是否全心全灵爱上主你们的天主。}\BibleRef{申 13:3} 在苦难与磨砺中，她转向依撒意亚的精彩话语：\BibleQuote{你们这些断奶、离怀的，要期待苦难上加苦难，希望上加希望；只因嘴唇的恶意和毒舌的缘故，这些事还要持续片刻。}\BibleRef{依 28:9} 她解释这段经文以自慰，意思是：断奶的，即那些已长大成人的人，必须忍受苦难上加苦难，才配得希望上加希望。她也想起宗徒的话：\BibleQuote{我们连在磨难中也欢跃，因为我们知道：磨难生忍耐，忍耐生老练，老练生希望，希望不使人蒙羞}\BibleRef{罗 5:3} 和\BibleQuote{我们外在的人虽在朽坏，内在的人却日日更新}\BibleRef{格后 4:16} 以及\BibleQuote{我们这暂时轻微的苦难，正为我们成就无比永远的荣耀；我们并不注目那看得见的，而是注目那看不见的；因为看得见的是暂时的，看不见的才是永远的。}\BibleRef{格后 4:17} 她常说，虽然对人类的急躁而言，时间似乎来得缓慢，但不会很久，帮助必从天主而来，祂说：\BibleQuote{在悦纳的时候，我俯允了你；在救恩的日子，我帮助了你。}\BibleRef{依 49:8} 她宣称，我们不应惧怕恶人的欺诈之唇和舌头，因为我们因主的帮助而喜乐，祂借先知警告我们：\BibleQuote{不要怕人的辱骂，也不要因他们的毁谤而惊惶；因为蛀虫将像咬衣服一样咬他们，像咬羊毛一样咬他们}\BibleRef{依 51:7} 她又引用主自己的话：\Quote{你们要凭着坚忍，保全你们的灵魂}：以及宗徒的话：\BibleQuote{现时的苦楚，与将来在我们身上要显示的光荣，是不能较量的}\BibleRef{罗 8:18} 在另一处：\Quote{我们应当受苦磨难}，好使我们在一切遭遇中忍耐，因为\BibleQuote{易怒的人，大显愚昧；宽缓的人，大有智慧。}\BibleRef{箴 14:29}\par

19. 她在频繁的病痛和软弱中常说：\BibleQuote{我软弱时，就是坚强}\BibleRef{格后 12:10} \BibleQuote{我们有这宝贝放在瓦器里}\BibleRef{格后 4:7} 直到\BibleQuote{这必朽坏的变成不朽坏的，这必死的变成不死的}\BibleRef{格前 15:54} 又：\BibleQuote{基督的苦难怎样在我们身上丰盛，藉着基督，我们的安慰也怎样丰盛}\BibleRef{格后 1:5} 然后：\BibleQuote{既然你们同受苦难，也必同受安慰}\BibleRef{格后 1:7} 忧伤时她常唱：\Quote{我的灵魂，你为何沮丧？为何在我内忧闷？应仰望天主，因为我仍要赞美他，他是我脸上的救援，是我的天主} 在危险时刻，她常说：\BibleQuote{若有人愿意跟随我，该否定自己，背起自己的十字架，跟随我}\BibleRef{路 9:23} 又：\BibleQuote{凡要救他生命的，必将丧失生命}和\BibleQuote{凡为我丧失生命的，必得救生命}\BibleRef{路 9:24} 当人通知她财产耗尽、家业毁坏时，她只说：\BibleQuote{人纵然赚得了全世界，却赔上了自己的灵魂，为他有什么益处？或者，人还能拿什么作为自己灵魂的代价？}\BibleRef{玛 16:26} 和\BibleQuote{我赤身出离母胎，也赤身归去。主赐予，主收回：愿主的名受赞美}\BibleRef{约 1:21} 以及圣若望的话：\BibleQuote{不要爱世界，也不要爱世界上的事。因为凡世界上的，肉身的贪欲、眼目的贪欲、生活的骄傲，都不是出于父，而是出于世界。世界和它的贪欲都要过去}\BibleRef{若一 2:15} 我知道，当人传话给她，说她的孩子们，特别是她深爱的\Person{Toxotius}病重时，她首先以自制力成就了这句话：\Quote{我忧闷，但我没有发言}；然后以圣经的话喊出：\BibleQuote{爱儿子或女儿超过我的，不配属于我}\BibleRef{玛 10:37} 她并向主祈祷说：主啊，\Quote{求你保护那些被定死者的儿女}，即那些为了你每天在身体上死去的人。我知道，有一个搬弄是非的人——这种人造成很大伤害——有一次出于好意告诉她，因她在德行上的极大热忱，有些人以为她疯了，并说应该治治她的头。她以宗徒的话回答说：\BibleQuote{我们成了一台戏，给世界、天使和世人看}\BibleRef{格前 4:9}，\BibleQuote{我们为了基督成了愚妄的人}\BibleRef{格前 4:10}，但\BibleQuote{天主的愚妄比人更智慧}\BibleRef{格前 1:25} 因此她说，连救主也对父说：\Quote{你知道我的愚妄}，又\Quote{我成了一惊奇于多人，但你是我的坚固避难所。}\Quote{我在你面前如同牲畜，但我仍常与你同在。} 在福音中我们读到，甚至连他的亲人都想绑住他，以为他疯了。\BibleRef{谷 3:21}他的对手也辱骂他说：\BibleQuote{你是撒玛黎雅人，又附了魔}\BibleRef{若 8:48} 另一次：\BibleQuote{他赖魔王贝耳则步驱魔}\BibleRef{路 11:15} 她又继续说：但让我们听从宗徒的劝勉：\BibleQuote{我们所夸耀的，就是我们良心的见证：我们以单纯和真诚……靠天主的恩宠，在世为人}\BibleRef{格后 1:12} 且让我们听主对他宗徒们说的话：\BibleQuote{你们若属于世界，世界必爱属于自己的人；但因你们不属于世界……因此世界就恨你们}\BibleRef{若 15:19} 然后她转向主自己说：\Quote{你知道人心的秘密}；\Quote{这一切都临到我们身上；我们却没有忘记你，也没有违背你的盟约；我们的心没有退缩。}\Quote{是的，我们终日为你被杀，被看作待宰的群羊。} 但\Quote{主在我身边，我无所畏惧，人能对我做什么。} 她读过撒罗满的话：\BibleQuote{我儿，要敬畏主，你必变得强壮；除主之外，别怕任何人}\BibleRef{箴 7:2} 这些段落以及其他类似的话，她用作天主的盔甲来抵御邪恶的攻击，特别是用来防御嫉妒的猛烈冲击；她以耐心承受不义，从而平复了最野蛮人心的暴怒。直到她去世的那一天，她一生中有两件事特别引人注目：一是她极大的忍耐，二是别人对她的嫉妒。嫉妒吞噬怀有它的人的心；它力图伤害对手之时，却以全部暴怒的威力冲向自身。\par

20. 如今我要描述她隐修院的秩序，以及她如何将圣洁灵魂的节制转化为自身益处的办法。她播下属肉的事物，为要收获属神的事物；\BibleRef{格前 9:11}她给出属地的事物，为要获得属天的事物；她放弃暂时的事物，为要换取永恒的事物。除了为男性建立一座隐修院（其管理交由男性负责），她还将从不同省份聚集的众多贞女分为三个团体和隐修院，其中一些出身高贵，另一些则属于中下层阶级。然而，虽然她们分开劳作和用餐，这三个团体却一起咏唱圣咏和祈祷。在咏唱了阿肋路亚——这是召集她们参加集合祈祷的信号——之后，不允许任何人逗留。但\Person{保拉}要么是第一个，要么是最先到达者之一，她常常等候其余的人到来，以自己谦逊的榜样而非恐惧的动机敦促她们勤勉。在黎明、第三时辰、第六时辰、第九时辰、傍晚和午夜，她们轮流诵念圣咏集。任何姐妹都不允许对圣咏一无所知，所有人每天都要学习一定数量的圣经段落。唯独在主日，她们前往所居住的教堂，每个团体跟随自己的院长。按同样顺序返回后，她们便投身于各自分配的任务，制作衣服，或为自己，或为他人。若有贞女出身高贵，她不得拥有自己家族的仆从，以免她的使女满心回忆往日的作为和童年的放纵，通过不断交谈揭开旧伤，重蹈覆辙。所有姐妹衣着相同。除了擦手之外，不使用亚麻布。\Person{保拉}如此严格地将她们与男性分开，甚至不允许阉人接近她们；以免给诽谤者（总是乐于挑剔修道者）提供借口，来为自己的恶行开脱。当一位姐妹在诵念圣咏时迟到或在工作中疏忽时，\Person{保拉}常以不同方式接近她。若她急躁，\Person{保拉}便劝慰她；若她迟钝，\Person{保拉}便责备她，效法宗徒的榜样，宗徒曾说：\BibleQuote{你们愿意怎样？我愿意带着棍棒到你们那里去呢，还是怀着慈爱和温柔的精神？}\BibleRef{格前 4:21}除了食物和衣物，她不允许任何人拥有任何可称为自己之物，因为\Person{保禄}曾说：\BibleQuote{只要有衣有食，就当知足。}\BibleRef{弟前 6:8}她害怕拥有更多的习惯会在她们心中滋生贪欲；这种欲望没有财富能够满足，因为拥有越多，要求越多，富裕与贫乏都无法削减它。当姐妹们彼此争吵时，她用温和的话语使她们和好。若较年轻的受肉体欲望困扰，她便加重斋戒来削弱其力量；因为她宁愿她的贞女身体不适，也不愿灵魂受苦。若她偶然注意到某位姐妹过于关注衣着，她便皱起眉头，以严厉的神态责备她的错误，说道：\Quote{干净的身体和干净的衣着意味着不洁的灵魂。贞女的嘴唇绝不应说出不当或不洁的话，因为这表明淫荡的心思，而外在的人显明了内在的过错。}当她看见一位姐妹多言多语、好辩，或喜欢争吵，并在屡次劝诫后发现犯错者毫无改进迹象时，她便把她放在姐妹中最末的位置，置于她们的团体之外，命令她在食堂门口而非与众姐妹一起祈祷，并命令她独自进食，希望哪里责备无效，羞耻或能带来悔改。她憎恶偷窃之罪如同亵渎；世人视为轻微或无所谓的，她宣称在隐修院中是极严重的罪行。我怎能描述她对病人的仁慈和关注，以及她照料他们时那奇妙的关心和奉献？然而，虽然别人生病时她慷慨地给予一切宽容，甚至允许他们吃肉；但当她本人患病时，她却对自己的软弱毫无让步，似乎不公平地将她常向他人显示的仁慈，在她自己身上变成了严苛。\par

21. 任何身体强健的年轻女子都难以坚持如此严格的斋戒，但年迈体衰的保辣却对自己施行了这样的纪律。我必须承认，她在这一点上过于固执，不肯宽待自己，也不听劝告。我将叙述一件我所知的事实：在七月酷暑中，她曾一度突发高烧，我们几乎不抱希望。然而，因天主的仁慈，她复苏了。医生们力劝她喝一点淡葡萄酒以加快康复，并说如果她继续只喝水，恐怕会患上水肿。我私下恳请可敬的教宗艾比法纽劝告她，甚至强迫她喝酒。但保辣以其惯有的敏锐和机智，立刻识破了这个计谋，微笑着让他明白，他给她的建议终究不是他自己的，而是我的。长话短说，这位可敬的主教在多次劝告后离开了她的房间。当我问他取得了什么成果时，他回答说：\Quote{只有这件事：我这般年迈，几乎被说服不再喝酒了。}我叙述此事，并非因为我赞同人轻率地将超出自己能力的重担加于自身（因为圣经岂不是说：\BibleQuote{不要担负超过你能力的重担}？\BibleRef{德 13:2} ），而是因为我希望从她这种坚持不懈的品质中，显示出她心灵的热忱和信德灵魂的渴望；这两者使她以达味的话歌唱：\BibleQuote{我的灵魂渴慕你，我的肉身切望你。}虽然避免极端总是困难的，但哲学家们认为美德是中庸、恶习是过度，这观点十分正确，或者如我们可用一句话概括：\Quote{万事不可过分。}保辣在轻看食物方面如此坚定，却容易因亲人的去世而悲伤，尤其是她的子女。当她的丈夫和女儿们相继安息时，每一次丧亲之痛都危及她的生命。尽管她以十字圣号划口和胸，试图缓解母亲的哀伤，但她的情感压倒了她，母性的本能战胜了她那信赖的心灵。因此，虽然她的理智保持掌控，但她却被身体的软弱所征服。有一次，疾病缠身，久久不愈，使我们焦虑，也使她危险。然而，即使在那时，她仍充满喜乐，时刻重复宗徒的话：\BibleQuote{我真不幸！谁能救我脱离这该死的肉身呢？}\BibleRef{罗 7:24}\par

细心的读者可能会说，我的话更像是抨击而非颂扬。我呼求她所侍奉、我所渴望侍奉的耶稣基督为我作证：我既没有过分赞扬她，也没有贬低她，而是作为一位基督徒真实地叙述另一位基督徒；我正在写一篇回忆录，而非一篇颂词；她的缺点，在那些不如她神圣的人身上，也许会成为德行。我如此谈论她的缺点，是为了满足我自己的情感，以及我们这些仍爱她、哀悼她的兄弟姐妹们的深切悲痛。\par

22. 然而，她已经跑完了赛程，守住了信德，如今获得了正义的冠冕。\BibleRef{弟后 4:7} 她跟随羔羊，无论祂往哪里去。\BibleRef{默 14:4} 如今她饱飨，因为她曾饥饿。\BibleRef{路 6:21} 她喜乐地歌唱：\BibleQuote{我们曾听说，如今又亲眼看见，在万军的上主的城里，在我们天主的城里。}哦，何等有福的改变！她曾哭泣，如今却永远欢笑。她曾轻视先知所说的破裂水池，\BibleRef{耶 2:13} 如今却在主内找到了活水的泉源。\BibleRef{若 4:14} 她曾穿麻衣，如今却身穿白衣，并能说：\BibleQuote{你脱去了我的苦衣，给我披上喜乐。}她曾以灰为食，以眼泪拌饮，说：\BibleQuote{我昼夜以眼泪为饮食}，但如今她永远享用天使的食粮，歌唱：\BibleQuote{请你们体验并观看，上主是何等美善}，以及\BibleQuote{我心中涌现优美的言辞，我要向君王述说我的作品。}如今她看见依撒意亚的话，或者更确切地说，上主通过依撒意亚说的话实现了：\BibleQuote{看，我的仆人将要吃喝，你们却要饥渴；看，我的仆人将要欢乐，你们却要蒙羞；看，我的仆人将因心中喜乐而欢呼，你们却因心中忧愁而哀哭，因精神沮丧而哀号。}\BibleRef{依 65:13} 我曾说，她总是避开破裂的水池；她这样做，是为了能在主内找到活水的泉源，并能喜乐地歌唱：\BibleQuote{天主，我的灵魂渴慕你，好像牝鹿渴望溪水。我何时才能来到天主面前呢？}\par

23. 我必须简要提及她如何避开异端者的污秽水池，她视他们与外邦人无异。有一个狡猾的恶徒，自视为博学而聪明，在我不知情的情况下开始向她提出诸如此类的问题：婴儿犯了什么罪，竟被魔鬼攫取？我们复活时将是年轻还是年老？如果我们年轻时死去并在年轻时复活，那么复活后我们需要奶娘。然而如果我们在年轻时死去却在年轻时复活，那么死者将根本不会复活：他们将转变为新的存在。在来世将有性别的区分吗？或者没有这种区分？如果区分继续存在，就会有婚姻、性的交合和生育子女。然而如果它不再继续，那么复活的身体将不再是相同的。因为，他辩称，\BibleQuote{这属土的帐幕，压住多虑的心智}，\BibleRef{智 9:15} 但是我们在天上将有的身体是微妙而属神的，按照宗徒的话：\BibleQuote{播下的是属生灵的身体，复活起来的是属神的身体。} \BibleRef{格前 15:44} 从所有这些思考中，他试图证明理性的受造物因其过错和从前的罪而屈服于身体的条件；并且根据他们过犯的本性和罪责，他们出生在这种或那种生活状态。有些，他说，拥有健康的身体和富足高贵的父母；另一些则分得病弱的躯体和贫困的家庭；并因囚于现世和身体中，为他们从前的罪偿付惩罚。\Person{保辣}听了，并把所听到的告诉了我，同时指出了那个人。因此，这任务落在我身上，要反对这最有害的毒蛇和致命的害虫。圣咏作者论到这样的人时写道：\BibleQuote{求你不要将你的斑鸠的灵魂交付给野兽}，以及\BibleQuote{斥责芦苇中的野兽}；这些家伙书写不义，说谎言反对主，高举口舌反对至高者。既然那家伙试图欺骗\Person{保辣}，我应她的请求去找他，通过问他几个问题使他陷入两难。我说：你相信死者复活，还是不相信？他答道：我相信。我接着说：复活的身体是相同的还是不同的？他说：相同的。然后我问：那么他们的性别呢？它是保持不变还是被改变？听到这个问题，他沉默了，像蛇一样左右摇晃头以躲避打击。因此我继续说：既然你无话可说，我就替你回答，并从你的前提中得出结论。如果女人不复活为女人，男人不复活为男人，就没有死者的复活。因为身体是由性别和肢体构成的。但如果没有性别和肢体，那么身体的复活又成了什么？身体没有性别和肢体就不能存在。如果没有身体的复活，就没有死者的复活。至于你从婚姻提出的反对意见——如果肢体保持不变，婚姻必然被允许——这被救主的话解决了：\BibleQuote{你们错了，既不明白圣经，也不晓得天主的能力；因为在复活的时候，他们也不娶也不嫁，而是像天使一样。} \BibleRef{玛 22:29} 当说不娶也不嫁时，性别的区分显示仍然存在。因为没有人对没有婚姻能力的东西，如一根棍子或一块石头，说他们不娶也不嫁；但对于那些虽然能结婚，却因自己的德行和基督的恩宠而戒绝结婚的人，这样说就很合适。但如果你对此吹毛求疵，说：那样我们怎能像没有男女性别的天使呢？请听我简短的答复如下。主所应许给我们的，不是天使的本性，而是他们的生活方式和福乐。因此，甚至在斩首之前，\Person{洗者若翰}就被称为天使，而所有天主的圣洁男女和贞女甚至在现世就显示出天使般的生活。当说\BibleQuote{你们将像天使}时，所应许的只是相似性，而不是本性的改变。\par

24. 现在，轮到你来回答我这些问题。你如何解释多默触摸了复活的主的手，并看见祂被长矛刺透的肋旁？\BibleRef{若 20:26} 以及伯多禄看见主站在岸上\BibleRef{若 21:4}，并吃一块烤鱼和一块蜜脾的事实？\BibleRef{路 24:42} 如果祂站着，祂必定有脚。如果祂指着被刺伤的肋旁，祂必定也有胸和腹，因为肋旁是附着于它们，没有它们便不能存在。如果祂说话，祂必定使用了舌头、上颚和牙齿。因为正如弓弦击弦，舌头接触牙齿才能发出声音。如果祂的手被触摸，那么祂也必定有手臂。因此，既然承认祂有构成身体的所有肢体，祂也必定有由它们构成的整个身体，而且那不是女人的身体，而是男人的；这就是说，祂以祂死时的性别复活了。如果你进一步吹毛求疵说：那么我猜想，我们复活后也要吃东西；或者，一个坚固而有形的身体怎能违反其本性穿过关闭的门？你将从我这里得到这样的答复：不要因食物的问题而责备对复活的信仰；因为我们的主在复活会堂长的女儿后，命令给她食物。\BibleRef{谷 5:43} 而死去了四天的拉匝禄被描述与祂同席，\BibleRef{若 12:2} 在这两个事例中，目的都是显示复活是真实的，而非仅仅表象。如果你从主穿过关闭的门\BibleRef{若 20:19} 试图证明祂的身体是属灵的和属空气的，那么祂在受苦之前必定已有这属灵的身体；因为——违反重体的本性——祂能在海面上行走。\BibleRef{玛 14:25} 宗徒伯多禄也必定被认为有属灵的身体，因为他也以轻盈的脚步在水面上行走。\BibleRef{玛 14:29} 真正的解释是，当某事违反自然时，这是天主大能与权能的彰显。为了清楚表明在这些伟大的记号中，我们被要求注意的不是本性的改变，而是天主的全能，那位因信在水上行走的人，因缺乏信德而开始下沉，若非主用责备的话扶起他，\Quote{小信德的人哪！你为什么怀疑？} \BibleRef{玛 14:31} 我惊讶你竟能如此厚颜，当主亲自说：\Quote{把你的指头伸到这里来，看看我的手；伸过你的手来，探入我的肋旁；不要作无信德的人，但要作有信德的人。} \BibleRef{若 20:27} 又在另一处说：\Quote{看我的手和脚，就知道是我自己；摸我看看，因为鬼神没有骨和肉，如同你们看我有的。说了这话，就把手和脚给他们看。} \BibleRef{路 24:39} 你听祂提到骨和肉、脚和手；而你却想用斯多葛派妄谈的气泡和虚空之物来搪塞我！\par

25. 再者，如果你问，一个从未犯罪的婴儿如何被魔鬼抓住，或者我们在不同年龄死去，将以什么年龄复活；我唯一的回答——我想是不受欢迎的——将是圣经的话：\Quote{天主的判断是多么高深！} 和 \Quote{啊，天主的财富、智慧和知识，是多么高深！他的判断，是多么不可测量；他的道路，是多么不可探察！有谁知道上主的心意？或者，有谁当过他的顾问？} \BibleRef{罗 11:33} 年龄的差异不能影响身体的真实性。虽然我们的身体不断变化，每日失去或获得，但这些变化并不使我们成为不同的个体。我不是十岁时一个人，三十岁另一个人，五十岁又另一个人；现在我满头白发时也不是另一个人。按照教会的传统和宗徒保禄的教导，答案必须是：我们将作为完美的人复活，达到基督圆满年龄的程度。\BibleRef{弗 4:13} 犹太人假定亚当是在这个年龄被造的，我们也读到主和救主在这个年龄复活。我还从新旧约中引用了许多其他论据，以压制这个异端者的叫嚣。\par

26. 从那天起，\Person{保拉} 开始如此深切地憎恶这个人——以及所有同意他学说的人——以至于她公开宣布他们是主的敌人。我叙述这件事，与其说是有意用几句话驳斥一个需要数卷才能驳倒的异端，不如说是为了表明这位圣善女性的伟大信德：她宁愿使自己承受人们永久的敌意，也不愿通过有害于自己的友谊来激怒或冒犯天主。\par

27. 现在，再回到我此前不久开始描述的她品格的篇章。没有比她的心灵更温良的了。她\BibleQuote{敏于听教，迟于发言}\BibleRef{雅 1:19}，牢记着那诫命：\Quote{以色列啊，你要静默聆听。}圣经她熟记于心，并称其中所载的历史是真理的基础；但即使她喜爱这些，她仍更倾向于寻求其内在的属灵意义，并将此作为在她灵魂内所建属灵殿宇的基石。她请求允许她和她的女儿在我的指导下通读新旧约。出于谦逊，我起初拒绝了，但她坚持要求，并屡次敦促我同意，我最终答应了，并教导她我所学到的——不是出于我自己，因为自信是最坏的老师——而是来自教会最著名的作家。每当我感到困难并诚实地承认自己有所不知时，她绝不满足，而是用新的问题逼我指出，在各种不同的解释中，哪一个在我看来最有可能。我还要在此提到另一件事实，对于那些嫉妒的人似乎难以置信。当我本人从年轻时开始，历经辛苦和努力，才部分地掌握了希伯来语，并不断学习，唯恐一旦放下它，它也会离我而去；而葆拉，一旦下定决心也要学习它，竟然学得如此成功，以致她能以希伯来语咏唱圣咏，而且说话时完全不带拉丁语特有的发音。同样的成就，今天在她的女儿欧斯托基的身上也能看到，她始终紧随母亲身旁，服从她的一切命令，从不与她分开而眠，从不离开她外出或用餐，没有一文钱是她自己的，她为母亲将她微薄的遗产施舍给穷人而欢喜，并深信在孝爱中她拥有最好的遗产和最真实的财富。我不能默然不提葆拉心中的喜悦，当她听到她的小孙女和同名人——莱塔和托克西乌斯的孩子，这孩子是在她父母奉献她守贞的誓愿后出生，我甚至可以说是在那誓愿中受孕的——当，我说，她听到小宝宝在摇篮里唱\Quote{阿肋路亚}，并结结巴巴地说出\Quote{祖母}和\Quote{姑姑}这些词时。只有一个愿望使她渴望重见故土：那就是希望她的儿子、儿媳和孩子弃绝世俗，事奉基督。这个愿望已部分地实现了。因为她的孙女注定要出家，而她的儿媳已发誓永守贞洁，并借着信德和施舍，效法母亲为她树立的榜样。她努力在罗马展现葆拉在耶路撒冷充分展现的美德。\par

28. 我的灵魂，你怎么了？你为什么畏惧提及她的死亡？我的信已经写得过于冗长；害怕走到尽头，徒然地以为，只要闭口不言、沉浸在赞美她之中，我就能推迟那不幸的日子。此前，一路顺风，我的龙骨平稳地犁过起伏的波涛。但现在，我的言辞正撞上礁石，波涛如山，你我二人即将船毁人亡。我们不得不呼喊：\BibleQuote{师傅，救我们，我们丧亡了！}以及\BibleQuote{上主，你为何沉睡，醒来吧！}谁能眼不干地讲述葆拉临终的故事呢？她患上了一场重病，因此获得了她最渴望的——离开我们，更完全地与主结合。欧斯托基对她母亲的爱，始终真诚可靠，在这次病痛中更向大家证明了自己。她坐在葆拉的床边，为她扇风，扶住她的头，整理枕头，按摩她的脚，揉她的肚子，抚平床单，加热水，拿毛巾。事实上，她抢在仆人们之前做了一切差事，而当其中一个仆人做了什么事，她就视之为剥夺了她自己的功劳。她祈祷是多么不间断，哭泣是多么充沛，在她卧病的母亲和主的洞穴之间来回奔跑是多么频繁！她恳求天主，不要让她失去如此亲爱的同伴，如果葆拉要死，她自己也不愿再活，希望同一副棺架能把她和母亲一起送去埋葬。唉，人性的脆弱和易朽！要不是我们的信仰在基督内提升我们到天堂，并应许永生给我们的灵魂，我们的生理状况与禽兽无异。\BibleQuote{善人与恶人、好人与坏人、洁净的人与不洁净的人、献祭的人与不献祭的人，都是一样：善人如同罪人，发誓的如同怕发誓的。}\BibleRef{训 9:2}人和禽兽都一样化为尘土和灰烬。\par

29. 我为何还在拖延，藉着推迟而延长我的痛苦？保拉的理智显示她的死亡临近了。她的身体和四肢渐渐冰冷，唯有在她圣善的胸怀中，仍存续着活泼灵魂的温暖搏动。然而，仿佛她是在离开陌生人、走向自己的家人，她低声吟诵着圣咏作者的诗句：\Quote{主，我爱慕你所住的殿，和你荣耀所居之处，} 以及 \Quote{万军的上主，你的居所多么可爱！我的灵魂渴慕、甚至憔悴于上主的庭院，} 还有 \Quote{我宁愿在我天主的屋中被弃，也不愿住在恶人的帐棚中。} 当我问她为何沉默不语，不回答我的呼唤，是否感到痛苦时，她以希腊语回答说，她没有痛苦，她眼中一切都是平静安宁的。此后她不再说话，只是闭上眼睛，仿佛已蔑视所有尘世之物，并不断重复方才引用的诗句，直到她呼出最后一口气，但声音极低，我们几乎听不清她说什么。她也曾将手指举到嘴边，在她的嘴唇上画了十字圣号。然后她的呼吸停止了，她喘息待死；然而，即使当她的灵魂渴望解脱时，她也将那临死时的喉鸣（最终人人都有的）转为对主的赞美。耶路撒冷的主教和来自其他城市的一些主教也在场，还有大量下级神职人员，包括司铎和助祭。整个隐修院挤满了童贞女和隐修士的团体。保拉一听到新郎说：\Quote{我的爱卿，我的佳丽，我的鸽子，起来，走吧！因为，冬天已过，雨止天晴，} 她便快乐地回答：\Quote{地上百花齐放，修剪的时期已来到，} 和 \Quote{我深信在活人的地上要见到上主的美物。}\par

30. 她死后没有哭泣或哀号，如同世上的习俗；但所有在场的人联合起来，用各自的方言咏唱圣咏。主教们亲手扶起已故的妇女，将她放在担架上，扛在肩上送到救主山洞的教堂，把她安放在教堂中央。其他主教在队伍中手持火炬和蜡烛，还有的带领唱诗班歌唱。巴勒斯坦各城全部居民都来参加她的葬礼。没有一个隐修士藏在旷野或停留在他的小室中。没有一个童贞女仍关在她房间的隐僻处。对每个人来说，若不向如此圣善的妇女致以最后的敬意，都像是亵渎。正如在多尔卡身上所发生的，\BibleRef{宗 9:39} 寡妇和穷人们展示了保拉施舍给他们的衣服；而赤贫者大声哭喊说，他们失去了一位母亲和乳母。奇怪的是，死亡的苍白并未改变她的面容；只有一种庄重和严肃笼罩着她的脸庞。你会以为她不是死了，而是睡着了。\par

他们一首接一首咏唱圣咏，时而希腊语，时而拉丁语，时而叙利亚语；这不仅在她被安葬在教堂之下、靠近主之山洞的三天之内，而是在整周余下的日子。所有聚集的人都觉得他们在帮助自己的葬礼，并如同自己死去一般流泪。保拉的女儿，可敬的童贞女欧斯托基亚，\Quote{如断奶的幼儿依偎在母亲怀中，} 无法从她母亲身边被拉开。她亲吻她的眼睛，将嘴唇贴在她的额头上，拥抱她的身体，别无他求，只愿与她同葬。\par

31. 耶稣作证，保拉没有给她的女儿留下一分钱，反而如我先前所言，欠下了一大笔债务；更糟的是，还有一群弟兄姊妹，她难以供养，但遗弃他们又是不孝。岂有比这位高贵夫人更光辉的克己实例？她因信德的热忱，施舍了如此多的财富，以致将自己降至赤贫的地步。别人可以自夸他们用于慈善的金钱、积存在天主宝库中的巨款、以金线悬挂的还愿供物。但没有人比保拉给予穷人更多，因为保拉没有为自己留下任何东西。但现在她享有了真正的财富，以及那些眼所未见、耳所未闻、人心所未想到的美好事物。\BibleRef{格前 2:9} 如果我们哀悼，那是为我们自己，而不是为她；但即便如此，如果我们坚持为一位与基督一同为王的人哭泣，我们似乎是在嫉妒她的光荣。\par

32. 欧斯托基亚，不要害怕：你继承了一份光辉的产业。主是你的福分；并且，为增加你的喜乐，你的母亲在长期殉道之后，如今赢得了她的冠冕。不仅流血才被认作宣认：虔敬心灵的完美服务本身就是每日的殉道。两者都获得冠冕；一种是玫瑰花和紫罗兰，另一种是百合花。因此在\BibleBook{}{雅歌}中写道：\Quote{我的爱人是洁白红润的。} \BibleRef{歌 5:10} 因为，无论胜利是在和平中还是在战争中赢得，天主给得胜者同样的赏报。如同亚巴郎，你的母亲听到了这些话：\Quote{离开你的故乡、你的家族，往我指给你的地方去。} \BibleRef{创 12:1} 不仅如此，还有主通过耶肋米亚发出的命令：\Quote{你们要从巴比伦中间逃出，各人拯救自己的灵魂。} \BibleRef{耶 51:6} 直到她去世之日，她从未返回加色丁，也不曾后悔埃及的肉锅或浓味的肉食。她由她的童贞女团体陪伴，成为救主的同乡；如今她从她的小白冷升到天上的国度，可以向真正的纳敖米说：\Quote{你的民族就是我的民族，你的天主就是我的天主。} \BibleRef{卢 1:16}\par

33. 我花了两夜工夫为您口述这篇论著；在这样做时，我感受到与您同样深切的悲痛。我说\Quote{口述}，因为我无法亲自书写。每次我拿起笔，试图履行诺言，手指便僵硬，手垂落，无法控制。文辞粗陋，全无优雅或魅力，正证明作者的感受。\par

34. 保拉，再会了，请以您的祈祷帮助您这位仆人的晚年。您的信德和善工使您与基督结合；这样，站在祂面前，您将更容易获得您所求的。在这封信中，我\Quote{建造了}为您纪念\Quote{比青铜更持久的纪念碑}，任何时间流逝都无法摧毁它。我还为您的坟墓刻了一段铭文，现附于下；这样，无论我的叙述传到何处，读者都会知道您安葬在\Place{白冷}，在那里并非无人纪念。\par

\Emph{保拉墓碑上的铭文。}\par

在这墓中躺着一个西庇阿的后裔，\newline{}一个闻名遐迩的保利乌斯家族的女儿，\newline{}一个格拉古的后代，出自\newline{}阿伽门农本人的血统，显赫：\newline{}这里安息着保拉夫人，深受\newline{}父母钟爱，有欧斯托基\newline{}为女儿；她是第一位罗马贵妇，\newline{}为基督选择了艰苦和白冷。\par

洞口前方还有另一段铭文如下：——\newline{}你可看见这里在岩石中凿出的坟墓，\newline{}这是保拉的坟墓；天堂拥有她的灵魂。\newline{}她为基督放弃了罗马、朋友、财富和家园，\newline{}在这偏僻之地寻求安息。\newline{}因为这里有基督的马槽，并且有贤士们\newline{}向祂——既是天主又是人——奉献了礼物。\par

35. 圣善而蒙福的保拉，在二月朔日前的第七天，即一周的第三天，日落后，安眠于主怀。同一朔日前的第五天安葬，值霍诺留皇帝第六次执政、阿里斯泰内特斯首次执政之年。她在罗马度修德生活五年，在白冷二十年。她一生总年岁为五十六岁零八个月零二十一天。\par

\SourceRecord{Translated by W.H. Fremantle, G. Lewis and W.G. Martley.；Nicene and Post-Nicene Fathers, Second Series；Vol. 6.；Edited by Philip Schaff and Henry Wace.；Buffalo, NY: Christian Literature Publishing Co.,；1893.；<http://www.newadvent.org/fathers/3001108.htm>.}

% structure: p0001 source=h1 target=section reason=page-title

\section{第一〇九封信}

\Emph{致里帕里乌斯}\par

\EditorialNote{里帕里乌斯，阿基坦的一位司铎，曾写信告知热罗尼莫，维吉兰提乌（参见第六十一封信）正在高卢南部宣扬反对敬礼圣髑和守夜礼，而且显然得到了其主教的同意。热罗尼莫现在回信，这封信以其尖刻而非逻辑而著称。尽管如此，他提出，如果里帕里乌斯将包含维吉兰提乌异端的书寄给他，他将写一篇完整的驳斥文。里帕里乌斯随后照办，于是热罗尼莫写了《驳维吉兰提乌》一文，这是他所有作品中最极端、最缺乏说服力的一篇。这封信的日期是公元404年。}\par

1. 既然我已收到你的来信，若不回信便是骄傲，若回信便是鲁莽。因你所咨询之事，无论言说或听闻都属亵渎。你告诉我，维吉兰提乌（其名\Emph{清醒}本身自相矛盾：更应称之为\Emph{睡眠}）再次张开其臭口，向圣殉道者的圣髑倾泻污秽的毒液，并称我们这些敬奉圣髑的人为灰烬贩子和偶像崇拜者，向死人骨头致敬。不幸的恶徒！应为所有基督徒所哀哭！难道他看不见，这样说就使自己与撒玛黎雅人和犹太人同流，他们认为尸体不洁，甚至认为与死者同屋的器皿也被玷污，随从那叫人死的文字，而不随从那叫人活的圣神。\BibleRef{格后 3:6} 我们确实拒绝敬拜或朝拜——我并非说殉道者的圣髑，而是连太阳和月亮、天使和总领天使、革鲁宾和色辣芬以及\BibleQuote{不拘是在今世或来世，凡被称呼的名号}\BibleRef{弗 1:21}。因为我们不可服务受造物而非造物主，祂是永受赞美者。\BibleRef{罗 1:25} 然而我们尊敬殉道者的圣髑，为能朝拜他们为之作证的那一位。我们尊敬仆人，使他们的荣耀反射到他们的主身上，主亲自说：\BibleQuote{谁接纳你们，就是接纳我}\BibleRef{玛 10:40}。我问维吉兰提乌：伯多禄和保禄的圣髑是不洁的吗？梅瑟的身体是不洁的吗？关于他，我们得知（依照正确的希伯来文本）是由主亲自埋葬的。\BibleRef{申 34:6} 难道我们每次进入宗徒、先知和殉道者的大殿，都是在敬拜偶像的神龛吗？在他们墓前点燃的蜡烛仅仅是偶像崇拜的标志吗？我要更进一步，问一个问题，使这个理论回击到发明者的头上，要么治愈他那妄想的头脑，要么杀死它，好让纯朴的灵魂不再被他亵渎的推理所颠覆。让他回答我：主的身体被放在坟墓里时是不洁的吗？难道穿白衣的天使只是看守一具死而污秽的尸体，好让许多代以后这个嗜睡的家伙沉迷于梦境，吐出他那污秽的饱嗝，以至于像迫害者尤利安一样，要么摧毁圣人的大殿，要么把它们变成异教的庙宇？\par

2. 我感到惊讶，那位可敬的主教，据说他辖区内的这位司铎，竟默许他疯狂的讲道，而非用宗徒的棍杖，甚至铁杖，击碎这无用的器皿，把他交于毁灭肉身，好使灵魂得救。\BibleRef{格前 5:5} 他应记得所说的这些话：\Quote{你见了盗贼，就与他同伙；又与淫乱者有分。}并在另一处：\Quote{我每日清早灭绝地上所有的恶人，为从上主的城中剪除所有作恶的人。}又：\Quote{上主，憎恨你的人，我岂不憎恨他们？起来攻击你的人，我岂不恼怒他们？我以极度的憎恨憎恨他们。}如果殉道者的圣髑不值得尊敬，那么我们读到的\Quote{在上主眼中，祂圣者的逝世是珍贵的}又作何解释？如果死人的骨头玷污触摸它们的人，那么已死的厄里叟如何使另一个死人复活，生命从这位先知的身体（按维吉兰提乌的说法必是不洁的）传到后者？若是这样，以色列军队和天主子民的每个营地都是不洁的，因为他们带着若瑟和众圣祖的尸体在旷野中行走，甚至把他们不洁的骨灰带入圣地。若是这样，若瑟——我们主救主的预像——就是一个恶人；因为他大张旗鼓地把雅各伯的骸骨运到赫贝龙，只是要把他不洁的父亲放在他不洁的祖父和曾祖父旁边，即一个尸体与其他尸体放在一起。这恶徒的舌头应被割掉，或者他应被送去治疗疯病。既然他不会说话，就该学会沉默。我本人以前见过这个怪物，尽力用圣经经文捆绑这个疯子，正如希波克拉底用锁链捆绑他的病人一样；但他\Quote{走了，离开了，逃脱了，挣脱了}，逃到亚得里亚海和科提乌斯王的阿尔卑斯山之间，转而攻击我。因为愚人所说的一切不过尽是喧嚷和空谈。\par

3. 你或许在内心深处责备我这样在背后攻击一个人。我坦率承认，我的义愤压倒了我的冷静；我无法耐心听取如此亵渎的见解。我读到丕乃哈斯的投枪\BibleRef{户 25:7}，厄里亚的严厉\BibleRef{列上 18:40}，热诚者西满的嫉妒愤怒，伯多禄处死阿纳尼雅和撒斐辣时的严苛\BibleRef{宗 5:1}，以及保禄的坚定：当术士厄里玛抵挡主的道时，他使他终身失明\BibleRef{宗 13:8}。为维护天主的尊荣，没有所谓的残酷。因此法律中也说：\BibleQuote{你的手应加在他们身上，你要流他们的血，这样你将由以色列中间除去那恶。}\BibleRef{申 13:6} 我再问一次：殉道者的圣髑是不洁的吗？若是这样，宗徒们为何允许自己在逝去的斯德望的遗体——不洁的遗体——前的送葬行列中行走？为何他们为他大声哀哭\BibleRef{宗 8:2}，好让他们的悲伤转为我们的喜乐？\par

你又告诉我，维吉兰修憎恶守夜。这当然与他的名字相反。守夜者想要睡觉，不愿听从救主的话：\BibleQuote{怎么，你们竟不能同我一起守候一个时辰吗？你们要醒寤祈祷，免得陷入诱惑：心神固然切愿，但肉体却软弱。}\BibleRef{玛 26:40} 在另一处，一位先知歌唱说：\BibleQuote{我半夜起来称谢你，因为你正义的判决。} 我们也读到福音中主如何整夜祈祷\BibleRef{路 6:12}，宗徒们被囚在监里时如何彻夜守候，唱圣咏直到大地震动，狱警相信了，长官和市民充满恐惧。\BibleRef{宗 16:25} 保禄说：\BibleQuote{你们要恒心祈祷，并在这事上\Emph{醒寤}}，\BibleRef{哥 4:2} 在另一处他自称\BibleQuote{屡次不眠}。\BibleRef{格后 11:27} 维吉兰修若愿意可以睡觉，并可能在睡眠中窒息，被毁灭埃及和埃及人的毁灭者所毁灭。但让我们同达味一起说：\BibleQuote{看，那保护以色列的，既不打盹，也不睡觉。} 这样圣者和守望者便会来到我们中间。若因我们的罪过，主睡着了，让我们向他说：\BibleQuote{主，醒来，你为什么睡着？} 当我们的小船被波浪颠簸时，让我们唤醒他说：\BibleQuote{师傅，拯救我们，我们丧亡了。}\par

4. 我本可多写一些，只是信件的长度限制了我，使我保持适度的沉默。若你寄来那人的狂言集，我便可知晓需要驳斥哪些要点，本可继续下去。但实际上我只是在打空气\BibleRef{格前 9:26}，更多地显露的不是他的不信——因为这对所有人都昭然若揭——而是我的信德。但若你希望我对他作更长的反驳，就请把他那些可怜的哼唱寄来，在我的答复中，他将听到若翰洗者话语的回响：\BibleQuote{现在斧子已经放在树根上；凡不结好果子的树，就砍下来丢在火里。}\BibleRef{玛 3:10}\par

\SourceRecord{Translated by W.H. Fremantle, G. Lewis and W.G. Martley.；Nicene and Post-Nicene Fathers, Second Series；Vol. 6.；Edited by Philip Schaff and Henry Wace.；Buffalo, NY: Christian Literature Publishing Co.,；1893.；<http://www.newadvent.org/fathers/3001109.htm>.}

% structure: p0001 source=h1 target=section reason=page-title

\section{书信第一一三号}

自\Person{德斐禄}致\Person{热罗尼莫}\par

\EditorialNote{亚历山大城主教\Person{德斐禄}曾编纂一篇抨击文，矛头指向君士坦丁堡主教\Person{金口若望}，此君（很大程度上归功于他的努力）被逐出教区，流放异乡。他现在将此文寄给\Person{热罗尼莫}，请求后者将其译成拉丁文，以便在西方传播。这篇抨击文（仅存少数片段）属最为激烈之类型。然而，\Person{热罗尼莫}仍将其连同此信一并译出，此信写作年份为公元405年。书信的后半部分已散佚。}\par

致亲爱又最亲爱的弟兄\Person{热罗尼莫}，\Person{德斐禄}在主内问安。\par

1. 起初，合乎真理的判决只能令少数人满意。但主借先知发言说：\BibleQuote{我的判决显明如光。}\BibleRef{欧 6:5}而那些被黑暗恐怖包围、不能清晰理解事物本质的人，将被永恒的羞辱所遮盖，并从自己行为的结局得知他们的努力是枉然的。因此，我们也一直期望那位曾一度治理君士坦丁堡教会的\Person{若望}能得到恩宠，使他能悦乐天主，我们迟迟不愿将导致他垮台的鲁莽行为归咎于他。但是，且不提他的其他恶行，他将奥力振主义者引为亲信，提拔了其中许多人担任司铎，犯下这一罪行，使得那位天主的人、可敬记忆的\Person{厄丕法尼}——他在全世界主教中如明星闪耀——深感悲痛。因此，他理所当然地听到了审判的话：\BibleQuote{巴比伦倾倒了，倾倒了。}\BibleRef{依 21:9}\par

2. 既然知道救主说过：\BibleQuote{你们不要按外貌断案，但要按公义断案。}\BibleRef{若 7:24}……\par

\SourceRecord{Translated by W.H. Fremantle, G. Lewis and W.G. Martley.；Nicene and Post-Nicene Fathers, Second Series；Vol. 6.；Edited by Philip Schaff and Henry Wace.；Buffalo, NY: Christian Literature Publishing Co.,；1893.；<http://www.newadvent.org/fathers/3001113.htm>.}

% structure: p0001 source=h1 target=section reason=page-title

\section{第一一四号书信}

\Emph{致\Person{德奥斐洛}}\par

\EditorialNote{热罗尼莫写信给\Person{德奥斐洛}，为自己迟迟未将其后者书信（第一一三号）及反对\Person{金口若望}的抨击文章翻译成拉丁文而道歉。不过，此处所指的可能并非这些作品，而是\Person{德奥斐洛}的其他著作（例如一封逾越节书函）。他将延误归咎于\Place{巴勒斯坦}的动荡局势、严冬、普遍的饥荒以及自己的健康状况不佳。他现在寄上自己完成的译文，并恳请对方不要批评自己的作品，同时赞扬\Person{德奥斐洛}的著作，特别赞同后者关于恭敬保存举行圣体圣事所用器皿的指示。书信的写作日期为公元四〇五年。}\par

致最荣福的教宗\Person{德奥斐洛}，\Person{热罗尼莫}敬上。\par

1. 我将尊驾的论著译成拉丁文后迟迟未寄回，其原因是遇到了许多的中断和阻碍。依索里亚人突然入侵；\Place{腓尼基}和\Place{加里肋亚}遭受蹂躏；\Place{巴勒斯坦}一片恐慌，\Place{耶路撒冷}尤甚；我们都在忙于建造围墙，而非书籍。此外，还有严冬和几乎无法忍受的饥荒；这些对我这个负责照顾许多弟兄的人打击尤为沉重。在如此困难之中，我利用夜间，尽量节省或挤出时间从事翻译工作。终于完成了，到四旬期时，只剩下将誊清本与原文核对。但就在那时，一场重病袭击了我，将我带到了死亡的边缘，我仅因天主的仁慈和你们的祈祷才得以幸免；也许正是为此，我才能完成你的嘱托，并以其作者的优雅，将你那以圣经最美丽的花朵所装饰的迷人著作译出。但身体虚弱和内心的忧伤无需我说，已使我理智的锋芒迟钝，语言流畅受阻。\par

2. 我欣赏你这著作的实用目的，它旨在以圣经的权威，教导各教会中无知的人，关于他们必须以何等敬畏处理圣物并在基督的祭台前服务；并使他们铭记，用于举行主苦难的圣爵、圣盖及其他附属物，并非无生命、无感觉、缺乏圣洁之物，而是因其与主的体血相结合，应以与体血同样的敬畏加以敬礼。\par

3. 请收回你的著作，不，是我的，或者更好说，是我们的；因为当你奉承我时，你只是在奉承你自己。我的智慧为你而辛劳；我以拉丁文贫乏的资源，努力寻找与希腊文雄辩相当的表达。我并没有像熟练的译者们那样逐字翻译，也没有将你给我的钱一枚一枚地数出来；但我给你的是足量。可能缺少一些词语，但意义毫无损失。此外，我已将你寄给我的那封信译成拉丁文，并置于本书卷首，以便所有读者都知道，我是遵照尊驾的命令而行，并非鲁莽和过分自信地承担了一项超出我能力的工作。我是否成功，只有留待你来判断。即使你可能会责备我的软弱，至少会认可我的善意。\par

\SourceRecord{Translated by W.H. Fremantle, G. Lewis and W.G. Martley.；Nicene and Post-Nicene Fathers, Second Series；Vol. 6.；Edited by Philip Schaff and Henry Wace.；Buffalo, NY: Christian Literature Publishing Co.,；1893.；<http://www.newadvent.org/fathers/3001114.htm>.}

% structure: p0001 source=h1 target=section reason=page-title

\section{\WorkTitle{第一一七封信}}

\Emph{致住在高卢的一位母亲和她的女儿}\par

一位高卢隐修士在拜访白冷时，请求热罗尼莫就以下情况提供建议。他的母亲是一位教会寡妇，他的妹妹是一位发愿守贞的童贞女，但两人无法和睦相处。她们虽然分开居住，但并非各自独居。因为每人都将一个隐修士接进自己家中，名义上是管家，实际上是姘夫。应这位访客的请求，热罗尼莫现在写信给母亲和女儿，敦促她们辞退陪伴者；或至少住在一起，并指出否则必然引起的严重恶表。\par

从论著\Emph{驳维吉兰提乌斯}（§3）我们得知，恶意的批评者坚持认为，信中所描述的人和事都是虚构的，热罗尼莫写这封信只是在练习他对一个合意主题的才智。\par

日期是公元405年。\par

% structure: p0006 source=h2 target=subsection reason=relative-heading-rank; base=h2

\subsection{引言}

1. 一位来自高卢的弟兄告诉我，他的守贞妹妹和寡居母亲，虽然住在同一城市，却各有住所，并且找了神职保护人，或以客人或以管家的身份；她们这样与陌生人交往，比分开居住引起了更多恶表。我叹息着，更多以沉默而非言语表达我的感受，他说：\Quote{我恳求你，写信责备她们，使她们重归于好；使她们重新成为母亲和女儿。}我回答他说：\Quote{你给我的任务真不错，让我一个陌生人去调和两个女人，而你作为儿子和兄弟却未能影响她们。你说得好像我坐在主教的椅子上，而不是被关在隐修小室里，远离尘世的喧嚣，哀悼过去的罪过，努力避免现在的诱惑。此外，一个人隐遁不见，却让自己的舌头在世上自由驰骋，这显然是不一致的。}他回答说：\Quote{你太过胆小了，你那旧日的胆识哪里去了？正如贺拉斯论卢奇利乌斯时所说，你曾\InnerQuote{用丰富的盐巴扫荡世界}。}我接口说：\Quote{正是这个使我不敢开口。因为指控罪行，我自己反被定为罪犯。人们争论并否认我的论断，以致如谚语所说，我几乎不知道自己是否还有耳朵或感觉。四面墙壁都回荡着针对我的咒骂，\InnerQuote{我成了醉汉的歌谣}。在不愉快经验的逼迫下，我学会了沉默，认为最好在口前设看守，守住嘴唇的门，而不让我的心偏向任何恶事，或在指责他人过错时，自己却陷入诽谤的罪。}他回答说：\Quote{说实话并非诽谤。你也不会因施行一次特定的责备而教训世界；因为很少有人，如果有的话，会面临这一特殊的指控。因此我请求你，既然我不辞辛劳长途跋涉，不要让我白来一趟。主知道，在看到圣地之后，我来的主要目的就是通过你的信治好我妹妹和我母亲之间的分裂。}我回答说：\Quote{好吧，我按你的意愿去做，毕竟信是写给海外的人，针对特定人物的言辞很少会冒犯他人。但我必须请你对我所说的话保密。你将带着我的建议上路，作为鼓励；如果被听从，我将与你一样欢喜；而如果，如我所想，被拒绝，我就白费了口舌，你也白走了远路。}\par

% structure: p0008 source=h2 target=subsection reason=relative-heading-rank; base=h2

\subsection{信函正文}

2. 首先，我的妹妹和我的女儿，我希望你们知道，我写信给你们，并非因为我怀疑你们有什么邪恶。相反，我恳求你们和睦相处，以免给人任何这种怀疑的借口。此外，倘若我以为你们被罪紧紧束缚——愿这离你们很远——我绝不会写信，因为我知道我的话是对聋子说的。再者，如果我写信给你们有些严厉，我请求你们不要归咎于我的苛刻，而是归咎于我所尝试治疗的疾病的性质。当坏疽一旦形成，烙铁和刀是唯一的疗法；毒药是已知的唯一解毒剂；巨大的痛苦只能通过施加更大的痛苦来缓解。最后，我必须说这一点：即使你们自己的良心宣告你们无罪，但这样的谣言本身也给你们带来耻辱。母亲和女儿是亲情的名称；它们意味着天然的联系和相互的责任；它们构成了人与人之间最亲密的关系，仅次于灵魂与天主的关系。如果你们彼此相爱，你们的行为不值得称赞；但如果你们彼此仇恨，你们就犯了罪。主耶稣曾服从祂的父母。\BibleQuote{\BibleRef{路 2:51}} 祂尊敬那位祂自己作为祂父母的母亲；祂尊重那位祂自己所抚养的养父；因为祂记得自己曾一度被怀在母腹中，一度被抱在怀中。因此，当祂悬在十字架上时，祂将母亲托付给祂的门徒，\BibleQuote{\BibleRef{若 19:26}} 这位母亲在祂受难前从未与祂分离过。\par

3. 好了，我不再对母亲多说什么了，因为也许年龄、衰弱和孤独足以成为她的借口；但对你，女儿，我要说：难道母亲的房子对你来说太小了，而你的子宫却从未觉得太小吗？既然你曾与她同住十个月，难道不能忍受与她同住一天吗？还是你无法面对她的目光？难道生你养你、抚育你、了解你的母亲，竟让你害怕作为你家庭生活的见证人？如果你是一个真正的童贞女，为什么害怕她细心的守护？而如果你已经失足，为什么不公开结婚？婚姻就像提供给遭遇海难之人的一块木板，你可以借此补救之前所犯的过错。我不是说你不必为你的罪悔改，也不是说你可以继续行恶；但是，一旦形成了这种关系，我就对完全打破它感到绝望。不过，回到你母亲身边，将使你更容易为你因离开她而失去的童贞哀痛。或者，如果你仍无玷污、未失贞洁，那就小心守护它，因为你可能会失去它。为什么你必须住在一个每天都要为生死挣扎的房子里？谁能与毒蛇同眠而安然无恙？不会的，因为即使它不攻击他，也必定会惊吓他。最好待在没有危险的地方，而不是身陷危险却逃脱。前者是平静，后者则需要小心掌舵。前者充满喜乐，后者只是避免悲伤。\par

4. 但你也许会回答：\Quote{我的母亲品行不端，她贪恋世俗的事物，她喜爱财富，她无视斋戒，她用锑剂染污自己的眼睛，她喜欢穿着艳丽的服装外出，她阻止我发隐修愿，所以我不能与她同住。}但是首先，即使她如你所说，你不愿因她的一切过错而离弃她，你将获得更大的赏报。她曾怀你在腹中，抚育你；她以温柔的情感忍受你童年的烦扰。她为你洗涤衣物，在你生病时照顾你，而做母亲的病痛不仅因你承受，更是由你引起。她将你养育成女子，教导你爱慕基督。你母亲既已奉献你为童贞女，事奉你的净配，你就不应对她的行为不满。然而，如果你无法忍受她的矫饰，觉得必须避开它们，而且你的母亲确实如人们常说的那样，是一个世俗女子，那么你还有其他人，像你一样的童贞女，贞洁的神圣同伴。为什么当你离弃母亲时，你却选择了一个也许已抛下自己姐妹和母亲的男人作伴？你告诉我她难相处，而他随和；她好争吵，而他友善。我要问你一个问题：你是直接从家投奔那人的，还是后来才遇到他的？如果你是直接去找他，你离开母亲的原因就很明显了。如果你是后来才遇到他，你的选择就表明你在母亲屋檐下缺失了什么。我施加的痛苦是剧烈的，我和你一样感到刀割。\BibleQuote{行走正直的，行路稳妥}\BibleRef{箴 10:9} 只因我的良心会责备我，我本应保持沉默，并在我自己并非无罪之处慎于指责他人。我眼中有大梁，就不愿看见邻人眼中的木屑。但事实上，我远居在基督徒弟兄中间；我的生活是光荣的，目击者可以作证；我很少见别人，也很少被人看见。因此，你既然声称以我为榜样，却拒绝在节制上效法我，这是最无耻的。如果你说：\Quote{我的良心也足以使我满足。天主是我的审判者，祂是我生命的见证。我不在乎人们说什么；}那么让我以宗徒的话敦促你：\BibleQuote{在人前做诚实的事}不仅在天主眼前，也在\BibleQuote{众人眼前。}\BibleRef{罗 12:17} 如果有人因你是基督徒和童贞女而挑剔你，不要在意；或许有人说你离开母亲住进隐修院的童贞女中，但这样的指责正是你的荣耀。当人们指责天主的侍女并非因为放纵，而只是不近人情时，他们所谴责的冷酷无情的弃亲行为，其实是对天主的热忱。因为你宁愿选择祂，超过你的母亲，而祂正是你被命令要优先于你自身灵魂去爱的。\BibleRef{路 14:26} 如果有一天她也这样优先选择祂，她将在你身上找到的不仅是女儿，还有姐妹。\par

5. \Quote{那么？}你会说，\Quote{与我同住一位修道人，难道是罪过吗？}你抓住我的衣领，拉我进法庭，要么认可我所不赞同的事，要么招致许多人的厌恶。修道人从不使女儿离开母亲。他欢迎两者并尊重两者。女儿可以随心所欲地虔诚；然而，一位寡妇母亲是她贞洁的保证。如果此人无论他是谁，与你年龄相仿，他应尊敬你的母亲如同自己的母亲；如果年长，他应爱你如女儿，并使你服从母亲的管教。他喜欢你超过你母亲，这对你的名声和他自己的名声都不好；因为他的情感可能看起来更多的是爱你的年轻，而不是爱你本人。如果那位隐修士不是你的兄弟，而你又没有能保护你的亲属，这就是我要说的话。但是，一个陌生人有什么理由介入既有母亲又有兄弟的家庭，母亲是寡妇，兄弟是隐修士？对你来说，意识到自己既是女儿又是姐妹，是好的。然而，如果你两者都做不到，而母亲又太难对付，至少你的兄弟应该使你满足。或者，如果他太严厉，那位生养你的母亲或许会更温柔。你为什么脸色苍白？为什么激动？为什么脸红，颤抖的嘴唇泄露了你内心的不安？只有一样东西能超越女子对母亲和兄弟的爱，那就是她对丈夫的热情。\par

6. 我更听说，你常与亲戚和密友一同前往郊区别墅及其怡人的花园。我毫不怀疑，是某位表姐妹或亲戚为了自己的满足，把你当作一种新奇的随从带着。我绝不愿意怀疑你渴望男人的陪伴，即使他们是你自己家族的成员。但是，请告诉我，姑娘，你是独自出现在亲戚面前吗？还是带着你的宠儿？即便你可能无耻，你也不敢公然在他人的眼前炫耀他。如果你真这么做，你整个圈子都会对你和他大嚷大叫；每个人都会指着你；你的表姐妹——她们在你面前为了取悦你而称他为隐修士和虔诚之人——会在背后嘲笑你，因为你竟有这样一个不合自然的丈夫。另一方面，如果你独自出门——这我更认为是事实——你会发现，自己穿着朴素的衣服，置身于少年奴仆、已婚或即将结婚的女子、放荡的少女以及长发紧身衣的花花公子之中。某个留胡子的花花公子会向你伸出手，你若感到疲倦，他会扶住你，他手指的紧握要么是对你的诱惑，要么表明你是对他的诱惑。再说，当你与已婚男女一同入席时，你将不得不看到与你无关的亲吻，品尝并非为你准备的菜肴。此外，看到其他女子身着丝绸衣裳和金色锦缎，不可能不给你带来伤害。在餐桌上，无论你愿不愿意，你都将被迫吃肉，而且是各种不同的肉。为了让你喝酒，他们会称赞酒是天主的受造物。为了诱使你沐浴，他们会厌恶地谈论污秽；而当你经过再三考虑依从时，他们会齐声称赞你纯洁、开朗，是一位十足的淑女。与此同时，某位歌者会向在座的人献上柔美流畅的曲调；由于他不敢直视他人的妻子，他会不断地把目光投向你——这个没有保护者的人。他会用点头示意，用语调传达他不敢说出口的话。置身于如此明显的感官诱惑中，即使是铁石心肠也容易被欲望征服；在贞女中这种欲望尤为强烈，因为她们认为最甜蜜的正是她们未曾经历过的。异教传说告诉我们，水手们竟会驾驶船只撞上岩石，只为聆听\OriginalTerm{塞壬}{Sirens}的歌声；\OriginalTerm{俄耳甫斯}{Orpheus}的竖琴能吸引树木和动物，软化燧石。在宴席厅中，贞洁难以持守。闪亮的肌肤显露出罪污的灵魂。\par

7. 作为一个学童，我曾读到过一个人——并且在大街上看到过他的逼真肖像——他即使在骨瘦如柴时仍继续怀有不法的情欲，他的疾病直到死亡才治愈。那么，你一个身体健全、娇嫩、健壮、红润的年轻女子，若让自己在肉食、美酒和沐浴中自由放任，更别提那些已婚男子和单身汉，你会变成什么样子？即使你在被请求时拒绝同意，你也会把人家来问你这件事本身当作你被认为漂亮的证据。贪欲的心灵追求不名誉的目标比追求名誉的目标更为热切；当一件事被禁止时，它更乐在其中。你穿的衣服，即使廉价而黯淡，也是你内心情感的标志。因为它没有褶皱，拖在地上，让你显得比实际更高。你的外衣特意撕裂以显露里面，在隐藏丑陋的同时展现美丽。你行走时，黑色闪亮的鞋子发出的吱嘎声吸引着年轻人的注意。你穿着束胸以保持胸部位置，一条紧绷的腰带紧束你的胸膛。你的头发遮住前额或耳朵。有时你还会让披肩滑落，露出白皙的肩膀；然后，仿佛不愿被人看见，你迅速遮住你故意显露的东西。在公共场合，你为了谦逊而遮住脸，但像老练的娼妓一样，你只展示可能取悦人的部分。\par

8. 你会喊道：\Quote{你怎么知道我的样子？或者，你离得那么远，怎么能看见我在做什么？}是你亲兄弟的眼泪和哭泣告诉我的，他时常、几乎无法忍受的悲痛爆发。但愿他说了谎，或者他的话只是出于担心，而非控告。但是，相信我，说谎者不会流泪。他愤慨的是，你偏爱一个年轻人——他确实不穿丝绸，不留长发，但在肮脏中却肌肉发达、娇嫩——而不偏爱他；这个家伙掌管钱袋，料理纺织，分配仆人任务，管理家务，从市场采购一切所需。他既当管家又当主人，由于他抢在仆人之前履行职责，被所有仆从指责。主母没有给他们的东西，他们都宣称被他偷走了。仆人这类人满腹牢骚；无论你给他们什么，总是不够。因为他们不考虑你有多少，只考虑你给了多少；他们只能以抱怨的方式弥补他们的懊恼。有人称他为食客，有人称他为骗子，有人称他为谋财者，还有人用他一时兴起的新奇绰号。他们四处张扬说，他总是在你的床边，当你生病时他跑去请护士，他端盆子、晾床单、为你折叠绷带。世人总是乐于相信丑闻，在家中编造的故事很快会在外面传播。如果连你的男仆和女仆都编造这样的故事，你也不必惊讶，因为连你的母亲和兄弟也在抱怨你的行为。\par

9. 因此，请照我所劝告并恳求你做的去做：如果可能，与你的母亲和解；或者，如果这不能做到，至少与你的兄弟达成协议。或者，如果你对通常如此亲密的亲属关系充满了无法缓和的仇恨，那么无论如何与那人分开，据说你偏爱此人胜过你的骨肉；如果这对你也不可能（因为如果你能离开他，你肯定会回到你的亲人身边），那么在与他作为同伴同居时，更要注意外表。住在一栋独立的建筑里，分开用餐；因为如果你与他同住一屋檐下，诽谤者会说你们共享床榻。这样，当你觉得需要时，你可以轻易地得到他的帮助，同时也在相当程度上避免公众的非议。然而，你必须小心不要沾染耶肋米亚告诉我们的那种污渍，他说任何碱或洗衣皂都不能洗掉它。\BibleRef{耶 2:22} 当你希望他来见你时，永远要有证人在场；无论是朋友、被释奴或奴隶。好的良心不惧怕任何人的目光。让他不受拘束地进来，自在地离开。让他沉默的眼神、未说出口的话语以及整个举止，虽然有时可能显得尴尬，却表明内心的平安。求你张开耳朵，聆听全城的呼声。你们俩都已经失去了自己的名字，彼此以对方的名字为人所知。你被称为他的，他被说成是你的。你的母亲和你的兄弟已经听说了，并准备将你接纳在他们中间。他们恳求你同意这个安排，以便你与此人亲密关系的丑闻——仅限于你个人——让位于大家共同的荣耀。你可以与你的母亲同住，他与你的兄弟同住。你可以更勇敢地尊重你兄弟的同伴；你的母亲会更恰当地珍视她儿子的朋友，而非女儿的朋友。但如果你皱眉并拒绝接受我的建议，这封信将公开与你争辩。它会说：\Quote{为什么你纠缠别人的仆人？为什么你使基督的仆役成为你的奴隶？看看人们，审视你面前每一张面孔。当他在教堂里诵读时，所有人的目光都聚焦于你；而你，以妻子的放纵，以你的耻辱为荣。秘密的丑行不再满足你；你称大胆为自由；\InnerQuote{你有一个娼妓的面容，不肯羞愧。}}\par

10. 你再次惊呼我多疑、心怀恶意、太容易听信谣言。什么？我多疑？我心地不善？我，正如开头所说，提笔是因为我并无疑虑？还是你粗心、放荡、傲慢？你，二十五岁的年纪，却把一位胡子尚未长成的青年网罗在你的怀抱中？他一定是个极好的导师，无疑能用他严厉的目光警告并恐吓你！没有年龄对情欲是安全的，但白发是体面行为的某种保障。总有一天会到来（因为时间悄然流逝），你那位英俊的年轻宠儿会找到一位更富有或更年轻的情妇。因为女人老得快，尤其是与男人同居时。你将会后悔你的决定，并在你的财富和名誉都已丧失、这段不幸的亲密关系终于幸福地结束时，为你当初的固执而懊悔。但也许你对自己的立场很有把握，认为在感情已经发展并成长了如此之久的情况下，无需担心破裂。\par

11. 对你，她的母亲，我也要说几句话。你的年龄使你超出丑闻的波及范围；不要利用这一点来放纵于罪。更恰当的是，你的女儿应该从你那里学习如何与同伴分离，而不是你从她那里学习如何放弃一个情夫。你有一个儿子，一个女儿，还有一个女婿，或者至少是女儿的伴侣。那么，为什么你要寻求其他社会关系，或者希望重新点燃即将熄灭的火焰？你遮掩女儿的过错，比以此为借口行你自己的不义更为得体。你的儿子是隐修士，如果他与你同住，会使你在你的宗教誓愿和守寡誓言上更加坚强。为什么你要接纳一个完全陌生的人，尤其是一个不足以容纳儿子和女儿的房子里？你已到了有孙辈的年龄。那么，把这对夫妻邀请回家吧。你的女儿独自离开了；让她与这个男人一起回来。我说\Quote{男人}而非\Quote{丈夫}，以免有人吹毛求疵。这个词描述他的性别，而非他与她的关系。或者，如果她羞于接受你的提议，或者觉得她出生的房子对她来说太狭窄，那么你们俩都搬到她的住处。无论其住所多么有限，容纳一位母亲和一位兄弟总比容纳一个陌生人好。事实上，如果她与一个男人住在同一所房子、占用同一个房间，她无法长久保持贞洁。两个女人和两个男人同住则不同。如果相关的第三人——我指的是那个支持你晚年的人——不愿加入，只会引起不和与混乱，那么你们俩可以不要他。但如果你们三人继续同住，那么你的兄弟和儿子将向他提供一位姐妹和一位母亲。其他人可能会称这两个陌生人为继父和女婿；但你的儿子必须称他们为养父和兄弟。\par

% structure: p0019 source=h2 target=subsection reason=relative-heading-rank; base=h2

\subsection{注释}

12. 我迅速工作，在一夜之内完成了这封信，既急于取悦一位朋友，也想在一个修辞主题上试试身手。然后清晨，他在将要出发时敲了我的门。我也想向我的诋毁者表明，像他们一样，我也可以随口说出想到的话。因此，我很少引用圣经，也没有像在我的大多数著作中那样，在论述中编织它的花朵。事实上，这封信是即兴口授，在灯光下如此迅速地倾泻而出，以至于我的舌头赶超了秘书的笔，我的流畅击败了速记的方法。我说这些，是为了让那些不宽容才能缺陷的人，也许能对时间的缺乏有所宽容。\par

\SourceRecord{Translated by W.H. Fremantle, G. Lewis and W.G. Martley.；Nicene and Post-Nicene Fathers, Second Series；Vol. 6.；Edited by Philip Schaff and Henry Wace.；Buffalo, NY: Christian Literature Publishing Co.,；1893.；<http://www.newadvent.org/fathers/3001117.htm>.}

% structure: p0001 source=h1 target=section reason=page-title

\section{第一一八封信}

\Emph{\Person{犹利安}}\par

\EditorialNote{热罗尼莫写信给\Person{犹利安}，一位显然来自\Place{达尔马提亚}的富裕贵族（§5），安慰他因最近妻子和两个女儿去世而遭受的损失。他提醒\Person{犹利安}记得\Person{约伯}的考验，并推荐他效法这位先祖的耐心。他也敦促他效法\Person{帕玛丘}和\Person{保利诺}所立的榜样，即为基督的缘故舍弃财富，成为隐修士。此信的日期是公元406年。}\par

1. 在\Person{奥索尼乌斯}——待我如子，亦为尔之兄弟——刚刚离开之际，他匆匆让我瞥见一眼，便又急忙离去，同时道别与问安。然而他以为，若不能带些我仓促写就的片言只语给尔，此行便空手而归。他身穿符合其身份的朱红服饰，已束好剑带，并下令备好驿站马匹。尽管如此，他仍让我召来秘书，向他口述一封信。我如此迅速行事，以至于他敏捷的手几乎跟不上我话语的节奏，无法记下我急促的句子。如此仓促的口述取代了精心书写；我若打破长久沉默，也不过是向尔表达善意。这是一封即兴的信件，既无逻辑顺序，亦无文采魅力。尔权且将我视作朋友；我保证，尔在此找不到任何演说家的痕迹。请记住，这信是灵机一动草就，作为旅途的干粮赠予急于上路之人。\par

圣经说：\Quote{不合时宜的故事如同哀悼中的音乐。}\BibleRef{德 22:6} 因此我鄙弃修辞的优雅和那些使少年着迷的口才魅力，而回归圣经的严肃语气。因为在其中可找到真正的疗伤药和忧愁的可靠治疗。在这些经文中，一位母亲甚至从担架上得回她的独子。\BibleRef{路 7:11} 在这些经文中，哀悼的人群听到：\Quote{这女孩不是死了，而是睡着了。}\BibleRef{玛 9:24} 在这些经文中，一个死了四天的人，因主的呼唤而裹着殓布出来。\par

2. 我听说在短时间内你遭受了几次丧亲之痛，你接连埋葬了两个年轻的未婚女儿，而\Person{福斯蒂纳}，你最忠贞贤淑的妻子，在信德热忱中如同你的姊妹，也是你失去子女后的唯一安慰，突然安眠主怀，被夺走了。你如同一个遭遇海难的人，刚上岸就落入强盗之手，或者用先知的雄辩话说，如同一个人\Quote{躲避狮子，却遇到了熊；或是进了房屋，以手扶墙，却被蛇咬。}\BibleRef{亚 5:19} 金钱损失紧随你的丧亲之痛；整个省份被蛮族敌人蹂躏，在普遍的毁坏中，你的私有财产被摧毁，你的牛羊被抢走，你可怜的奴隶或被俘或被杀害。最糟的是，你唯一的女儿（因失去其他孩子而更加珍贵）嫁给了一位年轻贵族，暂且不往坏处说，他给你的悲伤多于喜悦。这就是加在你身上的考验清单；这就是古老的仇敌向基督旗下新兵\Person{犹利安}发起的战斗。如果你只看自己，你的痛苦确实很大；但如果你看看那位强大的战士，\BibleRef{默 19:11} 它们不过是儿戏，战斗只是表面的。在无数考验之后，受祝福的\Person{约伯}身边仍留着一个恶妻，魔鬼希望他能从她那里学会亵渎天主。而你却失去了一位出色的妻子，为的是让你学会在不幸的时刻没有安慰。然而，忍受一个你不喜欢的妻子远比哀悼一个你所深爱的妻子更难。而且，当\Person{约伯}的子女去世时，他们在他房屋的废墟下找到共同的坟墓；他所能做的显示父爱的一切，就是撕裂外衣，俯伏在地朝拜，说：\Quote{我赤身脱离母胎，也要赤身归去；主赐的，主收回；主所喜悦的，就如此成就了；愿主的名受赞美。}\BibleRef{约 1:20} 但简而言之，你却得以为你所爱的人举行葬礼；这些葬礼有许多尊敬的亲戚和安慰的朋友参加。同样，\Person{约伯}立即失去了所有财富；当报信的使者一个接一个地展开新的灾祸时，他毫不畏缩，正如智者\Person{贺拉斯}所写：——\par

若你愿将世界击为碎片，\newline{}那被砸之人亦无所畏惧。\par

但你的情况不同。你的大部分财产得以保留，你的考验并未超过你的承受能力。因为你尚未达到如此的完美，以致于魔鬼需要调集所有力量来攻击你。\par

3. 很久以前，这位富有的家主和更为富有的父亲，因突如其来的打击而变得穷困和丧亡。但尽管遭遇这一切，他并没有在天主面前犯罪，也没有口出愚言，于是主——因祂仆人的胜利而欢欣，并将\Person{约伯}的忍耐视为自己的胜利——对魔鬼说：\Quote{你曾用心察看我的仆人约伯没有？地上再没有人像他那样完美正直，敬畏天主，远离邪恶。他仍然持守他的纯正。}\BibleRef{约 2:3}祂巧妙地加上最后一句，因为无辜者在遭受不幸压制时，很难不抱怨；也难避免在看到自己无故受苦时使信德失事。魔鬼回答主说：\Quote{以皮换皮，人肯舍去一切所有的保全性命。但你伸手伤他的骨和肉，他必当面诅咒你。}\BibleRef{约 2:4}你看仇敌何等狡猾，罪恶的日子使他多么刚硬！他知道外在事物与内在事物的区别。他知道连世上的哲学家也称前者为\OriginalTerm{无关紧要的}{\GreekText{ἀ} \GreekText{διάφορα}}，即无关紧要的，而德性的完美并不在于失去或轻视它们。后者才是内在的、更可取的，失去它们必然引起懊恼。因此他大胆反驳天主的话，宣称\Person{约伯}根本不值得称赞；因为他没有舍弃自己任何部分，只是舍弃了身外之物，他为了自己的皮而付出了儿女的皮，他只是放下钱袋来保全身体的健康。由此你的智慧可以看出，你的考验目前只达到了以皮换皮、以皮换皮的地步，你愿意为你的生命舍弃一切。主还没有向你伸手，没有伤你的肉，没有折断你的骨头。然而，当这样的苦难临到你时，很难不呻吟，不\Quote{祝福}天主，即当面咒骂祂。在\BibleBook{列王}{纪}中，\Quote{祝福}一词也有同样的用法，说\Person{纳波特}\Quote{祝福}天主和君王，因此被百姓用石头打死。但主认识祂的斗士，确信这位大英雄即使在这最后也是最严峻的搏斗中也会是不可战胜的。因此祂说：\Quote{看，他已经在你手中；只要保留他的性命。}\BibleRef{约 2:6}圣人的肉体被交在魔鬼手中，但他的生命力被扣留。因为如果魔鬼击打了感觉和心智判断所依赖的部分，那么滥用这些能力的罪责就不会落在犯罪者身上，而是落在扰乱他心智平衡者身上。\par

4. 别人若愿意，可以称赞你，颂扬你胜过魔鬼的胜利。他们可以歌颂你以微笑面对女儿们的丧失，或者歌颂你在她们安息四十天后，为参加殉道者遗骨的祝圣礼，毅然换上白衣，不再在意全城人都关注的丧亲之痛，而只关心分享殉道者的胜利。不，他们说，当你将妻子安葬时，你并非视她为死者，而是视她为踏上旅途的人。但我不会用谄媚之言欺骗你，也不会用滑溜的赞美使你失足。相反，我要说那对你有益的话：\Quote{我儿，你若要侍奉主，就当预备你的灵魂接受试探。}\BibleRef{德 2:1}以及\Quote{你们做完了一切所吩咐的事，只当说：\InnerQuote{我们是无用的仆人，所做的本是我们应当做的。}}要向天主说：\Quote{祢从我这里取去的儿女原是祢自己的恩赐。祢取去的那位婢女，也曾暂时借给我作安慰。我并不因祢收回而烦恼，反而感恩祢先前将她赐给我。}\par

从前有一个富有的少年人，自夸已满全了法律的一切要求，但主对他说（正如我们在福音中所读到的）：\BibleQuote{你还缺少一件：你若愿意成全，就去变卖你一切所有的，分给穷人，然后来跟随我。}\BibleRef{谷 10:21} 那宣称已做到一切的人，在财富的力量第一次冲击下就退让了。因此，富有的人难以进入天国，而天国渴望其子民拥有超脱一切束缚和障碍、高飞的灵魂。\BibleQuote{你去}，主说，\BibleQuote{变卖}不是一部分财产，而是\BibleQuote{你一切所有的，分给穷人}，不是给你的朋友、亲戚或亲属，也不是给你的妻子或儿女。我更要进一步说：不可因害怕日后贫穷而为自己留下任何东西，否则你就分担了\Person{阿纳尼雅}和\Person{撒彼辣}的谴责；\BibleRef{宗 5:1} 反而要将一切施舍给穷人，并用\OriginalTerm{不义的钱财}{mammon of unrighteousness}为自己结交朋友，以便他们接你们进入永久的帐幕。\BibleRef{路 16:9} 听从师傅的诫命\BibleQuote{跟随我}，\BibleRef{玛 9:9} 并以世界的主为你的产业；这样你才能与先知一同歌颂\BibleQuote{上主是我的福分}，并如同一位真正的肋未人，不拥有地上的产业。\BibleRef{户 18:20} 我不能不这样劝告你，如果你愿意成全，如果你想达到宗徒荣耀的顶峰，如果你想背起你的十字架来跟随基督。一旦你把手扶在犁上，就不可回头；\BibleRef{路 9:62} 一旦你站在屋顶上，就不要再想屋内的衣服；为了逃避你的埃及女主人，\BibleRef{创 39:12} 你必须丢弃属于这世界的外衣。连\Person{厄里亚}在迅速升天时，也不能带走他的外衣，而是将世上衣服留在了世上。你或许会反对说，这样的行为属于那些效法宗徒、追求成全的人。但你为何不追求成全呢？为何不呢？你在世上居于首位，为何不能也在基督的家庭中居于首位呢？难道因为你结了婚？\Person{伯多禄}也曾结婚，但他舍弃了船和网，也舍弃了妻子。主愿意所有的人都得救，并且更愿罪人悔改胜过死亡，祂以其全能的眷顾，为你除去了这个藉口。你的妻子不能再把你拉向尘世，但你却可以跟随她，因她把你拉向天国。为那些先你归主的孩子预备好处。不要让他们的一份增多姊妹的财富，而要用来赎回你的灵魂，并维持穷人的生活。这是你的女儿们期望从你那里得到的项链；这是她们希望看到在额头上闪耀的珠宝。那本会浪费在购买丝绸上的钱，若用来为穷人提供廉价的衣服，便可视为节省了。她们向你索要自己的那份。现在她们已与配偶结合，不愿显得贫穷卑微：她们渴望有配得起她们地位的装饰。\par

5. 你也不能以自己的高贵身份和财富责任为借口。请看\Person{Pammachius}和那位充满信德的司铎\Person{Paulinus}，他们二人不仅将自己的财富，连他们自身都献给了主。尽管魔鬼百般阻挠，他们绝没有\Quote{以皮换皮}，而是将自己的血肉、甚至灵魂都奉献给了主。他们确实能藉着榜样和宣讲——即言行——引导你走向更高的事物。你出身高贵，他们也是如此；但在基督内，他们变得更加高贵。你富有且受人敬重，他们一度也是如此；但如今，他们不是富有和受人敬重，而是贫穷和默默无闻；然而，正因为是为主基督的缘故，他们实际上更加富有、更加出名。诚然，你也表现出了慈善，据说你周济圣徒的需要，款待隐修士，并向教会大量捐献。但这只是你作为基督士兵的入门功夫。你轻视钱财；世上的哲学家也曾这样做。其中一位——且不提其余——将许多财产的价值抛入海中，说道：\Quote{沉下去吧，你们这些引发邪恶欲望的诱饵；我要将你们溺死在海中，免得你们把我溺死在罪里。}如果一位哲学家——一个虚荣之物，能为大众的掌声所收买和出卖——能一次卸下他所有的负担，那么当你只交出自己的一部分时，怎能认为自己已达到德行的顶峰呢？主所要的是你自己，是\BibleQuote{活生生的祭献……悦乐天主的}\BibleRef{罗 12:1}。我说的是你自己，而非你的财物。因此，正如他曾藉着许多灾祸苦难训练以色列人，如今他也藉着各种考验来警醒你。\BibleQuote{因为主所爱的，他必惩戒；他所接纳的每个儿子，他都鞭打。}\BibleRef{希 12:6}那穷寡妇只向银库投了两个小钱；但因为她投了她所有的一切，经上说她向主献的供物超过了所有富人\BibleRef{谷 12:43}。这样的供物，不是按重量，而是按奉献的善意来衡量的。你可能已为许多人花尽了资财，你的一部分同胞也有理由为你的慷慨而欢欣；但那些从未从你手中得到任何东西的人仍然更多。无论是\Person{Darius}的财富还是\Person{Crœsus}的财宝，都不足以满足世上穷人的需要。但是，一旦你把自己献给主，并决心以宗徒德性的圆满跟随救主，那么你就会明白你迄今为止的位置，以及你在基督军中是如何落后的。你刚为死去的女儿哀悼，对基督的敬畏就已擦干了你慈父面颊上的泪水。这确实是信德的一次伟大胜利。但\Person{Abraham}所赢得的胜利更大，他甘愿杀死自己的独生子——他曾被告知这儿子要继承世界，却仍希望\Person{Isaac}死后会复活\BibleRef{希 11:17}。\Person{Jephthah}也献上了自己的童贞女儿，为此他被宗徒列在圣人之中。因此，我不愿你只向主献上那些可能被盗贼偷走、被敌人夺取、被充公没收的财物，那些价值起伏不定、时涨时落的东西，那些属于一个又一个主人、如海浪般迅速更迭的东西——总而言之，那些无论你愿意与否，死时都必须留下的东西。相反，你要向主献上那没有敌人能夺走、没有暴君能强取、能随你进入坟墓、甚至进入天国和乐园之乐的东西。你已经建造隐修院，并在\Place{Dalmatia}的各岛屿上供养大批圣人。但如果你自己作为圣人生活在这些圣人中间，那就更好了。\BibleQuote{你们应是圣的，因为我是圣的。}宗徒们夸耀自己舍弃了一切跟随了救主\BibleRef{路 18:28}。我们没有读到他们舍弃了什么，除了他们的船和网；然而他们却得到了那要审判他们的主的赞许。为什么？因为当他们奉献自己时，他们确实舍弃了自己所有的一切。\par

6. 我说这一切，并非贬低你的善工，也非低估你慷慨施舍，而是不愿你在世俗人中作隐修士，而在隐修士中作世俗人。我要求你作出一切牺牲，因为我听说你的心意已奉献于服务天主。若有哪位朋友、追随者或亲属试图反驳我这劝言，诱你回到华筵之乐，要确知他关心的不是你的灵魂，而是自己的肚腹；并记住：死亡转瞬即终结优雅的宴席及财富所能提供的一切享乐。短短二十天内，你已失去两个女儿，一个八岁，一个六岁；难道你以为像你这样年事已高的人还能活很久吗？\Person{达味}告诉你，你能期望多久：\BibleQuote{我们一生的年日是七十岁，若是强壮可到八十岁，但其中所矜夸的不过是劳苦愁烦。} 有福的人，且堪受至福的，是老迈时仍为基督的仆人，末日时仍为救主的事业奋战的人。\BibleQuote{他在城门口与仇敌说话时，必不至羞愧。} 当他进入乐园时，有人会对他说：\BibleQuote{你一生中受了恶事，但如今你无处得安慰。} 主不会两次追讨同一罪过。从前那贫穷、满身疮痍的\Person{拉匝禄}，狗来舔他的疮，他只能靠从富人桌上掉下的碎屑勉强维生的可怜人，如今被接入\Person{亚巴郎}的怀抱，在这位大族长身上寻得慈父之乐。一个人既享受今世的美物，又享受来世的美物，既满足此世的肚腹，又满足彼世的精神，既从此生的欢乐转入彼世的欢乐，既在两个世界都居首位，又在地上有尊荣，在天上有荣耀，这是困难的，不，是不可能的。\par

7. 若你内心困恼，因为给你这劝言的我本人并非你所希望我成为的样子，又因为你也曾见过有些人开始得好，却在半途跌倒；我要简短辩护：我所说的不是我的话，而是主和救主的话；我敦促你的，不是我自己能达到的标准，而是每位真正的基督仆人所当渴望并实现的理想。运动员通常比他们的支持者更强壮；然而较弱的一方却催促较强的一方尽力而为。不要看\Person{犹达斯}否认他的主，而要看\Person{保禄}宣认他。\Person{雅各伯}的父亲家资巨富；然而，当\Person{雅各伯}前往\Place{美索不达米亚}时，他孤身一人、一无所有，仅靠手杖行走。当他疲惫时，他不得不在路边躺卧；虽然自幼受母亲\Person{黎贝加}的娇养，却不得不以石为枕。但正是在那时，他看到了从地通天的梯子，天神在梯子上上去下来，主在梯顶，向跌倒者伸出援手，并藉显现自己鼓励攀登者再接再厉。因此那地方被称为\Place{贝特耳}，即天主的殿；因为那里天天有上上下下。当圣人们疏忽时，也会失足；罪人若以泪水洗净污点，也能复得原位。我说这些，不是要让下来的人吓倒你，而是要让上去的人激励你。因为恶绝不能成为模范，即使在世俗事务中，德行的动力也总是来自更光明的一面。\par

但我忘了我的初衷和这封信的篇幅限制。我本愿多说许多。事实上，我所说的这一切，既不足以符合主题的严肃性，也不足以符合你的身份地位。但我们的\Person{奥索尼乌斯}已开始不耐烦地要纸，催促秘书们，并因他的马嘶鸣而不耐烦，指责我可怜的智慧迟钝。那么，请记住我，并在基督内昌盛。还有一事：效法家中圣\Person{维辣}所立的榜样，她如同基督的真正追随者，不惧承受朝圣的艰辛。在这崇高的事业中，从一个女子身上找到你的\Quote{在这崇高事业中的领袖}。\par

\SourceRecord{Translated by W.H. Fremantle, G. Lewis and W.G. Martley.；Nicene and Post-Nicene Fathers, Second Series；Vol. 6.；Edited by Philip Schaff and Henry Wace.；Buffalo, NY: Christian Literature Publishing Co.,；1893.；<http://www.newadvent.org/fathers/3001118.htm>.}

% structure: p0001 source=h1 target=section reason=page-title

\section{第一百一十九书}

\Emph{致\Person{米内尔维乌斯}与\Person{亚历山大}}\par

\BibleRef{格前 15:51}\par

\SourceRecord{Translated by W.H. Fremantle, G. Lewis and W.G. Martley.；Nicene and Post-Nicene Fathers, Second Series；Vol. 6.；Edited by Philip Schaff and Henry Wace.；Buffalo, NY: Christian Literature Publishing Co.,；1893.；<http://www.newadvent.org/fathers/3001119.htm>.}

% structure: p0001 source=h1 target=section reason=page-title

\section{\WorkTitle{第一二〇封书信}}

\Emph{致\Person{赫迪比亚}。}\par

应高卢一位热心研究\WorkTitle{圣经}的女士\Person{赫迪比亚}的请求，\Person{热罗尼莫}处理以下十二个问题。值得注意的是，其中几个问题属于我们今日的历史批判学。\par

(1) 人如何能成为完美的？以及没有子女的寡妇应如何为天主而生活？\par

(2) \BibleRef{玛 26:29} 的含义是什么？\par

(3) 福音记载的差异如何解释？如何调和\BibleRef{玛 28:1}与\BibleRef{谷 16:1-2}？\par

(4) 如何调和\BibleRef{玛 28:9}（星期六傍晚）与\BibleRef{若 20:1-18}（星期天早晨）？\par

(5) 如何调和\BibleRef{玛 28:9}与\BibleRef{若 20:17}？\par

(6) 若墓前有士兵看守，\Person{伯多禄}和\Person{若望}怎么能自由进入？\BibleRef{玛 27:66}\par

(7) 如何调和\Person{玛窦}和\Person{马尔谷}记载的宗徒被吩咐往\Place{加里肋亚}见耶稣，与\Person{路加}和\Person{若望}记载的主在\Place{耶路撒冷}显现给他们这二者？\par

(8) \BibleRef{玛 27:50-51} 的含义是什么？\par

(9) 如何调和\BibleRef{若 20:22}所记耶稣向宗徒嘘气赐予圣神，与\BibleRef{路 24:49}所记主将在升天后将圣神赐下二者？\par

(10) \BibleRef{罗 9:14-29} 这段经文含义是什么？\par

(11) \BibleRef{格后 2:16} 的含义是什么？\par

(12) \BibleRef{得前 5:23} 的含义是什么？\par

这封书信写于公元406年或407年。\par

\SourceRecord{Translated by W.H. Fremantle, G. Lewis and W.G. Martley.；Nicene and Post-Nicene Fathers, Second Series；Vol. 6.；Edited by Philip Schaff and Henry Wace.；Buffalo, NY: Christian Literature Publishing Co.,；1893.；<http://www.newadvent.org/fathers/3001120.htm>.}

% structure: p0001 source=h1 target=section reason=page-title

\section{第一二一号书信}

\Emph{致\Person{阿尔加西亚}}\par

\Person{热罗尼莫}写信给一位\Place{高卢}的女士，名叫\Person{阿尔加西亚}，以回答她向他提出的十一个问题。如下：——\par

(1) 如何将\BibleRef{路 7:18-19}与\BibleRef{若 1:36}调和？\par

(2) \BibleRef{玛 12:20}的含义是什么？\par

(3) \BibleRef{玛 16:24}的含义是什么？\par

(4) \BibleRef{玛 24:19-20}的含义是什么？\par

(5) \BibleRef{路 9:53}的含义是什么？\par

(6) 不义管家的比喻有何含义？\par

(7) \BibleRef{罗 5:7}的含义是什么？\par

(8) \BibleRef{罗 7:8}的含义是什么？\par

(9) \BibleRef{罗 9:3}的含义是什么？\par

(10) \BibleRef{哥 2:18}的含义是什么？\par

(11) \BibleRef{得后 2:3}的含义是什么？\par

这封信的日期是公元406年。\par

\SourceRecord{Translated by W.H. Fremantle, G. Lewis and W.G. Martley.；Nicene and Post-Nicene Fathers, Second Series；Vol. 6.；Edited by Philip Schaff and Henry Wace.；Buffalo, NY: Christian Literature Publishing Co.,；1893.；<http://www.newadvent.org/fathers/3001121.htm>.}

% structure: p0001 source=h1 target=section reason=page-title

\section{第一二二封信}

第一二二封信。致\Person{鲁斯蒂古斯}。\par

\EditorialNote{鲁斯蒂古斯及其妻阿尔泰米亚曾许下贞洁誓愿，却违背了誓愿。阿尔泰米亚前往巴勒斯坦为其罪过作补赎，鲁斯蒂古斯许诺随她前往。然而他未能履行诺言，因而请求热罗尼莫写这封信，希望促使他兑现承诺。日期约为公元408年。}\par

1. 我受基督的圣仆赫底比娅\Person{Hedibia}和我在信德中的女儿阿尔特弥亚\Person{Artemia}的恳求而写信给你，一个陌生人对另一个陌生人；她曾经是你的妻子，如今不再是你的妻子，而是你的姊妹和同仆。她不仅确保了自己的救恩，也寻求了你的救恩，先前在家中，如今在圣地。她渴望效法宗徒安德肋\Person{Andrew}和斐理伯\Person{Philip}的体贴；他们在被基督找到之后，转而渴望寻找：一位寻找他的兄弟西满\Person{Simon}，另一位寻找他的朋友纳塔乃耳\Person{Nathanael}。\BibleRef{若 1:41} 对前一人，衪说：\BibleQuote{你是若纳\Person{Jona}的儿子西满，你要叫刻法\Person{Cephas}，翻出来就是石头。} \BibleRef{若 1:42} 而对后者，他的名字纳塔乃耳意为天主的恩赐，基督因衪的见证安慰了他：\BibleQuote{看，这确是一个以色列人，在他内毫无诡诈。} \BibleRef{若 1:47} 同样，古时罗特\Person{Lot} 渴望救出他的妻子和两个女儿，\BibleRef{创 19:15} 他不愿意离开烈焰中的索多玛\Place{Sodom}和哈摩辣\Place{Gomorrha}，直到他自己几乎着火，试图领出那个被她过去罪孽捆绑的人。但她在绝望中失去镇定，回头看，成为不信灵魂的纪念碑。\BibleRef{智 10:7} 然而，似乎是为了弥补一个女人的丧失，罗特炽热的信德释放了整个左哈尔\Place{Zoar}城。事实上，当他离开索多玛所在的黑暗山谷来到山上时，他一进入左哈尔，也就是小城，太阳就照在他身上；之所以这样称呼，是因为罗特所拥有的微弱信德，虽然不能拯救更大的地方，但至少能保全较小的。因为一个如此迷途，以至住在哈摩辣的人，不能一下子到达那正午之境地，在那里，天主的朋友亚巴郎\Person{Abraham}\BibleRef{雅 2:23} 款待了天主和祂的天使。\BibleRef{创 18:1}（因为若瑟\Person{Joseph}是在埃及\Place{Egypt}供养他的兄弟；当新娘对新郎说话时，她呼喊：\BibleQuote{请告诉我，你在哪里牧羊，正午你在哪里使羊群休息？} \BibleRef{歌 1:7}）善人总是为别人的罪悲伤。古时撒慕尔\Person{Samuel}为撒乌耳\Person{Saul}哀哭，\BibleRef{撒上 15:35} 因为他忽略用补赎的香膏治疗骄傲的溃疡。保禄\Person{Paul}也为格林多人\BibleRef{格后 2:4} 流泪，因为他们拒绝用眼泪洗去淫乱的污点。为此，厄则克耳\Person{Ezekiel}吞下了那书卷，内外写着歌、哀叹与灾祸；歌是赞美义人，哀叹是为悔改者，灾祸是为那些经上记载的人：\Quote{当恶人陷入邪恶的深处时，他便充满轻视。} 以此，依撒意亚\Person{Isaiah}暗示说：\BibleQuote{那天，万军的上主天主呼吁哭泣、哀悼、剃头和披上麻布；看，却是欢喜和快乐，杀牛宰羊，吃肉喝酒，说：\InnerQuote{我们吃喝吧，因为明天我们就要死了。}} \BibleRef{依 22:12} 然而，对这些恶人，厄则克耳被命令这样说：\BibleQuote{人子，你要对以色列家说话：你们这样说：我们的过犯和罪恶在我们身上，我们因此消逝，我们怎能存活？你要对他们说：我指着我的永生说，上主天主说，我决不喜悦恶人之死，却喜悦恶人离开他的路而生活，} 又说：\BibleQuote{转回，转回离开你们的恶道；以色列家啊，你们为什么要死呢？}\BibleRef{则 33:10} 没有什么比人因对更好之事绝望而执着于更坏之事更使天主发怒；事实上，这种绝望本身就是不信的标记。一个对救恩绝望的人，不可能期待将来的审判。因为他若害怕审判，便会通过善工预备迎见他的审判者。让我们听天主通过耶肋米亚\Person{Jeremiah}所说的话：\Quote{你要把脚从粗糙的路上收回，把喉咙从干渴中收回。} 又说：\BibleQuote{他们跌倒，不再起来吗？他转离，不再回头吗？}\BibleRef{耶 8:4} 也让我们听天主通过依撒意亚所说的话：\Quote{当你回转并为自己哀悼时，你将得救，那时你将知道你从前在哪里。} 我们直到恢复健康才认识到疾病的痛苦。同样，罪成了美德蒙福的反衬；光在黑暗中照耀更显明亮。厄则克耳使用与其他先知类似的语言，因为他被相似的精神所感动。他呼喊说：\BibleQuote{悔改，转离你们的一切过犯；这样，罪孽就不会使你们毁灭。抛弃你们所犯的一切过犯，为自己造一个新的心和一个新的精神；以色列家啊，你们为什么要死呢？因为死者的死，我不喜欢，上主说。} \BibleRef{则 18:30} 因此，在随后的段落中，他说：\BibleQuote{我指着我的永生说，上主天主说，我决不喜悦恶人之死，却喜悦恶人离开他的道而生活。} \BibleRef{则 33:11} 这些话告诉我们，心灵不可因不信应许的福佑而陷入绝望；那曾注定丧亡的灵魂，不可因自己的创伤无法医治而拒绝应用疗法。厄则克耳描述天主起誓，如果我们不肯相信祂关于我们救恩的应许，至少可以相信祂的誓言。义人满怀信心地祈祷说：\Quote{我们救恩的天主，求你回转我们，向你止息你的忿怒。} 又说：\Quote{上主，因你的恩宠，你使我的山坚立；你掩面，我就惊慌。} 他的意思是说：\Quote{当我离弃我过犯的污秽，趋向美德的美善时，天主用祂的恩宠坚固了我的软弱。} 看哪，我听见祂的应许：\Quote{我要追赶我的仇敌，追上他们；不将他们消灭，我决不回转。} 因此，我从前是你的敌人，逃避你的人，将被你的手捉住。不要停止追赶我，直到我的邪恶被消灭，我回到我原来的丈夫那里，他会给我羊毛、细麻、油和细面，用最丰盛的食物喂养我。\BibleRef{欧 2:7} 正是祂围堵并封闭了我邪恶的道路，\BibleRef{欧 2:6} 好使我找到祂这真道，祂在福音中说：\BibleQuote{我就是道路、真理、生命。} \BibleRef{若 14:6} 听从先知的话：\Quote{流泪播种的，必欢呼收割。那带着种子流着泪出去的人，必会带着禾捆欢乐地回来。} 也与他一同说：\Quote{我整夜使我的床榻漂浮，用眼泪浇透我的床榻。} 又说：\Quote{我的心渴慕天主，如鹿渴慕溪水。我的灵魂渴念天主，生活的天主；我何时能来朝见天主？我的眼泪日夜成为我的食物。} 在另一处：\Quote{天主，你是我的天主，我急切寻求你；我的灵魂渴慕你，我的肉身切望你，在这干旱、疲乏、无水之地。因此我在圣所中仰望了你。} 因为虽然我的灵魂曾切望你，但我更通过我肉身的劳苦寻求你，却没能在你圣所中仰望你；至少在我先住在一个没有罪恶的干旱之地之前，那里疲惫的旅客不再受敌人攻击，也没有情欲的水池或河流。\par

救主也曾为耶路撒冷城哭泣，因为其中的居民没有悔改；\BibleRef{路 19:41} 伯多禄也以痛苦的眼泪洗去了他的三次否认，\BibleRef{路 22:62} 从而应验了先知的话：\Quote{我的眼睛流下如水般的眼泪。} 耶肋米亚也为他那不悔改的人民哀叹说：\BibleQuote{但愿我的头是水，我的眼是泪的泉源，好让我日夜为我的人民哭泣！} \BibleRef{耶 9:1} 接着他又解释他哀叹的原因：\BibleQuote{你们不要为死者哭泣，}他写道，\BibleQuote{也不要哀悼他：但要为那离去的人痛哭，因为他不再回来。} \BibleRef{耶 22:10} 因此，犹太人和外邦人都不值得哀悼，因为他们从未在教会内，并且已经一次而永远地死了（关于这些人，救主说：\BibleQuote{让死人埋葬他们的死人} \BibleRef{玛 8:22}）；要哭泣的是那些因自己的罪恶和过犯而离开教会的人，他们因自己的过错而受罚，不能再回到教会。正是这个意义上，先知对教会的牧者们说话，称他们为教会的墙和塔，并对每个人说：\BibleQuote{墙啊，让眼泪流下。} \BibleRef{哀 2:18} 这样，先知性地暗示，你们将完成宗徒的训诫：\BibleQuote{与喜乐的人一同喜乐，与哭泣的人一同哭泣。} \BibleRef{罗 12:15} 并且通过你们的眼泪，你们将融化罪人刚硬的心，直到他们也哭泣；然而，如果他们坚持作恶，他们将发现这些话适用于他们：\BibleQuote{我栽种你作上等的葡萄树，完全是纯正的种子；你怎么向我变成了外邦葡萄树的坏枝子呢？} \BibleRef{耶 2:21} 又说：\BibleQuote{对木头说：你是我的父亲；对石头说：你生了我；因为他们向我转背，而不是面向。} \BibleRef{耶 2:27} 他指的是，他们不愿悔改转向天主；反而在他们硬着心肠时背对着祂，羞辱祂。因此，主也对耶肋米亚说：\BibleQuote{你看见背道的以色列做了什么事？她上到各高山，在各青翠树下，在那里行淫。我说，在她行淫之后，} 她行了淫并\BibleQuote{做了这一切的事，你回转归向我吧。但她却不回转。} \BibleRef{耶 3:6}\par

2. 我们的心多么刚硬，而天主多么慈悲！祂甚至在我们犯了许多罪之后，仍敦促我们寻求救恩。然而即便如此，我们也不愿转向更好的事物。请听主的话：\Quote{若有人休了他的妻子，她离开他，成了另一个人的妻子，以后又想回到他那里，他还会接纳她吗？他岂不更厌恶她？但你们与许多情人行淫，仍要回到我这里来，这是主说的。} 最后一句，希伯来原文（希腊文和拉丁文译本未保存）是：\Quote{你离弃了我，但回来吧，我必接纳你，这是主说的。} 依撒意亚也以同样的意义说了几乎相同的话：\Quote{回来吧，以色列的子女们，你们这些深思熟虑而行恶的人。}他喊道，\BibleQuote{你归向我，我必救赎你。我是天主，除我之外没有别的神；公义的天主和救主，除我之外没有别的。大地四极的人，你们要仰望我，并得救。} \BibleRef{依 45:21} \Quote{你们要记念这事，显出你们是男子；你们这些背道的人，要把这事放在心上。回心转意，记念上古的事：因为我是天主，没有别的神。} 岳厄尔也写道：\BibleQuote{你们要全心归向我，禁食、哭泣、哀悼：要撕裂你们的心，而不是衣服，归向主你们的天主；因为他是宽仁的、慈悲的……并且对所降的灾祸回心转意。} \BibleRef{岳 2:12} 祂的慈悲多么伟大，多么过分——如果我可以这么说——而且祂的怜悯不可言喻，先知欧瑟亚告诉我们，当祂以主的名说话时：\BibleQuote{我怎能舍弃你，厄弗辣因？我怎能交出你，以色列？我怎能叫你如同阿德玛？我怎能待你如同责波殷？我的心在我内翻转，我的怜悯一同燃起。我不再施行我的烈怒。} \BibleRef{欧 11:8} 达味也在圣咏中说：\Quote{在死亡中无人记念你；在阴府里谁感谢你？} 在另一处：\Quote{我向你承认了我的罪，没有隐藏我的过犯。我说：我要向主承认我的过犯；你就赦免了我罪的过犯。为此，凡虔敬的人，在可寻获的时候，都向你祈祷；在大水泛滥时，他们必不能临近他。}\par

3. 想想那值得被比作洪水的哭泣有多么伟大。凡如此哭泣，并和先知耶肋米亚一起说\BibleQuote{愿我的眼珠不停息} \BibleRef{哀 2:18} 的人，将立刻发现这话应验在他身上：\Quote{慈爱和真理彼此相遇：正义与和平彼此相吻；} 这样，如果正义和真理惊吓他，慈爱和平安就能鼓励他寻求救恩。\par

罪人的全部悔改向我们展示在《圣咏》第五十一篇中，这篇圣咏是\Person{达味}在亲近了赫特人\Person{乌黎雅}的妻子\Person{巴特舍巴}之后，当他对先知\Person{纳堂}的斥责回答说\BibleQuote{我犯了罪}时写下的。他一承认自己的过失，就因这句话得到安慰：\BibleQuote{上主也赦免了你的罪。}\newline{}\BibleRef{撒下 12:13}他在奸淫上又加了谋杀；然而他泪流满面地说：\BibleQuote{天主，求你按你的仁慈怜悯我；按你丰厚的慈悲，涂抹我的过犯。}\newline{}一个如此重大的罪需要找到极大的慈悲。于是他继续说：\BibleQuote{求你将我的过犯洗尽，把我的罪恶除净，因为我认清了我的过犯；我的罪恶常在我的眼前。我得罪了你，惟独得罪了你}\newline{}——作为君王，他除了天主谁都不怕——\BibleQuote{在你眼前行了这恶；为叫你在说话时显正义，在审判时显公正。}\newline{}因为\BibleQuote{天主把众人都禁锢在悖逆之中，是为了要怜悯众人。}\newline{}\BibleRef{罗 11:32}达味取得了如此进步，以致这个曾经是罪人和忏悔者的人后来成为导师，能够说：\BibleQuote{我要教诲悖逆的人你的道路，罪人必将归向你。}\newline{}因为\Quote{在上主面前有认罪与美丽}，所以一个认罪并说\BibleQuote{因我的愚昧，我的伤口发臭流脓}的罪人，便失去其恶臭的伤口，得以痊愈而洁净。但\BibleQuote{掩饰自己罪过的，必不会顺利。}\newline{}\BibleRef{箴 28:13}\par

不虔的君王\Person{阿哈布}为夺取\Person{纳波特}的葡萄园流了他的血，他与比他床伴更残忍的同伴\Person{依则贝尔}一同受到\Person{厄里亚}的严厉斥责。\BibleQuote{上主这样说：你杀了人，还要占领吗？}\newline{}以及\BibleQuote{在狗舔了纳波特血的地方，狗也要舔你的血。}\newline{}和\BibleQuote{狗要在依次勒耳城墙旁吃掉依则贝尔。}\newline{}\BibleQuote{当阿哈布听见这些话}\newline{}——经文继续——\BibleQuote{就撕裂了自己的衣服，身穿麻衣，禁食，身穿麻衣睡觉……上主的话传给厄里亚说：因为阿哈布在我面前自卑，所以在他有生之日我不降此灾祸。}\newline{}\BibleRef{列上 21:27}阿哈布的罪和依则贝尔的罪相同；然而因为阿哈布悔改了，他的惩罚被推迟到他的儿子身上，而依则贝尔坚持她的邪恶，当场就遭遇了厄运。\par

此外，主在福音中告诉我们：\BibleQuote{尼尼微人在审判时，将同这一代人起来，并要定他们的罪，因为他们因\Person{约纳}的宣讲而悔改了；}\newline{}\BibleRef{玛 12:41}祂又说：\BibleQuote{我不是来召义人，而是来召罪人悔改。}\newline{}\BibleRef{玛 9:13}那丢失的银币被寻找，直到在泥中找到。\BibleRef{路 15:8}同样，那九十九只羊被留在旷野，而牧人将那只迷失的羊扛在肩上带回家。\BibleRef{路 15:4}因此，\BibleQuote{对于一个罪人悔改，在天主的天使前也是这样欢乐。}\newline{}\BibleRef{路 15:10}想到天上的生灵因我们的得救而欢喜，这是何等有福的念头！因为正是关于我们说了这话：\BibleQuote{你们悔改吧！因为天国临近了。}\newline{}\BibleRef{玛 3:2}死亡与生命彼此对立；没有中间状态。然而悔改能将死亡与生命连接起来。那个荡子，据记载，耗尽了他所有的财产，在遥远的地方离开父亲，\BibleQuote{恨不得拿猪吃的豆荚来果腹。}\newline{}但当他回到父亲那里时，肥牛犊被宰杀，给他穿上袍子，戴上戒指。\BibleRef{路 15:11}这就是说，他重新得到了他先前玷污了的基督的袍子，并听到令他安慰的训诫：\BibleQuote{你的衣服应当时常洁白。}\newline{}\BibleRef{训 9:8}他接受了天主的印玺，并向主呼喊：\Quote{父亲，我得罪了天，也得罪了你；}\newline{}并且接受了和好的亲吻，他对祂说：\Quote{上主，如今你脸上之光已经印在我们身上了。}\par

请听\Person{厄则克耳}的话：\BibleQuote{至于恶人的邪恶，在他转离他的邪恶之日，他必不因此而跌倒；义人在他犯罪之日，也不能因他的义而存活。}\newline{}\BibleRef{则 33:12}上主按祂找到每个人的情形审判各人。祂不看重过去，只看重现在。过去的罪可能存在，但更新和皈依能除去它们。我们读到：\BibleQuote{义人七次跌倒，仍能起来。}\newline{}\BibleRef{箴 24:16}如果他跌倒，怎么还是义人？如果他是义人，又怎么跌倒？答案是：一个罪人如果总是悔改并重新起来，就不会失去义人的称号。如果一个罪人悔改，他的罪不仅被赦免七次，而是七十个七次。\BibleRef{玛 18:21}那多得赦免的，爱得也多。\BibleRef{路 7:47}那个妓女用眼泪洗了救主的脚，用头发擦干；且对她说——她是教会从万民中聚集的预像——\BibleQuote{你的罪赦了。}\newline{}\BibleRef{路 7:48}那自以为义的\Person{法利塞人}因骄傲而丧亡，而谦卑的\Person{税吏}却因他的认罪而得救。\BibleRef{路 18:14}\par

天主藉先知耶肋米亚的口起誓说：\Quote{我一发言，论及一个民族和一个王国，要拔除、拆毁和毁灭它；如果我对它发言的那个民族，离弃自己的邪恶，我必后悔我原想加给它们的灾祸。我一发言，论及一个民族和一个王国，要建立和栽培它；如果它在我眼中行恶，不听从我的声音，我必后悔我原说要恩待它的那一切。}\newline{}紧接着祂又说：\Quote{看，我正向你们策划灾祸，设计谋略攻击你们；你们各人当离开自己的恶道，改善你们的行为和作为。他们却说：绝望了！我们仍要随从自己的计谋，各人照自己邪恶的心的想像而行。}\BibleRef{耶 18:7}义人西默盎在福音中说：\Quote{看，这孩子已被立定，使许多人跌倒和复起，}\BibleRef{路 2:34}跌倒，即罪人的跌倒；复起，即悔改者的复起。同样，宗徒致格林多人书说：\Quote{听说你们中间有淫乱的事，且是这样的淫乱，连在外邦人中也没有，竟有人占有自己父亲的妻子。你们还自高自大，而不哀痛，好叫那行这事的人从你们中间被除掉。}\BibleRef{格前 5:1}在致同一人的第二封书信中，他说：\Quote{免得这样的人被过分的忧愁所吞没，}\BibleRef{格后 2:7}他召他回来，并恳请他们向他确认他们的爱，好使那因乱伦而灭亡的人能因悔改得救。\par

\Quote{没有人是清洁无罪的，即使他只活了一天。}人的年岁为数众多。\Quote{星辰在他眼中也不纯洁，\BibleRef{约 25:5}且他指摘天使的不是。}若在天上有罪，地上岂不更有罪？若那些没有肉体诱惑的都被罪玷污，我们这些被脆弱的肉体包裹、被迫与宗徒同喊\Quote{我这个人真不幸呀！谁能救我脱离这该死的肉身呢？\BibleRef{罗 7:24}因为我的肉身内没有善住着}的人，岂不更是如此？\BibleRef{罗 7:18}因为我们所做的，不是我们所愿意的；我们所不愿意的，我们却去做；灵魂想做一事，肉体却被迫做另一事。若有人在圣经中被称作义人，且不但是义人，还是在天主眼中的义人，他们是按我引用的经文所说的义被称为义的：\Quote{义人七次跌倒，仍必兴起，}\BibleRef{箴 24:16}并基于所规定的原则：恶人的邪恶在他转向悔改之日，必不使他受害。\BibleRef{则 33:12}事实上，若翰的父亲匝加利亚被描述为义人，却因不信所传的信息而犯罪，并立即被罚喑哑。\BibleRef{路 1:20}甚至约伯，在他的历史开头被说成是完美、正直、无怨言的，后来却由天主的言语和他自己的认罪证明为罪人。如果亚巴郎、依撒格、雅各伯，先知们和宗徒们也绝非无罪，如果最好的麦子也掺杂着糠秕，对我们这些经上所记\Quote{糠秕怎能与麦子相比？上主说}的人，又能说什么呢？\BibleRef{耶 23:28}然而糠秕被保留用于未来的焚烧；同样，稗子如今也与生长的麦子混在一起。因为那一位要来，祂的手里拿着簸箕，要扬净祂的禾场，把麦子收进仓里，把糠秕用不灭的火烧掉。\BibleRef{玛 3:12}\par

4. 我如此漫游在圣经最美丽的田野中，采集了其中最可爱的花朵，为你编织一个悔改的花环；因为我的目的是，使你借着鸽子的翅膀飞翔，得享安息，并与仁慈的圣父和好。你从前的妻子，如今是你的姊妹和同仆，告诉我：根据宗徒的训示，\BibleRef{格前 7:5}你和她同意分居，以便专心祈祷；但过了一段时间，你的脚像踏在水上一样下沉，而且——坦白说——完全失足。至于她，她听见主向她说话，如同向梅瑟说：\Quote{至于你，你站在我身边}；\BibleRef{申 5:31}并同圣咏作者论及主说：\Quote{祂把我的脚放在磐石上。}但你的房屋——她继续说——因没有信德的稳固根基，在魔鬼的旋风前倒塌了。\BibleRef{玛 7:24}然而她的房屋仍然在主内屹立，并愿向你提供庇护；你仍可在精神上与她结合，你曾与她身体结合。因为正如宗徒说：\Quote{那与主结合的，便是与祂成为一神。}\BibleRef{格前 6:17}此外，当蛮族的暴怒和被掳的危险再次将你们分开时，你曾庄严起誓：如果她能前往圣地，你必立即或随后跟随她，并设法拯救你的灵魂——因你的疏忽而似乎已经失去的灵魂。现在，实践你当时在天主面前所许的愿吧。人的生命是不确定的。因此，免得你在履行诺言之前被夺去，效法她吧，你本应是她的导师。可耻！较弱的器皿战胜了世界，而较强的却被它战胜！\par

一个妇人引领着高超的勇力；\par

然而，当她的救恩引领你到达信仰的门槛时，你却不跟随她！但或许，你想要保全你的财产残余，并见你的朋友、同胞以及他们的城市和别墅最后一面。若是如此，在囚禁的恐怖中，在得意洋洋的敌人面前，在省城的沉船之际，至少紧紧抓住补赎的木板；并纪念你的同仆，他每日为你的救恩叹息，从不对此绝望。当你仍在自己的国家漂泊时（尽管，确实，你已不再拥有国家；你从前拥有的，如今已失去），她在那些目睹了我们的主和救主的诞生、被钉十字架与复活的可敬地点为你转求，而祂在其中第一个地点发出了祂婴孩的啼哭。她以祈祷吸引你归向她，好使你获得救恩，若非藉你自身的努力，至少也藉她的信德。从前有一个人瘫痪在床，全身关节都无力，既不能动脚行走，也不能动手祈祷；然而，当他被他人抬到主面前时，主彻底恢复了他的健康，使他能扛起那刚才还抬着他的床榻。\BibleRef{玛 9:1}你也是——身体缺席，但在她的信德中临在——你的同仆将你献给她的主和救主；并且她如同客纳罕妇人那样论你说道：\BibleQuote{我的女儿被魔鬼折磨得甚苦。}\BibleRef{玛 15:22}灵魂并无性别；因此我可以恰当地称你的灵魂为她的女儿。因为正如母亲哄诱她那尚不能吃固体食物的未断奶的孩子；同样，她也召唤你到适合婴孩的奶那里，并给你如同哺乳母亲所提供的养料。这样，你便能与先知一同说：\Quote{我如亡羊走迷了路；求你寻找你的仆人，因为我总不忘记你的诫命。}\par

\SourceRecord{Translated by W.H. Fremantle, G. Lewis and W.G. Martley.；Nicene and Post-Nicene Fathers, Second Series；Vol. 6.；Edited by Philip Schaff and Henry Wace.；Buffalo, NY: Christian Literature Publishing Co.,；1893.；<http://www.newadvent.org/fathers/3001122.htm>.}

% structure: p0001 source=h1 target=section reason=page-title

\section{Letter 123}

Letter CXXIII. To \Person{Ageruchia}.\par

\EditorialNote{呼吁高卢贵妇阿格鲁基亚不要再次结婚。应与写给富里亚（第54封）和萨尔维娜（第79封）的信比较。提及斯提利科与阿拉里克的条约可确定年代为公元409年。}\par

1. 我必须在旧路上寻觅新径，为这陈腐而磨损的题材构思出自然、相同又不相同的处理方式。的确，路只有一条；但常常可以抄近路达到目标。我曾多次写信给寡妇，为了教导她们，我从圣经中寻找榜样，将其五彩缤纷的花朵编成一个贞洁的花环。如今我致信给阿格鲁基亚；她的名字（由天主上智所赐）已成了她日后生活的预言。她身边有祖母、母亲和姑母；一群高贵、经受过考验的基督徒妇女。她的祖母梅特罗尼亚，如今守寡四十年，令我们想起福音中法奴耳的女儿亚纳。见：\BibleRef{路 2:36} 她的母亲贝尼尼亚，守寡现已十四年，身边有贞女环绕，她们的贞洁结出百倍的果实。阿格鲁基亚的父亲塞莱里努斯的姐妹，从侄女幼年起就抚养她，事实上她一出生就抱在膝上。她失去了丈夫的安慰，二十年来教导弟弟的孩子，将自己从母亲那里学到的功课传授给她。\par

2. 我简单提及这些，是为了让我的年轻朋友明白，决定不再婚只是尽了对家庭的责任；她若履行此责，并不值得称赞，但若失责，则理应受责。更何况她有一个遗腹子，以其父辛普利丘斯命名，因此不能以孤独或缺少继承人为借口。因为许多人的情欲常以这样的借口掩饰自己，仿佛放纵的冲动只是对后代的渴望。但为什么我对一个犹豫不决的人说话，而我听说阿格鲁基亚寻求教会的保护，对抗她在宫廷中遇到的众多求婚者？因为魔鬼煽动人互相竞争，来考验我们亲爱的寡妇的贞洁；地位、美貌、青春和财富使她成为众人追求的目标。但是，对她的贞节攻击越猛烈，胜利后的奖赏也越大。\par

3. 但我刚出港，就发现海路被礁石阻断。我面对的是宗徒保禄的权威，他在致弟茂德书中论及寡妇时这样说：\BibleQuote{所以我愿意年轻的寡妇嫁人，生养儿女，治理家务，不给敌人以辱骂的机会。因为有些人已经转去随从撒殚。} 见：\BibleRef{弟前 5:14} 因此我必须首先考虑这一声明的含义，并考察整段经文的前后文。然后我才能将脚置于宗徒的足迹，如俗语所说，不偏离一根头发之左右。他先前描述了他理想的寡妇：只做一个丈夫的妻子，养育儿女，有行善的美名，用自己的财物周济困苦的人。见：\BibleRef{弟前 5:9} 寄望于天主，昼夜祈祷。见：\BibleRef{弟前 5:5} 他将其对立者与之对比，说：\BibleQuote{但那宴乐度日的寡妇，虽活着也是死的。} 为了以一切必要的劝诫警告他的弟子弟茂德，他立即加上这些话：\BibleQuote{至于年轻的寡妇，你要拒绝；因为她们在基督内放纵情欲之后，就想嫁人，因而受罚，因为她们废弃了起初的信德。} 见：\BibleRef{弟前 5:11} 正是为了那些曾对基督——她们的净配——犯奸淫而出轨的人（因为这是希腊文 \OriginalTerm{放纵情欲}{\GreekText{καταστρηνιάσωσι}} 的意思）——正是为了这些人，宗徒希望她们再婚，认为按妇德再婚强于犯奸淫；但这再婚是许可，而非命令。\par

4. 我们还需逐句考察这段经文。他说：\BibleQuote{我愿意年轻的寡妇嫁人。} 请问为什么？因为我不愿年轻妇女犯奸淫。\BibleQuote{生养儿女；} 见：\BibleRef{弟前 5:14} 这是什么理由？免得她们因害怕后果而杀掉在奸淫中怀上的孩子。\BibleQuote{治理家务。} 请问为什么？因为妇女再嫁比卖淫要好得多，有第二个丈夫比有几个情夫要好。前者在悲惨处境中带来解脱，后者则涉及罪恶及其惩罚。他接着说：\BibleQuote{不给敌人以辱骂的机会。} 这是一条简短而全面的诫命，其中总结了许多劝诫。例如：寡妇不可因过分注重衣饰而玷污她守寡的宣认，不可因轻佻的微笑或传情的眼神吸引一群年轻人，不可言语是一套、行为是另一套，不可给他人以借口将那句著名诗句用到她身上：\par

\Quote{她意味深长地一瞥，狡黠地微笑。}\par

最后，为将这一切再婚的理由压缩成几句话，保禄表明他命令的动机，说：\BibleQuote{因为有些人已经转去随从撒殚。} 因此，他允许不能节制的妇女再婚，必要时甚至可以第三次婚姻，只是为了从撒殚手中救出她们，宁愿一个女人嫁给最坏的丈夫，也不愿她嫁给魔鬼。他对格林多人说了类似的话：\BibleQuote{我对那未结婚的和寡妇说：如果他们能和我一样，为他们也好。但如果他们不能节制，就让他们嫁娶，因为嫁娶比欲火中烧更好。} 见：\BibleRef{格前 7:8} 为什么，宗徒，嫁娶更好？他立即回答：因为欲火中烧更糟。\par

5. 除这些考量外，那绝对善而非仅相对善的，是如同宗徒一样，即无束缚而非受捆绑，自由而非被奴役，关心天主的事而非妻子的事。紧接着他加上：\BibleQuote{妻子当丈夫活着时，因法律被束缚；但如果丈夫长眠了，她便可自由嫁给她所愿意的人，只要在主内。但她如果按我的判断，仍这样守下去，便更有福；我想我也拥有天主的圣神。}这段经文在意义上与前文相符，因为两者的精神相同。虽然书信不同，但同出一位作者。当丈夫活着时，女人被束缚；当他死了，她便被释放。婚姻是一种束缚，而寡居是它的解除。妻子被束缚于丈夫，丈夫也被束缚于妻子；结合如此紧密，以致他们对自己的身体没有主权，而是彼此亏欠。在婚姻轭下的人没有选择节制的自由。当宗徒加上\Quote{只要在主内}这句话时，他排除了外邦人的婚姻，对此他在别处这样说：\BibleQuote{你们不要与不信者不平等的负轭：因为正义与不义有什么相交？光明与黑暗有什么相通？基督与\OriginalTerm{贝利亚尔}{Belial}有什么和谐？信者与不信者有什么相干？天主的圣殿与偶像有什么同意？}\BibleRef{格后 6:14}我们不可用牛和驴同轭耕地；\BibleRef{申 22:10}也不可将我们的婚礼外衣染成不同颜色。他随即收回所作的让步，好像后悔自己的意见，藉说\Quote{她如果这样守下去，便更有福}即不结婚，而撤回此说；并宣称按他的判断这条路更可取。为使这一点不被轻视为人言，他声称这具有圣神的权威，以致我们听见的不是同一个人向肉体的软弱让步，而是圣神藉着宗徒发言。\par

6. 再者，不可让任何年轻寡妇以他命令六十岁以下者不被登记入册\BibleRef{弟前 5:9}为借口，而平息良心的不安。他这样安排并非催促未婚女子或年轻寡妇结婚，因为论及已婚者他甚至说：\BibleQuote{时间是短促的：从今以后，有妻子的要像没有一样。}不，他指的是那些有亲属能供养的寡妇，有儿子或孙子负责她们的赡养。宗徒命令后者要在家里行孝，报答父母，并充分供应他们；好让教会不被累赘，能自由地救济那些真寡妇。他写道：\BibleQuote{尊敬那些真寡妇}，即那些孤单无靠、没有亲属帮助、不能亲手劳作、因贫穷而虚弱、被年岁压倒、倚靠天主、唯一工作就是祈祷的人。由此易于推断，年轻寡妇除非因健康不佳被豁免，否则要么自食其力，要么交由子女或亲属照顾。此处\Quote{尊敬}一词指施舍或馈赠，正如紧接的经文所说：\BibleQuote{长老们…应受加倍的尊敬，尤其是那些在言语和道理上劳碌的人。}\BibleRef{弟前 5:17}同样，在福音中，当主讨论那法律诫命——\BibleQuote{应孝敬你的父亲和你的母亲}\BibleRef{出 20:12}——时，祂声明这不应仅理解为空话（空话虽表面尊敬，却让父母的需要得不到满足），而应实际供应生活必需品。主命令儿女供养贫困的父母，在父母年老时回报他们幼时所受的恩惠。反之，经师和法利塞人教训儿女这样回答父母：\BibleQuote{这是\OriginalTerm{科尔班}{Corban}，即礼物\BibleRef{谷 7:11}，我已许给祭台并承诺献给圣殿：它在那里同样能救济你，好像我直接给你买食物一样。}于是常常发生父母穷困，而儿女却为司祭和经师献祭以供他们食用。如果宗徒迫使贫穷的寡妇——但只限于年轻且未因疾病衰弱的——亲手劳作，使教会不为他们的赡养所累，能支持年老的寡妇，那么一个今生富足、既能自给又能助人、且能用不义之财结交朋友\BibleRef{路 16:9}以便被接到永远帐幕里的人，还能提出什么借口呢？\par

也要想想，只有作过一个丈夫的妻子的女人，才能被选为寡妇。我们有时以为只有一夫一妻者才能被允许到祭台前是司铎职的独特标志。但不仅再婚者被排除在司铎职务外，他们也被禁止接受教会的施舍。一个已再婚的女子被认为不配得到信友的支持。就连平信徒也受司铎法律的约束，因为他的行为必须能够让他被选为司铎。如果他曾再婚，他就不能被选。因此，既然司铎是从平信徒中选出的，平信徒也受此诫命的约束，其实践对于达到司铎职是不可或缺的。\par

7. 我们必须区分宗徒自己所愿望的和祂不得不同意的事。如果他允许我再婚，这是由于我自己的不能自制，而非他的愿望。因为他愿意众人都像他一样，思念天主的事，一旦被释放，不再寻求束缚。但当他看到不稳定的人因不能自制而有陷入情欲深渊的危险时，他便向他们提供第二次婚姻的机会；这样，如果他们必须在泥沼中打滚，那也是与一个而非多个伴侣。第二次婚姻的丈夫不应认为这是苛刻的话，或与宗徒所定的规则相冲突。宗徒有两种心意：首先，他宣布一条命令，\Quote{因此我对未婚者和寡妇说：如果他们像我一样，保持现状，对他们是有益的。} 其次，他作出让步，\Quote{但如果他们不能自制，就让他们结婚，因为结婚比燃烧更好。} \BibleRef{格前 7:8} 他首先表明他自己所愿望的，然后是他不得不赞同的。他希望我们——在第一次婚姻之后——保持如同他一样，即未婚，并以他自身的宗徒榜样为我们树立了他所谈论的真福实例。然而，如果他发现我们不愿照他的愿望去做，他就对我们的不能自制作出让步。那么，我们为自己选择这两者中的哪一个？是他所偏好、本身是善的那一个？还是那与恶相比可容忍，但因其仅是恶的替代品而不完全善的那一个？假设我们选择了宗徒所不愿望、只是违心同意的那个方向，允许那些追求较低目标的人自行其是；在这种情况下，我们实现的是我们自己的愿望，而非宗徒的愿望。我们在旧约中读到，祭司的女儿若结过一次婚成了寡妇，可以吃祭司的食物，她们死后，葬礼与父母相同。反之，如果她们另嫁他人，则要与父亲和祭物隔离，被视为外人。\BibleRef{肋 22:12}\par

8. 这些婚姻上的限制甚至在外教人中也得到遵守；如果真信德不能为基督做到假信德为魔鬼所做的事——魔鬼以其自身的致罪的贞洁取代了福音的救恩的贞洁——那就是我们的罪过。雅典的\OriginalTerm{圣职者}{hierophant}否认自己的男性特质，并通过持久的节制削弱自己的情欲。火神的\OriginalTerm{祭司}{flamen}的圣职仅限于结过一次婚的人，而祭司妻子的随从也必须只嫁过一次。只有\OriginalTerm{一夫一妻者}{monogamists}被允许参与与埃及神牛相关的神圣仪式。我无需提及维斯塔贞女以及阿波罗、阿该亚的朱诺、狄安娜和弥涅耳瓦的贞女，她们都因履行圣召而终身童贞，日渐憔悴。我只想简要提及迦太基女王（\Person{狄多}），她宁愿自焚也不愿嫁给\Person{雅尔巴斯}国王；\Person{哈斯德鲁巴耳}的妻子，双手各抱一个孩子，跳入脚下的火焰以保全荣誉；\Person{卢克莱提亚}，在失去贞洁的珍宝后，不肯在灵魂受玷污后苟活。我不会再列举许多类似贞洁的实例而延长这封信——你可以在我的\WorkTitle{驳约维尼安}第一卷中读到，并得益。我只讲述一个发生在你本国的事例，这将向你表明，即使在野蛮、未开化、残忍的民族中，贞洁也受到高度尊崇。从前，条顿人从遥远的日耳曼海沿岸而来，蹂躏了高卢各地，直到几个罗马军团被他们击溃后，\Person{马略}终于在\Place{阿奎·塞克斯提埃}的一次遭遇战中打败了他们。根据投降条件，三百名他们的已婚妇女要交给罗马人。当条顿主妇们听到这一协议时，首先请求执政官允许她们被分别出来在\Person{刻瑞斯}和\Person{维纳斯}的神庙中服务；当请求被拒绝，并被侍卫带走时，她们杀死了自己的幼童，第二天早晨，人们发现她们互相拥抱而死，显然是在夜间自缢身亡。\par

9. 难道一位高贵的女士会做出这些野蛮妇女在作为战俘时都拒绝做的事吗？在失去第一位丈夫（无论好坏）之后，她要去尝试第二段婚姻，从而违背天主的审判吗？如果她立即又失去第二位，她会娶第三位吗？如果他也被召去安息，她会继续到第四、第五位，从而与娼妓为伍吗？不，寡妇必须万分小心，不能越出贞洁界限一步。因为一旦越界，冲破适合主妇的端庄，她很快就会放纵于各种无度；以致先知的话将应验在她身上：\Quote{你有娼妓的面容，你拒绝羞愧。} \BibleRef{耶 3:3}\par

那麼怎樣？我譴責第二次婚姻嗎？絕不；但我稱讚第一次婚姻。我將再婚者逐出教會嗎？決不；但我勸勉那些已婚一次的人過節制的生活。諾厄方舟內有不潔的動物，也有潔淨的動物。內中有爬蟲，也有人。在大戶人家中，有各式各樣的器皿，有些是貴重的，有些是卑賤的。\BibleRef{弟後 2:20} 在福音比喻中，撒在好地裡的種子結出果實，有的一百倍，有的六十倍，有的三十倍。那首先的一百倍象徵童貞的榮冠；其次的六十倍指的是寡婦的功績；而三十倍——以拇指和食指的指尖相合來表示——則指婚姻的結合。這樣，哪裡還有第二婚姻的餘地呢？沒有。它們不被計算在內。這等雜草不生長在好地裡，卻生長在荊棘和蒺藜中，那些狐狸最喜歡出沒的地方，主曾將這些狐狸比作惡王黑落德。\BibleRef{路 13:32} 一個多次結婚的女人自以為值得讚美，因為她不如妓女那麼壞，因為她與那些縱慾的犧牲品相比，只委身於一個人，而不是許多人，這使她顯得更出色。\par

10. 我即將要講述的故事難以置信，卻有許多證人擔保。多年前，我協助羅馬主教\Person{達瑪蘇}處理教會信函，為他撰寫答覆東西方議會向他提出的問題時，我看到一對夫婦，雙方都出身於社會的最底層。那男子已經埋葬了二十位妻子，那女子則有過二十二位丈夫。如今他們彼此結合，各自相信這是最後一次。男女之間都極度好奇，想看看這兩位老手誰能活到埋葬對方。丈夫勝利了，他在人群中走在他那多次結婚的妻子的棺材前面，從四面八方匯聚而來的一大群人簇擁著，他頭戴花冠，手持棕櫚枝，邊走邊向歡呼的人群撒播麥粒。對這樣一個女人，我們能說什麼呢？當然就是主對\Place{撒瑪黎雅}婦人所說的話：\Quote{你曾有過二十二個丈夫，而現在埋葬你的這個，不是你的丈夫。}\BibleRef{若 4:18}\par

11. 因此，我在基督內虔敬的女兒，我懇求你，不要沉溺於那些為放縱和不快樂者提供援助的經文，而要閱讀那些貞潔受到讚揚的經文。你失去了第一種和最高的類型，即童貞，又從第三種進入了第二種，這對你來說已足夠了；也就是說，你已履行了妻子的義務，如今作為寡婦守節度日。不要去想最低等的等級，不，是根本不算數的等級，我指第二次婚姻；也不要尋求牽強的先例來證明你再婚是正當的。你應儘可能仿效你的祖母、母親和姨母；她們的生活教導和勸告將為你樹立德行的準則。因為，如果許多婦女在丈夫在世時，就領悟了宗徒的話的真理：\Quote{凡事我都可行，但凡事不都有益處，}\BibleRef{格前 6:12} 並且為天國而自閹，\BibleRef{瑪 19:12} 無論是藉著聖洗池重生後表示同意，還是在婚後立即因信德的熱忱而如此行；那麼，一個因天主的旨意失去丈夫的寡婦，為何不該與約伯一同一次次歡呼：\Quote{主賜予，主收回，}\BibleRef{約 1:21} 然後抓住這個讓她有能力掌管自己身體的機會，而不是再次成為男人的奴僕呢？的確，克制不去享受擁有的東西，比後悔失去的東西要難得多。童貞更容易，因為童貞者對肉身的衝動一無所知；寡居更難，因為寡婦難免回想起過去曾享用的放縱。而如果她們認為丈夫是失落了，而不是先她們而去，那就更難了；因為前者帶來痛苦，後者卻帶來喜樂。\par

12. 第一个人的受造应教训我们拒绝多妻。只有一男亚当和一女厄娃；事实上，女人是由亚当的肋骨形成的。\BibleRef{创 2:21}如此分开后，他们随后在婚姻中结合在一起；按圣经的话：\BibleQuote{两人成为一体}，不是两个或三个。\BibleQuote{因此人要离开父亲和母亲，依附自己的妻子。}当然没有说\BibleQuote{依附自己的妻子们}。保禄在解释这段经文时，将其应用于基督和教会；\BibleRef{弗 5:31}使第一个亚当成为肉身的\OriginalTerm{一夫一妻者}{monogamist}，第二个成为精神的一夫一妻者。正如有一位厄娃是\BibleQuote{众生之母}，\BibleRef{创 3:20}同样有一位教会，是全体基督徒的母亲。正如被诅咒的拉默客将第一个厄娃分为两个妻子，\BibleRef{创 4:19}同样，异端者也把第二个（教会）分割成几个教会，按照\BibleBook{若望}{默示录}，这些教会应称为魔鬼的会堂，而非基督的会众。\BibleRef{默 2:9}在\BibleBook{}{雅歌}中，我们读到如下的话：\BibleQuote{有六十王后，八十妃嫔，有无数的童女。我的鸽子，我的完美人儿，只有一个；她是她母亲的独生女，是生养她者的特选者。}\BibleRef{歌 6:8}正是这位特选者，同一若望在书信中这样称呼：\BibleQuote{长老致蒙拣选的夫人和她的子女。}同样，在方舟的事上，伯多禄宗徒将方舟解释为教会的预像，\BibleRef{伯前 3:20}诺厄为他的三个儿子带进各一妻子，而不是两个。\BibleRef{创 7:13}同样，不洁的动物只取一对，一公一母，表明\OriginalTerm{再婚}{digamy}在野兽、爬虫、鳄鱼和蜥蜴中也没有立足之地。如果洁净的动物每种取七只，\BibleRef{创 7:2}即奇数，这指向等待着\OriginalTerm{童贞}{virginal chastity}的棕榈枝。因为出方舟时，诺厄向天主献了祭牲，\BibleRef{创 8:20}当然不是用成对取来的动物，因为这些是为繁衍物种的，而是用七只取来的动物，其中一些已被分别出来献祭。\par

13. 确实，圣祖们各有多妻，还有许多妾。仿佛他们的榜样还不够，达味有许多妻子，撒罗满多不胜数。犹大以为塔玛尔是妓女而与她同房；\BibleRef{创 38:12}并且按着那叫人死的字句，欧瑟亚先知不仅娶了妓女，而且娶了淫妇。\BibleRef{欧 1:2}如果这些事例要使我们自圆其说，让我们对我们所遇到的每个女人都像发情的公马；\BibleRef{耶 5:8}如同索多玛和哈摩辣人，让末日发现我们买卖、嫁娶；\BibleRef{路 17:27}让我们只在生命终结时才停止婚嫁。如果洪水前后都奉行这句格言：\BibleQuote{你们要生育繁殖，充满大地。}这对我们这些万世万代已经临到的人有什么相干呢？经上对我们说：\BibleQuote{时间短促了}，\BibleRef{格前 7:29}并且\BibleQuote{斧子已经放在树根上}；这就是说，婚姻和法律的丛林必须被福音的贞洁砍倒。有\BibleQuote{拥抱有时，戒避拥抱有时。}\BibleRef{训 3:5}由于被掳迫近，耶肋米亚被禁止娶妻。\BibleRef{耶 16:2}在巴比伦，厄则克耳说：\BibleQuote{我的妻子死了，我的口才张开。}无论那希望结婚的，或那已经结婚的，都不能在婚姻中自由地讲预言。在过去，人们欢喜听到关于他们的这句话：\BibleQuote{你的子女好似橄榄树枝围绕你的桌子}，和\BibleQuote{愿你见到你子女的子女。}但现在，对于那些度节制生活的人却说：\BibleQuote{那与主联合的，便是与主成为一神；}\BibleRef{格前 6:17}并且\BibleQuote{我的灵魂紧紧跟随你，你的右手扶持我。}那时曾说过：\BibleQuote{以眼还眼}；现在的诫命是：\BibleQuote{有人打你的右脸，连另一面也转给他。}\BibleRef{玛 5:38}在那些日子，人对战士说：\BibleQuote{大能者啊，腰间佩上你的刀！}现在却对伯多禄说：\BibleQuote{把你的刀收回原处；因为凡动刀的，必死在刀下。}\BibleRef{玛 26:52}\par

我说这话，并非要像\Person{马尔西翁}错误地那样，将法律与福音割裂。不，我在两者中承认同一天主，祂随着时期和对象的不同，既是起始也是终结；祂播种为了收割，祂种植为了有可砍伐的，祂奠基以便在时期圆满时加冕建筑。此外，若要处理将来之事的象征和预像，我们不应凭自己的意见，而应按宗徒的解释来判断。\Person{哈加尔}和\Person{撒辣}，或\Place{西乃山}和\Place{熙雍山}，是两约的预像\BibleRef{迦 4:22}。眼睛柔弱的\Person{肋阿}和\Person{雅各伯}所爱的\Person{辣黑耳}\BibleRef{创 29:17}象征会堂和教会。同样，\Person{亚纳}和\Person{培尼纳}也是如此：前者起初不育，后来在生育上超过后者。在\Person{依撒格}和\Person{黎贝加}身上，我们看到一夫一妻制的早期范例：只有\Person{黎贝加}，主在分娩时向她启示了自己，也只有她自己去求问主\BibleRef{创 25:22}。关于\Person{塔玛尔}，她生了双胞胎\Person{培勒兹}和\Person{则辣黑}，我该说什么呢？\BibleRef{创 38:27}在他们出生时，那隔断的中间墙被拆毁了，它预表两个民族之间的分裂\BibleRef{弗 2:14}；而用朱红线绑\Person{则辣黑}的手，甚至那时就以基督的血的污渍标记了犹太人的良心。我该如何谈论先知所娶的妓女呢？她预表从外邦人中聚集的教会，或者——更符合这段经文解释——预表会堂？起初由\Person{亚巴郎}和\Person{梅瑟}从偶像崇拜者中收养，如今她却否认了救主，对祂不忠。因此，她早已被剥夺了祭台、司祭和先知，必须等待多日才能回归她的第一任丈夫。因为当外邦人的数目满了以后，全以色列都将得救\BibleRef{罗 11:25}。\par

14. 我试图在有限的篇幅内压缩许多内容，就像绘图员在小型地图上描绘一个大国一样。因为我想处理其他问题，第一个我将用安娜对她姐姐狄多的言词来说明：\par

为什么在无尽的悲伤中独自浪费你的青春，\newline{}没有子女，爱的最甜蜜礼物？\newline{}你以为埋葬的死者需要这些？\par

受苦者如此简短地回答：\par

是你，我的妹妹，是你，首先\newline{}将我疯狂的心灵投进这痛苦。\newline{}我为何不能像某些无忧无虑的野物\newline{}一样过童贞的生活？\newline{}唉！我将我的灵魂许配给了死者：\newline{}我许了愿，却违背了它。\par

你在我面前摆上婚姻的喜乐。我呢，会提醒你狄多的剑、火葬堆和葬礼的火焰。在婚姻中，可期望的善比不上可能发生且必须畏惧的恶。情欲一旦放纵，总是带来悔恨；它永不满足，一旦熄灭，很快又重新燃起。它的增长或衰退是习惯的问题；它像被冲动俘虏，拒绝听从理性。但你会争辩说：\Quote{财富和财产的管理需要丈夫的监督。} 你的意思是说独身者的产业会荒废；除非你使自己像仆人一样为奴，否则就无法治理你的家吗？难道你的祖母、母亲和姑姑不是比以往享有更大的影响和尊重，被整个省份和教会领袖所景仰吗？难道士兵和旅行者没有妻子帮忙，不是照样管理家务、彼此款待吗？你为什么不能有严肃年长的仆人或自由人，比如那些从小养育你的人，来管理你的房子，回应公共事务，缴纳税款？这些人会视你为女庇护人，会爱你如乳儿，会尊敬你如圣人？\BibleQuote{你们先寻求天主的国，这一切自会加给你们。}\BibleRef{玛 6:33} 如果你关心衣服，福音吩咐你\BibleQuote{看看田野里的百合花}；如果关心食物，就去看看飞鸟，\BibleQuote{它们不播种，也不收割，你们的天父尚且养活它们。} 有多少童贞女和寡妇自己管理财产，却没有因此染上任何丑闻的污点！\par

15. 不要与年轻女子交往，也不可依恋她们，因为正是为了这样的人，宗徒才准许再婚；这样，你虽看似在平静的水面上，却可能遭遇海难。如果保禄可以对弟茂德说：\BibleQuote{至于年轻的寡妇，要拒绝，}\BibleRef{弟前 5:11}。又说：\Quote{要敬老年人如同母亲；待青年人如同姊妹，以完全的纯洁}，你还有什么借口不听我的劝诫呢？要远离一切可能沾染恶行嫌疑的人，不要满足于那句老套的回答：\Quote{我自己的良心就够了；我不在乎别人怎么说。}宗徒行事的原则并非如此。他不但行事端正在天主面前，也在众人面前行事端正，免得天主的名在外邦人中被亵渎。\BibleRef{罗 2:24}。虽然他有权柄带着一位姊妹、一位妻子同行，\BibleRef{格前 9:5}。但他没有这样做，因为他不想被不信者的良心论断。\BibleRef{格前 10:29}。而且，虽然他本可倚靠福音生活，\BibleRef{格前 9:14}。他却昼夜亲手劳作，以免成为信徒的负担。他说：\BibleQuote{若食物使我的弟兄跌倒，我就永远不吃肉。}\BibleRef{格前 8:13}。那么，让我们也说：如果一位姊妹或弟兄不是绊倒一两个人，而是绊倒整个教会，我就\Quote{不见那位姊妹或那位弟兄}。损失一部分财物，胜过危及灵魂的救恩。自由地舍弃那些迟早、无论我们愿意与否都会失去的东西，胜过失去那值得我们牺牲一切的东西。我们中谁能增加——我不说一肘，那已是极大的增加——就是一丝一毫的身量呢？我们为何忧虑吃什么、喝什么呢？让我们\Quote{不要为明天忧虑，因为明天自有明天的忧虑：一天的难处一天当就够了。}\par

雅各伯逃避哥哥时，将大量财富留在父亲家，空手前往美索不达米亚。而且，为向我们证明他的忍耐力，他拿石头当枕头。但他躺卧时，看见一个梯子立在地上，顶端直达天上，看，主站在梯子上，天使们在其上上下下。\BibleRef{创 28:11}。由此教导我们：罪人不该绝望于救恩，义人也不该因自己的德行而自以为安稳。省略大部分事迹（因为无暇详述全部细节），二十年后，那曾以杖渡过约但河的人，带着三群牲畜返回故乡，羊群牛群众多，儿女更丰盛。\BibleRef{创 32:7}。宗徒们同样周游世界，口袋里没有钱，手中没有杖，脚上没有鞋。\BibleRef{玛 10:9}。然而他们能自称：\BibleQuote{似乎一无所有，却拥有一切。}\BibleRef{格后 6:10}。他们说：\BibleQuote{金银我都没有，只把我所有的给你：因纳匝肋人耶稣基督的名，起来行走吧！}\BibleRef{宗 3:6}。因为他们没有被财富的重担压垮。因此，他们能像厄里亚一样站在岩石的裂缝里，能穿过针眼，并看见天主的背后。\par

但我们却燃烧着贪念，即使我们抨击爱财之心，仍张开衣襟接金子，永不知足。有句关于墨伽拉人的谚语，正好适用于所有受此激情折磨的人：\Quote{他们建造，如同要永远活着；他们生活，如同明天就要死去。}我们也一样，因为不相信主的话语。当我们达到人人渴望的年龄时，我们忘记了那作为人欠给自然的死亡近在眼前，并以虚假的希望许诺自己长寿。没有哪个老衰之人会以为自己活不过一年。一种对他真实状况的遗忘逐渐袭来；使他——这属土之人，且濒临解体——陷入骄傲，在想象中坐于天上。\par

16. 我这是在做什么？当我谈论货物时，船本身已沉没。那拦阻者已被移开，我们却仍不意识到敌基督近了。是的，敌基督近了，主耶稣基督将\BibleQuote{用口中的灵消灭他。}\BibleRef{得后 2:7}。他呼喊：\BibleQuote{那些日子，怀孕的和哺乳的有祸了！}\BibleRef{玛 24:19}。而这些正是婚姻的果实。\par

我现在要略言我们当前的苦况。我们中有少数人至今幸存，但这并非出于我们自己的作为，而是主的仁慈。无数野蛮部落侵占了高卢全境。阿尔卑斯山和比利牛斯山之间、莱茵河和大洋之间的整个地区，已被夸狄人、汪达尔人、萨尔马特人、阿兰人、格皮德人、赫鲁利人、萨克森人、勃艮第人、阿勒曼尼人以及——唉！为了国家！——甚至潘诺尼亚人蹂躏殆尽。因为\Quote{亚述也加入他们。}昔日荣耀的莫贡提亚库姆城已被攻陷摧毁。其教堂内成千上万人被屠杀。旺吉姆人经长期围攻后已被消灭。强盛的兰斯城、安比亚尼人、阿特雷巴特人、世界边缘的比利时人、图尔奈、施派尔和斯特拉斯堡都已落入日耳曼人之手：而阿基坦省、九族省、里昂省和纳博内省，除少数城市外，已是一片荒凉。刀剑未伤及之处，城内饥荒肆虐。我含泪提及图卢兹，它至今赖其可敬主教埃克苏佩里乌斯的功德得以保全。甚至西班牙也岌岌可危，每日因忆及辛布里人的入侵而战栗；而其他民族在实际经历一次灾祸后，他们却在持续的恐惧中遭受灾祸。\par

17. 我不提其他地方，免得我似乎对天主的仁慈绝望。我们现在从本都海到尤利安阿尔卑斯山的一切，在过去的岁月里曾一度不再属于我们。三十年来，蛮族冲破了多瑙河的屏障，在罗马帝国的中心作战。长久的习惯擦干了我们的眼泪。因为除了少数老人，所有人要么是在被俘期间，要么是在封锁期间出生的，因此他们并不怀念他们从未认识过的自由。然而日后谁会相信这一事实，或者什么样的历史著作会认真讨论：罗马不得不在自己的疆界内作战，不是为了荣耀，而是为了赤裸的生存？而且她甚至不作战，而是通过支付黄金、牺牲她所有的财富来购买生存的权利？这种屈辱并非由于她的两位最虔敬的皇帝们的过错，而是由于一个半蛮族叛徒的罪行，他用我们的金钱武装了我们的敌人来对抗我们。古时罗马帝国曾被永远打上耻辱的烙印，因为布伦努斯率领他的高卢人在洗劫了乡村并在阿利亚河击败罗马人之后，进入了罗马城本身。直到高卢（高卢人的出生地）和高卢希腊（他们在征服东西方之后定居的地方）被征服后，这一古老的污点才得以洗刷。即使是汉尼拔，像毁灭性的风暴一样从西班牙席卷到意大利，虽然他来到了城旁，却不敢围攻它。甚至皮洛士也完全被罗马名声的魔力所束缚，他摧毁了沿途的一切，却仍从城郊撤退，作为胜利者，他不敢注视他已知晓的这座众王之城。然而，作为对他们这样的侮辱——更不用说这种傲慢——的回报，这些都以对罗马的幸运结局而告终：一个被全世界放逐的人最终在比提尼亚服毒而死；另一个回到故乡，在自己的领土上被杀害。他们的国家都成了罗马人民的附属国。但现在，即使我们全副武装取得完全胜利，我们也只能从被击败的敌人那里夺回我们失去的东西，而无法得到更多。诗人卢坎在一段热情洋溢的文字中描述罗马城的力量时说：\par

如果罗马软弱，我们何处寻找力量？\par

我们可以改其词而说：\par

如果罗马沦陷，我们何处寻找援助？\par

或者借用\Person{维吉尔}的话：\par

即使我有百口铜喉，\newline{}难以诉说俘虏的苦难，\newline{}甚至历数所有被杀者的名字。\par

我所说的一切，对我这说话者以及所有听闻者都充满危险；因为我们甚至不能再自由地哀叹我们的命运，也不愿意——甚至可以说，不敢——为我们的苦难哭泣。\par

基督内最亲爱的女儿，请回答我这个问题：在如此场景中你会结婚吗？告诉我，你将嫁给什么样的丈夫？一个会逃跑的还是会战斗的？无论哪种情况，你都知道结果。代替费斯切尼歌谣，可怕的号角嘶哑声将震聋你的耳朵，你的伴娘甚至可能变成哀悼者。既然你已经失去了所有财产的收入，既然你看到你的小随从处于严密封锁之下，成为瘟疫和饥荒侵扰的猎物，你还能期望享受什么欢乐？但愿你绝不要这样低估你，也不要对这位将灵魂奉献给主的人怀有任何怀疑。虽然我的讲话名义上是针对你，但实际上是为了那些闲散、好奇、爱闲聊的其他人。这些人从一家游荡到另一家，从一个已婚妇人到另一个妇人，\BibleQuote{游手好闲}\BibleRef{弟前 5:13}，\BibleQuote{他们的神是自己的肚腹，他们的荣耀在于他们的羞耻}\BibleRef{斐 3:19}，他们对圣经一无所知，只知道那些有利于第二次婚姻的经文，但他们喜欢引用他人的例子来为自己的放纵辩护，并自我安慰说他们不比同罪的同伴更差。当你通过解释宗徒的真正含义而使这些女人的无耻提议哑口无言时；然后向你展示通过何种生活方式你可以最好地保持你的寡居，你可以阅读我所写的，这将对你大有裨益。我指的是写给欧斯托基姆的《论守童贞》以及我给富利亚和萨尔维娜的两封信。关于后两封信，你可能想知道：第一封是曾任执政官的普罗布斯的儿媳，第二封是前非洲总督吉尔多的女儿。这篇论不再婚的著作，我将以你的名字命名。\par

\SourceRecord{Translated by W.H. Fremantle, G. Lewis and W.G. Martley.；Nicene and Post-Nicene Fathers, Second Series；Vol. 6.；Edited by Philip Schaff and Henry Wace.；Buffalo, NY: Christian Literature Publishing Co.,；1893.；<http://www.newadvent.org/fathers/3001123.htm>.}

% structure: p0001 source=h1 target=section reason=page-title

\section{第一二四封书信}

\Emph{致\Person{阿维图斯}}\par

\EditorialNote{这封信所致之\Person{阿维图斯}大概就是促使\Person{热罗尼莫}写信给\Person{萨尔维纳}的同一个人（见前文\WorkTitle{第一二九封书信}，§I）。写信的背景如下：十年前（即主后399或400年），\Person{帕马丘}请求\Person{热罗尼莫}为他提供一份\Person{奥力振}\WorkTitle{论首要原理}的准确译本，以便他能发现\Person{鲁菲努斯}在其译本中引入的变动。\Person{热罗尼莫}欣然应允（见\WorkTitle{第一三三封书信}和\WorkTitle{第一三四封书信}），但当\Person{帕马丘}通读完整著作后，他认为其中教义错误极多，决定不予传播。然而，某位弟兄劝说他将手稿借出一小段时间；然后，当手稿到手后，此人匆匆抄写了一份不准确的副本，并立即发表，令\Person{帕马丘}大为懊恼。这份副本落入\Person{阿维图斯}手中后使他大为困惑，他似乎曾向\Person{热罗尼莫}请求解释。\Person{热罗尼莫}现在给予解答，同时附上他\WorkTitle{论首要原理}译本的权威版本。此信日期为主后409或410年。}\par

1. 大约十年前，圣人\Person{帕马丘}送给我一份某人对\Person{奥力振}\WorkTitle{论首要原理}的译作，或者更准确地说，误译；他请求我用拉丁文译出希腊文的真实含义，并记下作者的言论，不论好坏，不偏不倚。当我按要求做了并将书寄给他后，他读后大为震惊，将其锁在书桌里，以免流传出去伤害许多人的灵魂。然而，某位弟兄，他\Quote{对天主有热情，但不按真知识}，请求借阅手稿以便阅读；并承诺立即归还；帕马丘以为短时间不会出事，便轻信同意了。于是，那位借书阅读的人，在抄写员帮助下抄写了全书，并比承诺的更早归还。然后，他以同样的鲁莽——或者说用较轻的词语——轻率，将其偷来的东西透露给他人，使情况更糟。此外，一部关于深奥主题的巨著很难准确复制，尤其是当它必须秘密且匆忙记录时，许多地方丧失了条理和意义。因此，我亲爱的\Person{阿维图斯}，你请求我送给你一份我为\Person{帕马丘}而作而非为公众的译本副本，而上述弟兄发表的那个版本是篡改过的。\par

2. 那么，请取你所求；但要知道书中无数可憎之事，并且，正如主所说，你将不得不在蝎子和蛇中行走。它开篇说基督成为天主子，而非被生；天主圣父，因其本性不可见，甚至对圣子也不可见；圣子，作为不可见之父的肖像，与父相比不是真理，但与无法接受全能天父真理的我们相比，似乎是真理的肖像，以致我们在较小者中感知较大者的威严和伟大，在圣子中有限的父的光荣；天主圣父是不可思议的光，基督与之相比只是一丝微小的光明，尽管由于我们的无能，在我们看来他是巨大的。父与子被比作两尊雕像，一大一小；第一尊充满世界，因其大小在某种程度上不可见；第二尊可被人眼认知。写作人将全能的天主圣父称为善和完美之善；但关于圣子，他说：\Quote{他不是善，而是善的流出和肖像；不是绝对善，而只是有条件的善，如\InnerQuote{善牧}之类。}他将圣神置于父和子之后，作为第三位，在尊严和荣耀上。虽然他说他不知道圣神是受造还是非受造，但后来他给出了自己的意见：唯独天父是非受造的。\Quote{圣子}，他说，\Quote{次于父，因为他是第二位，父是第一位；而居于所有圣徒内的圣神次于圣子。同样，父的能力大于圣子和圣神。同样，圣子的能力大于圣神，因此圣神又比其它称为圣的事物有更大的美德。}\par

3. 然后，当他论及理性受造物，并描述它们因自己的疏忽而堕入尘世身体时，他接着说：\Quote{一个灵魂如此彻底地倒空自己，以致陷入罪恶，并被束缚在无理性畜生的物质身体上，这无疑是极大的疏忽和懒惰；}在后面的段落中：\Quote{这些推理使我认为，有些人是因自己的自由行为而被列入天主的圣者与仆人之列，而另一些人则是因自己的过错从圣洁中堕落到如此疏忽，以致变为对立的力量。}他主张每个终末之后都会产生新的开始，每个开始之后又有终末，而且可能发生彻底的变异；以致一个现在是人类的人，在另一个世界里可能变成魔鬼，而一个因自己的疏忽现在是魔鬼的人，将来可能被安置在更物质的身体中，从而变成人类。他如此推行这种转化过程，以至于根据他的理论，一个总领天使可能变成魔鬼，而魔鬼又可能变回总领天使。\Quote{那些摇摆或动摇但没有完全堕落的人，将受制于更好的事物——受制于宰制者和掌权者、王座和统治权，以接受统治、管理和指导；}或许这些人中会形成另一个人类种族，正如依撒意亚所说的，将有\BibleQuote{新天新地}\BibleRef{依 65:17}。但是那些不配通过人性回到先前状态的人，将变成魔鬼及其使者，即最坏的那种魔鬼；并且根据他们的行为，将在特定的世界中承担特殊的职责。\Quote{此外，在任何世界或多个世界中的魔鬼和黑暗的掌权者，若他们愿意转向更好的事物，也可以变成人类，并因此回到最初的起始。也就是说，在人类身体中承受了或长或短的惩罚和折磨的训练之后，他们可能再次达到他们堕落的那个天使顶峰。}因此，可以表明，我们人类可以变成任何其他有理性的存在，而且不止一次或紧急情况下，而是反复如此；我们和天使若忽视本分，将变成魔鬼；魔鬼若愿意接受德行，可以达到天使的品级。\par

4. 物质性的身体也将完全消逝，或在万物的终结时变得极其稀薄，如同天空、\newline{}\newline{}ewline{}\newline{}\OriginalTerm{以太}{æther}和其他精细物体。显然，这些原则必然影响作者对复活的看法。太阳、月亮以及其余星座都是活物。不仅如此，正如我们人类因自己的罪被包裹在物质而迟钝的身体中，天上的光体也因类似的原因接受了或明或暗的身体，而\OriginalTerm{魔鬼}{demon}因更严重的过错披上了星辰的外壳。他指出，这就是那位宗徒写作时所持的观点：—\newline{}\newline{}ewline{}\newline{}\Quote{受造之物被屈服于虚妄，并将在天主子女的显扬中获得释放。}\newline{}\newline{}ewline{}\newline{}为避免有人以为我把自己想法强加给他，我将引用他的原话。\newline{}\newline{}ewline{}\newline{}\Quote{在世界终结与完成之际，当灵魂和具有理性的存有被主从牢狱中释放时，他们将根据此前的怠惰或勤奋而缓慢移动或快速飞翔。由于他们都有\OriginalTerm{自由意志}{free will}，可自由选择德性或恶习，选择恶习者将比现今更糟；选择德性者将改善其状况。他们在此方向或彼方向的运动与决定将决定他们各异的未来：也就是说，天使会变成人还是魔鬼，\OriginalTerm{魔鬼}{demon}会变成人还是天使。}\newline{}\newline{}ewline{}\newline{}接着，在提出各种论据支持其论点并坚持\OriginalTerm{魔鬼}{the devil}并非不能行善却尚未选择行善之后，他最终以大量篇幅推论：天使、人的灵魂和魔鬼——按他所说，三者本性相同但意志不同——可能因巨大的疏忽或愚蠢而受罚变成野兽。此外，为了避免惩罚的痛苦和灼烧的火焰，较敏感者可能选择变成低级生物，居住在水中，变成这种或那种动物的形状；因此我们有理由恐惧不只变成四足动物，甚至鱼类的变形。然后，唯恐自己被指为与\Person{毕达哥拉斯}一同主张\OriginalTerm{灵魂转世}{transmigration of souls}，他用伤害读者的邪恶推理总结道：\newline{}\newline{}ewline{}\newline{}\Quote{不可将我视为将这些事定为教义；它们只是作为推测抛出，以表明它们未被全然忽略。}\par

5. 在他的第二卷书中，他主张\OriginalTerm{多重世界}{plurality of worlds}；但不像\Person{伊壁鸠鲁}所教导的那样，许多相似的世界同时存在，而是每当旧的世界结束，一个新的世界就开始。在这个世界之前有一个世界，之后也将先后有另一个世界，如此循环往复。他怀疑一个世界是否会与另一个世界如此完全相似，以至于两者之间没有任何差异，或者一个世界绝不会完全无法与另一个区分。紧接着他又写道：\newline{}\newline{}ewline{}\newline{}\Quote{如果按照讨论的进程所必需的，一切事物都可以没有身体而存活，那么所有物质性的存在都将被吞没，那曾经从无中造成的将再次归于无有。然而，将来会有一个时刻，它的用处将再次成为必要。}\newline{}\newline{}ewline{}\newline{}在同一上下文中：\newline{}\newline{}ewline{}\newline{}\Quote{但是如果按照理性和圣经的权威，这必朽坏的要穿上不朽坏，这必死的要穿上不死，那么死亡将在胜利中被吞没，腐朽在不朽坏中被吞没。\BibleRef{格前 15:53} 并且或许所有物质性的存在都将被移除，因为只有在物质中死亡才能运作。}\newline{}\newline{}ewline{}\newline{}稍后又写道：\newline{}\newline{}ewline{}\newline{}\Quote{如果这些事与信仰不相违背，或许我们将来某一天会生活在无身体的状态中。此外，如果只有脱离身体的人才能完全服从基督，而所有的一切都必须服从祂，那么我们也将失去身体，当我们完全服从祂时。}\newline{}\newline{}ewline{}\newline{}在同一段落中：\newline{}\newline{}ewline{}\newline{}\Quote{如果所有的一切都必须服从天主，那么所有的一切都将脱去身体；然后所有物质性的存在都将归于无有。但如果因理性存有的堕落而再次需要身体，它将重新出现。因为天主留下灵魂去奋斗和挣扎，以教导他们：完全而圆满的胜利不是靠他们自己的努力，而是靠祂的恩宠。因此，在我看来，世界因导致它们的罪而各不相同，而那些主张所有世界都相同的理论已经过时。}\newline{}\newline{}ewline{}\newline{}接着又说：\newline{}\newline{}ewline{}\newline{}\Quote{关于终结，我有三个猜测；由读者决定哪一个最接近真理。因为，要么我们在无身体的状态下，当我们服从基督时，我们将服从天主，祂将成为万有中的万有；要么，正如服从基督的事物将与基督本身一起服从天主，并被置于同一法律之下，所有实体都将被提炼成最完美的形式，并稀薄化为\OriginalTerm{以太}{æther}，它是一种纯粹而单纯的本质；要么，我称之为不动的球体及其所有内容将溶解成虚无，而包含\OriginalTerm{反区域}{antizone}本身的球体将被称为\InnerQuote{好地,}\BibleRef{玛 13:8} 而另一个在其旋转中环绕地球并被称为天堂的球体，将保留为圣徒的居所。}\par

6. 这样说，难道不是最清楚地追随异教徒的错误，把哲学的妄言强加给基督徒的纯朴信德吗？在同一本书中他写道：\Quote{天主是不可见的。但如果祂本性不可见，那么即使对于救主也必须是不可见的。}在后面他又说：\Quote{没有一个灵魂曾降入人体，而像基督的灵魂那样带有如此真实的先前品格的印记，关于这灵魂他说：\InnerQuote{没有人能夺去它，而是我自己舍弃它。}}在另一处又说：\Quote{我们必须仔细考虑，灵魂在获得救恩并达到永福之后，是否会不再是灵魂。因为主和救主来寻找并拯救那迷失的\BibleRef{路 19:10}，为叫它不再迷失；同样，那被主来拯救的迷失灵魂，得救后就不再是灵魂了。我们必须自问，正如那迷失之物一度不存在、将不再存在，灵魂是否也曾存在、又将不再是灵魂。}在对灵魂作了不少评论之后，他引进了以下的话：\Quote{\OriginalTerm{理智}{\GreekText{νοῦς}}或}理智因堕落而成为灵魂；又因获得德行而再次成为理智。这至少可以从厄撒乌的事例中公平推论出来，他因先前的罪而被罚过较低级的生活。关于天体，我们也必须作类似的承认：太阳的灵魂——或无论你怎样称呼它——并非自创世时起存在；它在进入那发光炽热的身体之前就已存在。月亮和星辰也是如此。出于先起的起因，它们被屈服于虚妄，不是出于自愿，而是为了将来的赏报\BibleRef{罗 8:20}，被迫不行自己的意志，而行那给它们分派各自领域的造物主的意志。\par

7. 此外，地狱的火和圣经恐吓罪人的酷刑，他解释为不是外在的惩罚，而是有罪良心的痛苦，当藉天主的能力将我们过犯的记忆呈现在眼前时。\Quote{我们罪孽的全部收成从灵魂中存留的种子重新生长出来，我们所有可耻和不孝的行为再次出现在眼前。因此，是良心的火和悔恨的刺痛折磨着心灵，当它回顾先前的放纵时。}又说：\Quote{但也许这粗俗和属地的身体应当被描述为云雾和黑暗；因为在世界终末，当需要过渡到另一个时，同样的黑暗将导致同样的肉体出生。}这样说，他显然支持毕达哥拉斯和柏拉图所教导的灵魂轮回。在第二卷结束时，在论及我们的成全时，他说：当我们取得如此进步，不仅不再成为肉体或身体，而且或许也不再成为灵魂时，我们完美的理智和知觉，不被任何情欲的云雾所遮蔽，将面对面地辨别理性的和可理解的实体。\par

8. 在第三卷中，包含以下错误陈述。\Quote{如果我们一旦承认，当一个器皿成为贵重的，另一个成为卑贱的\BibleRef{弟后 2:20}，这是由于先前的因由；那我们为何不回溯灵魂的奥秘，承认它在雅各伯身上被爱、在厄撒乌身上被恨，是因它先前的行为，在雅各伯成为夺权者之前，在厄撒乌被夺权之前？}\BibleRef{拉 1:2}又说：\Quote{灵魂有贵贱之分的事实，应由它们先前的历史来解释。}在同一处说：\Quote{按照我的假说，一个本是贵重的器皿，若没有完成其目的，在另一世界将变为卑贱的器皿；相反，一个因先前过错而被定为卑贱的器皿，若在此生接受改正，在新的造化中将变为\InnerQuote{圣洁、合乎主用、预备行各样善事的器皿}\BibleRef{弟后 2:21}。}他紧接着说：\Quote{我相信，那些从小过犯开始的人，可能在邪恶中变得如此刚硬，以致若不悔改归向更好的事，就必成为非人的能力；相反，敌对的魔鬼生灵，可能因时间而如此治愈他们的创伤，并遏制他们先前罪恶的洪流，以致能达到成全者的居所。如我常说的，在灵魂生活和存在的无数不息的世界里，有些越来越坏，直到达于堕落的极深处；而另一些在那些深处越来越好，直到达到德行的成全。}这样他试图表明人，或更确切地说人的灵魂，可能成为魔鬼；而魔鬼反过来可能恢复为天使的品级。在同一本书中他写道：\Quote{这一点也必须考虑：人的灵魂为何时而受这种影响，时而受那种影响。}他揣测这是由于出生前的行为。他认为，正因如此，若翰在母胎中跳跃，当玛利亚问候时，依撒伯尔自称不配受她的注意。\BibleRef{路 1:41}他立即补充说：\Quote{另一方面，刚断奶的婴孩被邪灵附着，成为占卜的和算命的；确实，有些从幼年起就被\OriginalTerm{皮同}{python}的灵内住。既然他们没有做什么事招致这些访临，那么凡认为没有天主的许可就没有事发生，且万事都由祂的正义所统辖的人，不能认为天主的眷顾无缘无故地遗弃了他们。}\par

9. 他又论及世界写道：\Quote{我自认为有这样的信念：在这个世界之前已有一个世界，在此之后还有另一个世界。你想知道这个世界消亡后会有新的世界吗？请听依撒意亚的话：\InnerQuote{我所要造的诸天和新地，将常在我面前。}\BibleRef{依 66:22} 你想知道在这个世界造成之前已有其他的世界吗？请听训道者的话：\InnerQuote{已有的事，后必再有；已行的事，后必再行。日光之下并无新事。岂有一件事人能指着说这是新的？哪知，在我们以前的世代早已有了。}\BibleRef{训 1:9} 这段话不仅证明曾有其他的世界，也将有其他的世界；然而并非同时并存，而是一个接一个的。}他随即又说：\Quote{我认为天是天主的居所，真正的安息之所；有理智的受造物在那里享受原始的幸福，然后降落到一个较低的层面，以可见的取代不可见的，坠落到地上，并需要物质的身体。既已堕落，创造者天主为他们造了适合其环境的身体，并塑造了这个可见的世界，派遣了其中使者以确保堕落者的救恩和矫正。这些使者中，有些持有指定的职位，服从于世界的必然律；而另一些则在天主计划所定时日，明智地履行赋予他们的职责。前者包括太阳、月亮、星辰，被宗徒称为\InnerQuote{受造物}，它们被分配在天的高处。受造物屈服于虚妄，因为它被物质身体所包裹，并为肉眼所见。然而它\InnerQuote{屈服于虚妄，并非出于自愿，而是出于那使它屈服的，带着希望}。其他第二类的使者，在只有其造物主知道的特定地方和时间，我们相信是天主派来治理世界的众天使。}稍后他又说：\Quote{世界的事由天主眷顾如此安排：一些天使从天上坠落，另一些则自由地滑向地上。前者被迫坠落，后者仅凭选择降临。前者被迫在一段时期内履行一项令人厌恶的职务；后者自发地承担帮助堕落者的任务。}他又写道：\Quote{由此可知，这些不同的运动导致不同世界的创造；我们这个世界将被一个完全不同的世界所接替。关于这种堕落与上升，对德行的赏报和对邪恶的惩罚，无论它们发生在过去、现在或将来，唯有创造者天主能分配应得并令万有归向同一终向。因为只有祂知道，为什么祂允许一些天使随其倾向，从高处降到最低；为什么祂眷顾另一些，伸手将他们拉回原状，并再次安置于天上。}\par

10. 在讨论世界的终结时，他使用了以下言辞。\Quote{由于，如我常说的，终结产生新的开始，可能会问：那时身体是否继续存在，或者，当身体被消灭后，我们是否将无形体地生活，并与我们所知的天主一样无形体？现在，毫无可疑，如果身体，或如宗徒所称的可见之物，只属于我们这个可感世界，那么离体后的生命将是无形的。}稍后他又说：\Quote{当宗徒写道：\InnerQuote{受造物将从朽坏的奴役中被释放，进入天主儿女光荣的自由}，我这样解释他的话。有理智且无形体的存有是受造物中的最高者，因为不披戴身体，他们不是朽坏的奴隶。因为哪里有身体，那里就一定会有朽坏。但将来\InnerQuote{受造物将从朽坏的奴役中被释放}，那时人们将接受天主儿女的光荣，天主将成为万有中的万有。}在同一段中他又写道：\Quote{最终状态将是无形的，这一点由我们救主的祈祷词所证实：\InnerQuote{父啊，你在我内，我在你内，愿他们也在我们内合而为一。}\BibleRef{若 17:21} 因为我们应当理解天主是什么，救主最终将是什么，以及这里应许给圣人们的与圣父圣子的相似如何在于：正如他们自身合一，我们也将在他们内合一。因为如果最终圣人们的生命要与天主的生命同化，我们或须承认宇宙的主宰披戴身体，并像我们被肉体包裹一样被物质包裹；或者，如果认为这样不妥——尤其对于那些仅凭微小线索推断天主的威严并猜测祂固有超然本性之荣耀的人——我们便陷入以下两难：要么我们永远有身体，那么就必须对肖似天主绝望；要么，如果天主的生命之福真的应许给我们，那么祂生命的条件必须是我们生命的条件。}\par

11. 这些段落证明了他对复活的看法。因为他显然主张所有肉身都将消亡，我们将成为无肉身的，如同他所说我们在接受现在的肉身之前那样。再者，当他论及世界的多样性并主张天使将变成魔鬼，魔鬼要么变成天使要么变成人，而人又变成魔鬼时；简言之，一切都将变成别的东西，他这样总结自己的意见：\Quote{无疑，经过一段时间后，物质将重新存在，肉身将被形成，不同的世界将被创造，以满足理性受造物变化无常的意志；这些受造物因失去持续到终点的完美福乐，逐渐陷入如此巨大的邪恶，以致改变其本性，拒绝持守他们最初那纯全福乐的状态。许多理性受造物，理当说，保持这状态直至第二、第三、第四个世界，并不给天主改变他们境况的理由。另一些堕落得如此轻微，似乎几乎毫无损失；又另一些则不得不被一头扔进深渊。唯有安排万有的天主知道如何按照各人的功过在适当的领域中使用每一等级；因为只有祂了解时机、动机以及世界必须被驾驭的航程。因此，一个夺得邪恶棕榈枝并沉入最底层堕落的人，将在以后被塑造的世界中成为一个魔鬼，成为主手工艺品的初果，成为那些失去最初美德的天使的笑柄。} 这不是在论证这个世界的罪人可能在另一个世界成为魔鬼或鬼魔；反之，现今的魔鬼日后可能变成人或天使吗？在对这一观点进行冗长讨论之后，他主张所有有形受造物都必须将其物质换成精微和属神的身体，所有实体都必须成为一个纯净无比、不可想象的光明身体，人类心智目前对此无法形成任何概念，他这样总结：— \Quote{\InnerQuote{天主将成为万物之中的万有}；这就是说，所有有形存在都将被造得尽可能地完美；它将被带入神性本质，没有比这更好的了。}\par

12. 在他著作的第四卷即最后一卷中，以下段落应受教会的谴责。\Quote{可能就如当人在此世界因灵魂与肉身分离而死后，他们按行为的差异被分配在地狱的不同位置；同样，当天使死亡时，他们离开天上耶路撒冷的体系，降到此世界如同地狱，并按功过被安置在地上。} 又说：\Quote{正如我们将从此世界进入地狱的灵魂与那些从天上来到我们这里、在某种意义上死亡的灵魂相比较；同样，我们必须仔细查问，这是否毫无例外地适用于所有灵魂。因为在这种情况下，生在地上的灵魂当渴望更好的事物时，就升出地狱并取得人的肉身，或者当渴望更坏的事物时，就从更好的世界降至我们这里；在我们以上的穹苍中，同样有灵魂正从我们的世界前往更高的世界，另有一些虽然从天上坠落，但并未犯罪到需被扔到地上的程度。} 他由此试图证明穹苍即天空相对于天堂是地狱；而这大地相对于穹苍是地狱；再者，我们的世界对于地狱是天国。换言之，对某些人是地狱的，对另一些人却是天国。他对此仍不满足，接着说：\Quote{在万物的终结，当我们返回天上的耶路撒冷时，敌对势力将向天主的子民宣战，以振奋他们的勇气、锻炼他们的毅力、坚定他们的决心。因为他们只有面对并挫败敌人之后才能获得这些；我们在\BibleBook{}{户籍纪}中读到，他们因理性、纪律和战术技巧而克服敌人。}\par

13. 在说明按\BibleBook{若望}{默示录}，那要在天上揭示的\Quote{永远的福音}\BibleRef{默 14:6}超越我们的福音，正如基督的宣讲超越古法之圣事；他竟断言，连想都是亵渎的事：基督将为拯救魔鬼而在天上再次受苦。尽管他没有明说，但他的话暗示：正如天主为了人而成为人以释放人，同样，为了拯救魔鬼，当祂来释放它们时，祂将成为一个魔鬼。为表明这不是我的曲解，我必须引用他自己的话：\Quote{正如基督，} 他写道，\Quote{以福音的影履行了法律的影，并且一切法律都是天上之事的模型和影，我们必须探究我们是否有理由推想，甚至天上的法律与天上敬拜的礼仪仍未完备，而需要那真正的福音，这福音在若望默示录中称为永远的，以区别于我们这仅属暂时的、在一个将要逝去的世界中所阐明的福音。现在我们若将探究延伸到我们的主和救主的苦难，那么猜测祂将在天上受苦或许过于鲁莽；然而，若在天上有属灵的邪恶\BibleRef{弗 6:12}，并且我们毫不羞愧地宣认主曾被钉十字架，为毁灭那些祂藉苦难所毁灭的事物；那么，我们何必害怕想像在圆满时期内，在上界发生类似的事，以使所有领域万民都因基督的苦难而得救？}\par

14. 以下是另一句他论及圣子所说的亵渎之言。\Quote{假定圣子认识圣父，那么似乎通过这种认识，圣子能领悟圣父，如同工匠能领悟他的技艺规则一样。而且，无疑地，如果圣父在圣子内，那么圣父也被圣子所包含。但如果我们所说的包含不仅仅指认识者通过感知和洞察力领会事物，而是指他凭借一种特殊的能力将其包含在自己内，那么在这个意义上，我们不能说圣子包含圣父。因为圣父包含万物，而圣子是其中之一；因此，圣父包含圣子。}而为了向我们说明为何圣父包含圣子而圣子不能包含圣父的理由，他补充道：\Quote{好奇的读者可能会问，圣父认识自己的方式是否与圣子认识圣父的方式相同。但如果他记起这句话：\InnerQuote{派遣我的父比我大}\BibleRef{若 14:28}，他会承认这话必须普遍真实，并会承认，在知识上如同在其他一切事上，圣父大于圣子，并且圣父认识自己比圣子所能认识的更完美、更直接。}\par

15. 以下段落是令人信服的证据，证明他主张灵魂转世和身体的消灭。\Quote{如果能证明一个无形的、有理性的存在体自身具有生命，独立于身体，而且它在身体内比在身体外更糟；那么，毫无疑问，身体只是次要的，并且不时出现，以满足理性受造物的变化条件。那些需要身体的人就穿上身体，相反，当堕落的灵魂提升到更好的事物时，它们的身体就会再次消灭。它们就这样不断消失，又不断重现。}而为了阻止我们减轻他先前言论的亵渎程度，他结束他的作品时坚持说，一切理性的存在体，即圣父、圣子、圣神、天使、掌权者、宰制者和异能者，甚至因灵魂尊严而享有权利的人，都是同一个本质。\Quote{天主}，他写道，\Quote{及其独生圣子和圣神意识到一种理智和理性的本性。但是天使、掌权者、异能者，以及那按天主的肖像和模样受造的内心人也是如此。由此我得出结论，天主和他们在某种意义上是同一个本质。}他加上\Quote{在某种意义上}以避免亵渎的指控；然而在另一处他不允许圣子和圣神与圣父是同一实体，以免显得使神圣本质可分割，在此他却将全能天主的本性赋予天使和人。\par

16. 这就是\Person{奥利振}的书的本质，那么，修改一些关于圣子和圣神的亵渎段落，然后不加改变地出版其余部分，并附上一份不知廉耻的颂词，难道不是疯狂之举吗？当不改的部分和修改的部分出自同一个极度不虔敬的源头时，这难道不是疯狂吗？现在不是详细反驳所有陈述的时候；事实上，那些曾反对\Person{亚略}、\Person{欧诺米}、\Person{摩尼}以及其他各种异端的人，想必也已经驳斥了这些亵渎之词。因此，若有人想读这部作品，让他穿着鞋子向应许之地走去；让他防备蛇的嘴和蝎子的弯嘴；让他先读这篇论文，在他踏入这条路之前，让他知道必须避免的危险。\par

\SourceRecord{Translated by W.H. Fremantle, G. Lewis and W.G. Martley.；Nicene and Post-Nicene Fathers, Second Series；Vol. 6.；Edited by Philip Schaff and Henry Wace.；Buffalo, NY: Christian Literature Publishing Co.,；1893.；<http://www.newadvent.org/fathers/3001124.htm>.}

% structure: p0001 source=h1 target=section reason=page-title

\section{\WorkTitle{第一二五封信}}

\Emph{致\Person{鲁斯蒂库斯}}\par

\EditorialNote{鲁斯蒂库斯，图卢兹的一位年轻隐修士（需与第一二二封信的收信人仔细区分），热罗尼莫建议他不要成为独修士，而应继续留在团体中。此处为修道生活提出了一些规则，并生动描绘了好隐修士与坏隐修士的区别。热罗尼莫顺便发泄了对已故对手鲁菲努斯的怨恨（§18）。该信的日期为公元411年。}\par

1. 没有比基督徒更幸福的人，因为天堂的王国已应许给他。没有比他更奋力拼搏的人，因为他每天都冒着生命的危险。没有比他更坚强的人，因为他战胜了魔鬼。没有比他更软弱的人，因为他被肉身所胜。这两组陈述都可以用许多例子证明。例如，强盗在十字架上信了主，随即听到安慰的话：\Quote{我实在告诉你，今天你就要同我一起在乐园里。}\BibleRef{路 23:43} 而犹大则从宗徒职位的巅峰跌入丧亡的深渊。无论是宴席上的亲密交情、蘸饼的举动\BibleRef{若 13:26}，还是主的亲吻\BibleRef{玛 26:49}，都不能阻止他将自己所认识的天主子当作人来出卖。有谁比撒玛黎雅妇人更卑贱呢？然而不仅她本人信了，在六个丈夫之后找到了一个主，还在井边认出了默西亚——就是那些犹太人在圣殿中都没有认出的那一位；她给许多人带来了救恩，并在宗徒们去买食物的时候，满足了救主的饥饿，解除了他的疲乏。有谁比撒罗满更智慧呢？然而对女人的爱竟连他也变愚昧了。盐是好的，每份祭品都必须用它调和\BibleRef{肋 2:13}。因此宗徒也命令说：\Quote{你们的言谈要常带着恩宠，用盐调和}\BibleRef{哥 4:6}。但\Quote{如果盐失了味}\BibleRef{玛 5:13}，它就被扔在外面。它完全失去了价值，连作肥料也不配\BibleRef{路 14:35}，而信徒却从那里取粪来肥沃自己灵魂贫瘠的土地。\par

我这样开始，我儿\Person{鲁斯蒂库斯}，是要教导你你的事业之伟大和你理想之崇高；并让你看到，只有把年轻的私欲踩在脚下，你才有希望攀登真正成熟的巅峰。因为你所走的路是滑溜的，成功的荣耀小于失败的耻辱。\par

2. 我现在无需将我的论述之流引入美德的花园，也无需费心向你展示其几朵花的美丽。我无需详述百合的纯洁、玫瑰的羞红、紫罗兰的皇家紫色，或它们各种颜色所预示的灿烂宝石的应许。因为靠着天主的仁慈，你已经手扶犁头\BibleRef{路 9:62}；你已经像宗徒\Person{伯多禄}一样上了屋顶\BibleRef{宗 10:3}。当他在犹太人中间感到饥饿时，\Person{科尔乃略}的信德满足了他的饥饿，并通过百夫长和其他外邦人的归化，平息了因犹太人的不信而产生的渴望。由天上降到地上的那件器皿，四角象征四福音，教导他所有人都可以得救。同样，他在异象中被收回的那块洁白的布，是教会的象征，它将信徒从地上带到天上，是主应许的保证：\Quote{心里纯洁的人是有福的，因为他们要看见天主。}\BibleRef{玛 5:8}\par

这一切意味着我牵着你的手，尽我所能让你牢记某些事实；就像一个经历过多次海难的老练水手，我急于提醒一个缺乏经验的乘客前方的风险。因为一边是贪婪的卡律布狄斯，\Quote{万恶的根源}\BibleRef{弟前 6:10}；另一边潜伏着毁谤的斯库拉，系着诽谤的猎犬，宗徒说：\Quote{如果你们彼此咬伤、吞食，你们要小心，免得彼此消灭。}\BibleRef{迦 5:15}\par

3. 那些航行在红海——我们必须祈祷真正的法郎连同他的全军都淹没在其中——的人，在到达奥克苏马城之前，会遇到许多困难和危险。两岸都是游牧的野蛮人和凶猛的野兽。因此，旅行者必须时刻武装着和警惕着，并且随身携带一整年的食物。此外，海水里满是暗礁和无法通行的浅滩，所以必须不断从桅顶瞭望，指引舵手如何掌舵。他们若能六个月后到达上述城市的港口，就算幸运了。从这里开始是大洋，要横渡它几乎需要一整年。然后到达印度和恒河——在圣经中称为丕雄河——\Quote{环绕哈威拉全地}\BibleRef{创 2:11}，据说从它天堂的源头携带下来各种染料和颜料。这里有红宝石和绿宝石，闪闪发光的珍珠和一级宝石，是高贵女士们热切渴望的。还有金山，但是因为守护着山的狮鹫、龙和巨大的怪物，人们无法靠近；因为这就是贪婪为其财宝所必需的守卫。\par

4. 你会问，这一切有什么意义？当然很清楚。因为如果世俗的商人为了获得不确定和暂时的利益而承受这样的艰辛，并且在历经许多危险得到后，还冒着生命危险来保有它；那么基督的商人，他\Quote{卖掉了他所有的一切}，为要获得\Quote{那颗重价的珍珠}，他用全部财产买下一块田地，为要在那里找到宝藏——这宝藏既不能被贼挖走，也不能被强盗夺走——他又该做什么呢？\par

我知道我必定会触怒许多人，他们会对我的批评感到愤怒，认为这些批评是针对他们自身的缺陷。然而，这样的愤怒恰恰显示出不安的良心，他们对自己的判决将远比对我严厉。因为我不会提及具体名字，也不会效仿旧喜剧中指名道姓批评个人的做法。明智的男女会努力隐藏，或者更确切地说，纠正他们自己觉察到的错误；他们更会对自己生气，而不是对我，并且不会诅咒他们的规劝者。因为规劝者虽然也犯同样的过失，但至少在这一方面胜出：他对自己不满。\par

我听说你的母亲是一位虔敬的妇人，寡居多年；你年幼时她亲自抚养并教导你。后来，你在高卢的繁荣学校学习了一段时间后，她又送你去了罗马，不惜花费，并以对你未来前景的考虑来安慰自己与你的分离。她希望看到你高卢式雄辩术的丰沛与华丽被罗马的稳重所调和，因为她看出你更需要缰绳而非马刺。所以，我们听说希腊最伟大的演说家们曾用雅典的盐来调和亚细亚的浮夸，并在葡萄树生长过快时修剪它们。因为他们希望用优秀判断力的丰沛葡萄汁而非单纯词语的卷须来装满雄辩的榨酒池。你的母亲也为你做了同样的事；因此，你应当视她为父母，敬爱她如慈爱的乳母，尊她为圣人。你不可效仿那些离开自己亲属而追求陌生女子的行为。他们的不名誉是显而易见的，因为他们以纯洁感情为借口所追求的只是非法交往。我认识一些年长的妇女，甚至相当多，她们在自己的年轻被释奴身上寻欢，将他们视作自己的神子，并以此假装是他们的母亲，逐渐克服了自己的羞耻感，放纵自己在婚姻的自由中。另一些妇女抛弃自己未婚的姐妹，而与陌生的寡妇结合。有些人憎恨自己的父母，对亲属没有亲情。她们的心态被一种不屑于找借口的焦躁不安所显露；她们撕裂贞洁的面纱，像拂去蛛网一样抛弃它。妇女们的行径就是如此；当然，男人也好不到哪里去。因为有些人可见（尽管他们束着腰、穿着晦暗的服装、留着长须）与妇女形影不离，与她们同住一屋、同席用餐，有年轻女仆服侍他们，除了不被称为丈夫，与已婚无异。但这并不是基督教的过错，那伪善者堕入罪中；相反，外邦人的迷惑在于，教会谴责那些所有善人都谴责的事。\par

但是，如果你自己愿意成为隐修士，而不只是看起来像，那么更要关心你的灵魂而非财产；因为你一旦选择度修会生活，就已永远放弃了财产。让你的衣服粗劣以表明你的心灵洁白；让你的长衣粗糙以证明你轻视世界。但不要骄傲，免得你的衣着和言语不相符合。沐浴刺激感官，因此必须避免；因为消解天然热力是寒冷守斋的目的。然而，这些也应适度，因为若过量，会削弱胃，从而需要更多食物，导致消化不良——它是不洁欲望的丰沛根源。俭朴适度的饮食对灵魂和身体都有益。\par

你可以随意探望你的母亲，但不要与其他妇女同去，因为她们的面容可能存留在你的思绪中，于是\par

\Quote{秘密的创伤可能在你胸中溃烂。}\par

侍奉她的那些婢女，你务必要视作许多张设下捕捉你的网罗；因为她们的地位越是低下，你反而更容易毁掉她们。\Person{若翰洗者}有一位虔诚的母亲，他的父亲是一位司铎。然而，无论是母亲的慈爱还是父亲的财富，都不能诱使他留居父母家中而冒失去贞洁的危险。他住在旷野里，双目寻找基督，拒绝观看别的任何事物。他粗糙的衣着、皮制的腰带、蝗虫和野蜜的食物\BibleRef{谷 1:6}，这一切都旨在鼓励德行和节制。先知们之子——他们就是旧约时代的隐修士——曾在\Place{约旦河}畔为自己建造茅屋，离弃拥挤的城市，以菜蔬和野菜为生。当你还在家中时，要使你的斗室成为你的乐园，在那里收集圣经的各种果实，让圣经成为你最喜爱的伴侣，并将它的训诫铭记于心。如果你的眼睛、你的脚或你的手使你跌倒，就把它砍下来丢掉\BibleRef{玛 18:8}。为了拯救你的灵魂，不要吝惜任何东西。主说：\BibleQuote{凡注视妇女，有意贪恋她的，他已在心里奸淫了她。}\BibleRef{玛 5:28}智慧者写道：\BibleQuote{谁能说：我洁净了我的心？}\BibleRef{箴 20:9}在主的眼中星辰也不纯洁；何况那些一生就是漫长诱惑的人呢？\BibleRef{约 25:5}我们每次怀有贪欲时就犯了奸淫，我们有祸了！天主说：\Quote{我的剑，}在天上已饱饮；\Quote{何况在这长满了荆棘和蒺藜的大地上呢？}\BibleRef{创 3:18}被拣选的器皿\BibleRef{宗 9:15}，即嘴边常挂着基督之名的那位，他痛击自己的身体，使它屈服\BibleRef{格前 9:27}。然而连他也受到肉欲的阻碍，不得不做他不愿做的事。他好像在受苦，喊道：\BibleQuote{我这个人真不幸呀！谁能救我脱离这该死的肉身呢？}\BibleRef{罗 7:24}那么，除非你全心谨慎持守己心\BibleRef{箴 4:23}，并与救主一起说：\BibleQuote{我的母亲和我的兄弟，就是那些听天主的话而遵行的人}\BibleRef{路 8:21}，否则你能毫无失足或受伤地通过吗？这也许看似残酷，但实际上是爱。确实，还有什么比为一个圣洁的母亲保护一个圣洁的儿子更显示爱的呢？她也希望你获得永生，并且甘愿暂时见不到你，这样她才可以永远与基督一起见到你。她像\Person{亚纳}，她生下\Person{撒慕尔}不是为了自己的安慰，而是为了会幕的服役。\par

据记载，\Person{约纳达布}的儿子们既不喝酒也不饮烈酒，住在帐篷里，夜晚在哪里就在哪里支搭。\BibleRef{耶 35:6} 根据\WorkTitle{圣咏集}，他们是第一批被掳掠的人；因为当加色丁人开始蹂躏\Place{犹大}时，他们被迫逃到城里避难。\par

8. 别人会怎么想且随他们去，各人随己意而行。但对我来说，城市是监狱，孤独是乐园。为什么我们这些名字本身就代表孤独的人，却渴望城市的喧嚣呢？为了使\Person{梅瑟}适合领导犹太人民，他在旷野中受了四十年的训练；\BibleRef{宗 7:29}直到这时，牧羊人才成为牧人。宗徒们原是\Place{革乃撒勒}湖上的渔夫，后来才成为\BibleQuote{捕人的渔夫}。\BibleRef{玛 4:19} 但一听到主的召唤，他们便舍弃了所有的一切——父亲、渔网和船，每日背起十字架，手中连一根棍杖也没有。\par

我说这些话是为了，假如你渴望加入圣职人员的行列，你可以学会你以后必须要教导的事，好向基督献上合理的祭献\BibleRef{罗 12:1}，不至于认为自己是个合格的战士而自己还是个新兵，或者以为自己已经精通而自己还是个初学者。以我卑微的才能，我不宜审判圣职人员，也不该说那些在教会中服务的人的坏话。他们有自己的等级和地位，应当持守。如果你有朝一日成为他们中的一员，我写给\Person{内波提安}的公开信将告诉你适合该圣召的生活方式。目前我所谈论的是一位隐修士的形成和训练，是一位在年轻时接受了博雅教育后，将基督的轭负在颈上的人。\par

9. 首先要考虑的是，你应当独居还是与其他人一起在隐修院内生活。就我而言，我希望你与圣人们为伴，以免完全依靠自己。因为如果你还没有向导就走上一条陌生的路，你必定会立即偏左或偏右，易受错误的攻击，走得太远或不够远，因跑得太快而疲惫，或在路上闲逛而睡着。孤独时，骄傲会迅速潜入人心：一个人若禁食了一会儿，又不见任何人，就自以为了不起；忘记自己是谁、从何而来、往何处去，内心放纵思绪，在外鲁莽说话。与宗徒的意愿相反，他判断别人的仆人\BibleRef{罗 14:4}，伸手抓取一切他欲望所贪求的，睡到任意时长，谁也不怕，随心所欲，以为所有人都比他低，花在城里时间比在斗室里更多，而在弟兄们面前他假装退隐，却在街头与人群摩肩接踵。那么，你将会说：我是在谴责独居生活吗？绝不：事实上我经常赞扬它。但我希望看到隐修院学校培养出的战士，他们不害怕旷野的艰苦训练，在相当长的时间内展示了圣洁生活的榜样，他们使自己成为最后的，好成为最先的，不为饥饿或饱足所胜，以贫穷为乐，以服装和神态教导德行，并且极其谨慎，不会像一些愚昧之人那样编造与魔鬼斗争的骇人故事，以在众人眼中抬高自己的英雄形象，并首先从中榨取钱财。\par

10. 近来我们痛心地看到，一位隐修士去世后暴露了堪比\Person{克罗伊斯}的财富，而一座城市为穷人募集的施舍，却被他立遗嘱留给自己的儿子和继承人。铁斧沉底后又浮上水面\BibleRef{列下 6:5}，人又在棕榈树间看到了\Place{玛辣}的苦水。然而这并不奇怪，因为此人有一个同伴和导师，他把穷人的饥饿变成自己的财源，把留给可怜人的钱款据为己有，最终自招痛苦。但他们的呼声上达于天，进入天主的耳朵，耗尽了他的忍耐。因此，他派了一个灾祸的天使对这个新的加尔默罗人、第二个\Person{纳巴耳}说：\BibleQuote{糊涂人哪，今夜就要索回你的灵魂；你所备置的，将归谁呢？}（\BibleRef{撒上 25:38}；\BibleRef{路 12:20}）\par

11. 因此，我不愿你与母亲同住，是基于上述理由，尤其是以下两点。如果她给你美味佳肴，你拒绝会令她伤心；你若接受，则会为已在你内燃烧的火焰添加燃料。再者，在那样一个有许多女子的家中，你白天会看到夜间诱惑你的景象。绝不要让你的手和眼睛离开书本；逐字学习圣咏，不断祈祷\BibleRef{得前 5:17}，时刻警醒，勿让虚妄的念头占据你。将身心都导向主，以忍耐克服愤怒，热爱圣经知识，你就不再喜爱肉体的罪。不要让你的心成为激动的猎物，因为一旦它在你的胸中立足，就会统治你，带你陷入大过犯。总要手上有活干，让魔鬼发现你忙碌。如果那些有权靠福音生活的宗徒\BibleRef{格前 9:14}尚且亲手劳作，以免拖累任何人，并且救济别人——他们本有权利收取别人属肉体的事物，因为他们为他们播下了属神的事物\BibleRef{格前 9:11}——你为何不为自己预备所需呢？用芦苇编篮子，或用柔韧的柳条编筐；锄地，把菜园规划整齐；种下卷心菜或移栽植物后，开渠引水浇灌，使你亲眼看到诗人的优美描绘：\par

技艺从附近的山顶引来清水，\newline{}溪流哗哗地穿过岩石，\newline{}使干涸的草地凉爽，解其干渴。\par

用芽和插条嫁接不结果的树木，使你不久便能从它们那里摘取甜苹果，酬报你的辛劳。还要为蜜蜂制作蜂箱，因为\WorkTitle{撒罗满的箴言}派遣你到它们那里，从这小生物身上，你可以学习管理隐修院和治理王国。还要搓线钓鱼，抄写书籍，使你的手挣得食物，你的心因阅读而满足。因为\Quote{游手好闲者必为虚妄的欲望所掠。}在埃及，隐修院有一条规矩：不接受不愿劳作的人；因为他们认为劳动不仅为维持身体所必需，也为灵魂的救恩所必需。不要让你的心迷失在有害的思想中，也不要像在淫行中的\Place{耶路撒冷}那样，向每一个过客敞开双足\BibleRef{则 16:25}。\par

12. 在我年轻时，旷野以孤寂将我围困，我仍不能忍受罪过的冲动和我血液的自然热力；虽然我试图以频繁的斋戒来打破两者的力量，我的内心仍然涌动着[邪恶]的念头。为了制服其骚动，我去找一位归化前曾是犹太人的弟兄，请他教我希伯来文。于是，在我熟悉了\Person{昆体良}的尖锐、\Person{西塞罗}的流畅、\Person{弗龙托}的庄重和\Person{普林尼}的温和之后，我开始重新学习字母，研读那些既粗厉又喉音重的词语。我为这项任务付出了多少辛劳，经历了多少困难，曾多少次绝望，多少次放弃，然后因学习的渴望又重新开始，这些都可以由我这个痛苦经历者以及与当时与我同住的人来作证。但我感谢主，从这苦涩中播下的学问种子，如今结出了甘甜的果实。\par

13. 我还要讲述一件在埃及见到的事。在一个团体中，有一位年轻的希腊人，其欲望的火焰既非不断的斋戒也非最严厉的劳作所能熄灭。他面临陷入大罪的极大危险，隐修院院长用以下方法救了他。他命令一位年长的弟兄用责骂和斥责追赶他，然后在这样冤屈了他之后，抢先向他提出控诉。当证人被叫来时，他们总是为攻击者说话。听到这些谎言，他常常哭泣，因为没有人相信事实；只有院长巧妙地为他辩护几句，免得他\BibleQuote{被过度的忧伤吞噬}\BibleRef{格后 2:7}。长话短说，这样过了一年，期满时，年轻人被问及他从前邪恶的念头，以及它们是否仍困扰他。\Quote{哎呀，}他回答说，\Quote{连活着都不被允许，我怎能从淫乱中找乐子？}如果他是一个独居隐修士，谁能帮助他克服袭来的诱惑呢？\par

14. 世上的哲人以灌输新的激情来驱除旧的；他们钉入一颗钉子，拔出一颗钉子。波斯七位大臣对国王薛西斯正基于此原则行事：他们诱使国王爱上别的少女，以消除他对瓦市提皇后的惋惜。\BibleRef{艾 2:1} 然而，他们以一种过失治疗另一种过失，以一种罪治疗另一种罪；而我们却必须以学习热爱相反的德行来克服我们的过失。圣咏作者说：\Quote{离弃邪恶，}又说\Quote{行善；寻求和平，追随不舍。}因为我们若不憎恨邪恶，就不能热爱良善。不仅如此，若我们要离弃邪恶，就必须行善。若我们要避免战争，就必须寻求和平。仅仅寻求还不够；当我们找到它，而它从我们面前逃逸时，我们必须全力以赴地追随它。因为\Quote{它超越一切理解；}\BibleRef{斐 4:7}它是天主的居所。正如圣咏作者所说：\Quote{在和平中也有他的居所。}追随和平是一个美妙的比喻，可与宗徒的话相比：\Quote{追求款待。}他的意思是，仅以口头邀请客人是不够的；我们应该像追赶抢夺我们积蓄的强盗一样急切地留住他们。\par

15. 没有老师，任何技艺都无法学到。即使是哑巴动物和野群也跟随自己的首领。蜜蜂有蜂王，鹤以V形随在一只鹤后面。只有一个皇帝，每个省只有一个法官。罗马由两兄弟建立，但由于不能同时有两个国王，便以弑兄之举开国。同样，厄撒乌和雅各伯在黎贝加腹中争斗。\BibleRef{创 25:22} 每个教会有一位主教、一位总司铎、一位总执事；每个教会级别都服从自己的主管。一艘船只有一个舵手，一个家只有一个主人，最大的军队也听命于一人。为免列举过多使您厌烦，我的要点很简单：不要依赖自己的判断，而要住在隐修院里。因为在那里，你虽受一位长上管制，却会有许多同伴；这些同伴中，一位教你谦卑，另一位教你忍耐，第三位教你缄默，第四位教你温和。你照别人的意愿行事；吃别人吩咐的食物；穿别人给你的衣服；执行派给你的工作；服从你不喜欢的人；疲倦地上床；站着入睡；在休息不足之前被迫起床。轮到你的班次时，你要诵念圣咏，这任务需要的不是优美的声调，而是真诚的情感。宗徒说：\Quote{我要以神魂祈祷，也要以理智祈祷，}\BibleRef{格前 14:15} 又对厄弗所人说：\Quote{在你们心中向主咏唱。}\BibleRef{弗 5:19} 因为他读过圣咏作者的训诫：\Quote{你们要明智地咏唱。}你要服事弟兄们，要洗客人的脚；若受委屈，要默默忍受；你要敬畏团体的长上如主人，爱他如父亲。无论他命令你做什么，你都要相信对你有益。你不可判断在你以上的人，因为你的职责是服从他们，做所吩咐的事，正如梅瑟说的话：\Quote{以色列，要静默聆听。}你会有许多任务占据你的心思，使你无暇顾及[邪恶]的念头；当你从一件事转到另一件事，做完一件工作又有新的工作，你只能想到当下负责的事。\par

16. 但我本人见过与这截然不同的隐修士：他们的弃世仅在于换衣和口头宣示，而实际生活和以往习惯丝毫未变。他们的财产非但没有减少，反而增加了。他们仍是原来的仆役，摆着同样的宴席。从廉价杯盏和普通陶器中他们吞下黄金。身边仆从如云，却自称隐士。另有些人，虽贫穷却自以为是，趾高气扬如游行队伍般走街串巷，碰见谁便咆哮一通。还有人耸耸肩，咕哝着只有自己知道的东西。他们眼盯地面，舌上却掂量着浮夸的言词。只差一个传令员来向您宣告：总督大人驾到了。还有一些人，因房间潮湿、守斋严厉、独处厌倦、读书过度，以致耳鸣日夜不停，转为忧郁疯狂，这时需要的更是希波克拉底的热敷，而不是我的劝诫。许多人无法摆脱先前从事的技艺和贸易。他们不再自称商人，却照样从事同样的买卖；宗徒指示他们寻求的并非\Quote{食物与衣服}\BibleRef{弟前 6:8}，而是金钱利润，且比世俗人所求的更大。过去，卖家的贪婪受到市政官（希腊人称之为市场监察员）行动的约束，人们不能逍遥法外地欺诈。但现在，自称虔诚的人却无耻地谋取不义之财，基督信仰的美名更常常是欺诈的掩护而非受害者。我羞于启齿，却仍不得不说——我们至少应为自己的丑行脸红——在公共场合我们伸手乞讨，却在破衣烂衫下藏着黄金；令众人惊异的是，我们生前过穷人的生活，死后却财大气粗，钱袋充实。\par

但你，因为不是独处而是团体的一员，就不会有如此行事的诱惑。起初勉强的事，将变成习惯。你会不待吩咐就动手劳作，并在劳作中找到乐趣。你将忘记背后，努力面前的。\BibleRef{斐 3:13} 你少想别人所做的恶，多想自己当行的善。\par

17. 不要被众多犯罪者所引导，也不要让那些丧亡者的队伍诱使你暗地里说：\Quote{什么？所有住在城市里的人都要丧亡吗？看，他们继续享受自己的财产，他们服侍教会，他们频繁出入浴场，他们不轻视化妆品，然而他们普遍被人称赞。} 对于这类言论，我以前已经回答过，现在再次简短地回答：这本小册子的目的不是讨论圣职人员，而是为隐修士制定规则。圣职人员是圣洁的人，他们的生活总是值得赞扬。那么，你当振奋起来，在隐修院中如此生活，好使你配得成为一位圣职人员，好使你保全青春不受玷污，好使你如同新娘从洞房出来一样，走向基督的祭台。要注意，外面对你口碑良好，妇女们熟悉你的名声却不熟悉你的面容。当你到了成熟的年龄（如果你活到那么久），并且被城市的百姓或主教选入圣职的行列时，要行事与圣职人员相称，并选择你兄弟中最优秀者作为榜样。因为在每一个阶层和境遇中，恶人都与善人混杂在一起。\par

18. 不要被某种疯狂的任性冲昏头脑，仓促投入写作。要长期而仔细地学习你打算教导的内容。不要相信一切奉承者对你说的话，或者我宁可说，不要轻易听从那些意图嘲笑你的人。他们会用谄媚的赞美讨好你，尽力蒙蔽你的判断力；然而，如果你突然回头看，你会发现他们正用手做出嘲弄的手势，或是模仿鹳的脖子，或是驴子的扇耳，或是口渴的狗伸出的舌头。\par

不可毁谤任何人，或者以为通过攻击别人的名誉就能使自己更显优越。我们对他人的指控常常反过来落在我们自己身上；我们谴责那些同样属于我们的缺点，因此我们的雄辩最终反而成了对自身的控诉。这就像哑巴批评演说家一样。那位咕噜人想要说话时，总是以蜗牛般的速度走上前来，断断续续地吐出一两个字，中间的停顿如此之长，让人感觉他与其说是演讲，不如说是在喘气。然后，当他摆好桌子，在上面堆起一摞书时，他会皱起眉头，缩起鼻孔，蹙起额头，打响手指——这些动作意在引起学生的注意。接着，他就会倾泻出一连串的废话，对每个人猛烈抨击，以至于你会以为他是像\Person{朗吉努斯}那样的批评家，或者幻想他是第二个\Person{监察官加图}，对罗马的雄辩术进行评判，并随心所欲地把某些人排除在学术元老院之外。由于他很有钱，通过举办宴会使自己更受欢迎。许多人共享他的盛情款待；因此，他出门时总是被一群嗡嗡作响的人群簇拥着，这并不奇怪。在家中，他像\Person{尼禄}一样的怪物；在外，他像\Person{加图}一样的完人。他由不同且对立的性格组成，整体上令人难以描述。你可能会说，他是由不协调的元素构成的，就像诗人告诉我们的那种不自然且闻所未闻的怪物：\Quote{前面是狮子，后面是龙，中间是它所命名的山羊。}\par

19. 这样的人，你决不可看他们一眼，也决不要与他们交往。也不可让你的心偏向邪恶的言语，以免圣咏作者向你说：\BibleQuote{你坐下毁谤你的兄弟，中伤你母亲的儿子。} 以免你成为\BibleQuote{人子，他们的牙齿是矛和箭，} 并如同那人，\BibleQuote{他的话比油更柔软，但却是出鞘的剑。} 训道者说得更清楚：\BibleQuote{蛇在没有符咒的地方必会咬人，毁谤者也好不了多少。} 但你会说：\Quote{我不惯于毁谤，但我怎能阻止别人这样做呢？} 如果我们提出这样的借口，那只可能是因为我们想\BibleQuote{与作恶的人一同行恶。} 然而，基督不会被这种诡计欺骗。不是我，而是一位宗徒说：\BibleQuote{不要自欺；天主是轻慢不得的。}\BibleRef{迦 6:7} \BibleQuote{人看外貌，上主却看人心。}\BibleRef{撒上 16:7} 在\BibleBook{}{箴言}中，撒罗满告诉我们：正如\BibleQuote{北风驱走雨水，怒容也驱散毁谤的舌头。}\BibleRef{箴 25:23} 有时会发生这样的事：一支箭射向硬物时会弹回射箭的人身上，伤了那想伤人的人，因此应验了这句话：\BibleQuote{他们如弯曲的弓一样偏离了目标，} 另一处也说：\BibleQuote{谁将石头抛向高处，就落在自己头上。}\BibleRef{德 27:25} 所以，当毁谤者看见听者的脸上有怒容——那听者不听他，反而塞住耳朵，不听流血的事，\BibleRef{依 33:15}——他立刻沉默，脸色发白，嘴唇紧闭，口舌干燥。因此，同一位智者说：\BibleQuote{不要与毁谤者交往，因为他们的灾祸必骤然兴起；谁知道他们两人的毁灭呢？} 即说话者和听者。真理不喜欢角落，也不寻求暗中低语。对于弟茂德，经上说：\BibleQuote{控告长老的事，除了有两三个见证人，不可受理；但犯罪的人，当在众人面前责备他，使其余的人也恐惧。} 当一个人年纪老迈时，你不应轻易相信关于他的邪恶；他过去的生活本身就是辩护，他的长老地位也是如此。然而，我们既是人，有时尽管年岁成熟，仍会陷入青年的罪过，如果我做错了事，而你希望纠正我，请公开指责我的过犯，不要暗中毁谤我。\BibleQuote{愿义人击打我，结乃是仁爱；愿他责备我，如同膏油，但不要让恶人的油肥甘我的头。} 因为宗徒说：\BibleQuote{主所爱的，他必管教，又鞭打凡所收纳的儿子。}\BibleRef{希 12:6} 上主藉依撒意亚的口这样说：\BibleQuote{我的百姓啊，那叫你有福的人使你走错路，毁坏你当行的道。} 你把我的过犯告诉别人，对我有什么帮助呢？你可能在不知不觉中，因叙述我的罪——或者更确切地说，是你毁谤地加给我的罪——而伤害了另一个人；当你急于到处传扬时，你可能假装向每个人吐露秘密，好像没有告诉过别人。这种做法的目的不是纠正我，而是放纵你自己的缺点。上主吩咐说：\BibleQuote{那些得罪我们的人，要私下或在证人面前控告他；如果他不听劝告，就将事情告诉教会；那些坚持行恶的人，应被视为外邦人和税吏。}\BibleRef{玛 18:15}\par

20. 我极力强调这些事，为使我所爱的青年免于口舌和耳朵的痒症；好使他既在基督内重生，我就能把他如同贞洁的童贞女一样，毫无瑕疵或皱纹地呈献给基督，\BibleRef{弗 5:27}\BibleRef{格后 11:2}——不仅在身体上，也在心灵上保持贞洁；使他引以为傲的童贞名副其实，免得新郎因他没有善功的油而灯熄了，将他拒之门外。\BibleRef{玛 25:1} 关于普罗库鲁斯，你有一位可敬且博学的主教，他能用他宣讲的声音为你做得比我写下的书页更多，而且必能每日以讲道引导你走在正道上。他不会容许你偏右或偏左，或离开王道；因为以色列在匆忙前往应许之地时，承诺要遵循此道。\BibleRef{户 20:17} 愿天主垂听教会的恳求之声。\Quote{主，求祢为我们奠定和平，因为祢也为我们成就了一切作为。} 愿我们自愿而非被迫地弃绝世俗！愿我们欣然追求贫穷以获得其光荣，而非因他人强加而痛苦不堪！此外，在我们当前的苦难中，四周刀剑肆虐，凡有足够的面包维持所需的人便是富足的，凡不沦为奴隶的人便是极其强大的。图卢兹的可敬主教埃克苏佩里乌斯，效法匝尔法特的寡妇，\BibleRef{列上 17:8} 自己忍饥挨饿却喂养他人。他的面容因禁食而苍白，然而别人的饥饿更使他痛苦。事实上，他已倾其所有以供应基督穷人的需要。然而无人比他更富足，因为他的柳条篮内含有主的身体，他的普通玻璃杯内盛有宝贵的血。如同他的主一样，他将贪婪逐出圣殿；既不用绳索鞭子，也不用斥责言语，就推翻了卖鸽子者的座椅——即圣神的恩赐。他推翻了玛门的桌子，散尽了钱庄的钱财；热心于使天主的殿称为祈祷的殿，而非贼窝。你要紧紧追随他的脚步，以及与他有相同德行的其他司铎的脚步，他们因司铎职而成为穷人，却更加谦卑。或者，你若愿成为成全的，就同亚巴郎一起离开你的故乡和亲戚，往你不知道的地方去。你若拥有财产，就变卖并施舍给穷人。你若一无所有，那么你就免去了一个重担。你自己贫穷，就跟随贫穷的基督。这事艰巨，伟大而艰难；但赏报也是伟大的。\par

\SourceRecord{Translated by W.H. Fremantle, G. Lewis and W.G. Martley.；Nicene and Post-Nicene Fathers, Second Series；Vol. 6.；Edited by Philip Schaff and Henry Wace.；Buffalo, NY: Christian Literature Publishing Co.,；1893.；<http://www.newadvent.org/fathers/3001125.htm>.}

% structure: p0001 source=h1 target=section reason=page-title

\section{第一二七封信}

\Emph{致\Person{普林基帕}}\par

\EditorialNote{本信实为\Person{玛尔策拉}（见第二十三封信注）的传记，致其挚友。在叙述其生平、性格及喜爱的研究之后，\Person{热罗尼莫}继续讲述她在正统信仰事业中的卓越服务，当时通过\Person{鲁菲努斯}的努力，奥利振主义似乎有可能在罗马盛行（第9—10节）。他简要叙述了罗马的陷落及其后的恐怖（第12—13节），这似乎是\Person{玛尔策拉}死亡的直接原因（第14节）。本信写于公元412年。}\par

1. 基督的贞女\Person{普林基帕}，你曾多次恳切地请求我，愿我将一封书信献给那位圣妇\Person{玛尔策拉}的纪念，并述说我们长久以来所享的美善，使他人知晓并效法。我自己也极愿赞美她的功德，因此你催促我，并以为你的恳求是必需的，这使我感到痛心，因为我即使在爱她方面也不输于你。当我记录她卓越的德行时，我自己将获得的益处远超过我能给予他人的。若我至今保持沉默，让两年过去而不作任何表示，这并非如你错误地认为的那样是出于忽略她的意愿，而是出于一种难以置信的悲伤，它如此压倒了我的心，以至于我认为宁可暂时沉默，也不要以不恰当的语言称颂她的美德。现在我也不愿遵循修辞学的规则来赞美那位对我们两人并对所有圣人都如此亲爱的\Person{玛尔策拉}，她乃是她故乡罗马的荣耀。我不会陈述她显赫的家庭和高贵的门第，也不会追溯她历代执政官和禁军长官的谱系。我只称赞她自身的美德，这美德更为高贵，因为她抛弃了财富和地位，寻求了贫穷和卑微的真贵族。\par

2. 她父亲去世使她成为孤儿，她结婚不足七个月，丈夫就离她而去。当时她年轻、出身高贵，并以美丽（这对男人总是有吸引力）和自制而著称，一位名叫\Person{塞雷亚利斯}的杰出执政官非常殷勤地向她求婚。作为一位老人，他提出将自己的财产转让给她，使她觉得自己更像他的女儿而非妻子。她的母亲\Person{阿尔比娜}尽力为这位年轻的寡妇争取如此显赫的保护人。但\Person{玛尔策拉}回答说：\Quote{若我有意结婚，而非立志终身守贞，我寻找的是丈夫，而非遗产。}当求婚者辩解说有时老人活得很长，而年轻人却早逝时，她机智地反驳：\Quote{年轻人固然可能早逝，但老人不可能长寿。}这一对\Person{塞雷亚利斯}的果断拒绝使其他人确信他们无望赢得她的手。\par

在路加福音中，我们读到以下段落：\BibleQuote{有一位女先知亚纳，是阿协尔支派法奴耳的女儿，年纪已经老迈，从童女时起与丈夫生活了七年，就寡居了，直到八十四岁，从不离开圣殿，禁食祈祷，昼夜侍奉天主。}\BibleRef{路 2:36}她如此热切地寻求救主，因而得以见到他，并非奇事。让我们将她与\Person{玛尔策拉}相比较，便可见后者在各方面都更优越。\Person{亚纳}与丈夫生活了七年，\Person{玛尔策拉}只有七个月。\Person{亚纳}只是盼望基督，\Person{玛尔策拉}却紧握了他。\Person{亚纳}在他初生时承认他，\Person{玛尔策拉}却信他被钉死的基督。\Person{亚纳}没有否认那婴孩，\Person{玛尔策拉}却以他为君王而欢欣。我不愿像某些人惯于对待圣男和教会领袖那样，以功德来区分圣妇；我的结论是：既然她们有同一任务，便有同一赏报。\par

3. 在罗马这样一个好诽谤的社群中，它从前充满各地来的人，并以各种邪恶著称，毁谤攻击正直的人，并试图玷污纯洁清白的人。在这样的环境中，很难避开毁谤的气息。无瑕的名声是困难乃至几乎不可能的；先知渴望它，但几乎不指望得到它：\Quote{他说：\InnerQuote{在道路上无瑕的人是有福的，他们行走在主法律中。}}在现世道路上无瑕的人，是那些从未被诽谤之风玷污美名，也未在邻舍口中招致责难的人。因此，救主在福音中说：\BibleQuote{与你的对头（或作：和好）在路上的时候，快点与他言和。}\BibleRef{玛 5:25}有谁听说过关于\Person{玛尔策拉}的诽谤而值得丝毫信任？或者有谁相信了而不使自己成为恶意和诋毁的罪人？没有；她通过展示基督徒寡妇生活的本质——她的良心和面容都证实了这一点——使外邦人羞愧。因为世上的妇女惯于用胭脂和铅粉涂脸，穿闪亮的丝绸长袍，用珠宝装饰自己，脖上挂金链，穿耳并挂上红海最昂贵的珍珠，用麝香熏身。她们为失去的丈夫哀悼，却为自己的解脱和选择新伴侣的自由而欢喜——不是如天主意愿要服从他们\BibleRef{弗 5:22}，而是要统治他们。\par

为此目的，她们选择贫穷的丈夫，这些人满足于丈夫的名分，更愿意容忍竞争对手，因为他们知道，若他们稍作抱怨，便会立刻被抛弃。我们的寡妇的衣着是为了御寒，而非显露身材。她连一个金戒指也不愿戴，宁可将钱财储存在穷人的肚子里，而非自己保管。她无论去哪里都带着母亲，除非有证人在场，她绝不接见隐修士和神职人员，尽管她那宽敞的住所需要她接见他们。她的随从总是由贞女和寡妇组成，这些妇女庄重而稳健；因为她深知，婢女的轻浮反映主母的品格，而一妇人的品性可由其选择同伴而显明。\par

4. 她对圣经的喜爱令人难以置信。她时常歌唱：\BibleQuote{我将你的话藏在我心里，免得我得罪你}，以及描述完美之人的话：\BibleQuote{他喜爱主的法律，昼夜默思他的律法}。这种对法律的默思，她不是理解为像法利塞人那样反复查看文字，而是按照宗徒的话去行动：\BibleQuote{所以你们或吃或喝，或无论做什么，一切都要为光荣天主而做}。\BibleRef{格前 10:31} 她也记得先知的话：\BibleQuote{藉着你的诫命我获得明智}，并确信只有当她实践了这些，才会获准理解圣经。在此意义上，我们在别处读到：\BibleQuote{耶稣开始传道和施教}。\BibleRef{宗 1:1} 因为当人的良心责备他时，教导便蒙羞；如果一个人像克罗伊斯那样财富滚滚，尽管穿着破旧的外衣，家中却有防蛀的丝绸袍子，那么他口头上宣讲贫穷或教导施舍也是徒劳的。\par

玛尔塞拉实行禁食，但有节制。她戒吃肉，她了解酒的香味甚于它的味道；仅仅为了胃和常常的虚弱才触酒。\BibleRef{弟前 5:23} 她很少在公共场合露面，并小心避免进入贵妇们的房屋，以免被迫看到自己已经彻底弃绝的东西。她频繁前往宗徒和殉道者的圣殿，好能避开人群，专心祈祷。她对母亲如此顺从，以至于为了母亲，她做了自己并不赞成的事。例如，当她的母亲不顾自身后代，要将所有财产从子女和孙辈转移给兄弟的家庭时，玛尔塞拉希望把钱捐给穷人，却又不忍违抗母亲。因此，她将自己的首饰和其他财物转交给了已经富有的人，甘心舍弃自己的钱财，而不愿让母亲伤心。\par

5. 在那些日子里，罗马没有一位贵族小姐宣布过修道生活，也没有人敢于——当时这显得如此奇怪、可耻和卑下——公开称自己为修女。正是从一些来自亚历山大的司铎、从教宗阿塔纳修，以及后来从为了躲避亚略异端迫害而逃到罗马寻求安全的港湾共融的伯多禄那里，玛尔塞拉听说了当时仍在世的真福安当的生平、帕科米乌在泰拜德建立的隐修院，以及为贞女和寡妇制定的纪律。她并不羞于宣认这她所学到的、使基督喜悦的生活。多年之后，索弗若尼亚首先效法了她的榜样，随后还有其他人，关于她们，可以用恩尼乌斯的话说是：\par

但愿佩利翁的森林中，\newline{}斧头从未砍倒这些松树。\par

我可敬的朋友保拉蒙受玛尔塞拉的友谊，正是在玛尔塞拉的斗室里，优斯托基乌姆，这位贞女中的典范，逐渐接受了训练。由此容易看出，找到如此门徒的那位导师是何种类型。\par

不信仰的读者或许会嘲笑我如此长久地停留在对区区妇女的赞美上；然而，他若记得圣洁的妇女如何跟随我们的主救主并用她们的财物服事他，记得三位玛利亚如何站在十字架前，尤其是玛利亚玛达肋纳——因她信德的炽热和光芒而被称为\Quote{塔}——如何比众宗徒更先蒙恩见到复活的基督，那么他将会更快地指责自己的骄傲，而不是我的愚拙。因为我们判断人的德行不是根据性别，而是根据品格，我们认为那些弃绝了地位和财富的人值得最高的荣耀。正是为此，耶稣喜爱福音的作者若望胜过其他门徒。因为若望出身高贵，为大司祭所认识，却丝毫不为犹太人的阴谋所惧，竟将伯多禄领进大司祭的庭院，并且是唯一敢于站在十字架前的宗徒。正是他，将救主的母亲接到自己家中；\BibleRef{若 19:26} 正是这位贞洁的儿子，从主那里继承了贞洁的母亲作为遗产。\par

6. 玛尔塞拉于是过了多年的苦修生活，在她意识到自己曾经年轻之前，已发现自己老了。她常以赞同的态度引用柏拉图的话：哲学在于默想死亡。这一真理也为我们的宗徒所认同，他说：\BibleQuote{为了你们的得救，我天天冒死。}实际上，根据古老的抄本，我们主自己说：\BibleQuote{凡不背着自己的十字架天天跟随我的，不能做我的门徒。}早在许多世代以前，圣神就已藉先知说：\BibleQuote{为了你，我们终日被判死；我们被视为待宰的羊群。}许多世代之后，又有话说：\BibleQuote{要记得你的结局，你就永远不会做错，}\BibleRef{德 7:36} 以及那位雄辩的讽刺作家的教诲：\Quote{要活在死亡的意识中；时光飞逝；我说的这句名言也从中取走。}正如我所说，她度过她的日子，常生活在必死的想法中。她的衣着本身就是对她的坟墓的提醒，她将自己作为生活的、合理的、可悦纳的祭献呈献给天主。\BibleRef{罗 12:1}\par

7. 当教会的需要最终使我与可敬的主教\Person{保利诺}和\Person{厄丕法尼}一同来到\Place{罗马}——前者治理\Place{叙利亚的安提约基雅}教会，后者主持\Place{塞浦路斯}的\Place{撒拉米斯}教会——我因谦逊，原想避开贵妇们的目光，但她却如此恳切地请求，\BibleQuote{无论得时或不得时} \BibleRef{弟后 4:2}，如宗徒所说，以至于她的坚持不懈终于克服了我的犹豫。那时，我作为圣经学者的名声颇受敬重，她每次来看我，总会问一些关于圣经的问题；她不会立即接受我的解释，反而会提出异议，但不是为了争辩，而是为了学习如何回答那些她认为可能对我的陈述提出的反对意见。我在她身上发现了多少德行和能力，多少圣德和纯洁，我不敢说；既怕超出人们的置信范围，又怕因提醒你已失去的福分而增加你的悲伤。我只说这一点：凡是我长期研究所得、通过不断默想成为我天性一部分的东西，她都品尝过，学习过，并化为己有。因此，在我离开\Place{罗马}后，若在圣经的见证上出现任何争议，人们都会向她请教以解决。她是如此智慧，如此懂得哲学家所谓\OriginalTerm{得体}{\GreekText{τὸ} \GreekText{πρέπον}}即在行为中恰当的分寸，以至于当她回答问题时，她不是作为自己的意见，而是作为来自我或他人的意见提出，从而承认她所教导的是她从别人那里学来的。因为她知道宗徒说过：\BibleQuote{我不许女人讲道} \BibleRef{弟前 2:12}，她不愿造成对男性的冒犯，因为许多男性（有时也包括司铎）向她请教隐晦可疑的问题。\par

8. 我听说你立刻接替了我在她那里的位置，你依恋她，而且如俗语所说，你与她之间从未有过一丝一毫的间隙。你们住在同一所房子里，共用同一个房间，以致城中人人都确信你找到了她作母亲，而她找到了你作女儿。在郊外，你们为自己找到了一处隐修的僻静之所，因着它的孤寂，你们选择了乡间而非城镇。你们共同生活了很长时间，许多妇女以你们为榜样塑造自己的行为，我欣喜地看到\Place{罗马}变成了另一个\Place{耶路撒冷}。为贞女而设的隐修院比比皆是，隐修士也数不胜数。事实上，天主的仆人如此之多，以至于隐修生活——以前是一个被人讥讽的称号——后来竟成了荣誉的象征。与此同时，我们以相互鼓励的话语来安慰彼此的分离，并在精神上偿还我们在肉身中无法偿还的债务。我们总是急切地等待彼此的信件，努力在关怀上超过对方，并在礼貌的问候中争先。这样，通过持续通信而有效地弥合的分离，并未损失多少。\par

9. 当圣玛策拉在神圣的安宁中如此事奉主时，这些省份兴起了一股异端的旋风，将一切搅得天翻地覆；其狂怒之甚，竟至不放过自身和一切善事。仿佛扰乱这里的一切还不够，它又将一艘满载亵渎之言的船驶入了\Place{罗马}本身的港口。碟子很快找到了盖子；异端者的泥脚玷污了罗马信仰的清水\BibleRef{厄 34:18}。难怪在街头和市场，占卜者可以击打愚人的后背，或者拿起棍棒敲碎非议者的牙齿；像这样恶毒污秽的教导在\Place{罗马}找到了可被诱骗的傻瓜。接着出现了奥力振\WorkTitle{论首要原理}的丑闻译本，还有那个\Quote{幸运}的门徒，他若从未遇到这样的老师，就真是幸运了。随后我的支持者们提出了驳斥，摧毁了法利塞人的论据，使他们陷入混乱。那时，圣玛策拉，她长久以来一直退让，以免被认为出于党派动机，挺身而出投入了缺口。她意识到罗马的信仰——曾受一位宗徒赞扬\BibleRef{罗 1:8}——如今正处在危险中，并且这个新异端不仅吸引了司铎和隐修士，还吸引了许多平信徒，此外还欺骗了那位自以为别人都像自己一样纯朴的主教，她便公开抵制其教师们，宁愿取悦天主而不取悦人。\par

10. 在福音中，救主称赞那不义的管家，因为他虽然欺骗了主人，却为自己的利益行事明智。\BibleRef{路 16:8} 异端者在此也采用了同样的策略；因为他们看见星星之火能燃起何等大的火灾，\BibleRef{雅 3:5} 见他们点燃的火苗已达屋顶，对许多人所行的欺骗再也无法隐藏，就向教会请求并获得了推荐信，以便显得直到他们离去之日都一直与教会保持圆满共融。不久之后，杰出的阿纳斯塔修斯继任教宗；但他很快被接走，因为世界之首不应在如此伟大之人的任期内被斩首。他无疑被移去，免得他祈求扭转天主一次而永远定下的判决。因为主对耶肋米亚论及以色列的话，也适用于罗马：\Quote{不要为这百姓祈求好处。他们禁食的时候，我不听他们的呼号；他们献全燔祭和素祭时，我也不悦纳；我却要用刀剑、饥荒和瘟疫消灭他们。}\BibleRef{耶 14:11} 你将会说，这与赞扬玛尔凯拉何干？我回答说：正是她发起了对异端者的谴责。正是她提供了先受他们教导、后为他们异端教导所迷惑的见证人。正是她指出了他们欺骗了多少人，并拿来了那本不敬的\WorkTitle{论首要原理}的书，这本书被那蝎子之手\Quote{改进}后在人手间流传。最后，正是她一封又一封地写信，呼吁异端者亲自出庭辩护。他们确实不敢来，因为他们良心不安，宁可缺席而败诉，也不愿面对控告者并被定罪。这光荣的胜利起源于玛尔凯拉，她是这伟大福佑的源头和原因。你们这些与她共享荣耀的人知道我说的是实话。你们也知道，在众多事件中我只提几件，以免冗长的复述使读者厌倦。若我多说，心怀恶意的人可能会以为我是在借赞扬妇女美德之名发泄自己的怨恨。我现在将讲下去。\par

11. 这旋风从西方刮到东方，在途中威胁要使许多高贵船只沉没。于是耶稣的话应验了：\Quote{当人子来时，能在世上找到信德吗？}\BibleRef{路 18:8} 许多人的爱德冷淡了。然而，少数仍爱慕真信德的人聚集到我身边。人们公开企图夺取他们的性命，用尽各种手段对付他们。迫害如此激烈，以致\Quote{连巴尔纳伯也被他们的伪善所带走；}不仅如此，他还犯了杀人罪，即便不是实际暴力，也是意志上的。然后，看，天主吹了一口气，风暴就过去了；如此，先知的预言应验了：\Quote{你收回他们的气，他们就死亡，归于灰土。当天他的思念就消灭。}以及福音的话：\Quote{糊涂人哪，今夜就要索回你的灵魂；你所预备的，将归谁呢？}\BibleRef{路 12:20}\par

12. 当这些事在耶布斯发生时，从西方传来可怕的消息：罗马被围，其公民被迫用黄金买命。在被这样洗劫后，他们再次被围，以致不仅失去财产，还失去生命。我的声音哽咽，当我口述时，哭泣使我不能言语。那曾征服全世界的城市竟被攻陷；更甚的是，饥饿先于刀剑，只有少数公民幸存下来成为俘虏。饥饿的狂人求助于可怕的食物，彼此撕咬以果腹。甚至母亲也不放过怀中的婴儿。\Quote{摩阿布在一夜之间被占领，在一夜之间墙垣倒塌。}\BibleRef{依 15:1}\Quote{天主啊，外邦人进入了你的产业，污秽了你的圣殿，使耶路撒冷成为果园。你仆人的尸首给了天空的飞鸟作食物，你圣徒的肉给了地上的野兽。他们在耶路撒冷周围如水流血，无人埋葬。}\par

谁能述说那夜的屠杀？\newline{}什么眼泪能等同其痛苦？\newline{}古老的首都倾覆了；\newline{}无数市民的尸体\newline{}横陈在街道和房屋中。\newline{}死亡以许多恐怖的形象出现。\par

13. 同时，在这种混乱的场面中，理所当然地，一位沾满血迹的胜利者来到\Person{玛尔塞拉}的屋子。现在让我来说我听到的，讲圣善的人所见的；因为当时有这样的人在场，他们说\Person{普林奇比亚}您也和她在危险中。据说当士兵进入时，她毫无惊恐地接待他们；当他们向她索要金子时，她指着自己粗糙的衣服，表明她并无埋藏的财宝。然而他们不信她自选的贫穷，便鞭打她，用棍棒击打她。据说她不觉痛苦，却俯伏在他们脚前，流着泪为您求情，以免您从她身边被带走，或因您年轻而不得不经受她作为老人无需害怕的事。基督软化了他们铁石的心肠，即使在染血的刀剑中，自然亲情也彰显了它的权能。蛮族把您和她都带到了\Place{圣保禄宗徒大殿}，好让您在那里找到安身之所，若不然，至少也是一处坟墓。于是\Person{玛尔塞拉}大喜，感谢天主因她的祈祷而保全了您毫发无损。她也感谢说，城陷时发现她贫穷，而不是使她变得贫穷；她如今缺少日用的食粮，基督满足了她的需要，使她不再感到饥饿；她在言语和行动上都能说：\Quote{我赤身出离母胎，也要赤身归去；主赐予，主收回；愿主的名受赞颂！}\par

14. 几天后，她在主内安眠，但直到最后她的精力未衰。您是她清贫的继承人，更确切地说，穷人藉着您继承。她闭目时，是在您的怀中；她呼出最后一口气时，您的嘴唇接住了它；您流泪，她却微笑，自觉度过了善生，并盼望来世的赏报。\par

在短短一夜之间，我笔录了这封信，以赞美您，可敬的\Person{玛尔塞拉}，和您，我的女儿\Person{普林奇比亚}；并非炫耀我的口才，而是表达我对您二位由衷的感激；我唯一的愿望是取悦天主和我的读者。\par

\SourceRecord{Translated by W.H. Fremantle, G. Lewis and W.G. Martley.；Nicene and Post-Nicene Fathers, Second Series；Vol. 6.；Edited by Philip Schaff and Henry Wace.；Buffalo, NY: Christian Literature Publishing Co.,；1893.；<http://www.newadvent.org/fathers/3001127.htm>.}

% structure: p0001 source=h1 target=section reason=page-title

\section{第一二八封书信}

致\Person{加温狄}\par

\EditorialNote{加温狄从罗马写信给\Person{热罗尼莫}，请教如何养育他的幼女；因她按照当时宗教习俗已被奉献于童贞生活。热罗尼莫的回信可与致\Person{莱塔}的书信（第一〇七封）相比较，两者极为相似。信中对公元410年\Person{阿拉里克}洗劫\Place{罗马}的生动描述也十分引人注目。这封信的日期为公元413年。}\par

1. 给一个小女孩写信很困难，她听不懂你说什么，你不知道她的心思，对她的倾向只能鲁莽地预言。用一位著名演说家的话说：\Quote{她应受的称赞更多是因她将来会成为什么，而不是因她现在是什么。}因为对于一个渴望糕饼、在母亲膝上咿咿呀呀、觉得蜂蜜比任何话语都甜美的小孩，你怎么能跟她谈节制呢？她会听宗徒的深奥道理吗，当她的全部乐趣都在童谣里？她会留意先知隐晦的话吗，当她的保姆一板脸就能吓到她？她怎能理解福音的威严，当它的光辉使最敏锐的才智也目眩？我要劝她服从父母吗，当她用胖胖的小手拍打她微笑的母亲？基于这些理由，我亲爱的\Person{帕卡图拉}，我此刻寄给她的信，她得改天再读。与此同时，让她学字母、拼写、文法和句法。为了诱导她用细细的嗓音重复功课，拿蛋糕、蜜酒和糖果作为奖赏。如果她希望之后得到一束鲜艳的花、一件闪亮的小玩意儿、一个迷人的娃娃，她就会赶快完成功课。她还得学纺线，用柔嫩的拇指搓出纱线来；即使她老是断线，总有不再断线的一天。然后，当她做完功课时，应该有些娱乐。那时她可以搂着母亲的脖子，或者从亲戚那里抢吻。她唱圣咏时给她奖励，这样她就会爱她所学的东西。这样，功课就会成为她的乐趣，不再需要强迫。\par

2. 有些母亲当她们把女儿奉献于童贞时，给她穿暗色的衣服，用深色披肩裹起来，不让她穿亚麻衣或戴金饰。她们明智地避免让她习惯于以后要放弃的东西。另一些母亲则持相反原则。\Quote{这有什么用？}她们说，\Quote{不让她接触这些？她不会看见别人有吗？女人喜爱装饰，许多节操无亏的女子打扮不是为了男人，而是为了自己。给她所要的，但要告诉她，最受称赞的是那些一无所求的人。与其因为得不到而想要，不如让她充分享受，从而学会鄙视它们。}她们继续：\Quote{这正是主对以色列子民采用的办法。当他们渴想在埃及的肉锅时，祂让成群的鹌鹑飞来，让他们吃到反胃。那些曾经过着世俗生活的人，更容易放弃感官的享乐，比那些从小不知道欲望的人更甚。}因为前者——她们论证——践踏已知的事物，而后者则被未知的事物所吸引。前者忏悔地躲避快感暗中的进袭，而后者则与感官的诱惑调情，以为它们像蜜一样甜，结果发现是致命的毒药。她们引用经文说\BibleQuote{淫妇的嘴唇滴下蜜糖，}\BibleRef{箴 5:3}这甜在吃的人口中但后来比胆汁更苦。\BibleRef{默 10:9}她们论证，这就是为何在祭祀上主时，既不献蜂蜜也不献蜡，\BibleRef{肋 2:11}而苦橄榄所出的油却在祂的圣殿中焚烧。\BibleRef{出 27:20}此外，逾越节吃的是苦菜，\BibleRef{出 12:8}，\BibleQuote{并用真诚和真理的无酵饼。}\BibleRef{格前 5:8}接受这些的人将在世上受苦。因此先知象征性地歌唱：\Quote{我独自静坐，因为我充满了苦涩。}\par

3. 那么，我回答，怎样呢？年轻人是否应当放纵，以便以后更坚决地拒绝放纵？绝不，他们反驳：\Quote{\BibleQuote{各人当在什么身份上蒙召，就在什么身份上安于现状。}\BibleRef{格前 7:24}\BibleQuote{有人蒙召时已受割礼，}——也就是说，作为童贞女？——\BibleQuote{就不要成为未受割礼的}\BibleRef{格前 7:18}——也就是说，不要寻求婚姻的外衣，这外衣是在亚当被逐出童贞乐园时赐给他的。\BibleRef{创 3:21}\BibleQuote{有人蒙召时未受割礼，}——也就是说，有妻子，被婚姻的皮肤包裹着？——就不要寻求童贞的赤身\BibleRef{创 3:25}和那永远失去的永恒贞洁。不，让他\BibleQuote{持守自己的身体，在圣化与尊荣中}\BibleRef{得前 4:4}，让他喝自己的井水，而不是娼妓放纵的水池，那些水池不能存留贞洁的清水。\BibleRef{箴 5:15}}同样，保禄在同一章讨论童贞和婚姻时，称那些已婚者为肉体的奴仆，而那些未受婚姻轭束缚的则是自由人，他们在完全的自由中事奉主。\BibleRef{格前 7:21}\par

我所说的并非普遍适用，只是部分地谈论这个主题。我只对一些人说话，不是对所有人。然而，我的话是针对男女两性的，而不仅仅是\BibleQuote{较软弱的器皿}\BibleRef{伯前 3:7}。你是守贞者吗？那你为何乐于与女子交往？为何将你脆弱的修补小舟驶入大海，如此自信地面对危险航程的巨险？你不知你想要什么，却依恋她，仿佛你曾渴望过她，或者最宽容地说，仿佛你将来会渴望她。你或许会说，女人比男人更适合做仆人。那就选一个畸形的老妇人，选一个在主内被认可节制的。你为何乐于与年轻、漂亮、肉感的女子相处？你频繁去浴场，油光满面、红润地行走，吃肉，财富丰裕，穿着最昂贵的衣服；而你竟以为能与一条致命毒蛇安全同眠。你或许告诉我，你并不与她同住一屋。这也只是晚上如此。但你整天与她交谈。为何独自与她相处？为何没有见证人？这样做，即使你没有实际犯罪，也显得在犯罪，（你的影响如此重要）你以榜样鼓励不幸的人去行恶。你，无论守贞者还是寡妇，为何允许一个男人长时间与你交谈？为何不害怕与他单独相处？至少为了自然的需要出门，并为此离开那个你对自己给予比对你兄弟更多自由、行为比对你丈夫更不端庄的男人。你说，你有关于圣经的问题要问。如果是，就公开问；让你的女仆和随从也听见。\Quote{凡显明出来的就是光}。只说该说的话的人不会寻找角落来说；他乐有听众，因为他喜欢被称赞。另一方面，他一定是个好老师，他不看重人，不关心兄弟，只在暗中劳作，仅仅教导一个软弱的女子！\par

3 \Emph{a.} 我暂时偏离了直接主题，讨论他人在类似案例中的做法；而我的目的是训练，更确切地说，哺育年幼的\Person{帕卡图拉}，却一时招致许多绝非和平之女的妇人的敌意。但现在我要回到正题。\par

一个女孩只应与女孩为伴，她不应认识男孩，甚至惧怕与他们玩耍。她绝不应听到不洁的话语，若在家中喧闹中偶然听到，也应不明白。她母亲的点头应如同口头命令一样对她有约束力。她应爱她如父母，顺服她如主母，尊敬她如老师。她现在还是个没有牙齿、没有思想的孩子，但一旦七岁，成为一个知道自己不该说什么、犹豫该说什么的腼腆女孩，直到她长大，她应背诵圣咏集和撒罗满之书；福音、宗徒和先知应是她心里的宝藏。她不应过于自由地在公共场合出现，或过于频繁地拥挤在教会中。她一切的快乐应在她内室。她绝不能看年轻人，也不能转眼向涂脂抹粉的纨绔子弟；而甜美嗓音女孩的淫荡歌曲会通过耳朵伤害灵魂，必须远离她。这类人越容易接近，她们来时就越难避开；她们自己一旦学会了什么，就会偷偷教她，从而以众人的闲言玷污我们幽居的\Person{达纳厄}。给她一位不贪酒、按宗徒所说不是闲荡和饶舌，而是清醒、庄重、勤于纺羊毛的女士兼家庭教师作为监护和同伴；她的言语会塑造孩子柔弱的心灵趋向美德。因为正如水随手指在沙中划动，柔软年幼的孩子也容易受善恶影响；你可以任意引导她。此外，衣着华丽的年轻男子常通过讨好保姆、仆从甚至贿赂来为自己寻求接近；当他们如此温和地接近后，就燃起情欲的第一星火花，直到爆发成火焰，并逐渐提出最无耻的要求。那时根本无法阻止他们，因为这句诗在他们身上应验了：\Quote{一旦允许根深蒂固，再斥责便为时已晚。}我羞于说却不得不说；出身高贵的女子拒绝了更高贵的求婚者，却与最低贱的人甚至奴隶同居。有时，她们以宗教和渴望独身之名，实际上为了这样的情夫而抛弃丈夫。你常能看到一位\Person{海伦}跟随她的\Person{帕里斯}，毫不畏惧\Person{墨涅拉俄斯}。我们看见这样的人并哀叹，却无法惩罚，因为罪人众多为罪行赢得了容忍。\par

4. 世界陷入毁灭：是的！但可耻的是，我们的罪仍旧活着并且繁盛。这著名的城市，罗马帝国的首都，在一场巨大的大火中被吞没；天下没有一处地方没有罗马人充军。一度被视为神圣的教堂，如今只剩下尘土和灰烬；然而我们仍一心贪求利益。我们生活得好像明天就要死去；但我们又大兴土木，好像要永远活在这世上。我们的墙壁金光闪闪，天花板和柱头也是如此；然而基督却在我们的门口，在祂的穷人身上，赤裸而饥饿地死去。我们读到，司祭亚郎面对熊熊火焰，把火放入他的提炉，阻止了天主的义怒。大司祭站在死者和生者之间，火焰不敢越过他的脚。另一次，天主对梅瑟说：\BibleQuote{你且由我……，让我消灭这人民}\BibleRef{出 32:10}，藉着\Quote{你且由我}这句话，显示祂可以因人的阻止而不施行所威胁的惩罚。祂仆人的祈祷阻碍了祂的大能。你认为，如今天下还有谁能阻止天主的义怒，面对祂审判的火焰，并同宗徒一起说：\BibleQuote{我宁愿自己受诅咒，为我的弟兄}\BibleRef{罗 9:3}？羊群与牧人一同灭亡，因为百姓怎样，司祭也怎样。\BibleQuote{百姓怎样，司祭也怎样}\BibleRef{依 24:2}。从前并非如此。那时梅瑟出于怜悯之情说：\BibleQuote{但如今，你若肯赦免他们的罪……；若不然，求你从你所写的册上涂抹我的名字。}\BibleRef{出 32:32} 他不满足于确保自己的救恩，他愿意与那灭亡的人一同灭亡。他这样做是对的，因为\BibleQuote{人民众多，是君王的光荣。}\BibleRef{箴 14:28}\par

我们的小\Person{帕卡图拉}就诞生在这样的时代。这就是包裹她的襁褓，她在这襁褓中第一次呼吸；她注定在欢笑之前先认识眼泪，在喜乐之前先感受悲伤。她刚登上舞台，就几乎被要求退场。那就让她以为世界从来就是现在这个样子吧。让她不知过去，让她躲避现在，让她渴望未来。\par

这些想法只是我仓促凑成的。因为我对逝去朋友的哀恸从未间断，直到最近我才恢复了足够的平静，写一封老人的信给一个小孩。我对你，兄弟\Person{高达狄乌斯}，的爱促使我尝试这么做，并且我认为说几句话总比什么也不说好。那使我意志麻木的哀痛将为我文字的简短开脱；然而，倘若我什么也不说，我友谊的真诚就很可能被怀疑。\par

\SourceRecord{Translated by W.H. Fremantle, G. Lewis and W.G. Martley.；Nicene and Post-Nicene Fathers, Second Series；Vol. 6.；Edited by Philip Schaff and Henry Wace.；Buffalo, NY: Christian Literature Publishing Co.,；1893.；<http://www.newadvent.org/fathers/3001128.htm>.}

% structure: p0001 source=h1 target=section reason=page-title

\section{书信129}

\Emph{致达尔达努斯}\par

\EditorialNote{作为对高卢总督达尔达努斯所提问题的答复，热罗尼莫论及预许之地，他将其认定为天堂而非客纳罕。随后他指出，犹太人现今所受的苦难完全是因他们犯下的将基督钉十字架的罪行所致。此信写于公元414年。}\par

\SourceRecord{Translated by W.H. Fremantle, G. Lewis and W.G. Martley.；Nicene and Post-Nicene Fathers, Second Series；Vol. 6.；Edited by Philip Schaff and Henry Wace.；Buffalo, NY: Christian Literature Publishing Co.,；1893.；<http://www.newadvent.org/fathers/3001129.htm>.}

% structure: p0001 source=h1 target=section reason=page-title

\section{第一三〇封信}

\Emph{致\Person{德梅特里亚斯}}\par

\EditorialNote{热罗尼莫写信给德梅特里亚斯，一位罗马的高贵女子，她最近接受了贞女的圣召。在叙述了她先在罗马、后在非洲的生活历史后，他继续制定规则和原则来指导她的新生活。这些规则涵盖了禁欲修行的全部领域，包括学习、祈祷、守斋、服从、为基督放弃财富以及持续工作的职责，实质上与三十年前热罗尼莫给欧斯托基乌姆建议的（第二十二封信）相似。然而，这封信的语气较为温和，不那么狂热；所推荐的苦修不那么严厉；少了狂喜，多了常识。这封信还应与白拉奇写给德梅特里亚斯的信相比较，该信收录于热罗尼莫著作集第十一卷（Migne's Patr. Lat. xxx. ed.）。日期是公元414年。}\par

1. 从我青年时期至今，凡我亲笔或由书记代笔所处理过的主题，没有比现在这一主题更困难的了。我将要写信给德梅特里亚斯，一位基督的贞女，一位在罗马世界中出身和财富无人能及的贵女。因此，如果我使用足以描述她德行的语言，人们会认为我是在奉承她；如果因某些细节似乎难以置信而略去不提，我的保留又会有亏于她无可置疑的长处。那么我该怎么办呢？我无法胜任当前的任务，却又不敢推辞。她的祖母和母亲都是杰出的妇女，她们同样有权命令，有信德寻求，以及有毅力获得她们所要求的东西。确实，她们向我要求的并非什么非常新颖或特殊的事；我的才智经常在类似的主题上得到锻炼。她们所希望的是，我应当提高嗓音，尽可能强烈地为一位——用那位著名演说家的话来说，她与其因现有的成就受赞扬，不如因她将来可能成就的受赞扬——之人的美德作证。然而，尽管她是个女孩，她有着超越年龄的热忱信德，并且从别人认为是显著美德终点的起点开始出发。\par

2. 让诋毁远离，让嫉妒退避；不要向我提出自私自利的指控。我作为一个陌生人写信给一个陌生人，至少就外貌而言是如此。因为内在的人借着那使宗徒\Person{保禄}认识哥罗森人以及许多他从未见过的其他信徒的知识，感到自己是众所周知的。我对这位贞女怀有多么崇高的敬意，更确切地说，我认为她是什么样的德行奇迹，你可以从以下事实判断：我正忙于解释\Person{厄则克耳}对圣殿的描述——这在全部圣经中是最难的部分——并且发现自己正处于圣殿中至圣所和香坛所在的那部分时，我选择作为一种短暂的休息，从那个祭台转到这个祭台，以便在这祭台上将活的祭献、蒙天主悦纳的奉献\BibleRef{罗 12:1}，毫无玷污地奉献给永恒的贞洁。我知道主教用祈祷的话语用贞女头纱覆盖了她神圣的头，同时朗诵了宗徒的庄严语句：\BibleQuote{我愿意把你们众人如同贞洁的童女献给基督。}\BibleRef{格后 11:2}她如同王后站在他的右边，身穿金线绣成的衣服和刺绣的袍子。那就是\Person{若瑟}所穿的彩色外衣，即由许多不同德行织成的；类似的服装古时是君王女儿们常穿的。于是新娘自己欢喜地说：\BibleQuote{君王带我进了他的内室}\BibleRef{歌 1:4}，她的同伴们回应说：\Quote{君王的女儿在内部全是荣耀的。}这样，她便成了一个公开献身的贞女。尽管如此，我的话仍不会没有用处。赛马的速度因观众的欢呼而加快；拳击手因支持者的呐喊而被激励做出更大努力；当军队列阵备战、刀剑出鞘时，将军的演说大大激发了士兵的勇气。现在也是如此。祖母和母亲已经栽种，但是我浇灌，主使之生长。\BibleRef{格前 3:6}\par

3. 修辞学家的惯例是通过提及被赞扬者的祖先和家族过去的尊贵来抬高他，以便肥沃的根可以弥补贫瘠的枝条，使你在树干上欣赏到在果实中没有的东西。因此，现在我应当回顾\Person{普罗布斯}家族和\Person{奥利布里乌斯}家族的高贵名字，以及那个著名的\Person{阿尼齐乌斯}家族，其代表人物很少或从未辜负过执政官之职。或者我应当提及我们贞女的父亲\Person{奥利布里乌斯}，他的英年早逝是\Place{罗马}不得不哀悼的。我不太敢多谈他，以免加深你圣洁母亲的痛苦，并且以免对他美德的纪念成为她悲伤的更新。他是一个孝子，一个可爱的丈夫，一个仁慈的主人，一个受欢迎的公民。他年幼时就被选为执政官；但他品格的善良使他作为一名元老更为显赫。他的死亡是幸福的，因为它使他免于看到祖国的毁灭；而他的后代更为幸福，因为他伟大的曾祖母\Person{德梅特里亚斯}那显赫的名字，如今因他的女儿\Person{德梅特里亚斯}发誓终身守贞而变得更加显赫。\par

4. 但我在做什么？我忘记了自己的目的，满怀对这位青年的钦佩，竟赞颂起纯粹世俗的好处来；而我本应称赞我们的童贞女，因为她拒绝了这一切，并决心不把自己视为富家或高贵的女士，而仅仅视自己为一名普通女子。她的意志力几乎令人难以置信。尽管她可以随意享用丝绸和珠宝，尽管她周围簇拥着成群的宦官和侍女，一个满是奉承和殷勤仆人的喧闹家庭，尽管大宅丰裕所能提供的精美筵席日日摆在她面前，她却宁愿选择严斋、粗衣和俭朴的生活。因为她曾读过主的话：\BibleQuote{那穿华丽衣服的人是在王宫里。}\BibleRef{玛 11:8} 她满心钦佩厄里亚和若翰洗者的生活方式；两人都用皮带束腰，克制并苦修自己的身体，而后者被说成是在厄里亚的精神和能力中来临，作为主的先驱。因此，他还在母胎中就预言了，\BibleRef{路 1:41} 并在审判之日以前赢得了审判者的称赞。她也钦佩法奴耳的女儿亚纳的热忱，她甚至在极其年老时仍以祈祷和斋戒在圣殿中事奉主。\BibleRef{路 2:36} 当她想到斐理伯的四个女儿，都是童贞女，\BibleRef{宗 21:9} 她就渴望加入她们的行列，与那些藉童贞纯洁获得预言恩宠的人同列。她用这些和类似的默想滋养自己的心灵，最怕的莫过于冒犯祖母和母亲。虽然她们的表样鼓励她，但她们明言表达的希望和愿望却使她气馁；并非她们不赞同她的神圣目的，而是这奖赏如此巨大，她们不敢期望或企及。因此，这位基督军中可怜的新手极度困惑。她开始憎恨自己所有华服，像艾斯德尔一样向主呼求：\BibleQuote{你知道我厌恶我尊位的记号}——即她作为王后所戴的王冠——\BibleQuote{我厌恶它如同月经的污布。}\BibleRef{艾 14:16} 在她见过并认识她的那些圣洁高贵的女士中，有些人曾被席卷非洲的风暴从高卢海岸驱赶，避难至圣地。她们告诉我，德默特利雅秘密地一夜又一夜（除在祖母和母亲随从中的献身天主的童贞女外无人知晓）拒绝亚麻床单和羽绒床铺，而将山羊毛的毯子铺在地上，以不断的泪水湿润自己的脸。一夜又一夜，她在思想中俯伏在救主的膝前，恳求他接纳她的选择，实现她的愿望，并软化祖母和母亲的心。\par

5. 我为何还要延迟叙述结局？当她的婚期已近，当新房正为新郎和新娘准备时，她秘密地，没有见证者，只有黑夜作伴，据说是这样鼓励自己：\Quote{德默特利雅，你怎么了？你为什么如此害怕捍卫你的贞洁？你需要的是自由和勇气。如果你在和平时期都如此惊慌，那么当被召去殉道时，你会怎样？如果你连亲人一个皱眉都受不了，又怎能面对迫害者的法庭？如果男子的榜样不能打动你，至少从真福殉道者亚格尼那里获得勇气和信心，她战胜了青春和暴君的诱惑，并以她的殉道圣化了贞洁之名。不幸的女子！你不知道，你不知道你的童贞应归于谁。不久之前，你还在蛮族人手中颤抖，紧贴着祖母和母亲，躲在她们的外衣下以求安全。你曾看见自己沦为俘虏，你的贞洁不在自己手中。你曾对敌人的凶狠目光战栗；你曾暗中痛苦地看见天主的童贞女被玷污。你的城市，曾是世界之都，如今成了罗马人民的坟墓；而你，在利比亚海岸上，自己是个流亡者，却要接受一个流亡者作丈夫？你到哪里找一个主妇来参加你的婚礼？你找谁来护送你回家？除了刺耳的布匿语舌头，没有人为你唱放荡的费斯钦诗歌。抛开一切犹豫！天主的\BibleQuote{完美爱情驱除恐惧}。\BibleRef{若一 4:18} \BibleQuote{拿起信德的盾牌、正义的胸甲、救恩的头盔}，\BibleRef{弗 6:14} 出去战斗吧！守护你的贞洁本身就包含一种殉道。你为什么害怕祖母？你为什么惧怕母亲？也许她们自己希望你所做之事，只是她们不认为你愿意为自己做。}当她以这些和其他的理由将自己鼓动到必要的决心程度时，她将所有的饰物和世俗衣裳都抛在一边，当作障碍物。她的贵重项链、昂贵的珍珠和闪闪发光的宝石都放回了盒子里。然后她穿上一件粗布长袍，披上一件更粗糙的外衣，在一个意想不到的时刻进来，突然跪倒在祖母膝前，泪流满面地显出她的真面目。那位端正圣洁的老妇看到孙女如此奇异的装扮，惊讶不已。她的母亲则喜不自胜。两位妇人都几乎不敢相信她们所渴望的是真实的。她们的声音哽在喉咙里，又是脸红，又是苍白，又是恐惧，又是喜悦，被许多矛盾的情感所折磨。\par

6. 我在此不得不让步，不再试图描绘那难以言喻之事。在试图解释那令人难以置信的喜乐的伟大时，\Person{图利乌斯}的雄辩之流会干涸，\Person{狄摩西尼}所投掷的雷霆也会耗尽而落空。人心所能想象、言语所能表达的一切人间欢乐，都在那里显现。母亲与女儿、祖母与外孙女一再彼此亲吻。两位年长妇女因喜乐而痛哭流涕，她们扶起倒在地上的少女，拥抱她那颤抖的身躯。在那心意中，她们认出了自己的心思，彼此祝贺如今一位贞女要以她的童贞使一个高贵的家族更加高贵。她们说\Quote{她已找到一条造福家庭、减轻罗马毁灭之灾的道路。}善哉耶稣！整个家中是何等的欢腾！如同从丰产的树干上瞬间萌发许多嫩枝，许多贞女随之而出，她们的女主人和夫人所立的榜样被一大群依附者和仆人效仿。童贞在每个家庭中被热切地接纳，虽然那些宣认童贞的人在肉体上地位低于\Person{德默特里亚}，但她们与她追求同一奖赏，即贞洁的奖赏。我的言语太无力了。\Place{阿非利加}的每一座教会都为之欢舞。消息不仅传到了城市、市镇和村庄，甚至传到了分散的茅舍。\Place{阿非利加}与\Place{意大利}之间的每一个岛屿都充满了这消息，这喜讯传遍四方，无人不喜。于是\Place{意大利}脱去了丧服，罗马那已毁的城墙部分恢复了昔日的辉煌；因为他们相信，他们养女的彻底归正是天主眷顾他们的标记。你会以为哥特人已被消灭，那群逃奴和奴隶已被主从天上的雷霆击倒。当罗马人在\Place{特雷比亚河}、特拉西梅诺湖和坎尼阵亡数千人之后，\Person{马塞卢斯}在\Place{诺拉}取得第一次胜利时，罗马的欢庆也没有如此热烈。当被困在\Place{卡皮托利山}上的贵族——罗马的未来系于他们——用黄金买命之后，听说高卢人终于被击溃时，他们的喜乐也没有如此之大。消息传到了东方海岸，这一基督信仰的胜利在遥远的内陆城市也被听闻。哪一位基督贞女不以\Person{德默特里亚}为伴为荣？哪一位母亲不称\Person{尤利亚纳}的胎是有福的？不信者或许会嘲笑未来赏报的不确定性。然而，就在成为贞女之时，你所获得的已超过你所舍弃的。你若成为人妻，只有一个省份会认识你；而作为基督贞女，你为整个世界所知晓。那些对基督信心不大的母亲们，不幸地常常只把畸形残疾、找不到合适丈夫的女儿奉献给童贞。正如俗语所说，\Quote{玻璃珠被当作珍珠。}那些以信仰自夸的人，给守贞女儿的钱财仅够维持生计，却把大部分财产留给生活在世俗中的儿女。最近在本城，一位富有的司铎留给他的两个发愿守贞的女儿微薄的资财，却为其他子女提供纵情享乐的充足钱财。我遗憾地说，许多度苦修生活的妇女也做了同样的事。但愿此类事例罕见，可惜并不罕见。然而，它们越是常见，那些拒绝效仿这众多先例的人就越有福。\par

7. 所有的基督徒都高声赞美基督神圣的同轸者，因为当\Person{德默特里亚}宣认自己为贞女时，他们将原本为她婚姻准备的嫁妆钱财给了她。他们不愿亏待她天上的新郎；事实上，他们希望她带着她从前的一切财富来到祂面前，以免这些财富浪费在世俗事物上，而能用来救济天主的仆人们的困苦。\par

谁会相信呢？那位\Person{普罗巴}——在罗马世界中所有高位和出身高贵者中，她拥有最显赫的名声，她的圣洁生活和普世爱德甚至为她赢得了蛮族的尊敬，她对她的三个儿子\Person{普罗比努斯}、\Person{奥利布里乌斯}和\Person{普罗布斯}所担任的常规执政官职务毫不在意——我说，那位\Person{普罗巴}，如今罗马已被攻陷，其财物或被烧或被抢，据说她正在变卖她所有的一切，\BibleQuote{为自己结交不义钱财的朋友，好使这些朋友接她进入永久的帐幕里。}\BibleRef{路 16:9}教会的各级牧者以及那些只在名义上是隐修士的人，真该羞愧脸红，因为他们正在购置地产，而这位贵妇却在变卖它们。\par

她刚刚从蛮族手中逃脱，刚刚停止为她怀中那些被蛮族夺走的贞女哭泣，就被一场突如其来的、无法承受的丧亲之痛压倒——这痛苦是她从未担心过的，即她爱子的去世。然而，作为一位将要成为基督贞女祖母的人，她承受住了这致命打击，怀着对未来的希望，以自身证实了那抒情诗人的话：\par

\Quote{即使整个宇宙碎裂崩塌，它的坠落\newline{} 可击中义人，可杀死他，却不能使他惊慌。}\par

我们在\WorkTitle{约伯传}中读到，当第一个报恶信的人还在说话时，另一个人也来了；\BibleRef{约 1:16}同一卷书上又记载：\BibleQuote{人岂不是一个考验}——或者按希伯来文更好的表达——\BibleQuote{在世上的争战吗？}\BibleRef{约 7:1}我们为此劳苦，为此在这世界的争战中冒生命危险，为要在来世得到冠冕。我们相信这一点适用于人并不奇怪，因为连主自己也受过试探，圣经也为亚巴郎作证说，天主试探了他。\BibleRef{创 22:1}宗徒也因此说：\BibleQuote{我们以困苦为荣……知道困苦生忍耐，忍耐生经验，经验生望德，望德不叫人蒙羞；}\BibleRef{罗 5:3}在另一处又说：\BibleQuote{谁能使我们与基督的爱分离？是困苦吗？是窘迫吗？是迫害吗？是饥饿吗？是赤身露体吗？是危险吗？是刀剑吗？正如所记载的：我们整天被置于死地，被视为待宰的羊。}\BibleRef{罗 8:35}依撒意亚先知这样安慰处境相同的人：\Quote{你们这些断奶的、离开乳房的，要期待困苦上加困苦，但也要期待望德上再加望德。}因为，正如宗徒所说：\BibleQuote{现今时代的苦楚，与将要显于我们身上的光荣相比，就不足以介意了。}\BibleRef{罗 8:18}我为何在此汇集这些经文，下文将会阐明。\par

普罗巴从海上望见故乡的烟火，将自己和亲人的安全托付给一叶扁舟，却发现阿非利加的海岸比她离开的更为残酷。因为有一个难以判断是更贪婪还是更冷酷的人在等候她，此人除酒与金钱外别无所求，表面上侍奉最温和的皇帝，实际上却成了最野蛮的暴君。若允许我引用诗人的寓言，他就像塔尔塔罗斯中的奥尔库斯。同样，他身边还有一只刻耳柏洛斯，不是三头，而是多头，随时准备攫取并撕碎触手可及的一切。他将许配给别人的女儿从母亲怀中夺走，把高贵的少女嫁给最贪婪的叙利亚商人。贫穷的恳求不能使他放过任何受监护的少女、寡妇或献身于基督的贞女。事实上，他看求助者的手多于看他们的脸。这就是这位夫人在逃避蛮族时所遭遇的可怖的卡律布狄斯和犬围的斯库拉；这些怪物既不怜悯船只遇难者，也不理会被俘之人的哭喊。残忍的恶徒！至少学学罗马帝国的敌人。我们时代的布伦努斯只取他所发现的，而你却搜求你所不能发现的。\par

事实上，德行常遭嫉妒，批评者可能惊异于普罗巴通过秘密协议买得众多同伴的贞洁。他们可能声称，那位伯爵本来可以拿走一切，却不会满足于一部分；而且她不能质疑他的要求，因为尽管她地位高贵，在他专横的手中却只是一个奴隶。我也意识到自己正招致敌人的攻击，似乎是在奉承一位出身最高贵、最显赫的夫人。然而，当他们得知我至今未提及她时，就不能指责我了。我从未在她丈夫生前或去世后称赞过她的古老家族或庞大的财富与权力——这些题材他人或许会用来作逢迎的演说。我的目的是以适合教会的方式称赞我这位贞女的祖母，并感谢她以善意帮助德默特利亚实现了愿望。至于其他方面，我的斗室、衣食、年迈和窘迫的境遇足以驳斥奉承的指责。在我余下的信中，我将把全部话语转向德默特利亚本人，她的圣德使她与她的地位一样高贵，可以说，她爬得越高，跌倒就越可怕。\par

此外\par

天主的女儿，这一件事我嘱咐你，\newline{}是的，首要的事，并多次敦促你：\par

要喜欢以阅读圣经来占据你的心灵。不要在你胸中的好土壤里只收获毒麦和野燕麦。不要当家主睡觉的时候，让仇敌把莠子撒在麦子中间\BibleRef{玛 13:25}（即当那常依附于天主的心灵未加提防之时）；而要常与\WorkTitle{雅歌}中的新娘一起说：\BibleQuote{夜间我寻找我心爱的。请告诉我：你在何处牧羊？正午你在何处使羊群歇息？}并与圣咏作者一起说：\BibleQuote{我的灵魂紧紧追随祢；祢的右手扶持了我。}并与耶肋米亚一起说：\BibleQuote{我追随祢并不觉得困难……}因为\BibleQuote{在雅各伯中没有悲伤，在以色列中也没有劳苦。}当你在世俗中时，你爱世俗之物。你用胭脂涂抹面颊，用铅粉改善肤色。你梳理头发，用不是你自己的发辫在头上筑起一座高塔。我不必提你昂贵的耳环、从红海深处采来的闪亮珍珠、鲜绿的翡翠、闪烁的缟玛瑙、清澈的蓝宝石——这些使贵妇们心动、渴望拥有的东西。因为你已经抛弃了世俗，并且在圣洗誓愿之外又立了一个新誓；你已与你的仇敌立约，并说：\Quote{我弃绝你，魔鬼，以及你的世俗、你的浮华和你的作为。}因此，要遵守你所立的盟约，并在这个世俗的旅途中与你的仇敌保持和约。否则，他某日可能将你交给审判者，证明你拿走了属于他的东西；然后审判者将把你交给差役——既是你的仇敌又是你的报复者——你就被投入监牢；投入那外边的黑暗中\BibleRef{玛 8:12}，这黑暗因将我们与基督——唯一的真光\BibleRef{若 8:12}——隔绝而更加可怕。你断不能从那里出来，直到你还清最后一文小钱\BibleRef{玛 5:25}，即直到你为最微小的罪做补赎；因为我们在审判之日要为每一句废话交账\BibleRef{玛 12:36}。\par

8. 我说这些，并非想对你发出不祥的预言，而只是想警告你，像一个忧虑而谨慎的劝诫者，在你的事情上连安全之处也提心吊胆。圣经说：\BibleQuote{若统治者的灵升起反对你，不要离开你的位置。}我们必须常全副武装、列阵以待，准备与仇敌交战。当他试图将我们击退、使我们后退时，我们必须稳稳站定，与圣咏作者同说：\BibleQuote{祂把我的脚立在磐石上，}以及\BibleQuote{岩石是野兔的藏身之处。}在后一句中，许多人将\Quote{野兔}读作\Quote{刺猬}。刺猬是一种小动物，非常胆怯，浑身布满刺。当耶稣戴上荆棘冠冕，背负我们的罪，为我们受苦时，是为了使我们从荆棘和蒺藜中生出童贞的玫瑰和贞洁的百合——自从厄娃听到\BibleQuote{你要在痛苦中生子；你的情欲要指向你丈夫，他要管辖你}\BibleRef{创 3:16}以来，荆棘和蒺藜就成了妇女的命运。我们被告知，新郎在百合花中牧羊\BibleRef{歌 2:16}，即在那些没有玷污自己衣裳的人中，因为他们保持了童贞\BibleRef{默 14:4}，并听从了训道者的诫命：\BibleQuote{你的衣裳要时常洁白}\BibleRef{训 9:8}。作为童贞的创始者和元首，祂大胆地说论自己：\BibleQuote{我是沙仑的玫瑰，谷中的百合}\BibleRef{歌 2:1}。所以\BibleQuote{岩石是野兔的藏身之处}，这些野兔在一城受迫害时就逃往另一城\BibleRef{玛 10:23}，并不惧怕先知的话\BibleQuote{救助离弃了我}会在他们身上应验。\BibleQuote{高山是野山羊的藏身之处，}它们的食物是小孩从洞中引出的蛇。同时，豹子与小山羊同卧，狮子像牛一样吃草\BibleRef{依 11:6}；当然不是为使牛从狮子学得凶残，而是为了使狮子从牛学得温顺。\par

但让我们回到最初引用的经文：\BibleQuote{若统治者的灵升起反对你，不要离开你的位置，}接着这句话还有：\BibleQuote{因为退让平息大过犯。}\BibleRef{训 10:4}意思是，如果蛇潜入你的思想，你必须\BibleQuote{全心谨守你的心}\BibleRef{箴 4:23}，并与达味同唱：\BibleQuote{求你洗净我隐密的过犯，也阻止你的仆人犯骄傲的罪，}不要陷入\Quote{极大过犯}，即行为上的罪。相反，要在恶的诱惑还只是思想时就杀死它们；将巴比伦女儿的婴孩摔在石头上，使蛇无法留下痕迹。要谨慎，向主许愿：\BibleQuote{求祢不让他们管辖我，然后我便正直，并从极大过犯中清白无辜。}因为圣经在其他地方也证明：\BibleQuote{我必追讨父亲的罪孽直到子孙三四代。}\BibleRef{户 14:18}也就是说，天主不会立刻因我们的思想和决心惩罚我们，但会将报应降到它们的后代，即从它们产生的恶行和罪习上。正如祂通过亚毛斯的口所说：\BibleQuote{为了某城的三种过犯，甚至四种，我必不收回惩罚。}\BibleRef{亚 1:3}\par

9. 我从圣经的美丽田野中掠过，采集了这些少许的花朵。它们足以提醒你，必须关闭你心扉的门户，并时常画十字圣号以坚固你的额角。唯有如此，埃及的毁灭者才找不到攻击你的地方；唯有如此，你灵魂的长子才能逃脱埃及人长子的命运；唯有如此，你才能与先知一同说：\BibleQuote{我的心灵已坚定，天主，我的心灵已坚定；我要歌唱，我要赞美。我的荣耀，醒来吧；醒来吧，琴瑟与竖琴。}因为，即使罪恶深重的提洛也被吩咐拿起她的竖琴\BibleRef{依 23:15}并做补赎；像伯多禄一样，她被告知要用痛苦的眼泪洗去她先前污秽的痕迹。然而，我们不要知道补赎，免得想到它而陷入罪恶。补赎是那些不幸遭遇船难者的木板；但一位贞洁无瑕的童贞女可望保全船本身。因为，寻求你所遗失的是一回事，保持你从未失去的则是另一回事。就连宗徒也克制自己的身体，使它服从，免得他给别人宣讲了，自己反而被淘汰。\BibleRef{格前 9:27}他被情欲的暴力燃烧着，成了人类族群的代言人：\BibleQuote{我这个人真悲惨啊！谁能救我脱离这死亡的身体呢？}又说：\BibleQuote{我知道，在我内，即在我的肉体内，没有善：因为愿意为善的能力已在我内，但如何去行善，我却找不到。因为我所愿意的善，我不去做；我所不愿意的恶，我却去做。}再次说：\BibleQuote{随从肉性的人，不能得天主的欢心。但你们已不属于肉性，而是属于圣神，只要天主的圣神住在你们内。}\BibleRef{罗 8:8}\par

10. 在你最审慎地留意你的思想之后，你必须穿上禁食的铠甲，与达味一同歌唱：\BibleQuote{我以禁食克制我的灵魂，}和\BibleQuote{我吃灰烬如同吃饭，}以及\BibleQuote{至于我，当他们扰乱我时，我的衣服是苦衣。}厄娃因吃了禁果而被逐出乐园。反之，厄里亚在禁食四十天后，被火马车载到天上。梅瑟四十昼夜借着与天主的亲密交谈而生活，从而在他自己身上完全证实了这句话：\BibleQuote{人生活不只靠饼，也靠上主口中所发的一切言语。}世界的救主，祂在德行和生活模式上给我们留下了效法的榜样，在祂受洗后，立即被圣神带到旷野，与魔鬼争斗，\BibleRef{玛 4:1}并在粉碎和推翻他之后，把他交给门徒们践踏。因为宗徒说什么？\BibleQuote{天主快要将撒殚踏碎在你们脚下。}\BibleRef{罗 16:20}然而，在救主禁食四十天后，那古老的仇敌借着食物向祂设下了陷阱，说：\BibleQuote{你若是天主子，就命这些石头变成饼吧。}\BibleRef{玛 4:3}在法律下，第七个月，在吹角节之后，当月第十日，为全体犹太人民宣布了禁食，那日在自我放纵而不克己的人，应从百姓中剪除。在《约伯传》中，论及贝希摩特写道：\BibleQuote{它的力量在它的腰间，它的能力在它肚腹的肚脐上。}我们的仇敌利用青春热情的热度来诱惑少男少女，\BibleQuote{点燃我们出生的轮子。}他如此应验了欧瑟亚的话：\BibleQuote{他们都是奸淫者，他们的心像火炉；}这火炉只有天主的仁慈和严厉的禁食才能熄灭。这些是\BibleQuote{火箭}\BibleRef{弗 6:16}，魔鬼用它们伤害人并点燃他们，巴比伦王曾用它们来对付三个青年。但是，当他使火高四十九肘时，他反而将自己毁灭，因为主曾指定七个七年为救恩时期。\BibleRef{肋 25:8}当时，第四位看似天主子的人缓解了可怕的炙热\BibleRef{达 3:25}，冷却了熊熊烈火，直到它们虽然看似可怖，却变得无害；如今对童贞女的灵魂也是如此。天上的露水和严厉的禁食扑灭了少女心中的情欲火焰，使她的灵魂即使在尘世的帐幕中也能度天使般的生活。因此，那蒙拣选的器皿\BibleRef{宗 9:15}声明，关于童贞女，他没有主的命令。\BibleRef{格前 7:25}因为，如果你要放弃你的自然功能，砍断你自己的根，不采集除了童贞之外的果实，弃绝婚床，避免与男人交往，并在肉身内如同出于肉身而生活，你就必须反乎自然，或者说是超越自然。\par

11. 然而，我并不强求你过度禁食或异常节食。此类做法往往在奠定圣善生活的基础之前，就拖垮了虚弱的体质，引发身体疾病。哲人们有一格言：德行在于中庸，一切极端皆为恶的性质；正是基于此意，七智者之一提出了那喜剧中引用的著名箴言：\Quote{凡事不可过分}。你不可禁食到心跳加速、呼吸艰难、需要他人搀扶或背负的地步。不，在克制肉身欲望的同时，你当保持足够的力量去阅读圣经、咏唱圣咏并守夜。因为禁食本身并非完整的德行，而只是其它德行得以建立的基础。圣化以及那非此不能见主的贞洁也是如此。每一步都是向上的阶梯，但没有一步能单独赢得童贞的荣冠。福音在智愚童女的比喻中教导了我们：前者进入新郎的洞房，后者却被关在门外，因为她们没有善行之油，让灯熄灭了。\BibleRef{玛 25:1} 禁食这个主题开启了一个广阔的领域，我本人常在其中漫游，许多作者也撰写了相关论著。若你想了解自律的益处以及暴食的害处，请参阅这些著作。\par

12. 要效法你的净配：\BibleRef{路 2:51} 服从你的祖母和母亲。除非在她们的陪伴下，绝不可注视男子，尤其是青年男子。绝不可结识她们所不认识的人。世俗的格言是：只有建立在同好同恶之上的友谊才是可靠的。你已通过她们的榜样以及家中的圣善生活得到教导，向往童贞，认识基督的诫命，知道何为你的益处，以及你当选择的道路。但不要把你自己的事物视为绝对属于自己。要记住，其中一部分属于那些将贞洁传递给你的人，而且你像一朵精选的花，从他们尊贵的婚姻和未受玷污的床榻中绽放出来。\BibleRef{希 13:4} 因为你注定要结出完美的果实，只要你谦卑自己于天主大能的手下，\BibleRef{伯前 5:6} 常常记着经上所记载的：\BibleQuote{天主拒绝骄傲的人，却赐恩宠给谦卑的人。}\BibleRef{伯前 5:5} 哪里有恩宠，就不是因行为而赐予，而是施予者的白白恩赐，如此便应验了宗徒的话：\BibleQuote{这不在于愿意的，也不在于奔跑的，而在于显示慈悲的天主。}\BibleRef{罗 9:16} 然而，愿意或不愿意是在于我们；同时，那属于我们的自由本身也完全是因天主的慈悲才属于我们。\par

13. 此外，在为自己挑选阉人、婢女和仆从时，要更看重他们的品行而非容貌；因为，无论他们的年龄或性别如何，甚至即使阉割保证了他们强制性的贞洁，你仍须考虑他们的性情，因为只有对基督的敬畏才能改变性情。当你面前时，戏谑和轻浮的言谈决无立足之地。你绝不该听到不当之言；若偶然听到，也不可为之所动。堕落之人常利用一句轻浮的话来试探贞洁的门户。把欢笑和被人取笑的特权留给世俗之人。处于你这种地位的人应当端庄持重。从前你城中的显要人物监察官加图（他在晚年转向研究希腊文学，既不因自己身为监察官而感到羞愧，也不因年迈而放弃希望），据卢基利乌斯说，他一生只笑过一次；同样的说法也适用于马库斯·克拉苏。这些人可能是刻意摆出这种严肃的态度，以博取名声和关注。因为只要我们仍居住在此身体的帐幕里，被这脆弱的肉身所包裹，我们只能约束和调节我们的情感和情绪，却不能完全根除它们。圣咏作者深知这一点，说道：\BibleQuote{你们愤怒，却不要犯罪。}宗徒这样解释：\BibleQuote{不要让太阳在你们的愤怒中落下。}\BibleRef{弗 4:26} 因为愤怒是人之常情，而平息愤怒则是基督徒的本分。\par

14. 我认为无需提醒你提防贪财，因为你的家族向来既拥有财富又轻视它。宗徒也告诉我们，贪财就是崇拜偶像，\BibleRef{弗 5:5}当有人问主说：\BibleQuote{善师，我该行什么善，才能得永生？}他回答说：\BibleQuote{你若愿意是成全的，去，变卖你所有的，施舍给穷人，你必有宝藏在天上；然后来跟随我。}这就是圆满的宗徒德行的顶峰——变卖一切所有的，分施给穷人，如此摆脱一切尘世束缚，与基督一同飞升天上。对我们，或者更确切地说，对你，托付了一项谨慎的管家职务，尽管在这类事情上，无论老少，每个人都享有完全的自由选择权。基督的话是：\BibleQuote{你若愿意是成全的。}他似乎说：我不强迫你，我不命令你，但我将棕榈枝放在你面前，我向你展示奖品；由你选择是否进入竞技场赢得桂冠。让我们想想智慧如何明智地发言。\BibleQuote{变卖你所有的。}这命令是给谁的？为什么？是给那人对他说：\BibleQuote{你若愿意是成全的。}不是变卖你的一部分财产，而是\BibleQuote{一切所有的。}变卖以后，然后呢？\BibleQuote{施舍给穷人。}不是给富人，不是给你的亲属，不是用来满足自我放纵，而是救济急需。不论一个人是司铎、亲戚或关系，你只应考虑他的贫困。让饥饿之人的肚子称赞你，而不是饱食者的盛宴。我们在\BibleBook{宗徒}{大事录}中读到，当主的血尚温热，信徒们处在初信的热忱中，他们都变卖了财产，把价银放在宗徒脚前（表明金钱应被践踏），\BibleQuote{按每人的需要分给每人}。\BibleRef{宗 4:34}但\Person{阿纳尼雅}和\Person{撒斐辣}却是胆怯的管家，更是欺骗的管家，因此他们招致了定罪。因为他们许了愿，就把钱献给天主，仿佛钱是他们自己的一样，而不是属于他们所许愿的那位；他们为自己留下属于别人的一部分，害怕饥荒——真正的信德从不害怕饥荒——于是突然招致了惩罚，这惩罚不是对他们残忍，而是作为对他人的警告。\BibleRef{宗 5:1}事实上，宗徒\Person{伯多禄}绝没有像\Person{波斐利}愚蠢所说的那样呼求死亡降临他们。他只是藉着预言的神恩宣布了天主的审判，以使两个人的惩罚成为许多人的教训。自从你献身于终身童贞之日起，你的财产就不再属于你；或者更确切地说，现在才真正属于你，因为已成了基督的财产。然而，只要你的祖母和母亲在世，你必须按照她们的心愿处理。但是，如果她们去世，安息于圣人的睡眠中（我知道她们希望你先她们而去）；当你的年岁更成熟，意志更坚定，决心更坚强时，你将按照你认为最好的——或者更确切地说，按照主的命令——处理你的钱财，因为你将知道，除了你在世上为善工所花费的，你将一无所有。他人可以建造教堂，用大理石装饰墙壁，购置巨大的柱子，用黄金和珍贵饰品装饰无知的柱头，用银包裹教堂门，用黄金和宝石装饰祭台。我不责备做这些事的人；我也不否定他们。每个人都应跟随自己的判断。这样花钱总比囤积和斤斤计较要好。然而，你的责任是另一种。你的责任是在穷人身上穿上基督，在病人中探望他，在饥饿中喂养他，在无家可归者中收留他，尤其是那些有信德的一家之人，\BibleRef{迦 6:10}支持贞女团体，照顾天主的仆人，那些神贫的人，他们日夜服侍同一位主，在地上度天使般的生活，只谈论对天主的赞美。他们有衣有食，就喜乐，自认为富有。他们不追求更多，只要能在其志向中持守就满足。因为一旦他们开始追求更多，就表明他们连那些必要的东西也不配拥有。\par

上述劝告是针对一位富有且地位尊贵的贞女。\par

15. 但现在我下面的话只对童贞女说。我要考虑的，不是你的处境，而是你本人。除了你必须始终遵守的咏唱圣咏和祈祷规则，即第三、第六、第九时辰，晚间、午夜和黎明，你当决定自己承诺多少时间用于学习和阅读圣经，目的在于使灵魂喜悦受教，而非使其负重。当你在这座研习中度过分配的时间后，常依灵魂的关切而跪下祈祷；手边常备羊毛，用拇指捻线成纱，挂于梭上投梭织布，或将他人纺好的线卷绕起来，或为织工备线。成品要检查，挑剔瑕疵，安排应做数量。若你忙于这许多工作，永不会觉得日子漫长；即便夏日日落再晚，若有未完之事，一日仍觉太短。遵此规则，你必能救自己并救他人，作为女主人立下好榜样，且将众人的贞洁归于你的功劳。因为经上说：\Quote{每个闲散者的灵魂充满欲望。}也不得以天主的恩赐使你不缺什么为由逃避劳作。不，你当与众同劳，如此常忙于工作，你便只思念事奉主。我要直言不讳。即使你将所有财产施舍穷人，基督认为珍贵无比的仍是你亲手所做的工作。你可以为自己劳作，或为你的贞女们作榜样；也可将手工赠予你的母亲和祖母，以使她们拿出更多款项救济穷人。\par

16. 我几乎忽略了最重要的一点。当你尚年幼时，蒙福而圣洁的教宗阿纳斯塔修斯治理罗马教会。在他的时代，一场可怕的异端风暴从东方而来，先试图败坏，继而瓦解宗徒所称赞的纯朴信仰。\BibleRef{罗 1:8}然而，这位主教虽贫穷却富足，像宗徒一样关切羊群，立刻击打毒物的头，止住九头蛇的嘶鸣。现在我有理由担心——事实上我已得如此报告——这异端的有毒胚芽至今仍活在有些人心中，并已发芽。因此，我以为当以全部仁慈与关爱告诫你：持守圣依诺森的信仰，他是阿纳斯塔修斯的神子，并在宗座上的继任者；不可接纳任何异端学说，无论你自以为多么智慧明达。这类人在角落窃窃私语，假装探究天主的公义。他们问：为什么某个灵魂生于某省？为何有些人生于基督徒父母，而他人则生于野兽与未识天主的野蛮部落中？他们能在何处以蝎子之针刺伤淳朴之人，形成合其目的的溃疡，就在何处散布毒液。\Quote{你以为这是毫无理由的吗，}——他们如此争辩说——\Quote{一个小孩子几乎不能以笑声或喜色认出母亲，未曾行善或作恶，就被魔鬼附身，或患上黄疸，或注定承受恶人逃过的苦难，而天主的仆人却必须承受这些？}现在，他们说，若天主的审判是\Quote{真实且全然公义}，若\Quote{在祂内毫无不义}，那么理性就强迫我们相信我们的灵魂曾在天上先存，因某些古罪而受罚，甚至，若我可这样说，埋葬在人体内，我们在这涕泣之谷为旧恶受惩罚。他们以此解释先知的话：\Quote{我受难之前曾走入歧途}，以及\Quote{将我的灵魂从监牢中领出。}他们同样解释福音中门徒的问题：\Quote{谁犯了罪，这个人还是他的父母，竟使他生来瞎眼？}\BibleRef{若 9:2}以及其他类似段落。\par

这种不虔不敬的邪恶教义，昔日盛行于埃及和东方，如今却像毒蛇藏在洞穴中，潜伏于那些地区的多人之中，玷污信仰的纯洁，像遗传病般逐渐蔓延，以致袭击众人。但我确信你若听闻，必不接受。因为你在天主之下有女导师，她们的信仰是纯正道理的准则。你会明白我的意思，因为天主必使你在一切事上领悟。你不可立刻要求我驳斥这可怕的异端以及更糟糕的其他异端；因为若我这样做，便成了\Quote{在应禁止之处指责}，而我现在的目的不是驳斥异端者，而是教导一位童贞女。不过，我已挫败他们的诡计，并在一篇靠天主帮助写成的论文中破坏了其颠覆真理的企图；你若想要此文，我必欣然速寄。我说，你若想要此文，因为如谚语所言：不请自来的货物不受珍视，供应充足使价格下降，而物以稀为贵。\par

17. 人们常常讨论独修生活与团体生活的优劣，通常前者优于后者。然而，即使对男人而言，也总有风险：脱离同伴的交往后，他们可能暴露于不洁和不信的想象中，在傲慢和轻蔑的满盈中，看不起除自己以外的任何人，并武装舌头以诋毁神职人员或那些与自己一样发愿独修的人。关于这样的人，圣咏作者说得很好：\BibleQuote{至于人子，他们的牙齿是长矛和箭矢，他们的舌头是利剑。} 如果这一切对男人是真的，那么对待女人更是如此，因为她们轻浮善变的头脑，若任其自便，很快就会堕落。我本人就认识男女\OriginalTerm{独修者}{anchorite}，他们因过度禁食而损害了心智，以至于不知道做什么、往哪里去，何时说话、何时沉默。最常见的是，受此影响的人曾居住在寒冷潮湿的独修小室中。此外，未经世俗学问训练的人阅读教会作家的作品，只能获得辞藻的流利，而不是（他们本可以获得的）对圣经的更全面认识。古老的谚语在他们身上应验了：虽然他们没有说话的智慧，却不能保持沉默。他们教导别人自己也不理解的圣经；如果侥幸说服了别人，就以博学者自居。事实上，他们自己还未曾上学，就以无知者的导师自居。因此，顺从比自己高明的人、服从管理自己的人、不仅从圣经也从别人的榜样学习如何安排自己的生活、不追随最坏的导师——自己的自信，是一件好事。对于这样僭越的女人，宗徒说她们\BibleQuote{为各种教义之风所飘荡，\BibleRef{弗 4:14} 常常学习，却总无法达到认识真理的地步。}\BibleRef{弟后 3:7}\par

18. 要避免与那些醉心于丈夫和世界的已婚妇女交往，免得你因听见丈夫对妻子或妻子对丈夫所说的话而心神不宁。这样的谈话充满致命的毒液。为表达对此的谴责，宗徒引用了一位世俗作家的诗句，将其纳入教会的服务。可以牺牲韵律而直译为：\Quote{恶友败坏善行。} 不，你应选择稳重而严肃的妇女作为伴侣，特别是寡妇和贞女，是那些言谈可嘉、寡言少语、圣洁谦逊的人。要躲避轻佻而轻率的少女，她们装饰头颅、刘海垂发，用化妆品改善肤色，穿紧袖、无褶的连衣裙和精致的半统靴，并假装是贞女，更轻易地卖身于毁灭。此外，女主人的性格和品味往往从其仆人的行为推断。认为那些不自觉地美丽、不修边幅、不在户外袒胸露背、不掀开外衣露出颈项、只用眼睛露在外面、且只用眼睛寻找道路的人，才是有魅力、可爱且合适的伴侣。\par

19. 我对我将要说的话犹豫不决，但正如常发生的，无论我喜欢与否，都必须说出来；并非我在你的情况中有什么可担忧的，因为可能你对此一无所知，甚至从未听说过，而是在给你建议的同时，也可警告他人。贞女应当避免所有卷发、涂麝香的卷发青年，他们如同贞洁的许多瘟疫和祸害；\OriginalTerm{培特罗尼乌斯·阿尔比特}{Petronius Arbiter}的话很适用于他们：\Quote{过多的香水制造出恶臭。} 我不必提及那些因固执地拜访贞女而给他们自己和贞女带来耻辱的人；因为，即使他们没做什么坏事，但比这更大的坏事莫过于令自己暴露于外教人的毁谤和攻击。我不是在此说所有人，而是仅指那些教会本身所责备的、有时被驱逐的、常常遭受主教和司铎谴责的人。因为，事实上，轻率的少女出现在宗教住所几乎比外出走动更危险。生活在团体中、大量聚集在一起的贞女，决不可独自外出，或没有母亲陪同。一只鹰常常从一群鸽子中单挑一只，扑上去撕开它，直到贪吃它的肉和血。生病的羊离群落入狼的嘴里。我认识一些圣洁的贞女，她们在圣日留在家里以避免人群，并拒绝外出，要么必须带着强大的护卫，要么完全避开所有公共场所。\par

大约三十年前，我出版了一篇题为\WorkTitle{论守贞}的论文，在文中我不得不反对某些恶习，并揭露魔鬼的诡计，以教导那位受书的贞女。当时我的言论冒犯了很多人，因为每个人都把我说的话应用到他自己身上，不是欢迎我的劝诫，而是远离我，把我当作他行为的指控者。你或许会问，这样激怒一群抗议者，并藉我的抱怨显示我良心所受的创伤，有什么用呢？是的，我回答说，因为他们虽然过去了，我的书却仍然存留。我也曾给几位贞女和寡妇写过简短的劝勉信，在这些较小的作品中，我汇集了关于这个主题所有可说的话。所以我处于两难境地：要么重复那些看似多余的劝勉，要么略去它们而严重损害这篇论文。真福西彼连留下了一部关于守贞的杰作；许多其他作者，无论希腊还是拉丁，也做了同样的事。事实上，贞女生活在各民族中，尤其在教会中，都受到口与笔的赞扬。然而，大多数写这个主题的人，都是针对那些尚未选择贞洁的人，以及那些需要帮助以便正确选择的人。但我和我所写信的对象已经做出了选择；我们唯一的目标就是坚持到底。因此，既然我们的道路在蝎子和毒蛇之间，在陷阱和祸害之间，让我们手杖在手，腰束带，脚穿鞋\BibleRef{出 12:11}，这样我们才能到达真正约旦河的甜水，进入应许之地，上到天主的殿。那时我们将与先知同唱：\BibleQuote{主，我喜爱你所居的殿宇和你荣耀的处所；}又说：\BibleQuote{有一件事我曾祈求上主，我仍要寻求：就是一生一世住在上主的殿里。}\par

有福的灵魂，有福的贞女，在她心里除了基督的爱之外，容不下别的爱。因为基督本身就是智慧、贞洁、忍耐、正义和一切其他美德。也有福的是那位能回忆起一个男人的面容而毫无叹息遗憾，并且不渴望看见他，因为看见后她可能不愿放弃。然而，有些贞女以恶劣的行为使贞洁的神圣职业，以及那属于天上和天使队伍的光荣蒙羞。对这些人必须直率地说：如果不能自持，就结婚；如果不愿结婚，就自持。还有一件事引人发笑，更确切地说，引人落泪：当主妇们外出时，前面走着比她们穿着更讲究的侍女；这已成为如此普遍的现象，以至于如果你看到两个女人，其中一个不如另一个整洁，你自然会以为那较差的才是主妇。然而这些侍女却是发愿的贞女。还有不少贞女选择僻静的住所，以免在别人眼下，以便能比她们本来能做的更自由地生活。她们洗澡，随心所欲，并尽可能避人耳目。我们看到这些，却容忍了；事实上，如果我们瞥见金银的光泽，我们就准备把它们当作善行。\par

20. 我以开始的方式结束，不满意只给你一次警告。热爱圣经，智慧就会爱你。热爱智慧，智慧就会保护你。尊重智慧，智慧就会拥抱你。让你胸前和耳上的珠宝成为智慧的宝石。让你的舌头只谈论基督，不让任何不圣洁的声音从你唇间发出，让你的言语常再现你的祖母和母亲为你树立的甜蜜榜样。效法她们，因为她们是德行的模范。\par

\SourceRecord{Translated by W.H. Fremantle, G. Lewis and W.G. Martley.；Nicene and Post-Nicene Fathers, Second Series；Vol. 6.；Edited by Philip Schaff and Henry Wace.；Buffalo, NY: Christian Literature Publishing Co.,；1893.；<http://www.newadvent.org/fathers/3001130.htm>.}

% structure: p0001 source=h1 target=section reason=page-title

\section{\WorkTitle{第一三三封信}}

\Emph{致\Person{特斯丰}}\par

\EditorialNote{特斯丰曾写信给\Person{热罗尼莫}，征求他对\Person{佩拉吉乌斯}教导中两个问题的意见：（一）他的寂静主义，（二）他对原罪的否认。\Person{热罗尼莫}现在驳斥了这两种教义，并指出\Person{佩拉吉乌斯}部分从哲学家、部分从异端者那里汲取了它们。他谴责\Person{鲁菲努斯}（五年前已去世），因为他将一本实为毕达哥拉斯派\Person{克西斯图斯}所著的书归于罗马主教\Person{西斯笃}，并将一篇原出于\Person{奥利振}友人\Person{尤西比乌斯}之手的\Person{奥利振}颂词冒充为殉道者\Person{庞斐禄}的作品。不过，\Person{热罗尼莫}在这两项断言中都是错多对少。（参见\WorkTitle{鲁菲努斯著作序言}。）这封信以承诺将来更充分地处理\Person{佩拉吉乌斯}异端而结束，这一承诺后来通过发表\WorkTitle{反佩拉吉乌斯对话录}得以实现。本信日期为公元415年。}\par

1. 你告知我取代旧争论的新论战，你认为自己行事鲁莽，这错了，因为你的行为是出于热心和友谊。早在你的信到达之前，东方有许多人已被骗入一种模仿谦卑的骄傲，并随同魔鬼说：\Quote{我要升到天上，我要高举我的宝座在天主的众星之上；我要与至高者同等。} \BibleRef{依 14:13} 还有比这不求肖似天主、但求与祂平等更大胆的事吗？并且将一切异端者的有毒教义——这些教义反过来又源于哲学家的陈述，尤其是\Person{毕达哥拉斯}和斯多葛学派创始人\Person{芝诺}——压缩成寥寥数语？因为那些希腊人称为\OriginalTerm{激情}{\GreekText{π}\GreekText{αθη}}、我们可以描述为\Quote{激情}的情绪状态，涉及现在或将来，如烦恼与快乐、盼望与恐惧——他们告诉我们，这些都能从我们心中根除；事实上，通过默想德行和不断实践，人身上的一切恶习都可以被连根拔除。他们采取的立场遭到追溯至\Person{亚里士多德}的逍遥学派以及\Person{西塞罗}为其门徒的新学园派的猛烈攻击；这些学派推翻的不是对手的事实——因为他们没有事实——而是对手用以替代事实的幻影和愿望。坚持这样的教义，就是剥夺人的本性，忘记人是由灵魂和肉体构成的，用单纯的愿望取代健康的教导。因为宗徒说：— \Quote{我这个人真不幸啊！谁能救我脱离这死亡的身体呢？} \BibleRef{罗 7:24} 但由于我无法在一封短信中说尽我所有想说的，我将简要触及你必须避免的几点。\Person{维吉尔}写道：—\par

\Quote{凡人如此恐惧与盼望，欢喜与悲伤\newline{}被幽禁在黑暗中，看不见天空。}\par

因为谁能逃脱这些情绪？当我们欢喜时，难道不都拍手，在悲伤临近时退缩？盼望岂不总是激励我们，恐惧岂不使我们惊骇？同样，在其一首\WorkTitle{讽刺诗}中，诗人\Person{贺拉斯}，其言辞如此沉重，写道：\par

\Quote{没有凡人能完全摆脱过错\newline{}过错最少者便是完美之人。}\par

2. 我们的一位作者说得很好：\Quote{哲学家是异端者的始祖。}正是他们用乖戾的教义玷污了教会的无瑕，却不知道论及人性的软弱时说过：\Quote{为什么尘土和灰烬骄傲？} \BibleRef{德 10:9} 同样，宗徒说：\Quote{我发觉在我的肢体内另有一条法律，与我理智的法律交战，并把我掳去。} \BibleRef{罗 7:23} 又说：\Quote{我所愿意的善，我不做；我所不愿意的恶，我倒去做。} 如果\Person{保禄}做他不愿意做的事，那么\Quote{人若愿意就能无罪}这一断言会变成什么？既然志愿已定，当宗徒告诉我们他无力做他所愿意的事时，志如何能得逞？此外，若问他们究竟认为哪些人是无罪的，他们就用一个新的借口来掩盖真理。他们说，他们并不宣称人现今或曾经无罪；他们所坚持的仅仅是，人有可能无罪。真是卓越的教师！他们主张一件事可能发生，而根据他们自己的说法，这件事从未发生过；然而圣经说：\Quote{将来所发生的事，早已在古时发生过。}\par

我不必逐一细数圣人的生平，或者指出在最皎洁的皮肤上的瑕疵和斑点。诚然，我们许多作者不明智地采取了这一方式；然而，几句圣经经文就足以同时驳斥异端者和哲学家。蒙选的器皿说：\Quote{天主把众人都禁锢在不信中，以便怜悯众人。} \BibleRef{罗 11:32} 另一处说：\Quote{众人都犯了罪，亏缺了天主的荣耀。} \BibleRef{罗 3:23} 代言神圣智慧的训道者也自由地宣告说：\Quote{世上没有行善而不犯罪的义人。} \BibleRef{训 7:20} 又说：\Quote{如果祢的百姓得罪了祢，因为没有人不犯罪。} \BibleRef{列上 8:46} 以及\Quote{谁能说：我洁净了我的心？} \BibleRef{箴 20:9} 和\Quote{没有不受玷污的，即使他在世上的生命仅有一天。} \Person{达味}坚持同样的事，他说：\Quote{看，我是在罪恶中成胎的，我母亲在罪孽中怀了我。} 在另一篇圣咏中：\Quote{在祢面前，凡有血肉的，没有一个能称为义人。} 最后这一节他们试图出于敬畏而解释掉，争辩说其意思是，与天主相比，没有人是完美的。然而圣经没有说\Quote{与祢相比，凡有血肉的没有一个能称为义人}，而是说\Quote{在祢面前，凡有血肉的，没有一个能称为义人。} 当它说\Quote{在祢面前}时，意思是那些在人看来圣洁的人，在天主更全面的认识中，绝非圣洁。因为\Quote{人看外表，上主却看人心。} \BibleRef{撒上 16:7} 但如果在察看万有、内心秘密向他敞开的天主面前，没有一个人是义的；那么这些异端者非但没有增加人的尊严，反而明显剥夺了天主的权能。我本可以汇集许多其他含义相同的经文；但如果这样做，我将超出不说是书信、而是整卷书的限度。\par

3. 他们并非以新教义，而是以自夸的背信弃义来欺骗单纯无知的人。然而，他们无法欺骗那些昼夜默思上主法律的经师。让那些声称人若愿意便可\Quote{无罪}，或如希腊人所称的 \OriginalTerm{无罪的}{\GreekText{αναμάρτητος}}，即\Quote{无瑕疵}的人，为他们的首领和同伴脸红吧。因为这样的说法在东方教会听来是难以容忍的，他们实际上只声称人可以是\Quote{无罪}，并不敢宣称他也可以是\Quote{无瑕疵的}。仿佛\Quote{无瑕疵}和\Quote{无罪}有不同的含义；而两者之间的唯一区别，不过是拉丁语需要用两个词来表达希腊语用一个词所表达的内容。若你采纳\Quote{无罪}而拒绝\Quote{无瑕疵}，那么就应谴责宣扬无瑕疵的人。但你不能这样做。你很清楚你在私下里对你的门徒传授什么；你嘴上说一套，心里却刻着另一套。对我们这些无知的局外人，你用比喻说话；但对你自己的追随者，你透露你的秘密含义。为此你引用圣经的权威说：\Quote{耶稣用比喻对群众说话；但对他的门徒，祂说：\InnerQuote{天国的奥秘是给你们知道的，但不是给他们知道的。}}\BibleRef{玛 13:10-11}\par

但言归正传；我将简要列出你们的首领和同伴的名字，以显示你们夸耀与之交往的是些什么人。摩尼论到他的\Quote{选民}——他将他们置于柏拉图的诸天轨道中——说他们完全无罪，甚至即使愿意也不能犯罪。他们在德行上已达到如此高度，以至于嘲笑肉身的行为。然后是西班牙的普里西利安，其恶名使他与摩尼一样坏，他的门徒们对你们推崇备至。这些人竟敢为自己声称兼备完美与智慧的双重荣誉。然而，他们却独自与女人关在一起，并用以下的诗句来为其罪恶的拥抱辩护：\par

全能的天父迎娶大地为妻；\newline{}将滋润的雨水倾注于她，\newline{}好从她腹中收获新的庄稼。\par

这些异端与诺斯替主义有关联，这可以追溯到巴西里德的邪恶教导。正是从他那裡你们得出断言：不熟悉法律就不可能避免罪恶。但我为何要谈论已被全世界谴责并被世俗之剑处死的普里西利安呢？本都伊贝拉的埃瓦格里乌斯写信给童贞女和隐修士，其中包括那位其名字证实其不忠之黑暗的女士，他出版了一本关于 \OriginalTerm{无感}{apatheia} 的格言集，即我们所说的\Quote{不动心}或\Quote{平静}；一种心灵不再动荡的状态——简单地说——变成石头或神。他的著作广为流传，东方有希腊文版，西方有其门徒鲁菲努斯翻译的拉丁文版。他还写了一本自称关于隐修士的书，其中包含许多并非隐修士的人，他宣称他们是奥利金主义者，并且确实已被主教们谴责。我指的是阿摩尼乌斯、优西比乌、欧提米乌斯、埃瓦格里乌斯本人、霍鲁斯、伊西多尔，以及许多其他难以一一列举的人。然而，他小心地效仿医生，如同卢克莱修所说：\par

他们给孩子们喂苦艾，\newline{}却在杯中抹上最甜的蜜汁。\par

也就是说，他把一位无可置疑的天主教的圣人若望放在书的前面，以此将后面提到的异端引入教会。但谁能充分描摹那种鲁莽或疯狂，导致他将毕达哥拉斯派哲学家、对基督一无所知的外教人克斯图斯的一本书，归于罗马教会的殉道者和主教息斯都？在这部作品中，广泛讨论了完美的问题，依据毕达哥拉斯学说，该学说使人等同于天主并与祂同一实体。因此，许多不知道其作者是哲学家、以为自己在读殉道者的言论的人，饮了巴比伦的金杯。此外，书中没有提及先知、圣祖、宗徒或基督；因此，根据鲁菲努斯，竟有一位主教和殉道者与基督毫不相干。你们和你们的追随者就是从那本书中引用反对教会的话。同样地，他随意玩弄神圣殉道者庞菲洛的名字，将凯撒勒雅的优西比乌所著、公认是亚略派人士的《为奥利金辩护》六卷中的第一卷归于他。他这样做当然是为了向拉丁语读者推荐奥利金那四本关于《原则》的美妙著作。\par

你要我提及你们另一位异端导师吗？你们的许多教导都可追溯到奥利金。例如，仅举一例，当他注释圣咏作者的话：\Quote{夜间，我的五内也教训我}\BibleRef{咏 16:7} 时，他坚持认为，当像你这样的一位圣人达到完善时，即使在夜间也免于人性的软弱，不受恶念的诱惑。你们不必羞于承认自己是这些人的追随者；否认他们的名字而无用，既然你们采纳了他们的亵渎。此外，你们的教导与约维尼安的第二个立场一致。因此，你必须把我对他的回答视为同样适用于你。凡意见相同的人，他们的命运也难以不同。\par

4. 情形既是如此，那么\Quote{\InnerQuote{愚蠢的妇人，被罪恶压制，被各样的教训之风摇荡，常常学习，终久不能明白真道}}有什么用处呢？那些趋炎附势、耳朵发痒\BibleRef{弟后 4:3}、既不会听也不会说的人，又有什么帮助呢？他们将旧泥与新灰泥混合，如厄则克耳所说，用未调和的灰泥涂抹墙壁，以致当真理如雨水降临时，一切都被冲毁。\BibleRef{则 13:10} 西满术士正是借助娼妓海伦创立了他的教派。安提约基雅的尼阁老，那个一切污秽的发明者，身边总跟着一群妇女。玛尔基翁派遣一个女人先他前往罗马，为使人落入他的圈套。阿培肋有一个斐洛美纳作为他谬误教义的同伴。蒙达诺，那不洁之灵的喉舌，利用两位富有而高贵的女子普里斯卡和玛克西米拉，先是贿赂，然后败坏了许多教会。暂且不论古代历史，让我转向较近的时代。亚略一心要迷惑世界，先从迷惑皇帝的姐妹开始。多纳图借助路基拉的财富，在非洲以他那污秽的圣洗使许多不幸的人受玷污。在西班牙，盲女阿加贝将盲人厄尔比迪乌领入了沟渠。普里西利安紧随其后，他原是一个狂热的琐罗亚斯德崇拜者和魔术师，后来成为主教。一个名叫加拉的女人助长了他的势力，并留下一个游荡的姐妹继续传播另一种类似的异端。如今，不法的奥秘正在运行。\BibleRef{得后 2:7} 男人和女人互相设下陷阱，以致我们不禁想起先知的话：\Quote{鹧鸪高声叫唤，它聚集了不是自己所生的幼雏，它得财富不依正道；在它日子过半时，它要离弃它们，在它的结局，它将成为一个愚拙的人。}\par

5. 为了更好地欺骗人，他们在上述格言之外加上了保留条件：\Quote{\InnerQuote{但不是没有天主的恩宠}}；乍一看，这可能会蒙蔽一些读者。然而，当仔细筛选并加以思考时，它就不能欺骗任何人了。因为他们虽然承认天主的恩宠，却告诉我们我们的行为不依赖于祂的帮助。相反，他们将天主的恩宠理解为自由意志和法律的诫命。他们引用依撒意亚的话：\Quote{天主赐下法律以帮助人}，并说我们应当感谢祂，因为祂把我们造得能够凭自由意志择善避恶。他们没有看出，在这样主张时，魔鬼正用他们的嘴唇嘶出无法容忍的亵渎。因为如果天主的恩宠仅限于祂赋予我们自己的意志，而我们满足于自由意志，不寻求神圣的帮助，以免损害这自由，那么我们就应停止祈祷；因为我们不能祈求天主的仁慈每天赐给我们那已经一次赐予我们、已经在手的恩赐。这样想的人使祈祷成为不可能，并夸耀自由意志不但使他们成为自己的主宰，而且如同天主一样有权能。因为他们不需要任何外在的帮助。禁食吧，放弃一切形式的克己吧！因为我何必费尽劳苦去争取那已一次放在我手边的东西呢？我用的论点不是我的，而是白拉奇的一个门徒提出的，或者更确切地说，是他的整个军队的教师和指挥官。此人，与保禄相反，因为他是一个丧亡之器，在荆棘丛中徘徊——并非如他的党羽所说，是诡辩的荆棘，而是文法错误的荆棘——并如此推理：\Quote{如果我没有天主的帮助什么也不做，我的一切行为都是祂的行动，那么我就不再劳苦，我所赢得的冠冕将不属于我，而属于天主的恩宠。祂赐予我抉择的能力，若没有祂不断的帮助我就不能使用它，那是徒然的。因为需要外在支持的意志就不再是意志。天主赐予我自由，但如果我不能随心所欲，这自由还有什么意义？}因此他提出如下的两难之说：\Quote{要么我一次使用所赐予我的能力，从而保持我意志的自由；要么我需要他人的帮助，这样我的意志自由就完全被废除了。}\par

6. 说这话的人确实不是普通的亵渎者；他的异端毒害也不是普通的毒害。他们论说：既然我们的意志是自由的，我们就再不需要依赖天主；他们忘记了宗徒的话：\Quote{你有什么不是你领受的呢？既然是领受的，为什么你还夸耀，好像不是领受的呢？}\BibleRef{格前 4:7}人为了维护自己意志的自由而反抗天主，这真是对天主绝妙的回报！至于我们，我们欣然拥抱这自由，却从不忘记感谢施予者，因为我们知道，除非祂不断在我们内保存祂自己的恩赐，我们便是软弱的。正如宗徒所说：\Quote{所以，不在乎那愿意的，也不在乎那奔跑的，只在乎显示慈悲的天主。}\BibleRef{罗 9:16}愿意和奔跑是我的事，但除非天主不断赐予我祂的帮助，它们就不再是我的事。因为同一宗徒说：\Quote{是天主在你们内运作，使你们愿意，并使你们实行。}\BibleRef{斐 2:13}福音中救主说：\Quote{我父一直工作，我也一直工作。}\BibleRef{若 5:17}祂总是施予者，总是赏赐者。祂一次赐给了我恩宠并不足够；祂必须不断赐给我。我寻求是为了获得，获得之后我又寻求。我贪求天主的慷慨；正如祂从不怠于施予，我也从不倦于领受。我喝得越多，就越渴。因为我读了圣咏作者的话：\Quote{请你们体验并观看，上主是何等的甘饴。}我们拥有的一切美好事物，都是对主的体验。当我自以为已经完成了德行的书卷时，我才刚刚开始。因为\Quote{敬畏上主是智慧的开始}，而这敬畏又被爱所驱逐。\BibleRef{若一 4:18}人只有在认识到自己不完全时才是完全的。基督说：\Quote{同样，你们做完了一切所吩咐的，要这样说：\InnerQuote{我们是无用的仆人，我们只做了我们应做的事。}}\BibleRef{路 17:10}如果那做完一切的人尚且是无用的，那么那没有做完的又该怎么说呢？这就是为什么宗徒声明他已经部分获得了、部分领悟了，他还不完全，并且忘记背后，向着前面奔跑。\BibleRef{斐 3:12}那总是忘记过去、渴望未来的人，表明他对现状并不满足。\par

他们一直反对我们，说我们摧毁了自由意志。不，我们回答说，是你们摧毁了它；因为你们滥用它，并否认其施予者的慷慨。到底是谁摧毁了自由？是那常常感谢天主、将自己涓滴的源流追溯到祂的人？还是那说：\Quote{不要靠近我，因为我是圣洁的；我不需要你。你已经一次赐给了我按自己意愿行事的自由选择。那么你为什么还要干涉，阻止我行事，除非你首先使你的恩赐在我内有效？}对这样的人，我要说：\Quote{你对天主恩宠的信仰表白是虚伪的。因为你把恩宠仅仅解释为人的受造状态，并不指望天主在人的行动中帮助他。依你的说法，这样就会放弃人的自由。因此，你蔑视天主的援助，不得不转而求助于人。}\par

7. 听着，只要听听这个亵渎者的话。他断言：\Quote{假设我想弯曲手指，或移动手，或坐、立、行、跑，或吐痰、擤鼻涕，或行自然之事，难道天主的帮助对我总是不可或缺的吗？}忘恩负义、不，亵渎的可怜人，听听宗徒的宣告：\Quote{所以，你们或吃或喝，或无论做什么，一切都要为光荣天主而做。}\BibleRef{格前 10:31}也听听雅各伯的话：\Quote{注意！你们说：今天或明天，我们往某城去，在那里住一年，作买卖获利。其实你们不知道明天会发生什么；你们的生命是什么？只是片刻出现，随后就消失的雾气。你们倒应当说：\InnerQuote{如果主愿意，我们就可以活着，做这事或那事。}}\BibleRef{雅 4:13}你以为，如果你被迫以天主为你一切行动的动因，被迫依赖祂的旨意，被迫和圣咏作者一样说：\Quote{我的眼目常仰望上主，因为祂必把我的脚从网中拉出}，那么你就受到了伤害，你的自由选择就被摧毁了。于是你鲁莽地断言，每个人都受自己的选择支配。但如果他受自己的选择支配，那么天主的帮助在哪里呢？如果他不需要基督管理他，为什么耶肋米亚写道：\Quote{人的道路不由自己}\BibleRef{耶 10:23}，以及\Quote{上主引导他的脚步}？\BibleRef{箴 16:9}\par

你说天主的诫命是容易的，然而你却找不出一个完全遵守了它们的人。回答我这个问题：它们是容易还是困难？如果容易，那么请找出一个完全遵守了它们的人。也解释圣咏作者的话：\Quote{你以法律使我劳苦}，以及\Quote{因了你嘴唇的话，我遵守了艰难的道路。}并且解释主在福音中的话：\Quote{你们要从窄门进去}；\BibleRef{玛 7:13}以及\Quote{爱你们的仇敌}；\Quote{为那迫害你们的祈祷。}\BibleRef{玛 5:44}另一方面，如果诫命是困难的，并且没有人完全遵守了它们，你怎么敢说它们是容易的呢？你没看到你自相矛盾吗？因为要么它们容易，那么无数人已经遵守了；要么它们困难，那么你称它们容易就是过于仓促了。\par

你们一派常以天主的诫命或可行或不可行为论据。凡可行的，你们承认天主加给我们是正确的；凡不可行的，你们便声称罪责不在领受诫命的我们，而在加给诫命的天主。怎么！难道天主命令我成为祂那样，与造物主毫无分别，比最高的天使还大，拥有天使都没有的能力？\OriginalTerm{无罪}{sinlessness}被视为基督的特色，正如经上记载：\BibleQuote{他没有犯过罪，他口中也找不出诡诈。}\BibleRef{伯前 2:22}但如果我和祂一样无罪，无罪又如何再是祂的特色？因为若真有这种区别，你们的理论便自相矛盾了。\par

你们断言，人若愿意就可以无罪；然后，仿佛从沉睡中醒来，你们又试图欺骗不警惕的人，加上这个保留条件：\Quote{但没有天主的恩宠却不行。}因为如果人凭自己的力量就能保持无罪，他何必需要天主的恩宠？反之，如果离了恩宠他就一事无成，说什么他能做他不能做的事，又有什么用处？有人争辩说，人只要愿意就能无罪且成全。哪个基督徒不愿无罪？又有谁拒绝成全，如其如你们所说，愿意就能得到，而且意愿必定伴随能力？没有一个基督徒不愿无罪；因此，他们既然愿意，就都会无罪。不管你们喜欢与否，你们必然陷入两难：你们找不出一个（或几乎找不出）无罪的人，却又不得不承认人人都可以愿意无罪就无罪。有人争辩说，天主的诫命是可能遵守的。谁否认呢？但如何理解这一真理，蒙选之器皿清楚解释：\BibleQuote{法律因了肉性的软弱所不能行的，天主却做到了：祂派遣了自己的儿子，带着罪恶肉身的形状，当作赎罪祭，就在肉身中定罪了罪恶；}\BibleRef{罗 8:3}又说：\BibleQuote{凡有血气的，没有一个因行法律而成义。}\BibleRef{罗 3:20}为了表明所指的不只是梅瑟法律或所有总称为法律的规条，同一位宗徒写道：\BibleQuote{按我的内心，我是喜悦天主的法律；但我发觉在我的肢体内，另有一条法律，与我理智的法律交战，并把我掳去，叫我隶属于那在我肢体内的罪恶的法律。我这个人真不幸呀！谁能救我脱离这该死的肉身呢？感谢天主，借着我们的主耶稣基督。}他的其他话进一步解释他的意思：\BibleQuote{我们知道法律是属神的，但我是属血肉的，已被卖给罪恶。因为我所做的，我自己不明白；我所愿意的，我偏不做；我所憎恨的，我反而去做。如果我所不愿意的，我反而去做，我便同意法律是美好的。那么，现在做那事的已不是我，而是住在我内的罪恶。我知道，在我内，即在我的肉性内，没有善；因为我有心行善，但实际行善的能力却找不到。我所愿意的善，我不去行；我所不愿意的恶，我反而去做。如果我所不愿意的，我反而去做，那么做那事的已不是我，而是住在我内的罪恶。}\BibleRef{罗 7:14}\par

9. 但你将反对这一点，并说我追随\OriginalTerm{摩尼派}{Manichæans}和其他反对教会教义之人的教导，因为他们相信有两种互不相同的本性，并且有一种邪恶的本性，绝不能改变。但你这指责不应针对我，而应针对那位宗徒，他深知天主是一回事、人是另一回事，肉情是软弱的、心神是坚强的。\Quote{因为肉情和圣神相争，圣神和肉情相争，这两个彼此敌对，致使你们不能做你们所愿意的事。}\BibleRef{迦 5:17} 但你永远不会从我口中听到任何本性本质上是邪恶的。那么，让我们从告诉我们肉情软弱的这位那里学习。问他为何说：\Quote{我所愿意的善，我不去做；我所不愿意的恶，我却去做。}\BibleRef{罗 7:19} 什么必然性束缚了他的意志？什么强迫命令他做他所厌恶的事？为什么他必须不做他所愿意的，而做他所厌恶和不愿意的？他将这样回答你：\Quote{人哪，你是谁，竟敢向天主抗辩？受造物岂能对造物主说：你为什么这样造了我？难道陶工不能随意用一团泥，做这个器皿为尊贵，做那个器皿为卑贱吗？}\BibleRef{罗 9:20} 再向天主提出更严重的指控，问祂为什么当厄撒乌和雅各伯还在母胎中时，祂说：\Quote{我爱了雅各伯，却恨了厄撒乌。} 控告祂不义，因为当加尔米的儿子阿干偷窃了耶里哥战利品的一部分时，祂因一人的过犯屠杀了成千上万的人。\EditorialNote{若苏厄书第七章} 问祂为什么因厄里儿子的罪过，百姓几乎被消灭，约柜被掳走。\EditorialNote{撒慕尔纪上第四章} 又为什么，当达味因统计百姓而犯罪时，成千上万的人丧命。\EditorialNote{撒慕尔纪下第二十四章} 或者最后，你拿起你的同伙波菲利最喜欢的挑剔，问天主怎能被描述为慈悲和宽仁的，当从亚当到梅瑟、从梅瑟到基督来临，祂让所有民族在不知晓法律和祂的诫命中死去。因为不列颠，那个充满暴君的省份，苏格兰部落，以及直到大洋的所有蛮族，同样都不认识梅瑟和先知们。为什么基督的来临要延迟到最后时刻？为什么祂不在如此众多的人丧亡之前来临？对这最后的问题，蒙福的宗徒在写给罗马人的信中极其明智地处理了，承认他不知道，只有天主知道。那么，你也要谦卑地对你所探究的保持无知。让天主保留对自己所有之物的权能；祂不需要你为祂的行为辩护。我是那个不幸的人，你应当把辱骂指向我，我常常读到这些话：\Quote{由于恩宠，你们得救了}\BibleRef{弗 2:5} 和\Quote{罪恶得赦免、过犯得遮盖的人是有福的。} 然而，为了暴露我自己的软弱，我知道我渴望做许多我应当做却做不到的事。因为我的心神是坚强的，引领我走向生命，但我的肉情是软弱的，将我拖向死亡。我耳中常听到主的警告：\Quote{你们醒寤祈祷罢！免陷于诱惑；心神固然切愿，但肉体却软弱。}\BibleRef{玛 26:41}\par

10. 你歪曲我的意思，并试图让无知者相信我在谴责自由意志，这是徒劳的。愿谴责自由意志的人自己受谴责。我们被造时赋有自由意志；尽管如此，并非这使我们有别于禽兽。因为人的自由意志，如我以前所说，依赖于天主的助佑，并时刻需要祂的帮助——这是你和你们的人不愿承认的。你们的立场是：若一个人一旦有了自由意志，他就不再需要天主的帮助。确实，自由意志带来决定的自由。然而，人不是立即凭其自由意志行事，而是需要天主的援助，天主自己却不需要援助。你自己夸口说，人的义德可以是完美的，与天主的义德相等；然而你却承认自己是个罪人。那么，回答我：你愿意，还是不愿意摆脱罪恶？若你愿意，为什么按你的原则你不实现你的愿望？若你不愿意，难道你不是证明自己轻视天主的诫命？你若轻视，那你就是罪人。若你是罪人，那么经上说：\Quote{天主对恶人说：你怎么胆敢传述我的诫命？你的口怎敢提及我的盟约？你原是恼恨教训的，将我的话置诸脑后。}\BibleRef{咏 50:16-17} 只要你不愿做天主所命令的，你就将祂的话置诸脑后。然而，你像一位新的宗徒，为世界规定该做什么和不该做什么。但你的言行完全不一致。因为当你说自己是罪人——却又说人若愿意可以无罪——你是希望人们理解你是圣人，摆脱一切罪恶。你称自己为罪人只是出于谦卑；为了让你有机会在贬低自己的同时赞扬他人。\par

11. 你的另一个论点也是不能容忍的，它这样说：\Quote{无罪是一回事，能够无罪是另一回事。前者不在我们能力之内，后者通常是在的。因为虽然从未有人是无罪的，但是，如果一个人愿意无罪，他就能无罪。} 请问，这是什么推理？一个人能成为从来没有人成为过的样子！一件事情是可能的，而按照你自己的承认，还没有人做到！你是在对人赋予一种品质，就你所知，他可能永远不会拥有！而且你把一种恩宠归于任何一个人，而你却不能显示这恩宠曾标记圣祖、先知或宗徒。请听教会的言语，尽管在你看来可能平淡、粗鲁或无知。说出你的想法；公开宣讲你私下告诉门徒的。你自称有自由选择，为何不自由说出你的思想？你的密室听到一种教义，平台周围的人群听到另一种。我猜想，未受教育的群众不能消化你的深奥教导。他们满足于婴儿的奶食，不能吃固体食物。\par

我尚未写任何东西，而你仍用答复的雷霆威胁我；我猜想，你希望我被你的恐吓吓住，不敢开口。你看不到我写作的目的是迫使你回答，并明确承诺那些你现在根据时间、地点和人物而坚持或忽略的教义。有一种自由我必须否认你，就是否认你曾经写过的东西的自由。你公开承认你所持有的意见，将是教会的胜利。因为要么你答复的言辞与我的相合，那样我就不再把你们视为反对者，而是朋友；要么你将反驳我的教义，那样你的意见被所有教会知晓，对我就是胜利。将你的信条带到光中就是克服它们。亵渎写在他们脸上，一个在其陈述本身即为亵渎的教义不需要驳斥。你以答复威胁我，但没有人能逃脱这一点，除非根本不写的人。你怎么知道我要说什么而谈答复？也许我会接受你的观点，那样你就白费了心机。欧诺米派、亚略派、马其顿尼派——所有这些，名称不同，不虔敬相同，但我不在乎。因为他们说出他们所想的。唯你的异端公开羞于坚持私下却敢于教导的。但门徒的疯狂暴露了师傅的沉默；因为他们从他们那里在密室听到的，他们在房顶上宣讲。如果听众喜欢他们所说的，师傅得到称赞；如果讨厌，只有门徒受责备，师傅得以脱身。这样，你的异端成长了，你欺骗了许多人；特别是那些依恋妇女且确信自己不能犯罪的人。你总是教导，你总是否认；你应得先知的话用在你身上：\Quote{上主，当他们在产痛和劳苦中时，求你把荣耀赐给他们。上主，求你将什么赐给他们？你赐给他们什么呢？求你赐给他们流产的子宫和干瘪的乳房。} 我的怒气上升，我不能控制我的言语。一封信的篇幅不容许冗长的讨论。我在这里没有指名攻击任何人。我只是针对歪曲教义的教师说话。如果怨恨促使他回复，他只会像老鼠一样暴露自己的踪迹，而当真正动手时，将受到更重的伤。\par

12. 从我年轻时直到现在，我花费多年写作各种著作，并始终努力教导听众我在教会中公开领受的教义。我没有追随哲学家的讨论，而是宁愿顺从宗徒的朴素言语。因为我知道经上记载：\BibleQuote{我要摧毁智慧人的智慧，废除明达人的明达}，以及\BibleQuote{天主的愚笨也比人明智。}\BibleRef{格前 1:25} 既然如此，我挑战我的反对者彻底审查我过去所有的著作，如果他们在其中能找到任何错误，就把它暴露在光中。两件事之一必然发生：要么我的作品被证明是有益的，我将驳斥对我的虚假指控；要么它们被证明是有过错的，我将承认我的错误。因为我宁愿纠正一个错误，也不愿坚持一个被证明是错误的意见。至于你，尊贵的师傅，你也照样去做：要么为你所发表的言论辩护，用相应的雄辩支持你的巧妙理论，不要一时兴起就否认你自己的话；要么，如果作为一个常人你可能犯了错误，就坦白承认，并恢复存在分歧之处的和谐。请回忆起连士兵也没有撕裂救主的衣服。\BibleRef{若 19:23} 当你看见弟兄们争吵时，你发笑；并且乐于有些人以你的名字称呼，另一些人以基督的名字称呼。更好的是效法约纳说：\BibleQuote{如果这狂风是因我的缘故临到你们，请把我举起来，抛入海中。}\BibleRef{纳 1:12} 他因谦卑被抛入深渊，为要光荣地复活，成为主的预象。\BibleRef{玛 12:39} 但你因骄傲被高举到星辰，只是为要你也听见耶稣说：\BibleQuote{我看见撒殚像闪电一样从天坠落。}\BibleRef{路 10:18}\par

13. 诚然，在圣经中许多被称为义人，如\Person{匝加利亚}和\Person{依撒伯尔}、\Person{约伯}、\Person{约沙法特}、\Person{约史雅}以及圣经中提到的许多其他人。关于这一事实，如果天主赐我恩宠，我将在已承诺的著作中充分解释；在这封信中，只需说明他们被称为义人，并非因为他们无罪，而是因为他们的过错被德行所掩盖。事实上，\Person{匝加利亚}因哑巴受罚\BibleRef{路 1:20}，\Person{约伯}被自己的口定罪\BibleRef{约 42:6}，而\Person{约沙法特}和\Person{约史雅}虽毫无疑义地被描述为义人，却记载他们做了主所不悦的事。前者与不虔敬的\Person{阿哈布}结盟，招致\Person{米加雅}的责备\BibleRef{列上 22:19}；后者——虽然主藉\Person{耶肋米亚}发言禁止他——仍去攻打\Place{埃及}王\Person{法郎乃苛}，并被他所杀\BibleRef{编下 35:20}。然而，他们两者都称为义人。其余的人，现在不是写作的时候；因为你向我要求的不是一篇论文，而是一封信。要彻底驳斥，我需要闲暇，然后希望藉基督的帮助摧毁他们所有的吹毛求疵。为此，我将依赖圣经，天主在其中每日向信者说话。这也是我要藉你向你圣洁而显赫的家中所有聚集之人提出的警告：他们不应让一个或至多三个小人物用如此众多异端的残渣和至少与之相关的恶名玷污自己。一个曾以德行和圣洁闻名的地方，决不能被魔鬼的狂妄和不洁的交往所玷污。让那些向这种人提供金钱的人知道，他们正在壮大异端的队伍，兴起基督的敌人，供养他公开的反对者。他们口里说一套，而行动证明他们想的是另一套，这是徒劳的。\par

\SourceRecord{Translated by W.H. Fremantle, G. Lewis and W.G. Martley.；Nicene and Post-Nicene Fathers, Second Series；Vol. 6.；Edited by Philip Schaff and Henry Wace.；Buffalo, NY: Christian Literature Publishing Co.,；1893.；<http://www.newadvent.org/fathers/3001133.htm>.}

% structure: p0001 source=h1 target=section reason=page-title

\section{第一三五封信}

\Emph{自\Person{教宗依诺增爵}致\Person{奥雷留斯}}\par

\EditorialNote{在迪奥斯波利斯会议后不久，伯拉纠派因成功而得意忘形，袭击了\Person{耶柔米}在白冷的隐修院，将其洗劫并部分焚毁。这为他赢得了\Person{依诺增爵}的同情，后者现在（公元417年）请求\Person{奥雷留斯}将以下信件转交给他。}\par

\Person{依诺增爵}致他最敬爱的朋友与兄弟\Person{奥雷留斯}。\par

我们的同僚司铎\Person{耶柔米}告知我们，您最恳切的愿望是来探望我们。我们与他一同受苦，如同他是我群中的一员。我们也迅速采取了我们认为合乎时宜且可行的措施。由于您视自己为我们中的一员，最亲爱的兄弟，请速将以下信件转交上述的\Person{耶柔米}。\par

\SourceRecord{Translated by W.H. Fremantle, G. Lewis and W.G. Martley.；Nicene and Post-Nicene Fathers, Second Series；Vol. 6.；Edited by Philip Schaff and Henry Wace.；Buffalo, NY: Christian Literature Publishing Co.,；1893.；<http://www.newadvent.org/fathers/3001135.htm>.}

% structure: p0001 source=h1 target=section reason=page-title

\section{第一三六封信}

\Person{教宗依诺森}致\Person{热罗尼莫}\par

\EditorialNote{依诺森对热罗尼莫表达同情，并承诺如果热罗尼莫能提出明确指控，他将采取严厉措施惩罚其反对者。该信的日期为公元417年。}\par

宗徒\BibleRef{铎 3:10}作证说，争论在教会中从未带来好处；因此他指示，应在其异端初起时劝诫异端者一两次，而不应长期反复责备。哪里疏于遵守这一规则，所要防范的祸患非但不能避免，反而更加剧。\par

你的悲痛与哀叹如此触动我们，以致我们既不能行动也不能进谏。\par

然而，首先，我们称赞你信德之恒毅。引用你多次在众人耳中说过的话：\Quote{人若盼望来世的真福，必会为真理欣然面对毁谤甚至个人危险。}我们提醒你你自己曾宣讲过的，虽然我们确信你并不需要提醒。这些可怕祸患的景象彻底激怒了我们，以致我们急忙运用宗座的权威来抑制这瘟疫的一切表现；但是，由于你的信件未指名道姓，也未提出具体指控，目前我们无法对任何人采取行动。但我们已尽力而为；我们深表同情。并且，如果你愿意对任何特定的人提出明确而清楚的指控，我们将任命合适的法官审理他们的案件；或者，你，我们极为尊敬的儿子，认为我们需要采取更严重而紧急的行动，我们也决不会迟缓。同时，我们已写信给我们兄弟\Person{若望主教}，劝告他行事更为慎重，以免在他受委托的教会中发生任何他应当预见并防止的事，也不发生任何日后可能给他带来麻烦的事。\par

\SourceRecord{Translated by W.H. Fremantle, G. Lewis and W.G. Martley.；Nicene and Post-Nicene Fathers, Second Series；Vol. 6.；Edited by Philip Schaff and Henry Wace.；Buffalo, NY: Christian Literature Publishing Co.,；1893.；<http://www.newadvent.org/fathers/3001136.htm>.}

% structure: p0001 source=h1 target=section reason=page-title

\section{第一三七封书信}

\Emph{教宗\Person{依诺增爵}致耶路撒冷主教\Person{若望}}\par

\EditorialNote{依诺增爵谴责若望允许白冷（Bethlehem）的骚乱（如前面两封信所述）由白拉奇派（Pelagians）煽动起来，并劝勉他日后对其教区更警醒。此信日期为公元417年。这一年若望和依诺增爵均去世，很可能若望从未收到此信。}\par

\Person{依诺增爵}致其最敬爱的兄弟\Person{若望}。\par

童贞女\Person{欧斯托基}和\Person{保辣}向我哀诉，她们说魔鬼在属于她们教会的地域内施行了各种蹂躏、谋杀、纵火和暴行；因为她们以极妙的仁慈和慷慨，隐去了其人间代理人的名字和动机。如今，虽然对谁是犯罪者无疑义，但我的兄弟，你本应预先防范，更加留意你的羊群，以免发生这类骚乱；因为他人因你的疏忽而受苦，而你的疏忽鼓励人摧残主的羊群，直到祂柔弱的羔羊被火、剑和迫害剪毛并削弱，她们的亲属被杀害、死亡，据我们所知，她们自己也几乎丧命。难道这没有触及你作为司铎的神圣责任，即魔鬼在你和你的人面前显得如此强大？我说在你面前，因为如此可怕的罪行竟然发生在你的教会管辖范围内，这确实说明你的司铎能力不足。你的防范措施在哪里？打击之后，你的救济努力在哪里？你的安慰话语又在哪里？这些女士告诉我，迄今为止，她们过于恐惧，不敢抱怨已经遭受的苦楚。倘若她们对我更自由地谈及此事，我会对这件事作出更严厉的判断。因此，兄弟，要警惕古仇的诡计，以善牧的精神警醒，或改正或压制这类恶行。因为它们是以传闻而非具体指控的形式传到我的耳中。如果无所作为，关于伤害的教会法可能迫使未能保护其羊群的人说明其疏忽的原因。\par

\SourceRecord{Translated by W.H. Fremantle, G. Lewis and W.G. Martley.；Nicene and Post-Nicene Fathers, Second Series；Vol. 6.；Edited by Philip Schaff and Henry Wace.；Buffalo, NY: Christian Literature Publishing Co.,；1893.；<http://www.newadvent.org/fathers/3001137.htm>.}

% structure: p0001 source=h1 target=section reason=page-title

\section{\WorkTitle{Letter 138}}

\Emph{To \Person{Riparius}}\par

\EditorialNote{热罗尼莫赞扬利帕里乌斯为维护天主教信仰的热忱以及为打击白拉奇派所做的努力。他接着描述了这些异端分子对白冷隐修院的攻击。现在他很高兴地说，他们终于被驱逐出巴勒斯坦。但大部分人，也就是说，有些仍滞留在约帕，包括他们的一位主要领袖。日期为公元417年。}\par

你自己的来信以及许多人的报告都已告诉我，你在为\Person{基督}作战，对抗天主教信仰的敌人。但我听说，你遇到了逆风，而那些本应成为世界战士的人，却为了彼此的毁灭而支持了丧亡的事业。你要知道，在世界的这一部分，没有借助任何人力，仅凭\Person{基督}的命令，\Person{卡提里纳}不仅被赶出京城，而且已被逐出\Place{巴勒斯坦}的边界。然而，\Person{伦图路斯}以及他的许多同谋者仍存留在\Place{约帕}，这是我们感到悲伤的。我自己认为，改变住处比放弃真信仰更好；我宁愿离开我舒适的家，也不愿被异端玷污。因为我不能与那些要么坚持要我立即服从，要么就要以刀剑来支持他们意见的人交往。我敢说，许多人已经告诉了你我所受的苦难，以及\Person{基督}高举右手为我向敌人复仇的事。因此，我请求你完成你所承担的任务，并且在你身在其中之时，不要让\Person{基督}的教会没有防卫者。人人都知道在这场战争中必须使用的武器；我相信你也不会要求别的。你必须全力与敌人抗争；但必须不是用武力，而是用那永不失败的爱德。与我同在的尊敬弟兄们——虽然我不配——热忱地问候你。执事\Person{阿冷提乌斯}尊兄一定会给你——我可敬的朋友——忠实地叙述所有事实。愿我们的主\Person{基督}以祂全能的大能护佑你，使你平安且记念我。真正可敬的先生和尊敬的弟兄。\par

\SourceRecord{Translated by W.H. Fremantle, G. Lewis and W.G. Martley.；Nicene and Post-Nicene Fathers, Second Series；Vol. 6.；Edited by Philip Schaff and Henry Wace.；Buffalo, NY: Christian Literature Publishing Co.,；1893.；<http://www.newadvent.org/fathers/3001138.htm>.}

% structure: p0001 source=h1 target=section reason=page-title

\section{第一三九封书信}

\Emph{致阿普罗尼乌斯}\par

\EditorialNote{关于阿普罗尼乌斯，我们一无所知；但根据他对依诺森（参见第一四三封书信）的提及，可以合理推断他住在西方。热罗尼莫在此祝贺他信仰坚定，并劝他前往白冷。随后他提及贝拉基造成的危害，并抱怨自己的隐修院被他或其同伙摧毁。此信写于公元417年。}\par

我不知道魔鬼用了何种诡计，致使你的一切辛劳、可敬的司铎依诺森的努力以及我本人的祈祷和心愿，眼下似乎都未生效。感谢天主，你安然无恙，且即使身陷魔鬼的诡计之中，信德之火仍在你内燃烧。我最大的喜乐是听说我的神子们为基督的事业而争战；我们信仰的那一位必定会激发我们的热忱，使我们欣然为捍卫祂的信德而流血。\par

我悲伤地听说一个显赫的家族被颠覆，原因我无从得知；因为送信人无法告诉我任何情况。我们理应为失去共同的朋友而悲伤，并恳求基督，那唯一的掌权者和主\BibleRef{弟前 6:15}，怜悯他们。同时，我们理应受到天主的惩罚，因为我们庇护了主的敌人。\par

你最好的选择是放下一切，来到东方，尤其是圣地；因为这里如今一切都平静。诚然，异端者们尚未清除胸中的毒液，但他们不敢张开不敬的口。他们\Quote{像那堵住耳朵的聋虺}。请替我向你可敬的弟兄们致意。\par

至于我们的住所，就世俗财富而言，已被异端者的攻击彻底摧毁；但藉基督的仁慈，它仍充满属灵的财富。靠面包而活，胜过失去信德。\par

\SourceRecord{Translated by W.H. Fremantle, G. Lewis and W.G. Martley.；Nicene and Post-Nicene Fathers, Second Series；Vol. 6.；Edited by Philip Schaff and Henry Wace.；Buffalo, NY: Christian Literature Publishing Co.,；1893.；<http://www.newadvent.org/fathers/3001139.htm>.}

% structure: p0001 source=h1 target=section reason=page-title

\section{第一百四十封书信}

\Emph{致长老西彼廉}\par

\EditorialNote{西彼廉曾到白冷拜访热罗尼莫，并请他写一篇对圣咏第九十篇的注释，要求语言简单易懂。热罗尼莫现在应此请求，逐节详释该圣咏，倾其学识，尤其运用希伯来文知识。他断言作者为梅瑟，但谨慎地补充说，天主的人可能不是为自己说话，而是以犹太人民的名义。他提到圣咏集可分为五卷，并指出将所有圣咏归于达味是错误的。一处对白拉奇主义的暗示表明，这封信必定属于热罗尼莫晚年，瓦拉西将其定于公元四一八年，很可能是正确的。}\par

\SourceRecord{Translated by W.H. Fremantle, G. Lewis and W.G. Martley.；Nicene and Post-Nicene Fathers, Second Series；Vol. 6.；Edited by Philip Schaff and Henry Wace.；Buffalo, NY: Christian Literature Publishing Co.,；1893.；<http://www.newadvent.org/fathers/3001140.htm>.}

% structure: p0001 source=h1 target=section reason=page-title

\section{第一四四封信}

第一四四封信。\Person{奥思定}致\Person{奥普塔图斯}。\par

\EditorialNote{\Person{奥思定}写信给\Place{米莱维}主教\Person{奥普塔图斯}，说他不能寄给他关于灵魂起源的致\Person{热罗尼莫}书信（\WorkTitle{第一三一封信}）的副本，因为没有\Person{热罗尼莫}的回复；他后来批评\Person{奥普塔图斯}反驳灵魂传生说的论据，指出其推理不充分。写信日期是公元420年。该信在翻译中有所压缩：原文错综复杂的句子已被简化，冗词也进行了削减。}\par

致蒙福的主内兄弟、真诚敬爱与渴慕的、同为主教的\Person{奥普塔图斯}，\Person{奥思定}在主内致候。\par

1. 藉着可敬的司铎\Person{撒图尔尼诺}之手，我收到了您的来信，可敬的阁下。您在信中恳切地向我索取我尚未得到的东西，这清楚地表明您相信关于这个问题我已收到答复。但愿如此！深知您殷切的期望，我决不会梦想隐瞒您所分享的恩赐；但如果您相信我，亲爱的兄弟，事实并非如此。虽然从我向东方的教会寄出我的信（那是一封询问而非断言的信）以来已过五年，但我迄今未收到任何回覆，因此无法按您的期望解开这个结。如果我有这两封信，我很乐意将它们一并寄给您；但我认为最好不要单独流传我的信，以免收信人（他可能仍按我所愿地回答我）感到不悦。如果我发表如此详尽的长篇大论而没有他的回覆，他可能会理所当然地愤怒，并认为我更多意在展示才华，而非促进某种有益的目的。这看起来好像我执意要提出一些他难以解决的问题。最好等他可能打算寄来的回覆。因为我很清楚，他有其他更严肃和紧迫的事务占据他的时间，超过我这个疑问。如果您读到他一年后在我送信人返回时所写给我的回信，您就会明白这一点。以下是他信中的摘要：\par

\Quote{我们遭遇了一个极其艰难的时刻，我发现保持沉默比说话更好。因此，我的研究已经停止，以免给\Person{阿皮乌斯}所谓的\InnerQuote{狗的巧辩}提供机会。为此，我未能对您两封博学而精彩的信函作出任何答复。当然，并非我认为它们有任何需要修正之处，而是因为我回忆起宗徒的话：\BibleQuote{有人这样判断，有人那样判断；让每个人充分表达自己的意见。}一个崇高的才智所能从圣经宝库中汲取的一切，您已经汲取了。这些，阁下必须允许我这样赞美您的才能。但是，尽管在我们之间的任何讨论中，我们的共同目标是促进学识，我们的对手，尤其是异端者，却会将我们之间的任何意见分歧归因于相互嫉妒。而就我而言，我决意爱您、敬仰您、尊崇和钦佩您，并像维护自己的意见一样维护您的意见。此外，在我最近发表的一段对话中，我以适当的方式提到了您的圣洁。因此，我们不如竭尽全力，将那个最有害的异端（它假装忏悔以便获得在教会中教导的自由）从教会中驱逐出去。因为如果它暴露在光天化日之下，它就会被驱逐而灭亡。}\par

2. 您可以看出，敬爱的兄弟，从这一回覆中，我的朋友并未拒绝回答我的询问；他是因为被迫将时间用于更紧迫的事务而推迟了。此外，他对我的善意从他友好的警告中可见一斑：我们之间出于爱德和为了学识利益而开始的辩论，可能会被嫉妒和异端者误解为出于相互敌意。不；最好让公众同时拥有两者——他的解释和我的询问。因为，如果他能够解释这个问题，我将感谢他的教导，那么这场讨论在世人知晓时将大有裨益。后世的人不仅会知道他们应当对这样一个经过充分辩论的话题持何种看法，而且会在天主恩慈的庇佑下，学习到在关爱中的兄弟们如何对困难问题进行辩论，同时又保持相互尊重。\par

3. 另一方面，如果我发表了提出这晦涩问题的信，而没有附上可以平息问题的答复，它可能会广泛流传，到达那些如宗徒所说的\Quote{彼此比较}、\Quote{和自己相比}\BibleRef{格后 10:12}的人手中，他们会误解他们无法理解的动机，并会将我对我所爱戴和尊敬、因其巨大贡献而敬重之人的情感，解释为不是他们所见的样子（因为这对他们是不可见的），而是任由他们自己的幻想和恶意所主导。现在，这是一种危险，在我能力范围内，我必须防范。但如果一份我不愿发表的文件未经我同意就被发表，并落入了我不愿让其接触的人手中，那么我将不得不顺从天主的旨意。事实上，如果我曾希望我的文字永远不为人知，我就绝不该把它们寄给任何人。因为如果（尽管我希望不会如此）偶然或必要阻止我得到任何答复，我的问询信仍迟早会公之于众。而且它对读者也不是毫无用处的；因为，虽然他们不会找到他们所寻求的，但他们将学到，当人不知情时，提出疑问远比作出断言要好得多；同时，他们所咨询的人将会阐明我提出的问题，放下争论，为了学识和爱德的益处，努力获得关于这些问题的正确见解。这样，他们要么得到他们渴望的解答，要么他们的能力会被激发，并从探究中了解到进一步的追问是无用的。然而现在，由于我没有理由对我的朋友的回信感到绝望，我决定不发表我已寄给他的信，并且我深信，我亲爱的同道，这个决定会得到你的认同。它理应如此，因为你所要求的不是我那封信，而是对它的答复；如果我有答复可以寄给你，我会很高兴地寄给你。的确，在你的信中，你提到了\Quote{我的智慧之清晰展示，那是赐光明者因我的生命而赐予我的}；如果你所指的不是我陈述问题的方式，而是我就所问之点获得的解答，我乐意满足你的愿望。但我必须承认，直到现在我尚未发现：灵魂如何能从\Person{亚当}承袭罪（这是一个不容质疑的真理），而灵魂本身却并非来自\Person{亚当}。目前，我认为最好进一步探究此事，而不是轻率地作出断言。\par

4. 你的信提到\Quote{许多老人和由博学司铎培养的人，你未能使他们回到你谦逊的思维方式，以及陈述真理本身的案情。}然而，你没有解释这种表达方式是什么。如果你的老人们坚持他们从博学司铎那里领受的，那么你为什么会被一群粗野无学识的圣职人员所困扰？另一方面，如果老人们和无学识的圣职人员已经邪恶地偏离了司铎的教导，那么后者——指司铎——当然应当纠正他们，并制止他们争论过度。再者，当你说\Quote{你作为一个初出茅庐、缺乏经验的教师，不敢篡改由伟大的著名主教们传下来的教义，并且你不愿把人引上更好的道路，以免使亡者蒙羞}时，你岂不是暗示，拒绝同意你的你所关切的对象，只不过是宁愿选择伟大著名主教们的传统，而不是一个初出茅庐、缺乏经验的教师的观点吗？关于他们在这件事上的行为，我什么都不说，但我非常渴望了解那\Quote{真理本身的表达方式}，不是被表达的事物，而是表达方式。\par

5. 因为你已经向我充分表明，你不赞成那些主张人的灵魂源自始祖的灵魂、并一代一代传下去的人；但你的信没有告诉我，你是根据什么道理和圣经的哪些章节证明这一观点是错误的。从你给\Place{凯撒勒雅}弟兄们的信，以及你最近写给我的信中，都没有明确你所赞同的是什么。我只看到你相信并写道：\Quote{天主过去是、现在是、将来也是人的创造者，天上地下没有任何事物不是完全因祂而存在。}这当然是一个无人质疑的常识。但既然你肯定灵魂不是传递的，你就应当解释天主用什么造它们。是用某种预先存在的材料，还是从无中造出？因为你不可能持有\Person{奥力振}、\Person{普里西连}和其他异端者的意见，\Emph{即}灵魂因先前生命中的行为而被困于尘世和可朽的身体中。这个意见确实被宗徒直接驳斥，他论及\Person{雅各伯}和\Person{厄撒乌}时，说他们在出生之前没有行善也没有作恶\BibleRef{罗 9:11}。那么，你对这件事的看法我只知道一部分，但你认为它正确的理由我一无所知。这就是为什么我在前一封信中请你寄给我你的信仰宣认，就是那篇你发现你的一位司铎不诚实签署而恼怒的那篇。我现在再次请求你寄来这个，以及你用于讨论这个问题的任何圣经章节。因为你在给\Place{凯撒勒雅}弟兄们的信中写道：\Quote{我已决定让平信徒法官（凭公开邀请列席，调查所有涉及信仰的问题）审查所有教义的界定。}你继续说：\Quote{天主的仁慈使他们能够以积极而确定的形式提出他们的观点，而你那谦逊的能力以大量证据加强了这些观点。}现在正是这\Quote{大量证据}我如此渴望得到。因为据我所见，你的唯一目的是反驳那些否认我们的灵魂是天主的化工的对手。如果他们持有这样的观点，你认为应该谴责是对的。如果他们对我们身体也说同样的话，他们将被迫收回，否则将受到咒骂。因为哪个基督徒能否认每一个人的身体是天主的化工？然而，当我们承认它们源于天主时，我们并不否认它们是人为生育的。因此，当有人主张我们的灵魂是由某种非物质种子繁殖而来，像我们的身体一样从父母而来，但藉天主的化工成为灵魂时，不是用人的猜测来反驳这一主张，而是用神圣圣经的见证。确实可以从具有正典权威的圣书中引用大量章节来证明我们的灵魂是天主的化工。但这些章节只反驳了那些否认每一个人的灵魂由天主造的人，而不是那些承认这一点却主张灵魂像身体一样通过父母的工具由天主形成的人。要反驳这些人，你必须寻找确凿的经文；或者，如果你已经发现这样的经文，请凭你的爱心告诉我。因为我虽然最勤勉地寻找，却未能找到。\par

正如你在信（给\Place{凯撒勒雅}弟兄们的信的末尾）中简短陈述的，你的困境如下：\Quote{因为我是你的儿子和门徒，最近才靠天主的帮助开始思考这些奥秘，我恳请你以你的司铎智慧教导我，在两个相反的观点中我应当持守哪一个。我应当坚持灵魂由世代传递，并以某种神秘的方式从我们的始祖\Person{亚当}而来\BibleRef{智 10:1}？还是我应当与你的兄弟和这里的司铎们一同坚持，天主过去是、现在是、将来也是万有和万人的作者和创造者？}\par

6. 在你提出的两个选项中，你希望被催促选择其中一个；如果这两个选项如此对立，以至于选择其中一个就意味着拒绝另一个，那么这是明智的做法。但事实上，一个人可能不说选择其中一个，而说两者都是真的。他可能坚持所有人的灵魂都来自我们的始祖\Person{亚当}，同时又相信并断言天主过去是、现在是、将来也是万有和万人的作者和创造者。按照你的原则，这样的人如何被反驳？我们是否可以说：\Quote{如果灵魂由世代传递，天主就不是它们的作者，因为祂不造它们？}那样他会回答：\Quote{身体也是生养出来的，不是天主造的；按你的说法，那么祂不是它们的作者。}有谁能坚持天主只造了\Person{亚当}的身体（祂用尘土造的）和\Person{厄娃}的身体（祂从\Person{亚当}的肋骨形成的），而其他身体不是祂造的，因为它们由人类父母生养？\par

7. 如果您的对手在坚持灵魂的起源问题上走得如此之远，以致否认灵魂是由天主所造和形成的，您就可以运用这一论点，在天主的帮助下尽可能反驳他们。但是，如果他们在断言灵魂的起源首先来自\Person{亚当}，然后来自人的父母的同时，却又坚持每一个人的灵魂都是由万有的创造者天主所创造和形成的，那么他们就只能通过圣经来加以反驳。因此，请您继续寻找，直到找到一段既不晦涩也不可能有两种含义的经文；或者如果您已经找到了一段，就请如我所恳求的那样告诉我。但是，如果您像我一样，迄今未能找到这样的经文，您仍然必须竭尽全力反驳那些声称灵魂绝非天主亲手所造的人。这似乎是您的对手的立场，因为您在您的第一封信中写道：\Quote{他们暗中散布丑闻性的教义，并因此愚昧、甚至是不敬的见解而离开了您的共融和教会的服从。} 针对这样的人，请尽一切可能的方式维护和坚持您在同一封信中所阐述的教义，即天主过去是、现在是、将来也是灵魂的创造者；天地间的一切存在完全归功于祂。因为这对每一个受造物都是真实的；因此应当相信、断言、维护和证明。天主过去是、现在是、将来也是一切事物和所有人的创造者和制造者，正如您对\Place{凯撒勒雅}省的同道主教们所说的那样，并以您的弟兄和同僚司铎的榜样劝勉他们采纳这一教义。但这里有两个截然不同的困境：（1）天主是否是一切灵魂和身体的创造者和制造者（这是正确的观点），还是自然界中有某种东西不是祂所造的（这是完全错误的观点）？（2）如果灵魂无疑是天主的亲手所造，祂是直接创造它们，还是通过传递间接创造它们？在处理第二个困境时，我希望您保持清醒和警醒。否则，在反驳传递论时，您可能会不慎陷入\Person{伯拉纠}的异端。人人都知道人的身体是通过生殖传递的；然而，如果我们说所有人类的灵魂——不仅\Person{亚当}和\Person{厄娃}的——都是由天主创造的，那么断言它们是通过生殖传递的并不等于否认它们的神圣起源。因为在这种观点下，天主创造灵魂正如祂创造身体一样，是通过生殖过程间接创造的。如果真理谴责这是一种错误，就必须寻找新的论据来反驳它。在这方面，没有人能比那些已故的贤哲更好地给您建议（只要他们能够接触得到），您曾担心因引导人们离开他们走上更好的道路而使他们丧失信誉。您说，他们是伟大而著名的主教，而您当时是一位初出茅庐、缺乏经验的教师；因此您不愿轻易改动他们的教义。我多么希望能知道这些伟人赖以支持灵魂传递论观点的经文！因为您在写给\Place{凯撒勒雅}弟兄们的信中，完全不顾他们的权威，将他们的观点说成是新奇的发明、闻所未闻的教义；尽管我们都知道，这或许是一种错误，但并非新奇，而是古老而历史悠久。\par

8. 当我们对某一点有理由怀疑时，我们无需怀疑我们正在怀疑是正确的。毫无疑问，我们应当怀疑那些可疑之事。例如，宗徒被提到三层天上时，对于自己是在身内还是在身外，他毫无怀疑地怀疑。\BibleRef{格后 12:4} 是这种情况还是那种情况，他说，我不知道；天主知道。那么，既然我没有光照，我为什么不能怀疑我的灵魂是通过生殖而来还是非生殖而来？既然我不怀疑无论哪种情况它都是至高天主的作为，我为什么不能对此怀疑？为什么我不能说：\Quote{我知道我的灵魂的存在归于天主，完全是祂的亲手所造；但它是如同身体一样通过生殖而来，还是如同\Person{亚当}的灵魂一样非生殖而来，我不知道；天主知道。} 您希望我明确断言一种观点或另一种。如果我知道哪个正确，我或许会这样做。您可能对此有一些光照，如果是这样，您会发现我更热衷于学习我所不知道的，而不是教导我所知道的。但是，如果您像我一样处于黑暗之中，您应该像我一样祈祷，要么通过祂的仆人之一，要么通过祂自己的口，教导我们那位曾对祂的门徒说：\BibleQuote{你们不要被称为师傅，因为你们的师傅只有一位，就是基督。} \BibleRef{玛 23:10} 然而，这样的知识只有在我们知道祂既知道祂必须教导什么，也知道我们应当学习什么时，才对我们有益。尽管如此，我亲爱的朋友，我承认我的渴望。尽管我多么渴望知道您所寻求的这一点，但我宁愿知道万民的渴望何时到来、圣人的国度何时建立，而不是我的灵魂如何来到它地上的居所。但当祂的门徒（即我们的宗徒）向全知的基督提出这个问题时，他们被告知：\BibleQuote{父以自己的权柄所定的时候和日期，不是你们应当知道的。} \BibleRef{宗 1:7} 如果基督知道什么对我们有益，却知道这知识无益，那又怎样？通过祂，我知道父以自己的权柄所定的时间不是我们应当知道的；但关于灵魂的起源，我不知道它是否属于我们应当知道的。如果我能确定这样的知识不属于我们，我不仅会停止武断，甚至会停止探究。事实上，尽管这个主题如此深奥和黑暗，以至于我害怕成为一个鲁莽的教师几乎超过了我渴望学习真理，但我仍然希望如果可能的话知道它。也许圣咏作者所祈求的知识：\BibleQuote{上主，求你让我知道我的结局，} 更为必要；然而我也希望我的起源也能启示给我。\par

9. 但即便在此事上，我不得对我的主忘恩负义。我知道人的灵魂是精神的而非物质的，具有理性和智力，并非天主的本质，而是一个受造物。它既有死，又是不朽的：前者因为它可朽坏，且与天主的生命分离——只有在其中才是幸福的；后者因为它的意识必永远持续，并成为其快乐或痛苦之源。确实，它的沉入肉身并非源于肉身之前的行为；然而在人身上，它从不离罪，甚至当\Quote{它的生命只有一天}时也是如此。在亚当子孙所生的人中，无人是无罪而生的，就连婴孩也必须藉基督的再生恩宠而重生。所有这些都是我关于灵魂所知道的，这已经很多；其中大部分不仅是知识，也是信德之事。我欣喜地学到这一切，并能真实地说我知道。如果还有我所不知道的事（例如天主是藉生育还是独立地创造灵魂——因为祂确实创造灵魂，我毫不怀疑），我宁愿知道真相而不愿无知。但只要我无法知道，我宁可暂不判断，也不断言明显违背不可争辩的真理的事。\par

10. 我的兄弟，你要求我为你决定，人的灵魂是像他们的身体那样藉生育从亚当而来，还是像他的灵魂那样没有生育而分别为每个人创造的。因为无论哪种方式，我们都承认它们是天主的工程。那么请允许我反过来问你一个问题。一个灵魂能从一个它本身并非由之而来的根源获得原罪吗？除非我们要陷入可憎的佩拉季异端，否则我们俩都必须承认，所有灵魂都从亚当获得原罪。如果你不能回答我的问题，那么请你允许我承认我同样不知道你的问题和我的问题。但如果你已经知道我所问的，请教导我，然后我将教导你所想知道的。请不要因我采取这个方式而不高兴，因为虽然我没有对你的问题给出明确的答案，但我已经向你展示了应该如何提出它。一旦你清楚这一点，你就可以对自己所怀疑的地方非常肯定。\par

我认为应当向你的圣座写这些，因为你如此确信灵魂的传递是一种应被拒绝的教义。如果我是写给该教义的支持者，我或许会表明他们对自己以为知道的事是多么无知，以及他们应当多么谨慎，不要妄下断语。\par

也许令你困惑的是，在我于这封信中引用的我朋友的答复里，他提到他无暇回复的\Emph{两}封我的信。其中一封涉及灵魂的问题；在另一封信中，我请求解决另一个难题。再者，当他敦促我要更加努力从教会中除去一个最有害的异端时，他指的是佩拉季人的错误，我恳切地请求你，我的兄弟，无论如何要避免它。在探讨或辩论灵魂的起源时，你绝不可给这个异端及其阴险的暗示留有余地。因为除了一位中保的灵魂之外，没有一个灵魂不从亚当获得原罪。原罪就是灵魂出生时附着在其上的，只有通过重生才能摆脱。\par

\SourceRecord{Translated by W.H. Fremantle, G. Lewis and W.G. Martley.；Nicene and Post-Nicene Fathers, Second Series；Vol. 6.；Edited by Philip Schaff and Henry Wace.；Buffalo, NY: Christian Literature Publishing Co.,；1893.；<http://www.newadvent.org/fathers/3001144.htm>.}

% structure: p0001 source=h1 target=section reason=page-title

\section{第一百四十五封信}

\Emph{致\Person{Exuperantius}}\par

\EditorialNote{热罗尼莫劝告\Person{Exuperantius}，一位罗马士兵，来到\Place{白冷}，和他的兄弟\Person{Quintilian}成为隐修士。根据\Person{Palladius}（\WorkTitle{拉乌息亚史}第\OriginalTerm{八十章}{c. lxxx.}），\Person{Exuperantius}来到热罗尼莫那里，但又离开了，\Quote{不能忍受他的暴力和恶意}。这封信的日期不详。}\par

在我与可敬的兄弟\Person{Quintilian}的友谊所给予我的一切恩惠中，最大的恩惠是，他让我在精神上认识了您，尽管我并不认识您本人。谁不会爱一个这样的人呢？他身穿士兵的斗篷和制服，却履行着先知的职责；他的外表让人以为他是另一种人，但他内在的人却克服了这一切，这内在的人是按照造物主的肖像形成的。因此，我出面挑战您，要与您交换书信，并恳求您常常给我回复的机会，这样我将来写信时就会少些拘束。\par

目前，我将满足于向您的明智建议：您应牢记宗徒的话：\BibleQuote{你束缚于妻子吗？不要寻求解脱。你解脱于妻子吗？不要寻求妻子；}\BibleRef{格前 7:27}也就是说，不要寻求那种与解脱相悖的束缚。谁若承担了婚姻的义务，就是被束缚的；被束缚的就是奴隶；另一方面，被解脱的就是自由的。既然你因基督的自由而喜乐，既然你的生活优于你的职业，既然你几乎已在救主所说的屋顶上，你就不应下来拿你的衣服。\BibleRef{玛 24:17}你不应回头，你既已手扶犁，就不应再放手。\BibleRef{路 9:62}相反，如果你能，就效法若瑟，把你的衣服留在埃及主母的手中，\BibleRef{创 39:12}这样你就可以赤身跟随你的主和救主。因为在福音中祂说：\BibleQuote{凡不抛弃他所有的一切，不背起自己的十字架，并跟随我的，不能做我的门徒。}\BibleRef{路 14:26}抛开这世界之事的重担，不要追求那些在福音中被比作骆驼峰一样的财富。赤身且无挂碍地飞升天堂；大堆的黄金只会阻碍你德行的翅膀。我这样说并非因为我知道你贪财，而是因为我猜想你长期留在军队中的目的是为了填满那个主曾命令你倒空的钱包。因为那些拥有财产和财富的人被告知要卖掉他们所有的一切，分给穷人，然后跟随救主。\BibleRef{玛 19:21}因此，如果阁下您已经富有，就应履行命令，卖掉你的财富；或者如果你仍然贫穷，就不应积攒那些你将来必须施舍的东西。基督接受为祂所作的祭献，\BibleRef{格后 8:12}按照奉献者所有的心愿。从未有人比宗徒们更贫穷；但也从未有人为主留下更多。福音中那位只投了两个小钱在银库里的穷寡妇，被置于所有富人之前，因为她把她所有的一切都投进去了。\BibleRef{路 21:1}你也应如此。不要追求那些你将来必须施舍的财富；相反，要放弃你已经获得的，这样基督就会知道祂的新兵是勇敢而坚定的，然后，当你还在远处时，祂的父亲就会欢喜地跑来迎接你。祂会给你一件袍子，给你戴上戒指，为你宰杀肥牛犊。\BibleRef{路 15:20}那时，当你摆脱一切牵挂，天主很快就会为你开辟一条路，让你与可敬的兄弟\Person{Quintilian}渡海来到我这里。我现在已经叩响了友谊之门：如果你为我开门，你会发现我是一位常客。\par

\SourceRecord{Translated by W.H. Fremantle, G. Lewis and W.G. Martley.；Nicene and Post-Nicene Fathers, Second Series；Vol. 6.；Edited by Philip Schaff and Henry Wace.；Buffalo, NY: Christian Literature Publishing Co.,；1893.；<http://www.newadvent.org/fathers/3001145.htm>.}

% structure: p0001 source=h1 target=section reason=page-title

\section{第一四六封书信}

致\Person{爱凡格路}\par

\EditorialNote{哲罗姆驳斥那些使执事与长老平等的人，但这样做时他自己使长老与主教平等。信件的日期未知。}\par

1. 我们在\BibleBook{依撒意亚}{先知书}中读到这句话：\Quote{愚昧人必说愚昧话，} 我听说有人竟疯狂到将执事置于长老之上，即置于主教之上。因为当宗徒清楚教导长老与主教是同一回事时，一个仅管理饭食和寡妇的人，竟狂妄地凌驾于那些藉其祈祷而产生基督体血的人之上，难道不是疯狂吗？\BibleRef{宗 6:1} 你要我证明我所说的话吗？请听这段经文：\Quote{\Person{保禄}与\Person{弟茂德}，耶稣基督的仆人，致书给在\Place{斐理伯}的、在基督耶稣内的众圣徒，连同主教与执事。} 你还要另一个例子吗？在\BibleBook{}{宗徒大事录}中，\Person{保禄}这样对单一教会的司铎们说：\Quote{圣神既立你们为主教，牧养天主用自己的血所取得的教会，你们要留心自己和整个羊群。} 为避免有人好争论，说那时一个教会必有不止一位主教，下面这段经文清楚证明主教与长老是同一回事。宗徒致\Person{弟铎}书说：\Quote{我从前留你在\Place{克里特}，是要你整顿那尚未完成的事，并照我所吩咐你的，在各城设立长老：若有无可指责的人，只作一个妻子的丈夫，有信德的子女，没有放荡或不服管教的指控。因为主教，作为天主的管家，必须无可指责。} \BibleRef{铎 1:5} 又致书\Person{弟茂德}说：\Quote{不要忽略你所得的恩赐，就是从前藉预言、并藉长老团的覆手赐给你的。} \BibleRef{弟前 4:14} \Person{伯多禄}也在他的第一封书信中说：\Quote{我这同为长老的，为基督苦难的见证人，并分享那将要显示的光荣的人，劝勉你们中间的众长老：牧放基督的羊群……以天主所喜悦的心情，并非出于勉强，而是出于甘心。} 在希腊文中，意思更加明显，因为所用的词是\OriginalTerm{监督}{\GreekText{επισκοποῦντες}}，即\InnerQuote{监督}，这正是监督或主教的名称来源。但或许这些伟人的见证在你看来不够充分。若是这样，那么就听一听福音号角的声音，那位雷霆之子，\BibleRef{谷 3:17} \Person{耶稣}所爱的门徒，\BibleRef{若 13:23} 他曾靠在救主的怀中，畅饮纯正教义的活水。他的一封书信这样开头：\Quote{长老致蒙选的主母和她的儿女，就是我在真理中所爱的；} 另一封这样说：\Quote{长老致可爱的\Person{加约}，就是我在真理中所爱的。} 后来，选立一位长老主持其余长老，这是为补救分裂，防止各人因拉拢信友而撕裂基督的教会。因为甚至在\Place{亚历山大里亚}，从福音圣史\Person{马尔谷}的时代直到\Person{赫拉克拉斯}和\Person{狄约尼削}担任主教的时候，长老们总是从他们自己当中推选一位，立为主教，并置于更高的地位，如同军队选举将军一样，或者如同执事们从他们当中推举一位他们知道是勤勉的人，称他为总执事。除了授予圣秩之外，有什么职能是属于主教而不也是属于长老的呢？并非\Place{罗马}有一个教会，而世界上其他地方有另一个教会。\Place{高卢}与\Place{不列颠}、\Place{非洲}与\Place{波斯}、\Place{印度}与\Place{东方}，都敬拜同一个基督，遵守同一真理的规则。若你问权威，全世界大过它的首都。无论何处有主教，无论是在\Place{罗马}还是在\Place{恩古比姆}，无论是在\Place{君士坦丁堡}还是在\Place{勒基乌姆}，无论是在\Place{亚历山大里亚}还是在\Place{左安}，他的尊位是同一的，他的司铎职是同一的。财富的显赫或贫穷的卑微，都不会使他成为更高级或更低级的主教。他们都同样是宗徒的继承人。\par

2. 但你或许会说：那么，为什么在\Place{罗马}，长老只有在执事的推荐下才被授予圣秩？对此我回答如下：你为什么提出一个只存在于一城的惯例？你为什么以一个微不足道的例外来反对教会的法律，而这例外已引发傲慢与骄傲？东西越稀有，越被人追求。在\Place{印度}，薄荷比胡椒更昂贵。因为执事人数稀少，他们便成了重要人物；而长老因人数众多，反而不受重视。但即使在\Place{罗马}教会，执事站立，长老就座，虽然不良习惯已逐渐蔓延，我甚至看见过一位执事，在主教不在时，坐在长老中间，并在社交聚会中向他们赐福。如此行事的人必须认识到他们的错误，并留心宗徒的话：\Quote{我们放弃天主的圣言而管理饭食，实在不相宜。} \BibleRef{宗 6:2} 他们必须思考当初设立执事的理由，必须阅读\BibleBook{}{宗徒大事录}，并牢记他们真正的地位。\par

关于长老和主教这两个名称，前者表示年龄，后者表示品级。宗徒在写给\Person{弟铎}和\Person{弟茂德}的书信中，只谈到主教和执事的授予圣秩，却只字不提长老的授予圣秩；事实上，主教的名称也包括了长老。再者，人的晋升是从较低的地位到较高的地位。那么，要么长老被授予执事圣秩，从较低的职务升到更重要的职务，以证明长老低于执事；要么，另一方面，如果执事被授予长老圣秩，那么这位长老应当认识，尽管他可能比执事的报酬低，但他藉着司铎职而优越于执事。事实上，好像是要告诉我们，宗徒所传下来的传统出自旧约，主教、长老和执事在教会中所占的地位，与\Person{亚郎}、他的儿子们和\Person{肋未}人在圣殿中所占的地位相同。\par

\SourceRecord{Translated by W.H. Fremantle, G. Lewis and W.G. Martley.；Nicene and Post-Nicene Fathers, Second Series；Vol. 6.；Edited by Philip Schaff and Henry Wace.；Buffalo, NY: Christian Literature Publishing Co.,；1893.；<http://www.newadvent.org/fathers/3001146.htm>.}

% structure: p0001 source=h1 target=section reason=page-title

\section{第一四七号书}

\Person{致萨比尼亚努斯}\par

\EditorialNote{热罗尼莫以严厉但温和的措辞写信给执事萨比尼亚努斯，呼吁他悔改自己的罪。信中详细列举了他两项最严重的罪行：在罗马犯下的奸淫罪和在白冷试图引诱一位修女。此信的写作日期不确定。}\par

1. 昔日，当主后悔立了撒乌耳作以色列王时，我们读到撒慕尔为他哀悼；同样，当保禄听说在格林多人中间有淫乱的事，并且是这样的淫乱，连外邦人中也没有时，\BibleRef{格前 5:1} 他含泪劝他们悔改，说：\BibleQuote{怕我再来的时候，我的天主使我在你们面前惭愧，并且我要为许多从前犯罪、没有悔改他们所有的不洁、淫乱和邪荡的人哀哭。} \BibleRef{格后 12:21} 如果一位宗徒或先知，自己无玷，能用此等包容一切的宽仁说话，那么像我这样的罪人，更应该多么恳切地为你这样的罪人求情！你跌倒了，却不愿起来；你甚至不举目望天；你浪费了父亲的财产，却喜爱猪吃的豆荚；你登上骄傲的悬崖，却一头栽进深渊。你以肚腹为天主，而不是基督；你作情欲的奴仆；你的光荣在于你的羞辱；\BibleRef{斐 3:19} 你养肥自己如同待宰的牲口，效法恶人的生活，不顾他们的命运。\BibleQuote{你不知道天主的恩慈是领你悔改吗？但你竟任着你那顽固而不悔改的心，为自己积蓄忿怒，直到忿怒的日子。} \BibleRef{罗 2:4} 难道你的心刚硬了，如同法郎的心，因为你的惩罚被延迟，没有立即受击打吗？十灾降在法郎身上，不是出于发怒的天主，而是出于警告的父亲，他的恩宠之日被延长，直到他后悔了自己的悔改。然而，当他追赶那被他释放的百姓，在旷野中执着地追赶，并擅自进入海中时，灾祸终于临到他。因为只有这样才能让他学到，连元素都服从的那一位是可畏的。他曾说：\BibleQuote{我不认识主，也不让以色列去；} \BibleRef{出 5:2} 而你效仿他，说：\BibleQuote{他所见的异像是关于许多日子以后的事，他预言的是遥远的时期。} \BibleRef{则 12:27} 然而，同一位先知用这些话驳斥你：\BibleQuote{主上主这样说：我的话不再迟延，我所说的话必然成就。} 达味也论到不敬虔的人（你证明自己不是轻微、而是显著的不敬虔之例），说他们在今世享福，说：\BibleQuote{天主怎能知道？至高者岂有知识？看，这些是不敬虔的人，他们在世亨通，财宝增加。} 然后他几乎失脚，站立不稳，抱怨说：\BibleQuote{我实在徒然洁净了我的心，徒然洗手表明无辜。} 因为他先前说：\BibleQuote{我嫉妒愚妄人，我见恶人享平安。他们至死没有痛苦，他们的力气也壮实。他们不像别人受苦，也不像别人遭灾。所以骄傲如链子戴在他们颈上，强暴像衣裳遮住他们。他们的眼睛因体胖而凸出；他们所得的，过于心里所想的。他们讥笑，凭恶意说欺压人的话，说话高傲。他们的口亵渎上天，他们的舌在地上来往。} \BibleRef{咏 73:3-9}\par

2. 这整篇圣咏难道不是为你写的吗？确实，你强健有力；仿佛一位新的敌基督的宗徒，当你在一个城市被识破时，你就转到另一个城市。\BibleRef{玛 10:23} 你无需钱财，没有致命的打击击倒你，你也不像别人那样遭灾，他们不像你仅仅是牲畜。因此你狂傲自大，情欲如衣裳遮盖你。从你肥胖臃肿的躯体里，你吐出死亡的话语。你从不思考你必须死，也绝不在满足情欲之后感到丝毫悔改。你所得的，过于心里所想的；并且，为了不独自作恶，你捏造关于天主仆人的丑闻。虽然你不知道，其实你是对至高者说不义的话，口对天说狂傲的话。难怪天主的仆人，无论大小，都受你的亵渎，因为你的父辈曾毫不犹豫地称家主为\Quote{贝耳则步}。\BibleQuote{徒弟不能超过师傅，奴仆不能超过主人。} \BibleRef{玛 10:24} 他们若对青绿的树这样行，对我这枯干的树又会怎样呢？\BibleRef{路 23:31} 同样，在玛拉基亚书中，那些被冒犯的信徒也表达了与你相似的情感；他们说：\BibleQuote{事奉天主是徒然的；遵守他的吩咐，在主万军的上主面前哀痛而行，有什么益处呢？如今我们称骄傲的人为有福；行恶的人被建立；试探天主的人竟然得救。} 然而，主后来用审判的日子警告他们，并预先宣告那时要在义人和不义之人之间作出区分，对他们这样说：\BibleQuote{你们回转吧，又要分别义人和恶人，事奉天主和不事奉天主的。}\par

3. 这一切在你看来或许都是笑谈，因为你如此喜爱喜剧、抒情诗和\Person{楞图禄}那样的摹拟剧；尽管你的心智如此迟钝，我甚至不认为你能理解这样简单的语言。你可以藐视先知的话，但\Person{亚毛斯}仍要回答你：\Quote{主这样说：\InnerQuote{为了三件事，为了四件事，我岂能不收回成命？}} 因为大马士革、迦萨、提洛、厄东、阿孟人和摩阿布人，以及犹太人和以色列子民，尽管天主常常预言要他们悔改，他们却不肯听从祂的声音，主要显示祂将要降下的忿怒有最正当的理由，就用了上面引用的那句话：\Quote{\InnerQuote{为了三件事，为了四件事，我岂能不收回成命？}} 天主说，怀有恶念是邪恶的；但我仍容许他们如此。付诸实行更为邪恶；但我因慈悲和仁慈，甚至也容许了这一点。但罪恶的念头是否已成了罪恶的行为？人是否因骄傲而如此践踏我的温柔？然而，\BibleQuote{我不喜欢恶人丧亡，却喜欢恶人离开他的道路而生活；}\BibleRef{则 33:11} 既然\BibleQuote{健康的人不需要医生，有病的人才需要，}\BibleRef{路 5:31} 即使在犯罪之后，我仍向跌倒的罪人伸出手，劝告他——即使他被自己的血所污染\BibleRef{则 16:6} ——以忏悔之泪洗去污点。但如果他仍不愿悔改，如果他在遭遇船难后拒绝抓住唯一能拯救他的木板，我最终不得不这样说：\Quote{主这样说：\InnerQuote{为了三件事，为了四件事，我岂能不收回成命？}} 因为天主把这种\Quote{转离}视为惩罚，因为罪人被任凭自生自灭。祂就是这样\BibleQuote{追讨父亲的罪债到儿子身上直至三四代。}\BibleRef{出 20:5} 祂不是立刻惩罚那些直接犯罪的人，而是宽恕他们最初的过犯，只到最后才对他们宣判。因为如果不是这样，如果天主立刻成为罪恶的报复者，教会将失去许多圣人；当然也会失去宗徒\Person{保禄}。上面引用过的先知\Person{厄则克耳}重复天主对他所说的话，这样说：\Quote{开口吃我给你的东西。看，}他说，\Quote{一只手向我伸出，里面有一卷书；他将其展开在我面前，内外部写着字，上面写着哀歌、颂歌和祸患。} 这三种中的第一种属于你，如果你愿意作为一个罪人，忏悔你的罪过。第二种属于那些圣洁的人，他们被召唤来颂扬天主；因为赞美不出于罪人之口。第三种属于像你这样因绝望而放纵于不洁、奸淫、肚腹和最低下的情欲之人；这些人以为死亡就是终结，没有来世；他们说：\BibleQuote{当泛滥的灾祸经过时，不会落到我们身上。}\BibleRef{依 28:15} 先知所吃的书卷就是整部圣经，它依次哀悼忏悔者、颂扬正义者、诅咒绝望者。因为没有什么比一颗不悔改的心更令天主不悦。不悔改是唯一不可宽恕的罪。因为如果犯罪后停止犯罪的人得到赦免，如果祈祷有能力使法官改变心意，那么每一个不悔改的罪人必然激起法官的愤怒。因此，绝望是唯一无药可救的罪。通过顽固地拒绝天主的恩宠，人们将祂的慈悲变为严厉和苛刻。然而，为让你知道天主天天召唤罪人悔改，请听\Person{依撒意亚}的话：\Quote{在那一天，}他说，\Quote{万军的上主呼唤哭泣、哀号、剃头、腰束苦衣；但看，欢乐和喜乐，宰牛杀羊，吃肉喝酒：\InnerQuote{我们吃喝吧，因为明天我们就要死了。}} 在这些充满绝望的轻率之语之后，圣经接着说：\BibleQuote{万军的上主亲自启示给我：这罪孽决不得从你们身上赦免，直到你们死。}\BibleRef{依 22:12} 只有当他们死于罪时，他们的罪才会被赦免。因为只要他们仍在罪中生活，罪就不能被除去。\par

4. 求祢怜悯，我恳求祢，怜悯祢的灵魂。请思量天主的审判终有一天会临到祢身上。记住祢是由怎样的一位主教祝圣的。那位圣人在选择上犯了错误；但这情有可原。因为就连天主也曾后悔傅油立\Person{撒乌耳}为王。\BibleRef{撒上15:11} 即使在十二宗徒中，也发现了\Person{犹达斯}这个叛徒。而\Place{安提约基雅}的\Person{尼苛劳}——像祢一样的执事\BibleRef{宗6:5}——散播了尼苛劳党人的异端和种种不洁。\BibleRef{默2:6} 我现在不向祢提起祢据说曾诱惑的许多贞女，或因被祢玷污而蒙受死亡的尊贵主妇，以及祢在罪恶巢穴中贪婪放荡的行径。因为这类罪过本身虽严重而沉重，但与我即将叙述的相比，实则微不足道。比诱惑与奸淫更微不足道的罪，该是何等重大？祢这可怜虫！当祢进入天主圣子诞生的山洞，即真理从大地生出、土地产出果实之处，却是为了幽会。难道祢不怕圣婴从马槽中哭喊？不怕刚分娩的贞女看见祢？不怕主的母亲注视祢？天使高呼，牧人奔跑，星光从天照耀，贤士朝拜，\Person{黑落德}惊恐，耶路撒冷骚动，而此时祢却潜入贞女的斗室，诱惑属于她的贞女。当我试图将祢所作的行径呈现在祢眼前时，我充满惊骇，灵魂与身体一阵战栗。整个教会彻夜守夜，宣认基督为主；虽语言不同，却同心同德颂扬天主的赞美诗。然而祢却将祢的情书塞进如今是祭台、曾是主马槽的缝隙中，选择此处，好让祢不幸的牺牲品前来跪拜时发现并阅读。随后祢混入歌咏者中，以无耻的点头向她传递祢的情欲。\par

5. 啊！可耻的哭泣！我无法继续。因为哭泣抢在话语之前，愤怒与悲伤哽住了我的发声。啊！愿有\Person{图理乌斯}那雄辩之海！啊！愿有\Person{狄摩斯提尼}那激烈抨击之流！然而在此案中，我相信祢们都会哑口无言；祢们的口才会失效。一种任何修辞都无法解释的行为已被揭露；一种任何哑剧都无法表现、任何小丑都无法扮演、任何喜剧演员都无法描述的罪行已被发现。\par

在\Place{埃及}和\Place{叙利亚}的隐修院中，惯常做法是：那些已向天主发愿、弃绝世俗并践踏其欢乐的贞女与寡妇，请求她们团体的母亲为她们剪发；并非此后她们就违抗宗徒的命令而露头行走，\BibleRef{格前11:5} 因为她们戴着紧贴头部的帽子和面纱。除了剪发者与受剪者之外，无人知晓此事；但因这做法普遍，几乎尽人皆知。事实上，这习俗已成为第二天性。其目的在于避免那些不沐浴、头脸从未接触任何油膏的人积累污垢，以及有时在发根滋生的小生物。\par

6. 那么，我的好友，让我们看看祢在此等环境中是如何行事的。祢许诺要娶祢不幸的牺牲品；随后在那可敬的山洞中，祢从她那里取走了一些发缕、几方手帕和一条腰带，作为她忠贞的保证或婚约的信物，同时发誓祢绝不会再爱别人像爱她一样。然后祢跑到牧人听见天使在上空歌唱时看守羊群的地方，再次盟誓。我不再多说；我不指控祢亲吻或拥抱她。虽然我相信祢无恶不作，但马厩与田野的神圣性质阻止我假设祢有罪，除非是意愿与决心。不幸的人！当祢最初站在山洞中的贞女身旁时，想必迷茫遮蔽了祢的眼睛，祢的舌头瘫痪，双臂垂落，胸膛起伏，步履蹒跚。她已在宗徒\Person{伯多禄}的大殿中戴上了基督的新娘面纱，并立誓从此住在隐修院中，即被主的十字架、复活与升天祝圣的地点；然而在这之后，祢竟敢接受那缕头发——她在基督的命令下于祂诞生的山洞中剪下的——作为她愿意与祢同寝的信物。此外，祢常从傍晚至清晨坐在她的窗下；因窗户太高无法接近，祢便用绳索传递物品，她亦如此回传。女主人的谨慎程度，由祢只在教堂中见过贞女一事可见一斑；而且，尽管祢俩有同样的倾向，却除夜晚在窗下外无法交谈。后来有人告诉我，祢在日出时总是十分遗憾。祢的面色苍白、憔悴、毫无血色；为消除一切猜疑，祢常诵读基督的福音，仿佛真是执事。我和其他人曾将祢的苍白归因于禁食，并钦佩祢那通常鲜艳的口唇竟变得毫无血色，以为是频繁守夜所致。祢已预备好梯子，要将那不幸的贞女从她的斗室中劫出；祢已安排路线，订购船只，定下日期，谋划逃亡细节。然而，看，那位守护\Person{玛利亚}内室的天使，那位看守主马槽、怀抱婴孩基督的天使——在祂面前祢犯下了这些大罪——正是他，而非他人，揭露了祢。\par

7. 嗟乎！我可怜的眼目！嗟乎！堪受最严厉诅咒的日子，在那日我读你的信时，惊惶失措，其内容我至今仍被迫记忆！信中竟有何等秽语！何等谄媚！何等洋洋得意于贞女的受辱！一个执事甚至不应知晓这些事，更遑论谈论它们。不幸的人！你从何处学得这些？你曾夸口在教会中长大。然而，这些信上你发誓从未度贞洁生活，且实际上你并非真正的执事。你若试图否认，你的笔迹将定你罪，连这些信也会向你呼喊。但此刻，你可自得其罪，因你写下的污秽如此不堪，我无法将它们当作证据来指控你。\par

8. 你俯伏在我膝前，五体投地，恳求我——我用你的原话——饶恕\Quote{你的半品脱血}。嗟乎！可怜虫！你全然不顾天主的审判，只怕我的报复。我宽恕了你，我承认；身为基督徒，我还能做什么？我催促你悔改，披麻蒙灰，滚于灰中，寻求隐居，住在隐修院，以不断的泪水恳求天主的仁慈。然而你却显为自信的柱石，被毒蛇之刺所激，你向我成了诡诈的弓；你向我射出辱骂的箭。我因告诉了你真理，就成了你的仇敌。\BibleRef{迦 4:16}我不抱怨你的诽谤；人人都知道你只赞美与你同样声名狼藉的人。我所哀叹的是你不哀叹自己，不意识到你已经死了，像那为利比蒂娜预备好的角斗士，你竟为自己葬礼打扮。你穿的并非麻衣，而是细麻；你手指戴满戒指，你用牙粉刷牙，你精心梳理你棕色头颅上的乱发。你的牛颈因肥肉突显，毫无下垂，因它已向情欲让步。此外，你香气四溢，从一个浴室到另一个浴室，你与那不顾你意愿而生长之头发作战，你走过广场和街道，一个整洁滑面的浪子。你的脸成了娼妓的脸：你不知羞耻。\BibleRef{耶 3:3}不幸的人，归向主，祂必归向你。\BibleRef{拉 3:7}悔改吧，祂必回心转意，不降祂所定意要降于你的灾祸。\par

9. 你为何无视自己的伤疤，却要诽谤他人？为何我给你最好的劝告，你却像疯子一样攻击我？或许我正如你公开宣称的那样声名狼藉；那样的话，你至少可以像我一样诚心悔改。或许我正如你所说的那样是个大罪人；若是，你至少可以效法一个罪人的眼泪。难道我的罪成了你的美德？或者，许多人和你处境相同就能减轻你的痛苦吗？让几滴眼泪落在那使你光彩夺目的丝绸和细麻上吧。意识到你是赤身露体、破烂不堪、不洁、是个乞丐。\BibleRef{默 3:17}悔改永不嫌晚。你或许已从耶路撒冷下来，在路上受了伤；然而撒玛黎雅人必把你放在他的牲口上，带你到客栈，照顾你。\BibleRef{路 10:30}即使你躺在坟墓里，主也必使你复活，虽然你的肉身可能已经发臭。\BibleRef{若 11:39}至少效法那些瞎子，救主为他们离开了自己的家园和产业，来到耶里哥。他们坐在黑暗和死影中，当光照耀他们时，\BibleRef{路 1:79}因为他们得知是主经过，便开始喊叫说：\Quote{达味之子，可怜我们吧！}你也要恢复视力，只要你向祂呼求，并听从祂的呼唤，丢掉你的污秽衣服。\BibleRef{谷 10:50}\Quote{当你回头哀哭自己时，你必得救，那时你将知道你从前在哪里。}让祂触摸你的伤疤，将手放在你的眼球上；虽然你可能是生来瞎眼，你的母亲可能犯罪怀了你，祂必用牛膝草净化你，你必洁净；祂必洗涤你，你必比雪更白。你为何佝偻着弯向地面？为何仍匍匐在泥沼中？那个被撒殚捆绑了十八年的妇女来到救主面前；祂治好了她，使她挺直，得以再次仰望天堂。\BibleRef{路 13:11}天主对你说的话，正如祂对加音说的：\Quote{你犯了罪，静默吧！}你为何逃避天主的面，住在诺得地方？你为何在波浪中挣扎，而你能把脚立在磐石上？当你在与米德杨女子行淫时，留心丕乃哈斯不将他的长矛刺透你。\BibleRef{户 25:6}阿默农没有饶恕塔玛尔，\BibleRef{撒下 13:14}而你，她信仰中的弟兄和亲属，竟没有怜悯这位贞女。但为何在你玷污她之后，你变成了阿贝沙隆，想要杀死那哀悼你叛乱和灵性死亡的大卫？纳波特的鲜血\BibleRef{列上 21:13}向你喊冤。依则勒耳的葡萄园，即天主的种子，也向你要求应得的报复，因你把它变成了享乐园，使它成了情欲的苗床。天主派厄里亚告诉你折磨和死亡。因此你要俯首片刻，穿上麻衣；也许主会像论及阿哈布那样论及你：\Quote{你看见阿哈布怎样在我面前自卑吗？因为他在我面前自卑，\BibleRef{列上 21:29}我不在他活着的时候降这灾祸。}\par

10. 但也许你自诩，既然那位擢升你为执事的主教是位圣人，他的功德会弥补你的过犯。我已告诉过你，父亲不为儿子受罚，儿子也不为父亲受罚。\Quote{犯罪的灵魂，它必要死亡。} \BibleRef{则 18:4} 撒慕尔也有儿子，他们离弃了对上主的敬畏，\Quote{转向了贪财和不义。} \BibleRef{撒上 8:3} 厄里也是一位圣洁的司铎，但他的儿子，我们在希伯来文中读到，他们与聚集在会幕门口的妇女同寝，并且像你一样，无耻地为自己要求在圣所服务的权利。因此，会幕本身倾覆，圣地因那些作为天主司铎者的罪恶而沦为荒凉。甚至厄里本人也因对儿子过分宽纵而得罪了天主；因此，你的主教的义德非但不能拯救你，反倒更可惧的是，你的邪恶可能将他从座位上推下，使他像厄里一样仰面跌倒，无可挽回地灭亡。\BibleRef{撒上 4:18} 如果肋未人乌匝只是因为试图扶住那本应由他特殊负责抬运的约柜不倒下而受击打；\BibleRef{撒下 6:6} 那么你认为，对于那些试图在主的约柜稳固时将其推翻的人，将施以何种惩罚？那位祝圣你的主教越值得尊敬，你就越可憎，因为你辜负了如此善人的期望。他长久不知你的恶行，这确实容易理解；因为通常我们总是最后才知道影响我们自家丑闻的事，并且对子女和妻子的罪过茫然不知，即使邻人谈论的别无其他。无论如何，全意大利都知悉你的邪恶生活；各处都在哀叹你竟仍站在基督的祭台前。因为你既无诡计也无远见来隐藏你的恶习。你如此狂热、如此好色、如此放荡，完全被各式各样的任性放纵所支配，以至于不满足于满足你的情欲，你还夸耀每一次勾引是胜利，并从中取得胜利的棕枝。\par

11. 再次，不洁的火焰攫住了你，这次是在野蛮人的刀剑之间，在一个有势力有权力、已婚的蛮族人寓所中。你竟敢在受辱的丈夫无需借助法官就能惩罚你的屋子里犯奸淫。你被郊外的园林和花园吸引而前往；在丈夫不在时，你行为如此大胆疯狂，仿佛你以为你的同伴不是你的情妇，而是你的妻子。她最终被捉住，而你却通过一条地下通道逃脱，秘密地来到罗马。你在那里藏在一些萨莫奈强盗中间；当一听到那受辱的丈夫如同新的汉尼拔正从阿尔卑斯山下来搜寻你时，你便认为唯有登船避难才安全。你的逃亡如此仓促，以至于你宁愿面对海上的风暴，也不愿承受留在岸上的后果。你不知怎地到了叙利亚，到达后表示想去耶路撒冷，在那里事奉主。谁会拒绝欢迎一个自称是隐修士的人呢？尤其是如果他不了解你那悲惨的经历，并且读过你的主教写给其他主教们的推荐信。不幸的人！你把自己装扮成光明的天使；\BibleRef{格后 11:14} 而实际上你是撒殚的仆役，却假装是正义的仆役。你只是披着羊皮的狼；\BibleRef{玛 7:15} 既然你对一个人的妻子犯了一次奸淫，如今你便想对基督的新娘犯奸淫。\par

12. 我叙述这些事的目的是为你勾勒出你邪恶生活的图像，并将你的恶行清楚地摆在你眼前。我希望阻止你以天主的仁慈和祂丰厚的温柔作为借口去犯新的罪，并阻止你再次将天主子钉在十字架上，明明地羞辱祂。因为如果你不读我引用的那段经文之后的话，你可能就会这样做。这些话是：\Quote{那吸收那屡次降在它上面的雨水，并生出合于耕种者使用的蔬菜的土地，从天主领受祝福；但那生出荆棘和蒺藜的，便被废弃，近于诅咒，它的结局就是焚烧。}\par

\SourceRecord{Translated by W.H. Fremantle, G. Lewis and W.G. Martley.；Nicene and Post-Nicene Fathers, Second Series；Vol. 6.；Edited by Philip Schaff and Henry Wace.；Buffalo, NY: Christian Literature Publishing Co.,；1893.；<http://www.newadvent.org/fathers/3001147.htm>.}

% structure: p0001 source=h1 target=section reason=page-title

\section{第一四八书函}

\Emph{致贵妇\Person{塞兰提亚}}\par

\EditorialNote{这是一封写给一位贵妇的有趣信函，讨论圣洁生活的原则和方法。然而，它并非热罗尼莫的作品，其风格几乎没有体现他的痕迹。这封信曾被归为诺拉的保利诺和苏尔皮基乌斯·塞维鲁斯所作。}\par

\SourceRecord{Translated by W.H. Fremantle, G. Lewis and W.G. Martley.；Nicene and Post-Nicene Fathers, Second Series；Vol. 6.；Edited by Philip Schaff and Henry Wace.；Buffalo, NY: Christian Literature Publishing Co.,；1893.；<http://www.newadvent.org/fathers/3001148.htm>.}

% structure: p0001 source=h1 target=section reason=page-title

\section{第一四九封书信}

\Emph{论犹太人的庆节}\par

\EditorialNote{这封信的主题是福音律法废止了犹太人的庆节。它不能被视为热罗尼莫的作品。}\par

\SourceRecord{Translated by W.H. Fremantle, G. Lewis and W.G. Martley.；Nicene and Post-Nicene Fathers, Second Series；Vol. 6.；Edited by Philip Schaff and Henry Wace.；Buffalo, NY: Christian Literature Publishing Co.,；1893.；<http://www.newadvent.org/fathers/3001149.htm>.}

% structure: p0001 source=h1 target=section reason=page-title

\section{\WorkTitle{第150封信}}

\Emph{自\Person{普罗科比乌斯}致\Person{热罗尼莫}}\par

此信亦见存于\Person{迦萨的普罗科比乌斯}文集中，原当归其名下。因该普罗科比乌斯活跃于热罗尼莫后一世纪，故不可能致函于他。\par

\SourceRecord{Translated by W.H. Fremantle, G. Lewis and W.G. Martley.；Nicene and Post-Nicene Fathers, Second Series；Vol. 6.；Edited by Philip Schaff and Henry Wace.；Buffalo, NY: Christian Literature Publishing Co.,；1893.；<http://www.newadvent.org/fathers/3001150.htm>.}
